\documentclass[11pt,openany]{book}
\setcounter{secnumdepth}{3}
\usepackage[margin=1.1in]{geometry}
\usepackage[T2A]{fontenc}
\usepackage[utf8]{inputenc}
\usepackage[english,ukrainian]{babel}
\usepackage{amsmath,amssymb,amsthm}
\usepackage{graphicx}
\usepackage{tikz}
\usetikzlibrary{arrows.meta,decorations.markings,decorations.pathmorphing,decorations.pathreplacing,calc}
\input{figures/icons}
\usepackage[unicode,colorlinks=true,linkcolor=blue!60!black,citecolor=blue!60!black]{hyperref}
\usepackage{microtype}
\setlength{\emergencystretch}{3em} % long names (licences, references) in narrow lines
\usepackage{empheq}
\usepackage{array}
\usepackage{booktabs}
\usepackage{multirow}
\usepackage{longtable}
\usepackage{circuitikz}
\definecolor{keyyellow}{rgb}{1.0,0.97,0.78}
% key equations: \begin{empheq}[box=\keybox]{equation} ... \end{empheq} -- light-yellow box, number stays outside
\newcommand{\keybox}[1]{{\setlength{\fboxsep}{7pt}\colorbox{keyyellow}{#1}}}
\usepackage[most]{tcolorbox}
% key statements (propositions etc.): the same light-yellow background, breakable across pages
\newtcolorbox{keyblock}{enhanced,breakable,colback=keyyellow,colframe=keyyellow,boxrule=0pt,arc=0pt,left=6pt,right=6pt,top=2pt,bottom=2pt,boxsep=0pt}
\renewcommand{\chaptermark}[1]{\markboth{\small\chaptername~\thechapter.\ #1}{}}
\renewcommand{\sectionmark}[1]{\markright{\small\thesection.\ #1}}
\renewcommand{\topfraction}{0.95}
\renewcommand{\textfraction}{0.05}
\renewcommand{\floatpagefraction}{0.75}

% part pages: title at the top third, and text may follow on the same page (used for the part introductions)
\makeatletter
\renewcommand\part{\if@openright\cleardoublepage\else\clearpage\fi\thispagestyle{plain}\if@twocolumn\onecolumn\@tempswatrue\else\@tempswafalse\fi\null\vspace{0.18\textheight}\secdef\@part\@spart}
\renewcommand\@endpart{\par\vspace{3em}\if@tempswa\twocolumn\fi}
\makeatother
\newtheorem{theorem}{Теорема}[chapter]
\newtheorem{proposition}{Твердження}[chapter]
\newtheorem{lemma}{Лема}[chapter]
\theoremstyle{remark}
\newtheorem{remark}{Зауваження}[chapter]
% exercises: \begin{exercise}[\lvlB] ... \end{exercise}, numbered "Exercise 4.3"; difficulty marks
\theoremstyle{definition}
\newtheorem{exercise}{Задача}[chapter]
\newcommand{\lvlA}{$\star$}
\newcommand{\lvlB}{$\star\star$}
\newcommand{\lvlC}{$\star\star\star$}
% answer to an exercise in the Answers appendix: \answer{ex:NN:k}{text}
\newcommand{\answer}[2]{\par\noindent\textbf{\ref{#1}.}~#2\par\smallskip}
\newcommand{\dd}{\mathrm{d}}
% neuron type code: first four letters of the primary mediator
\newcommand{\nt}[1]{\textsf{\textbf{#1}}}
\newcommand{\mV}{\,\text{мВ}}
\newcommand{\ms}{\,\text{мс}}
\newcommand{\mM}{\,\text{мМ}}
\newcommand{\um}{\,\text{мкм}}
\newcommand{\ion}[2]{\mathrm{#1}^{#2}}
\newcommand{\Na}{\ion{Na}{+}}
\newcommand{\K}{\ion{K}{+}}
\newcommand{\Cl}{\ion{Cl}{-}}
\newcommand{\Ca}{\ion{Ca}{2+}}
\newcommand{\conc}[1]{[\mathrm{#1}]}

\title{\Huge Комп'ютер на 20 ватах\\[16pt] \Large Як обчислює мозок: кількісний підхід для інженерів}
\hypersetup{pdftitle={Комп'ютер на 20 ватах. Як обчислює мозок: кількісний підхід для інженерів},pdfauthor={Костянтин Кладько}}
\author{\Large Костянтин~Кладько}
\date{}

\begin{document}
\frontmatter
\maketitle
\thispagestyle{empty}
\vspace*{\fill}
\noindent{\small \copyright~2026 Костянтин~Кладько.\\[4pt]
Текст, рисунки, задачі та відповіді поширюються за ліцензією Creative Commons «Зазначення авторства --- Некомерційна --- Поширення на тих самих умовах» 4.0 (CC BY-NC-SA 4.0): \url{https://creativecommons.org/licenses/by-nc-sa/4.0/deed.uk}. Книгу можна вільно копіювати, використовувати у викладанні, перекладати й доопрацьовувати з некомерційною метою за умови зазначення автора та поширення похідних творів на тих самих умовах.\\[4pt]
Програми моделей і збирання поширюються за ліцензією MIT.\\[4pt]
Вихідні тексти, моделі та нові версії: \url{https://github.com/kladkogex/20-watt-computer}.}
\clearpage

\tableofcontents
\chapter*{Як читати цю книгу}
\addcontentsline{toc}{chapter}{Як читати цю книгу}

Книгу не обов'язково читати підряд. Розділи~1--4 --- спільна основа: типи нейронів, що обчислюють \nt{GLUT} і \nt{GABA} та якою мовою вони говорять. Далі книга розгалужується. Карта на рис.~\ref{fig:roadmap} показує, які розділи на які спираються: стрілка від розділу $A$ до розділу $B$ означає, що в $B$ використано поняття й формули з $A$. На карті лише головні залежності; дрібні посилання назад читанню не заважають.

\begin{figure}[h]
\centering
\begin{tikzpicture}[>=Stealth,scale=0.95,
  ch/.style={draw,thick,rounded corners=3pt,align=center,font=\scriptsize,text width=1.85cm,minimum height=0.62cm,inner sep=2pt},
  base/.style={ch,fill=blue!8}, bus/.style={ch,fill=orange!14}, sum/.style={ch,fill=gray!18},
  io/.style={ch,fill=green!10}, sort/.style={ch,fill=cyan!8}, lab/.style={ch,fill=violet!10},
  dep/.style={->,thick,gray!60!black}]
% foundation
\node[base] (c1) at (0,0) {1 Типи нейронів};
\node[base] (c2) at (2.3,0) {2 \nt{GLUT}};
\node[base] (c3) at (4.6,0) {3 \nt{GLUT} і \nt{GABA}};
\node[base] (c4) at (6.9,0) {4 Абетка Морзе};
% broadcast channels and the synthesis
\node[bus] (c5) at (4.6,-1.5) {5 \nt{SERO}};
\node[bus] (c6) at (7.3,-1.5) {6 \nt{DOPA}};
\node[bus] (c7) at (9.2,-0.9) {7 Теорії дофаміну};
\node[bus] (c8) at (9.2,-2.1) {8 \nt{NORA}};
\node[bus] (c9) at (11.5,-2.1) {9 \nt{ACET}};
\node[sum] (c10) at (13.8,-0.9) {10 Уся машина};
% inputs and outputs
\node[io] (c12) at (2.3,-4.1) {12 Розпі-\\знавання};
\node[io] (c11) at (4.6,-4.1) {11 Входи};
\node[io] (c13) at (6.9,-4.1) {13 М'язи};
\node[io] (c14) at (9.2,-3.05) {14 Біль};
\node[io] (c15) at (9.2,-4.1) {15 Навчання\\рухів};
\node[io] (c16) at (9.2,-5.15) {16 Хімія\\тіла};
% lab, sorting, sleep
\node[lab] (c17) at (11.5,-4.1) {17 Налагодження};
\node[lab] (c21) at (13.8,-4.1) {21 Живі машини};
\node[lab] (c22) at (13.8,-5.15) {22 DishBrain};
\node[sort] (c18) at (11.5,-5.15) {18 Нейрони-\\прилади};
\node[sort] (c19) at (13.8,-6.2) {19 Кількість\\дендритів};
\node[bus] (c20) at (11.5,-6.2) {20 Сон};
% dependencies
\draw[dep] (c1) -- (c2); \draw[dep] (c2) -- (c3); \draw[dep] (c3) -- (c4);
\draw[dep] (c1) -- (c5); \draw[dep] (c4) -- (c6); \draw[dep] (c5) -- (c6);
\draw[dep] (c6) -- (c7); \draw[dep] (c6) -- (c8); \draw[dep] (c8) -- (c9);
\draw[dep] (c7) -- (c10); \draw[dep] (c9) -- (c10);
\draw[dep] (c4) -- (c11); \draw[dep] (c11) -- (c12); \draw[dep] (c11) -- (c13); \draw[dep] (c5) -- (c13);
\draw[dep] (c13) -- (c14); \draw[dep] (c13) -- (c15); \draw[dep] (c13) -- (c16);
\draw[dep] (c15) -- (c17); \draw[dep] (c9) -- (c17); \draw[dep] (c17) -- (c21); \draw[dep] (c21) -- (c22);
\draw[dep] (c16) -- (c18); \draw[dep] (c18) -- (c19); \draw[dep] (c16) -- (c20);
% legend
\begin{scope}[yshift=-7.25cm]
\foreach \s/\t [count=\i] in {base/основа, bus/канали й режими, sum/підсумок, io/входи й виходи, sort/сортування нейронів, lab/лабораторія}{
  \pgfmathtruncatemacro{\col}{mod(\i-1,3)}\pgfmathtruncatemacro{\row}{(\i-1)/3}
  \node[\s,text width=0.3cm,minimum height=0.32cm] at ({\col*4.8+0.2},{-\row*0.55}) {};
  \node[anchor=west,font=\scriptsize] at ({\col*4.8+0.55},{-\row*0.55}) {\t};}
\end{scope}
\end{tikzpicture}
\caption{Карта розділів. Стрілка веде від розділу до того, що на нього спирається. Інші посилання між розділами --- вказівники, а не залежності.}
\label{fig:roadmap}
\end{figure}

\section*{Шляхи крізь книгу}

\begin{center}
\footnotesize
\setlength{\tabcolsep}{4pt}
\renewcommand{\arraystretch}{1.2}
\begin{tabular}{>{\raggedright}p{3.0cm}>{\raggedright}p{3.9cm}>{\raggedright\arraybackslash}p{6.4cm}}
\toprule
Шлях & Для кого & Розділи за порядком \\
\midrule
курс на одну чверть & викладач, 10 тижнів & тиждень 1: розд.~1--2; 2: розд.~3; 3: розд.~4; 4: розд.~5--6; 5: розд.~7--8; 6: розд.~9; 7: розд.~10; 8: розд.~11--12; 9: розд.~13 і~15; 10: розд.~17 \\
ШІ та машинне навчання & той, хто знає нейромережі й хоче зрозуміти, як навчається мозок & 1--4, 5 (побіжно), 6, 7, 8, 9, 11 (побіжно), 12, 13 (побіжно), 15 \\
електроніка й залізо & схемотехнік, розробник нейроморфних чипів & 1--4, 5, 11, 13, 16, 18, 19, 17 (розділи 9 і 15 побіжно) \\
нейроінтерфейси й живі обчислювачі & той, хто будує інтерфейси мозок---комп'ютер і чипи з живими нейронами & 1, 2, 4, 11, 13, 17 (розділи 9 і 15 побіжно), 21, 22 (розділ 12 побіжно) \\
швидкий огляд & інженер, який має вихідні & 1, 2, 3, 6, 10; посилання розділу~6 на розділи~4--5 зрозумілі з контексту \\
\bottomrule
\end{tabular}
\end{center}

«Побіжно» означає: досить прочитати початок розділу, його твердження в жовтих рамках і підсумкову таблицю.

Наприкінці кожного розділу є задачі: оцінки, виведення, моделювання й проєктування; короткі відповіді зібрано в додатку~\ref{sec:answers}. Усі позначення --- у додатку~\ref{sec:notation}, а досліди, на яких тримаються головні твердження кожного розділу, --- у додатку~\ref{sec:supplement}. Числа в книзі отримано з невеликих програм у теці \texttt{models/}; їх можна запустити й змінити.

\mainmatter

\chapter{Типи нейронів}
\label{sec:neurontypes}

\section{Нейрон як чорна скринька}

Припустімо, вам дали мішок із вісімдесятьма шістьма мільярдами нейронів~\cite{Azevedo2009} і попросили розкласти їх на купки. Як би ви це зробили?

Інженер, якому дали мішок незнайомих мікросхем, не став би сортувати їх за формою корпусу. Він запитав би: скільки в деталі входів, скільки виходів і що вона видає на вихід. Намалюймо нейрон так само --- як блок на схемі (рис.~\ref{fig:blackbox}).

\begin{figure}[t]
\centering
\begin{tikzpicture}[>=Stealth,x=1cm,y=1cm]
% the box
\node[draw,very thick,minimum width=2.6cm,minimum height=2.6cm,fill=blue!6] (N) at (0,0) {};
\node at (0,0.55) {\small нейрон};
\node at (0,-0.15) {$\displaystyle u=\sum_i w_i x_i$};
\node at (0,-0.8) {\small $y=\Theta(u-\theta)$};
% inputs
\foreach \k/\y/\lab in {1/1.05/{x_1},2/0.55/{x_2},3/0.05/{x_3}}{
  \draw[->,thick,blue!60!black] (-4.2,\y) -- (-1.3,\y);
  \node[left] at (-4.2,\y) {$\lab$};
  \node[above,blue!60!black] at (-2.1,\y-0.06) {\scriptsize $w_{\k}$};}
\node at (-4.45,-0.4) {$\vdots$};
\node at (-2.1,-0.4) {$\vdots$};
\draw[->,thick,blue!60!black] (-4.2,-1.05) -- (-1.3,-1.05);
\node[left] at (-4.2,-1.05) {$x_n$};
\node[above,blue!60!black] at (-2.1,-1.11) {\scriptsize $w_{n}$};
\draw[decorate,decoration={brace,amplitude=5pt},gray!60!black] (-4.9,-1.2) -- (-4.9,1.2);
\node[left,align=right,gray!40!black] at (-5.1,0) {\small $n\sim10^3$--$10^5$\\[-2pt]\small входів};
% single output wire, then fan-out
\draw[very thick,red!70!black] (1.3,0) -- (3.2,0);
\foreach \x in {1.6,2.2,2.34,2.9}{\draw[thick,red!70!black] (\x,0) -- (\x,0.4);}
\node[above right,font=\tiny\itshape,gray!30!black,inner sep=1pt] at (1.42,0.45) {клац\dots\ клац-клац\dots};
\node[below,red!70!black,align=center] at (2.25,-0.05) {\scriptsize один вихід $y$};
\fill[red!70!black] (3.2,0) circle (0.06);
\foreach \y in {1.3,0.65,0,-0.65,-1.3}{
  \draw[->,thick,red!70!black] (3.2,0) -- (5.0,\y);}
\draw[decorate,decoration={brace,amplitude=5pt},gray!60!black] (5.3,1.45) -- (5.3,-1.45);
\node[right,align=left,gray!40!black] at (5.5,0) {\small копії $y$\\[-2pt]\small на $10^3$--$10^4$\\[-2pt]\small інших нейронів};
\end{tikzpicture}
\caption{Нейрон очима інженера. Багато входів $x_i$, кожен зі своєю вагою $w_i$; усередині --- зважена сума $u$ і порівняння з порогом $\theta$ ($\Theta$ --- сходинка: $1$, якщо аргумент додатний, і $0$ інакше); один вихід $y$ --- послідовність імпульсів, розмножена розгалуженням на тисячі отримувачів.}
\label{fig:blackbox}
\end{figure}

На виході блока --- не число, а послідовність коротких імпульсів напруги: «клац», пауза, «клац-клац», довга пауза. Усі імпульси однакової висоти, тож уся інформація --- в тому, \emph{коли} вони відбуваються. Усередині блока, в першому грубому наближенні, стоять суматор і компаратор: входи додаються з вагами, і якщо сума перевищила поріг, блок видає імпульс. У наступному розділі ми замінимо його моделлю, яка чесно працює в неперервному часі.

Вихід один, але провід розгалужується, і кожна гілка несе \emph{ту саму} копію сигналу. В електроніці це називають навантажувальною здатністю за виходом, fan-out. У нейрона кори вона становить кілька тисяч; логічний вентиль у мікросхемі зазвичай навантажений на одиниці інших вентилів.

\emph{Що} ж передається вихідним проводом до отримувача? Імпульс добігає до кінця кожної гілки й там перетворюється на хімічний сигнал: гілка викидає у вузьку щілину між клітинами дрібку певної речовини, \emph{нейромедіатора}, яка змінює вхідний струм отримувача. Ось і ознака для сортування: \emph{який медіатор викидає вихід нейрона}.

\section{Один вихід --- один тип сигналу}

Сортування має сенс, лише якщо в кожного нейрона тип виходу один. Чи так це? Дослід каже, що майже так.

\begin{keyblock}
\begin{proposition}[Один нейрон --- один тип виходу]
\label{prop:dale}
Майже кожен нейрон має один головний медіатор і викидає його в усіх гілках свого виходу. Тому тип нейрона задається його головним медіатором~\cite{Strata1999}.
\end{proposition}
\end{keyblock}

Домовмося позначати тип нейрона \emph{першими чотирма літерами англійської назви його головного медіатора}. Так, нейрон, головний медіатор якого глутамат, --- це \nt{GLUT}-нейрон. Усі типи зібрано в таблиці~\ref{tab:types}.

\begin{table}[t]
\centering
\footnotesize
\setlength{\tabcolsep}{3pt}
\renewcommand{\arraystretch}{1.15}
\begin{tabular}{>{\raggedright}p{1.0cm}>{\raggedright}p{2.05cm}>{\raggedright}p{1.75cm}>{\raggedright}p{1.6cm}>{\centering}p{0.7cm}>{\raggedright}p{1.55cm}>{\raggedright}p{1.8cm}>{\raggedright\arraybackslash}p{3.2cm}}
\toprule
Тип & Головний медіатор & Шина & Швидкість & Знак & Нейронів у мозку & Виходів у нейрона & Роль в обчисленні \\
\midrule
\nt{GLUT} & глутамат & адресна & швидко й повільно & $+$ & $\sim8\cdot10^{10}$ & $\sim10^4$ & головна додатна вага \\
\nt{GABA} & ГАМК & адресна & швидко й повільно & $-$ & $\sim3\cdot10^{9}$ & $10^3$--$10^4$ & головна від'ємна вага \\
\nt{GLYC} & гліцин & адресна & швидко & $-$ & $10^6$--$10^7$ & $10^2$--$10^3$ & від'ємна вага в спинному мозку та стовбурі; другий ключ для одного з рецепторів глутамату \\
\nt{ACET} & ацетилхолін & обидві & швидко й повільно & $\pm$ & $\sim10^6$ & м'яз: $10$--$10^3$; мозок: $\sim10^5$ & вихід на м'яз; у мозку --- швидкість навчання (гіпотеза) \\
\nt{DOPA} & дофамін & широко\-мовна & повільно & $\pm$ & $\sim5\cdot10^{5}$ & $10^5$--$10^6$ & помилка передбачення винагороди $\delta$ \\
\nt{SERO} & серотонін & широко\-мовна & повільно (один тип рецептора швидкий) & $\pm$ & $\sim3\cdot10^{5}$ & $\sim10^5$ & горизонт планування (гіпотеза) \\
\nt{HIST} & гістамін & широко\-мовна & повільно & $+$ & $\sim6\cdot10^{4}$ & немає оцінки & глобальний сигнал «не спи» \\
\nt{NORA} & норадреналін & широко\-мовна & повільно & $\pm$ & $\sim5\cdot10^{4}$ & $\sim10^5$ & коефіцієнт підсилення (гіпотеза) \\
\nt{ADRE} & адреналін & широко\-мовна & повільно & $\pm$ & $\sim10^4$ & немає оцінки & мало нейронів; роль у мозку неясна \\
\nt{ASPA} & аспартат & адресна & швидко & $+$ & немає оцінки & немає оцінки & слабка додатна вага; роль спірна \\
\bottomrule
\end{tabular}
\caption{Типи нейронів, за спаданням кількості нейронів у мозку. \emph{Шина}: адресна --- медіатор викидається у вузьку щілину до певного отримувача; широкомовна --- від малої групи нейронів-джерел на великі області (підрозділ~\ref{sec:buses}). \emph{Швидкість}: швидко --- через іонотропні рецептори, мілісекунди; повільно --- через метаботропні, сотні мілісекунд і довше. \emph{Знак} --- звичайний; остаточно його задає отримувач (підрозділ~\ref{sec:receiver}), і $\pm$ означає, що трапляються обидва знаки. \emph{Нейронів у мозку} --- скільки нейронів цього типу в усьому мозку людини, за порядком величини; загалом нейронів близько $8.6\cdot10^{10}$~\cite{Azevedo2009}. Майже всі \nt{GLUT}-нейрони --- близько $7\cdot10^{10}$ --- це дрібні вхідні нейрони мозочка; числа для \nt{GLUT} і \nt{GABA} отримано з цього підрахунку та частки \nt{GABA}-нейронів у корі~\cite{Hendry1987}. З нечисленних типів безпосередньо підраховано лише \nt{HIST}~\cite{Airaksinen1991} і головну з двох груп \nt{DOPA}-нейронів~\cite{Pakkenberg1991}, тож для \nt{DOPA} це нижня межа; для \nt{SERO} --- сума двох часткових підрахунків~\cite{Baker1990,Baker1991}. Числа для \nt{NORA}, \nt{GLYC}, \nt{ACET} і \nt{ADRE} --- грубі оцінки за порядком величини, без прямого підрахунку. \emph{Виходів у нейрона} --- скільки в одного нейрона синаптичних закінчень, тобто місць викиду медіатора на гілках вихідного проводу (рис.~\ref{fig:blackbox}), за порядком величини. Для широкомовних типів закінчення безпосередньо не рахували: для \nt{DOPA} число виведено з того, яку частку своєї мішені покриває розгалуження одного вихідного проводу~\cite{Matsuda2009}, для решти оцінки ще грубіші. Для \nt{ACET} два випадки: нейрон, що керує м'язом, і широкомовний нейрон у мозку.}
\label{tab:types}
\end{table}

% the pie is a table-type float with a figure caption, so it can never float ahead of Table~\ref{tab:types}
\begin{table}[tb]
\makeatletter\def\@captype{figure}\makeatother
\centering
\definecolor{vizglut}{HTML}{2A78D6}
\definecolor{vizgaba}{HTML}{E34948}
\definecolor{vizrest}{HTML}{8A8986}
\begin{tikzpicture}[>=Stealth]
% pie: GLUT 96.4 %, GABA 3.6 % (13.0 degrees), the rest 0.006 % (a hairline at 0 degrees)
\fill[vizglut] (0,0) -- (13:1.9) arc (13:360:1.9) -- cycle;
\fill[vizgaba] (0,0) -- (0:1.9) arc (0:13:1.9) -- cycle;
\draw[white,line width=1pt] (0,0) -- (13:1.9);
\draw[vizrest,line width=1.2pt] (0,0) -- (0:1.9);
\node[white,align=center,font=\small] at (190:0.95) {\nt{GLUT}\\$96.4\,\%$};
\draw[gray!60!black] (6.5:1.75) -- (2.05,2.1);
\node[anchor=west,font=\small] at (2.05,2.1) {\nt{GABA} $3.6\,\%$};
% zoom box for the remaining types
\draw[densely dashed,vizrest] (1.9,0) -- (3.1,1.65);
\draw[densely dashed,vizrest] (1.9,0) -- (3.1,-2.2);
\draw[vizrest,rounded corners=3pt] (3.1,-2.2) rectangle (10.1,1.65);
\node[anchor=west,font=\scriptsize] at (3.2,1.4) {усі інші разом: $\approx0.006\,\%$; логарифмічна шкала};
\foreach \code/\lg/\y/\val/\lx in {GLYC/6.477/1.0/{$10^6$--$10^7$}/4.05,ACET/6.0/0.6/{$\sim10^6$}/3.0,DOPA/5.699/0.2/{$\sim5\cdot10^5$}/2.699,SERO/5.477/-0.2/{$\sim3\cdot10^5$}/2.477,HIST/4.778/-0.6/{$\sim6\cdot10^4$}/1.778,NORA/4.699/-1.0/{$\sim5\cdot10^4$}/1.699,ADRE/4.0/-1.4/{$\sim10^4$}/1.0}{
  \node[anchor=east,font=\scriptsize] at (4.35,\y) {\nt{\code}};
  \fill[vizrest] (4.45,\y-0.13) rectangle ({4.45+\lg-3},\y+0.13);
  \node[anchor=west,font=\scriptsize] at ({4.5+\lx},\y) {\val};}
\draw[vizrest,thick] (7.45,1.0) -- (8.45,1.0);
\draw[vizrest,thick] (7.45,0.92) -- (7.45,1.08);
\draw[vizrest,thick] (8.45,0.92) -- (8.45,1.08);
\draw[gray!60!black] (4.45,-1.72) -- (8.45,-1.72);
\foreach \e in {3,4,5,6,7}{
  \draw[gray!60!black] ({4.45+\e-3},-1.72) -- ({4.45+\e-3},-1.66);
  \node[below,font=\scriptsize] at ({4.45+\e-3},-1.72) {$10^{\e}$};}
\end{tikzpicture}
\caption{Частки типів нейронів у мозку за оцінками таблиці~\ref{tab:types}. Коло майже цілком зайняте \nt{GLUT}; \nt{GABA} --- вузький сектор, а всі інші типи разом --- волосина на межі. Праворуч цю волосину збільшено, за логарифмічною шкалою кількості нейронів; для \nt{GLYC} стовпчик сягає середини діапазону $10^6$--$10^7$, відрізок показує весь діапазон. Для \nt{ASPA} оцінки немає, і на рисунку його немає.}
\label{fig:typepie}
\end{table}

Наскільки нерівні ці купки, видно на рис.~\ref{fig:typepie}: з кожної тисячі нейронів $964$ --- \nt{GLUT}, $36$ --- \nt{GABA}, а на всі інші типи разом припадає близько шести нейронів на сто тисяч.

Назва GABA (англійське скорочення ГАМК) і так складається з чотирьох літер. Речовини, які майже ніколи не бувають головним медіатором, --- нейропептиди, гази, ліпіди, пурини, --- власного коду не отримують; про них --- у додатку~\ref{sec:othermediators}. Слово «майже» у твердженні~\ref{prop:dale} важливе. Багато нейронів викидають разом із головним медіатором допоміжні речовини~\cite{Yao2023}. Деякі викидають і другий швидкий медіатор, навіть протилежного знака: частина \nt{DOPA}-нейронів викидає ще й ГАМК~\cite{Tritsch2012}. А в деяких нейронів різні медіатори виходять із різних гілок одного й того самого виходу~\cite{Zhang2015}. Усе це виміряно, тож «один нейрон --- один медіатор» --- правило з винятками, а не закон. Але як перший поділ воно витримує перевірку: у клітинному атласі мозку миші, де нейрони розкладено за генами, що в них працюють, на понад п'ять тисяч груп, нейрони так само діляться на великі класи насамперед за головним медіатором~\cite{Yao2023}.

Із цього правила випливає наслідок, який одразу відрізняє мозок від штучної нейронної мережі. У штучній мережі той самий нейрон може мати додатну вагу до одного отримувача й від'ємну до іншого: ваги $w_{ij}$ --- просто числа в матриці. У мозку знак дії здебільшого визначається медіатором, а медіатор у нейрона один. Отже, якщо нейрон $j$ збуджує одного отримувача, він збуджує й усіх інших. Увесь \emph{стовпець} матриці ваг, що виходить із нейрона $j$, одного знака:
\begin{empheq}[box=\keybox]{equation}
\operatorname{sign} w_{ij} \;=\; s_j \quad\text{для всіх отримувачів } i, \qquad s_j\in\{+1,-1\}.
\label{eq:dalesign}
\end{empheq}
Нейрони поділяються на збуджувальні ($s_j=+1$) і гальмівні ($s_j=-1$), і це властивість самого нейрона, а не зв'язку. Якщо мережі потрібен від'ємний зв'язок від збуджувального нейрона, його доводиться робити через посередника: збуджувальний нейрон збуджує гальмівний, а той гальмує ціль: від'ємна вага із затримкою на одну ланку.

Медіаторів відомо понад сотню, але майже вся робота мозку тримається на півдюжині, і вони поділяються на два зовсім різні класи.

\section{Дві шини: адресна й широкомовна}
\label{sec:buses}

У комп'ютерних мережах є два способи доставити повідомлення. Перший --- \emph{адресний}, unicast: пакет іде від одного відправника до певного отримувача, швидко, і більше його ніхто не бачить. Другий --- \emph{широкомовний}, broadcast: відправник передає одне повідомлення відразу всім вузлам мережі. Нейрони користуються обома способами, і медіатори поділяються саме за цією ознакою (рис.~\ref{fig:phone_radio}).

\begin{figure}[t]
\centering
\begin{tikzpicture}[>=Stealth,scale=0.95]
% ---- unicast: point-to-point wiring
\node[above] at (2.0,3.3) {\small \textbf{Адресна}: глутамат і ГАМК};
\foreach \i in {0,...,4}{
  \fill[blue!60!black] (0.2+0.9*\i,2.5) circle (0.1);
  \fill[blue!60!black] (0.2+0.9*\i,0.6) circle (0.1);}
\foreach \a/\b in {0/0,0/1,1/1,1/3,2/2,3/2,3/4,4/4,2/0}{
  \draw[->,blue!60!black] (0.2+0.9*\a,2.38) -- (0.2+0.9*\b,0.72);}
\fill[red!70!black] (1.1,1.55) circle (0.09);
\pic[scale=1.2] at (1.75,1.9) {letter};
\draw[->,red!70!black,thick] (1.1,1.55) -- (0.28,0.7);
\draw[->,red!70!black,thick] (1.1,1.55) -- (1.95,0.72);
\node[align=center,below] at (2.0,0.2) {\small мільярди нейронів,\\[-2pt]\small у кожного свої адресати,\\[-2pt]\small час: мілісекунди};
% ---- broadcast: one small source, global targets
\begin{scope}[xshift=7.2cm]
\node[above] at (1.8,3.3) {\small \textbf{Широкомовна}: дофамін, \dots};
\foreach \i in {0,...,8}{\fill[blue!60!black] (-0.2+0.5*\i,2.75) circle (0.08);}
\fill[orange!85!black] (1.8,0.35) ellipse (0.35 and 0.18);
\pic[scale=1.5] at (2.7,0.75) {mast};
\foreach \x in {-0.2,0.3,0.8,1.3,1.8,2.3,2.8,3.3,3.8}{
  \draw[->,orange!85!black] (1.8,0.5) .. controls (1.8,1.3) and (\x,1.6) .. (\x,2.62);}
\node[align=center,below] at (1.8,0.1) {\small тисячі --- сотні тисяч нейронів,\\[-2pt]\small одне повідомлення всім,\\[-2pt]\small час: сотні мілісекунд і довше};
\end{scope}
\end{tikzpicture}
\caption{Два способи передачі. Ліворуч --- адресна шина, схожа на пошту з адресою на конверті; червоним виділено один нейрон і двох його отримувачів. Праворуч --- широкомовна, як радіостанція: невелика група нейронів-джерел, і те саме повідомлення отримують усі.}
\label{fig:phone_radio}
\end{figure}

\emph{Адресна шина} --- це глутамат і ГАМК. Їх викидає переважна більшість нейронів. Кожен вихід веде до певних клітин, медіатор діє лише у вузькій щілині контакту, і діє швидко: за мілісекунду відкриваються іонні канали отримувача, за кілька мілісекунд усе закінчується. Це шина \emph{даних}: нею передається те, що мозок обчислює.

\emph{Широкомовна шина} --- дофамін, серотонін, норадреналін і почасти ацетилхолін. Їх викидають крихітні групи нейронів, але вихід кожного з них розгалужується неймовірно широко. Медіатор часто викидається не у вузьку щілину, а в міжклітинний простір, де він розпливається й діє на всіх, хто поруч. Діє повільно: не відкриває канали безпосередньо, а запускає всередині отримувача ланцюжок хімічних реакцій. Нею йдуть не дані, а \emph{керувальні параметри}, спільні для великих ділянок мережі, які змінюють режим її роботи: коефіцієнт підсилення, швидкість навчання, сигнал «це було добре». Тому такі медіатори називають \emph{модуляторами}. Довго вважали, що кожен такий канал --- одне число на весь мозок. Сучасні вимірювання уточнили картину: канал --- не один провід, а невеликий джгут паралельних ліній. Нейрони-джерела одного медіатора поділяються на підсистеми, що мають своїх отримувачів і своє повідомлення, а скільки медіатора викинути на місці, підлаштовують ще й місцеві сигнали~\cite{Ren2018,Poe2020}. Таких ліній --- одиниці чи десятки, а не мільярди, тож канал лишається широкомовним.

Якщо вам потрібна одна картинка з цього розділу --- ось вона. Мозок --- величезна мережа з адресною передачею даних, над якою висить кілька широкомовних каналів із керувальними сигналами. Програміст упізнає знайомий устрій: обчислення плюс невеликий набір керувальних змінних, кожна з яких спільна для великої ділянки мережі.

\section{\nt{GLUT}: додатна вага}

Глутамат --- головний збуджувальний медіатор мозку. Коли вихід нейрона викидає глутамат, в отримувача відкриваються канали, через які всередину входить натрій, напруга на мембрані отримувача зростає, і його шанс видати власний імпульс збільшується. Мовою рис.~\ref{fig:blackbox} це вхід із \emph{додатною} вагою, $w_i>0$.

Хто видає глутамат? Майже всі, хто передає дані на великі відстані: від трьох чвертей до чотирьох п'ятих нейронів кори й більшість нейронів, що з'єднують різні області мозку. Мережа мозку --- це насамперед мережа глутаматних зв'язків.

Одна інженерна деталь. Після викиду концентрація глутамату в щілині злітає в тисячі разів, приблизно до мілімоля, і спадає зі сталою часу близько мілісекунди~\cite{Clements1992}: глутамат розпливається з щілини, а спеціальні білки-насоси відкачують його назад у клітини. Навіщо? Із тієї самої причини, з якої в лінії зв'язку після кожного символу сигнал повертають до нуля. Якщо медіатор не прибрати, наступний імпульс прийде на «зайняту» лінію, й імпульси з інтервалом у кілька мілісекунд зіллються.

\section{\nt{GABA}: від'ємна вага}

Уявіть мережу, у якій усі ваги додатні. Якщо в середньому кожен імпульс породжує більше ніж один наступний імпульс, активність зростає експоненційно й за кілька кроків охоплює всю мережу. Електронник упізнає петлю додатного зворотного зв'язку з підсиленням більше одиниці: схема йде в насичення. Щоб цього не сталося, потрібні від'ємні ваги. Їхнє головне джерело --- гамма-аміномасляна кислота, скорочено ГАМК.

ГАМК відкриває в отримувача канали, що пропускають іони хлору. Рівноважний потенціал хлору лежить біля потенціалу спокою або трохи нижче, тому відкриті хлорні канали утримують мембрану біля спокою й не дають підняти напругу до порога. Це вхід із \emph{від'ємною} вагою, $w_i<0$. ГАМКергічні --- від п'ятої частини до чверті нейронів кори~\cite{Hendry1987}. Їхні виходи короткі, і гальмують вони своїх найближчих сусідів.

Гальмування в мозку робить набагато більше, ніж просто віднімання. Воно працює як від'ємний зворотний зв'язок, що тримає всю мережу в робочій точці. Вважають, що в нормально працюючій корі збуджувальний і гальмівний струми, які надходять на нейрон, майже точно врівноважують один одного: такий режим передбачено теорією~\cite{vanVreeswijk1996} і видно у вимірюваннях струмів у живій корі~\cite{Haider2006}, і нейрон весь час сидить біля порога. Схема, стабілізована сильним від'ємним зворотним зв'язком біля нестійкої точки, швидко й чутливо реагує на малі сигнали --- давній інженерний прийом.

\section{\nt{DOPA}: сигнал помилки}

Перший широкомовний канал --- дофамін. Дофамінових нейронів у людини близько пів мільйона, приблизно один на сто сімдесят тисяч; стільки підраховано в головній із двох їхніх груп~\cite{Pakkenberg1991}, тож це нижня межа. Їхні виходи розходяться до тих областей, які обирають дії.

Що передає цей канал? Відповідь дають досліди, у яких реєструють імпульси дофамінових нейронів у тварини, що отримує винагороду~\cite{Schultz1997}. Тварині час від часу несподівано дають краплю соку, і дофамінові нейрони відповідають на неї короткою пачкою імпульсів. Потім перед соком починають запалювати лампочку, завжди за секунду до соку. Коли тварина вивчає цей зв'язок, нейрони починають відповідати пачкою на \emph{лампочку}, а на сам сік, що прийшов точно вчасно, не відповідають зовсім. А якщо після лампочки соку не дають, то в момент, коли він мав прийти, нейрони \emph{замовкають}, і їхня частота падає нижче звичайної.

Правило, що пояснює всі три випадки (рис.~\ref{fig:rpe}): нейрони повідомляють не про те, наскільки добре, а про те, наскільки \emph{краще чи гірше, ніж очікувалося}.

\begin{figure}[t]
\centering
\begin{tikzpicture}[>=Stealth,x=3.4cm,y=1cm]
\foreach \y/\lab in {2.8/{несподіваний сік},1.4/{сік після лампочки},0/{лампочка, а соку немає}}{
  \draw[gray!70!black] (0,\y) -- (3,\y);
  \draw[densely dotted,gray!60!black] (0,\y+0.3) -- (3,\y+0.3);
  \node[left,align=right,font=\small] at (-0.03,\y+0.35) {\lab};}
% row 1: juice without a cue -> burst at the juice
\draw[very thick,orange!85!black] plot[domain=0:3,samples=120] (\x,{2.8+0.3+0.75*exp(-((\x-2.08)/0.07)^2)});
\pic[scale=0.8] at (2,2.58) {juice};
% row 2: cue then juice -> burst at the cue, nothing at the juice
\draw[very thick,orange!85!black] plot[domain=0:3,samples=120] (\x,{1.4+0.3+0.75*exp(-((\x-1.08)/0.07)^2)});
\pic[scale=0.85] at (1,1.15) {lamp}; \pic[scale=0.85] at (2,1.15) {juice};
% row 3: cue, juice omitted -> burst at the cue, dip when the juice was due
\draw[very thick,orange!85!black] plot[domain=0:3,samples=120] (\x,{0.3+0.75*exp(-((\x-1.08)/0.07)^2)-0.28*exp(-((\x-2.1)/0.12)^2)});
\pic[scale=0.85] at (1,-0.25) {lamp}; \pic[scale=0.85] at (2,-0.25) {nojuice};
\node[right,font=\scriptsize] at (2.36,0.1) {провал};
\node[right,font=\scriptsize] at (2.13,3.75) {сюрприз!};
\node[left,font=\scriptsize] at (0.97,2.25) {сплеск на лампочку};
\node[above,font=\scriptsize] at (2,1.75) {нічого};
\draw[->,gray!70!black] (0,-0.75) -- (3.05,-0.75) node[right,black,font=\small] {$t$};
\draw[gray!60!black] (1,-0.8) -- (1,-0.7) node[below=2pt,font=\scriptsize] {лампочка, $1\,\text{с}$};
\draw[gray!60!black] (2,-0.8) -- (2,-0.7) node[below=2pt,font=\scriptsize] {сік, $2\,\text{с}$};
\end{tikzpicture}
\caption{Частота імпульсів \nt{DOPA}-нейронів у трьох дослідах (схематично, за~\cite{Schultz1997}). Пунктир --- рівна фонова частота. Лампочка --- сигнал, крапля --- сік, перекреслена крапля --- сік, який обіцяли, але не дали. Відгук --- помилка передбачення~\eqref{eq:rpe}.}
\label{fig:rpe}
\end{figure}

Програміст, знайомий із навчанням із підкріпленням~\cite{SuttonBarto2018}, упізнає цю величину відразу. Агент, який вчиться обирати дії, зберігає оцінку $V(s)$ --- скільки винагороди він очікує, перебуваючи в стані $s$. Після кожного кроку він порівнює очікування з тим, що отримав:
\begin{empheq}[box=\keybox]{equation}
\delta_t \;=\; r_t \;+\; \gamma\,V(s_{t+1}) \;-\; V(s_t) ,
\label{eq:rpe}
\end{empheq}
і виправляє оцінку: $V(s_t)\leftarrow V(s_t)+\alpha\,\delta_t$. Тут $r_t$ --- отримана винагорода, $\gamma<1$ --- коефіцієнт дисконтування (наскільки майбутня винагорода дешевша за негайну), $\alpha$ --- швидкість навчання. Величину $\delta_t$ називають \emph{помилкою часових різниць}.

Вона поводиться точнісінько як дофамінові нейрони. Несподіваний сік: $V$ ще нуль, $r>0$, $\delta>0$. Вивчений сік: винагороду передбачила лампочка, $V(s_t)$ перед соком уже дорівнює $r$, і $\delta=0$. Сік не прийшов: $V(s_t)=r$, а $r_t=0$, $\delta<0$. А сплеск переїжджає на лампочку, бо саме в момент лампочки очікування стрибком зросло: $\gamma V(s_{t+1})-V(s_t)>0$. Кроки $t,t+1,\dots$ тут --- умовність алгоритму: в організмі немає годинника, що відлічує кроки, і ту саму величину в неперервному часі ми отримаємо в розділі~\ref{sec:dofa}.

Отже, дофамін --- це широкомовна розсилка скалярного сигналу помилки $\delta$ усім, хто має за ним навчатися. У першому наближенні --- один провід на весь мозок і одне число в кожен момент часу; наскільки цей провід насправді джгут, розглянуто в розділі~\ref{sec:dofaalt}.

\section{Інші широкомовні канали}

Для решти модуляторів є лише робочі гіпотези, які перекладають їхні ролі на ту саму мову глобальних параметрів~\cite{Doya2002}. Сприймайте їх саме як гіпотези.

\emph{\nt{NORA}, норадреналін: коефіцієнт підсилення.} Норадреналінових нейронів ще менше, ніж дофамінових, за грубою оцінкою кілька десятків тисяч, а їхні виходи розходяться практично по всьому мозку. Вони різко активуються, коли стається щось несподіване або важливе. Норадреналін змінює крутизну передавальної характеристики нейронів-отримувачів: слабкі входи стають слабшими, сильні --- сильнішими. Вимірюють це так: фонові імпульси отримувача пригнічуються сильніше, ніж відповідь на вхід, і відношення сигналу до шуму зростає~\cite{Foote1975NE}. У простій моделі це описують одним числом: якщо нейрон відповідає частотою $f=F\bigl(g\cdot u\bigr)$ на вхідний струм $u$, то норадреналін збільшує коефіцієнт $g$~\cite{AstonJones2005} (докладніше --- у розділі~\ref{sec:nora}). Мережа з великим $g$ працює «контрастніше» й швидше ухвалює рішення; з малим --- «розмито» й більше перебирає варіанти, як агент навчання з підкріпленням у режимі пошуку.

\emph{\nt{ACET}, ацетилхолін: швидкість навчання.} Ацетилхолін керує тим, наскільки мережа слухає поточний вхід, а наскільки --- власні збережені дані. За високого рівня ацетилхоліну входи від органів чуття посилюються, а внутрішні, «пригадувальні» зв'язки слабшають, і водночас полегшується зміна ваг~\cite{Hasselmo2006}. Грубо кажучи, це перемикач «запис/читання» і разом із ним --- швидкість навчання $\alpha$.

\emph{\nt{SERO}, серотонін: горизонт планування.} Що передає серотонін, зрозуміло найгірше. Одна з гіпотез пов'язує його з коефіцієнтом дисконтування $\gamma$ з~\eqref{eq:rpe}: високий рівень серотоніну відповідає готовності почекати на велику винагороду, а не хапати малу зараз. Серотонінові нейрони працюють рівно й повільно, кілька імпульсів за секунду в неспанні, і майже замовкають в одній із фаз сну~\cite{JacobsAzmitia1992}.

Усі шість медіаторів зібрано в таблиці~\ref{tab:transmitters}.

\begin{table}[t]
\centering
\small
\begin{tabular}{>{\raggedright}p{2.4cm}>{\raggedright}p{3.0cm}>{\raggedright}p{2.0cm}>{\raggedright\arraybackslash}p{5.2cm}}
\toprule
Тип нейрона & Частка нейронів & Шина & Роль в обчисленні \\
\midrule
\nt{GLUT} (глутамат) & більшість & адресна & додатна вага, $w>0$ \\
\nt{GABA} (ГАМК) & 20--25\% у корі & адресна & від'ємна вага, $w<0$; стабілізація \\
\nt{DOPA} (дофамін) & $\sim 6\cdot10^{-6}$ & широко\-мовна & помилка передбачення $\delta$ \\
\nt{NORA} (норадреналін) & $\sim 6\cdot10^{-7}$ & широко\-мовна & коефіцієнт підсилення $g$ (гіпотеза) \\
\nt{ACET} (ацетилхолін) & мала & обидві & швидкість навчання $\alpha$; «запис/читання» (гіпотеза) \\
\nt{SERO} (серотонін) & $\sim 3\cdot10^{-6}$ & широко\-мовна & дисконтування $\gamma$ (гіпотеза) \\
\bottomrule
\end{tabular}
\caption{Головні медіатори мозку мовою обчислень. Правий стовпець --- сильне спрощення: він надійний для глутамату, ГАМК і дофаміну, для решти це гіпотези.}
\label{tab:transmitters}
\end{table}

\section{Знак задає отримувач}
\label{sec:receiver}

Я трохи вас обдурив, коли казав: глутамат --- додатна вага, ГАМК --- від'ємна. Жодна молекула не несе в собі знака. Молекула --- це просто ключ. Що станеться, коли її впіймають, вирішує \emph{замок} --- рецептор на клітині-отримувачі.

Найкращий приклад --- дофамін. Він має рецептори двох родин~\cite{Beaulieu2011}. На одних він полегшує збудження клітини, на інших --- утруднює. В області, яка обирає дії, одні нейрони несуть рецептори першого типу, інші --- другого, і той самий дофаміновий сигнал водночас підсилює одну групу й послаблює іншу. Це як один широкомовний пакет, який різні вузли мережі інтерпретують по-різному.

\begin{keyblock}
\begin{proposition}[Зміст повідомлення задає отримувач]
\label{prop:receptor}
Знак і швидкість дії медіатора визначає не сам медіатор, а рецептор, з яким він зв'язується. Рецептори бувають двох родів. \emph{Іонотропні} самі є іонними каналами й відкриваються за частки мілісекунди: це пряме керування струмом, як затвор транзистора. \emph{Метаботропні} запускають усередині клітини ланцюжок хімічних реакцій і діють за сотні мілісекунд і довше: це радше запис у регістр налаштувань, який змінює подальшу роботу клітини. Адресна шина говорить здебільшого через перші, широкомовна --- здебільшого через другі.
\end{proposition}
\end{keyblock}

Чому ж тоді взагалі можна сказати «глутамат збуджує»? Тому що майже всі рецептори глутамату, які трапляються на практиці, збуджувальні, а майже всі рецептори ГАМК --- гальмівні. Це не закон природи, а дуже надійна домовленість, і правило~\eqref{eq:dalesign} на ній і тримається.

\section{Що взяти з собою}

Переважна більшість нейронів --- адресна шина даних із вагами одного знака на нейрон; мізерна меншість --- широкомовні канали, що передають на великі ділянки мозку кілька спільних керувальних сигналів. Близько двох нейронів із кожних ста тисяч --- і від них залежить, як навчається і в якому режимі працює вся решта мережі. Для інженера це не дивно: у будь-якому комп'ютері регістрів керування незрівнянно менше, ніж комірок пам'яті, і все ж без них машина не працює.

У наступному розділі ми перевіримо, що може обчислити мережа з самих лише \nt{GLUT}-нейронів, найчисленнішого типу, яка працює в неперервному часі.

\section{Задачі}

\begin{exercise}[\lvlA]
\label{ex:01:1}
\emph{Купки й проводи.} За таблицею~\ref{tab:types} (середина діапазону за логарифмічною шкалою; \nt{ACET} --- $10^5$ виходів): (а)~якщо нейрон --- людина, скільки населень Землі становитимуть \nt{GLUT}-нейрони і місто якого розміру --- усі широкомовні? (б)~У скільки разів частка закінчень \nt{DOPA}, \nt{SERO}, \nt{NORA}, \nt{ACET} більша за частку їхніх нейронів?
\end{exercise}

\begin{exercise}[\lvlA]
\label{ex:01:2}
\emph{Повернення до нуля.} Глутамат у щілині спадає як $c_0e^{-t/\tau_c}$, $c_0=1\mM$, $\tau_c=1\ms$; отримувач помічає концентрацію вище $10$~мкМ. (а)~Скільки часу лінія «зайнята» після одного викиду? (б)~Насоси частково заблоковано, $\tau_c=10\ms$. До якої частоти викидів лінія ще встигає звільнитися?
\end{exercise}

\begin{exercise}[\lvlB]
\label{ex:01:3}
\emph{Куди переїжджає сплеск.} У досліді рис.~\ref{fig:rpe} після лампочки йдуть стани $s_0\to\dots\to s_{K-1}$ із кроком $\Delta$, $T=K\Delta$; винагорода $r$ приходить на останньому переході, момент лампочки непередбачуваний. Знайдіть вивчені $V(s_k)$ і відповідь $\delta$ на лампочку. Поклавши $\gamma=e^{-\Delta/\tau_\gamma}$, знайдіть $\tau_\gamma$ при $T=1$~с, $\Delta=0.1$~с, $\gamma=0.95$.
\end{exercise}

\begin{exercise}[\lvlB]
\label{ex:01:4}
\emph{Моделювання: навчання за помилкою.} Змоделюйте ланцюжок задачі~\ref{ex:01:3} при $K=10$, $\gamma=0.95$, $r=1$ із правилом $V(s_t)\leftarrow V(s_t)+\alpha\delta_t$. У якому випробуванні $n^*$ відповідь на лампочку вперше переганяє відповідь на сік при $\alpha=0.05$, $0.1$, $0.2$, $0.5$? Як $n^*$ залежить від $\alpha$?
\end{exercise}

\begin{exercise}[\lvlB]
\label{ex:01:5}
\emph{Проєктування: розсилка одного числа.} Число $\delta$ треба доставити всім $10^{10}$ нейронам; кожен отримувач, включно з ретрансляторами, має отримати $10$ закінчень від попереднього шару, ланка додає $1\ms$. Скільки нейронів і ланок у найощаднішому дереві ретрансляторів при $10^4$ виходах на нейрон (\nt{GLUT}) і при $10^{5.5}$ (\nt{DOPA})?
\end{exercise}

\begin{exercise}[\lvlB]
\label{ex:01:6}
\emph{Чому рівновага чутлива.} У мережі $80\,\%$ \nt{GLUT} і $20\,\%$ \nt{GABA}, кожен зв'язаний із кожним вагами $J_E/N$ і $-J_I/N$; $\tau\dot r=-r+Wr+h$. При $J_E=10$ єдине ненульове власне значення $W$ дорівнює $0.5$. Знайдіть відношення гальмівного струму до збуджувального і на скільки відсотків досить послабити гальмування, щоб мережа пішла в лавину.
\end{exercise}

\begin{exercise}[\lvlC]
\label{ex:01:7}
\emph{Дослідження: \nt{SERO} --- це $\gamma$?} За задачею~\ref{ex:01:3} відповідь \nt{DOPA}-нейронів на лампочку дорівнює $re^{-T/\tau_\gamma}$. Затримки $T=1,\dots,5$~с по $n$ пред'явлень, $\ln\delta$ вимірюється з розкидом $0.5$. Скільки потрібно $n$, щоб відрізнити $\tau_\gamma=2$~с від $2.4$~с (рівень $0.05$, потужність $0.8$)? Який механізм, крім $\gamma$, дав би такий самий спад із $T$?
\end{exercise}

\chapter{Що обчислюють \nt{GLUT}-нейрони}
\label{sec:glutcomputer}

\section{Правила гри}

У розділі~\ref{sec:neurontypes} ми побачили, що переважна більшість нейронів мозку --- \nt{GLUT}: вони лише збуджують, їхні ваги лише додатні. Візьмімо скільки завгодно таких нейронів і запитаймо: що може обчислити така мережа, якщо дозволяти собі лише те, що буває в живому організмі?

А в живому організмі немає тактового генератора, який перемикав би всі елементи одночасно. Кожен нейрон працює сам по собі, в неперервному часі: імпульси приходять і йдуть коли випаде, з різними затримками. Є \emph{входи} --- чутливі нейрони, що спрацьовують від світла, звуку чи тиску, --- і \emph{виходи} --- нейрони, що керують м'язами. Між ними --- мережа, і ніхто не каже їй, коли починати й коли закінчувати обчислення. Для електронника це \emph{асинхронна} схема з елементами, затримки яких заздалегідь точно не відомі.

Модель нейрона візьмімо найпростішу з тих, що чесно працюють у неперервному часі. Напруга $V$ на мембрані нейрона в спокої дорівнює нулю. Кожен імпульс, що надійшов від нейрона $j$, стрибком піднімає її на $w_j\ge0$, після чого напруга сама стікає назад до нуля зі сталою часу $\tau$:
\begin{empheq}[box=\keybox]{equation}
\tau\,\frac{\dd V}{\dd t} \;=\; -V \;+\; \tau\sum_j w_j \sum_k \delta\bigl(t-t_j^{(k)}-d_j\bigr), \qquad w_j\ge 0 .
\label{eq:glutneuron}
\end{empheq}
Тут $t_j^{(k)}$ --- моменти імпульсів нейрона $j$, $d_j$ --- затримка на шляху від нього, $\delta$ --- дельта-функція, яка й означає миттєвий стрибок. Коли $V$ досягає порога $\theta$, нейрон видає імпульс, $V$ скидається в нуль, і приблизно мілісекунду нейрон не може спрацювати знову. Типові числа: $\tau$ --- від часток мілісекунди до двадцяти мілісекунд у різних нейронів, затримки --- від часток мілісекунди до десятків мілісекунд. Це \emph{інтегрувальний нейрон із витоком}: RC-коло з порогом. Це свідоме спрощення: у нейронів кори гілки входів самі видають місцеві імпульси~\cite{Smith2013}, і один нейрон поводиться радше як невелика багатошарова мережа (розділ~\ref{sec:synthesis}). Для питань цього розділу досить і одного RC-кола.

Головна відмінність від логічного вентиля --- витік: нейрон пам'ятає імпульс не вічно, а приблизно $\tau$, і додається лише те, що надійшло в межах цього вікна.

\section{АБО та І}

Почнімо з вентилів. Якщо один імпульс піднімає напругу вище порога, $w\ge\theta$, нейрон спрацьовує на кожен імпульс будь-якого входу. Це \emph{АБО}.

З І цікавіше. Нехай $w<\theta<2w$: одного імпульсу замало, потрібні два. Але тепер питання «чи прийшли обидва» не має сенсу, доки не сказано, \emph{у межах якого часу}. Нехай перший імпульс прийшов у момент $0$, а другий --- через час $\Delta$. На момент приходу другого від першого залишилося $we^{-\Delta/\tau}$, і нейрон спрацює, якщо $we^{-\Delta/\tau}+w\ge\theta$. Звідси вікно збігу:
\begin{empheq}[box=\keybox]{equation}
\Delta \;\le\; \Delta_{\max} \;=\; \tau\,\ln\frac{w}{\theta-w} .
\label{eq:window}
\end{empheq}
Це не логічне І, а \emph{детектор збігів}: І в часі (рис.~\ref{fig:coincidence}). При $\theta=1.5w$ вікно дорівнює $\tau\ln2\approx0.7\tau$. У нейрона кори з $\tau=10\ms$ це близько семи мілісекунд. У нейронів слухової системи, де $\tau$ менша за мілісекунду, вікно стискається до часток мілісекунди.

\begin{figure}[t]
\centering
\begin{tikzpicture}[>=Stealth,x=0.9cm,y=1.6cm]
\foreach \dx/\lab/\c [count=\i] in {0.5/{збіглися}/red!70!black, 2.6/{не збіглися}/blue!60!black}{
 \begin{scope}[xshift={(\i-1)*7.2cm}]
  \draw[->,gray!70!black] (0,0) -- (6.3,0) node[below left,black] {\small $t$};
  \draw[->,gray!70!black] (0,0) -- (0,1.75) node[left,black] {\small $V$};
  \draw[densely dashed,gray!60!black] (0,1.5) -- (6.0,1.5) node[right,black] {\small $\theta$};
  \draw[very thick,\c] (0,0) -- (1,0) -- (1,1)
      plot[domain=1:1+\dx,samples=30] (\x,{exp(-(\x-1)/1.5)});
  \pgfmathsetmacro{\v}{exp(-\dx/1.5)+1}
  \draw[very thick,\c] (1+\dx,{exp(-\dx/1.5)}) -- (1+\dx,{min(\v,1.5)});
  \ifdim\v pt<1.5pt
    \draw[very thick,\c] plot[domain=1+\dx:6,samples=40] (\x,{\v*exp(-(\x-1-\dx)/1.5)});
  \else
    \draw[very thick,\c] (1+\dx,1.5) -- (1+\dx,0) -- (6,0);
    \draw[->,thick,\c] (1+\dx,1.55) -- (1+\dx,1.72) node[above,font=\scriptsize] {імпульс};
  \fi
  \draw[->,thick] (1,-0.35) -- (1,-0.08); \draw[->,thick] (1+\dx,-0.35) -- (1+\dx,-0.08);
  \draw[decorate,decoration={brace,amplitude=3pt,mirror}] (1,-0.42) -- (1+\dx,-0.42) node[midway,below=3pt] {\scriptsize $\Delta$};
  \node[\c] at (3.5,-0.95) {\small \lab};
 \end{scope}}
\end{tikzpicture}
\caption{Детектор збігів, поріг $\theta=1.5w$. Ліворуч другий імпульс прийшов через $\Delta<\Delta_{\max}$, сума досягла порога, нейрон спрацював. Праворуч $\Delta>\Delta_{\max}$: від першого імпульсу майже нічого не лишилося, і нейрон змовчав.}
\label{fig:coincidence}
\end{figure}

Інший спосіб отримати І --- через тривалість, а не через точний час. Поки горить світло, чутливий нейрон видає імпульси часто й рівно; якщо порогові замало одного такого входу, а двох досить, нейрон працює, поки активні обидва. Так мережа обчислює за \emph{рівнями}, і точний час приходу імпульсів неважливий. Більша частина того, що далі, буде саме так.

\section{Чого зробити не можна}

Тепер спробуймо зробити НЕ: нейрон, який працює, коли вхід мовчить. Не вийде: без вхідних імпульсів у~\eqref{eq:glutneuron} немає жодного стрибка, і напруга стоїть на нулі, нижче порога. А мережа з багатьох нейронів? Теж ні, і це доводиться відразу для всіх мереж.

Найзручніше говорити про частоти. Нехай $r_i$ --- частота імпульсів нейрона $i$, $I_i$ --- приплив від входів, а $f_i$ --- його передавальна характеристика: як частота залежить від сумарного припливу. В інтегрувального нейрона ця характеристика неспадна: більший приплив --- частіші імпульси. Частоти в мережі, що дійшла рівноваги, задовольняють рівняння
\begin{equation}
r_i \;=\; f_i\Bigl(\sum_j w_{ij}\,r_j + I_i\Bigr), \qquad w_{ij}\ge 0 .
\label{eq:ratenet}
\end{equation}

\begin{keyblock}
\begin{proposition}[Монотонність]
\label{prop:monotone}
Нехай мережа~\eqref{eq:ratenet} складається лише з \nt{GLUT}-нейронів і починає з повної тиші. Якщо посилити входи, то усталена частота жодного нейрона не зменшиться. Усе, що обчислює така мережа, --- \emph{монотонні} функції входів.
\end{proposition}
\end{keyblock}

Доведення. Почнімо з тиші, $r^{(0)}=0$, і підставляймо частоти в праву частину~\eqref{eq:ratenet}: $r^{(k+1)}=F\bigl(r^{(k)}\bigr)$. Права частина $F$ не спадає за жодним аргументом, бо ваги невід'ємні, а $f_i$ не спадають. Оскільки $r^{(1)}\ge0=r^{(0)}$, за індукцією $r^{(k+1)}\ge r^{(k)}$: частоти лише зростають і приходять до найменшого розв'язку~\eqref{eq:ratenet}. Якщо входи $I$ замінити на більші $I'$, то на кожному кроці $r^{(k)}(I)\le r^{(k)}(I')$, і в границі теж. Справжня мережа приходить до рівноваги не кроками, а неперервно, але для неї справджується те саме: за додатних зв'язків нерівність між двома прогонами, виконана на початку, зберігається весь час.

Для окремих імпульсів твердження правдиве не буквально: зайвий вхідний імпульс може змусити нейрон спрацювати трохи раніше, і через мілісекунду несприйнятливості пропаде його наступний імпульс. Але ефект слабкий і ненадійний, обчислювати на ньому не можна. Для частот монотонність строга.

Наслідки ті самі, що в будь-якої схеми без інверторів. Немає НЕ. Немає виключного АБО: додавання другого входу не може погасити вихід. І найнеприємніше: \emph{активність не можна зупинити активністю}, можна лише прибрати те, що її підтримує.

\section{Два проводи: вмикання й вимикання}

Вихід із глухого кута підказує сам організм. У багатьох органах чуття подразник передається двома наборами нейронів: у зорі одні клітини відповідають на посилення світла, інші --- на послаблення~\cite{Schiller1992}. Назвімо їх \emph{вмикальними} та \emph{вимикальними}. Замість одного проводу $x$ мережа отримує пару $(x,\bar x)$, і відсутність, яку вона не вміє обчислити, їй повідомляє сам датчик.

З парою проводів НЕ безплатне: досить поміняти проводи місцями. І та АБО робляться кожне двома нейронами:
\begin{empheq}[box=\keybox]{equation}
\begin{aligned}
x\wedge y \;&\to\; \bigl(\;x\wedge y,\;\; \bar x\vee\bar y\;\bigr),\\
x\vee y \;&\to\; \bigl(\;x\vee y,\;\; \bar x\wedge\bar y\;\bigr).
\end{aligned}
\label{eq:dualrail}
\end{empheq}
Праворуч усі операції монотонні, тож твердження~\ref{prop:monotone} не порушено.

Для асинхронної схеми двопровідний код має й другу перевагу, навіть важливішу. Пара проводів буває в трьох станах: «так», «ні» і, коли обидва мовчать, \emph{«відповіді ще немає»}. Отримувач дізнається, що обчислення закінчено, не за годинником, а за самим сигналом: у кожній парі запрацював рівно один провід. Між подразниками датчики мовчать, і мережа повертається в тишу, готова до наступного разу. В асинхронній електроніці на цьому прийомі будують схеми, які працюють правильно за будь-яких затримок елементів~\cite{Sparso2001}.

\begin{keyblock}
\begin{proposition}[Асинхронне обчислення]
\label{prop:dualrail}
Якщо кожен вхід надходить парою проводів «вмикання---вимикання», то мережа з \nt{GLUT}-нейронів може обчислити будь-яку логічну функцію входів. Відповідь теж надходить парами проводів, а її готовність видно з того, що в кожній парі запрацював один провід. Годинник для цього не потрібен. Ціна --- приблизно вдвічі більше нейронів.
\end{proposition}
\end{keyblock}

Приклад --- виключне АБО: «змінився рівно один із двох подразників». Прямий провід відповіді --- $(x\wedge\bar y)\vee(\bar x\wedge y)$, зворотний --- $(x\wedge y)\vee(\bar x\wedge\bar y)$. Це шість нейронів, і жоден із них не потребує гальмування (рис.~\ref{fig:dualxor}).

\begin{figure}[t]
\centering
\begin{tikzpicture}[>=Stealth,x=1cm,y=1cm,
  nrn/.style={circle,draw,very thick,fill=blue!6,minimum size=0.75cm,inner sep=0pt,font=\small},
  inp/.style={font=\small}]
\node[inp] (X)  at (0,3.3) {$x$};
\node[inp] (Xb) at (0,2.2) {$\bar x$};
\node[inp] (Y)  at (0,1.1) {$y$};
\node[inp] (Yb) at (0,0)   {$\bar y$};
\node[nrn] (A1) at (3.2,3.3) {$2$};
\node[nrn] (A2) at (3.2,2.2) {$2$};
\node[nrn] (A3) at (3.2,1.1) {$2$};
\node[nrn] (A4) at (3.2,0)   {$2$};
\node[nrn] (O)  at (7.2,2.75) {$1$};
\node[nrn] (Ob) at (7.2,0.55) {$1$};
\foreach \s/\t in {X/A1,Yb/A1,Xb/A2,Y/A2,X/A3,Y/A3,Xb/A4,Yb/A4}{\draw[->,thick,blue!60!black] (\s) -- (\t);}
\draw[->,thick,blue!60!black] (A1) -- node[pos=0.45,above,sloped,font=\scriptsize,black] {$x\wedge\bar y$} (O);
\draw[->,thick,blue!60!black] (A2) -- node[pos=0.45,below,sloped,font=\scriptsize,black] {$\bar x\wedge y$} (O);
\draw[->,thick,blue!60!black] (A3) -- node[pos=0.45,above,sloped,font=\scriptsize,black] {$x\wedge y$} (Ob);
\draw[->,thick,blue!60!black] (A4) -- node[pos=0.45,below,sloped,font=\scriptsize,black] {$\bar x\wedge\bar y$} (Ob);
\draw[->,very thick,red!70!black] (O) -- (8.6,2.75) node[right,black] {$x\oplus y$: «так»};
\draw[->,very thick,red!70!black] (Ob) -- (8.6,0.55) node[right,black] {$\overline{x\oplus y}$: «ні»};
\node[below,font=\small,gray!40!black] at (3.2,-0.55) {І: поріг $2$};
\node[below,font=\small,gray!40!black] at (7.2,-0.55) {АБО: поріг $1$};
\node[below,font=\small,gray!40!black] at (0,-0.55) {пари проводів};
\end{tikzpicture}
\caption{Виключне АБО на двопровідному коді, лише з \nt{GLUT}-нейронів. Число в кружечку --- поріг в одиницях ваги одного входу. Чотири нейрони І впізнають чотири комбінації входів, два нейрони АБО збирають їх у пару проводів відповіді.}
\label{fig:dualxor}
\end{figure}

\section{Затримки замість годинника}

Для обчислень із часом досить затримок у проводах і детекторів збігів. Найкращий приклад --- те, як мозок визначає, звідки прийшов звук.

Звук від джерела, розташованого збоку, надходить у ближнє вухо раніше, ніж у дальнє. Різниця залежить від напрямку: від нуля, коли джерело прямо попереду, до $0.22\,\text{м}/343\,\text{м/с}\approx0.6\ms$ у людини, коли воно точно збоку. Мозок розрізняє різницю близько десяти \emph{мікросекунд}~\cite{KlumppEady1956,Grothe2010}, хоча самі нейрони спрацьовують із точністю не кращою за частки мілісекунди. Як?

Проведімо від лівого вуха провід уздовж ряду детекторів збігів зліва направо, а від правого --- справа наліво (рис.~\ref{fig:itd}). Сигнал біжить проводом зі швидкістю $v$. До детектора, що стоїть на відстані $x$ від лівого краю ряду довжини $L$, лівий сигнал іде час $x/v$, правий --- $(L-x)/v$. Якщо в праве вухо звук прийшов пізніше на $\delta t$, два сигнали зустрінуться одночасно в тій точці, де
\begin{empheq}[box=\keybox]{equation}
\frac{x}{v} \;=\; \frac{L-x}{v} + \delta t
\quad\Longrightarrow\quad
x \;=\; \frac{L}{2} + \frac{v\,\delta t}{2} .
\label{eq:itd}
\end{empheq}
Спрацює лише той детектор, до якого обидва сигнали прийшли в межах вікна~\eqref{eq:window}. Номер детектора, що спрацював, і є напрямок на джерело. Різниця в часі перетворилася на \emph{місце}.

\begin{figure}[t]
\centering
\begin{tikzpicture}[>=Stealth,
  nrn/.style={circle,draw,very thick,fill=blue!6,minimum size=0.7cm,inner sep=0pt}]
\foreach \k in {0,...,6}{\node[nrn] (D\k) at (1.4*\k,0) {\scriptsize $2$};}
\node[nrn,fill=red!15] at (D4) {\scriptsize $2$};
\draw[->,thick,blue!60!black] (-1.4,0.9) node[left,black] {\small ліве вухо} -- (9.2,0.9);
\draw[->,thick,orange!85!black] (9.8,-0.9) node[right,black] {\small праве вухо} -- (-0.8,-0.9);
\foreach \k in {0,...,6}{
  \draw[->,blue!60!black] (1.4*\k,0.9) -- (D\k.north);
  \draw[->,orange!85!black] (1.4*\k,-0.9) -- (D\k.south);}
\draw[->,thick,red!70!black] (D4.north east) ++(0.1,0.05) -- ++(0.5,0.45) node[above right,font=\small,black] {спрацював};
\node[font=\small,gray!40!black] at (4.2,-1.5) {затримка зростає $\longrightarrow$ \hspace{2.2cm} $\longleftarrow$ затримка зростає};
\pic[scale=1.4] at (-2.3,1.75) {speaker};
\node[font=\scriptsize,align=center] at (-2.3,2.35) {джерело ліворуч};
\end{tikzpicture}
\caption{Визначення напрямку на звук. Сигнали від двох вух біжать назустріч один одному; кожен кружечок --- детектор збігів із порогом $2$. Якщо звук прийшов у праве вухо пізніше, сигнали зустрічаються правіше від середини, за формулою~\eqref{eq:itd}. Вихід мережі --- номер нейрона, що спрацював. Тут джерело ліворуч: у ліве вухо звук прийшов раніше.}
\label{fig:itd}
\end{figure}

Тут немає ні годинника, ні гальмування, ні навіть двопровідного коду: мережа не відповідає «так» чи «ні», а вказує на один нейрон із багатьох. Така схема з ліній затримки й детекторів збігів, схоже, справді працює у слуховому стовбурі птахів~\cite{Carr1990}. Мікросекундна точність виникає не з точності кожного нейрона, а з того, що детекторів багато й кожен дивиться на свою затримку. У ссавців ряду таких затримок не знайшли: там ту саму задачу, схоже, розв'язують детектори збігів, налаштовані точно розрахованим у часі гальмуванням від \nt{GLYC}-нейронів~\cite{Brand2002,Grothe2010}. Як гальмування звужує вікно збігу, ми розберемо в розділі~\ref{sec:gaba} на прикладі \nt{GABA}; гліцин гальмує так само, відкриваючи в отримувача хлорні канали.

\section{Пам'ять}

Комп'ютерові потрібна пам'ять. У такій мережі пам'ять --- це активність, яка підтримує сама себе. Візьмімо групу \nt{GLUT}-нейронів, сильно зв'язаних одне з одним, і опишімо її однією частотою $r$ (рис.~\ref{fig:recurrentgroup}а). У рівновазі $r=f(wr+I)$, де $w$ --- сила зворотного зв'язку групи на саму себе, а $I$ --- приплив ззовні. Розв'яжімо це рівняння графічно, перетинаючи криву $f(wr+I)$ з прямою $r$ (рис.~\ref{fig:bistable}).

\begin{figure}[t]
\centering
% (b): tau dr/dt = -r + f(w r + I), f as in fig:bistable, w=1, I=0.6; Euler dt=0.001 tau.
\begin{tikzpicture}[>=Stealth,
  nrn/.style={circle,draw,thick,fill=blue!6,minimum size=0.5cm,inner sep=0pt}]
\begin{scope}
\node[font=\small] at (-0.5,1.55) {(а)};
\foreach \k in {1,...,6}{\node[nrn] (n\k) at ({1.4+1.0*cos(60*\k)},{1.0*sin(60*\k)}) {};}
\foreach \a/\b in {1/2,2/3,3/4,4/5,5/6,6/1,1/4,2/5,3/6}{\draw[<->,blue!60!black] (n\a) -- (n\b);}
\foreach \k in {2,3,4}{\draw[->,thick,green!45!black] ($(n\k)+(-0.95,0)$) -- (n\k);}
\node[font=\small,green!45!black] at (-0.35,-0.45) {$I$};
\node[font=\Large] at (3.1,0) {$\approx$};
\node[nrn,minimum size=0.85cm,font=\small] (R) at (4.35,-0.1) {$r$};
\draw[->,thick,blue!60!black] (R) edge[loop above,looseness=6] node[above,font=\small] {$w$} (R);
\draw[->,thick,green!45!black] (3.55,-0.95) node[left,font=\small] {$I$} -- (R);
\end{scope}
\begin{scope}[xshift=6.3cm,y=2.6cm,x=1.05cm]
\node[font=\small] at (-0.6,1.25) {(б)};
\draw[->,gray!70!black] (0,0) -- (7.3,0) node[below left,font=\small,black] {$t/\tau$};
\draw[->,gray!70!black] (0,0) -- (0,1.12) node[left,font=\small,black] {$r$};
\foreach \t in {0,2,4,6}{\draw[gray!60!black] (\t,-0.02) -- (\t,0.02); \node[below,font=\scriptsize] at (\t,0) {\t};}
\draw[densely dotted,gray!60!black] (0,0.5) -- (7,0.5) node[right,font=\scriptsize,black,align=left] {поріг\\запису};
\fill[green!45!black,opacity=0.25] (1,-0.2) rectangle (2,-0.13);
\fill[gray!60!black,opacity=0.35] (1,-0.29) rectangle (1.5,-0.22);
\draw[very thick,blue!60!black] plot coordinates {(0.0,0.002) (0.1,0.002) (0.2,0.002) (0.3,0.002) (0.4,0.002) (0.5,0.002) (0.6,0.002) (0.7,0.002) (0.8,0.002) (0.9,0.002) (1.0,0.002) (1.1,0.083) (1.2,0.165) (1.3,0.242) (1.4,0.313) (1.5,0.378) (1.6,0.437) (1.7,0.491) (1.8,0.539) (1.9,0.583) (2.0,0.623) (2.1,0.643) (2.2,0.665) (2.3,0.688) (2.4,0.71) (2.5,0.732) (2.6,0.753) (2.7,0.773) (2.8,0.792) (2.9,0.81) (3.0,0.826) (3.2,0.855) (3.4,0.88) (3.6,0.9) (3.8,0.917) (4.0,0.931) (4.4,0.953) (4.8,0.967) (5.2,0.977) (5.6,0.984) (6.0,0.988) (6.5,0.992) (7.0,0.994)};
\draw[very thick,gray!60!black,densely dashed] plot coordinates {(0.0,0.002) (0.1,0.002) (0.2,0.002) (0.3,0.002) (0.4,0.002) (0.5,0.002) (0.6,0.002) (0.7,0.002) (0.8,0.002) (0.9,0.002) (1.0,0.002) (1.1,0.083) (1.2,0.165) (1.3,0.242) (1.4,0.313) (1.5,0.378) (1.6,0.358) (1.7,0.336) (1.8,0.313) (1.9,0.291) (2.0,0.269) (2.2,0.228) (2.4,0.191) (2.6,0.159) (2.8,0.133) (3.0,0.11) (3.2,0.091) (3.4,0.076) (3.6,0.063) (3.8,0.052) (4.0,0.043) (4.4,0.03) (4.8,0.021) (5.2,0.015) (5.6,0.011) (6.0,0.008) (6.5,0.006) (7.0,0.004)};
\node[font=\scriptsize,blue!60!black,anchor=west] at (3.3,0.8) {приплив $1\tau$: записано};
\node[font=\scriptsize,gray!50!black,anchor=west] at (2.3,0.3) {приплив $0.5\tau$: згасло};
\end{scope}
\end{tikzpicture}
\caption{Група, сильно зв'язана сама з собою. (а)~Багато \nt{GLUT}-нейронів, що збуджують одне одного, поводяться як один нейрон із частотою $r$, який збуджує сам себе із силою $w$; ззовні надходить приплив $I$. (б)~Частота групи в часі при $w=1$ і тій самій кривій $f$, що на рис.~\ref{fig:bistable}. Приплив $I=0.6$ подається на час, позначений смужкою під віссю. Якщо він триває $1\tau$, група встигає перейти нестійку точку $r=0.5$, розгоряється й лишається горіти, коли приплив закінчився. Якщо лише $0.5\tau$, вона не дотягує до цієї точки й згасає назад: щоб записати біт, приплив має тривати довше приблизно $0.7\tau$.}
\label{fig:recurrentgroup}
\end{figure}

\begin{figure}[t]
\centering
\begin{tikzpicture}[>=Stealth,x=5cm,y=3.2cm]
\draw[->,gray!70!black] (0,0) -- (1.1,0) node[below left,black] {\small $r$};
\draw[->,gray!70!black] (0,0) -- (0,1.1);
\draw[thick,gray!60!black] (0,0) -- (1.05,1.05) node[right,black] {\small $r$};
\draw[very thick,blue!60!black] plot[domain=0:1.05,samples=80] (\x,{1/(1+exp(-(\x-0.5)/0.08))});
\draw[very thick,red!70!black,densely dashed] plot[domain=0:1.05,samples=80] (\x,{1/(1+exp(-(\x+0.3-0.5)/0.08))});
\fill[blue!60!black] (0.002,0.002) circle (0.018);
\draw[blue!60!black,fill=white,thick] (0.5,0.5) circle (0.018);
\fill[blue!60!black] (0.998,0.998) circle (0.018);
\node[below right,font=\small] at (0.02,0.0) {вимк.};
\node[right,font=\small] at (0.52,0.46) {нестійко};
\node[below,font=\small] at (0.99,0.96) {увімк.};
\node[anchor=east,font=\small,red!70!black] at (0.19,0.62) {$f(wr+I)$};
\node[anchor=west,font=\small,blue!60!black] at (0.45,0.1) {$f(wr)$};
\end{tikzpicture}
\caption{Пам'ять із групи \nt{GLUT}-нейронів із сильним зворотним зв'язком. Суцільна крива --- без зовнішнього припливу: два стійкі перетини з прямою $r$, «вимкнено» й «увімкнено», і нестійкий між ними. Короткий приплив $I$ (штрихова крива) залишає лише верхній.}
\label{fig:bistable}
\end{figure}

Якщо зворотний зв'язок сильний, перетинів три: нижній --- тиша, верхній --- група горить на повну, середній нестійкий. Вийшов елемент із двома стійкими станами, тригер. Записати в нього одиницю легко: короткий приплив піднімає криву, нижній перетин зникає, і група розгоряється --- і лишається горіти, коли приплив закінчився (рис.~\ref{fig:recurrentgroup}б). Надто короткий приплив не встигає довести групу до нестійкої точки, і вона згасає назад. Таку самопідтримувану активність, що триває секунди після того, як подразник зник, справді спостерігають у мозку й вважають основою короткочасної пам'яті~\cite{Wang2001}. Але це не єдиний спосіб пам'ятати секунди. Запис можна зберігати й у повільній змінній самих синапсів: після пачки імпульсів полегшення, як у синапсу-«тире» розділу~\ref{sec:code}, триває близько секунди, і мережа між короткими спалахами може майже мовчати --- не тригер, а заряджений конденсатор~\cite{Mongillo2008}. Такі «тихі» проміжки під час утримання в пам'яті спостерігають~\cite{Stokes2015}, і зберігання без постійних імпульсів дешевше за енергією. Який із двох способів кора використовує в якому випадку, обговорюється.

А як стерти? За твердженням~\ref{prop:monotone} активністю --- ніяк; лишається щось \emph{прибрати}. Перший спосіб --- прибрати підтримку: якщо група горить лише разом із припливом від іншої групи, пам'ять згасає разом із ним. Другий --- утома. У справжніх синапсів за тривалої роботи запас медіатора виснажується~\cite{Tsodyks1997}, а в нейронів наростають повільні струми, що знижують збудливість. Зворотний зв'язок слабшає, верхній перетин зникає, і група згасає сама через секунди. Виходить пам'ять, яка вмикається сигналом, а вимикається лише часом.

\section{Послідовності}

Замість годинника можна мати \emph{таймер, що запускається подією}. З'єднаймо групи \nt{GLUT}-нейронів у ланцюжок: перша збуджує другу, друга --- третю, і так далі. Зовнішній сигнал запалює першу групу, і хвиля активності пробігає ланцюжком, кожна ланка через одну-дві затримки після попередньої і лише один раз: відразу після імпульсу нейрони несприйнятливі, а хвиля вже пішла далі.

Такий ланцюжок відмірює час від події: спрацювала ланка номер $k$ --- отже, від моменту запуску минуло $k$ затримок. Його виходи можна подати на м'язи й отримати послідовність рухів, яка розгортається сама. У співочих птахів під час співу окремі \nt{GLUT}-нейрони одного з відділів мозку спрацьовують суворо по черзі, кожен у свій момент пісні, --- дуже схоже на такий ланцюжок~\cite{Hahnloser2002}.

На відміну від годинника, ланцюжок мовчить, доки його не запустили, відпрацьовує один раз і відмірює час лише для тих, хто до нього під'єднаний. Спільного такту в мережі й далі немає.

\section{Чого бракує}

З самих лише \nt{GLUT}-нейронів, без гальмування, виходить чимало. Але трьох речей бракує, і всіх трьох через твердження~\ref{prop:monotone}.

\emph{Вибір.} Мережа має обирати один напрямок, одну дію. Якщо два подразники майже рівні, у мережі рис.~\ref{fig:itd} спрацюють обидва виходи, і придушити слабший на користь сильнішого нічим.

\emph{Скидання.} Стерти пам'ять сигналом не можна.

\emph{Стійкість.} У мережі, де зв'язки виникають не за планом інженера, петлі додатного зворотного зв'язку неминучі, й активність у них може наростати лавиною (розділ~\ref{sec:neurontypes}). Єдине гальмо --- та сама втома, повільна й груба.

Усі три біди лікуються одними ліками --- нейроном, що гасить імпульс імпульсом. Це \nt{GABA}, і потрібен він не для того, щоб обчислювати, а для того, щоб обирати, скидати й тримати обчислення під контролем.

\section{Задачі}

\begin{exercise}[\lvlA]
\label{ex:02:1}
\emph{Ряд детекторів для звуку.} Сигнали від вух біжать проводами рис.~\ref{fig:itd} зі швидкістю $5$~м/с, найбільша різниця приходу $0.64\ms$. Детектори --- нейрони~\eqref{eq:glutneuron} з $\tau=0.5\ms$, $\theta=1.5w$ --- стоять так, що сусідні розрізняють $10$~мкс. Скільки детекторів у ряду і скільки з них спрацює на один звук?
\end{exercise}

\begin{exercise}[\lvlB]
\label{ex:02:2}
\emph{Нерівні входи зсувають вікно.} Детектор збігів~\eqref{eq:glutneuron} з $\tau=0.5\ms$, $\theta=1.5$ отримує по імпульсу від двох входів із вагами $1$ і $0.8$. Знайдіть вікно збігу та його центр. На скільки мікросекунд помилиться ряд рис.~\ref{fig:itd}, якщо вхід від одного вуха на $20\,\%$ слабший, ніж від іншого?
\end{exercise}

\begin{exercise}[\lvlB]
\label{ex:02:3}
\emph{Які мережі не коливаються.} Мережа $\tau\dot r_i=-r_i+f_i\bigl(\textstyle\sum_jw_{ij}r_j+I_i\bigr)$ з неспадними $f_i$ і $w_{ij}\ge0$ при $i\ne j$ монотонна й стійко не коливається. Доведіть: заміною $r_i\to\pm r_i$ до неї зводиться мережа з парною кількістю гальмівних зв'язків у кожному циклі. Чи може коливатися взаємне гальмування двох груп? Петля \nt{GLUT}--\nt{GABA}?
\end{exercise}

\begin{exercise}[\lvlB]
\label{ex:02:4}
\emph{Проєктування: двопровідний суматор.} Кожен біт передається парою проводів $(x,\bar x)$. Побудуйте повний суматор трьох бітів, що видає пари $(S,\bar S)$ суми та $(C,\bar C)$ перенесення, з \nt{GLUT}-нейронів, указавши ваги й пороги. Вкладіться у вісім нейронів. Скільки потребувала б схема з елементів І та АБО, як на рис.~\ref{fig:dualxor}?
\end{exercise}

\begin{exercise}[\lvlB]
\label{ex:02:5}
\emph{Моделювання: скільки триває запис.} Група рис.~\ref{fig:recurrentgroup}: $\tau\dot r=-r+f(r+I)$, $f(u)=1/\bigl(1+e^{-(u-0.5)/0.08}\bigr)$, до запису --- у нижньому стані; приплив $I$ подається на час $T$. Знайдіть чисельно найменше $T$, що записує біт при $I=0.6$, і поріг $I_c$, нижче якого біт не записується за жодного $T$.
\end{exercise}

\begin{exercise}[\lvlA]
\label{ex:02:6}
\emph{Ланцюжок як таймер.} Ланка ланцюжка --- $100$ \nt{GLUT}-нейронів, затримка $2\ms$ з незалежним розкидом $0.2\ms$. (а)~Скільки нейронів потрібно, щоб відміряти $1$~с, і яка відносна похибка? Як вона залежить від інтервалу? (б)~У справжнього таймера похибка --- $5\,\%$ за будь-якого інтервалу. Яке джерело розкиду це дає?
\end{exercise}

\begin{exercise}[\lvlC]
\label{ex:02:7}
\emph{Дослідження: тригер чи конденсатор.} $100$ нейронів тримають біт $10$~с. Тригер: кожен дає $20$ імпульсів за секунду. Конденсатор: полегшення $u$ спадає з $\tau_F=1$~с, біт читається при $u\ge u_{\max}/2$, освіження --- пачка з трьох імпульсів на нейрон. У скільки разів конденсатор ощадливіший і що буде з бітом після глушіння групи на $200\ms$?
\end{exercise}

\chapter{\nt{GLUT} і \nt{GABA} разом}
\label{sec:gaba}

\section{Правила гри}

Мережа з самих лише \nt{GLUT}-нейронів не вміє обирати, скидати пам'ять і втримувати активність від лавини, і все через монотонність (твердження~\ref{prop:monotone}). Додаймо до неї другий за чисельністю тип нейронів --- \nt{GABA}. Правила ті самі: лише те, що буває в живому організмі, і жодного годинника.

Модель нейрона та сама, що~\eqref{eq:glutneuron}, але імпульс від \nt{GABA}-нейрона не піднімає напругу отримувача, а опускає. Ваги тепер бувають двох знаків, і знак задає відправник, правило~\eqref{eq:dalesign}.

Кілька параметрів схеми. \nt{GABA}-нейронів у корі від п'ятої частини до чверті~\cite{Hendry1987}. Виходи в них короткі: вони гальмують сусідів, а не далекі області. Гальмування приходить через мілісекунду й триває кілька мілісекунд. Без входу \nt{GABA}-нейрон мовчить: щоб загальмувати когось, він сам мусить отримати збудження, зазвичай від \nt{GLUT}-нейронів тієї самої мережі.

Тому від'ємний зв'язок від \nt{GLUT}-нейрона $A$ до нейрона $B$ доводиться робити через посередника: $A$ збуджує \nt{GABA}-нейрон, а той гальмує $B$. Зміна знака коштує одного нейрона й однієї затримки, близько мілісекунди.

Для гальмування важливо не лише, \emph{від кого} йде зв'язок, а й \emph{куди} він сідає: місце посадки багато в чому визначає, що зв'язок робить~\cite{Markram2004}. \nt{GLUT}-виходи сідають здебільшого на вхідні гілки отримувача, його \emph{дендрити}, і несуть дані: кожен такий вхід --- один доданок суми. Вихід на тіло нейрона діє на суму цілком: так багато швидких \nt{GABA}-нейронів ділять відповідь отримувача й задають, коли йому спрацювати. Вихід на те місце, де народжується імпульс отримувача, --- останнє слово, заборона на весь вихід одразу. А вихід на закінчення чужого проводу змінює ймовірність викиду $p$ однієї вхідної лінії, не чіпаючи решти~\cite{Rudomin1999}; так, зокрема, працюють ворота болю в розділі~\ref{sec:pain}. За межами мозку виходи сідають на м'язові волокна (розділ~\ref{sec:muscles}) і на органи (розділ~\ref{sec:chemoutput}), а деякі не сідають нікуди й викидають медіатор в об'єм тканини або в кров (розділи~\ref{sec:serocomputer} і~\ref{sec:chemoutput}).

\section{НЕ і заборона}

Чи можна тепер зробити НЕ? Майже. Нейрон, що працює при мовчазному вході, мусить звідкись брати збудження. У живому мозку воно є: на кожен нейрон постійно надходить фонова активність від тисяч інших. Нехай фон тримає нейрон $Y$ у роботі, а вхід $x$ через \nt{GABA}-посередника його гальмує. Це НЕ.

Корисніша, однак, \emph{заборона}: «$a$, але не при $b$»:
\begin{empheq}[box=\keybox]{equation}
y \;=\; a \wedge \bar b .
\label{eq:veto}
\end{empheq}
Нейрон $Y$ отримує збудження від $a$ і гальмування від \nt{GABA}-нейрона, якого збуджує $b$. Фон тут не потрібен: збудження дає сам $a$. Із заборон і АБО збирається все. Наприклад, виключне АБО: $x\oplus y=(x\wedge\bar y)\vee(\bar x\wedge y)$ --- два нейрони-заборони, два \nt{GABA}-посередники й один нейрон АБО.

\begin{keyblock}
\begin{proposition}[Повнота \nt{GLUT}+\nt{GABA}]
\label{prop:gabacomplete}
Мережа з \nt{GLUT}- і \nt{GABA}-нейронів може обчислити будь-яку логічну функцію входів, і кожен біт передається одним проводом. Датчики «вмикання---вимикання» з розділу~\ref{sec:glutcomputer} для цього більше не потрібні. Якщо функція має давати одиницю при всіх мовчазних входах, потрібне джерело фонової активності.
\end{proposition}
\end{keyblock}

Доведення: заборона~\eqref{eq:veto} разом з АБО функціонально повна, якщо є стала одиниця, бо НЕ $x$ --- це заборона «одиниця, але не при $x$». А функції, які при всіх мовчазних входах дають нуль, збираються із заборон і АБО й без одиниці.

\section{Напрямок руху}

Заборона в часі, як і І в часі в розділі~\ref{sec:glutcomputer}, дає детектор напрямку руху.

Нехай два сусідні датчики, $A$ і $B$, стоять на відстані $\Delta x$ один від одного. Вихідний нейрон $Y$ отримує збудження від $B$ і гальмування від $A$ через \nt{GABA}-посередника, із затримкою $d$ (рис.~\ref{fig:direction}). Пляма світла, що біжить зі швидкістю $v$ від $A$ до $B$, проходить $A$ у момент $0$, і гальмування приходить до $Y$ у момент $d$; $B$ вона проходить у момент $\Delta x/v$, і тоді ж приходить збудження. Якщо ці моменти збігаються з точністю до вікна~\eqref{eq:window}, гальмування гасить збудження, і $Y$ мовчить. Це відбувається за швидкості близько
\begin{empheq}[box=\keybox]{equation}
v^* \;=\; \frac{\Delta x}{d} .
\label{eq:vstar}
\end{empheq}
За руху від $B$ до $A$ збудження від $B$ приходить у момент $0$, а гальмування від $A$ --- лише в момент $\Delta x/v+d$, коли $Y$ уже спрацював.

\begin{figure}[t]
\centering
\begin{tikzpicture}[>=Stealth,
  nrn/.style={circle,draw,very thick,fill=blue!6,minimum size=0.9cm,inner sep=0pt},
  gab/.style={nrn,fill=red!10},
  inh/.style={-{Bar[width=7pt]},thick,red!70!black}]
\node[nrn] (A) at (0,2) {$A$};
\node[nrn] (B) at (3,2) {$B$};
\node[gab] (G) at (0.9,0.6) {\small \nt{GABA}};
\node[nrn] (Y) at (3,-0.6) {$Y$};
\draw[->,thick] (A) -- (G);
\draw[inh] (G) -- node[below left=-2pt] {\scriptsize затримка $d$} (Y);
\draw[->,thick] (B) -- (Y);
\draw[->,very thick,red!70!black] (Y) -- (4.6,-0.6) node[right,black] {\small вихід};
\draw[->,thick,orange!85!black] (-0.6,2.9) -- (3.6,2.9) node[right,black] {\small мовчить};
\draw[<-,thick,orange!85!black] (-0.6,3.4) -- (3.6,3.4) node[right,black] {\small відповідає};
\foreach \yy/\dd in {2.9/-1,3.4/1}{
  \foreach \k in {1,2,3}{\draw[orange!70!yellow,line width=0.8pt] (1.5+\dd*0.12+\dd*0.13*\k,\yy+0.1-0.1*\k+0.1) -- ++(\dd*0.25,0);}
  \fill[yellow!70!orange,draw=orange!70!black] (1.5,\yy) circle (0.14);}
\node[font=\scriptsize,align=center] at (1.5,3.95) {пляма світла};
\end{tikzpicture}
\caption{Детектор напрямку руху. $B$ збуджує $Y$, $A$ гальмує його через \nt{GABA}-посередника із затримкою $d$. За руху від $A$ до $B$ зі швидкістю близько $\Delta x/d$ гальмування гасить збудження; за зворотного руху воно запізнюється. Жовта пляма з рисками --- світло, що біжить над датчиками.}
\label{fig:direction}
\end{figure}

У сітківці вибірковість до напрямку руху, схоже, створюється саме так: асиметричним гальмуванням із того боку, звідки рух не повинен викликати відповіді~\cite{Barlow1965}.

\section{Віднімання чи ділення}

Досі гальмування в нас віднімало. Насправді воно влаштоване тонше.

Синапс відкриває канал, тобто вмикає провідність. У збуджувального синапсу рівноважний потенціал лежить далеко вище порога, приблизно на $70\mV$ вище спокою: відкритий канал тягне напругу вгору. У гальмівного \nt{GABA}-синапсу канал пропускає хлор, рівноважний потенціал якого лежить біля спокою, і відкритий канал тягне напругу не вниз, а \emph{до спокою}. Якщо відлічувати напругу від спокою, усталена напруга мембрани з провідністю витоку $g_L$, збуджувальною $g_E$ і гальмівною $g_I$ дорівнює
\begin{empheq}[box=\keybox]{equation}
V_\infty \;=\; \frac{g_E\,E_E}{g_L+g_E+g_I} ,
\label{eq:shunt}
\end{empheq}
де $E_E\approx70\mV$ --- рівноважний потенціал збуджувального синапсу. Це закон Ома для подільника напруги: $g_E$ тягне до $E_E$, а $g_L$ і $g_I$ --- до нуля.

$g_I$ стоїть не в чисельнику зі знаком мінус, а в знаменнику: гальмування не \emph{віднімає} від збудження, а \emph{ділить} його (рис.~\ref{fig:shunt}). Таке гальмування називають шунтувальним: відкриті хлорні канали закорочують мембрану. Для мережі це регулювання підсилення. Якщо $g_I$ пропорційна сумарній активності сусідів, кожен нейрон відповідає не на абсолютну силу свого входу, а на її частку в загальній активності. Таке ділення на загальну активність трапляється в багатьох відділах мозку й вважається однією з його стандартних операцій~\cite{Carandini2012}: та сама мережа працює і за слабкого, і за сильного входу.

Застереження: формула~\eqref{eq:shunt} --- про нейрон, який не видає імпульсів. У нейрона, що спрацьовує, напруга весь час тримається біля порога, струм через гальмівну провідність майже сталий, і на \emph{частоту} імпульсів шунт діє здебільшого відніманням~\cite{HoltKoch1997}. Частота ділиться, якщо нейронові одночасно додати врівноважені збуджувальну й гальмівну фонові провідності, як у тремтливому, врівноваженому режимі живої кори~\cite{Chance2002}. Тож ділення мозок отримує не від одного шунта, а від гальмування разом із фоновим шумом~\cite{Carandini2012}.

\begin{figure}[t]
\centering
\begin{tikzpicture}[>=Stealth,x=1.25cm,y=0.05cm]
\draw[->,gray!70!black] (0,0) -- (5.4,0) node[right,black] {\small $g_E/g_L$};
\draw[->,gray!70!black] (0,0) -- (0,68) node[above,black] {\small $V_\infty$, мВ};
\foreach \v in {20,40,60}{\draw[gray!60!black] (-0.05,\v) -- (0.05,\v) node[left=3pt,font=\scriptsize] {\v};}
\foreach \x in {1,2,3,4,5}{\draw[gray!60!black] (\x,-1.5) -- (\x,1.5) node[below=5pt,font=\scriptsize] {\x};}
\draw[very thick,blue!60!black] plot[domain=0:5,samples=60] (\x,{70*\x/(1+\x)});
\draw[very thick,red!70!black] plot[domain=0:5,samples=60] (\x,{70*\x/(3+\x)});
\draw[thick,densely dashed,gray!50!black] plot[domain=0.56:5,samples=60] (\x,{70*\x/(1+\x)-25});
\node[right,font=\small,blue!60!black] at (5.02,58.3) {без гальмування};
\node[right,font=\small,red!70!black] at (5.02,45.5) {ділить: $g_I=2g_L$};
\node[right,font=\small,gray!40!black] at (5.02,32.5) {віднімає $25\mV$};
\end{tikzpicture}
\caption{Шунтувальне гальмування ділить напругу, а не віднімає від неї. Синя крива --- усталена напруга~\eqref{eq:shunt} без гальмування, червона --- з гальмівною провідністю $g_I=2g_L$. Червона --- та сама синя, розтягнута по горизонталі втричі: вхід ніби поділили на $1+g_I/g_L=3$. Штрихова --- гальмування, що віднімає сталі $25\mV$: крива опускається цілком, нахил не змінюється.}
\label{fig:shunt}
\end{figure}

Є й другий наслідок. Стала часу мембрани $\tau_{\text{еф}}=C/(g_L+g_E+g_I)$, і гальмування її вкорочує. А вікно збігу~\eqref{eq:window} пропорційне $\tau$, тож гальмування, що прийшло разом зі збудженням, \emph{звужує вікно}. Схема, у якій вхід збуджує нейрон безпосередньо й водночас, із затримкою в мілісекунду, гальмує його через \nt{GABA}-посередника, --- одна з найпоширеніших у мозку: нейрон встигає відповісти лише на те, що прийшло в першу мілісекунду-другу~\cite{Pouille2001}.

\section{Вибір}

Тепер головне, чого бракувало \nt{GLUT}-мережі: вибір одного з кількох. Візьмімо дві групи \nt{GLUT}-нейронів із входами $I_1$ і $I_2$. Кожна через своїх \nt{GABA}-посередників гальмує іншу із силою $\beta$ (рис.~\ref{fig:wta}). Для частот $r_1,r_2$ отримуємо
\begin{equation}
\tau\,\frac{\dd r_1}{\dd t} = -r_1 + \bigl[I_1-\beta r_2\bigr]_+ ,\qquad
\tau\,\frac{\dd r_2}{\dd t} = -r_2 + \bigl[I_2-\beta r_1\bigr]_+ ,
\label{eq:wta}
\end{equation}
де $[u]_+=\max(u,0)$: частота не буває від'ємною.

\begin{figure}[t]
\centering
\begin{tikzpicture}[>=Stealth,
  nrn/.style={circle,draw,very thick,fill=blue!6,minimum size=0.9cm,inner sep=0pt},
  gab/.style={nrn,fill=red!10,minimum size=0.75cm},
  inh/.style={-{Bar[width=7pt]},thick,red!70!black}]
\node[nrn] (E1) at (0,0) {$r_1$};
\node[nrn] (E2) at (4,0) {$r_2$};
\node[gab] (G1) at (1.4,1.1) {};
\node[gab] (G2) at (2.6,-1.1) {};
\draw[->,thick] (E1) -- (G1); \draw[inh] (G1) -- (E2);
\draw[->,thick] (E2) -- (G2); \draw[inh] (G2) -- (E1);
\draw[->,thick,blue!60!black] (-1.6,0) node[left,black] {$I_1$} -- (E1);
\draw[->,thick,blue!60!black] (5.6,0) node[right,black] {$I_2$} -- (E2);
\end{tikzpicture}
\caption{Вибір із двох. Дві групи \nt{GLUT}-нейронів (сині) гальмують одна одну через \nt{GABA}-посередників (червоні).}
\label{fig:wta}
\end{figure}

Поки працюють обидва нейрони, система~\eqref{eq:wta} лінійна, і її поведінку задають два власні значення: $-(1+\beta)/\tau$ для однакових відхилень обох частот і $-(1-\beta)/\tau$ для протилежних. За слабкого гальмування, $\beta<1$, обидва значення від'ємні: обидва нейрони працюють, гальмування лише приглушує слабшого. За сильного гальмування, $\beta>1$, друге значення стає додатним: будь-яка різниця між $r_1$ і $r_2$ наростає, доки один нейрон не замовкне зовсім.

\begin{keyblock}
\begin{proposition}[Переможець забирає все]
\label{prop:wta}
Якщо взаємне гальмування сильніше за одиницю, $\beta>1$, стан, у якому працюють обидві групи, нестійкий. Мережа приходить у стан, де працює одна група, а інша мовчить. Переможець із частотою $r_1=I_1$ утримує перемогу, поки $I_2<\beta I_1$.
\end{proposition}
\end{keyblock}

Ми перевірили це на моделі. При $\beta=0.5$ і входах $1.0$ та $0.9$ обидві групи працюють, із частотами $0.73$ і $0.53$. При $\beta=2$ і тих самих входах друга група мовчить зовсім. Такий вибір має пам'ять: якщо потім поміняти входи місцями, колишній переможець лишається переможцем, бо $1.0<2\cdot0.9$.

Із сотнею груп усе так само: кожна гальмує решту, виживає одна.

\section{Скидання пам'яті}

Пам'ять із розділу~\ref{sec:glutcomputer} вимикалася лише втомою. Додаймо до неї \nt{GABA}-нейрон, який збуджується сигналом «скидання» і гальмує групу. Гальмування опускає криву $f(wr+I)$ на рис.~\ref{fig:bistable}, верхній перетин зникає, і група згасає --- і лишається в нижньому стані, коли скидання закінчилося. Тепер пам'ять записується одним сигналом і стирається іншим, як у тригера зі входами встановлення й скидання.

Ще красивіше з'єднати дві такі групи взаємним гальмуванням, як на рис.~\ref{fig:wta}, і дати кожній зворотний зв'язок на саму себе. Сигнал на вхід однієї з груп переводить комірку в її стан, і комірка пам'ятає останнє рішення. Це прямий аналог засувки з двох перехресно з'єднаних елементів АБО-НЕ.

\section{Стійкість і ритм}

Лишається третя біда \nt{GLUT}-мережі --- лавина. Візьмімо групу \nt{GLUT}-нейронів зі зворотним зв'язком на себе $w_{EE}$ і групу \nt{GABA}-нейронів, яких перша збуджує із силою $w_{IE}$ і які гальмують першу із силою $w_{EI}$. Для частот $E$ та $I$ отримуємо
\begin{equation}
\tau_E\frac{\dd E}{\dd t} = -E + w_{EE}E - w_{EI}I + h, \qquad
\tau_I\frac{\dd I}{\dd t} = -I + w_{IE}E ,
\label{eq:ei}
\end{equation}
де $h$ --- зовнішній вхід~\cite{WilsonCowan1972}. Без гальмування при $w_{EE}>1$ активність зростає без меж: кожен імпульс породжує більше ніж один наступний. З гальмуванням усталена частота
\begin{equation}
E^* \;=\; \frac{h}{1-w_{EE}+w_{EI}w_{IE}}
\label{eq:eistar}
\end{equation}
скінченна, якщо знаменник додатний. Але цього замало: рівновага має ще й притягувати. Із лінійного аналізу~\eqref{eq:ei} випливають дві умови.

\begin{keyblock}
\begin{proposition}[Умови стійкості]
\label{prop:ei}
Рівновага~\eqref{eq:eistar} стійка, якщо гальмування
\begin{enumerate}
\item[(i)] \emph{досить сильне}: $w_{EI}\,w_{IE} > w_{EE}-1$;
\item[(ii)] \emph{досить швидке}: $\tau_I < \tau_E/(w_{EE}-1)$.
\end{enumerate}
Якщо порушено (i), активність наростає лавиною. Якщо виконано (i), але порушено (ii), гальмування запізнюється, й активність починає коливатися.
\end{proposition}
\end{keyblock}

Умова~(i) --- це визначник матриці системи~\eqref{eq:ei}, умова~(ii) --- її слід. Зміст другої зрозумілий кожному, хто налаштовував регулятор: запізнілий зворотний зв'язок не гасить відхилення, а розгойдує його.

\begin{figure}[t]
\centering
\begin{tikzpicture}[>=Stealth]
\draw[->,gray!70!black] (0,0) -- (10.9,0) node[below left,black] {\small $t$, мс};
\draw[->,gray!70!black] (0,0) -- (0,3.3) node[left,black] {\small $E$};
\foreach \t in {0,50,100,150}{\draw[gray!70!black] (\t*0.07,0.06) -- (\t*0.07,-0.06) node[below,font=\scriptsize] {\t};}
\input{figures/ei_traces}
\node[font=\small,gray!40!black,anchor=west] at (0.9,3.15) {без гальмування: лавина};
\node[font=\small,blue!60!black,anchor=west] at (7.4,1.25) {швидке гальмування};
\node[font=\small,red!70!black,anchor=west] at (7.4,2.55) {повільне гальмування};
\end{tikzpicture}
\caption{Частота \nt{GLUT}-групи зі зворотним зв'язком $w_{EE}=1.5$, $\tau_E=10\ms$, після ввімкнення входу. Без гальмування (штрих) активність зростає без меж. З гальмуванням сили $w_{EI}w_{IE}=2$ і $\tau_I=2\ms$ (синя крива) вона виходить на рівень $E^*=h/(1-1.5+2)=0.67h$. З тим самим гальмуванням, але повільним, $\tau_I=30\ms$ (червона крива), активність коливається.}
\label{fig:ei}
\end{figure}

Вважають, що кора працює саме в такому режимі: її власний зворотний зв'язок досить сильний, щоб без гальмування піти в лавину, і тримає її в робочій точці лише швидке гальмування~\cite{Tsodyks1997isn}.

Цей режим має парадоксальний підпис, за яким його можна впізнати ззовні~\cite{Tsodyks1997isn}. Додаймо до~\eqref{eq:ei} зовнішній вхід $h_I$ безпосередньо \nt{GABA}-групі. Усталені частоти стануть
\begin{equation}
E^* \;=\; \frac{h-w_{EI}\,h_I}{D},\qquad
I^* \;=\; \frac{w_{IE}\,h+(1-w_{EE})\,h_I}{D},\qquad
D \;=\; 1-w_{EE}+w_{EI}w_{IE} .
\label{eq:isn}
\end{equation}
Якщо $w_{EE}>1$, коефіцієнт при $h_I$ у другій формулі від'ємний: підштовхнеш \nt{GABA}-нейрони --- і їхня частота \emph{впаде}. Причина в петлі. Зайве гальмування приглушує \nt{GLUT}-групу, та менше збуджує \nt{GABA}-нейрони, і вони втрачають більше, ніж отримали. За чисел рис.~\ref{fig:ei} ($w_{EE}=1.5$, $w_{EI}w_{IE}=2$, $D=1.5$) кожна одиниця доданого входу знижує $I^*$ на третину одиниці. Цей підпис і знайшли в корі: коли відповідь нейронів зорової кори пригнічується подразником навколо, гальмівний струм, який вони отримують, не зростає, а падає разом зі збуджувальним~\cite{Ozeki2009}. Пізніше той самий підпис знайшли в різних областях і шарах кори~\cite{Sanzeni2020}.

Коливання в разі порушення умови~(ii) у мозку теж трапляються. Петля «\nt{GLUT} збуджує \nt{GABA}, \nt{GABA} гальмує \nt{GLUT}» із затримкою в кілька мілісекунд породжує ритм із періодом порядку $12$--$50\ms$ (частота $20$--$80$\,Гц), і такі швидкі ритми в корі добре відомі~\cite{Cardin2009,Buzsaki2012}. Це не годинник комп'ютера: ритм місцевий і виникає лише тоді, коли мережа активна. Але протягом кількох циклів він ділить час на вікна, у яких нейрони можуть спрацювати.

\section{Підсумок}

\begin{table}[t]
\centering
\small
\begin{tabular}{>{\raggedright}p{3.3cm}>{\raggedright}p{4.4cm}>{\raggedright\arraybackslash}p{4.4cm}}
\toprule
& лише \nt{GLUT} & \nt{GLUT} + \nt{GABA} \\
\midrule
НЕ & лише через датчики «вмикання---вимикання» & заборона $a\wedge\bar b$; НЕ за фонової активності \\
проводів на біт & два & один \\
напрямок руху & немає & заборона із затримкою \\
регулювання підсилення & немає & ділення: шунт разом із фоновим шумом \\
вікно збігу & $\sim\tau$ & звужується гальмуванням \\
вибір одного з багатьох & немає & взаємне гальмування, $\beta>1$ \\
скидання пам'яті & лише втомою & сигналом через \nt{GABA} \\
стійкість & лише втомою & швидкий від'ємний зворотний зв'язок \\
ритм & не потрібен & виникає за повільного гальмування \\
\bottomrule
\end{tabular}
\caption{Що додає \nt{GABA} до мережі з \nt{GLUT}-нейронів.}
\label{tab:gaba}
\end{table}

\nt{GABA} майже не додає мережі нових \emph{обчислень} --- усе це можна було обчислити й із двопровідними датчиками, --- а додає \emph{керування} (таблиця~\ref{tab:gaba}): вибір, скидання, регулювання підсилення й стійкість. Збудження в мозку несе дані, а гальмування вирішує, яким даним пройти, коли й наскільки гучно. Для кожного четвертого-п'ятого нейрона кори це дуже розумний поділ праці.

\section{Задачі}

\begin{exercise}[\lvlA]
\label{ex:03:1}
\emph{Шунт у числах.} $E_E=70\mV$. За~\eqref{eq:shunt} знайдіть $V_\infty$ при $g_E=g_L$ і $4g_L$, без гальмування й із шунтом $g_I=2g_L$. Яке $V_\infty$ при $g_E=4g_L$ дало б віднімання, що знімає при $g_E=g_L$ стільки ж, скільки шунт? У скільки разів шунт звужує вікно збігу?
\end{exercise}

\begin{exercise}[\lvlB]
\label{ex:03:2}
\emph{Виведення умов стійкості.} Знайдіть слід і визначник матриці лінійної системи~\eqref{eq:ei} і виведіть із них обидві умови твердження~\ref{prop:ei}. На межі умови~(ii) рівновага втрачає стійкість коливально. Знайдіть період цих коливань за чисел рис.~\ref{fig:ei}.
\end{exercise}

\begin{exercise}[\lvlB]
\label{ex:03:3}
\emph{Ритм із затримки.} Замінімо повільне гальмування моделі~\eqref{eq:ei} швидким із затримкою $d$: $\tau_E\dot E=-E(t)-GE(t-d)+h$. При $\tau_E=10\ms$ знайдіть найменшу затримку $d_c$, що дає коливання, і їхній період для $G=2$ і $G=5$. Чи потрапляють вони у швидкі ритми кори, $12$--$50\ms$, на відміну від задачі~\ref{ex:03:2}?
\end{exercise}

\begin{exercise}[\lvlB]
\label{ex:03:4}
\emph{Моделювання: вибір із пам'яттю.} Проінтегруйте~\eqref{eq:wta} при $\beta=2$, $\tau=1$. (а)~При $I_1=1$ повільно піднімайте $I_2$ від $0$ до $3$ і опускайте назад: за яких $I_2$ мережа перемикається? (б)~На скільки зростає час рішення з тиші, коли різниця входів $\varepsilon$ зменшується вдесятеро?
\end{exercise}

\begin{exercise}[\lvlB]
\label{ex:03:5}
\emph{Проєктування: детектор на смугу швидкостей.} Детектор рис.~\ref{fig:direction} має мовчати за руху від $A$ до $B$ із часом пробігу $5$--$20\ms$. \nt{GABA}-посередник із затримкою $d_k$ забороняє $Y$ на відрізку $[d_k,\,d_k+5\ms]$ після імпульсу $A$; збудження від $B$ приходить одразу. Скільки потрібно посередників і з якими затримками?
\end{exercise}

\begin{exercise}[\lvlB]
\label{ex:03:6}
\emph{Заборони й фон.} Скільки нейронів потребує парність $x\oplus y\oplus z$ у загальній схемі «\nt{GABA}-посередник на вхід, заборона на кожен набір з $f=1$, спільне АБО»? Побудуйте її з двох нейронів, \nt{GABA} і \nt{GLUT}, що отримують входи безпосередньо, указавши ваги й пороги.
\end{exercise}

\begin{exercise}[\lvlC]
\label{ex:03:7}
\emph{Проєктування: вибір зі ста.} Кожна зі $100$ груп \nt{GLUT} має однаково гальмувати всі інші, але не себе. Лінійні \nt{GABA}-посередники збуджуються групами й гальмують їх; групи можуть збуджувати одна одну, але не себе. Яка найменша кількість посередників? А якщо прямих збуджувальних зв'язків немає зовсім?
\end{exercise}

\begin{exercise}[\lvlC]
\label{ex:03:8}
\emph{Дослідження: коли гальмування парадоксальне.} У моделі рис.~\ref{fig:ei} ($w_{EI}=2$, $w_{IE}=1$) \nt{GLUT}-група отримала передавальну функцію $f(u)=[u]_+^2$, \nt{GABA}-група лишилася лінійною. Вище якого входу $h_c$ добавка на \nt{GABA}-групу знижує її ж частоту? Перевірте моделюванням.
\end{exercise}

\chapter{Абетка Морзе: що ми знаємо про мову, якою говорять \nt{GLUT}- і \nt{GABA}-нейрони}
\chaptermark{Абетка Морзе}
\label{sec:code}

\section{Алфавіт з одного знака}

В абетці Морзе два знаки, крапка й тире, і паузи трьох довжин між ними. У нейрона, як ми бачили в розділі~\ref{sec:neurontypes}, алфавіт ще бідніший: знак один. Імпульс триває близько мілісекунди, і всі імпульси одного нейрона однакові за висотою й формою. Тому все повідомлення --- в паузах, тобто в послідовності моментів $t_1,t_2,t_3,\dots$ Це абетка без тире, у якій зміст несе лише ритм крапок (рис.~\ref{fig:morse}).

\begin{figure}[t]
\centering
\begin{tikzpicture}[>=Stealth,x=0.08cm,y=1cm]
% Morse letter: dot, dash, dot
\node[left,font=\small] at (-6,2.55) {Морзе:};
\fill[black] (0,2.45) rectangle (6,2.65);
\fill[black] (14,2.45) rectangle (32,2.65);
\fill[black] (40,2.45) rectangle (46,2.65);
\node[right,font=\small,gray!40!black] at (52,2.55) {крапка, тире, крапка --- зміст у тривалостях і паузах};
% stimulus onset
\node[left,font=\small] at (-6,0.4) {нейрон:};
\draw[->,thick,green!45!black] (0,-0.55) -- (0,-0.05);
\node[below,font=\scriptsize,green!45!black] at (0,-0.55) {подразник};
\draw[gray!70!black] (0,0) -- (152,0) node[right,black] {\small $t$};
% spikes: single ones blue, the burst red
\foreach \s in {14,26,41,88,112,131}{\draw[very thick,blue!60!black] (\s,0) -- (\s,0.8);}
\foreach \s in {60,63,66,69}{\draw[very thick,red!70!black] (\s,0) -- (\s,0.8);}
% latency
\draw[<->,thick] (0,1.1) -- node[above,font=\scriptsize] {$t_1$: час першого імпульсу} (14,1.1);
% burst
\draw[decorate,decoration={brace,amplitude=4pt},red!70!black] (58,0.95) -- (71,0.95) node[midway,above=4pt,font=\scriptsize] {пачка --- «тире»};
% counting windows
\foreach \a/\b/\n in {0/50/3,50/100/5,100/150/2}{
  \draw[decorate,decoration={brace,amplitude=4pt,mirror},gray!50!black] (\a+1,-0.2) -- (\b-1,-0.2) node[midway,below=4pt,font=\small] {$n=\n$};}
\node[font=\scriptsize,gray!40!black] at (75,-1.05) {частотний код: кількість імпульсів $n$ у кожному вікні};
\foreach \x in {0,50,100,150}{\draw[gray!60!black] (\x,-0.06) -- (\x,0.06);}
\end{tikzpicture}
\caption{Абетка Морзе й абетка нейрона. Ту саму послідовність імпульсів можна прочитати по-різному: порахувати імпульси у вікнах по $50\ms$ (підрозділ~\ref{sec:rate}), виміряти затримку $t_1$~\eqref{eq:latency} або помітити пачку --- «тире»~\eqref{eq:burst}.}
\label{fig:morse}
\end{figure}

Які паузи взагалі можливі? Знизу їх обмежує несприйнятливість: після імпульсу нейрон близько мілісекунди не може спрацювати знову, тож більше кількох сотень імпульсів за секунду не буває. Згори не обмежує ніщо: нейрон може мовчати хвилинами. У \nt{GLUT}-нейронів кори частоти лежать у межах від часток імпульсу до кількох десятків імпульсів за секунду, і більшість нейронів більшу частину часу працює біля нижньої межі.

У розділах~\ref{sec:glutcomputer} і~\ref{sec:gaba} ми мимохідь приймали то один, то інший спосіб запису чисел імпульсами. Тепер запитаймо прямо: якою мовою \nt{GLUT}- і \nt{GABA}-нейрони говорять одне з одним? Повної відповіді немає, цю мову не розшифровано. Але дещо про неї відомо твердо, і майже все відоме можна отримати, міркуючи як інженер зв'язку: пропускна здатність каналу, шум, приймач, ціна передачі.

\section{Скільки можна передати}

Почнімо з верхньої межі. Нехай отримувач розрізняє моменти імпульсів із точністю $\Delta t$: зсуви, менші за $\Delta t$, для нього нерозрізненні. Розріжмо час на клітинки довжини $\Delta t$, менші за час несприйнятливості, тож у кожній клітинці або один імпульс, або жодного (рис.~\ref{fig:cells}). За середньої частоти $r$ імпульс потрапляє в клітинку з імовірністю $p=r\Delta t$. Найбільше інформації потік несе, коли клітинки незалежні, і тоді кожна клітинка дає ентропію $-p\log_2p-(1-p)\log_2(1-p)\approx p\log_2(e/p)$ біт при $p\ll1$. Поділивши на $\Delta t$, отримаємо граничну швидкість передачі та її частку на один імпульс~\cite{MacKay1952}:
\begin{empheq}[box=\keybox]{equation}
R_{\max} \;\approx\; r\,\log_2\frac{e}{r\,\Delta t}\ \ \text{біт/с},
\qquad
\frac{R_{\max}}{r} \;\approx\; \log_2\frac{e}{r\,\Delta t}\ \ \text{біт на імпульс}.
\label{eq:spikeentropy}
\end{empheq}
При $r=10$ імпульсів за секунду й $\Delta t=1\ms$ це близько восьми біт на імпульс і вісімдесяти біт за секунду.

\begin{figure}[t]
\centering
\begin{tikzpicture}[>=Stealth,x=0.5cm,y=1cm]
% time cut into cells of width Delta t; a cell holds one impulse or none
\foreach \k in {0,...,23}{\draw[gray!55] (\k,0) rectangle (\k+1,0.7);}
\foreach \k in {2,7,8,15,21}{
  \fill[red!10] (\k,0) rectangle (\k+1,0.7);
  \draw[very thick,red!70!black] (\k+0.5,0.05) -- (\k+0.5,0.65);}
\foreach \k in {0,...,23}{
  \pgfmathtruncatemacro{\bit}{(\k==2)||(\k==7)||(\k==8)||(\k==15)||(\k==21)}
  \ifnum\bit=1 \def\bitcol{red!70!black}\else \def\bitcol{gray!60!black}\fi
  \node[font=\scriptsize,text=\bitcol] at (\k+0.5,-0.25) {\bit};}
\draw[<->,gray!60!black] (0,0.95) -- (1,0.95) node[midway,above,font=\scriptsize,black] {$\Delta t$};
\node[font=\small,anchor=east] at (-0.3,0.35) {імпульси};
\node[font=\small,anchor=east] at (-0.3,-0.25) {біти};
\draw[decorate,decoration={brace,amplitude=5pt,mirror},gray!60!black] (0,-0.55) -- (24,-0.55)
  node[midway,below=6pt,font=\scriptsize,black,align=center] {отримувач, який лише рахує: «$n=5$»;\\ отримувач, який бачить клітинки: які саме $5$ із $24$ --- це $\log_2\binom{24}{5}\approx15$ біт};
\end{tikzpicture}
\caption{Звідки береться пропускна здатність~\eqref{eq:spikeentropy}. Час нарізано на клітинки довжини $\Delta t$, у кожній клітинці імпульс або є ($1$), або немає ($0$): потік імпульсів --- це двійковий рядок. Імпульс потрапляє в клітинку з імовірністю $p=r\Delta t$. Що дрібніші клітинки, то довший рядок і то більше біт; лічильник імпульсів втрачає майже все, що записано в тому, \emph{де} стоять одиниці.}
\label{fig:cells}
\end{figure}

Біти на імпульс залежать від точності часу лише логарифмічно: за точності $10\ms$ залишиться близько п'яти біт на імпульс. Але якщо отримувач лише рахує імпульси за секунду, то при $r=10$ він отримає менше двох біт за всю секунду (підрозділ~\ref{sec:rate}).

\begin{keyblock}
\begin{proposition}[Пропускна здатність імпульсного каналу]
\label{prop:capacity}
Потік імпульсів із середньою частотою $r$, прочитаний із точністю часу $\Delta t$, несе не більше~\eqref{eq:spikeentropy}. Майже вся ця пропускна здатність міститься в \emph{часі} імпульсів: отримувач, який лише рахує імпульси, видобуває з того самого потоку в десятки разів менше, ніж той, хто розрізняє їхні моменти з точністю до мілісекунди.
\end{proposition}
\end{keyblock}

Це верхня межа, а не вимірювання. Вимірювання на чутливих нейронах, для яких відомий подразник, дають від одного до кількох біт на імпульс; у найкращих випадках це до половини межі, обчисленої для їхньої точності~\cite{Rieke1997}. Отже, принаймні на вході в нервову систему час імпульсів використовується, а не пропадає.

\section{Шумний канал}

Пропускну здатність~\eqref{eq:spikeentropy} пораховано без шуму. Про те, де він виникає, твердо відомі три факти.

\emph{Сам генератор імпульсів точний.} Якщо в нейрон кори, вийнятий із мозку разом зі шматочком тканини, багато разів подати той самий струм, що швидко змінюється в часі, імпульси повторюються з точністю близько мілісекунди~\cite{Mainen1995}. Якщо струм сталий, моменти імпульсів від разу до разу розповзаються: точний час породжують швидкі перепади входу, а не сам нейрон.

\emph{Синапс ненадійний.} Імпульс, що добіг до \nt{GLUT}-синапсу, викидає порцію медіатора лише з певною ймовірністю $p$. На синапсах гіпокампа понад половина імпульсів до отримувача просто не доходить, і причина --- саме випадковість викиду, а не збої проведення~\cite{AllenStevens1994}. Для зв'язківця це канал зі стираннями; кілька синапсів між двома нейронами почасти рятують справу.

\emph{У живій корі потік імпульсів виглядає випадковим.} Інтервали між імпульсами розкидані приблизно так само, як у випадкового пуассонівського потоку, а при повторенні того самого подразника кількість імпульсів за вікно змінюється так, що її дисперсія близька до середнього~\cite{Softky1993}.

Як точний елемент породжує випадковий потік? Нейрон отримує тисячі входів, кожен ненадійний, а збудження й гальмування врівноважені, як у розділі~\ref{sec:gaba}. Напруга мембрани тремтить біля порога, і моменти імпульсів задають ці тремтіння. Шум це чи повідомлення, яке ми не вміємо прочитати, --- багато в чому відкрите питання. Частина «шуму» вже виявилася повідомленням: у зоровій корі значна частка розкиду слідує за рухами самої тварини~\cite{Stringer2019}. Складність тут принципова: добре стиснене повідомлення для того, хто не знає коду, не відрізнити від випадкового потоку.

\section{Частота: повільно, зате надійно}
\label{sec:rate}

Найпростіший код --- \emph{частотний}: число записано в тому, скільки імпульсів нейрон видає за вікно $T$. Ні стирання, ні тремтіння моментів такому кодові не страшні. Але він має ціну. Якщо потік пуассонівський, кількість імпульсів $n$ за вікно має середнє $rT$ і розкид $\sqrt{rT}$:
\begin{empheq}[box=\keybox]{equation}
\frac{\sigma_n}{\bar n} \;=\; \frac{1}{\sqrt{r\,T}} .
\label{eq:rateerror}
\end{empheq}
Щоб дізнатися частоту з точністю в десять відсотків, потрібно сто імпульсів, за двадцяти імпульсів на секунду --- п'ять секунд. Сусідні рівні розрізненні, якщо вони віддалені на пару розкидів, тому до частоти $r$ за час $T$ розрізняється близько $\sqrt{rT}$ рівнів: при $r=10$ і $T=1\,\text{с}$ це три-чотири рівні, менше двох біт.
Мозок так повільно не працює. Людина впізнає тварину на картинці, показаній на мить, і відповідь у мозку з'являється приблизно через $150\ms$~\cite{Thorpe1996}. За грубою оцінкою, за цей час сигнал проходить близько десятка послідовних ступенів обробки, по $10$--$20\ms$ на ступінь. За звичайних частот кори кожен нейрон встигає видати на ступені нуль або один імпульс, і його частота в такому вікні не визначена. Виходів два: взяти багато нейронів або дивитися на час.

\emph{Багато нейронів.} Нехай $N$ нейронів несуть те саме число, і отримувач додає їхні імпульси. Якщо їхні шуми незалежні, у~\eqref{eq:rateerror} замість $rT$ стає $NrT$: сто нейронів при $r=20$ за $15\ms$ дають тридцять імпульсів і точність близько двадцяти відсотків. Простір замінює час. Однак шуми сусідніх нейронів кори не незалежні: коефіцієнт кореляції кількостей імпульсів у пари сусідів зазвичай близько $0.1$~\cite{Zohary1994}. Якщо він дорівнює $c$, дисперсія середнього за $N$ нейронами дорівнює
\begin{empheq}[box=\keybox]{equation}
\sigma^2_N \;=\; \sigma^2_1\,\frac{1+(N-1)\,c}{N}
\;\xrightarrow[N\to\infty]{}\; c\,\sigma^2_1 .
\label{eq:correlated}
\end{empheq}

\begin{keyblock}
\begin{proposition}[Межа усереднення]
\label{prop:pooling}
За корельованого шуму усереднення за групою зменшує похибку не більше ніж у $1/\sqrt{c}$ разів. При $c=0.1$ це втричі, і такого виграшу досягають уже на кількох десятках нейронів; додавати більше марно.
\end{proposition}
\end{keyblock}

Не кожна кореляція однаково шкідлива. Обмежує точність лише та частина спільного шуму, яка \emph{схожа на сам сигнал}, тобто зсуває всю групу так само, як її зсунула б зміна вимірюваної величини~\cite{MorenoBote2014}. Скільки такого шуму в мозку насправді, ще з'ясовують.

Є й інший спосіб використати багато нейронів: записати число в тому, \emph{який} нейрон спрацював, як у визначнику напрямку на звук (рис.~\ref{fig:itd}). Це код \emph{місцем}, найнадійніший із відомих: у чутливих нейронів номер проводу каже, якої точки шкіри торкнулися чи якої висоти звук, і це встановлено багаторазово. Він дорогий за кількістю нейронів, зате швидкий: на повідомлення йде одна затримка.

\section{Час: швидко, але від чого відлічувати}

Другий вихід --- писати число в час імпульсу. Візьмімо нейрон~\eqref{eq:glutneuron} і подаймо на нього з моменту $t=0$ сталий вхід, який підняв би напругу до рівня $U$. Напруга зростає як $V=U\bigl(1-e^{-t/\tau}\bigr)$ і досягає порога $\theta$ у момент
\begin{empheq}[box=\keybox]{equation}
t_1 \;=\; \tau\,\ln\frac{U}{U-\theta}, \qquad U>\theta .
\label{eq:latency}
\end{empheq}
Сильний вхід --- ранній імпульс, слабкий --- пізній, підпороговий --- жодного. При $\tau=10\ms$ вхід $U=4\theta$ дає перший імпульс через $2.9\ms$, $U=2\theta$ --- через $6.9\ms$, $U=1.2\theta$ --- через $17.9\ms$. Один-єдиний імпульс несе число.

Але від якого моменту відлічувати $t_1$? Годинника немає, і отримувач не знає, коли почався подразник. Є три відповіді.

\emph{Відносна затримка.} Два нейрони бачать той самий подразник, і різниця їхніх затримок $t_1^A-t_1^B$ залежить лише від їхніх входів, а не від моменту початку. Порівнюють моменти детектори збігів із затримками (рис.~\ref{fig:itd}). Якщо $N$ нейронів видали по одному імпульсу, сам \emph{порядок} їхнього приходу несе до $\log_2N!$ біт: для десяти нейронів близько двадцяти двох біт, тоді як відповідь «спрацював чи ні» дала б не більше десяти. У сітківці показано, що за відносними затримками перших імпульсів її вихідних нейронів картинку можна відновити швидко й надійно~\cite{Gollisch2008}.

\emph{Залп як початок слова.} Різка зміна подразника змушує спрацювати майже одночасно безліч нейронів. Такий залп сам слугує позначкою початку, як довга пауза між словами в абетці Морзе, і отримувач відлічує затримки від нього.

\emph{Місцевий ритм.} У розділі~\ref{sec:gaba} петля \nt{GLUT}--\nt{GABA} породжувала місцевий ритм із періодом $12$--$50\ms$. Це не годинник, але всередині його циклу нейрон із сильним входом встигає спрацювати раніше, ніж нейрон зі слабким, і~\eqref{eq:latency} перетворюється на код \emph{фазою}: число записано в тому, в якій частині циклу прийшов імпульс. Такий код спостерігали: нейрони, що відповідають за певне місце в просторі, спрацьовують дедалі раніше в циклі повільного місцевого ритму в міру того, як тварина проходить через «своє» місце~\cite{OKeefe1993}. Наскільки широко мозок користується фазою, поки невідомо.

\section{Код обирає отримувач}

У послідовності імпульсів є і частота, і часи, і порядок. Що з цього буде прочитано, вирішує отримувач, і вирішує одним параметром --- своєю сталою часу $\tau$.

Усереднімо рівняння~\eqref{eq:glutneuron} за часом. Якщо на нейрон надходять незалежні потоки з частотами $r_j$, то середня напруга та її відносне тремтіння дорівнюють
\begin{empheq}[box=\keybox]{equation}
\bar V \;=\; \tau\sum_j w_j\,r_j ,
\qquad
\frac{\sigma_V}{\bar V} \;\approx\; \frac{1}{\sqrt{2\,r_{\text{вх}}\,\tau}} ,
\label{eq:reader}
\end{empheq}
де друга рівність записана для однакових ваг, а $r_{\text{вх}}$ --- сумарна частота всіх входів. Коли $\tau$ велика, напруга майже не тремтить і повторює зважену суму частот: отримувач --- \emph{інтегратор}, він читає частоти й забуває часи. Коли $\tau$ мала, напруга встигає впасти між імпульсами, і поріг досягається лише тоді, коли кілька імпульсів прийшли в межах вікна~\eqref{eq:window}: отримувач --- \emph{детектор збігів}, він читає часи (рис.~\ref{fig:reader}). Тисяча входів по п'ять імпульсів за секунду при $\tau=10\ms$ дають тремтіння близько десяти відсотків, майже чисте інтегрування. Якщо гальмування, як у розділі~\ref{sec:gaba}, вкоротило $\tau$ до $2\ms$, тремтіння зростає до двадцяти з гаком відсотків, і нейрон починає помічати збіги.

\begin{figure}[t]
\centering
\begin{tikzpicture}[>=Stealth,x=0.115cm,y=1cm]
% common input: the same spike train for both receivers
\foreach \s in {6,21,22,38,52,55.5,59,62.5,66,69.5,73,76.5,80,83.5,87,90.5,94,97.5}{\draw[thick,gray!40!black] (\s,5.25) -- (\s,5.65);}
\draw[gray!70!black] (0,5.25) -- (105,5.25);
\node[left,font=\small] at (-1,5.45) {вхід};
\node[above,font=\scriptsize,gray!40!black] at (13.5,5.65) {рідко};
\node[above,font=\scriptsize,gray!40!black] at (21.5,6.0) {пара};
\node[above,font=\scriptsize,gray!40!black] at (75,5.65) {часто};
% axes and thresholds
\foreach \y/\th in {2.7/2.5,0/1.5}{
  \draw[gray!70!black] (0,\y) -- (106,\y) node[right,black] {\small $t$};
  \draw[densely dashed,gray!60!black] (0,\y+0.6*\th) -- (104,\y+0.6*\th) node[right,black,font=\small] {$\theta$};
}
\node[anchor=south west,font=\small] at (0,4.3) {$\tau=10\ms$: інтегратор};
\node[anchor=south east,font=\small] at (104,1.0) {$\tau=2\ms$: детектор збігів};
% integrator: tau=10 ms, theta=2.5w
\begin{scope}[yshift=2.7cm]
\foreach \a/\b/\v in {0/6/0.000,6/21/1.000,21/22/1.223,22/38/2.107,38/52/1.425,52/55.5/1.351,55.5/59/1.952,59/62.5/2.376,62.5/66/0.000,66/69.5/1.000,69.5/73/1.705,73/76.5/2.201,76.5/80/0.000,80/83.5/1.000,83.5/87/1.705,87/90.5/2.201,90.5/94/0.000,94/97.5/1.000,97.5/104/1.705}{\draw[very thick,blue!60!black] plot[domain=\a:\b,samples=14] (\x,{0.6*\v*exp(-(\x-\a)/10)});}
\foreach \s/\p/\q in {6/0.000/1.000,21/0.223/1.223,22/1.107/2.107,38/0.425/1.425,52/0.351/1.351,55.5/0.952/1.952,59/1.376/2.376,62.5/1.674/2.500,62.5/2.500/0.000,66/0.000/1.000,69.5/0.705/1.705,73/1.201/2.201,76.5/1.551/2.500,76.5/2.500/0.000,80/0.000/1.000,83.5/0.705/1.705,87/1.201/2.201,90.5/1.551/2.500,90.5/2.500/0.000,94/0.000/1.000,97.5/0.705/1.705}{\draw[very thick,blue!60!black] (\s,0.6*\p) -- (\s,0.6*\q);}
\foreach \s in {62.5,76.5,90.5}{\draw[->,thick,red!70!black] (\s,1.58) -- (\s,1.95);}
\end{scope}
% coincidence: tau=2 ms, theta=1.5w
\begin{scope}[yshift=0cm]
\foreach \a/\b/\v in {0/6/0.000,6/21/1.000,21/22/1.001,22/38/0.000,38/52/1.000,52/55.5/1.001,55.5/59/1.174,59/62.5/1.204,62.5/66/1.209,66/69.5/1.210,69.5/73/1.210,73/76.5/1.210,76.5/80/1.210,80/83.5/1.210,83.5/87/1.210,87/90.5/1.210,90.5/94/1.210,94/97.5/1.210,97.5/104/1.210}{\draw[very thick,blue!60!black] plot[domain=\a:\b,samples=14] (\x,{0.6*\v*exp(-(\x-\a)/2)});}
\foreach \s/\p/\q in {6/0.000/1.000,21/0.001/1.001,22/0.607/1.500,22/1.500/0.000,38/0.000/1.000,52/0.001/1.001,55.5/0.174/1.174,59/0.204/1.204,62.5/0.209/1.209,66/0.210/1.210,69.5/0.210/1.210,73/0.210/1.210,76.5/0.210/1.210,80/0.210/1.210,83.5/0.210/1.210,87/0.210/1.210,90.5/0.210/1.210,94/0.210/1.210,97.5/0.210/1.210}{\draw[very thick,blue!60!black] (\s,0.6*\p) -- (\s,0.6*\q);}
\foreach \s in {22}{\draw[->,thick,red!70!black] (\s,0.98) -- (\s,1.35);}
\end{scope}
\foreach \x in {0,50,100}{\draw[gray!60!black] (\x,-0.06) -- (\x,0.06) node[below,font=\scriptsize] {\x\,мс};}
\end{tikzpicture}
\caption{Той самий вхід --- два різні отримувачі. Кожен імпульс піднімає напругу на $w$, після чого вона стікає, як у моделі нейрона~\eqref{eq:glutneuron}; червоні стрілки --- імпульси отримувача. Повільний отримувач ($\tau=10\ms$, $\theta=2.5w$) спрацьовує там, де імпульсів \emph{багато}, і не помічає пари. Швидкий ($\tau=2\ms$, $\theta=1.5w$) спрацьовує лише на пару в межах вікна~\eqref{eq:window}, а частий, але не збіжний потік пропускає.}
\label{fig:reader}
\end{figure}

Вимірювання показують, що в активній корі провідність мембрани зростає в кілька разів порівняно зі спокоєм~\cite{Destexhe2003}. Отже, $\tau$ там коротша, і нейрони кори в робочому режимі ближчі до детекторів збігів, ніж до інтеграторів.

\begin{keyblock}
\begin{proposition}[Код задає отримувач]
\label{prop:reader}
Та сама послідовність імпульсів для повільного отримувача несе частоту, для швидкого --- часи, для об'єму, у який викинуто медіатор, --- рівень концентрації, згладжений за секунди (розділ~\ref{sec:serocomputer}). Який код використовується, визначає не відправник, а стала часу отримувача, і її підлаштовує гальмування.
\end{proposition}
\end{keyblock}
\section{Крапки й тире: синапс як фільтр}

Досі синапс був у нас проводом із вагою. Насправді його відповідь залежить від попередніх імпульсів, і це дає мові нейронів другий знак.

\emph{Виснаження: синапс передає зміни.} Кожен викид витрачає частку $U$ запасу медіатора, готового до викиду~\cite{Tsodyks1997}, а запас відновлюється зі сталою часу $\tau_{\text{відн}}$ порядку сотень мілісекунд. За частоти $r$ амплітуда відповіді на імпульс установлюється приблизно на рівні $A(r)=A_0/(1+U r\tau_{\text{відн}})$, де $A_0$ --- відповідь відпочилого синапсу. Струм, який отримувач отримує за одиницю часу, дорівнює
\begin{empheq}[box=\keybox]{equation}
r\,A(r) \;=\; \frac{A_0\,r}{1+U r\,\tau_{\text{відн}}}
\;\xrightarrow[r\to\infty]{}\; \frac{A_0}{U\tau_{\text{відн}}} .
\label{eq:depression}
\end{empheq}
При $U=0.5$ і $\tau_{\text{відн}}=0.8\,\text{с}$ насичення починається вже з $1/(U\tau_{\text{відн}})=2.5$ імпульсу за секунду. Якщо частота зросла з п'яти до двадцяти, усталений струм зросте лише на третину. Зате перші імпульси на новій частоті приходять, поки запас ще колишній, і струм на мить підскакує вчетверо, а потім за $\tau_{\text{відн}}/(1+Ur\tau_{\text{відн}})\approx90\ms$ спадає (рис.~\ref{fig:depression}).

Такий синапс передає не частоту, а її \emph{зміну}, по суті --- відносну зміну $\Delta r/r$~\cite{Abbott1997depression}. Для електронника це фільтр верхніх частот, поставлений просто в провід.

\begin{figure}[t]
\centering
\begin{tikzpicture}[>=Stealth,x=12.5cm,y=1cm]
% input rate
\draw[->,gray!70!black] (0,2.6) -- (0.9,2.6) node[right,black] {\small $t$};
\draw[->,gray!70!black] (0,2.6) -- (0,4.3) node[above,black] {\small $r$};
\draw[very thick,blue!60!black] (0,2.9) -- (0.2,2.9) -- (0.2,3.8) -- (0.8,3.8);
\node[left,font=\scriptsize] at (0,2.9) {$5$};
\node[left,font=\scriptsize] at (0,3.8) {$20$};
\node[right,font=\small,blue!60!black] at (0.4,4.1) {частота на вході синапсу};
% output current
\draw[->,gray!70!black] (0,0) -- (0.9,0) node[right,black] {\small $t$};
\draw[->,gray!70!black] (0,0) -- (0,2.45) node[left,black] {\small $rA$};
\draw[densely dashed,gray!60!black] (0,0.667) -- (0.8,0.667);
\draw[very thick,red!70!black] (0,0.5) -- (0.2,0.5) -- (0.2,2.0)
   plot[domain=0.2:0.8,samples=60] (\x,{0.667+(2.0-0.667)*exp(-(\x-0.2)/0.089)});
\node[left,font=\scriptsize] at (0,0.5) {$1.7$};
\node[left,font=\scriptsize] at (0,2.0) {$6.7$};
\node[right,font=\scriptsize] at (0.8,0.667) {$2.2$};
\node[right,font=\small,red!70!black] at (0.28,1.6) {струм до отримувача: підскок і спад за $\approx90\ms$};
\draw[gray!60!black] (0.2,-0.08) -- (0.2,0.08); \node[below,font=\scriptsize] at (0.2,0) {$0.2\,\text{с}$};
\draw[gray!60!black] (0.8,-0.08) -- (0.8,0.08); \node[below,font=\scriptsize] at (0.8,0) {$0.8\,\text{с}$};
\end{tikzpicture}
\caption{Синапс, що виснажується, за формулою~\eqref{eq:depression} при $U=0.5$, $\tau_{\text{відн}}=0.8\,\text{с}$; струм в одиницях $A_0$ за секунду, штрихова лінія --- усталений струм: він зріс лише з $1.7$ до $2.2$, а в момент стрибка --- до $6.7$. Синапс повідомляє отримувачеві про те, що частота змінилася, а не про те, яка вона.}
\label{fig:depression}
\end{figure}

\emph{Полегшення: пачка як тире.} В інших синапсів імовірність викиду мала, але після кожного імпульсу на десятки мілісекунд зростає. Одиночний імпульс через такий синапс найчастіше губиться. А \emph{пачка} --- кілька імпульсів поспіль з інтервалами в кілька мілісекунд --- проходить майже напевно. Навіть без полегшення ймовірність, що пропадуть усі $k$ імпульсів пачки, дорівнює
\begin{empheq}[box=\keybox]{equation}
P_{\text{втрати}} \;=\; (1-p)^k ,
\label{eq:burst}
\end{empheq}
при $p=0.2$ це $0.8$ для одиночного імпульсу й $0.41$ для пачки з чотирьох; полегшення зменшує її ще більше. Так у нейрона з'являється другий знак: одиночний імпульс --- крапка, пачка --- тире. Синапс, що виснажується, найкраще передає перший імпульс після паузи; полегшуваний пропускає пачки й глушить одиночні крапки. Те, що пачки передаються надійніше, виміряно; що мозок справді користується ними як окремим знаком, --- правдоподібна гіпотеза~\cite{Lisman1997}.

\section{Як говорять \nt{GABA}-нейрони}

Найчисленніші з них у корі --- швидкі~\cite{Tremblay2016}. Їхні імпульси коротші, ніж у \nt{GLUT}-нейронів, вони можуть видавати їх рівно й довго з частотою в сотні за секунду, майже не втомлюючись. Їхні синапси надійніші: зв'язок із кожним отримувачем складається з багатьох місць викиду, й імпульс майже ніколи не губиться. Їхні виходи сідають поруч із тим місцем отримувача, де народжується його імпульс, тож вони керують не тим, як отримувач додає входи, а тим, чи спрацює він і коли.

Самі вони отримують збудження від багатьох сусідніх \nt{GLUT}-нейронів і повідомляють про сумарну активність навколо --- саме те, що потрібно для ділення з розділу~\ref{sec:gaba}. Нарешті, швидкі \nt{GABA}-нейрони зв'язані один з одним і легко спрацьовують разом; саме вони задають місцевий ритм, про який ішлося вище~\cite{Cardin2009}.

Мовою абетки Морзе \nt{GABA}-нейрони відповідають не за літери, а за \emph{паузи}. Вони ріжуть неперервний потік \nt{GLUT}-імпульсів на вікна, усередині яких отримувач може спрацювати, звужують вікна збігу й приглушують увесь потік, коли він стає надто гучним.

\begin{keyblock}
\begin{proposition}[Поділ ролей]
\label{prop:punctuation}
\nt{GLUT}-нейрони несуть зміст повідомлення: \emph{що} і \emph{скільки}. Швидкі \nt{GABA}-нейрони несуть його розмітку: \emph{коли} можна говорити й \emph{наскільки гучно}; ще один клас, гальмуючи гальмівних, вирішує, \emph{для кого} ворота відчинені. Окремі частини цього поділу виміряно; як загальне правило це робоча гіпотеза.
\end{proposition}
\end{keyblock}

Є й інші, повільніші \nt{GABA}-нейрони, виходи яких сідають на окремі входи отримувача. Вони працюють як заборона~\eqref{eq:veto} з розділу~\ref{sec:gaba}: закривають один потік \nt{GLUT}-імпульсів, не чіпаючи решти.

Поділ на швидкі й повільні \nt{GABA}-нейрони --- стара картина, і вона неповна. Зараз у корі розрізняють щонайменше три великі класи~\cite{Tremblay2016}, і третій влаштований несподівано: він гальмує не \nt{GLUT}-нейрони, а інші \nt{GABA}-нейрони, насамперед ті, що закривають входи~\cite{Pfeffer2013}. Гальмування гальмування --- подвійне заперечення. Такий нейрон \emph{відчиняє} ворота, не посилаючи отримувачеві жодного збуджувального імпульсу. Вмикають його сигнали з інших областей кори та широкомовні канали, тож він працює як дозвільний вхід цілої ділянки мережі: поки він активний, заборони знято, і потік проходить~\cite{Pi2013}. У додатку~\ref{sec:othermediators} цей прийом названо розгальмовуванням.
\section{Ощадлива мова}

Останнє, що відомо про мову нейронів твердо, --- вона дорога. Імпульс разом із роботою синапсів, яку він викликає, обходиться приблизно в $10^9$ молекул АТФ (оцінка для кори гризунів~\cite{Attwell2001}), а одна молекула дає близько $10^{-19}$ джоуля. Отже, імпульс коштує порядку $E_{\text{імп}}\sim10^{-10}$ джоуля. Увесь мозок споживає близько $20$ ватів, і якщо на передачу сигналів іде половина, $P\sim10$ ватів, то середня частота одного нейрона
\begin{empheq}[box=\keybox]{equation}
\bar r \;\approx\; \frac{P}{N\,E_{\text{імп}}}
\;\sim\; \frac{10\ \text{Вт}}{8.6\cdot10^{10}\cdot10^{-10}\ \text{Дж}}
\;\approx\; 1\ \text{імпульс за секунду}.
\label{eq:energy}
\end{empheq}
Докладніші оцінки для кори дають ще менше~\cite{Lennie2003}. Код мозку мусить бути \emph{розрідженим}: швидко працювати може лише мала частка нейронів.

З~\eqref{eq:spikeentropy} видно, що це не так уже й погано. За точності $1\ms$ імпульс нейрона, що працює з частотою $100$ за секунду, несе не більше п'яти біт, а імпульс нейрона, що працює раз на секунду, --- до одинадцяти. Якщо платити доводиться за імпульси, вигідно говорити рідко. Кора й справді міняє точність на енергію: за нестачі їжі нейрони зберігають частоту імпульсів, але витрачають менше на синаптичні струми, і їхні відповіді стають менш точними~\cite{Padamsey2022}.

В абетці Морзе часті літери короткі: найчастіша літера англійського тексту E --- одна крапка. Код нейронів, наскільки відомо, дотримується того самого правила в іншій формі: те, що очікуване, коштує мало імпульсів~\cite{Barlow1961}. Нейрон за сталого подразника поступово знижує частоту, датчики розділу~\ref{sec:glutcomputer} відповідають на зміни, а не на рівні, а \nt{DOPA}-нейрони розділу~\ref{sec:neurontypes} повідомляють лише про помилку передбачення винагороди.

\begin{keyblock}
\begin{proposition}[Ощадливість]
\label{prop:economy}
Енергія обмежує середню частоту нейрона величиною порядку одного імпульсу за секунду. Тому мова \nt{GLUT}-нейронів розріджена й витрачає імпульси на новини: на те, що змінилося або не було передбачене. Енергетичну межу виміряно; що код мозку підлаштований під найбільшу кількість біт на джоуль, --- добре обґрунтована гіпотеза.
\end{proposition}
\end{keyblock}

\section{Підсумок}

\begin{table}[t]
\centering
\footnotesize
\setlength{\tabcolsep}{4pt}
\renewcommand{\arraystretch}{1.15}
\begin{tabular}{>{\raggedright}p{2.6cm}>{\raggedright}p{2.3cm}>{\raggedright}p{3.0cm}>{\raggedright}p{2.0cm}>{\raggedright\arraybackslash}p{3.6cm}}
\toprule
Спосіб запису & Що несе & Хто читає & Час на повідомлення & Що відомо \\
\midrule
номер нейрона, що спрацював (код місцем) & яке значення & будь-який отримувач; детектори збігів & одна затримка & встановлено в чутливих нейронів і у визначенні напрямку на звук \\
частота одного нейрона & величину & інтегратор із великою $\tau$; об'єм (розд.~\ref{sec:serocomputer}) & сотні мс -- секунди & встановлено; для ступеня обробки в $10$--$20\ms$ надто повільно \\
частота групи & величину & нейрон із тисячами входів & $10$--$50\ms$ & встановлено; виграш обмежений кореляціями~\eqref{eq:correlated} \\
час першого імпульсу, порядок & величину, порядок & детектори збігів із затримками & одна затримка & встановлено на вході нервової системи; у корі обговорюється \\
фаза в місцевому ритмі & величину, положення & нейрони зі спільним ритмом & цикл $12$--$50\ms$ і повільніше & спостерігається в окремих областях; загальна роль --- гіпотеза \\
пачка («тире») & «важливо»: надійна доставка & полегшуваний синапс & $\sim10\ms$ & надійність виміряно; окремий знак --- гіпотеза \\
зміна частоти & новину & синапс, що виснажується & $\sim0.1\,\text{с}$ & встановлено \\
імпульси швидких \nt{GABA}-нейронів & коли й наскільки гучно & усі сусідні \nt{GLUT}-нейрони & $1$--$30\ms$ & ритм і ділення виміряно; «пунктуація» як правило --- гіпотеза \\
\bottomrule
\end{tabular}
\caption{Способи, якими \nt{GLUT}- і \nt{GABA}-нейрони записують повідомлення імпульсами. Правий стовпець відділяє виміряне від припущеного.}
\label{tab:code}
\end{table}

Таблиця~\ref{tab:code} збирає те, що ми знаємо; видно з неї й те, чого ми не знаємо. Ми не знаємо, наскільки кора користується точним часом імпульсів, а не лише частотами груп. Ми не знаємо, чи є в нейронів «слова» --- стійкі поєднання імпульсів багатьох нейронів, які читаються як ціле. І ми не знаємо, чи однакова мова в різних частинах мозку. Абетку Морзе можна вивчити за вечір, бо її таблиця відома. Таблицю мови нейронів ще належить скласти, і складатимуть її, найімовірніше, інженери зв'язку.

\section{Задачі}

\begin{exercise}[\lvlA]
\label{ex:04:1}
\emph{Пропускна здатність у числах.} $\Delta t=1\ms$. (а)~Скільки біт на джоуль дає нейрон із частотою $1$ імпульс за секунду за ціни імпульсу з~\eqref{eq:energy}? (б)~При $r=10$ у скільки разів менше за межу~\eqref{eq:spikeentropy} видобуває отримувач, який лише рахує імпульси за секунду?
\end{exercise}

\begin{exercise}[\lvlB]
\label{ex:04:2}
\emph{Оптимальна частота.} Нейрон витрачає потужність $P_0+rE_{\text{імп}}$, де $P_0$ --- друга половина з $20$~Вт, поділена на всі нейрони, $E_{\text{імп}}=10^{-10}$~Дж. За якої частоти $r$ кількість біт на джоуль $R_{\max}(r)/(P_0+rE_{\text{імп}})$ за~\eqref{eq:spikeentropy} з $\Delta t=1\ms$ найбільша?
\end{exercise}

\begin{exercise}[\lvlB]
\label{ex:04:3}
\emph{Яка кореляція шкідлива.} $N=1000$ нейронів, дисперсія $\sigma^2=1$, попарна кореляція $c=0.1$. Обчисліть інформацію Фішера $\mathcal I=f'^{\mathsf T}\Sigma^{-1}f'$ ($\Sigma$ --- матриця коваріацій) для однакових нахилів $f'=\mathbf 1$ і для нахилів $\pm1$ з $\mathbf 1^{\mathsf T}f'=0$. У скільки разів вони різняться?
\end{exercise}

\begin{exercise}[\lvlB]
\label{ex:04:4}
\emph{Моделювання: детектор збігів.} Нейрон~\eqref{eq:glutneuron} отримує $1000$ пуассонівських входів по $5$ імпульсів за секунду, $w=1$; поріг $\theta$ такий, що вільна напруга перетинає його раз на секунду. Моделлю знайдіть, скільки одночасних входів $m$ запалюють нейрон з імовірністю $1/2$ при $\tau=10$ і $2\ms$.
\end{exercise}

\begin{exercise}[\lvlB]
\label{ex:04:5}
\emph{Проєктування: детектор відносних змін.} Доберіть синапс~\eqref{eq:depression}: усталений струм при $40$ імпульсах за секунду не більш ніж на $10\,\%$ вищий, ніж при $10$, а після стрибка до $40$ струм виходить на новий рівень зі сталою часу не більше $50\ms$. Яке найменше $\tau_{\text{відн}}$ і за якого $U$?
\end{exercise}

\begin{exercise}[\lvlA]
\label{ex:04:6}
\emph{Час першого імпульсу й пачки.} (а)~За~\eqref{eq:latency} знайдіть вхід, за якого перший імпульс приходить рівно через одну сталу часу. Від чого він залежить? (б)~За ймовірності викиду $p=0.2$ скільки імпульсів потрібно в пачці, щоб утратити всі з імовірністю нижче $5\,\%$?
\end{exercise}

\begin{exercise}[\lvlC]
\label{ex:04:7}
\emph{Дослідження: скільки біт у розташуванні.} Нейрон --- незалежні клітинки~\eqref{eq:spikeentropy}, $r=10$, $\Delta t=1\ms$. (а)~Розкладіть ентропію «слова» з $20$ клітинок на ентропію кількості імпульсів і залишок. (б)~Змоделюйте оцінку ентропії за частотами слів із $N=30$ і $10^4$ записів. Наскільки вона занижена?
\end{exercise}

\chapter{Що обчислюють \nt{SERO}-нейрони}
\label{sec:serocomputer}
\begin{chaptergoals}
\item як ритмоводій із гальмівним рецептором дає НЕ та АБО-НЕ з окремих \nt{SERO}-нейронів;
\item чому викид в об'єм лишає лише кілька незалежних проводів, заданих $\ell\sim\sqrt{D\tau}$;
\item як концентрація згладжує імпульси в рівень, а від'ємний зворотний зв'язок міг би його тримати;
\item чому \nt{SERO} --- кандидат на роль горизонту планування і як горизонт задає терпіння.
\end{chaptergoals}

\section{Перший широкомовний канал}

Досі всі нейрони книги говорили адресною шиною: \nt{GLUT} і \nt{GABA} надсилали імпульси певним отримувачам, і ми з'ясували, що обчислює така мережа і якою мовою вона говорить. Тепер переходимо до другої шини з підрозділу~\ref{sec:buses} --- широкомовної. Її перший канал --- \nt{SERO}. Як і раніше, дозволимо собі лише те, що буває в живому організмі.

У цьому розділі з'являться три пристрої, якими згодом користуватимуться всі широкомовні канали: нейрон, який за рівного фонового притоку працює сам, доки його не зупинять; провід, який насправді є об'ємом тканини; і повільна концентрація замість окремих імпульсів. \nt{DOPA}, \nt{NORA} і \nt{ACET} у розділах~\ref{sec:dofa}--\ref{sec:acet} буде зібрано з тих самих деталей.

З погляду інженера \nt{SERO}-нейрон має чотири важливі відмінності.

\emph{Перша: знак вибирає отримувач.} Серотонін має рецептори обох знаків, і правило~\eqref{eq:dalesign} тут не діє: знак визначає не відправник, а рецептор отримувача (твердження~\ref{prop:receptor})~\cite{BarnesSharp1999}.

\emph{Друга: \nt{SERO}-нейрон працює сам, якщо є рівний фон.} У стані неспання він рівно, кілька разів на секунду, видає імпульси, і поштовх на кожен імпульс йому не потрібен: це ритмоводій. Але йому потрібен сталий фоновий приток. У зрізі мозку, де цього притоку немає, більшість таких клітин мовчить; рівний ритм повертається, якщо подати норадреналін через $\alpha_1$-рецептори, а їхня блокада його гасить~\cite{VandermaelenAghajanian1983}. В організмі цей фон надходить від \nt{NORA}. Для схеми це генератор, якому потрібне живлення: поки воно є, він працює сам, доки його не зупинить гальмування.

\emph{Третя: він повільний.} Майже всі рецептори серотоніну метаботропні: вони не відкривають канал самі, а запускають у клітині ланцюжок реакцій. Відповідь повільна: гальмівний струм через головний рецептор наростає і спадає в межах приблизно секунди~\cite{CourtneyFord2016} --- у десятки разів довше, ніж відповідь швидкого синапсу \nt{GLUT}.

\emph{Четверта: він викидає медіатор не в провід, а в об'єм.} В областях, куди йдуть виходи \nt{SERO}-нейронів, серотонін часто виходить не у вузьку щілину, а в міжклітинний простір і розпливається навколо~\cite{Bunin1998}. Отримувач відчуває концентрацію і не може знати, від якого відправника прийшла молекула.

За таких часів окремі імпульси вже неважливі, тому описуватимемо мережу частотами $r_j$ і концентраціями. Концентрація $c$ серотоніну в об'ємі, куди викидають медіатор нейрони $j$, зростає від викиду і спадає, бо медіатор відкачують назад:
\begin{empheq}[box=\keybox]{equation}
\tau_c\,\frac{\dd c}{\dd t} \;=\; -c \;+\; \kappa\sum_j r_j ,
\qquad
r_i \;=\; f_i\bigl(b_i + s_i\,g\,c\bigr), \quad s_i=\pm1 .
\label{eq:seroneuron}
\end{empheq}
Це ще один спосіб прочитати потік імпульсів, що доповнює твердження~\ref{prop:reader}: об'єм не бачить ні окремих імпульсів, ні їхнього порядку, а лише частоту, усереднену за час $\tau_c$. Перше рівняння --- те саме RC-коло, що й мембрана: $\tau_c$ --- час життя медіатора в об'ємі; прямо його не виміряно, і ми вважаємо його порядку секунди, $\kappa$ --- скільки медіатора дає один імпульс. Друге описує отримувача: $b_i$ --- його власний приток, у ритмоводія додатний (до нього входить і рівний фоновий вхід від \nt{NORA}); $g$ --- чутливість рецептора; знак $s_i$ вибирає отримувач; $f_i$ --- неспадна передавальна характеристика.

\section{НЕ одним нейроном}

Візьмімо ритмоводій із гальмівним рецептором, $s=-1$. Поки серотоніну навколо немає, він працює на власному притоці $b$. Коли серотонін з'являється, приток зменшується на $gc$, і якщо $gc>b$, нейрон замовкає. Це НЕ: вихід є, коли входу немає. На нього пішов один нейрон (рис.~\ref{fig:serogates}). Твердження~\ref{prop:monotone} тут незастосовне: воно вимагало додатних ваг.

Якщо на той самий ритмоводій діє серотонін із двох різних об'ємів, він мовчить, коли серотонін є хоча б в одному з них. Це АБО-НЕ, а з АБО-НЕ можна зібрати будь-яку логічну функцію. Двопровідний код \nt{GLUT}-мережі тут не потрібен: відсутність сигналу обчислюють нейрони, які працюють, доки їх не зупинять.

\begin{figure}[t]
\centering
\begin{tikzpicture}[>=Stealth,
  nrn/.style={circle,draw,very thick,fill=blue!6,minimum size=0.95cm,inner sep=0pt},
  pace/.style={nrn,double,double distance=1.2pt},
  inh/.style={-{Bar[width=7pt]},thick,red!70!black}]
\node[pace] (N1) at (0,0) {};
\draw[inh] (-1.6,0) node[left,black] {$c$} -- (N1);
\draw[->,very thick,red!70!black] (N1) -- (1.3,0) node[right,black] {$\bar c$};
\node[below=0.55cm] at (0,0) {\small НЕ};
\node[pace] (N2) at (5,0) {};
\draw[inh] (3.4,0.45) node[left,black] {$c_1$} -- (N2);
\draw[inh] (3.4,-0.45) node[left,black] {$c_2$} -- (N2);
\draw[->,very thick,red!70!black] (N2) -- (6.3,0) node[right,black] {$\overline{c_1\vee c_2}$};
\node[below=0.55cm] at (5,0) {\small АБО-НЕ};
\end{tikzpicture}
\caption{Логіка з \nt{SERO}-нейронів. Подвійний кружок --- ритмоводій: за рівного фонового притоку він працює сам, доки його не загальмують. Лінія з рискою --- гальмівний рецептор, $s=-1$; $c$, $c_1$, $c_2$ --- концентрації серотоніну в об'ємах навколо отримувача. Ліворуч --- НЕ, праворуч --- АБО-НЕ.}
\label{fig:serogates}
\end{figure}

\begin{keyblock}
\begin{proposition}[Повнота \nt{SERO}-логіки]
\label{prop:serocomplete}
Якщо в мережі є ритмоводії з гальмівними рецепторами і викиди в різні об'єми не змішуються, мережа~\eqref{eq:seroneuron} може обчислити будь-яку логічну функцію, і кожен біт передається одним проводом.
\end{proposition}
\end{keyblock}

Логічно \nt{SERO}-мережа сильніша за \nt{GLUT}, але умова «викиди в різні об'єми не змішуються» в мозку не виконується (підрозділ~\ref{sec:volume}).

\section{Від'ємний зворотний зв'язок}

Серотонінові нейрони мають гальмівні рецептори на собі самих~\cite{Sprouse1987}. Серотонін, який вони викинули, повертається до них же і пригальмовує їх. Для інженера це знайома схема: підсилювач із від'ємним зворотним зв'язком (рис.~\ref{fig:serofeedback}).

Що тут виміряно, а що --- модель. У самому ядрі шва, де розташовані \nt{SERO}-нейрони, серотонін гальмує сусідів не через спільний об'єм, а точка-точка, через синапси: переносник швидко прибирає його і не дає накопичитися поза щілиною. Поодинокий викид дає лише коротку паузу в розряді, а середня частота від цього не змінюється~\cite{CourtneyFord2016}. Тому петля нижче, замкнена через спільний об'єм, --- модель: ми рахуємо, що робив би такий зворотний зв'язок, якби він працював на рівні середніх частот.

\begin{figure}[t]
\centering
\begin{tikzpicture}[>=Stealth,
  nrn/.style={circle,draw,very thick,fill=blue!6,minimum size=0.95cm,inner sep=0pt},
  pace/.style={nrn,double,double distance=1.2pt},
  blk/.style={draw,very thick,rounded corners=2pt,fill=orange!10,minimum height=0.9cm,minimum width=2.2cm,align=center,font=\small},
  inh/.style={-{Bar[width=7pt]},thick,red!70!black}]
\node[pace] (N) at (0,0) {};
\node[blk] (V) at (3.6,0) {об'єм\\$\tau_c$};
\draw[->,very thick] (N) -- node[above] {\small $r$} (V);
\draw[->,very thick] (V) -- (6.0,0) node[right] {\small $c$ --- усім отримувачам};
\draw[inh] (V.south) -- ++(0,-0.9) -| node[pos=0.25,below] {\small $-g\,c$} (N.south);
\draw[->,thick,blue!60!black] (-1.7,0) node[left,black] {\small $b$} -- (N);
\end{tikzpicture}
\caption{\nt{SERO}-нейрон зі зворотним зв'язком через власний медіатор: концентрація $c$ іде до всіх отримувачів і гальмує сам нейрон. У моделі це регулятор рівня; у мозку така роль петлі --- гіпотеза.}
\label{fig:serofeedback}
\end{figure}

Порахуймо, що робить ця петля. Нехай передавальна характеристика лінійна, $f(u)=au$ при $u>0$. У рівновазі $c=\kappa r$ і $r=a(b-gc)$. Підставивши одне в інше, отримаємо
\begin{empheq}[box=\keybox]{equation}
r \;=\; \frac{a\,b}{1+G}, \qquad G \;=\; a\,g\,\kappa .
\label{eq:feedback}
\end{empheq}
Величина $G$ --- підсилення петлі. Це та сама формула, що й для будь-якого підсилювача з від'ємним зворотним зв'язком, і вона дає ті самі два виграші.

\emph{Стійкість до розкиду.} Нехай збудливість нейрона $a$ змінилася на частку $\varepsilon$. Без зворотного зв'язку на ту саму частку змінилася б і частота. Зі зворотним зв'язком, при $G\gg1$, частота $r\approx b/(g\kappa)$ майже не залежить від $a$: її зміна в $1+G$ разів менша. При $G=9$ розкид пригнічується вдесятеро.

\emph{Швидкодія.} Підставмо $r=a(b-gc)$ у перше рівняння~\eqref{eq:seroneuron}: отримаємо $\tau_c\,\dd c/\dd t=-(1+G)\,c+\kappa ab$. Концентрація наближається до свого рівня зі сталою часу $\tau_c/(1+G)$ --- в $1+G$ разів швидше, ніж без зворотного зв'язку.

Ось перше обчислення, яке могла б виконувати \nt{SERO}-система: тримати загальний рівень серотоніну в мозку на заданій позначці, як термостат тримає температуру. Це гіпотеза: у досліді автогальмування поки видно лише як короткі паузи, що не змінюють середньої частоти.

\section{Від імпульсів --- до рівня}

Друге обчислення сховане в першому рівнянні~\eqref{eq:seroneuron}. Нейрони видають окремі імпульси, а отримувач відчуває концентрацію, яка згладжує їх за час $\tau_c$. Якщо частота джерела змінюється, концентрація йде за нею із запізненням:
\begin{equation}
c(t) \;=\; \frac{\kappa}{\tau_c}\int_{-\infty}^{t} e^{-(t-t')/\tau_c}\,\sum_j r_j(t')\,\dd t' .
\label{eq:lowpass}
\end{equation}
Це ковзне середнє активності джерел за останні кілька $\tau_c$ --- фільтр нижніх частот (рис.~\ref{fig:lowpass}). Так найпростіший цифро-аналоговий перетворювач пропускає імпульси через RC-коло і перетворює їхню частоту на рівень. \nt{SERO}-система робить те саме із сотнями тисяч нейронів одночасно.

Ця величина водночас і пам'ять. Після того як джерела замовкли, концентрація тримається ще час порядку $\tau_c$. Це не біт, а слід, що плавно тане.

\begin{figure}[t]
\centering
\begin{tikzpicture}[>=Stealth,x=1.45cm,y=1cm]
\draw[gray!70!black] (0,3.2) -- (8.1,3.2);
\node[left,font=\small] at (-0.05,3.4) {імпульси};
\node[above,font=\scriptsize,gray!40!black] at (1,3.65) {$2$ на секунду};
\node[above,font=\scriptsize,gray!40!black] at (3.5,3.65) {$8$ на секунду};
\node[above,font=\scriptsize,gray!40!black] at (6.5,3.65) {$2$ на секунду};
\draw[->,gray!70!black] (0,0) -- (8.3,0) node[right,black] {\small $t$};
\draw[->,gray!70!black] (0,0) -- (0,2.75) node[left,black] {\small $c$};
\foreach \s in {0.136,0.529,0.760,1.223,1.714,1.748,1.755,2.663,2.701,2.734,3.414,3.493,3.719,3.800,3.928,3.948,4.074,4.327,4.420,4.589,4.728,4.736,4.914,5.025,5.205,5.220,6.224,6.544,7.178}{\draw[thick,blue!60!black] (\s,3.2) -- (\s,3.6);}
\foreach \a/\b/\v in {0.000/0.136/0.000,0.136/0.529/1.000,0.529/0.760/1.675,0.760/1.223/2.330,1.223/1.714/2.466,1.714/1.748/2.509,1.748/1.755/3.425,1.755/2.663/4.402,2.663/2.701/2.775,2.701/2.734/3.672,2.734/3.414/4.553,3.414/3.493/3.306,3.493/3.719/4.055,3.719/3.800/4.235,3.800/3.928/4.905,3.928/3.948/5.316,3.948/4.074/6.211,4.074/4.327/6.476,4.327/4.420/6.028,4.420/4.589/6.493,4.589/4.728/6.483,4.728/4.736/6.642,4.736/4.914/7.589,4.914/5.025/7.351,5.025/5.205/7.579,5.205/5.220/7.331,5.220/6.224/8.221,6.224/6.544/4.012,6.544/7.178/3.914,7.178/8.000/3.076}{\draw[very thick,orange!85!black] plot[domain=\a:\b,samples=10] (\x,{0.25*\v*exp(-(\x-\a))});}
\foreach \s/\p/\q in {0.136/0.000/1.000,0.529/0.675/1.675,0.760/1.330/2.330,1.223/1.466/2.466,1.714/1.509/2.509,1.748/2.425/3.425,1.755/3.402/4.402,2.663/1.775/2.775,2.701/2.672/3.672,2.734/3.553/4.553,3.414/2.306/3.306,3.493/3.055/4.055,3.719/3.235/4.235,3.800/3.905/4.905,3.928/4.316/5.316,3.948/5.211/6.211,4.074/5.476/6.476,4.327/5.028/6.028,4.420/5.493/6.493,4.589/5.483/6.483,4.728/5.642/6.642,4.736/6.589/7.589,4.914/6.351/7.351,5.025/6.579/7.579,5.205/6.331/7.331,5.220/7.221/8.221,6.224/3.012/4.012,6.544/2.914/3.914,7.178/2.076/3.076}{\draw[very thick,orange!85!black] (\s,0.25*\p) -- (\s,0.25*\q);}
\draw[thick,densely dashed,black] plot[domain=0:2,samples=30] (\x,{0.25*2*(1-exp(-\x))})
  plot[domain=2:5,samples=40] (\x,{0.25*(8-6.271*exp(-(\x-2)))})
  plot[domain=5:8,samples=40] (\x,{0.25*(2+5.688*exp(-(\x-5)))});
\foreach \x in {0,2,5,8}{\draw[gray!60!black] (\x,-0.05) -- (\x,0.05) node[below=3pt,font=\scriptsize] {\x\,с};}
\end{tikzpicture}
\caption{Від імпульсів --- до рівня. Угорі --- імпульси \nt{SERO}-нейронів: дві секунди рідко, три секунди часто, потім знову рідко. Унизу --- концентрація серотоніну~\eqref{eq:lowpass} при $\tau_c=1\,\text{с}$. Штрихова лінія --- те саме, якби імпульси йшли ідеально рівно.}
\label{fig:lowpass}
\end{figure}

\section{Провід --- це об'єм}
\label{sec:volume}

Повернімося до умови твердження~\ref{prop:serocomplete}. Серотонін, викинутий поблизу різними відправниками, складається в одну концентрацію.

\emph{Провід тут --- не волокно, а об'єм тканини.} Два сигнали лишаються різними, тільки якщо їх викидають в області, рознесені далі, ніж встигає розпливтися медіатор. За час $\tau$ молекула з коефіцієнтом дифузії $D$ відходить на
\begin{empheq}[box=\keybox]{equation}
\ell \;\sim\; \sqrt{D\,\tau}\, .
\label{eq:diffusionlength}
\end{empheq}
Звивистість міжклітинних щілин сповільнює дифузію невеликої молекули приблизно в $2.6$ раза~\cite{Sykova2008}, а коефіцієнт дифузії самого серотоніну в тканині не виміряно; за розміром молекули оцінимо його в кілька одиниць на $10^{-6}$~см$^2$/с. Якщо сигнал живе $\tau\sim1$~с (теж оцінка), то $\ell\sim\sqrt{3\cdot10^{-6}}$~см $\approx20$~мкм. Адреса в такій мережі --- це \emph{місце}: отримувач дізнається, від кого сигнал, лише з того, де він сам розташований.

Справжні \nt{SERO}-нейрони влаштовані так, що цих окремих адрес майже не лишається. Вихід кожного з них розкиданий по величезних ділянках мозку: гілки однієї клітини сягають багатьох далеких одна від одної областей~\cite{GagnonParent2014}. Місць викиду на клітину, за оцінкою, близько $10^5$ (таблиця~\ref{tab:types}); прямо їх не підраховано. «Проводи» різних \nt{SERO}-нейронів перекриваються і змішуються. Проте зовсім в один провід вони не зливаються: \nt{SERO}-нейрони поділяються на кілька підсистем, які надсилають вихід у різні області і по-різному відповідають на ту саму подію~\cite{Ren2018}. Але таких підсистем одиниці, а не сотні тисяч.

\begin{keyblock}
\begin{proposition}[Кількість незалежних проводів]
\label{prop:volumewires}
У мережі з об'ємною передачею кількість незалежних сигналів обмежена не кількістю нейронів, а кількістю неперетинних областей розміром $\ell\sim\sqrt{D\tau}$, у які вони викидають медіатор. Якщо вихід кожного нейрона розкиданий по великому об'єму, уся мережа передає лише кілька спільних змінних.
\end{proposition}
\end{keyblock}

Тому логічні схеми твердження~\ref{prop:serocomplete} в мозку нема з чого зібрати. А регуляторові~\eqref{eq:feedback} і фільтру~\eqref{eq:lowpass} окремі проводи не потрібні: їм досить однієї спільної величини. Фільтр прямо випливає з виміряного викиду в об'єм в областях призначення~\cite{Bunin1998}; регулятор рівня --- гіпотеза.

\section{Горизонт планування}
\label{sec:horizon}

Для чого може знадобитися одна спільна повільна величина? Добрий кандидат --- параметр, який має бути однаковим для всієї мережі і змінюватися не швидше, ніж за секунди чи хвилини. У розділі~\ref{sec:neurontypes} такий параметр уже з'являвся: коефіцієнт $\gamma$ у формулі помилки~\eqref{eq:rpe}, який показує, наскільки майбутня винагорода дешевша за негайну. У неперервному часі його місце займає \emph{горизонт} $\tau_\gamma$: винагорода $R$, на яку треба чекати час $D$, коштує зараз $R\,e^{-D/\tau_\gamma}$.

Горизонт вирішує, чи варто чекати. Нехай можна одразу взяти малу винагороду $r_0$ або дочекатися великої $R$. Чекати вигідно, доки
\begin{empheq}[box=\keybox]{equation}
D \;<\; \tau_\gamma\,\ln\frac{R}{r_0} .
\label{eq:patience}
\end{empheq}
При $\tau_\gamma=2\,\text{с}$ і втричі більшій відкладеній винагороді варто чекати до $2.2\,\text{с}$; якщо горизонт удвічі довший, --- до $4.4\,\text{с}$. Терпіння пропорційне горизонту.

Формула~\eqref{eq:patience} спирається на експоненційне знецінення. Виміряне знецінення на затримках у секунди ближче до гіперболи $R/(1+D/\tau_\gamma)$: так у мавп спадають і цінність відкладеної винагороди у виборі, і відповідь \nt{DOPA}-нейронів на її провісник~\cite{KobayashiSchultz2008}. Для гіперболи умова~\eqref{eq:patience} замінюється на $D<\tau_\gamma\,(R/r_0-1)$: залежність від відношення винагород уже не логарифмічна, а лінійна, і за тих самих чисел межа очікування зсувається майже вдвічі. Гіпербола виходить із тієї самої експоненти, якщо затримка випадкова (задача~\ref{ex:05:7}), а пропорційність терпіння масштабу часу зберігається.

За гіпотезою, горизонт і задає \nt{SERO}~\cite{Doya2002}. Для такої ролі в нього є все: один спільний рівень на всю мережу, що змінюється за секунди. Є й прямі досліди: якщо штучно ввімкнути серотонінові нейрони, тварина довше чекає на відкладену винагороду, перш ніж облишити очікування~\cite{Miyazaki2014}, і довше намагається здобути винагороду там, де вона стала рідкісною~\cite{Lottem2018}. Як саме концентрація серотоніну перетворюється на $\tau_\gamma$ усередині мережі, невідомо. У наступному розділі горизонт увійде в помилку, яку розсилає \nt{DOPA}.

\section{Підсумок}

\begin{table}[t]
\centering
\small
\begin{tabular}{>{\raggedright}p{3.6cm}>{\raggedright}p{4.3cm}>{\raggedright\arraybackslash}p{4.3cm}}
\toprule
& \nt{GLUT} & \nt{SERO} \\
\midrule
час реакції & мілісекунди & частки секунди --- секунда \\
знак дії & лише $+$ & $+$ і $-$, вибирає отримувач \\
без входу & мовчить & працює сам за рівного фонового притоку \\
адресація & точка --- точка & об'єм, адреса --- місце \\
НЕ & лише через датчики «вимкнення» & один нейрон \\
пам'ять & самопідтримна активність, гасне від утоми & концентрація, тане за час $\tau_c$ \\
головні обчислення & збіги, затримки, логіка, послідовності & згладжування; регулювання рівня і горизонт $\tau_\gamma$ (гіпотези) \\
\bottomrule
\end{tabular}
\caption{Що обчислюють мережі лише з \nt{GLUT}-нейронів і лише з \nt{SERO}-нейронів.}
\label{tab:glutvssero}
\end{table}

Як логічний елемент (таблиця~\ref{tab:glutvssero}) \nt{SERO}-нейрон багатший за \nt{GLUT}, але як система зв'язку він --- широкомовна розсилка: повільна і майже без адрес. Тому він не намагається бути процесором, а згладжує і розсилає кілька спільних величин, про які йшлося в розділі~\ref{sec:neurontypes}, --- за гіпотезою, зокрема й горизонт планування. Ритмоводій, об'єм і повільна концентрація трапляться нам ще тричі: у \nt{DOPA}, \nt{NORA} і \nt{ACET}.

\section{Задачі}

\begin{exercise}[\lvlA]
\label{ex:05:1}
\emph{Провід --- це об'єм.} Візьміть для серотоніну $D=3\cdot10^{-6}$~см$^2$/с (підрозділ~\ref{sec:volume}). Знайдіть за~\eqref{eq:diffusionlength} $\ell$ для сигналу, що живе $1\,\text{с}$, і час, за який медіатор розпливеться на $5$~см. Чим тоді забезпечено широкомовність \nt{SERO} --- дифузією чи розгалуженням проводу з $\sim10^5$ місць викиду?
\end{exercise}

\begin{exercise}[\lvlA]
\label{ex:05:2}
\emph{Регулятор рівня.} Покажіть, що в~\eqref{eq:seroneuron} при $f(u)=au$ відносна чутливість $S=(a/r)\,\partial r/\partial a$ дорівнює $1/(1+G)$ точно, а не лише при $G\gg1$. При $G=9$ збудливість $a$ зросла на $20\,\%$. На скільки відсотків зміниться $r$ без лінеаризації?
\end{exercise}

\begin{exercise}[\lvlB]
\label{ex:05:3}
\emph{Повільний рецептор у петлі.} Рецептор \nt{SERO} відповідає із затримкою: $r(t)=a\bigl(b-gc(t-T)\bigr)$, об'єм --- за~\eqref{eq:seroneuron}. Покажіть, що при $G>1$ петля стійка лише при $T<T^*=\tau_c\arccos(-1/G)/\sqrt{G^2-1}$. Знайдіть $T^*$ при $\tau_c=1\,\text{с}$ для $G=2$ і $G=9$. Яке пригнічення розкиду досяжне за затримки в сотні мілісекунд?
\end{exercise}

\begin{exercise}[\lvlB]
\label{ex:05:4}
\emph{Дробовий шум рівня.} Кожен імпульс додає до $c$ за~\eqref{eq:lowpass} експоненту з $\tau_c=1\,\text{с}$. Знайдіть коефіцієнт варіації $c$, якщо в один об'єм викидають пуассонівські імпульси всі $3\cdot10^5$ \nt{SERO}-нейронів по $2$ на секунду. Яке $\tau_c$ дало б $\mathrm{CV}=10\,\%$ одному нейрону з частотою $4$ імпульси на секунду?
\end{exercise}

\begin{exercise}[\lvlB]
\label{ex:05:5}
\emph{Моделювання: зворотний зв'язок проти шуму.} Змоделюйте $N=50$ пуассонівських ритмоводіїв із частотами $r_i=a_i[b-gc]_+$, $a_i$ рівномірні в $[2.8,\,5.2]$, і об'єм~\eqref{eq:seroneuron} з $\kappa=1/N$, $\tau_c=1\,\text{с}$. У скільки разів зворотний зв'язок ($b=1$, $g=2.25$) знижує CV рівня $c$ порівняно з $b=0.1$, $g=0$? Чи вирівнює він частоти нейронів?
\end{exercise}

\begin{exercise}[\lvlB]
\label{ex:05:6}
\emph{Проєктування: XOR на \nt{SERO}-логіці.} Вентиль --- ритмоводій, що видає $r_0$, якщо $b-g\sum_kc_k>0$, і $0$ інакше, у свій об'єм $\tau_c\,\dd c/\dd t=-c+\kappa r$; він обчислює АБО-НЕ. Зберіть XOR із п'яти таких вентилів і знайдіть час його встановлення при $\tau_c=1\,\text{с}$ і $G'=g\kappa r_0/b=2$.
\end{exercise}

\begin{exercise}[\lvlB]
\label{ex:05:7}
\emph{Терпіння за невідомої затримки.} (а)~Тварина погоджується чекати втричі більшу винагороду, якщо очікування не довше за $6\,\text{с}$. Який горизонт $\tau_\gamma$ це означає? (б)~Нехай затримка $D$ розподілена експоненційно із середнім $\bar D$. Покажіть, що чекати вигідно при $\bar D<\tau_\gamma(R/r_0-1)$, і порівняйте з фіксованою затримкою при $\tau_\gamma=2\,\text{с}$, $R=3r_0$.
\end{exercise}

\begin{exercise}[\lvlC]
\label{ex:05:8}
\emph{Як концентрація задає горизонт.} \nt{DOPA}-нейрон отримує $V$ напряму з вагою $k_1$ і через \nt{GABA}-посередника, що згладжує з $\tau_d=0.1\,\text{с}$, з вагою $-k_2$; для повільних $V$ сума дорівнює $(k_1-k_2)V+k_2\tau_d\dot V$. Доберіть $k_1$, $k_2$ під~\eqref{eq:ctd} з $\tau_\gamma=2\,\text{с}$. \nt{SERO} множить $k_1$ на $m$: наскільки треба змінити $m$, щоб подвоїти горизонт?
\end{exercise}

\chapter{Нейронні мережі, які вміють навчатися: додаємо \nt{DOPA}}
\chaptermark{Мережі, які навчаються}
\label{sec:dofa}
\begin{chaptergoals}
\item чому синапсу доступні лише локальні сигнали і чому це виключає зворотне поширення помилки;
\item чого навчає правило Гебба: напрямку, в якому входи змінюються найсильніше;
\item як слід і широкомовний сигнал $\delta$ перетворюють шум на кроки вгору градієнтом винагороди;
\item як критик дає $\delta$ у неперервному часі і чому один спільний провід навчає повільно.
\end{chaptergoals}

\section{Правила гри}

В усіх схемах попередніх розділів ваги ставив інженер. В організмі інженера немає. Ваги мають установлюватися самі, під час роботи, і так, щоб мережа робила щось корисне. Це й називають навчанням.

До попередніх правил додається ще одне. Синапс змінює свою вагу сам, і знати він може лише те, що відбувається \emph{поруч із ним}: активність відправника $x$, активність отримувача $y$ і ті речовини, що до нього доходять. Ніхто не може надіслати синапсу персональну поправку.

Це виключає головний метод штучних нейронних мереж --- зворотне поширення помилки. Там помилка виходу проходить мережу у зворотному напрямку тими самими вагами, в окрему фазу після прямого проходу, і кожна вага отримує свою поправку. Для цього потрібні зворотні проводи, що точно повторюють прямі, і спільний такт. Ні того, ні іншого в мозку не знайдено, принаймні в тій формі, що в штучних мережах~\cite{Lillicrap2020,WhittingtonBogacz2019}.

Зміну ваги записуватимемо як повільний неперервний процес:
\begin{equation}
\frac{\dd w}{\dd t} \;=\; \eta\,\Phi(\text{те, що відомо синапсу}),
\label{eq:plasticity}
\end{equation}
де $\eta$ --- швидкість навчання, настільки мала, що за час одного імпульсу вага майже не змінюється. Увесь розділ --- про те, якою може бути функція $\Phi$.

\section{Де зберігається вага}
\label{sec:weightstore}

Що таке синапс, який «змінює свою вагу»? Зв'язок має два боки: закінчення проводу відправника (\emph{пресинаптичний} бік) і ділянку мембрани отримувача з рецепторами (\emph{постсинаптичний} бік), а між ними --- вузька щілина. У \nt{GLUT}-зв'язків мембрана отримувача зазвичай розташована на \emph{шипику} --- маленькому виступі вхідної гілки отримувача. Вагу зв'язку зручно розкласти на три множники~\cite{delCastillo1954}:
\begin{empheq}[box=\keybox]{equation}
w \;=\; N_s\,p\,A_1 .
\label{eq:quantal}
\end{empheq}
Тут $N_s$ --- кількість місць викиду між двома нейронами, $p$ --- імовірність, що імпульс викине порцію медіатора в одному місці (та сама $p$ з розділу~\ref{sec:code}), $A_1$ --- відповідь отримувача на одну порцію, тобто скільки рецепторів її вловлює. Множник $p$ живе у відправника, $A_1$ --- в отримувача, а $N_s$ потребує обох боків: нові місця викиду виростають, старі зникають, і це триває все життя.

\emph{Вирішує отримувач.} У типового \nt{GLUT}-синапсу один із рецепторів глутамату проводить струм, лише якщо збіглися дві умови: глутамат прийшов від відправника, і мембрана отримувача вже збуджена~\cite{Nowak1984}. Це правило Гебба, вбудоване в одну молекулу, --- детектор збігів із розділу~\ref{sec:glutcomputer}, поставлений на вхід самого синапсу. Спрацювавши, рецептор впускає в отримувача кальцій, і кальцій запускає зміну ваги. Отримувач --- єдине місце, де видно водночас і вхід, і вихід, тому рішення ухвалюється в нього.

\emph{Виконати рішення може будь-який бік.} Найбільш вивчене довготривале посилення відбувається в отримувача: у мембрану вбудовуються нові рецептори~\cite{Malinow2002}, і шипик росте~\cite{Matsuzaki2004} --- росте $A_1$. Але отримувач може змінити і $p$: він викидає речовину, яка йде через щілину назад, до відправника, --- зворотний канал із додатка~\ref{sec:othermediators}, --- і той починає викидати менше медіатора~\cite{WilsonNicoll2001}. А короткочасні зміни ваги, виснаження і полегшення з розділу~\ref{sec:code}, --- це зміни $p$, які відправник веде сам за своєю недавньою активністю.

Для інженера регістр ваги розділено між двома мікросхемами. Правило навчання виконує отримувач, а записати результат він може і в себе, і у відправника, через зворотний канал. Тому слово «локально» в правилах цього розділу означає не один бік зв'язку, а весь зв'язок цілком.

\section{Навчатися без учителя}

Найпростіше, що може зробити синапс, --- посилитися, коли відправник і отримувач активні разом:
\begin{equation}
\frac{\dd w}{\dd t} \;=\; \eta\,x\,y .
\label{eq:hebb}
\end{equation}
Це правило Гебба~\cite{Hebb1949}. Воно локальне і не потребує годинника. Але в нього та сама біда, що в \nt{GLUT}-мережі в розділі~\ref{sec:glutcomputer}: додатний зворотний зв'язок. Сильна вага робить $y$ більшим, великий $y$ ще посилює вагу, і ваги ростуть без меж. Ліки --- від'ємний зворотний зв'язок за вагою: нехай вага ще й спадає пропорційно до $y^2$. Для нейрона з входами $\mathbf x$ і лінійним виходом $y=\mathbf w\cdot\mathbf x$ отримуємо
\begin{empheq}[box=\keybox]{equation}
\frac{\dd\mathbf w}{\dd t} \;=\; \eta\,y\,\bigl(\mathbf x - y\,\mathbf w\bigr) .
\label{eq:oja}
\end{empheq}
Правило, як і раніше, локальне: синапс знає свій вхід $x_i$, вихід $y$ і вагу $w_i$.

\begin{keyblock}
\begin{proposition}[Що знаходить правило Гебба]
\label{prop:oja}
Правило~\eqref{eq:oja} зводить вектор ваг до одиничної довжини вздовж напрямку, в якому входи змінюються найсильніше, --- до головного власного вектора матриці кореляцій входів $C=\langle\mathbf x\mathbf x^{\mathsf T}\rangle$~\cite{Oja1982}. Нейрон навчається видавати найвиразнішу одну величину, в яку можна стиснути його входи.
\end{proposition}
\end{keyblock}

Доведення. Усереднимо~\eqref{eq:oja} за входами: $\dd\mathbf w/\dd t=\eta\bigl(C\mathbf w-(\mathbf w^{\mathsf T}C\mathbf w)\,\mathbf w\bigr)$. Нерухомі точки --- власні вектори $C$ одиничної довжини. Якщо $\mathbf w$ близький до власного вектора $\mathbf e_k$ із числом $\lambda_k$, мала добавка вздовж іншого вектора $\mathbf e_j$ росте як $e^{\eta(\lambda_j-\lambda_k)t}$. Стійка лише та точка, для якої всі $\lambda_j<\lambda_k$, тобто головний вектор.

Є й варіант правила Гебба, який дивиться на час імпульсів. Якщо імпульс відправника прийшов за $s$ мілісекунд \emph{до} імпульсу отримувача, вага росте на $A_+e^{-s/\tau_+}$; якщо через $s$ мілісекунд \emph{після} --- зменшується на $A_-e^{-s/\tau_-}$, з $\tau_\pm$ порядку двадцяти мілісекунд. Таке правило виміряно на синапсах \nt{GLUT}-нейронів~\cite{BiPoo1998}. Воно посилює зв'язки, які \emph{могли бути причиною} відповіді, і послаблює ті, що запізнилися (рис.~\ref{fig:stdp}). Проженіть через мережу багато разів одну послідовність збуджень, і в ній самі виростуть зв'язки від кожної ланки до наступної --- ланцюжок із розділу~\ref{sec:glutcomputer}, зібраний без інженера.

\begin{figure}[t]
\centering
\begin{tikzpicture}[>=Stealth,x=0.07cm,y=1.3cm]
\draw[->,gray!70!black] (-66,0) -- (68,0) node[right,black] {\small $s$};
\draw[->,gray!70!black] (0,-1.2) -- (0,1.35) node[above,black] {\small зміна ваги};
\draw[very thick,blue!60!black] plot[domain=0.5:60,samples=50] (\x,{exp(-\x/20)});
\draw[very thick,red!70!black] plot[domain=-60:-0.5,samples=50] (\x,{-0.6*exp(\x/20)});
\draw[blue!60!black,thick] (0.5,0) -- (0.5,1);
\draw[red!70!black,thick] (-0.5,0) -- (-0.5,-0.6);
\foreach \x in {-40,-20,20,40}{\draw[gray!60!black] (\x,-0.04) -- (\x,0.04);}
\node[below,font=\scriptsize] at (20,-0.04) {$20\ms$};
\node[above,font=\scriptsize] at (-20,0.04) {$-20\ms$};
\node[align=left,font=\small,blue!60!black,anchor=west] at (22,0.95) {відправник раніше:\\ зв'язок посилюється};
\node[align=right,font=\small,red!70!black,anchor=east] at (-12,-0.95) {відправник пізніше:\\ зв'язок слабшає};
\end{tikzpicture}
\caption{Правило Гебба за часом. По горизонталі --- $s=t_{\text{отрим}}-t_{\text{відпр}}$, на скільки імпульс отримувача відстав від імпульсу відправника; $\tau_\pm=20\ms$, $A_-=0.6A_+$. Вікно завширшки кілька десятків мілісекунд --- того самого порядку, що й вікно збігу~\eqref{eq:window}.}
\label{fig:stdp}
\end{figure}

Але всі правила цього підрозділу вчать того, що \emph{часто}, а не того, що \emph{корисно}. Синапс, якому відомі лише $x$ і $y$, не знає, чи принесла дія сік, чи біль. Щоб навчатися на результаті, синапсу потрібен третій сигнал.

\section{Третій множник}

Третій сигнал приносить \nt{DOPA}: дофамінові нейрони розсилають широкомовно одне число --- помилку передбачення винагороди $\delta$~\eqref{eq:rpe} (розділ~\ref{sec:neurontypes}). Нехай зміна ваги пропорційна добутку трьох множників: активності відправника, активності отримувача і $\delta$.

Складність у тому, що винагорода надходить не тоді, коли мережа вибирала дію, а через сотні мілісекунд або секунди. На момент надходження $\delta$ добуток $xy$ давно дорівнює нулю, і синапсу потрібна пам'ять про те, що він брав участь. Найпростіша пам'ять --- знайоме нам RC-коло:
\begin{empheq}[box=\keybox]{equation}
\tau_e\,\frac{\dd e}{\dd t} \;=\; -e \;+\; x\,y ,
\qquad
\frac{\dd w}{\dd t} \;=\; \eta\,\delta(t)\,e(t) .
\label{eq:threefactor}
\end{empheq}
Величину $e$ називають \emph{слідом}~\cite{Fremaux2016}. Спільна активність лишає в синапсі мітку, яка тане за час $\tau_e$. Якщо за цей час надходить дофамін, вага змінюється зі знаком $\delta$; якщо ні, мітка зникає безслідно (рис.~\ref{fig:trace}). На \nt{GLUT}-синапсах виміряно: дофамін, що надійшов протягом секунди-двох після спільної активності, перетворює її на посилення зв'язку, а той, що надійшов пізніше, --- ні~\cite{Yagishita2014}. Вікна пластичності завдовжки в секунди знайшли й там, де третім сигналом слугує не дофамін, а сильне збудження самого отримувача~\cite{Bittner2017}.

\begin{figure}[t]
\centering
\begin{tikzpicture}[>=Stealth,x=7cm,y=1cm]
\fill[orange!12] (0.7,-0.2) rectangle (0.8,5.2);
% rows: x*y, delta, traces, weights
\foreach \y/\lab in {4.4/{$x\,y$},3.1/{$\delta$},1.4/{слід $e$},0/{вага $w$}}{
  \draw[gray!70!black] (0,\y) -- (1.42,\y);
  \node[left,font=\small] at (-0.02,\y+0.3) {\lab};}
\draw[very thick,blue!60!black] (0,4.4) -- (0,4.9) -- (0.2,4.9) -- (0.2,4.4);
\node[right,font=\scriptsize,blue!60!black] at (0.21,4.8) {відправник і отримувач активні разом};
\draw[very thick,orange!85!black] (0,3.1) -- (0.7,3.1) -- (0.7,3.7) -- (0.8,3.7) -- (0.8,3.1) -- (1.4,3.1);
\node[right,font=\scriptsize,orange!85!black] at (0.81,3.6) {надійшов дофамін};
% traces, each in units of its own maximum
\draw[very thick,blue!60!black] plot[domain=0:0.2,samples=20] (\x,{1.4+1.1*(1-exp(-\x))/(1-exp(-0.2))})
  plot[domain=0.2:1.4,samples=60] (\x,{1.4+1.1*exp(-(\x-0.2))});
\draw[thick,densely dashed,gray!50!black] plot[domain=0:0.2,samples=30] (\x,{1.4+1.1*(1-exp(-\x/0.05))/(1-exp(-4))})
  plot[domain=0.2:1.4,samples=80] (\x,{1.4+1.1*exp(-(\x-0.2)/0.05)});
\node[right,font=\scriptsize,blue!60!black] at (1.0,2.15) {$\tau_e=1\,\text{с}$};
\node[right,font=\scriptsize,gray!40!black] at (0.3,1.65) {$\tau_e=50\ms$};
% weights
\draw[very thick,blue!60!black] (0,0.25) -- (0.7,0.25) -- (0.8,0.8) -- (1.4,0.8) node[right,font=\scriptsize] {$\tau_e=1\,\text{с}$: вага зросла};
\draw[thick,densely dashed,gray!50!black] (0,0.2) -- (1.4,0.2) node[right,font=\scriptsize,gray!40!black] {$\tau_e=50\ms$: без змін};
% time axis marks
\foreach \x/\l in {0/0,0.2/0.2,0.7/0.7,1.4/1.4}{\draw[gray!60!black] (\x,-0.05) -- (\x,0.05) node[below=3pt,font=\scriptsize] {\l\,с};}
\draw[<->,thick,gray!50!black] (0.2,5.35) -- node[above,font=\scriptsize,black] {затримка винагороди $0.5\,\text{с}$} (0.7,5.35);
\end{tikzpicture}
\caption{Слід за правилом~\eqref{eq:threefactor}. Спільна активність триває $0.2\,\text{с}$, дофамін надходить через півсекунди після її кінця. Сліди показано в частках свого максимуму. Повільний слід ($\tau_e=1\,\text{с}$) на цей момент ще живий, швидкий ($\tau_e=50\ms$) давно розтанув.}
\label{fig:trace}
\end{figure}

\begin{keyblock}
\begin{proposition}[Навчання широкомовним сигналом]
\label{prop:threefactor}
Правило~\eqref{eq:threefactor} змінює лише ті синапси, які брали участь у роботі мережі за останні $\tau_e$, у бік, заданий спільним сигналом $\delta$. Широкомовний сигнал разом із локальним слідом дає адресну поправку. Затримка між дією і результатом допустима, доки вона не більша за $\tau_e$.
\end{proposition}
\end{keyblock}

\section{Чому це працює: шум як пошук}
\label{sec:dofagradient}

Чому правило~\eqref{eq:threefactor} веде до кращого? Відповідь дає шум із розділу~\ref{sec:code}. Нехай вихід нейрона $y=f(u)+\xi$, де $u=\sum_jw_jx_j$, а $\xi$ --- випадкове тремтіння з нульовим середнім і дисперсією $\sigma^2$. Винагорода $R$ залежить від $y$. Візьмімо як слід не просто $x_jy$, а його випадкову частину $x_j\xi$, а як третій множник --- помилку $\delta=R-\bar R$, де $\bar R$ --- очікувана винагорода. За малого шуму $R\approx R\bigl(f(u)\bigr)+R'\,\xi$, тому
\begin{empheq}[box=\keybox]{equation}
\bigl\langle \delta\,\xi\,x_j \bigr\rangle \;=\; \sigma^2\,R'\,x_j \;=\; \frac{\sigma^2}{f'(u)}\,\frac{\partial\langle R\rangle}{\partial w_j} .
\label{eq:perturbation}
\end{empheq}
Середній крок іде за градієнтом очікуваної винагороди~\cite{Williams1992}. Ніхто не обчислював похідних: мережа випадково відхилилася, дофамін повідомив, чи стало краще, і синапси, що брали участь, зсунулися у вигідний бік (рис.~\ref{fig:perturbation}). Шум, який у розділі~\ref{sec:code} заважав передавати повідомлення, тут стає пошуком.

\begin{figure}[t]
\centering
\begin{tikzpicture}[>=Stealth,x=2cm,y=3cm]
\draw[->,gray!70!black] (0,0) -- (4.2,0) node[below left,font=\small,black] {вихід $y$};
\draw[->,gray!70!black] (0,0) -- (0,1.2) node[left,font=\small,black] {винагорода $R$};
% reward hill
\draw[very thick,blue!60!black] plot[domain=0:4,samples=60] (\x,{max(0,1-((\x-3)/2.2)^2)});
\node[font=\scriptsize,blue!60!black,anchor=south] at (3,1.02) {найкращий вихід};
% noise around f(u)
\draw[thick,gray!60!black] plot[domain=0.6:2.4,samples=50] (\x,{0.28*exp(-((\x-1.5)/0.3)^2/2)});
\draw[densely dashed,gray!60!black] (1.5,0) -- (1.5,0.535);
\node[below,font=\small] at (1.5,0) {$f(u)$};
\node[font=\scriptsize,gray!50!black] at (1.5,-0.2) {тремтіння $\xi$};
% expected reward
\draw[densely dotted,gray!60!black] (0.5,0.535) node[left,font=\scriptsize,black] {$\bar R$} -- (2.1,0.535);
% two random deviations
\fill[green!45!black] (2.1,0.833) circle (0.035);
\draw[very thick,green!45!black] (2.1,0.535) -- (2.1,0.833);
\node[font=\scriptsize,green!45!black,anchor=west,align=left] at (2.18,0.6) {$\xi>0$: краще,\\ $\delta>0$};
\fill[red!70!black] (0.9,0.089) circle (0.035);
\draw[very thick,red!70!black] (0.9,0.535) -- (0.9,0.089);
\node[font=\scriptsize,red!70!black,anchor=east,align=right] at (0.83,0.27) {$\xi<0$: гірше,\\ $\delta<0$};
% resulting step
\draw[->,very thick,orange!85!black] (1.55,0.4) -- (1.95,0.4) node[right,font=\scriptsize,black] {крок};
\end{tikzpicture}
\caption{Шум як пошук~\eqref{eq:perturbation}. Вихід нейрона тремтить навколо $f(u)$. Випадкове відхилення праворуч (зелене) дало більше за очікуване, $\delta>0$; ліворуч (червоне) --- менше, $\delta<0$. В обох випадках добуток $\delta\,\xi$ додатний, і вага зсувається праворуч, туди, де винагорода росте. У середньому крок пропорційний нахилу $R'$: на вершині гірки відхилення в обидва боки однаково погані, і кроки гасять один одного.}
\label{fig:perturbation}
\end{figure}

\begin{keyblock}
\begin{proposition}[Спуск за градієнтом без зворотних проводів]
\label{prop:gradient}
Якщо в мережі є випадкові відхилення активності, а широкомовний сигнал дорівнює помилці передбачення винагороди і в кожній обстановці в середньому дорівнює нулю, правило~\eqref{eq:threefactor} у середньому змінює ваги за градієнтом очікуваної винагороди. Зворотні проводи й окрема фаза навчання для цього не потрібні.
\end{proposition}
\end{keyblock}

Тепер зрозуміло, чому дофамін повідомляє не про винагороду, а про \emph{помилку} її передбачення. Якби третім множником була сама винагорода $R\ge0$, середнє в~\eqref{eq:perturbation} не змінилося б, але з'явилися б дві біди. По-перше, до корисної частини додався б доданок $\bar R\,\xi x_j$: у середньому нуль, але сильний шум. По-друге, що гірше, справжній синапс не вміє відокремлювати в сліді $x_jy$ випадкову частину від невипадкової. За даних входів $x_jf(u)$ --- те саме число від спроби до спроби; якщо $\delta$ у цій обстановці в середньому дорівнює нулю, ця частина сліду нічого не змінює, а при $\delta\ge0$ вона раз у раз посилює все, що мережа робила, --- правильне і неправильне. Тому $\bar R$ має бути очікуванням \emph{у даній обстановці}, а не середнім за все життя. Обчислювати його --- справа критика, про якого нижче.

Ми перевірили це на моделі. Дві групи \nt{GLUT}-нейронів вибирають одну з двох дій через взаємне гальмування, як на рис.~\ref{fig:wta}; ваги від сигналу «почали» до груп спочатку рівні, а шум вирішує, хто переможе. Дія $A$ приносить сік з імовірністю $0.8$, дія $B$ --- з імовірністю $0.2$, і сік надходить через півсекунди після вибору. При $\tau_e=1\,\text{с}$ і $\delta=R-\bar R$ частка виборів $A$ за п'ятсот спроб зросла з $0.65$--$0.8$ у першій сотні до $0.96$--$1.0$ в останній, а ваги лишилися в межах $0.85$--$1.35$. При $\tau_e=50\ms$, меншому за затримку винагороди, мережа не навчилася нічого: частка виборів $A$ лишилася близько половини. При $\delta=R$ без віднімання очікування мережа теж стала вибирати $A$, але вага переможця за ті самі п'ятсот спроб зросла в тисячу разів і продовжувала рости; а при тому самому $\delta=R$ і вдесятеро більшій швидкості навчання мережа в одному з трьох прогонів назавжди закріпила свій перший випадковий вибір --- гірший.

\section{Ціна одного проводу}

Сигнал $\delta$ один на всіх, а шум, який він оцінює, --- сума тремтінь усіх $N$ нейронів, що впливають на результат. Кожен синапс бачить у $\delta$ свою корисну частину і чужий шум $N-1$ сусідів. Щоб виокремити своє, доводиться усереднювати за багатьма спробами, і час навчання росте пропорційно до $N$~\cite{Werfel2005}. Зворотне поширення такої плати не потребує.

\begin{keyblock}
\begin{proposition}[Ціна широкомовлення]
\label{prop:broadcastcost}
Час навчання за одним спільним сигналом росте пропорційно до кількості нейронів, чий шум впливає на результат. Таке навчання придатне для невеликих кіл, що вибирають із небагатьох варіантів, але не для налаштування мільярдів ваг.
\end{proposition}
\end{keyblock}

Мозок, схоже, так і розподіляє роботу: виходи \nt{DOPA}-нейронів ідуть насамперед до областей, які вибирають дії (розділ~\ref{sec:neurontypes}), а величезна \nt{GLUT}-мережа кори, де переважна більшість ваг, навчається здебільшого місцевими правилами, без зовнішнього вчителя. Раніше їх зводили до правил Гебба~\eqref{eq:oja}. Тепер дедалі більше даних свідчить, що кора має власну мету --- \emph{передбачати} свій наступний вхід, і тоді помилка обчислюється на місці, як різниця між передбаченим і тим, що надійшло~\cite{Keller2018}: учитель вбудований у саму мережу. Вікна пластичності завдовжки в секунди, де третім сигналом слугує сильне збудження самого отримувача~\cite{Bittner2017}, показують, що місцевий сигнал помилки в синапсів справді буває. Наскільки це загальне правило кори, --- гіпотеза.

\section{Звідки береться помилка: критик у неперервному часі}

Лишається питання, як самі \nt{DOPA}-нейрони обчислюють $\delta$. У програмах навчання з підкріпленням помилку передбачення рахують кроками $t,t+1,\dots$ В організмі кроків немає, тому запишімо її в неперервному часі~\cite{Doya2000}. Нехай $r(t)$ --- потік винагороди, а $V(t)$ --- очікування всієї майбутньої винагороди, в якій віддаленіша цінується менше:
\begin{equation}
V(t) \;=\; \int_t^{\infty} e^{-(s-t)/\tau_\gamma}\,r(s)\,\dd s .
\end{equation}
Продиференціювавши, отримаємо $\dd V/\dd t=V/\tau_\gamma-r$. Нев'язка цієї рівності і є помилкою передбачення:
\begin{empheq}[box=\keybox]{equation}
\delta(t) \;=\; r(t) \;+\; \frac{\dd V}{\dd t} \;-\; \frac{V}{\tau_\gamma} .
\label{eq:ctd}
\end{empheq}
Стала $\tau_\gamma$ --- горизонт планування, неперервний аналог коефіцієнта $\gamma$; за гіпотезою, його задає \nt{SERO} (підрозділ~\ref{sec:horizon}), і тоді~\eqref{eq:patience} --- правило, скільки варто чекати. Перевірмо три випадки (рис.~\ref{fig:ctdcases}). Засвітилася лампочка, що обіцяє сік: $V$ стрибком зросло, $\dd V/\dd t$ велике, сплеск. Надійшов обіцяний сік: $r>0$, але очікування в ту саму мить вичерпалося, $\dd V/\dd t\approx-r$, і сплеску немає. Сік не надійшов: очікування зникло без винагороди, $\dd V/\dd t<0$, провал.

\begin{figure}[t]
\centering
\begin{tikzpicture}[>=Stealth,x=1.15cm,y=1cm]
\foreach \col/\ttl/\lampon/\juice in {0/{несподіваний сік}/0/1, 1/{обіцяний сік}/1/1, 2/{сік не надійшов}/1/0}{
 \begin{scope}[xshift=\col*4.6cm]
  \node[font=\small] at (1.5,3.55) {\ttl};
  \foreach \yy in {2.4,1.2,0}{\draw[gray!60!black] (0,\yy) -- (3.1,\yy);}
  % reward r
  \ifnum\juice=1
    \fill[green!45!black] (1.95,2.4) rectangle (2.05,2.85);
    \pic[scale=0.55] at (2.0,3.1) {juice};
  \else
    \pic[scale=0.55] at (2.0,3.1) {nojuice};
  \fi
  % expectation V and error delta
  \ifnum\lampon=1
    \pic[scale=0.55] at (1.0,3.1) {lamp};
    \draw[very thick,blue!60!black] (0,1.2) -- (1,1.2) -- (1,1.45)
      plot[domain=1:2,samples=20] (\x,{1.2+0.25*exp((\x-1)/1.5)}) -- (2,{1.2+0.25*exp(1/1.5)}) -- (2,1.2) -- (3.1,1.2);
    \draw[very thick,red!70!black] (0,0) -- (0.95,0) -- (0.95,0.5) -- (1.05,0.5) -- (1.05,0) -- (3.1,0);
    \ifnum\juice=0
      \draw[very thick,red!70!black] (1.95,0) -- (1.95,-0.45) -- (2.05,-0.45) -- (2.05,0);
      \node[font=\scriptsize,red!70!black,anchor=west] at (2.1,-0.35) {провал};
    \else
      \node[font=\scriptsize,gray!50!black,anchor=north,align=center] at (2.0,-0.08) {$r$ і спад $V$\\ гасять одне одного};
    \fi
  \else
    \draw[very thick,blue!60!black] (0,1.2) -- (3.1,1.2);
    \draw[very thick,red!70!black] (0,0) -- (1.95,0) -- (1.95,0.5) -- (2.05,0.5) -- (2.05,0) -- (3.1,0);
  \fi
  \ifnum\col=0
    \node[font=\small,anchor=east] at (-0.1,2.55) {$r$};
    \node[font=\small,anchor=east] at (-0.1,1.35) {$V$};
    \node[font=\small,anchor=east] at (-0.1,0.1) {$\delta$};
  \fi
 \end{scope}}
\end{tikzpicture}
\caption{Три випадки для помилки~\eqref{eq:ctd}; згори донизу --- винагорода $r$, очікування $V$ і помилка $\delta$. Ліворуч очікування немає, і сік цілком --- несподіванка. Посередині лампочка стрибком піднімає $V$: це стрибок $\dd V/\dd t$, сплеск $\delta$ припадає на лампочку. Далі $V$ росте рівно так, як вимагає горизонт, і $\delta=0$; у момент соку винагорода $r$ і спад $V$ гасять одне одного. Праворуч $V$ падає без винагороди --- провал. Це ті самі сплеск, перенесення і провал, що на рис.~\ref{fig:rpe}, але тепер видно, звідки вони беруться.}
\label{fig:ctdcases}
\end{figure}

Що потрібно для~\eqref{eq:ctd} зі схем, які в нас уже є? Три речі.

\emph{Сама оцінка $V$.} Її видає група \nt{GLUT}-нейронів --- критик, --- ваги якої навчаються тим самим правилом~\eqref{eq:threefactor} з тим самим $\delta$: якщо $\delta>0$, очікування було занижене, і зв'язки, активні перед цим, посилюються.

\emph{Похідна $\dd V/\dd t$.} Її обчислює будь-яка схема, що реагує на зміни, а не на рівень: пряме збудження разом із затриманим гальмуванням через \nt{GABA}-посередника, як у детекторі напрямку з розділу~\ref{sec:gaba}, або синапс із виснаженням~\eqref{eq:depression} з розділу~\ref{sec:code}.

\emph{Знак в одному проводі.} Помилка буває і додатною, і від'ємною, а частота імпульсів --- лише додатною. Розв'язок той самий, що в \nt{SERO}-нейрона в розділі~\ref{sec:serocomputer}: \nt{DOPA}-нейрон --- ритмоводій. Його рівні кілька імпульсів на секунду означають $\delta=0$, пачка --- $\delta>0$, пауза --- $\delta<0$. Від'ємній помилці лишається менше місця: частоту не можна опустити нижче нуля. У справжніх \nt{DOPA}-нейронів провали справді мілкіші за сплески~\cite{Bayer2005}.

Що дофамін поводиться як $\delta$, виміряно. Як саме мозок обчислює $\dd V/\dd t$, --- гіпотеза.

\section{Актор і критик}

Зберімо все разом (рис.~\ref{fig:actorcritic}). Нейрони обстановки збуджують групи дій, що вибирають переможця через \nt{GABA} (твердження~\ref{prop:wta}), --- це \emph{актор}, --- і критика, що видає $V$~\cite{SuttonBarto2018}. \nt{DOPA}-нейрон отримує винагороду $r$ і $V$, обчислює~\eqref{eq:ctd} і розсилає $\delta$ усім синапсам актора і критика.

\begin{figure}[t]
\centering
\begin{tikzpicture}[>=Stealth,
  nrn/.style={circle,draw,very thick,fill=blue!6,minimum size=0.9cm,inner sep=0pt,font=\small},
  gab/.style={circle,draw,very thick,fill=red!10,minimum size=0.6cm,inner sep=0pt},
  dofa/.style={circle,draw,very thick,double,double distance=1.2pt,fill=orange!15,minimum size=1.35cm,inner sep=0pt,font=\scriptsize},
  inh/.style={-{Bar[width=7pt]},thick,red!70!black},
  bc/.style={->,thick,densely dashed,orange!85!black}]
\node[nrn] (S1) at (0,2.4) {$x_1$};
\node[nrn] (S2) at (0,1.2) {$x_2$};
\node[nrn] (S3) at (0,0) {$x_3$};
\node[left=0.15cm,align=right,font=\small] at (-0.5,1.2) {обстановка};
\node[nrn] (A) at (4,2.6) {$A$};
\node[nrn] (B) at (4,1.0) {$B$};
\node[gab] (G) at (5.2,1.8) {};
\draw[->,thick] (A) -- (G); \draw[->,thick] (B) -- (G);
\draw[inh] (G) to[bend left=35] (A);
\draw[inh] (G) to[bend right=35] (B);
\node[above,font=\small] at (4.6,3.05) {актор: вибір дії};
\node[nrn] (V) at (4,-1.1) {$V$};
\node[below,font=\small] at (4,-1.6) {критик};
\foreach \s in {S1,S2,S3}{\foreach \t in {A,B,V}{\draw[->,gray!60!black] (\s) -- (\t);}}
\node[dofa] (D) at (8.0,-1.1) {\nt{DOPA}};
\draw[->,thick] (V) -- node[above,font=\scriptsize] {$V,\ \dd V/\dd t$} (D);
\draw[->,thick,green!45!black] (10.0,-1.1) node[right,black,font=\small] {винагорода $r$} -- (D);
\pic[scale=1.1] at (10.6,-0.5) {juice};
\draw[bc] (D.north) -- (8.0,4.0) -- node[above,font=\small,black] {$\delta$ --- усім синапсам актора і критика} (2.2,4.0) -- (2.2,3.0);
\draw[bc] (D.south) -- (8.0,-2.4) -- (2.2,-2.4) -- (2.2,-0.6);
\end{tikzpicture}
\caption{Мережа, яка навчається вибирати дії. Сірі стрілки --- \nt{GLUT}-синапси зі змінними вагами; червоний кружок --- \nt{GABA}-нейрон; подвійний кружок --- \nt{DOPA}-ритмоводій, що обчислює $\delta$ за~\eqref{eq:ctd}; штрихові стрілки --- широкомовна розсилка $\delta$.}
\label{fig:actorcritic}
\end{figure}

Знак дії медіатора вибирає отримувач (твердження~\ref{prop:receptor}), і в області вибору дій частина нейронів несе дофамінові рецептори однієї родини, частина --- іншої. У дослідах на зрізах дофамін разом зі спільною активністю в перших посилює синапси, у других --- послаблює~\cite{Shen2008}; у моделі це означає, що в других знак $\delta$ у~\eqref{eq:threefactor} перевернутий. Перші навчаються того, що \emph{варто робити}, другі --- того, чого \emph{робити не варто}.

\section{Підсумок}

\begin{table}[t]
\centering
\footnotesize
\setlength{\tabcolsep}{4pt}
\renewcommand{\arraystretch}{1.15}
\begin{tabular}{>{\raggedright}p{2.6cm}>{\raggedright}p{2.7cm}>{\raggedright}p{3.1cm}>{\raggedright\arraybackslash}p{5.0cm}}
\toprule
Правило & Що знає синапс & Чого навчається мережа & Що відомо \\
\midrule
Гебб за частотами~\eqref{eq:oja} & $x$, $y$, свою вагу & стискання входів: головних напрямків & посилення за спільної активності виміряно; механізм обмеження ваг --- предмет досліджень \\
Гебб за часом & порядок імпульсів $x$ і $y$ у межах $\sim20\ms$ & причинних зв'язків, послідовностей & виміряно на \nt{GLUT}-синапсах \\
три множники~\eqref{eq:threefactor} & $x$, $y$, слід $e$, спільний сигнал $\delta$ & дій, які приносять винагороду & вплив дофаміну протягом секунд після активності виміряно; як головний механізм навчання вибору --- добре обґрунтована гіпотеза \\
критик~\eqref{eq:ctd} & те саме & очікування винагороди $V$ & схожість дофаміну з $\delta$ виміряно; схема, що обчислює $\dd V/\dd t$, --- гіпотеза \\
зворотне поширення помилки & персональну поправку для кожної ваги & усього, але потрібні зворотні проводи і такт & у мозку в цій формі не знайдено; шукають наближення до нього \\
\bottomrule
\end{tabular}
\caption{Правила навчання в мережі з \nt{GLUT}-, \nt{GABA}- і \nt{DOPA}-нейронів.}
\label{tab:learning}
\end{table}

Правила навчання зведено в табл.~\ref{tab:learning}. \nt{GLUT} несе дані, \nt{GABA} керує потоком, а \nt{DOPA} вирішує, які зміни в мережі залишити.

\section{Задачі}

\begin{exercise}[\lvlA]
\label{ex:06:1}
\emph{Скільки живе слід.} Спільна активність $xy=1$ триває $T_a=0.2\,\text{с}$, починаючи з $e=0$, а дофамін надходить через $d=0.5\,\text{с}$ після її кінця. (а)~У скільки разів слід~\eqref{eq:threefactor} у цей момент більший при $\tau_e=1\,\text{с}$, ніж при $\tau_e=50\ms$? (б)~За якого $\tau_e$ він найбільший?
\end{exercise}

\begin{exercise}[\lvlA]
\label{ex:06:2}
\emph{Дрейф правила Гебба.} Відправник і отримувач видають незалежні пуассонівські потоки по $10$ імпульсів на секунду, і кожна пара імпульсів змінює вагу за вікном рис.~\ref{fig:stdp} з $\tau_\pm=20\ms$, $A_-=0.6A_+$. На скільки $A_+$ зросте вага за годину цілком випадкової активності? За якого $\tau_-$ дрейф зникає?
\end{exercise}

\begin{exercise}[\lvlB]
\label{ex:06:3}
\emph{Кореляції, а не коваріації.} Правило твердження~\ref{prop:oja} знаходить головний власний вектор $C=\langle\mathbf x\mathbf x^{\mathsf T}\rangle$. Частоти додатні: $\mathbf x=\mathbf m+\mathbf n$, $\mathbf m=(\mu,0)$, коваріація шуму $\operatorname{diag}(s_1^2,s_2^2)$. За якої умови нейрон вивчить $\mathbf e_1$? Що він вивчить при $\mu=1$, $s_1=0.1$, $s_2=0.5$?
\end{exercise}

\begin{exercise}[\lvlA]
\label{ex:06:4}
\emph{Горизонт у сплеску.} Лампочка засвічується в непередбачуваний момент, сік розміру $R$ надходить через $L=1\,\text{с}$. Покажіть, що для вивченого очікування за~\eqref{eq:ctd} між лампочкою і соком $\delta\equiv0$, і знайдіть площу сплеску на лампочку при $\tau_\gamma=2\,\text{с}$ і $\tau_\gamma=4\,\text{с}$.
\end{exercise}

\begin{exercise}[\lvlB]
\label{ex:06:5}
\emph{Звідки береться ціна широкомовлення.} $N$ нейронів тремтять незалежно, $\xi_i\sim\mathcal N(0,\sigma^2)$, помилка $\delta=a\sum_i\xi_i$, синапс $j$ оцінює градієнт як $\delta\,\xi_jx_j$. Знайдіть відношення сигнал/шум однієї спроби. Один нейрон навчається за $100$ спроб по $5\,\text{с}$: скільки триватиме навчання при $N=10^4$ і $N=10^{10}$?
\end{exercise}

\begin{exercise}[\lvlB]
\label{ex:06:6}
\emph{Проєктування: схема, що обчислює $\delta$.} \nt{DOPA}-нейрон отримує $r$, $V$ з вагою $k_1$ і копію $V$, згладжену з $\tau_d$, з вагою $-k_2$. (а)~Доберіть $k_1$, $k_2$ так, щоб для повільних $V$ вихід був~\eqref{eq:ctd}. (б)~При $V=V_0e^{t/\tau_\gamma}$ істинна $\delta=0$. За якого $\tau_d$ хибна помилка схеми не перевищує $5\,\%$ від $V/\tau_\gamma$, якщо $\tau_\gamma=2\,\text{с}$?
\end{exercise}

\begin{exercise}[\lvlB]
\label{ex:06:7}
\emph{Моделювання: гранична затримка винагороди.} Додайте в \texttt{run()} копії \texttt{models/\allowbreak dofa\_learning.py} затримку винагороди $d$. Знайдіть частку виборів $A$ в останній сотні з $500$ спроб при $d=0.5\,\text{с}$ і $\tau_e=0.05$, $0.2$, $1$, $4\,\text{с}$ та при $\tau_e=1\,\text{с}$, $d=3\,\text{с}$. За якої частки сліду, що доживає до винагороди (задача~\ref{ex:06:1}), навчання зникає?
\end{exercise}

\begin{exercise}[\lvlC]
\label{ex:06:8}
\emph{Чи можна обійти ціну широкомовлення.} $N$ нейронів: $y_i=w_i+\xi_i$, $\xi_i\sim\mathcal N(0,\sigma^2)$, винагорода $R=-\sum_i(y_i-w_i^*)^2$, правило $\Delta w_i=\eta\,(R-\langle R\rangle)\,\xi_i$. За скільки спроб за найкращого $\eta$ помилка $\sum_i(w_i-w_i^*)^2$ спадає в $e$ разів? А якщо тремтять лише $m$ випадкових нейронів з $N$? Чи допомагає розріджений пошук?
\end{exercise}

\chapter{Альтернативні сучасні теорії навчання на основі дофаміну}
\chaptermark{Інші теорії \nt{DOPA}}
\label{sec:dofaalt}
\begin{chaptergoals}
\item як \nt{DOPA}-нейрони з асиметричною відповіддю на помилки записують увесь розподіл винагороди;
\item як частина \nt{DOPA}-нейронів повідомляє важливість чи небезпеку, а не знак помилки;
\item як один провід несе швидку помилку, що вчить, і повільний рівень, що задає швидкість дій;
\item чому сигнал \nt{DOPA} --- джгут, а не одне число, і яка теорія ставить під сумнів саму помилку.
\end{chaptergoals}

\section{Що насправді йде проводом}

У розділі~\ref{sec:dofa} \nt{DOPA}-система була в нас одним проводом, яким іде одне число --- помилка передбачення винагороди $\delta$~\eqref{eq:ctd}. Ця картина пояснює багато, але що точніше вимірюють дофамін, то більше знаходять того, що в неї не вкладається. Одні нейрони відповідають на неприємне, як на винагороду; концентрація дофаміну росте, поки тварина біжить до мети, хоча нічого несподіваного не відбувається; різні \nt{DOPA}-нейрони поводяться по-різному.

Для інженера всі ці знахідки --- питання про одне: \emph{який формат має сигнал на широкомовній шині?} Одне число чи вектор? Один сигнал на проводі чи кілька, рознесених за частотою? Чи говорить він про майбутнє, як $\delta$, чи про минуле? Нижче --- шість сучасних відповідей; для кожної ми відокремимо виміряне від припущеного.

\section{Не одне число, а розподіл}

Помилка~\eqref{eq:ctd} вчить критика \emph{середнього} очікування винагороди. Але пів краплі соку напевно і крапля з імовірністю одна друга мають однакове середнє. Щоб їх розрізняти, потрібен увесь розподіл винагороди.

Візьмімо найпростішу задачу: після провісника надходить винагорода випадкового розміру $r$. Нехай \nt{DOPA}-нейронів багато, і кожен, $i$-й, обслуговує свою оцінку $V_i$, тож у момент винагороди його помилка дорівнює $\delta_i=r-V_i$. І нехай додатну помилку кожен враховує з вагою $\alpha_i^+$, від'ємну --- з вагою $\alpha_i^-$. Після кожної винагороди оцінка зсувається на
\begin{equation}
\Delta V_i \;=\; \eta\,\bigl(\alpha_i^+\,[\delta_i]_+ \;-\; \alpha_i^-\,[-\delta_i]_+\bigr).
\label{eq:asymmetric}
\end{equation}
Оцінка перестає змінюватися, коли в середньому додатні поправки врівноважують від'ємні:
\begin{empheq}[box=\keybox]{equation}
\alpha_i^+\,\bigl\langle[r-V_i]_+\bigr\rangle \;=\; \alpha_i^-\,\bigl\langle[V_i-r]_+\bigr\rangle ,
\qquad
\kappa_i \;=\; \frac{\alpha_i^+}{\alpha_i^++\alpha_i^-} .
\label{eq:expectile}
\end{empheq}
При $\kappa_i=1/2$ це звичайне середнє. При $\kappa_i>1/2$ нейрон --- «оптиміст»: його оцінка зсунута до верхнього краю розподілу, при $\kappa_i<1/2$ --- «песиміст».

Приклад: сік надходить з імовірністю $1/2$, розмір краплі $1$. Тоді~\eqref{eq:expectile} дає $\kappa_i\cdot\tfrac12(1-V_i)=(1-\kappa_i)\cdot\tfrac12V_i$, тобто $V_i=\kappa_i$ (рис.~\ref{fig:expectiles}). Набір нейронів із різними $\kappa_i$ записує розподіл винагороди так само, як набір квантилів записує розподіл у статистиці.

\begin{figure}[t]
\centering
\begin{tikzpicture}[>=Stealth,x=7cm,y=3cm]
\draw[->,gray!70!black] (-0.08,0) -- (1.15,0) node[right,black] {\small винагорода};
\draw[->,gray!70!black] (-0.08,0) -- (-0.08,0.75) node[left,black] {\small імовірність};
\fill[blue!25] (-0.03,0) rectangle (0.03,0.5);
\fill[blue!25] (0.97,0) rectangle (1.03,0.5);
\node[below,font=\small] at (0,0) {$0$};
\node[below,font=\small] at (1,0) {$1$};
\node[left,font=\scriptsize] at (-0.08,0.5) {$1/2$};
\foreach \k/\c/\lab/\h in {0.2/blue!60!black/{песиміст, $\kappa=0.2$}/0.62, 0.5/gray!50!black/{середнє, $\kappa=0.5$}/0.42, 0.8/red!70!black/{оптиміст, $\kappa=0.8$}/0.22}{
  \draw[very thick,\c] (\k,0) -- (\k,\h);
  \node[above,font=\scriptsize,\c] at (\k,\h) {\lab};}
\end{tikzpicture}
\caption{Розподіл винагороди записано набором \nt{DOPA}-нейронів. Винагорода $0$ або $1$ з імовірністю $1/2$ (сині стовпчики); нейрон із часткою $\kappa$ приходить до оцінки $V=\kappa$~\eqref{eq:expectile}.}
\label{fig:expectiles}
\end{figure}

Що виміряно: у різних \nt{DOPA}-нейронів різна «точка перевертання» --- розмір винагороди, за якого відповідь змінюється зі сплеску на провал, --- і нейрони з більшою асиметрією відповідей на додатні й від'ємні помилки перевертаються за більших винагород, як і вимагає~\eqref{eq:expectile}~\cite{Dabney2020}. З відповідей групи нейронів вдається відновити форму розподілу. Що це головна функція розкиду, а не побічний ефект, --- гіпотеза.

\begin{keyblock}
\begin{proposition}[Розподіл з асиметрії]
\label{prop:expectile}
Якщо \nt{DOPA}-нейрони різняться часткою $\kappa_i$, з якою вони враховують додатні помилки, їхні оцінки сходяться до різних точок~\eqref{eq:expectile} розподілу винагороди. Група нейронів передає не середнє, а розподіл, а отже --- ризик.
\end{proposition}
\end{keyblock}

\section{Не лише помилка, а й важливість}

За формулою помилки~\eqref{eq:ctd} неприємна подія --- удар, гіркий смак --- має спричиняти провал: вийшло гірше, ніж очікувалося. Частина \nt{DOPA}-нейронів так і робить, а інша відповідає на неприємне \emph{сплеском}, як на винагороду~\cite{Matsumoto2009}. Ці нейрони повідомляють не знак помилки, а те, що сталося щось \emph{важливе}: несподіване, сильне, що потребує уваги~\cite{BrombergMartin2010}.

Інженерові зручно думати про це як про два сигнали. Перший --- знакова помилка $\delta$: куди змінювати ваги. Другий --- її модуль $|\delta|$ або новизна події: \emph{наскільки сильно} їх змінювати, тобто швидкість навчання; після несподіваної події навчатися треба швидше. За змістом він близький до норадреналінового коефіцієнта підсилення з розділу~\ref{sec:neurontypes}.

Є й третій варіант: окремий провід для окремої осі. \nt{DOPA}-нейрони, що йдуть до однієї з областей вибору дій, відповідають не на винагороду, а на новизну і силу подразника і потрібні, щоб тварина навчилася уникати небезпечного~\cite{Menegas2018}: проводом іде не «краще чи гірше», а «небезпечно чи ні».

Що виміряно: різні групи \nt{DOPA}-нейронів по-різному відповідають на неприємне і на нове, і різні виходи вчать різного. Як саме мозок користується сигналом важливості --- як швидкістю навчання, як сигналом уваги чи як окремою помилкою за віссю небезпеки, --- обговорюється.

\section{Не лише вчити, а й підганяти}

Дофамін має й негайну дію: що його більше, то охочіше і швидше тварина діє. Як поєднати це з навчанням?

Інженерна відповідь --- \emph{розділення за частотою}. Швидкі сплески і провали тривалістю в сотні мілісекунд несуть $\delta$ і вчать. Повільний рівень, що змінюється за хвилини, несе іншу величину --- середню винагороду за одиницю часу $\bar R$~\cite{Niv2007}. Нехай дія, виконана за час $\tau_a$, потребує зусилля $C/\tau_a$: що швидше, то дорожче. Натомість кожна секунда зволікання коштує $\bar R$ --- стільки винагороди можна було б за неї отримати. Сумарна ціна $C/\tau_a+\bar R\,\tau_a$ мінімальна при
\begin{empheq}[box=\keybox]{equation}
\tau_a^{*} \;=\; \sqrt{\frac{C}{\bar R}} .
\label{eq:vigor}
\end{empheq}
Що багатша обстановка, то дорожчий час і то швидше вигідно діяти: якщо $\bar R$ подвоїлася, час дії зменшується в $\sqrt2\approx1.4$ раза. Зауважте, що $\bar R$ --- та сама очікувана винагорода, яку віднімали в розділі~\ref{sec:dofa}.

Викид дофаміну в місці призначення може змінюватися і без зміни частоти імпульсів: його підлаштовують місцеві сигнали біля самих виходів~\cite{Mohebi2019}. Але повільні складові є і в самих імпульсах: у дослідах наступного підрозділу плавне зростання під час наближення до винагороди видно і в частоті \nt{DOPA}-нейронів~\cite{Kim2020}. Отже, повільний рівень складається з двох джерел --- частоти імпульсів і місцевого регулювання викиду, --- і яке з них головне, схоже, залежить від виходу і від задачі.

\begin{keyblock}
\begin{proposition}[Два сигнали на одному проводі]
\label{prop:multiplex}
\nt{DOPA}-система передає щонайменше дві величини, рознесені в часі: швидку помилку $\delta$, яка вчить, і повільний рівень, який задає ціну часу~\eqref{eq:vigor} і тим самим швидкість дій. Виміряно, що швидка і повільна складові поводяться по-різному; що повільна дорівнює саме $\bar R$, --- гіпотеза.
\end{proposition}
\end{keyblock}

\section{Зростання під час наближення}

Якщо тварина біжить коридором до далекої винагороди, концентрація дофаміну в області вибору дій плавно росте протягом секунд, поки мета наближається~\cite{Hamid2016}. За розділом~\ref{sec:dofa} цього не мало б бути: усе йде за планом, і $\delta$ має дорівнювати нулю.

Перше пояснення: це зростання --- не $\delta$, а саме очікування $V$, сигнал «мета близько, варто постаратися» з попереднього підрозділу~\cite{Hamid2016}. Друге пояснення зберігає $\delta$~\cite{Kim2020}. Якщо очікування правильне, то за горизонту $\tau_\gamma$ на шляху до винагороди $R$ воно росте як $V=Re^{-(T-t)/\tau_\gamma}$, і за формулою помилки~\eqref{eq:ctd} $\dd V/\dd t-V/\tau_\gamma=0$. Але нехай здалеку очікування занижене, а в міру наближення оцінка уточнюється. Тоді очікування росте крутіше, скажімо як $V=Re^{-(T-t)/\tau_v}$ з $\tau_v<\tau_\gamma$, і
\begin{empheq}[box=\keybox]{equation}
\delta(t) \;=\; \frac{\dd V}{\dd t}-\frac{V}{\tau_\gamma} \;=\; V\Bigl(\frac{1}{\tau_v}-\frac{1}{\tau_\gamma}\Bigr) \;>\; 0 .
\label{eq:ramp}
\end{empheq}
Помилка додатна і росте разом із $V$ --- те саме плавне зростання (рис.~\ref{fig:ramp}). При $\tau_\gamma=4\,\text{с}$ і $\tau_v=1\,\text{с}$ виходить $\delta=0.75\,V$ на секунду. Цю версію перевіряли у віртуальному коридорі, миттєво переносячи тварину ближче до мети: дофамін відповідав сплеском, як і належить похідній, а не плавному рівню~\cite{Kim2020}.

\begin{figure}[t]
\centering
\begin{tikzpicture}[>=Stealth,x=1.6cm,y=2.6cm]
\draw[->,gray!70!black] (-4.2,0) -- (0.5,0) node[below left,black] {\small $t$};
\draw[->,gray!70!black] (-4.2,0) -- (-4.2,1.2);
\foreach \x/\l in {-4/{$T-4$},-2/{$T-2$},0/{$T$}}{\draw[gray!60!black] (\x,-0.03) -- (\x,0.03); \node[below,font=\scriptsize] at (\x,0) {\l\,с};}
\draw[thick,densely dashed,gray!60!black] plot[domain=-4:0,samples=40] (\x,{exp(\x/4)});
\draw[very thick,blue!60!black] plot[domain=-4:0,samples=60] (\x,{exp(\x)});
\draw[very thick,red!70!black] plot[domain=-4:0,samples=60] (\x,{0.75*exp(\x)});
\node[anchor=south east,font=\small,gray!50!black] at (-1.3,0.77) {$V$ при $\tau_v=\tau_\gamma$: $\delta=0$};
\node[right,font=\small,blue!60!black] at (0.05,1.0) {$V$ при $\tau_v=1\,\text{с}$};
\node[right,font=\small,red!70!black] at (0.05,0.72) {$\delta$ за~\eqref{eq:ramp}};
\end{tikzpicture}
\caption{Зростання дофаміну під час наближення до винагороди як помилка передбачення, $\tau_\gamma=4\,\text{с}$. Штрихова лінія --- очікування, що росте рівно так, як вимагає горизонт ($\delta=0$); синя --- очікування, що росте крутіше; червона --- помилка~\eqref{eq:ramp}.}
\label{fig:ramp}
\end{figure}

Що виміряно: плавне зростання є --- і в концентрації дофаміну, і в частоті імпульсів \nt{DOPA}-нейронів~\cite{Kim2020}, --- і миттєві стрибки обстановки спричиняють стрибки дофаміну. Яке з двох пояснень правильне --- чи обидва правильні на різних виходах, --- обговорюється.

\section{Погляд назад}

Усі теорії вище дивляться \emph{вперед}. Одна із сучасних теорій перевертає напрямок~\cite{Jeong2022}. Нехай мозок навчається не передбачати винагороду за провісником, а, отримавши винагороду, \emph{озиратися}: яка подія частіше за інші їй передувала? Міра такого зв'язку --- різниця двох імовірностей:
\begin{empheq}[box=\keybox]{equation}
M \;=\; P(\text{провісник перед винагородою}) \;-\; P(\text{провісник у випадковому вікні}) .
\label{eq:retro}
\end{empheq}
Якщо провісник часто буває і без винагороди, другий доданок великий, і зв'язок слабкий. У цій теорії дофамін повідомляє не помилку передбачення, а те, наскільки дана подія --- імовірна \emph{причина} винагороди.

Механізм озирання потребує того самого, що й правило трьох множників розділу~\ref{sec:dofa}: кожна подія має лишати слід, що тане, щоб у момент винагороди можна було прочитати, що було перед нею. Різниця не в синапсі, а в тому, що передає \nt{DOPA}.

У дослідах, де теорія~\eqref{eq:retro} і помилка~\eqref{eq:ctd} передбачають різне, викид дофаміну в низці випадків ішов за першою~\cite{Jeong2022}. Але в інших дослідах відповідь дофаміну під час навчання поступово зсувалася від моменту винагороди до моменту провісника --- як вимагає навчання за помилкою~\eqref{eq:ctd}~\cite{Amo2022}. Яка теорія описує мозок, поки невідомо.

\section{Помилка не лише у винагороді}

Останнє розширення --- про те, \emph{про що} помилка. Якщо замінити очікуваний сік іншим, так само цінним, але іншого смаку, цінність не змінилася, і помилка~\eqref{eq:ctd} дорівнює нулю. Проте \nt{DOPA}-нейрони на таку заміну відповідають~\cite{Takahashi2017}. А штучно викликані сплески дофаміну можуть навчити тварину зв'язку між двома нейтральними подразниками, взагалі без винагороди~\cite{Sharpe2017}. Схоже, \nt{DOPA} повідомляє про помилку передбачення не лише винагороди, а й того, \emph{що} станеться.

Формально це просте узагальнення~\eqref{eq:ctd}: замість одного потоку винагороди $r$ береться набір ознак того, що відбувається, $\varphi_k$ --- смак, місце, звук, --- і для кожної своє очікування $M_k$ і своя помилка~\cite{Gardner2018}:
\begin{empheq}[box=\keybox]{equation}
\delta_k(t) \;=\; \varphi_k(t) \;+\; \frac{\dd M_k}{\dd t} \;-\; \frac{M_k}{\tau_\gamma} ,
\qquad k=1,\dots,K .
\label{eq:vectortd}
\end{empheq}
Винагорода --- лише одна з ознак, і вектор помилок учить не тільки того, що добре, а й того, що за чим іде, --- моделі світу.

З цим узгоджується неоднорідність \nt{DOPA}-нейронів: в одній задачі різні нейрони несуть різні величини --- одні пов'язані з винагородою, інші з рухом, треті з тим, куди спрямована увага~\cite{Engelhard2019}.

\begin{keyblock}
\begin{proposition}[Джгут замість проводу]
\label{prop:bundle}
Сигнал \nt{DOPA}-системи --- не одне число. Він розкладений за нейронами (розподіл~\eqref{eq:expectile}, різні ознаки~\eqref{eq:vectortd}, окрема вісь небезпеки) і за часом (швидка помилка і повільний рівень~\eqref{eq:vigor}). Скалярна помилка~\eqref{eq:ctd} --- перше наближення до цього сигналу, як суматор із порогом --- перше наближення до нейрона.
\end{proposition}
\end{keyblock}

\section{Підсумок}

\begin{table}[t]
\centering
\footnotesize
\setlength{\tabcolsep}{4pt}
\renewcommand{\arraystretch}{1.15}
\begin{tabular}{>{\raggedright}p{2.5cm}>{\raggedright}p{3.2cm}>{\raggedright}p{3.6cm}>{\raggedright\arraybackslash}p{4.2cm}}
\toprule
Теорія & Що йде проводом & Головне вимірювання & Що відомо \\
\midrule
помилка передбачення (розд.~\ref{sec:dofa}) & скаляр $\delta$~\eqref{eq:ctd} & сплеск, перенесення на провісник, провал & виміряно; перше наближення \\
розподіл & набір $\delta_i$ з різною асиметрією & різні точки перевертання в різних нейронів & узгодженість із дослідом є; роль розкиду --- гіпотеза \\
важливість & $|\delta|$, новизна; вісь небезпеки & сплески на неприємне в частини нейронів & виміряно; як використовується --- обговорюється \\
ціна часу & повільний рівень $\bar R$ & дофамін пришвидшує дії; повільна складова --- і у викиді біля виходів, і в частоті імпульсів & розділення швидкого і повільного виміряно; $\bar R$ --- гіпотеза \\
зростання під час наближення & $V$ або $\delta$ при уточненні оцінки~\eqref{eq:ramp} & плавне зростання, стрибки при перенесенні в коридорі & зростання виміряно; пояснення обговорюється \\
погляд назад & міра причинного зв'язку~\eqref{eq:retro} & досліди, де передбачення двох теорій розходяться & результати суперечать одне одному \\
вектор помилок & $\delta_k$ за ознаками~\eqref{eq:vectortd} & відповідь на зміну смаку; навчання без винагороди & виміряно; векторна форма --- гіпотеза \\
\bottomrule
\end{tabular}
\caption{Що передає \nt{DOPA}-система: стандартна картина розділу~\ref{sec:dofa} та її сучасні розширення.}
\label{tab:dofaalt}
\end{table}

Теорії таблиці~\ref{tab:dofaalt} не виключають одна одну: майже всі зберігають помилку передбачення як основу і розширюють її формат. Усерйоз конкурує з нею лише погляд назад, і цю суперечку розв'яжуть досліди. Для інженера висновок простий: ціна широкомовлення з твердження~\ref{prop:broadcastcost} падає, якщо провід насправді --- багато проводів із різними адресатами.

\section{Задачі}

\begin{exercise}[\lvlA]
\label{ex:07:1}
\emph{Точки розподілу для однієї краплі.} Сік розміру $1$ надходить з імовірністю $p=0.2$, інакше винагорода $0$. Знайдіть за~\eqref{eq:expectile} $V_i$ для $\kappa_i=0.2$, $0.5$, $0.8$. Чому дорівнює медіана цієї винагороди, і чим точка~\eqref{eq:expectile} відрізняється від квантиля?
\end{exercise}

\begin{exercise}[\lvlB]
\label{ex:07:2}
\emph{Розподіл за двома нейронами.} Винагорода дорівнює $0$ або $x$, імовірність $x$ дорівнює $p$. Два нейрони з правилом~\eqref{eq:asymmetric} повідомили $V(1/2)=0.25$ і $V(0.8)=0.5$. Відновіть $p$, $x$ і стандартне відхилення винагороди. Що повідомить нейрон із $\kappa=0.2$?
\end{exercise}

\begin{exercise}[\lvlB]
\label{ex:07:3}
\emph{Моделювання: розподіл з асиметрії.} Змоделюйте дев'ять \nt{DOPA}-нейронів з $\kappa_i=0.1,\dots,0.9$ і правилом~\eqref{eq:asymmetric}, $\alpha_i^+=2\kappa_i$, $\alpha_i^-=2(1-\kappa_i)$, за винагороди, рівномірної на $[0,1]$. Як розкид $V_i$ залежить від $\eta$, і яке $\eta$ потрібне, щоб відрізнити цей розподіл від двох рівноймовірних крапель $0.25$ і $0.75$?
\end{exercise}

\begin{exercise}[\lvlA]
\label{ex:07:4}
\emph{Ціна часу.} (а)~Виведіть~\eqref{eq:vigor} і знайдіть сумарну ціну в оптимумі. (б)~Мозок оцінив $\bar R$ з помилкою в $k$ разів. Знайдіть відносну переплату при $k=2$ і $k=4$. Наскільки точно повільний рівень дофаміну має повідомляти $\bar R$?
\end{exercise}

\begin{exercise}[\lvlB]
\label{ex:07:5}
\emph{Сплеск при перенесенні.} Очікування росте за~\eqref{eq:ramp} з $\tau_v=1\,\text{с}$, $\tau_\gamma=4\,\text{с}$; тварину миттєво переносять із точки $T-2\,\text{с}$ у $T-1\,\text{с}$. Знайдіть площу сплеску $\delta$ у частках $R$ і скільки секунд рампи до перенесення він замінює. Що покаже перенесення, якщо дофамін повідомляє саме $V$?
\end{exercise}

\begin{exercise}[\lvlB]
\label{ex:07:6}
\emph{Проєктування: два сигнали на проводі.} Проводом \nt{DOPA} ідуть рівень, що росте до $s=1/60$ імпульсу на секунду за секунду, і події по $A=3$ імпульси. Синапс зі слідом $\tau_e=1\,\text{с}$ віднімає від частоти її копію, згладжену з $\tau_f$. За яких $\tau_f$ втрачається не більше $10\,\%$ події, а внесок рівня за час сліду менший за $10\,\%$ внеску події?
\end{exercise}

\begin{exercise}[\lvlB]
\label{ex:07:7}
\emph{Проєктування досліду: погляд назад проти помилки.} У вікні провісник буває з імовірністю $0.1$, після нього винагорода --- з імовірністю $0.8$. Порівняйте $\Delta P=P(\text{вин.}\mid\text{пров.})-P(\text{вин.}\mid\text{нема})$ і $M$ з~\eqref{eq:retro}, якщо (а)~додати винагороди без провісника, $0.08$ на вікно; (б)~подвоїти частоту провісника без винагород.
\end{exercise}

\begin{exercise}[\lvlC]
\label{ex:07:8}
\emph{Джгут і ціна широкомовлення.} У моделі задачі~\ref{ex:06:8} $N$ нейронів поділено на $K$ груп із власними проводами, і в кожен провід просочується частка $\rho$ чужих помилок. Покажіть, що стала навчання дорівнює $N/K+2+\rho^2N(1-1/K)$. Скільки спроб вона дає при $K\to\infty$, $N=10^{10}$ і $\rho=0.01$?
\end{exercise}

\chapter{Коли шукати, а коли вирішувати: додаємо \nt{NORA}}
\chaptermark{Додаємо \nt{NORA}}
\label{sec:nora}

\section{Чого бракує}

У розділі~\ref{sec:dofa} мережа навчалася завдяки шуму: випадкові відхилення активності були пошуком, а \nt{DOPA} повідомляв, чи стало від них краще. Але там само лишився незаданим один параметр --- \emph{скільки} шукати. Якщо шум надто сильний, мережа метається і не користується тим, що вже знає. Якщо надто слабкий, вона раз у раз вибирає те саме і не помічає, що десь поруч стало краще. У навчанні з підкріпленням це називають вибором між використанням і пошуком, і кожен алгоритм має для нього ручку. У розділі~\ref{sec:neurontypes} ми припустили, що в мозку цю ручку крутить \nt{NORA}.

Норадреналінових нейронів у людини кілька десятків тисяч (таблиця~\ref{tab:types}), майже всі вони зібрані в одній маленькій групі, а їхні виходи розходяться практично по всьому мозку, хоча різні частини цієї групи, як з'ясувалося, частково обслуговують різні області~\cite{Poe2020}. Це широкомовний канал, ще вужчий, ніж \nt{DOPA}. Працює він у двох режимах~\cite{AstonJones2005}. У \emph{фоновому} нейрони видають імпульси рівно, з частотою від часток до кількох імпульсів на секунду, і ця частота повільно змінюється разом із рівнем неспання. У \emph{уривчастому} вони відповідають короткою пачкою на важливу для поточної задачі подію, приблизно в момент ухвалення рішення. Зручна, але неточна контрольна точка --- діаметр зіниці: він змінюється разом з активністю цих нейронів, але також разом з активністю інших областей~\cite{Joshi2016} і холінергічних виходів у корі~\cite{Reimer2016}. Зіниця показує загальний рівень збудження, а не лише \nt{NORA}.

\section{Коефіцієнт підсилення}

Виміряно ось що: норадреналін пригнічує фонову активність нейронів сильніше, ніж їхню відповідь на стимул, і відношення сигналу до фону зростає~\cite{Foote1975NE}. Що крутість характеристики окремого нейрона при цьому буквально множиться на коефіцієнт, з досліду не випливає. Як і в розділі~\ref{sec:neurontypes}, візьмімо просту модель, яка передає цей ефект: норадреналін змінює крутість передавальної характеристики отримувача~\cite{ServanSchreiber1990}. Слабкі, підпорогові входи (фон) тоді дають ще менше, а сильні --- ще більше. Нехай нейрон відповідає частотою
\begin{empheq}[box=\keybox]{equation}
r \;=\; F\bigl(g\,(u-\theta)\bigr), \qquad F(z)=\frac{r_{\max}}{1+e^{-z}} ,
\label{eq:gain}
\end{empheq}
де $u$ --- вхідний струм, $\theta$ --- поріг, а $g$ --- коефіцієнт підсилення, яким у моделі керує \nt{NORA}. За малого $g$ частота плавно йде за входом у широкому діапазоні: нейрон --- \emph{аналоговий підсилювач}. За великого $g$ майже будь-який вхід вище порога дає повну частоту, а нижче --- нуль: нейрон --- \emph{компаратор}, майже цифровий елемент (рис.~\ref{fig:gaincurves}). Одна ручка перемикає той самий нейрон між двома режимами, які в електроніці виконують різні схеми.

\begin{figure}[t]
\centering
\begin{tikzpicture}[>=Stealth,x=1.1cm,y=2.4cm]
\draw[->,gray!70!black] (-3.3,0) -- (3.4,0) node[below left,black] {\small $u$};
\draw[->,gray!70!black] (-3.3,0) -- (-3.3,1.2) node[left,black] {\small $r$};
\draw[densely dashed,gray!60!black] (0,0) -- (0,1.1);
\node[below,font=\small] at (0,0) {$\theta$};
\draw[densely dotted,gray!60!black] (-3.3,1) -- (3.2,1) node[right,black,font=\small] {$r_{\max}$};
\draw[very thick,blue!60!black] plot[domain=-3.2:3.2,samples=60] (\x,{1/(1+exp(-0.8*\x))});
\draw[very thick,orange!85!black] plot[domain=-3.2:3.2,samples=60] (\x,{1/(1+exp(-2*\x))});
\draw[very thick,red!70!black] plot[domain=-3.2:3.2,samples=120] (\x,{1/(1+exp(min(-8*\x,9)))});
\node[blue!60!black,font=\small,anchor=west] at (0.5,0.12) {мале $g$: підсилювач};
\node[red!70!black,font=\small,anchor=east] at (-0.3,0.85) {велике $g$: компаратор};
\end{tikzpicture}
\caption{Передавальна характеристика~\eqref{eq:gain} за трьох значень коефіцієнта підсилення $g$.}
\label{fig:gaincurves}
\end{figure}

Чи допомагає велике підсилення проти шуму? Не завжди. Нехай два входи відрізняються на $\Delta u$, і треба сказати, який більший. Якщо шум $\sigma_{\text{до}}$ доданий до входу \emph{до} підсилювача, то підсилюється і сигнал, і шум, і їхнє відношення не змінюється. Якщо ж шум $\sigma_{\text{після}}$ додається \emph{після} підсилювача --- від ненадійних синапсів і випадкових імпульсів самої мережі, як у розділі~\ref{sec:code}, --- то відношення сигналу до шуму
\begin{empheq}[box=\keybox]{equation}
d' \;=\; \frac{g\,\Delta u}{\sqrt{g^2\sigma_{\text{до}}^2+\sigma_{\text{після}}^2}}
\label{eq:gainsnr}
\end{empheq}
зростає з $g$, доки $g\sigma_{\text{до}}$ не зрівняється з $\sigma_{\text{після}}$~\cite{ServanSchreiber1990}.

\begin{keyblock}
\begin{proposition}[Коли підсилення допомагає]
\label{prop:gainnoise}
Збільшуючи коефіцієнт підсилення, \nt{NORA} покращує розрізнення входів лише проти того шуму, що виникає \emph{після} підсилювача, у самій мережі. Шуму, що прийшов разом із входом, підсилення не прибере.
\end{proposition}
\end{keyblock}

\section{Підсилення у виборі з двох}

Повернімося до вибору з розділу~\ref{sec:gaba}: дві групи \nt{GLUT}-нейронів із входами $I_1$, $I_2$ гальмують одна одну із силою $\beta$. Тепер кожна група має коефіцієнт підсилення $g$: у рівняннях вибору~\eqref{eq:wta} вираз у дужках множиться на $g$. Поки працюють обидві групи, усталені частоти знаходимо з $r_1=g(I_1-\beta r_2)$, $r_2=g(I_2-\beta r_1)$. Додавши і віднявши ці рівності, отримаємо суму $r_1+r_2=g(I_1+I_2)/(1+g\beta)$ і різницю $r_1-r_2=g(I_1-I_2)/(1-g\beta)$. Їхнє відношення --- наскільки різко мережа віддає перевагу одному входу над іншим:
\begin{empheq}[box=\keybox]{equation}
\frac{r_1-r_2}{r_1+r_2} \;=\; \varepsilon\,\frac{1+g\beta}{1-g\beta}, \qquad \varepsilon=\frac{I_1-I_2}{I_1+I_2}, \quad g\beta<1 .
\label{eq:wtagain}
\end{empheq}
Мала різниця входів $\varepsilon$ підсилюється в $(1+g\beta)/(1-g\beta)$ раза, і цей множник необмежено зростає, коли $g\beta$ наближається до одиниці. При $g\beta>1$ стан, де працюють обидві групи, нестійкий, як у твердженні~\ref{prop:wta}: перемагає одна, інша мовчить (рис.~\ref{fig:pitchfork}).

\begin{figure}[t]
\centering
\begin{tikzpicture}[>=Stealth,x=3.2cm,y=1.5cm]
\draw[->,gray!70!black] (0,0) -- (2.15,0) node[below left,black] {\small $g\beta$};
\draw[->,gray!70!black] (0,-1.25) -- (0,1.35) node[above,black] {\small $(r_1-r_2)/(r_1+r_2)$};
\draw[gray!60!black] (1,-0.04) -- (1,0.04); \node[below right,font=\small] at (1,0) {$1$};
\node[left,font=\scriptsize] at (0,1) {$1$}; \node[left,font=\scriptsize] at (0,-1) {$-1$};
\draw[very thick,blue!60!black] plot[domain=0:0.905,samples=60] (\x,{0.05*(1+\x)/(1-\x)}) -- (1,1) -- (2.05,1);
\draw[very thick,blue!60!black] (1.105,-1) -- (2.05,-1);
\draw[thick,densely dashed,gray!60!black] (1,0) -- (2.05,0);
\node[font=\small,align=center] at (0.45,-0.62) {розмірковує:\\обидві групи\\працюють};
\node[font=\small,align=center] at (1.55,0.45) {вирішила:\\працює одна};
\node[font=\small,blue!60!black,anchor=south west] at (1.05,-0.97) {слабкий вхід переміг раніше --- тримається};
\end{tikzpicture}
\caption{Вибір із двох за різного коефіцієнта підсилення, входи відрізняються на $\varepsilon=0.05$. При $g\beta>1$ стійкі обидва рішення (суцільні лінії; перемога слабкого входу --- при $g\beta>(1+\varepsilon)/(1-\varepsilon)\approx1.1$), і мережа тримає те, яке вже ухвалила; симетричний стан (штрихова лінія) нестійкий.}
\label{fig:pitchfork}
\end{figure}

\begin{keyblock}
\begin{proposition}[Підсилення перемикає режим вибору]
\label{prop:gainwta}
У мережі вибору із взаємним гальмуванням $\beta$ коефіцієнт підсилення $g$ переводить мережу через межу $g\beta=1$. Нижче неї мережа «розмірковує»: обидві групи працюють, а різниця входів підсилюється за~\eqref{eq:wtagain}. Вище --- мережа «вирішила»: працює одна група, і рішення тримається. Змінюючи $g$, \nt{NORA} може перемикати мережу між цими режимами, не чіпаючи жодної ваги.
\end{proposition}
\end{keyblock}

\section{Підсилення як температура}

Тепер додамо до вибору шум, що виникає в самій мережі, після підсилювача. Яка група переможе, вирішує знак величини $g(u_A-u_B)+\eta$, де $u_A-u_B$ --- різниця входів, тобто різниця вивчених ваг із розділу~\ref{sec:dofa}, а $\eta$ --- шум із масштабом $\sigma$. Якщо шум розподілений за логістичним законом (для гаусового крива майже та сама), імовірність вибрати $A$ дорівнює
\begin{empheq}[box=\keybox]{equation}
P(A) \;=\; \frac{1}{1+e^{-g\,(u_A-u_B)/\sigma}} .
\label{eq:softmax}
\end{empheq}
Програміст упізнає тут вибір за softmax із температурою $T=\sigma/g$. Коефіцієнт підсилення --- обернена температура. За малого $g$ навіть помітно кращий варіант вибирається ненабагато частіше за гірший: мережа багато шукає. За великого $g$ найкращий відомий варіант вибирається майже завжди: мережа користується тим, що знає.

Уточнімо, що таке $g$ у~\eqref{eq:wtagain} і~\eqref{eq:softmax}: це \emph{діюче} підсилення різниці входів поточної задачі в момент вибору. Два режими \nt{NORA} діють на нього протилежно. Уривчаста пачка в момент рішення підвищує його, і мережа закріплює вибір. Висока фонова активність, навпаки, іде разом із пошуком: у людей перед вибором навмання базова зіниця ширша~\cite{JepmaNieuwenhuis2011}, а в щурів вхід \nt{NORA} у передню поясну кору переводить вибір у випадковий режим~\cite{Tervo2014stochastic}. Отже, високій фоновій \nt{NORA} відповідає \emph{мале} діюче $g$, а пачці --- велике.

\section{Скільки шукати}

Яке підсилення краще? Ми перевірили це на моделі. Мережа вибирає одну з двох дій за~\eqref{eq:softmax}; кожна приносить винагороду з певною ймовірністю, і мережа оцінює її за винагородами вибраної дії, як у розділі~\ref{sec:dofa}. Час від часу світ змінюється: обидві ймовірності розігруються заново. Змінитися може і вибрана дія, і та, яку мережа давно не пробувала; про другу мережа може дізнатися лише пошуком. На рис.~\ref{fig:invertedU} показано, яку частку найкращої можливої винагороди отримує мережа за різних сталих $g$.

\begin{figure}[t]
\centering
\begin{tikzpicture}[>=Stealth,x=1.2cm,y=1cm]
\draw[->,gray!70!black] (0,0) -- (7.6,0) node[below left,black] {\small $g$};
\draw[->,gray!70!black] (0,0) -- (0,4.4) node[above,black,font=\small] {частка найкращої винагороди};
\foreach \x/\l in {0/0.5,1/1,2/2,3/4,4/8,5/16,6/32,7/64}{\draw[gray!60!black] (\x,-0.06) -- (\x,0.06); \node[below,font=\scriptsize] at (\x,0) {\l};}
\foreach \v/\y in {0.8/0.8,0.9/2.4,1.0/4.0}{\draw[gray!60!black] (-0.06,\y) -- (0.06,\y); \node[left,font=\scriptsize] at (0,\y) {\v};}
\draw[very thick,blue!60!black] plot[mark=*,mark size=1.3pt] coordinates {(0,0.368) (1,0.880) (2,1.728) (3,2.784) (4,3.456) (5,3.616) (6,3.488) (7,3.264)};
\draw[very thick,red!70!black] plot[mark=*,mark size=1.3pt] coordinates {(0,0.368) (1,0.672) (2,1.184) (3,1.840) (4,2.176) (5,1.920) (6,1.456) (7,1.104)};
\node[blue!60!black,font=\small,anchor=south] at (5,3.7) {спокійний світ};
\node[red!70!black,font=\small,anchor=north] at (5.1,1.25) {світ, що швидко змінюється};
\draw[->,thick,blue!60!black] (5,3.25) -- (5,3.55);
\draw[->,thick,red!70!black] (4,1.8) -- (4,2.1);
\node[font=\small\itshape,gray!35!black,anchor=west] at (0.15,2.3) {шукає наосліп};
\node[font=\small\itshape,gray!35!black] at (6.45,2.55) {не помічає змін};
\end{tikzpicture}
\caption{Частка найкращої можливої винагороди за сталого коефіцієнта підсилення $g$ (шкала за $g$ логарифмічна). Синя крива --- світ змінюється в середньому раз на тисячу спроб, червона --- раз на двадцять. Стрілки позначають найкраще $g$: у світі, що швидко змінюється, воно вдвічі менше. Середнє за 60 прогонами по 4000 спроб.}
\label{fig:invertedU}
\end{figure}

При $g=0.5$ мережа майже не користується знанням і отримує близько трьох чвертей можливого, за дуже великого $g$ пропускає зміни. Найкраще значення: $g=16$, коли світ змінюється раз на тисячу спроб (частка $0.976$), і $g=8$, коли раз на двадцять ($0.886$).

\begin{keyblock}
\begin{proposition}[Перевернуте U]
\label{prop:explore}
Винагорода, яку збирає мережа, залежить від коефіцієнта підсилення як перевернуте U: надто мале $g$ витрачає спроби на пошук, надто велике не помічає змін. Найкраще $g$ тим менше, чим швидше змінюється світ.
\end{proposition}
\end{keyblock}

Перевернуте U спостерігається і в дослідах: тварини розв'язують задачу найкраще за помірної фонової активності \nt{NORA}-нейронів із чіткими уривчастими відповідями, а за високої фонової активності відволікаються і перемикаються між задачами~\cite{AstonJones2005}. Це узгоджується з різницею режимів із попереднього підрозділу: уривчаста пачка в момент рішення дає велике діюче $g$ саме тоді, коли треба вибрати, а висока фонова активність без пачок підсилює все підряд, зокрема й сторонні входи, тож діюче $g$ для поточної задачі падає, --- це пошук.

\section{Хто крутить ручку}

Якщо найкраще підсилення залежить від того, як швидко змінюється світ, його добре було б підлаштовувати. Але звідки \nt{NORA}-нейронам знати, що світ змінився? Природна ознака --- \emph{несподіваність}: помилки передбачення стали більшими, ніж зазвичай. Важливо розрізняти два види невизначеності. Якщо винагорода просто випадкова, помилки великі завжди, і це очікувано. Якщо ж помилки раптом зросли порівняно з їхньою звичайною величиною, отже, щось змінилося. За однією з гіпотез, саме цю неочікувану невизначеність і повідомляє \nt{NORA}, а очікувану --- \nt{ACET}~\cite{YuDayan2005}.

Що робити з цим сигналом? Можна знизити підсилення і більше шукати. Можна підвищити швидкість навчання, щоб швидше забути застарілі оцінки: досліди показують, що після різкої зміни люди навчаються швидше, і зіниця при цьому розширюється~\cite{Nassar2012}; нагадаємо, що зіниця --- показник не лише \nt{NORA}, а й \nt{ACET}~\cite{Reimer2016}. А можна зробити і те, й інше одразу: за ще однією гіпотезою, уривчаста пачка \nt{NORA} працює як \emph{переривання} --- сигнал, за яким мережа скидає поточний стан і перебудовується під нову обстановку~\cite{BouretSara2005}, як спільна лінія переривання процесора.

Чесно скажемо, чого ми не знайшли. У моделі ми пробували обидва прості способи підлаштування: знижувати $g$ або підвищувати швидкість навчання, коли згладжена помилка перевищує своє довготривале середнє. У світі, який то довго стоїть, то часто змінюється, найкращі налаштування обох способів дали $0.926$--$0.927$ від найкращої винагороди --- майже стільки ж, скільки найкраще стале $g$ ($0.925$ при $g=8$) або стала, але досить велика швидкість навчання ($0.928$), а невдалі налаштування --- менше. Криві на рис.~\ref{fig:invertedU} біля вершини пологі, і в такому світі підлаштовувати майже нічого. Підлаштування має окупатися там, де різні режими вимагають дуже різних налаштувань. Який саме закон керування реалізує \nt{NORA}, поки не встановлено.

\section{Підсумок}

\begin{table}[t]
\centering
\footnotesize
\setlength{\tabcolsep}{4pt}
\renewcommand{\arraystretch}{1.15}
\begin{tabular}{>{\raggedright}p{3.0cm}>{\raggedright}p{4.4cm}>{\raggedright\arraybackslash}p{6.0cm}}
\toprule
Що робить \nt{NORA} & Як & Що відомо \\
\midrule
змінює крутість відповіді нейронів & коефіцієнт підсилення $g$ у~\eqref{eq:gain} & підвищення відношення сигналу до шуму виміряно; мультиплікативна модель --- спрощення \\
перемикає вибір & переводить мережу через $g\beta=1$ & випливає з моделі~\eqref{eq:wtagain}; прямих вимірювань порога немає \\
задає, скільки шукати & $g$ --- обернена температура~\eqref{eq:softmax}; фон --- мале діюче $g$ (пошук), пачка --- велике (рішення) & перевернуте U і два режими роботи виміряно; зв'язок із пошуком --- добре обґрунтована гіпотеза \\
повідомляє про зміни & неочікувана невизначеність & зіниця і швидкість навчання зростають після змін --- виміряно; що це сигнал \nt{NORA} --- гіпотеза \\
перериває і скидає & уривчаста пачка на несподівану подію & пачки на важливі події виміряно; «переривання» --- гіпотеза \\
\bottomrule
\end{tabular}
\caption{Що додає \nt{NORA} до мережі з \nt{GLUT}, \nt{GABA} і \nt{DOPA}.}
\label{tab:nora}
\end{table}

\nt{DOPA} каже мережі, \emph{куди} змінювати ваги. \nt{NORA} не змінює жодної ваги, але вирішує, \emph{як} мережа користується тим, що вже знає (таблиця~\ref{tab:nora}): одна ручка --- коефіцієнт підсилення --- перетворює нейрон із підсилювача на компаратор, мережу вибору --- з тієї, що «розмірковує», на ту, що «вирішила», а вибір --- з пошуку на використання. Як \nt{NORA} дізнається, наскільки швидко змінюється світ, --- відкрите питання.

\section{Задачі}

\begin{exercise}[\lvlA]
\label{ex:08:1}
\emph{Одна ручка на весь мозок.} (а)~За таблицею~\ref{tab:types} оцініть, скільки нейронів мозку припадає на один \nt{NORA}-нейрон. (б)~Шум до підсилювача $\sigma_{\text{до}}=0.5$, після --- $\sigma_{\text{після}}=4$. За якого $g$ величина $d'$ з~\eqref{eq:gainsnr} досягає $90\,\%$ своєї границі?
\end{exercise}

\begin{exercise}[\lvlB]
\label{ex:08:2}
\emph{Вибір із коефіцієнтом підсилення.} $\tau\,\dd r_1/\dd t=-r_1+g[I_1-\beta r_2]_+$, $\tau\,\dd r_2/\dd t=-r_2+g[I_2-\beta r_1]_+$, $\tau=20\ms$. (а)~За який час загасає $r_1-r_2$ при $g\beta=0.5$ і $0.9$? (б)~При $\varepsilon=0.05$ знайдіть $g\beta$, нижче якого працюють обидві групи і вище якого стійка перемога слабкого входу.
\end{exercise}

\begin{exercise}[\lvlA]
\label{ex:08:3}
\emph{Softmax із шуму.} Кожна з $K$ груп отримує свій шум Гумбеля, $P(\eta_k<z)=\exp(-e^{-z/\sigma})$, і перемагає найбільше $gu_k+\eta_k$. Знайдіть $P(k)$. За якого $g$ кращий із двох варіантів з $u_A-u_B=0.1$, $\sigma=1$, вибирається в $90\,\%$ випадків?
\end{exercise}

\begin{exercise}[\lvlB]
\label{ex:08:4}
\emph{Оцінка найкращого підсилення.} У моделі рис.~\ref{fig:invertedU} імовірності винагород після зміни рівномірні на $[0,1]$. Мережа пробує гіршу дію раз на $e^{g\Delta}$ спроб; прирівнявши це до $1/h$, отримаємо $g^*\approx\ln(1/h)/\bar\Delta$. Знайдіть $\bar\Delta=\langle|p_A-p_B|\rangle$ і $g^*$ при $h=0.001$, $0.01$, $0.05$.
\end{exercise}

\begin{exercise}[\lvlB]
\label{ex:08:5}
\emph{Моделювання: підсилення і швидкість змін.} Знайдіть $g^*(h)$ при $h$ від $0.001$ до $0.2$ копією \texttt{models/\allowbreak nora\_explore.py} (вона пише файли в поточну теку). Де правильна оцінка задачі~\ref{ex:08:4} і чому за швидких змін $g^*$ перестає спадати?
\end{exercise}

\begin{exercise}[\lvlB]
\label{ex:08:6}
\emph{Проєктування: переривання для мережі вибору.} Мережа задачі~\ref{ex:08:2} з $\beta=2$, $\tau=20\ms$, $g=1$ тримає $r_1=0.9$, $r_2=0$, хоча входи тепер $I_1=0.9$, $I_2=1.0$. \nt{NORA} на час $T_p$ знижує $g$ до $g_{\text{lo}}$. Знайдіть найменше $T_p$ при $g_{\text{lo}}=0$ аналітично і моделлю при $g_{\text{lo}}=0.25$.
\end{exercise}

\begin{exercise}[\lvlC]
\label{ex:08:7}
\emph{Коли підлаштування окупається.} Оракул знає, спокійна епоха чи мінлива, і ставить у кожній своє $g$. (а)~Моделлю знайдіть його виграш над найкращим сталим $g$ у світі розділу (епохи по $1000$ спроб, $h=0.001$ і $0.05$). (б)~Побудуйте світ, де виграш більший, і знайдіть, яку його частину забирає підлаштування за несподіваністю.
\end{exercise}

\chapter{Коли записувати, а коли читати: додаємо \nt{ACET}}
\chaptermark{Додаємо \nt{ACET}}
\label{sec:acet}

\section{Чого бракує}

Мікросхема пам'яті, крім проводів адреси і даних, має ще одну лінію --- дозвіл запису. Поки вона вимкнена, пам'ять лише читається; коли ввімкнена, те, що є на лініях даних, записується. Без такої лінії пам'ять не працює: якщо кожне читання водночас щось записує, прочитане, зокрема прочитане з помилкою, одразу ж записується назад.

У мозку ця задача постає гостріше, ніж у мікросхемі. У розділі~\ref{sec:glutcomputer} пам'яттю була група \nt{GLUT}-нейронів, яку тримають увімкненою її власні зв'язки (рис.~\ref{fig:bistable}). У розділі~\ref{sec:dofa} ми бачили, що ці зв'язки навчаються правилом Гебба просто під час роботи. Отже, ті самі синапси і зберігають старе, і приймають нове, а окремої фази запису, як і годинника, немає. Чим мережа відрізняє момент, коли треба запам'ятати те, що бачиш, від моменту, коли треба згадати те, що знаєш? У розділі~\ref{sec:neurontypes} ми припустили, що цю лінію тримає \nt{ACET}. У цьому розділі ми перевіримо, як така лінія має працювати і чому без неї не обійтися.

\section{Три дії одного проводу}

Ацетилхолінових нейронів, що працюють усередині мозку, небагато, і зібрані вони в кількох невеликих групах, а їхні виходи розходяться по всій корі. Це широкомовний канал, такий самий, як \nt{DOPA} і \nt{NORA}. Рівень ацетилхоліну в корі високий в уважному неспанні і низький у повільному сні~\cite{Hasselmo1999}. Як і в дофаміну, зміст сигналу задає рецептор отримувача (твердження~\ref{prop:receptor}): ацетилхолін має і швидкі рецептори, що відкривають канал, і повільні, що запускають реакції всередині клітини. Для нас важливі три дії.

\emph{Перша: послабити внутрішні зв'язки.} Досліди на зрізах нюхової кори показали, що ацетилхолін сильно пригнічує передачу зв'язками між нейронами однієї області і майже не зачіпає входи, що надходять у неї ззовні~\cite{HasselmoBower1992}.

\emph{Друга: посилити входи.} Ацетилхолін підвищує збудливість нейронів і полегшує передачу деякими вхідними шляхами; сигнал від органів чуттів стає гучнішим за власну активність мережі~\cite{Hasselmo2006}.

\emph{Третя: полегшити зміну ваг.} За високого рівня ацетилхоліну зв'язки змінюються легше~\cite{Hasselmo2006}.

Для інженера всі три дії --- одна операція: перемикання з режиму «читати» в режим «записувати». Мережа перестає слухати саму себе, починає слухати вхід і запам'ятовує те, що чує.

\section{Мережа, яка дописує}

Візьмімо $N$ \nt{GLUT}-нейронів, з'єднаних один з одним. Збережімо в їхніх зв'язках кілька образів --- наборів по $pN$ одночасно активних нейронів --- правилом Гебба~\cite{Hebb1949}. Якщо подати на таку мережу половину образу, зв'язки збудять відсутню половину: мережа \emph{дописує} образ за частиною, як група на рис.~\ref{fig:bistable} підтримувала сама себе. Це асоціативна пам'ять: адресою слугує частина вмісту. \nt{GABA}, як у розділі~\ref{sec:gaba}, тримає загальну кількість активних нейронів близько $pN$, щоб разом не могли засвітитися два образи.

Позначмо рівень ацетилхоліну $a$, від $0$ до $1$, і запишімо три його дії в модель прямо:
\begin{empheq}[box=\keybox]{equation}
\tau\frac{\dd r_i}{\dd t} = -r_i + F\Bigl(\tfrac{1+a}{2}\,x_i + (1-a)\sum_j W_{ij}\,r_j - g\Bigr),
\qquad
\frac{\dd W_{ij}}{\dd t} = \eta\,a\,(r_i-p)(r_j-p) .
\label{eq:acetnet}
\end{empheq}
Тут $r_i$ --- частота нейрона $i$, $x_i$ --- його вхід від органів чуттів, $F$ --- порогова передавальна характеристика, $g$ --- загальне гальмування від \nt{GABA}, що зростає, коли активних нейронів більше за $pN$. Множник $\frac{1+a}{2}$ --- посилення входу, $1-a$ --- послаблення внутрішніх зв'язків, $\eta a$ --- швидкість навчання. Друге рівняння --- правило Гебба~\eqref{eq:hebb} у формі, зручній для рідкісних образів~\cite{Tsodyks1988}: вага росте, коли обидва нейрони активніші, ніж зазвичай, і зменшується, коли активний лише один. За від'ємною частиною ваг, як завжди, стоять \nt{GABA}-посередники. Годинника немає: і нейрони, і ваги змінюються неперервно.

\begin{figure}[t]
\centering
\begin{tikzpicture}[>=Stealth,
  nrn/.style={circle,draw,very thick,fill=blue!6,minimum size=0.8cm,inner sep=0pt,font=\small},
  acet/.style={circle,draw,very thick,double,double distance=1.2pt,fill=orange!15,minimum size=1.35cm,inner sep=0pt,font=\scriptsize},
  bc/.style={->,thick,densely dashed,orange!85!black}]
\foreach \i/\y in {1/2.4,2/1.2,3/0}{
  \node[nrn] (N\i) at (3,\y) {$r_\i$};
  \draw[->,thick,blue!60!black] (0,\y) node[left,black] {\small $x_\i$} -- (N\i);}
\node[left,align=right,font=\small] at (-0.5,1.2) {від органів\\чуттів};
\draw[->,thick,gray!60!black] (N1) to[bend left=30] (N2);
\draw[->,thick,gray!60!black] (N2) to[bend left=30] (N1);
\draw[->,thick,gray!60!black] (N2) to[bend left=30] (N3);
\draw[->,thick,gray!60!black] (N3) to[bend left=30] (N2);
\draw[->,thick,gray!60!black] (N1.east) .. controls (4.6,2.0) and (4.6,0.4) .. (N3.east);
\node[right,font=\small,gray!40!black,align=left] at (4.7,1.2) {внутрішні\\зв'язки $W$};
\node[acet] (A) at (1.2,-1.9) {\nt{ACET}};
\draw[bc] (A) -- (1.6,-0.15);
\node[font=\small,orange!60!black,anchor=east] at (0.95,-0.9) {$+$ вхід};
\draw[bc] (A) -- (4.3,0.6);
\node[font=\small,orange!60!black,anchor=west] at (4.3,0.1) {$-$ внутрішні зв'язки, $+$ швидкість навчання};
\node[right,font=\small,align=left] at (2.2,-1.9) {високий рівень: «записувати»\\низький рівень: «читати»};
\end{tikzpicture}
\caption{\nt{ACET} як лінія дозволу запису. Нейрони $r_i$ отримують вхід $x_i$ від органів чуттів (сині стрілки) і збуджують один одного через внутрішні зв'язки $W$ (сірі), у яких зберігаються образи. Ацетилхолін (штрихові стрілки) посилює вхід, послаблює внутрішні зв'язки і полегшує зміну ваг, як у~\eqref{eq:acetnet}.}
\label{fig:acetnet}
\end{figure}

\section{Запис і читання заважають одне одному}

Перевірмо модель~\eqref{eq:acetnet} на задачі, в якій запис і читання справді конфліктують. Мережа з $200$ нейронів уже зберігає п'ять образів по $20$ активних нейронів. Тепер треба запам'ятати шостий, який наполовину збігається з одним зі старих: десять нейронів у них спільні. Так буває постійно: нова кімната схожа на знайому, нове обличчя --- на старе.

\emph{Запис за низького $a$.} Спочатку відокремимо перші дві дії ацетилхоліну від третьої: залишимо швидкість навчання повною, а змінюватимуться лише посилення входу і послаблення внутрішніх зв'язків. При $a=0$ вхід запалює спільні десять нейронів, а вони через внутрішні зв'язки запалюють решту половини \emph{старого} образу. \nt{GABA} не дає горіти обом одразу, і старий образ, підтриманий власними зв'язками, витісняє новий, який тримається лише на вході. Мережа бачить не те, що перед нею, а те, що пам'ятає, --- і правило Гебба старанно записує побачене.

У підсумку новий образ узагалі не запам'ятався: за його відмінною половиною мережа нічого не згадує. А старий зіпсований: викликаний своєю відмінною половиною, він тягне за собою в середньому шість чужих нейронів із десяти. При $a=0.1$ новий образ уже запам'ятовується, але зливається зі старим, і псування навіть сильніше: кожен із двох, викликаний своєю половиною, тягне за собою близько восьми чужих нейронів; починаючи з $a=0.2$ запис чистий: зайвих нейронів під час згадування менше за один (усереднення за десятьма прогонами).

\emph{Запис зі справжньою швидкістю навчання.} Тепер нехай ацетилхолін, як у~\eqref{eq:acetnet}, задає і швидкість навчання. За одне пред'явлення новий образ запам'ятовується цілком лише при $a\ge0.9$; при $a=0.75$ --- на вісім десятих, при $a\le0.5$ не запам'ятовується зовсім.

\emph{Читання.} Подамо на мережу з п'ятьма чисто записаними образами половину образу, а потім приберемо і її. При $a\le0.2$ мережа дописує образ цілком і тримає його сама; при $a=0.3$ --- майже цілком; при $a\ge0.4$ внутрішні зв'язки надто послаблені, і після того як підказка зникла, активність гасне (рис.~\ref{fig:acetsim}).

\begin{figure}[t]
\centering
\begin{tikzpicture}[>=Stealth,x=9cm,y=3.2cm]
\draw[->,gray!70!black] (0,0) -- (1.1,0) node[right,black] {\small $a$};
\draw[->,gray!70!black] (0,0) -- (0,1.2);
\foreach \x in {0,0.2,0.4,0.6,0.8,1.0}{\draw[gray!60!black] (\x,-0.02) -- (\x,0.02); \node[below,font=\scriptsize] at (\x,0) {$\x$};}
\foreach \y in {0,0.5,1}{\draw[gray!60!black] (-0.01,\y) -- (0.01,\y); \node[left,font=\scriptsize] at (0,\y) {$\y$};}
\node[rotate=90,font=\small] at (-0.1,0.6) {частка відновленого образу};
\draw[very thick,blue!60!black] plot[mark=*,mark size=1.6pt] coordinates {(0,1.00) (0.1,1.00) (0.2,1.00) (0.3,0.96) (0.4,0.04) (0.5,0.00) (0.75,0.00) (1.0,0.00)};
\draw[very thick,red!70!black] plot[mark=*,mark size=1.6pt] coordinates {(0.1,0.00) (0.2,0.00) (0.3,0.00) (0.5,0.00) (0.75,0.80) (0.9,1.00) (1.0,1.00)};
\node[font=\small,blue!60!black,anchor=west,align=left] at (0.03,0.5) {читання:\\образ\\за половиною};
\node[font=\small,red!70!black,anchor=east,align=right] at (1.0,0.35) {запис:\\новий образ\\за раз};
\end{tikzpicture}
\caption{Модель~\eqref{eq:acetnet}: $200$ нейронів, образи по $20$ активних, середнє за десятьма прогонами. Сині точки --- читання: яку частку образу відновлено за його половиною після того, як підказку прибрано. Червоні точки --- запис: яка частка нового образу, наполовину схожого на старий, згадується після одного пред'явлення, якщо ацетилхолін задає і швидкість навчання. Читання працює при $a\le0.3$, запис --- при $a\ge0.9$.}
\label{fig:acetsim}
\end{figure}

\begin{keyblock}
\begin{proposition}[Запис і читання потребують різних режимів]
\label{prop:writeread}
У моделі~\eqref{eq:acetnet}, де ті самі зв'язки зберігають старе і навчаються нового, а ацетилхолін задає і швидкість навчання, немає рівня $a$, за якого однаково добре йдуть і запис, і читання. Для запису внутрішні зв'язки треба послабити, інакше мережа запише не вхід, а згадане; для читання їх треба посилити, інакше вона не допише образ за частиною. Потрібен перемикач, і ним може слугувати один широкомовний провід.
\end{proposition}
\end{keyblock}

Якби ацетилхолін не змінював швидкості навчання, для обох справ годилася б вузька смуга $0.2\le a\le0.3$. Але тоді мережа навчалася б і під час кожного читання, тобто заново записувала б прочитане разом із його помилками. Саму ідею --- пригнічувати внутрішні зв'язки під час навчання --- взято з моделей, побудованих за вимірюваннями на зрізах~\cite{HasselmoAndersonBower1992}. Числа вище стосуються нашої спрощеної моделі; де саме лежать межі режимів у мозку, невідомо.

\section{Наскільки вірити очам}

Перемикач «записувати/читати» --- крайній випадок загальнішого питання: наскільки вірити входу і наскільки --- пам'яті? На це питання інженери мають точну відповідь, і вона пояснює, чому один провід керує водночас і режимом, і швидкістю навчання.

Нехай мережа стежить за величиною $s(t)$, яка повільно блукає: за секунду її дисперсія зростає на $q$. Органи чуттів повідомляють $y(t)=s(t)+$шум з інтенсивністю $R$. Найкращу оцінку $\hat s$ у неперервному часі дає фільтр Калмана--Б'юсі~\cite{KalmanBucy1961}:
\begin{empheq}[box=\keybox]{equation}
\frac{\dd\hat s}{\dd t} \;=\; \frac{P}{R}\,\bigl(y-\hat s\bigr),
\qquad
\frac{\dd P}{\dd t} \;=\; q-\frac{P^2}{R} ,
\label{eq:kalman}
\end{empheq}
де $P$ --- дисперсія помилки оцінки, тобто наскільки мережа сама не впевнена в тому, що знає. Перше рівняння --- те саме RC-коло, що й у мембрани в моделі нейрона~\eqref{eq:glutneuron}, але зі сталою часу $R/P$, яка змінюється. Коли мережа не впевнена, $P$ велике, стала часу мала, і оцінка швидко йде за входом: це «записувати». Коли впевнена, $P$ мале, і оцінка майже не слухає входу, тримаючись за пам'ять: це «читати».

В усталеному режимі $P=\sqrt{qR}$ і стала часу пам'яті дорівнює $\sqrt{R/q}$: якщо світ відходить на одиницю за сто секунд, $q=0.01$, а шум органів чуттів $R=1$, пам'ять має пам'ятати близько десяти секунд.

Головне тут те, що одне число $P/R$ задає водночас дві речі: і вагу входу порівняно з пам'яттю, і швидкість, з якою змінюється збережена оцінка. Це саме те, що робить ацетилхолін у~\eqref{eq:acetnet}. За гіпотезою~\cite{YuDayan2005}, ацетилхолін і повідомляє мережі величину, подібну до $P$, --- \emph{очікувану} невизначеність, ту, про яку мережа знає заздалегідь. Повільним цей канал буває не завжди: частина \nt{ACET}-нейронів відповідає на винагороду і покарання за два десятки мілісекунд, тим сильніше, чим несподіваніше підкріплення~\cite{Hangya2015}.

\begin{keyblock}
\begin{proposition}[Одна ручка для ваги входу і швидкості навчання]
\label{prop:kalman}
В оптимальному фільтрі~\eqref{eq:kalman} вага, з якою вхід змішується з пам'яттю, і швидкість навчання --- та сама величина $P/R$. Тому розумно керувати ними одним сигналом. За незмінних властивостей світу цей сигнал виходить на сталий рівень $\sqrt{q/R}$, і підлаштовувати його треба, лише коли властивості світу змінюються.
\end{proposition}
\end{keyblock}

Друге речення твердження пояснює результат розділу~\ref{sec:nora}, де підлаштування швидкості навчання за несподіваністю не виграло в найкращої сталої: у світі з незмінними властивостями змін найкраща швидкість навчання і є сталою. Виграш з'являється, коли світ раптом стає іншим. Тоді оцінка несподівано опиняється далеко від входу, а $P$ про це не знає. Правильна дія --- різко підняти $P$ і тим самим перейти в режим запису. За гіпотезою, яку ми обговорювали в розділі~\ref{sec:nora}, про таку \emph{неочікувану} зміну повідомляє \nt{NORA}~\cite{YuDayan2005}. Розподіл праці виходить природним: \nt{NORA} помічає, що світ змінився, \nt{ACET} тримає мережу в режимі запису, доки нову обстановку не вивчено.

\section{Сон як режим читання}

Якщо високий рівень ацетилхоліну --- режим запису, то що робить мережа за низького? У повільному сні рівень ацетилхоліну падає, внутрішні зв'язки звільняються від пригнічення, входу від органів чуттів майже немає. Мережа в режимі читання без підказки програє те, що в ній записано. Реєстрація активності показує, що нейрони, які вдень працювали разом, у повільному сні знову спрацьовують разом~\cite{WilsonMcNaughton1994}, і навіть у тому самому порядку~\cite{Skaggs1996}. За гіпотезою~\cite{Hasselmo1999}, ці повтори переписують швидко вивчене вдень у довготривалу пам'ять (штучне подовження коротких спалахів повторів покращує пам'ять~\cite{FernandezRuiz2019}): удень --- запис із входу, уночі --- перезапис зсередини. Для інженера це схоже на запис у швидкий буфер удень і його скидання в повільну постійну пам'ять уночі.

\section{Підсумок}

\begin{table}[t]
\centering
\footnotesize
\setlength{\tabcolsep}{4pt}
\renewcommand{\arraystretch}{1.15}
\begin{tabular}{>{\raggedright}p{3.2cm}>{\raggedright}p{4.2cm}>{\raggedright\arraybackslash}p{6.0cm}}
\toprule
Що робить \nt{ACET} & Як & Що відомо \\
\midrule
послаблює внутрішні зв'язки & множник $1-a$ у~\eqref{eq:acetnet} & виміряно на зрізах \\
посилює вхід & множник $\frac{1+a}{2}$ & підвищення збудливості і передачі вхідними шляхами виміряно \\
полегшує навчання & швидкість $\eta a$ & виміряно \\
перемикає «записувати/читати» & твердження~\ref{prop:writeread} & випливає з моделей; прямого вимірювання режимів немає \\
повідомляє очікувану невизначеність & $P$ у~\eqref{eq:kalman} & гіпотеза \\
звільняє мережу для повторів у сні & низьке $a$ в повільному сні & низький рівень уві сні і повтори виміряно; перезапис у довготривалу пам'ять --- гіпотеза \\
\bottomrule
\end{tabular}
\caption{Що додає \nt{ACET} до мережі з \nt{GLUT}, \nt{GABA}, \nt{DOPA} і \nt{NORA}.}
\label{tab:acet}
\end{table}

\nt{DOPA} каже мережі, куди змінювати ваги, \nt{NORA} --- наскільки рішуче користуватися тим, що вона знає, а \nt{ACET} --- чи слухати їй світ, чи себе (таблиця~\ref{tab:acet}). Це остання широкомовна лінія керування з таблиці~\ref{tab:types}: дозвіл запису, без якого пам'ять, що зберігає образи в тих самих зв'язках, якими вона їх читає, псувала б себе під час кожного читання.

\section{Задачі}

\begin{exercise}[\lvlA]
\label{ex:09:1}
\emph{Скільки пам'ятати.} Фільтр~\eqref{eq:kalman} з $q=0.01$, $R=1$ пам'ятає близько $10$~с. Знайдіть $P_\infty$ і пам'ять $\sqrt{R/q}$, якщо (а)~шум датчика зріс удвічі за амплітудою; (б)~світ став блукати в сто разів швидше. (в)~У скільки разів має зрости шум, щоб мережа пам'ятала хвилину?
\end{exercise}

\begin{exercise}[\lvlB]
\label{ex:09:2}
\emph{Режим запису після зміни.} Світ змінився, і невизначеність фільтра~\eqref{eq:kalman} з $q=0.01$, $R=1$ підскочила до $P_0\to\infty$. (а)~З якою сталою часу $P$ повертається до $P_\infty$? (б)~Через скільки секунд швидкість навчання перевищуватиме усталену не більше ніж на $10\,\%$?
\end{exercise}

\begin{exercise}[\lvlB]
\label{ex:09:3}
\emph{Стала швидкість навчання.} Мережа навчається як $\dd\hat s/\dd t=(y-\hat s)/T$; світ блукає, $\dd s/\dd t=w$, $y=s+n$, де $w$ і $n$ --- білі шуми з інтенсивностями $q$ і $R$. Знайдіть найкраще $T$ і дисперсію помилки за нього. У скільки разів вона зросте, якщо $T$ помилкове в $2$ і в $10$ разів?
\end{exercise}

\begin{exercise}[\lvlB]
\label{ex:09:4}
\emph{Де межа запису.} У \texttt{models/\allowbreak acet\_memory.py} поріг $\theta=0.35$, вага всередині образу $w=0.041$ (в ідеалі $0.045$). Новий образ ділить $m=10$ нейронів зі старим. За якого $a$ у~\eqref{eq:acetnet} решта нейронів старого образу не запалюється від цих $m$ (поріг різкий)?
\end{exercise}

\begin{exercise}[\lvlB]
\label{ex:09:5}
\emph{Моделювання: ємність пам'яті.} Змініть у \texttt{models/\allowbreak acet\_memory.py} функцію \texttt{read\_sweep}: при $a=1$ запишіть $M$ образів ($M$ від $5$ до $100$), потім при $a=0$ прочитайте перші десять за половиною. За якого $M$ з'являються помилки і як граничне $M_{\max}$ залежить від $N$ і $p$?
\end{exercise}

\begin{exercise}[\lvlB]
\label{ex:09:6}
\emph{Проєктування: запис без витоку.} Зі швидкістю навчання $\eta a$ мережа навчається і під час читання. Запропонуйте монотонну $\eta(a)\le\eta$, що дорівнює нулю при $a\le0.3$ і не гірша за $\eta a$ при $a\ge0.9$. Перевірте: після $20$ читань при $a=0.2$ з підказкою, де три чужі нейрони, скільки їх увійшло в образ?
\end{exercise}

\begin{exercise}[\lvlC]
\label{ex:09:7}
\emph{Як переписати буфер уночі.} (а)~Уві сні $a=0.05$, повтор триває $0.1\,\text{с}$ проти $0.2\,\text{с}$ удень при $a=1$. Скільки повторів накопичать зміну ваг одного денного пред'явлення за швидкості $\eta a$ і за $\eta(a)$ задачі~\ref{ex:09:6}? (б)~Сховище навчається в $100$ разів повільніше за буфер і чує образ $20$ разів за ніч по $0.1\,\text{с}$. Скільки потрібно ночей?
\end{exercise}

\chapter{Уся машина}
\chaptermark{Уся машина}
\label{sec:synthesis}

\section{Що ми зібрали}

Ми брали з таблиці~\ref{tab:types} по одному типу нейронів і питали, що він додає до обчислювальної машини. Щоразу правила були ті самі: лише те, що буває в живому організмі, і жодного спільного тактового сигналу. Тепер зберемо частини в одну машину й подивимося, як вони працюють разом.

Картина з розділу~\ref{sec:neurontypes} підтвердилася і стала точнішою. Величезна адресна мережа \nt{GLUT}-нейронів несе дані. \nt{GABA}-нейрони керують потоком цих даних. А над мережею висять чотири широкомовні канали, кожен зі своєю ручкою: \nt{DOPA} каже, які зміни ваг залишити, \nt{SERO} тримає загальний рівень і, за гіпотезою, горизонт планування, \nt{NORA} вибирає між пошуком і рішенням, \nt{ACET} --- між записом і читанням (рис.~\ref{fig:machine}).

\begin{figure}[t]
\centering
\begin{tikzpicture}[>=Stealth,
  blk/.style={draw,very thick,rounded corners=3pt,align=center,font=\small},
  reg/.style={draw,very thick,double,double distance=1.2pt,circle,minimum size=1.25cm,inner sep=0pt,font=\scriptsize,fill=orange!15},
  bc/.style={->,thick,densely dashed,orange!85!black}]
% sensors, bus, actions
\node[blk,fill=green!8,text width=1.7cm] (S) at (0.9,0) {датчики\\\scriptsize увімк./вимк.};
\draw[very thick,rounded corners=4pt,fill=blue!5] (2.6,-1.35) rectangle (9.0,1.35);
\node[font=\small] at (5.8,0.95) {шина \nt{GLUT}: дані};
\node[font=\scriptsize,align=center] at (4.3,0.25) {збіги, затримки,\\логіка (розд.~\ref{sec:glutcomputer})};
\node[font=\scriptsize,align=center] at (7.4,0.25) {пам'ять, код\\(розд.~\ref{sec:code}, \ref{sec:acet})};
\draw[thick,red!70!black,fill=red!8,rounded corners=2pt] (2.8,-1.15) rectangle (8.8,-0.55);
\node[font=\scriptsize,red!60!black] at (5.8,-0.85) {\nt{GABA} (розд.~\ref{sec:gaba}): заборона, вибір, ділення, ритм};
\node[blk,fill=blue!5,text width=2.0cm] (A) at (10.9,0) {вибір\\дії};
\draw[->,very thick] (S) -- (2.6,0);
\draw[->,very thick] (9.0,0) -- (A);
\draw[->,very thick] (A) -- (12.6,0) node[right,font=\small] {м'язи};
\pic[scale=1.1] at (13.2,-0.5) {spring};
\pic[scale=1.15] at (0.4,0.85) {eyepic};
\pic[scale=1.05] at (1.35,0.85) {speaker};
% registers
\node[reg] (D) at (3.4,3.2) {\nt{DOPA}};
\node[reg] (C) at (5.6,3.2) {\nt{SERO}};
\node[reg] (N) at (7.8,3.2) {\nt{NORA}};
\node[reg] (H) at (10.0,3.2) {\nt{ACET}};
\node[font=\scriptsize,align=center,above=1pt] at (D.north) {помилка $\delta$};
\node[font=\scriptsize,align=center,above=1pt] at (C.north) {рівень, $\tau_\gamma$};
\node[font=\scriptsize,align=center,above=1pt] at (N.north) {підсилення $g$};
\node[font=\scriptsize,align=center,above=1pt] at (H.north) {запис $a$};
\draw[bc] (D.south) -- (3.4,1.35);
\draw[bc] (C.south) -- (5.6,1.35);
\draw[bc] (N.south) -- (7.8,1.35);
\draw[bc] (H.south) -- (10.0,2.1) -- (8.8,1.35);
\draw[bc] (H.south east) to[bend left=20] (A.north east);
\draw[->,thick] (2.85,1.35) -- (2.85,2.75) -- (D.west) node[pos=0.25,left,font=\scriptsize] {$V$};
\draw[->,thick,green!45!black] (0.4,3.2) node[left,black,font=\small] {винагорода $r$} -- (2.2,3.2) -- (D.west);
\pic[scale=0.9] at (-0.45,3.7) {juice};
\draw[->,thick,densely dotted,gray!50!black] ([yshift=0.45cm]D.north) to[bend left=18] node[above,font=\scriptsize,black] {несподіванка} ([yshift=0.45cm]N.north);
\end{tikzpicture}
\caption{Машина цілком. Шина \nt{GLUT} несе дані від датчиків до вибору дії; \nt{GABA} керує потоком усередині неї. Подвійні кружки --- широкомовні канали; штрихові стрілки --- розсилання їхніх сигналів усій мережі, у кожного своя ручка. \nt{DOPA} отримує винагороду $r$ і очікування $V$ від критика всередині шини (розд.~\ref{sec:dofa}); точкова стрілка --- гіпотеза, що про несподіванку \nt{NORA} дізнається з незвично великих помилок (розд.~\ref{sec:nora}).}
\label{fig:machine}
\end{figure}

\section{Одна формула для всіх ручок}

Ручки можна зібрати в одну схематичну формулу. Це не нова модель, а зведення моделей розділів~\ref{sec:glutcomputer}--\ref{sec:acet}: кожен символ узято зі свого розділу.
\begin{empheq}[box=\keybox]{equation}
\begin{aligned}
\tau\,\frac{\dd r_i}{\dd t} \;&=\; -r_i + F\Bigl(g\,\Bigl[\tfrac{1+a}{2}\,x_i + (1-a)\sum_j W_{ij}\,r_j - \sum_k M_{ik}\,q_k - \theta\Bigr]\Bigr),\\
\frac{\dd W_{ij}}{\dd t} \;&=\; \eta\,a\,\delta(t)\,e_{ij}(t),
\qquad \tau_e\,\frac{\dd e_{ij}}{\dd t} = -e_{ij} + r_i\,r_j,
\qquad \delta = r + \frac{\dd V}{\dd t} - \frac{V}{\tau_\gamma} .
\end{aligned}
\label{eq:machine}
\end{empheq}
Тут $r_i$ --- частоти \nt{GLUT}-нейронів, $x_i$ --- вхід від датчиків, $W_{ij}\ge0$ --- їхні взаємні ваги (розділ~\ref{sec:glutcomputer}); $q_k$ --- частоти \nt{GABA}-нейронів, $M_{ik}\ge0$ --- сила їхнього гальмування (розділ~\ref{sec:gaba}); $e_{ij}$ --- слід, $\delta$ --- помилка \nt{DOPA} (розділ~\ref{sec:dofa}); $\tau_\gamma$ --- горизонт, який, за гіпотезою, задає \nt{SERO}; $g$ --- підсилення \nt{NORA} (розділ~\ref{sec:nora}); $a$ --- рівень \nt{ACET} (розділ~\ref{sec:acet}).

Погляньте, чого в~\eqref{eq:machine} немає. Немає такту: усе записано диференціальними рівняннями, і кожен нейрон змінюється сам. Немає окремої пам'яті: пам'ятають ті самі $W_{ij}$, якими йде обчислення, і та сама активність $r_i$. Немає персональних поправок для переважної більшості синапсів: їхнє навчання --- добуток місцевого сліду й одного спільного числа; окремий провід помилки до кожного виходу є лише в колі, що навчає рухів (розділ~\ref{sec:motorlearning}). І керувальних сигналів лише чотири види: $\delta$, $\tau_\gamma$, $g$, $a$ на вісімдесят мільярдів нейронів. Кожен із них насправді не одне число, а невеликий джгут паралельних ліній (твердження~\ref{prop:bundle}), але ліній у джгуті лічені одиниці.

\begin{keyblock}
\begin{proposition}[Архітектура]
\label{prop:architecture}
Машину мозку, наскільки ми її зараз розуміємо, можна описати трьома шарами. Дані йдуть адресною шиною \nt{GLUT}. Потік даних спрямовують, обмежують і нормують адресні \nt{GABA}-нейрони. Режим роботи задають кілька широкомовних каналів, кожен --- невеликий джгут паралельних ліній, спільних для великих ділянок мережі. Перші два шари встановлено твердо; розподіл ролей між широкомовними каналами --- здебільшого гіпотеза.
\end{proposition}
\end{keyblock}

\section{Усі типи в одній таблиці}

Таблиця~\ref{tab:synthesis} зводить усі типи нейронів книги. Чотири типи, яким не дістався власний розділ, посіли в ній по рядку; двоє з них знайшли роль у розділах про виходи машини: \nt{GLYC} --- у петлях спинного мозку (розділ~\ref{sec:muscles}), \nt{ADRE} --- у хімічному виході (розділ~\ref{sec:chemoutput}). Роль решти двох в обчисленні або проста, або невідома. Речовини, що не задають тип нейрона, зібрано в додатку~\ref{sec:othermediators}.

\begin{table}[t]
\centering
\footnotesize
\setlength{\tabcolsep}{3pt}
\renewcommand{\arraystretch}{1.15}
\begin{tabular}{>{\raggedright}p{1.3cm}>{\raggedright}p{1.0cm}>{\raggedright}p{3.9cm}>{\raggedright}p{1.5cm}>{\raggedright}p{1.8cm}>{\raggedright\arraybackslash}p{3.8cm}}
\toprule
Тип & Розділ & Що обчислює & Час & Шина & Що відомо \\
\midrule
\nt{GLUT} & \ref{sec:glutcomputer}, \ref{sec:code}, \ref{sec:sensors}, \ref{sec:motorlearning}, \ref{sec:dendrites} & дані: збіги, затримки, логіка, пам'ять, послідовності & мс & адресна & встановлено \\
\nt{GABA} & \ref{sec:gaba}, \ref{sec:code}, \ref{sec:pain}, \ref{sec:motorlearning} & керування потоком: заборона, вибір, ділення, стійкість, ритм; ворота болю & мс & адресна & встановлено; ділення частоти --- разом із фоновим шумом \\
\nt{SERO} & \ref{sec:serocomputer}, \ref{sec:pain}, \ref{sec:sleep} & регулятор рівня, згладжування; горизонт $\tau_\gamma$; гучність болю; режими сну & секунди & об'єм & самогальмування виміряно; горизонт --- гіпотеза \\
\nt{DOPA} & \ref{sec:dofa}, \ref{sec:dofaalt} & помилка передбачення $\delta$; повільний рівень --- ціна часу & $0.1$\,с і хвилини & широкомовна & $\delta$ виміряно; розширення --- розд.~\ref{sec:dofaalt} \\
\nt{NORA} & \ref{sec:nora}, \ref{sec:pain}, \ref{sec:chemoutput}, \ref{sec:sleep} & підсилення $g$: шукати чи вирішувати; несподіванка; «газ» для органів & секунди & широкомовна & зростання відношення сигнал/шум виміряно; роль у пошуку --- гіпотеза \\
\nt{ACET} & \ref{sec:acet}, \ref{sec:muscles}, \ref{sec:chemoutput}, \ref{sec:sleep} & запис чи читання; швидкість навчання; вихід на м'яз; «гальмо» для органів & секунди & обидві & три ефекти виміряно; перемикання режимів --- гіпотеза \\
\nt{GLYC} & \ref{sec:muscles}, \ref{sec:pain} & від'ємна вага в спинному мозку й стовбурі: заборона м'яза-антагоніста & мс & адресна & встановлено \\
\nt{HIST} & --- & глобальний сигнал «не спи» & повільно & широкомовна & роль в обчисленні не вивчено \\
\nt{ADRE} & \ref{sec:chemoutput} & «газ» для всього тіла через кров; у мозку невідомо & повільно & широкомовна & поза мозком встановлено; роль у мозку неясна \\
\nt{ASPA} & --- & слабка додатна вага & мс & адресна & роль спірна \\
\bottomrule
\end{tabular}
\caption{Усі типи нейронів книги: що обчислює кожен і наскільки це відомо.}
\label{tab:synthesis}
\end{table}

\section{Як ручки працюють разом}

Досі кожен розділ повертав одну ручку. Простежмо тепер одну подію через усю машину (рис.~\ref{fig:timeline}). Тварина давно вивчила, що дія $A$ приносить сік. Якоїсь миті світ змінюється: $A$ більше не приносить соку, а приносить його дія $B$, яку тварина майже не пробує.

\emph{Помилка.} Сік після $A$ не прийшов, і \nt{DOPA}-нейрони відповідають провалом, $\delta<0$ за формулою помилки~\eqref{eq:ctd}. Синапси, які щойно вибрали $A$, несуть слід і слабшають.

\emph{Несподіванка.} Один провал --- звичайна випадковість. Але провали повторюються, і помилки стають більшими за звичайні. За гіпотезою розділу~\ref{sec:nora}, це і є сигнал \nt{NORA}: світ змінився. Підсилення $g$ падає, вибір стає м'яким, і шум частіше приводить до спроби $B$.

\emph{Запис.} За гіпотезою розділу~\ref{sec:acet}, після зміни \nt{ACET} зростає й переводить мережу в режим запису: внутрішні зв'язки, що зберігають стару звичку, послаблено, вхід від датчиків підсилено, ваги змінюються легко.

\emph{Нова звичка.} Спроба $B$ приносить сік, якого ніхто не чекав: сплеск, $\delta>0$, і синапси, що вибрали $B$, посилюються. Кілька таких спроб --- і очікування $V$ наздоганяє новий світ, помилки знову малі.

\emph{Повернення.} Несподіванка минула, \nt{NORA} повертає високе підсилення, вибір знову жорсткий; \nt{ACET} спадає, мережа в режимі читання користується вивченим. Увесь цей час \nt{SERO}, за гіпотезою, задає горизонт $\tau_\gamma$: наскільки далекий сік ще варто чекати.

\begin{figure}[t]
\centering
\begin{tikzpicture}[>=Stealth,x=1.1cm,y=0.95cm]
\foreach \y/\lab in {4/{$\delta$ (\nt{DOPA})},3/{$g$ (\nt{NORA})},2/{$a$ (\nt{ACET})},1/{вибір $B$}}{
  \draw[gray!60!black] (0,\y-0.35) -- (10,\y-0.35);
  \node[left,font=\scriptsize] at (0,\y) {\lab};}
\draw[densely dashed,gray!60!black] (2,0.4) -- (2,4.55) node[above,font=\scriptsize,black] {світ змінився};
\draw[densely dashed,gray!60!black] (5,0.4) -- (5,4.55) node[above,font=\scriptsize,black] {перший сік за $B$};
\pic[scale=0.85] at (2,5.25) {nojuice};
\pic[scale=0.85] at (5,5.25) {juice};
% delta: dips after the change, burst at the first juice for B, then back to zero
\draw[thick,red!70!black] (0,3.65) -- (2.3,3.65) -- (2.3,3.3) -- (2.5,3.3) -- (2.5,3.65) -- (3.2,3.65) -- (3.2,3.3) -- (3.4,3.3) -- (3.4,3.65) -- (4.1,3.65) -- (4.1,3.35) -- (4.3,3.35) -- (4.3,3.65) -- (5.0,3.65) -- (5.0,4.35) -- (5.2,4.35) -- (5.2,3.65) -- (5.8,3.65) -- (5.8,4.1) -- (6.0,4.1) -- (6.0,3.65) -- (6.6,3.65) -- (6.6,3.85) -- (6.8,3.85) -- (6.8,3.65) -- (10,3.65);
% gain: high, drops after surprise, recovers
\draw[thick,blue!60!black] (0,3.2) -- (3.0,3.2) -- (3.4,2.75) -- (5.8,2.75) -- (7.2,3.2) -- (10,3.2);
% ACh: low, rises after the change, falls after learning
\draw[thick,green!45!black] (0,1.75) -- (2.8,1.75) -- (3.3,2.25) -- (6.8,2.25) -- (7.8,1.75) -- (10,1.75);
% choice of B
\draw[thick,black] (0,0.7) -- (3.0,0.7) plot[domain=3:10,samples=40] (\x,{0.7+0.6/(1+exp(-(\x-5.3)*1.8))});
\draw[->,gray!70!black] (0,0.2) -- (10.2,0.2) node[below left,font=\scriptsize,black] {час};
\end{tikzpicture}
\caption{Одна подія, простежена через усю машину. Це якісна схема за моделями й гіпотезами розділів~\ref{sec:dofa}--\ref{sec:acet}, а не результат розрахунку. Після зміни \nt{DOPA} відповідає провалами; повторні провали, за гіпотезою, піднімають сигнал несподіванки, \nt{NORA} знижує підсилення, \nt{ACET} вмикає запис; перший сік за $B$ дає сплеск, вибір переходить на $B$, і ручки повертаються.}
\label{fig:timeline}
\end{figure}

Зауважте, що жодна ручка не знає, що роблять інші. Кожна відгукується на свої ознаки --- помилку, її незвичну величину, невизначеність, --- і порядок дій складається сам, як в асинхронній схемі, де кожен блок спрацьовує, коли готові його входи. Таку злагоджену роботу загалом, наскільки нам відомо, ще не перевіряли: окремо ручки виміряно, а їхня спільна робота --- гіпотеза.

\section{Чим ця машина не схожа на комп'ютер}

\begin{table}[t]
\centering
\footnotesize
\setlength{\tabcolsep}{4pt}
\renewcommand{\arraystretch}{1.15}
\begin{tabular}{>{\raggedright}p{2.2cm}>{\raggedright}p{4.2cm}>{\raggedright\arraybackslash}p{6.6cm}}
\toprule
& Звичайний комп'ютер & Мозок \\
\midrule
час & спільний такт, близько гігагерца & спільного такту немає; кожен нейрон змінюється сам, вікна задають події й місцеві ритми (розд.~\ref{sec:glutcomputer}, \ref{sec:gaba}, \ref{sec:code}) \\
елементи & надійні двійкові вентилі & аналогові порогові елементи; синапс губить більшість імпульсів (розд.~\ref{sec:code}) \\
знак зв'язку & будь-який & задається нейроном-відправником (розд.~\ref{sec:neurontypes}) \\
пам'ять & окремо від процесора & у тих самих зв'язках, що й обчислення, і в самопідтримній активності (розд.~\ref{sec:glutcomputer}, \ref{sec:acet}) \\
навчання & зворотне поширення помилки & місцеві правила, одне широкомовне число $\delta$ (розд.~\ref{sec:dofa}) і проводи помилки до виходів кола рухів (розд.~\ref{sec:motorlearning}) \\
керування & регістри й шина команд & чотири широкомовні канали, кожен --- джгут із кількох ліній (розд.~\ref{sec:serocomputer}, \ref{sec:dofa}--\ref{sec:acet}) \\
енергія & десятки--сотні ватів на один процесор & близько $20$ ватів на весь мозок, у середньому близько імпульсу за секунду на нейрон (розд.~\ref{sec:code}) \\
\bottomrule
\end{tabular}
\caption{Мозок і звичайний комп'ютер.}
\label{tab:computer}
\end{table}

Таблиця~\ref{tab:computer} зводить головні відмінності. Кожна з них --- не вада, яку природа не встигла виправити, а частина одного рішення. Без спільного такту потрібні датчики «увімкнення---вимкнення» і детектори збігів. Ненадійним елементам потрібна надлишковість багатьох нейронів. Пам'яті, суміщеній з обчисленням, потрібен \nt{ACET}, щоб запис не псував читання. Навчанню без зворотних проводів потрібні \nt{DOPA} і шум. А енергетична межа в один імпульс за секунду змушує всю мову бути ощадливою й витрачати імпульси на новини.

\section{Чого ми не знаємо}

Найчесніше завершити це зведення переліком того, чого ми не знаємо. Кожен пункт --- задача для інженера.

\emph{Код.} Ми знаємо ємність імпульсного каналу і знаємо, що одержувач вибирає, яку частину повідомлення читати (розділ~\ref{sec:code}). Але наскільки кора користується точним часом імпульсів і чи є в нейронів «слова», невідомо.

\emph{Навчання через багато шарів.} Широкомовний сигнал навчає повільно, і що більше нейронів впливає на результат, то повільніше (твердження~\ref{prop:broadcastcost}). Як тоді кора налаштовує мільярди ваг у багатошарових колах, невідомо; є моделі, в яких помилку між шарами розносять пачки імпульсів~\cite{Payeur2021}. Можливо, частина відповіді --- в обчисленнях усередині гілок самого нейрона, де один нейрон поводиться як маленька багатошарова мережа: щоб повторити відповідь моделі одного нейрона кори, штучній мережі потрібно п'ять--вісім шарів~\cite{Beniaguev2021}, а гілки нейронів людини вміють обчислювати навіть виключне АБО~\cite{Gidon2020}. Інший кандидат --- самонавчання: якщо кора вчиться передбачати власні входи, помилка виникає на місці, в кожному шарі, і спільний сигнал для цього не потрібен (розділ~\ref{sec:dofa})~\cite{Keller2018}.

\emph{Що насправді передають широкомовні канали.} Для \nt{DOPA} є стандартна картина і шість її розширень (розділ~\ref{sec:dofaalt}); для \nt{NORA} і \nt{ACET} --- правдоподібні гіпотези (розділи~\ref{sec:nora}, \ref{sec:acet}); для \nt{SERO} --- здебільшого здогадки.

\emph{Узгодження в часі.} Ми бачили, як обчислити функцію, відміряти час від події й нарізати потік на вікна місцевим ритмом. Але як величезна мережа без спільного такту узгоджує роботу далеких ділянок, зрозуміло лише почасти.

\emph{Енергія.} Двадцять ватів на все. Ми розуміємо, чому код має бути розрідженим (твердження~\ref{prop:economy}), але збудувати машину, яка робила б те саме за такої самої потужності, поки що не вміє ніхто~\cite{MehonicKenyon2022}.

Мозок --- обчислювальна машина, зібрана не інженером, але за законами, які впізнає будь-який інженер: RC-кола, пороги, зворотні зв'язки, фільтри, шини й широкомовні регістри. Її схему ми прочитали лише частково. Дочитуватимуть її ті, хто вміє читати схеми.

\section{Що далі в книзі}

Цей розділ зібрав обчислювальне ядро машини. Решта розділів виходить за його межі. Розділи~\ref{sec:sensors} і~\ref{sec:recognition} --- про входи: як датчики перетворюють світло, звук і молекули на імпульси і як машина впізнає в них предмети. Розділи~\ref{sec:muscles}--\ref{sec:chemoutput} --- про виходи й те, що їх обслуговує: м'язи, біль як сигнал тривоги, навчання рухів окремими проводами помилки і повільний хімічний вихід у тіло. Розділ~\ref{sec:debugging} --- про те, як машину налагоджують ззовні, зокрема інтерфейсами мозок---комп'ютер. Розділи~\ref{sec:eetypes} і~\ref{sec:dendrites} ще раз сортують нейрони, тепер за електричною функцією і за кількістю входів. Розділ~\ref{sec:sleep} --- про те, що машина робить, коли відключена від світу, а розділи~\ref{sec:livingmachine} і~\ref{sec:dishbrain} --- про машини, зібрані з живих нейронів на чипі.

\section{Задачі}

Задачі цього розділу поєднують результати різних розділів: кожна спирається щонайменше на два з них.

\begin{exercise}[\lvlA]
\label{ex:10:1}
\emph{Шина і регістри.} Чотири широкомовні канали --- джгути приблизно по п'ять ліній (твердження~\ref{prop:bundle}); кожна лінія змінюється раз на секунду й розрізняє чотири рівні. Скільки років знадобилося б керуванню, щоб передати те, що шина \nt{GLUT} ($8.6\cdot10^{10}$ нейронів із частотою~\eqref{eq:energy}, точність $1\ms$ за~\eqref{eq:spikeentropy}) передає за секунду?
\end{exercise}

\begin{exercise}[\lvlB]
\label{ex:10:2}
\emph{Дві ручки на одній петлі.} У петлі~\eqref{eq:ei} ручки~\eqref{eq:machine} стоять на збудженні: $\tau_E\dot E=-E+g[(1-a)w_{EE}E-w_{EI}I+h]$, \nt{GABA} ними не керується. За чисел рис.~\ref{fig:ei} пачка \nt{NORA} піднімає $g$ до $6$. За якого найменшого $a$ петля стійка? Чи можна крутити ручки \nt{NORA} і \nt{ACET} незалежно?
\end{exercise}

\begin{exercise}[\lvlB]
\label{ex:10:3}
\emph{Звідки береться ціна широкомовлення.} Виходи $N$ нейронів $y_k=f(u_k)+\xi_k$ з незалежним гауссовим шумом дисперсії $\sigma^2$; $R\approx R_0+\sum_kR'_k\xi_k$ з однаковими $R'_k$, \nt{DOPA} розсилає $\delta=R-R_0$. Знайдіть відношення сигнал/шум поправки $\delta\,\xi_jx_j$ за одну спробу. Як із $N$ зростає кількість спроб?
\end{exercise}

\begin{exercise}[\lvlB]
\label{ex:10:4}
\emph{Зв'язати звук і спалах.} Звук приходить на шину через $5\ms$, спалах --- через $30\ms$ плюс затримка першого імпульсу~\eqref{eq:latency} ($\tau=10\ms$, $U$ від $1.2\theta$ до $4\theta$); фон на кожному вході --- $5$ імпульсів за секунду. Виберіть затримку звукової лінії й найменше вікно збігів для всіх контрастів. Скільки хибних зв'язувань за секунду дасть фон?
\end{exercise}

\begin{exercise}[\lvlB]
\label{ex:10:5}
\emph{Моделювання: хто помічає зміну.} Критик бачить $y=s+\text{шум}$ ($R=1$); з імовірністю $1/200$ $s$ стрибком набуває нового значення з дисперсією $J^2=4$. Змоделюйте фільтр Калмана з $q=0$, що піднімає $P$ до $J^2$, коли $E\leftarrow0.9E+\delta^2/(P+R)$ перевищує $20$. На скільки його помилка менша, ніж за найкращого сталого $\alpha$?
\end{exercise}

\begin{exercise}[\lvlA]
\label{ex:10:6}
\emph{Скільки спроб на зміну звички.} Мережа вивчила $Q_A=0.8$, $Q_B=0.2$ і вибирає за~\eqref{eq:softmax}, $\sigma=1$, $g=8$. Тепер $A$ дає сік з імовірністю $0.2$; $Q_A\leftarrow Q_A+\alpha(r-Q_A)$, $Q_B$ не змінюється. Через скільки виборів $A$ імовірність вибрати $A$ впаде до $0.6$ при $\alpha=0.1$ і $0.5$? А якщо \nt{NORA} знизить $g$ до $2$?
\end{exercise}

\begin{exercise}[\lvlC]
\label{ex:10:7}
\emph{Джгут замість проводу.} Виходи $y_k=w_k+\xi_k$, винагорода $R=-\sum_ky_k^2$; лінія джгута розсилає групі з $G$ нейронів помилку її виходів, $w_k\leftarrow w_k+\eta\,\delta_l\xi_k$. Покажіть, що за найкращого $\eta$ до $\sum w_k^2=\varepsilon\sum w_k^2(0)$ потрібно $\approx(G+2)\ln(1/\varepsilon)$ спроб. Скільки за $\varepsilon=0.01$ вчиться коло з $10^6$ нейронів на п'яти лініях, якщо спроба --- раз на три секунди?
\end{exercise}

\chapter{Входи машини: як до неї під'єднано очі, вуха та інші датчики}
\chaptermark{Входи машини}
\label{sec:sensors}
\begin{chaptergoals}
\item чому кожне чуття --- перетворювач: рецептор-затвор, регулювання підсилення, кодер імпульсів;
\item чому дробовий шум фотонів обмежує зір у темряві і чому в сутінках ми бачимо повільніше;
\item як регулювання підсилення вміщує величезний діапазон в імпульси, тож датчики передають зміни;
\item як око стискає зображення в сто разів, а вухо працює як банк фільтрів.
\end{chaptergoals}

\section{Датчик як перетворювач}

На рис.~\ref{fig:machine} дані входили в машину зліва, з блока «датчики». Тепер відкриймо цей блок. Для інженера кожен орган чуття --- це \emph{перетворювач} з аналоговим вхідним каскадом і аналого-цифровим перетворювачем наприкінці, тільки цифри на виході --- не числа, а імпульси.

Схема в усіх органів чуття та сама (рис.~\ref{fig:transducer}). Фізична величина --- світло, тиск, коливання, молекула --- діє на білок-рецептор у мембрані чутливої клітини. Цей білок працює як затвор: він сам або через короткий ланцюжок реакцій відкриває іонний канал. Виникає \emph{рецепторний струм}, а з ним --- аналогова напруга на мембрані. Далі напруга проходить регулювання підсилення й потрапляє в кодер імпульсів --- уже знайомий інтегрувальний нейрон~\eqref{eq:glutneuron}, який перетворює рівень на частоту й моменти імпульсів. Вихід майже завжди той самий: чутливі нейрони ока, вуха, нюху й шкіри виділяють глутамат, тож входи під'єднуються просто до шини \nt{GLUT}.

\begin{figure}[t]
\centering
\begin{tikzpicture}[>=Stealth,
  blk/.style={draw,very thick,rounded corners=2pt,align=center,font=\footnotesize,minimum height=1.0cm,text width=1.9cm,fill=blue!5}]
\node[align=center,font=\small] (P) at (0.1,0) {світло, звук,\\тиск,\\молекула};
\node[blk] (R) at (3.0,0) {рецептор-\\затвор};
\node[blk] (G) at (6.5,0) {регулювання\\підсилення};
\node[blk] (E) at (10.0,0) {кодер\\імпульсів};
\node[align=center,font=\small] (B) at (13.4,0) {шина\\\nt{GLUT}};
\draw[->,very thick] (P) -- (R);
\foreach \ic/\xx in {sun/-1.1,speaker/-0.3,press/0.5,molecule/1.3}{\pic[scale=0.85] at (\xx,1.05) {\ic};}
\draw[->,very thick] (R) -- node[above,font=\scriptsize] {струм} (G);
\draw[->,very thick] (G) -- node[above,font=\scriptsize] {рівень} (E);
\draw[->,very thick,red!70!black] (E) -- node[above,font=\scriptsize,black] {імпульси} (B);
\draw[decorate,decoration={brace,amplitude=5pt,mirror},gray!60!black] (1.8,-0.8) -- (7.7,-0.8) node[midway,below=5pt,font=\scriptsize,black] {аналогова частина};
\draw[decorate,decoration={brace,amplitude=5pt,mirror},gray!60!black] (8.8,-0.8) -- (11.2,-0.8) node[midway,below=5pt,font=\scriptsize,black] {імпульси};
\end{tikzpicture}
\caption{Будь-який орган чуття як ланцюг перетворень. Білок-рецептор відкриває канал і дає струм; регулювання підсилення пристосовує його до обставин; кодер~\eqref{eq:glutneuron} перетворює рівень на імпульси, які йдуть на шину \nt{GLUT}. В оці й вусі перші каскади взагалі не видають імпульсів: сигнал лишається аналоговим, доки не доводиться передавати його далеко.}
\label{fig:transducer}
\end{figure}

Одна деталь добре знайома електронникові. В оці й у вусі найперші клітини імпульсів не видають зовсім: вони передають сусідам плавну напругу, як аналоговий каскад на платі. Імпульси з'являються лише там, де сигнал треба везти далеко. Аналоговий сигнал на відстані міліметрів загасає й зашумлюється, а імпульс, як ми бачили в розділі~\ref{sec:neurontypes}, доходить до кінця проводу тієї самої висоти. Оцифровують там, де починається довга лінія.

\section{На межі фізики}

Датчики мозку працюють біля фізичних меж чутливості. Датчик світла в темряві відповідає на \emph{один фотон}: поглинання однієї частинки світла дає струм близько пікоампера, помітний на тлі власного шуму~\cite{Baylor1979}. Датчик звуку помічає зміщення свого чутливого пучка менш ніж на нанометр~\cite{Hudspeth1989}.

Коли світла мало, точність обмежує не датчик, а саме світло: фотони надходять випадково, пуассонівським потоком. Якщо за час накопичення $T$ у датчик потрапляє в середньому $N=\Phi T$ фотонів, їхня кількість тремтить на $\sqrt N$. Щоб приріст $\Delta N$ був помітний, він має в кілька разів перевищувати це тремтіння, і розрізненна частка зміни яскравості дорівнює
\begin{empheq}[box=\keybox]{equation}
\frac{\Delta I}{I} \;\approx\; \frac{k}{\sqrt{\Phi\,T}} .
\label{eq:photonnoise}
\end{empheq}
Це та сама формула, що й для кількості імпульсів~\eqref{eq:rateerror}: дробовий шум фотонів і дробовий шум імпульсів підкоряються одному закону. У темряві око може поліпшити точність лише одним способом --- довше накопичувати фотони, і час накопичення справді зростає до десятих часток секунди. Тому в сутінках ми бачимо повільніше.

\section{Динамічний діапазон: регулювання підсилення}

Найважча інженерна задача для датчиків --- діапазон. Освітленість від зоряної ночі до сонячного снігу змінюється приблизно на десять порядків, інтенсивність звуку від порога чутності до больового --- на дванадцять. А частота імпульсів змінюється від нуля до кількох сотень за секунду, менш ніж на три порядки. Лінійно такий вхід у такий вихід не вкласти.

Рішення те саме, що в будь-якого приймача: \emph{автоматичне регулювання підсилення}. Відповіді чутливих клітин ока добре описує формула
\begin{empheq}[box=\keybox]{equation}
r \;=\; r_{\max}\,\frac{I}{I+\sigma} ,
\label{eq:nakarushton}
\end{empheq}
де $\sigma$ --- яскравість, за якої відповідь сягає половини~\cite{NakaRushton1966}. Сама собою~\eqref{eq:nakarushton} охоплює лише два-три порядки. Але $\sigma$ не стала: за тривалої роботи на яскравому тлі вона зростає слідом за середньою яскравістю. Крива відповіді зсувається вздовж логарифмічної осі яскравості й щоразу ставить свою круту ділянку туди, де зараз перебуває вхід (рис.~\ref{fig:adaptation}). Це те саме ділення на загальну активність, що й шунтувальне гальмування розділу~\ref{sec:gaba}~\cite{Carandini2012}, тільки знаменник тут --- середня яскравість за останні секунди.

\begin{figure}[t]
\centering
\begin{tikzpicture}[>=Stealth,x=1.25cm,y=2.8cm]
\draw[->,gray!70!black] (-2,0) -- (6.4,0) node[below left,font=\small,black] {$\log_{10} I$};
\draw[->,gray!70!black] (-2,0) -- (-2,1.15) node[left,font=\small,black] {$r/r_{\max}$};
\foreach \u in {-2,0,2,4,6}{\draw[gray!60!black] (\u,-0.02) -- (\u,0.02); \node[below,font=\scriptsize] at (\u,0) {$\u$};}
\draw[densely dashed,gray!60!black] (-2,0.5) -- (6.2,0.5);
\foreach \s/\c/\lab in {0/blue!60!black/{темне тло},2/green!45!black/{середнє},4/red!70!black/{яскраве}}{
  \pgfmathsetmacro{\lo}{max(-2,\s-3)}
  \draw[very thick,\c] (-2,0) -- (\lo,0) plot[domain=\lo:6.2,samples=80] (\x,{1/(1+10^(\s-\x))});
  \node[font=\scriptsize,\c,anchor=west] at (\s+0.15,0.42) {\lab};}
\foreach \s/\ic in {0/moon,2/cloud,4/sun}{\pic at (\s+0.6,0.64) {\ic};}
\end{tikzpicture}
\caption{Регулювання підсилення датчика світла за~\eqref{eq:nakarushton}. Кожна крива охоплює два-три порядки яскравості; при переході на яскравіше тло $\sigma$ зростає, і крива зсувається логарифмічною віссю. Разом криві, що зсуваються, покривають увесь діапазон.}
\label{fig:adaptation}
\end{figure}

Регулювання має ціну, знайому з розділу~\ref{sec:code}: датчик повідомляє не яскравість, а яскравість \emph{відносно тла}. Сталий вхід поступово перестає викликати відповідь, а кожна зміна викликає. Входи машини від самого початку говорять про зміни, а не про рівні, і в цьому та сама економія імпульсів, що й у твердженні~\ref{prop:economy}~\cite{Barlow1961}. Звідси й датчики увімкнення та вимкнення розділу~\ref{sec:glutcomputer}~\cite{Schiller1992}: одні повідомляють, що стало яскравіше, інші --- що темніше.

Інженери повторили цей прийом у \emph{подієвих камерах}. Така камера не знімає кадрів за тактовим сигналом: кожен її піксель сам, без спільного такту, видає подію «яскравіше» або «темніше», коли логарифм яскравості змінився на задану частку~\cite{Lichtsteiner2008}. Це точнісінько двопровідний код увімкнення---вимкнення з розділу~\ref{sec:glutcomputer}, втілений у кремнії, і він дає камері діапазон і швидкодію, недосяжні для звичайної матриці~\cite{Gallego2022}.

\section{Око: стиснення до проводу}

В оці близько $10^8$ датчиків світла~\cite{Curcio1990}, а вихідних нейронів, які везуть зображення в мозок, --- близько мільйона. Перш ніж сигнал залишить око, його стискають приблизно в сто разів. Головні прийоми стиснення нам знайомі. Кожен вихідний нейрон відповідає на різницю між яскравістю в маленькій плямі та яскравістю навколо неї --- це віднімання сусідів через гальмування, як у розділі~\ref{sec:gaba}. Рівні ділянки зображення майже нічого не коштують, а краї й зміни передаються. Різні вихідні нейрони спеціалізовані: одні повідомляють про посвітлішання, інші про потемніння, ще інші про рух у певний бік (рис.~\ref{fig:direction}); у миші таких вихідних каналів налічили понад тридцять~\cite{Baden2016}.

Яка пропускна здатність отриманого кабелю? У морської свинки за записами вихідних нейронів виміряли близько $875$ тисяч біт за секунду на $\sim10^5$ таких нейронів; перерахунок на мільйон вихідних нейронів людини дає порядку $10^7$ біт за секунду~\cite{Koch2006} --- стільки ж, скільки старий мережевий кабель на десять мегабіт. За середньої частоти в кілька імпульсів за секунду на нейрон це кілька бітів на імпульс, як і виходило в розділі~\ref{sec:code}.

\section{Вухо: банк фільтрів}

Вухо розв'язує іншу задачу: звук треба розкласти за частотами. Інженер поставив би банк смугових фільтрів, і вухо робить саме це механічно. Звукове коливання біжить пружною перегородкою, властивості якої плавно змінюються вздовж її довжини, і кожна ділянка резонує на свою частоту. Датчики, що сидять на ділянці, відповідають лише на свою смугу (рис.~\ref{fig:filterbank}). Частоти розкладено за місцем приблизно логарифмічно: кожній октаві дістається схожа частка датчиків. Номер датчика, що спрацював, --- це частота, код місцем із розділу~\ref{sec:code}.

\begin{figure}[t]
\centering
\begin{tikzpicture}[>=Stealth,x=1.35cm,y=2.2cm]
\draw[->,gray!70!black] (-0.3,0) -- (7.6,0) node[above left,font=\small,black] {частота, Гц};
\draw[->,gray!70!black] (-0.3,0) -- (-0.3,1.15) node[left,font=\small,black] {відповідь};
\foreach \k/\f in {0/125,1/250,2/500,3/1000,4/2000,5/4000,6/8000}{
  \draw[gray!60!black] (\k,-0.02) -- (\k,0.02); \node[below,font=\scriptsize] at (\k,0) {\f};
  \draw[thick,blue!60!black] plot[domain={max(\k-1.2,-0.3)}:{\k+0.7},samples=60] (\x,{1/(1+((\x-\k)/(\x<\k?0.45:0.2))^4)});}
% a piano keyboard: seven white keys per octave, like the octaves on the axis
\foreach \i in {0,...,48}{\draw[gray!50!black,fill=white,line width=0.3pt] ({\i/7},-0.44) rectangle ({(\i+1)/7},-0.26);}
\foreach \o in {0,...,6}{\foreach \j in {0,1,3,4,5}{\fill[black] ({\o+(\j+1)/7-0.035},-0.35) rectangle ({\o+(\j+1)/7+0.035},-0.26);}}
\end{tikzpicture}
\caption{Вухо як банк смугових фільтрів, схема. Кожна крива --- відповідь датчиків однієї ділянки на звук різної частоти; центральні частоти йдуть через октаву, як клавіші на логарифмічній шкалі. У справжніх фільтрів крутий високочастотний схил і пологий низькочастотний. Унизу --- клавіатура: на ній, як і у вусі, кожній октаві дістається однакове місце.}
\label{fig:filterbank}
\end{figure}

Резонатори вуха не пасивні. Частина датчиків сама розгойдує перегородку в такт звуку й підсилює слабкі коливання в тисячі разів за амплітудою (на $66$--$76$\,дБ), а зі зростанням гучності підсилення падає~\cite{Ruggero1997,Robles2001}. Для інженера це підсилювач із додатним зворотним зв'язком, поставлений біля самої межі самозбудження, --- те саме життя на межі, що й у мережі розділу~\ref{sec:gaba}: біля межі система особливо чутлива до малих сигналів. Стиснення гучності тут робить той самий підсилювач: тихе підсилюється сильно, гучне --- слабко.

У вуха є й другий код. На низьких частотах імпульси нейронів ідуть у фазі зі звуком: нейрон спрацьовує на певній ділянці кожної хвилі, хоч і не на кожній хвилі. У морської свинки прив'язка до фази починає слабшати близько $600$\,Гц і ще помітна приблизно до $3.5$\,кГц~\cite{PalmerRussell1986}. Так у моментах імпульсів зберігається точний час звуку, а з нього --- різниця приходу звуку у два вуха, за якою мережа розділу~\ref{sec:glutcomputer} знаходить напрямок~\cite{Grothe2010}. На високих частотах імпульси за хвилею не встигають, і лишається тільки код місцем. Одне вухо користується обома кодами розділу~\ref{sec:code}: часом там, де нейрони встигають, і місцем там, де ні.

\section{Інші датчики}

\begin{table}[t]
\centering
\footnotesize
\setlength{\tabcolsep}{3pt}
\renewcommand{\arraystretch}{1.15}
\begin{tabular}{>{\raggedright}p{1.9cm}>{\raggedright}p{2.6cm}>{\raggedright}p{4.0cm}>{\raggedright\arraybackslash}p{4.6cm}}
\toprule
Чуття & Що вимірює & Як улаштовано вхід & Прийом із книги \\
\midrule
зір & світло & $\sim10^8$ датчиків, стиснення в $\sim100$ разів, $\sim10^7$ біт/с & регулювання підсилення, віднімання сусідів, два проводи, напрямок руху \\
слух & тиск повітря & механічний банк фільтрів, активний підсилювач & код місцем і код часом, затримки \\
дотик & тиск, вібрація & датчики, що звикають швидко й повільно & пара «фільтр верхніх частот --- фільтр нижніх» \\
температура & тепло & окремі датчики тепла й холоду & двопровідний код \\
нюх & молекули & $\sim400$ типів рецепторів, у кожного датчика один тип & комбінаторний код: запах --- набір типів, що спрацювали \\
смак & речовини в їжі & п'ять якостей: солодке, гірке, солоне, кисле, умамі & мічені лінії; частина клітин виділяє АТФ або серотонін, а не глутамат \\
положення тіла & розтяг м'язів & датчики довжини й швидкості розтягу & зворотний зв'язок для регулятора руху \\
\bottomrule
\end{tabular}
\caption{Як під'єднано датчики. Усі пристрої в третьому стовпці виміряно; правий стовпець пов'язує їх із прийомами з попередніх розділів.}
\label{tab:sensors}
\end{table}

Інші чуття зібрано в таблиці~\ref{tab:sensors}, і майже в кожному рядку впізнається прийом із попередніх розділів. Дотик користується парою фільтрів: одні датчики тиску відповідають на дотик і відразу замовкають, інші тримають відповідь, доки тиск є. Температуру передають два проводи, окремі для тепла й для холоду. Найнезвичайніший вхід --- нюх. Типів нюхових рецепторів у людини близько чотирьохсот, і кожен нюховий нейрон несе лише один тип~\cite{Mombaerts2004}. Запах активує не один тип, а набір, і повідомлення --- це \emph{які} типи спрацювали. Для інженера це двійковий вектор довжиною близько чотирьохсот, кількість можливих значень якого астрономічна.

\begin{keyblock}
\begin{proposition}[Входи машини]
\label{prop:sensors}
Кожне чуття --- перетворювач, який працює біля фізичної межі чутливості, стискає величезний динамічний діапазон автоматичним регулюванням підсилення й передає шиною \nt{GLUT} насамперед \emph{зміни}. Для цього датчики користуються тими самими прийомами, що й сама машина: двопровідним кодом, кодом місцем, кодом часом і діленням на тло. Будову датчиків виміряно; те, що їхні коди підлаштовано під найбільшу інформацію на імпульс, --- добре обґрунтована гіпотеза.
\end{proposition}
\end{keyblock}

Входи, як і все інше, працюють асинхронно. Кожен датчик видає імпульс, коли йому є що повідомити, а різні чуття надходять із різними затримками: звук --- через мілісекунди, світло --- через десятки мілісекунд. Як машина зводить їх в одну картину світу, якщо спільного такту немає, --- одна із задач узгодження в часі з розділу~\ref{sec:synthesis}.

\section{Задачі}

\begin{exercise}[\lvlA]
\label{ex:11:1}
\emph{Скільки накопичувати фотони.} У сутінках у датчик потрапляє $100$ фотонів за секунду; приріст помітний, якщо він у $3$ рази більший за тремтіння~\eqref{eq:photonnoise}. (а)~Скільки часу потрібно одному датчикові, щоб помітити зміну яскравості на $10\,\%$? (б)~Скільки датчиків треба підсумувати, щоб устигнути за $0.2\,\text{с}$, і в скільки разів упаде роздільність за віссю?
\end{exercise}

\begin{exercise}[\lvlA]
\label{ex:11:2}
\emph{Скільки бітів в імпульсі.} (а)~Око передає $\sim10^7$ біт/с через $10^6$ вихідних нейронів із частотою $3$--$5$ імпульсів за секунду. Яку частку пропускної здатності~\eqref{eq:spikeentropy} при $\Delta t=1\ms$ воно використовує? (б)~Запах вмикає $20$ з $\approx400$ типів рецепторів. Скільки бітів несе такий набір і скільки спільних типів у середньому мають два випадкові запахи?
\end{exercise}

\begin{exercise}[\lvlB]
\label{ex:11:3}
\emph{Що вміє одна крива~\eqref{eq:nakarushton}.} (а)~У якому діапазоні яскравостей відповідь зростає від $10\,\%$ до $90\,\%$ від $r_{\max}$? Скільки це порядків? (б)~Скільки положень зсувної кривої потрібно, щоб покрити десять порядків яскравості? дванадцять порядків гучності?
\end{exercise}

\begin{exercise}[\lvlB]
\label{ex:11:4}
\emph{Межа коду часом у вусі.} Різниця приходу звуку у два вуха --- до $0.22\,\text{м}/343\,\text{м/с}$. (а)~Імпульси йдуть у фазі з тоном: до якої частоти напрямок за ними визначається однозначно? (б)~Імпульс тремтить на $0.5\ms$, а розрізняти треба $20$~мкс. Скільки нейронів по $100$ імпульсів за секунду треба усереднити за $50\ms$?
\end{exercise}

\begin{exercise}[\lvlB]
\label{ex:11:5}
\emph{Подієва камера на кабель ока.} Камера з $10^6$ пікселів із виходом $10^7$ біт/с надсилає подію (адресу й знак), коли $\ln I$ у пікселі змінився на $\ln1.2$. З якою найбільшою швидкістю може бігти край довжиною $500$ пікселів із перепадом яскравості $4$? Скільки кадрів за секунду дала б через той самий вихід звичайна камера по $8$ бітів на піксель?
\end{exercise}

\begin{exercise}[\lvlB]
\label{ex:11:6}
\emph{Моделювання: регулювання підсилення.} Датчик~\eqref{eq:nakarushton} підлаштовує $\sigma$ з $\tau=1\,\text{с}$ у логарифмах, $\tau\,\dd\ln\sigma/\dd t=\ln I-\ln\sigma$, або лінійно, $\tau\,\dd\sigma/\dd t=I-\sigma$; тло стрибає $1\to100\to1$. Через скільки секунд відповідь на пробу $+10\,\%$ відновлюється до половини звиклої після стрибка вгору і вниз за кожного закону?
\end{exercise}

\begin{exercise}[\lvlC]
\label{ex:11:7}
\emph{Під який світ підлаштовано криву.} За малого сталого шуму виходу інформація про яскравість найбільша, коли $r(I)=r_{\max}\Phi_I(I)$, де $\Phi_I$ --- функція розподілу яскравостей. (а)~Для якого розподілу оптимальна крива~\eqref{eq:nakarushton} і який його розкид $\log_{10}I$? (б)~Яка крива оптимальна за пуассонівського шуму виходу?
\end{exercise}

\chapter{Розпізнавання зображень і відео}
\chaptermark{Розпізнавання}
\label{sec:recognition}

\section{Задача: те саме за різних пікселів}

Розділ~\ref{sec:sensors} довів зображення до входу машини: близько мільйона вихідних нейронів ока, стиснений потік, у якому передаються здебільшого краї й зміни. Але між датчиком і дією є крок, про який ми досі мовчали. Кішка на картинці --- це не набір пікселів. Зсуньте її на пів кадру, поверніть, віддаліть, вимкніть половину світла --- майже всі пікселі стануть іншими, а відповідь «кішка» має лишитися тією самою. Інженер називає це \emph{інваріантністю}: відповідь класифікатора не повинна змінюватися від зсуву, масштабу, повороту й освітлення.

Труднощі добре видно на простому досліді. Візьмімо одновимірну «сітківку» з $24$ точок і дві фігури по п'ять точок, які можуть стояти будь-де. Класифікатор, що порівнює вхід зі зразком, запам'ятаним в одному місці, вгадує в $49$--$52\%$ випадків --- стільки ж, скільки монетка. Зсув цілком ламає порівняння за пікселями. Потрібна інша схема.

\section{Конвеєр зі ступенів}

Як швидко мозок це робить, ми вже знаємо з розділу~\ref{sec:code}: тварину на миттєво показаній картинці впізнають приблизно за $150\ms$, і сигнал проходить близько десятка ступенів, по $10$--$20\ms$ на кожен~\cite{Thorpe1996}. На кожному ступені нейрон встигає видати нуль або один імпульс. Отже, швидке розпізнавання --- це один прохід конвеєром, без довгих роздумів і повторів. Питання в тому, що робить кожен ступінь.

Відповідь, яку підказали вимірювання на перших ступенях зору~\cite{HubelWiesel1962}, складається з двох почергових операцій (рис.~\ref{fig:conveyor}).

\emph{Вибірковість.} Нейрон ступеня $S$ дивиться на невелику ділянку входу й спрацьовує, лише якщо там є певне поєднання частин: ось цей край, а поруч той. Це І за частинами --- пороговий нейрон розділу~\ref{sec:glutcomputer}, поріг якого вищий, ніж дає будь-яка одна частина.

\emph{Інваріантність.} Нейрон ступеня $C$ збирає виходи багатьох $S$-нейронів, що шукають те саме поєднання в різних місцях, і спрацьовує, якщо воно знайшлося хоч десь. Це АБО за положеннями --- а точніше, вибір найбільшого.

Для одновимірної моделі обидві операції записуються так:
\begin{empheq}[box=\keybox]{equation}
s_{f,p} \;=\; \Bigl[\sum_i t_{f,i}\,x_{p+i} - \theta\Bigr]_+ ,
\qquad
c_f \;=\; \max_p\, s_{f,p} ,
\label{eq:scstage}
\end{empheq}
де $x$ --- вхід, $t_{f,i}$ --- зразок ознаки $f$, $p$ --- положення, $\theta$ --- поріг. Перший вираз робить відповідь вибірковою, другий --- незалежною від положення. Найбільше з багатьох входів мережа отримує засобами розділу~\ref{sec:gaba}: взаємне гальмування залишає найсильніший вхід (твердження~\ref{prop:wta}), а ділення на загальну активність вирівнює відповіді за різної яскравості~\cite{Carandini2012}.

\begin{figure}[t]
\centering
\begin{tikzpicture}[>=Stealth,
  px/.style={draw,minimum size=0.36cm,inner sep=0pt},
  on/.style={px,fill=blue!55!black},
  snode/.style={circle,draw,very thick,minimum size=0.62cm,inner sep=0pt,font=\scriptsize,fill=blue!6},
  cnode/.style={circle,draw,very thick,minimum size=0.8cm,inner sep=0pt,font=\scriptsize,fill=red!10}]
% two inputs: the same shape at two positions
\foreach \row/\sh/\lab in {0/1/{кадр 1},1/7/{кадр 2}}{
  \begin{scope}[yshift=-\row*1.0cm]
  \node[left,font=\scriptsize] at (-0.25,0) {\lab};
  \foreach \i in {0,...,13}{\node[px] at (\i*0.4,0) {};}
  \foreach \d in {0,1,3,4}{\node[on] at ({(\sh+\d)*0.4},0) {};}
  \end{scope}}
% S stage: detectors of the shape at several positions
\foreach \k in {0,...,5}{
  \node[snode] (S\k) at ({0.8+0.8*\k},1.7) {$S$};}
\node[font=\scriptsize,left] at (-0.25,1.7) {ступінь $S$: І за частинами};
\draw[->,thick,blue!60!black] (1.2,0.25) -- (S1.south);
\draw[->,thick,orange!85!black] (3.6,0.25) -- (S4.south);
\node[font=\scriptsize,blue!60!black,anchor=east] at (1.25,0.85) {кадр 1};
\node[font=\scriptsize,orange!85!black,anchor=west] at (3.9,0.85) {кадр 2};
% C stage
\node[cnode] (C) at (2.8,3.25) {$C$};
\node[font=\scriptsize,left] at (-0.25,3.25) {ступінь $C$: АБО за положеннями};
\foreach \k in {0,...,5}{\draw[->,gray!60!black] (S\k.north) -- (C);}
\draw[->,very thick,red!70!black] (C.north) -- ++(0,0.6) node[above,font=\scriptsize,black] {«фігура є» --- в обох кадрах};
% next stages
\draw[decorate,decoration={brace,amplitude=5pt},gray!60!black] (6.3,1.4) -- (6.3,-1.3);
\node[right,font=\scriptsize,align=left] at (6.5,0.05) {той самий\\об'єкт, інші\\пікселі};
\draw[->,thick,densely dashed,gray!60!black] (3.6,3.25) -- (5.9,3.25) node[right,font=\scriptsize,black,align=left] {наступні пари\\$S$, $C$ --- більші,\\складніші, інваріантніші};
\end{tikzpicture}
\caption{Одна пара ступенів конвеєра. Нейрони $S$ шукають те саме поєднання частин у різних місцях; нейрон $C$ відповідає, якщо воно знайшлося хоч десь~\eqref{eq:scstage}. Фігура в кадрі 1 і в кадрі 2 збуджує різні $S$-нейрони, але той самий $C$. Наступні пари ступенів повторюють те саме на дедалі більших і складніших поєднаннях.}
\label{fig:conveyor}
\end{figure}

У тій самій одновимірній моделі пара ступенів~\eqref{eq:scstage} упізнає фігуру в будь-якому місці без помилок, а якщо кожну точку входу випадково перевернути з імовірністю $0.1$ --- у $86\%$ випадків, з імовірністю $0.2$ --- у $70\%$. Монетка, як і раніше, дає половину. Покладіть кілька таких пар одну над одною: перша шукає краї, наступна --- поєднання країв, ще вище --- частини предметів, нагорі --- предмети цілком. З кожною парою поєднання більші, а відповідь дедалі менше залежить від того, де і як лежить предмет~\cite{Riesenhuber1999}.

\begin{keyblock}
\begin{proposition}[Конвеєр вибірковості й інваріантності]
\label{prop:conveyor}
Швидке розпізнавання --- один прохід десятком ступенів. На кожній парі ступенів І за частинами~\eqref{eq:scstage} робить відповідь вибірковою, а АБО за положеннями --- незалежною від зсуву. Чергування цих двох операцій дає відповідь, яка розрізняє предмети й майже не залежить від того, де і як їх видно. Будову перших ступенів і швидкість проходу виміряно; те, що весь ланцюжок зводиться до цих двох операцій, --- спрощення.
\end{proposition}
\end{keyblock}

Якщо ця будова здається знайомою, так і має бути. Саме її чергування --- пошук ознаки в усіх місцях, потім об'єднання за сусідніми місцями --- інженери перенесли в штучні \emph{згорткові} мережі~\cite{Fukushima1980}. А нещодавно пішли у зворотному напрямку: згорткові мережі, навчені розпізнавати предмети, виявилися найкращою з відомих моделей відповідей нейронів на верхніх ступенях зору~\cite{Yamins2014}. Результат нагорі читається просто: за сумарними частотами близько сотні нейронів останнього ступеня лінійний вирішальний блок упізнає предмет через десяту частку секунди після показу~\cite{Hung2005}. Це частотний код групи з розділу~\ref{sec:code}.

\section{Учитель, якого немає: вчитися на відео}

Згорткову мережу навчають на мільйонах картинок із підписами. Мозку підписів майже ніхто не дає: дитина бачить предмети годинами, а слово «чашка» чує зрідка. Звідки тоді ступені $C$ знають, які $S$-нейрони об'єднувати? Адже щоб об'єднати «цю фігуру тут» із «цією фігурою там», треба вже знати, що це та сама фігура.

Підказку дає саме відео. Світ змінюється неперервно: предмет у полі зору лишається тим самим предметом, доки він зсувається, повертається й змінює освітлення, і лише зрідка погляд переходить на інший. Усе, що змінюється \emph{повільно}, найімовірніше, належить до предмета, а все, що змінюється швидко, --- до його положення й вигляду~\cite{WiskottSejnowski2002}. Отже, $C$-нейронові досить правила Гебба з розділу~\ref{sec:dofa} зі слідом: пов'язувати те, що йшло поспіль, а не лише те, що йшло одночасно~\cite{Foldiak1991}:
\begin{empheq}[box=\keybox]{equation}
\tau_{\text{сл}}\,\frac{\dd\bar y_k}{\dd t} \;=\; -\bar y_k + y_k ,
\qquad
\frac{\dd\mathbf w_k}{\dd t} \;=\; \eta\,\bar y_k\,\bigl(\mathbf x - \mathbf w_k\bigr) .
\label{eq:traceinvariance}
\end{empheq}
Тут $y_k$ --- активність $C$-нейрона $k$, який виграв змагання через гальмування, $\bar y_k$ --- її слід, те саме RC-коло, що й у правилі~\eqref{eq:threefactor}, а $\mathbf x$ --- виходи $S$-нейронів. Нейрон, що виграв на першому вигляді предмета, ще пам'ятає це, коли надходить другий вигляд, і підтягує до себе й його.

\begin{figure}[t]
\centering
\begin{tikzpicture}[>=Stealth,x=1.0cm,y=1.0cm]
\foreach \i/\lab in {0/1,1/2,2/3,3/4}{
  \node[draw,thick,fill=blue!8,minimum width=0.9cm,minimum height=0.55cm,font=\scriptsize] at (\i+0.5,2.2) {$A$ у \lab};}
\foreach \i/\lab in {0/4,1/3,2/2,3/1}{
  \node[draw,thick,fill=orange!12,minimum width=0.9cm,minimum height=0.55cm,font=\scriptsize] at (\i+6.5,2.2) {$B$ у \lab};}
\node[font=\scriptsize,gray!40!black] at (5.0,2.2) {пауза};
\draw[->,gray!70!black] (0,0) -- (10.4,0) node[below left,font=\scriptsize,black] {час};
% traces
\draw[thick,blue!60!black] (0,0.05) .. controls (0.4,1.2) and (0.8,1.3) .. (1.2,1.35) -- (4,1.4) .. controls (4.6,0.5) and (5.2,0.15) .. (6,0.08) -- (10,0.05);
\draw[thick,orange!85!black] (0,0.05) -- (6,0.05) .. controls (6.4,1.2) and (6.8,1.3) .. (7.2,1.35) -- (10,1.4);
\node[font=\scriptsize,blue!60!black,anchor=west] at (1.4,1.65) {слід $\bar y_1$ --- тримається весь прохід $A$};
\node[font=\scriptsize,orange!85!black,anchor=west] at (7.2,1.65) {слід $\bar y_2$};
\draw[decorate,decoration={brace,amplitude=4pt},gray!60!black] (4.0,2.6) -- (0,2.6);
\draw[decorate,decoration={brace,amplitude=4pt,mirror},gray!60!black] (0,-0.35) -- (4,-0.35) node[midway,below=4pt,font=\scriptsize,black] {усі вигляди $A$ пов'язуються з нейроном 1};
\end{tikzpicture}
\caption{Як відео навчає інваріантності, схема. Предмет $A$ проходить через чотири положення, потім пауза, потім предмет $B$. Слід~\eqref{eq:traceinvariance} нейрона, що виграв на першому вигляді $A$, тримається весь прохід, і всі вигляди $A$ виявляються пов'язаними з одним нейроном. Пауза стирає слід, і $B$ дістається іншому нейронові.}
\label{fig:tracevideo}
\end{figure}

Ми перевірили це на моделі (рис.~\ref{fig:tracevideo}). Два $C$-нейрони змагаються за входи від $16$ $S$-нейронів: дві фігури у восьми положеннях. Фігура проїжджає через усі положення, по $50\ms$ на кожне, потім пауза $0.3\,\text{с}$ і наступна випадкова фігура. Жодних підписів. Якщо слід короткий, $\tau_{\text{сл}}=5\ms$, нейрони ділять вхід за \emph{положенням}, і відповідь збігається з фігурою не краще, ніж у монетки: $52\%$ у середньому за десять прогонів. Якщо слід порівнянний із часом показу одного положення, від $\tau_{\text{сл}}\approx20\ms$, кожен нейрон збирає \emph{всі} положення своєї фігури (за $20$, $50$ і $200\ms$ --- в усіх десяти прогонах): інваріантність вивчено на самому лише відео. Пауза між предметами важлива: без неї слід перетікає на наступний предмет і плутає їх.

Те саме видно й у досліді. Якщо в експерименті непомітно підміняти предмет, доки око перескакує з одного місця на інше, то нейрони верхніх ступенів за годину-дві починають відповідати на підмінений предмет як на той самий: вони пов'язують те, що йшло поспіль~\cite{LiDiCarlo2008}. Інваріантність у мозку справді вчиться на неперервності часу. Наскільки цим правилом пояснюється все навчання зору, --- гіпотеза.

\section{Відео --- це рух}

Відео --- не лише потік кадрів для навчання, а й сама задача: що рухається, куди і як швидко. Перший крок нам знайомий --- детектор напрямку розділу~\ref{sec:gaba} (рис.~\ref{fig:direction}), у якому запізніле гальмування гасить рух в один бік. Такі детектори стоять по всьому полю зору, і далі з ними чинять так само, як з ознаками форми: ступені, що збирають багато детекторів одного напрямку в різних місцях, відповідають на рух великого предмета незалежно від того, де він. Це та сама пара І та АБО~\eqref{eq:scstage}, тільки ознака не «край», а «край, що зсунувся за час $d$».

Відео ще й передбачуване: наступний кадр майже такий самий, як попередній. За гіпотезою \emph{передбачувального кодування} верхні ступені надсилають нижнім передбачення, а вгору йде головно помилка --- те, чого передбачення не врахувало~\cite{RaoBallard1999}. Це та сама економія імпульсів, що в твердженні~\ref{prop:economy}: платити за новини, а не за те, що й так відомо. Окремі спостереження з цією гіпотезою узгоджуються, але як загальну будову конвеєра її не доведено.

\section{Мозок і згорткова мережа}

Наскільки далеко сягає схожість? Для швидкого розпізнавання одного предмета вона справді глибока~\cite{Yamins2014}. Але в мозку є й те, чого немає в простій згортковій мережі.

\emph{Зворотні зв'язки.} Конвеєр мозку не лише прямий: у кожного ступеня є зв'язки назад і вбік. Для складних картинок --- з перекриттям, за поганого світла --- відповідь верхніх ступенів складається пізніше, і ці пізні відповіді краще описують мережі зі зворотними зв'язками, ніж прямі~\cite{Kar2019}. Один прохід розв'язує легкі випадки, складні потребують кількох кіл.

\emph{Стійкість.} Штучну мережу можна обдурити непомітною для ока підробкою пікселів~\cite{Szegedy2014}: зміна, якої людина не бачить, перетворює «панду» на «гібона»~\cite{Goodfellow2015}. Зір людини до таких підробок набагато стійкіший, але не невразливий: картинки, що обманюють одразу багато мереж, за короткого показу зсувають і вибір людини~\cite{Elsayed2018}. Чому стійкість так відрізняється --- питання відкрите; кандидати --- шум, ділення на загальну активність і зворотні зв'язки.

\emph{Учитель.} Мережі потрібні мільйони підписаних прикладів, мозку --- здебільшого відео без підписів~\eqref{eq:traceinvariance} і трохи підказок від \nt{DOPA} про те, що важливо.

\emph{Енергія.} Увесь мозок --- близько $20$ ватів (твердження~\ref{prop:economy}), а розпізнавання займає лише частину цього; навчена згорткова мережа на графічному прискорювачі витрачає сотні ватів.

\section{Ілюзії: де машина помиляється і чому}
\label{sec:illusions}

Інженер, який хоче зрозуміти чужу схему, подає на неї незвичні сигнали й дивиться, де вона помиляється. Для зору такі сигнали давно придумано --- це зорові ілюзії. Ілюзія --- не поломка. Кожна вказує на певний механізм машини, який у звичайних умовах працює правильно, а на спеціально дібраній картинці себе виказує.

\emph{Розумні очікування.} Вхід шумний, і машина доповнює його тим, що зазвичай буває у світі. Повільні рухи трапляються частіше за швидкі; коли вхід ненадійний --- наприклад, картинка малоконтрастна, --- оцінка швидкості зсувається до повільної, і блідий предмет здається рухомим повільніше за яскравий~\cite{Weiss2002}. Так само пояснюють багато ілюзій яскравості: ми бачимо не яскравість пікселя, а те, якій поверхні така яскравість найчастіше відповідала в минулому~\cite{Purves2011}. Фотографія сукні, яку одні бачать біло-золотою, а інші синьо-чорною, --- той самий механізм: картинка двозначна, і люди розходяться в тому, яким вони неусвідомлено вважають освітлення~\cite{LaferSousa2015,Wallisch2017}. Відмінності між людьми виміряно; те, що мозок тут поєднує вхід з очікуванням за правилами теорії ймовірностей, --- добре обґрунтована гіпотеза.

\emph{Регулювання підсилення.} Подивіться хвилину на водоспад, потім на нерухоме каміння: каміння попливе вгору. Канали «вниз» і «вгору» --- пара проводів, як у~\eqref{eq:dualrail}; регулювання підсилення~\eqref{eq:nakarushton} приглушило втомлений канал, і рівновага пари зсунулася. Ілюзії нахилу й контрасту, в яких оточення змінює вигляд центру, пояснюються діленням на активність сусідів~\cite{Schwartz2009}, як у розділі~\ref{sec:gaba}. Є й повчальний приклад обережності: темні плями на перехрестях білих смуг ґратки довго пояснювали простим гальмуванням сусідів, але досліди зі зміненою ґраткою показали, що це пояснення не працює~\cite{SchillerCarvey2005}.

\emph{Компенсація затримок.} Шлях від ока до верхніх ступенів займає десятки мілісекунд, і рухомий предмет за цей час устигає зсунутися. Якщо поруч із ним спалахне нерухома точка, предмет здається попереду спалаху, хоча вони збігалися~\cite{Nijhawan1994}. За одним із пояснень машина продовжує рух уперед, щоб компенсувати затримку, --- те саме, що в розділі~\ref{sec:muscles}: за запізнілого зворотного зв'язку доводиться передбачати. Пояснень у цієї ілюзії кілька, і суперечку не розв'язано.

\emph{Помилка в коді часу.} Нерухома картинка з кілець із різкими колірними переходами здається обертовою. Контрастні ділянки обробляються швидше за бліді, і різниця затримок~\eqref{eq:latency} читається детекторами напрямку як рух~\cite{Conway2005}. Запускають ілюзію дрібні мимовільні стрибки ока й кліпання: кожен стрибок оновлює вхід, і детектори знову бачать «рух»~\cite{OteroMillan2012}. Це код часом із розділу~\ref{sec:code}, прочитаний неправильно.

\emph{Дві картинки в одній.} Каркас куба, що дивиться на нас то однією гранню, то іншою, або різні картинки у двох очах: сприйняття перескакує кожні кілька секунд. Найкращі моделі --- взаємне гальмування двох тлумачень, повільна втома й шум~\cite{MorenoBote2007}, тобто та сама схема~\eqref{eq:halfcentre}, яка в розділі~\ref{sec:muscles} породжувала ритм кроку. Час зміни задають шум і втома; за моделлю~\cite{MorenoBote2007} головну роль відіграє шум, а втома лише допомагає.

А штучні мережі? Мережа, навчена передбачати наступний кадр відео, бачить обертання на тих самих нерухомих кільцях і не бачить його на контрольних картинках~\cite{Watanabe2018}. Мережі, навчені очищати зображення від шуму й підвищувати їхню різкість, піддаються тим самим ілюзіям кольору та яскравості, що й люди~\cite{GomezVilla2020}. Отже, частина ілюзій --- наслідок самої задачі, якої вчиться машина, а не особливостей нервової тканини. Які ілюзії властиві лише мозкові, --- добра перевірка для будь-якої моделі зору.

\begin{keyblock}
\begin{proposition}[Ілюзії як відбитки механізмів]
\label{prop:illusions}
Ілюзія виникає, коли механізм, корисний у звичайному світі, отримує незвичний вхід: очікування перемагає ненадійний вхід, регулювання підсилення зсуває пару каналів, компенсація затримки виносить предмет уперед, різниця затримок читається як рух, взаємне гальмування із шумом і втомою перескакує. Кожна ілюзія --- випробувальний сигнал, що вказує на свій механізм. Самі ілюзії виміряно; їхні пояснення --- від добре обґрунтованих до спірних.
\end{proposition}
\end{keyblock}

\section{Підсумок}

\begin{table}[t]
\centering
\footnotesize
\setlength{\tabcolsep}{4pt}
\renewcommand{\arraystretch}{1.15}
\begin{tabular}{>{\raggedright}p{3.0cm}>{\raggedright}p{4.6cm}>{\raggedright\arraybackslash}p{5.2cm}}
\toprule
Що робить машина & Як & Що відомо \\
\midrule
упізнає за $150\ms$ & один прохід десятком ступенів & швидкість виміряно \\
робить відповідь вибірковою & І за частинами, ступінь $S$~\eqref{eq:scstage} & виміряно на перших ступенях \\
робить відповідь інваріантною & АБО за положеннями, ступінь $C$; гальмування й ділення & виміряно на перших ступенях; для верхніх --- спрощення \\
видає відповідь & частоти близько сотні нейронів останнього ступеня, лінійне читання & виміряно \\
вчиться без підписів & правило Гебба зі слідом~\eqref{eq:traceinvariance} на неперервному відео & підміна предмета змінює відповіді --- виміряно; як головне правило --- гіпотеза \\
бачить рух & детектори напрямку, далі та сама пара І/АБО & виміряно \\
заощаджує на передбачуваному & угору йде помилка передбачення & гіпотеза \\
розв'язує складні випадки & зворотні зв'язки, кілька кіл & пізні відповіді й роль зворотних зв'язків виміряно \\
помиляється передбачувано & ілюзії: очікування, регулювання підсилення, затримки, взаємне гальмування із шумом і втомою & ілюзії виміряно; пояснення --- від обґрунтованих до спірних \\
\bottomrule
\end{tabular}
\caption{Як машина розпізнає зображення і відео.}
\label{tab:recognition}
\end{table}

Розпізнавання зібрано з деталей, які в нас уже були (таблиця~\ref{tab:recognition}): поріг дає І за частинами, взаємне гальмування --- АБО за положеннями, ділення --- незалежність від яскравості, правило Гебба зі слідом --- учителя, якого немає. Інженери взяли в зору чергування вибірковості й інваріантності та збудували на ньому згорткові мережі; навзаєм мозок поки не віддав головного --- вміння вчитися на відео без підписів майже без енергії.

\section{Задачі}

\begin{exercise}[\lvlA]
\label{ex:12:1}
\emph{Частота на одному ступені.} Ступінь конвеєра триває $15\ms$, нейрони працюють із частотою $20$ імпульсів за секунду, шуми корельовані з $c=0.1$~\eqref{eq:correlated}. Яка відносна похибка частоти групи зі ста нейронів і її межа для як завгодно великої групи? Скільки рівнів частоти розрізненні на ступені?
\end{exercise}

\begin{exercise}[\lvlA]
\label{ex:12:2}
\emph{Енергія одного впізнавання.} За один прохід за $150\ms$ десять ступенів по $10^7$ нейронів видають не більше одного імпульсу ціною $10^{-10}$~Дж. (а)~Яку частку енергії всього мозку за ці $150\ms$ коштує впізнавання? (б)~У скільки разів більше витрачає прискорювач потужністю $300$~Вт, що впізнає картинку за $10\ms$?
\end{exercise}

\begin{exercise}[\lvlB]
\label{ex:12:3}
\emph{Пороги ступеня $S$.} «Сітківка» з $24$ точок, фігури $A=11011$ і $B=10101$, зразок~\eqref{eq:scstage} дорівнює $+1$ на одиницях фігури й $-1$ на нулях. (а)~За яких порогів $S$-нейрон прощає одну перевернуту точку, але мовчить на чужій фігурі? (б)~Скільки $S$-нейронів потрібно для десяти ознак $5\times5$ на картинці $100\times100$?
\end{exercise}

\begin{exercise}[\lvlB]
\label{ex:12:4}
\emph{Скільки триває одна картинка з двох.} У схемі перескоків~\eqref{eq:halfcentre} при $\tau\ll\tau_a$, доки перемагає група~1, $r_2=0$ і $r_1=I-a_1$. (а)~Знайдіть тривалість перемоги при $I=1$, $\beta=2$, $b=2$, $\tau_a=0.4\,\text{с}$. (б)~Яке $\tau_a$ дає зміну картинки кожні три секунди?
\end{exercise}

\begin{exercise}[\lvlB]
\label{ex:12:5}
\emph{Моделювання: ціна інваріантності.} Функцією \texttt{two\_stage} з \texttt{models/\allowbreak recognition\_models.py} ($4000$ спроб) знайдіть, за якої ймовірності перевертання точки $p$ пара ступенів $S$, $C$ при $N=24$ помиляється в кожному четвертому випадку. Як змінюється частка правильних відповідей при $p=0.1$ від $N=8$ до $N=192$?
\end{exercise}

\begin{exercise}[\lvlB]
\label{ex:12:6}
\emph{Моделювання: вікно сліду.} Десять прогонів \texttt{trace\_learning} з \texttt{models/\allowbreak recognition\_models.py} з паузою $0.3\,\text{с}$: знайдіть, за яких $\tau_{\text{сл}}$ від $5\ms$ до $2\,\text{с}$ усі прогони вивчають інваріантність без помилок. Звідки в цього вікна верхня межа?
\end{exercise}

\begin{exercise}[\lvlB]
\label{ex:12:7}
\emph{Рух у будь-якому місці.} На сусідніх точках ряду з кроком $0.1^\circ$ стоять детектори напрямку розділу~\ref{sec:gaba}: гальмування від $A$ із затримкою $d$ гасить збудження від $B$, якщо вони розійшлися не більш ніж на $5\ms$; над ними --- ступінь $C$. Які затримки потрібні, щоб $C$ мовчав за зворотного руху зі швидкістю від $5$ до $20^\circ$ за секунду?
\end{exercise}

\begin{exercise}[\lvlC]
\label{ex:12:8}
\emph{Підробка пікселів.} Скільки перевернутих точок потрібно, щоб змінити відповідь моделі \texttt{two\_stage} ($N=24$) на чистій фігурі? А класифікатора «понад $3.5$ одиниці на вході --- це $A$», теж інваріантного до зсуву? Яка його частка правильних відповідей при $p=0.1$ проти $0.86$ у пари ступенів?
\end{exercise}

\chapter{Виходи машини: як мозок керує м'язами}
\chaptermark{Виходи машини}
\label{sec:muscles}

\section{Швидкий вихід}

Усе, що машина робить у світі швидко, вона робить м'язами. Слово, погляд, крок, усмішка --- це скорочення м'язів. Другий, повільний вихід --- хімію тіла --- розібрано в розділі~\ref{sec:chemoutput}. На рис.~\ref{fig:machine} вихід був стрілкою «м'язи»; тепер подивімося, що за цією стрілкою.

Остання ланка кола --- \nt{ACET}-нейрони, ті самі, в яких у таблиці~\ref{tab:types} записано «вихід на м'яз». Кожне м'язове волокно отримує вхід рівно від одного такого нейрона, а один нейрон керує групою волокон --- від кількох сотень до пари тисяч~\cite{Feinstein1955mu}. У м'язах, яким потрібне тонке підлаштування, одиниці дрібніші, у великих м'язах ніг --- більші. Нейрон разом зі своїми волокнами називають \emph{руховою одиницею}: це найменший шматок м'яза, який мозок може ввімкнути окремо.

Синапс нейрона на м'язі влаштований зовсім не так, як синапси кори з розділу~\ref{sec:code}. Там більшість імпульсів губилася. Тут кожен імпульс вивільняє в кілька разів більше ацетилхоліну, ніж потрібно, щоб волокно спрацювало~\cite{WoodSlater2001}. Це синапс із \emph{запасом надійності}: один імпульс нейрона --- рівно одне скорочення волокон. Усередині машини можна дозволити собі ненадійні зв'язки й виправляти помилки надлишковістю; на виході помилка коштує руху, і надійність закладено просто в залізо.

\section{Сила: частота й кількість одиниць}

Один імпульс дає поодиноке скорочення --- посмикування: сила наростає за кілька десятків мілісекунд і спадає. Добре наближення для нього~\cite{Fuglevand1993}
\begin{equation}
h(t) \;=\; F_1\,\frac{t}{T_c}\,e^{\,1-t/T_c},
\label{eq:twitch}
\end{equation}
де $T_c$ --- час наростання, порядку $50\ms$, а $F_1$ --- сила одного посмикування. Якщо імпульси йдуть частіше, посмикування додаються:
\begin{empheq}[box=\keybox]{equation}
F(t) \;=\; \sum_k h\bigl(t-t_k\bigr) .
\label{eq:forcesum}
\end{empheq}
Це той самий фільтр нижніх частот, що перетворював імпульси на концентрацію в розділі~\ref{sec:serocomputer}, тільки тепер на виході не концентрація, а сила. М'яз --- цифро-аналоговий перетворювач з імпульсів у силу (рис.~\ref{fig:tetanus}). У нашій моделі за п'яти імпульсів на секунду посмикування майже не перекриваються й сила стрибає від $0.2F_1$ до $1.1F_1$; за п'ятнадцяти вона вже тримається між $1.8F_1$ і $2.2F_1$; за сорока стає майже рівною, близько $5.4F_1$. Справжній м'яз за частих імпульсів насичується, тож лінійна сума~\eqref{eq:forcesum} правильна лише до кількох десятків імпульсів за секунду.

\begin{figure}[t]
\centering
\begin{tikzpicture}[>=Stealth,x=20cm,y=0.55cm]
\draw[->,gray!70!black] (0,0) -- (0.66,0) node[right,font=\small,black] {$t$, с};
\draw[->,gray!70!black] (0,0) -- (0,6.4) node[left,font=\small,black] {$F/F_1$};
\foreach \t in {0,0.2,0.4,0.6}{\draw[gray!60!black] (\t,-0.1) -- (\t,0.1); \node[below,font=\scriptsize] at (\t,0) {\t};}
\foreach \f in {1,3,5}{\draw[gray!60!black] (-0.004,\f) -- (0.004,\f); \node[left,font=\scriptsize] at (0,\f) {\f};}
\draw[thick,blue!60!black] plot file {figures/muscle_twitch_5.dat};
\draw[thick,green!45!black] plot file {figures/muscle_twitch_15.dat};
\draw[thick,red!70!black] plot file {figures/muscle_twitch_40.dat};
\node[font=\scriptsize,blue!60!black,anchor=west] at (0.605,0.9) {5 імп/с};
\node[font=\scriptsize,green!45!black,anchor=west] at (0.605,2.1) {15 імп/с};
\node[font=\scriptsize,red!70!black,anchor=west] at (0.605,5.4) {40 імп/с};
\end{tikzpicture}
\caption{М'яз як цифро-аналоговий перетворювач: сила~\eqref{eq:forcesum} за регулярних імпульсів однієї рухової одиниці, $T_c=50\ms$. Рідкі імпульси дають окремі посмикування, часті зливаються в рівну силу --- так само, як часті імпульси зливаються в рівну концентрацію на рис.~\ref{fig:lowpass}.}
\label{fig:tetanus}
\end{figure}

У мозку дві ручки сили: частота імпульсів кожної одиниці й кількість увімкнених одиниць. І друга ручка влаштована дуже винахідливо. Одиниці в м'язі різні: найслабші в сотню разів слабші за найсильніші. Коли потрібна дедалі більша сила, одиниці вмикаються \emph{за порядком розміру}, від малих до великих~\cite{Henneman1957}. Ніхто цього порядку не обчислює: малі нейрони просто легше збуджуються від того самого спільного входу, тож порядок увімкнення задано самим залізом.

Що це дає? Нехай кожна наступна за порядком одиниця сильніша за попередню в $q$ разів, де $q$ трохи більше за одиницю. Сила всіх увімкнених одиниць --- сума геометричної прогресії, і майже вся вона припадає на останні одиниці. Тоді чергова ввімкнена одиниця додає
\begin{empheq}[box=\keybox]{equation}
\frac{\Delta F}{F} \;\approx\; q-1 \;=\; \text{const} ,
\label{eq:sizeprinciple}
\end{empheq}
тобто крок сили пропорційний самій силі. За слабкого зусилля мозок дозує його дрібними порціями, за сильного --- великими. Це \emph{число з рухомою комою}: порядок задає кількість увімкнених одиниць, а мантису --- частота їхніх імпульсів. Та сама логарифмічна шкала, що в датчиків розділу~\ref{sec:sensors}: вхід і вихід машини міряють світ у відносних частках.

Рухома кома має зворотний бік. Розкид сили теж зростає пропорційно їй самій: сильний рух неминуче менш точний, ніж слабкий. За однією впливовою гіпотезою, саме цей шум, що залежить від сигналу, пояснює, чому наші рухи плавні: плавна команда дає найменший розкид у кінцевій точці~\cite{HarrisWolpert1998}.

\section{Тягнути, але не штовхати: два проводи на суглоб}

М'яз уміє лише тягнути. Штовхати він не може, як не може \nt{GLUT}-нейрон видати від'ємний сигнал. Рішення знайоме з розділу~\ref{sec:glutcomputer}: \emph{два проводи}. Кожен суглоб обслуговує пара м'язів, що тягнуть у протилежні боки (рис.~\ref{fig:antagonists}). Якщо їхні сили $F_1$ і $F_2$, а плече $\ell$, то обертальний момент і жорсткість суглоба дорівнюють
\begin{empheq}[box=\keybox]{equation}
M \;=\; \ell\,(F_1-F_2), \qquad K \;\approx\; \kappa_m\,(F_1+F_2),
\label{eq:antagonists}
\end{empheq}
де $\kappa_m$ --- коефіцієнт, що пов'язує силу м'яза з його жорсткістю: напружений м'яз пружинить сильніше за розслаблений. Різниця сил задає момент, сума --- жорсткість~\cite{Hogan1984}. Двома проводами мозок керує двома незалежними величинами: куди повернути суглоб і наскільки твердо його тримати. Коли треба втримати чашку, не розхлюпавши, ми напружуємо обидва м'язи одразу: момент той самий, жорсткість вища, і випадковий поштовх менше збиває руку.

\begin{figure}[t]
\centering
\begin{tikzpicture}[>=Stealth,
  nrn/.style={circle,draw,very thick,minimum size=0.95cm,inner sep=0pt,font=\scriptsize},
  inh/.style={-{Bar[width=7pt]},thick,red!70!black}]
% the joint: a bar on a pivot, two springs pulling down on either side
\fill[gray!30] (4.6,-0.25) -- (5.4,-0.25) -- (5,0.15) -- cycle;
\draw[line width=3pt,gray!50!black] (2.4,0.2) -- (7.6,0.2);
\fill[black] (5,0.2) circle (0.08);
\draw[decorate,decoration={coil,segment length=5pt,amplitude=4pt},very thick,blue!60!black] (3.0,0.2) -- (3.0,-1.6);
\draw[decorate,decoration={coil,segment length=5pt,amplitude=4pt},very thick,orange!85!black] (7.0,0.2) -- (7.0,-1.6);
\draw[very thick] (2.5,-1.6) -- (3.5,-1.6); \draw[very thick] (6.5,-1.6) -- (7.5,-1.6);
\node[font=\small,blue!60!black] at (2.35,-0.75) {$F_1$};
\node[font=\small,orange!85!black] at (7.65,-0.75) {$F_2$};
\draw[->,thick] (5.9,0.75) arc (60:120:1.8) node[above,font=\scriptsize,pos=0.5] {$M=\ell(F_1-F_2)$};
% motor neurons and reflex circuit
\node[nrn,fill=green!12] (M1) at (1.6,-3.0) {\nt{ACET}};
\node[nrn,fill=green!12] (M2) at (8.4,-3.0) {\nt{ACET}};
\node[nrn,fill=red!10] (G) at (5.0,-3.0) {\nt{GLYC}};
\node[nrn,fill=blue!6] (S) at (1.6,-5.0) {\nt{GLUT}};
\draw[->,very thick] (M1) -- (3.0,-1.7);
\draw[->,very thick] (M2) -- (7.0,-1.7);
\draw[->,thick] (S) -- (M1);
\draw[->,thick] (S) -- (G);
\draw[inh] (G) -- (M2);
\draw[->,thick,densely dashed,gray!60!black] (3.15,-1.2) to[bend left=25] node[pos=0.35,right,font=\scriptsize,black] {розтяг} (S.north east);
\draw[->,thick,blue!60!black] (-0.3,-3.0) node[left,font=\scriptsize,black,align=right] {команда\\мозку} -- (M1);
\draw[->,thick,blue!60!black] (10.3,-3.0) node[right,font=\scriptsize,black,align=left] {команда\\мозку} -- (M2);
\end{tikzpicture}
\caption{Два м'язи на один суглоб і найшвидша петля зворотного зв'язку. \nt{ACET}-нейрони отримують команди мозку й тягнуть суглоб у різні боки; різниця сил повертає його, сума робить жорсткішим. Датчик розтягу (\nt{GLUT}) лівого м'яза збуджує його власний \nt{ACET}-нейрон і через \nt{GLYC}-посередника гальмує нейрон м'яза-антагоніста --- заборона з розділу~\ref{sec:gaba}.}
\label{fig:antagonists}
\end{figure}

\section{Зворотний зв'язок із затримкою}

М'язи не працюють наосліп. У кожному сидять датчики розтягу, про які йшлося в таблиці~\ref{tab:sensors}, і найкоротша петля замикається просто в спинному мозку (рис.~\ref{fig:antagonists}). М'яз несподівано розтягнули --- датчик розтягу збуджує його власний \nt{ACET}-нейрон, м'яз скорочується й повертає суглоб. Водночас через гальмівний нейрон-посередник вимикається м'яз-антагоніст, щоб не заважав. Цим посередником працює \nt{GLYC} --- той самий тип із таблиці~\ref{tab:types}, якому не дістався власний розділ: його місце саме тут, у швидких петлях спинного мозку.

Кожна петля має затримку $d$: спинномозкова --- близько $25$--$30\ms$ для руки, петля через кору --- у півтора--три рази більшу, петля через око --- $100$--$200\ms$. А затриманий зворотний зв'язок, як ми бачили в твердженні~\ref{prop:ei}, розгойдує систему. Для найпростішого регулятора, що виправляє відхилення $x$ зі швидкістю $G$ за відомостями $d$-секундної давності, $\dd x/\dd t=-G\,x(t-d)$, умова стійкості така~\cite{Hayes1950}:
\begin{empheq}[box=\keybox]{equation}
G\,d \;<\; \frac{\pi}{2} ,
\label{eq:delaystable}
\end{empheq}
а на межі система коливається з частотою $1/(4d)$. Ми перевірили це на моделі: при $d=30\ms$ відхилення загасає при $G=40$ за секунду, а при $G=55$ за секунду, трохи вище за межу $52$, розгойдується.

Звідси два наслідки. Перший: що довша петля, то повільніше вона мусить виправляти. Через око із затримкою в півтори сотні мілісекунд руку не можна поправляти швидше, ніж приблизно за десяту частку секунди, інакше вона затремтить. Другий: межа стійкості спинномозкової петлі, $1/(4\cdot30\ms)\approx8$ коливань за секунду, близька до частоти дрібного тремтіння, яке є в кожної здорової людини, --- вісім-дванадцять коливань за секунду. Петля розтягу --- одне з імовірних джерел цього тремтіння~\cite{McAuleyMarsden2000}.

Як тоді мозок рухається швидко й точно, якщо зворотний зв'язок запізнюється? За поширеною гіпотезою --- \emph{передбачає}: тримає внутрішню модель тіла й обчислює, яким буде результат команди, не чекаючи на датчики~\cite{WolpertGhahramani2000}. Датчики потрібні, щоб виправляти модель, а не кожен рух. Інженер назвав би це керуванням із випередженням і спостерігачем стану.

\section{Генератори ритму рухів}

Ходьба, дихання, жування --- ритмічні рухи. Чи потрібен для них тактовий генератор? Ні. У спинному мозку й стовбурі є невеликі мережі, які самі породжують ритм, навіть якщо на них не надходить нічого ритмічного~\cite{MarderBucher2001}. Найпростішу таку мережу зібрано з деталей, які в нас уже є: дві групи \nt{GLUT}-нейронів, що гальмують одна одну через \nt{GABA}- або \nt{GLYC}-посередників, як на рис.~\ref{fig:wta}, і повільно втомлюються, як пам'ять розділу~\ref{sec:glutcomputer}:
\begin{empheq}[box=\keybox]{equation}
\tau\,\frac{\dd r_i}{\dd t} = -r_i + \bigl[I - \beta\,r_j - a_i\bigr]_+ ,
\qquad
\tau_a\,\frac{\dd a_i}{\dd t} = -a_i + b\,r_i ,
\label{eq:halfcentre}
\end{empheq}
де $a_i$ --- втома групи $i$, $j$ --- інша група. Взаємне гальмування з $\beta>1$ вибирає переможця (твердження~\ref{prop:wta}). Переможець працює і втомлюється; коли втома з'їдає його вхід, переможений, уже відпочилий, виривається вперед і сам починає втомлюватися. Групи працюють по черзі, і кожна керує своїм м'язом пари (рис.~\ref{fig:halfcentre}).

\begin{figure}[t]
\centering
\begin{tikzpicture}[>=Stealth,x=3.2cm,y=2.4cm]
\draw[->,gray!70!black] (0,0) -- (4.2,0) node[right,font=\small,black] {$t$, с};
\draw[->,gray!70!black] (0,0) -- (0,1.2) node[left,font=\small,black] {$r$};
\foreach \t in {0,1,2,3,4}{\draw[gray!60!black] (\t,-0.03) -- (\t,0.03); \node[below,font=\scriptsize] at (\t,0) {\t};}
\draw[thick,blue!60!black] plot file {figures/halfcentre1.dat};
\draw[thick,orange!85!black] plot file {figures/halfcentre2.dat};
\node[font=\scriptsize,blue!60!black,anchor=west] at (1.6,1.28) {група 1: «вперед»};
\node[font=\scriptsize,orange!85!black,anchor=west] at (2.7,1.28) {група 2: «назад»};
\end{tikzpicture}
\caption{Генератор ритму за моделлю~\eqref{eq:halfcentre}: $I=1$, $\beta=2$, $b=2$, $\tau=20\ms$, $\tau_a=0.4\,\text{с}$. Дві групи із взаємним гальмуванням і втомою працюють по черзі з періодом $0.84\,\text{с}$, хоча на вхід подано сталий сигнал.}
\label{fig:halfcentre}
\end{figure}

У нашій моделі за сталого входу групи змінюють одна одну з періодом $0.84\,\text{с}$, приблизно вдвічі більшим за час утоми $\tau_a$. Стала команда згори --- «іди» --- перетворюється на ритм на місці. Це ще один місцевий ритм, як ритм петлі \nt{GLUT}--\nt{GABA} у розділі~\ref{sec:gaba}: він виникає, лише коли мережа працює, і задає такт тільки тим м'язам, якими керує.

\begin{keyblock}
\begin{proposition}[Вихід машини]
\label{prop:muscles}
Мозок керує м'язами як ієрархія регуляторів. Нагорі вибирається дія (розділ~\ref{sec:dofa}); нижче місцеві генератори перетворюють сталі команди на ритм, а швидкі петлі спинного мозку тримають суглоби. Сила задається з рухомою комою: порядок --- кількістю одиниць, що вмикаються за розміром, мантиса --- частотою імпульсів. Кожним суглобом керує пара тягнучих м'язів, різниця сил яких задає момент, а сума --- жорсткість. Ці пристрої виміряно; що мозок спирається на внутрішню модель тіла, --- добре обґрунтована гіпотеза.
\end{proposition}
\end{keyblock}

\section{Підсумок}

\begin{table}[t]
\centering
\footnotesize
\setlength{\tabcolsep}{4pt}
\renewcommand{\arraystretch}{1.15}
\begin{tabular}{>{\raggedright}p{3.0cm}>{\raggedright}p{4.6cm}>{\raggedright\arraybackslash}p{5.2cm}}
\toprule
Що робить вихід & Як & Що відомо \\
\midrule
перетворює імпульси на силу & сума посмикувань~\eqref{eq:forcesum}, фільтр нижніх частот & виміряно; лінійна сума --- спрощення \\
дозує силу & увімкнення одиниць за розміром, рухома кома~\eqref{eq:sizeprinciple} & порядок увімкнення виміряно \\
штовхає, уміючи лише тягнути & пара м'язів: момент --- різниця, жорсткість --- сума~\eqref{eq:antagonists} & виміряно \\
тримає суглоб & петля розтягу, заборона антагоніста через \nt{GLYC} & виміряно \\
не тремтить & $Gd<\pi/2$~\eqref{eq:delaystable} & випливає з теорії регулювання; внесок у тремтіння --- імовірна причина \\
рухається швидко & передбачення за внутрішньою моделлю & добре обґрунтована гіпотеза \\
ходить і дихає ритмічно & генератор ритму~\eqref{eq:halfcentre} & генератори виміряно; найпростіша схема --- модель \\
\bottomrule
\end{tabular}
\caption{Як машина керує м'язами.}
\label{tab:muscles}
\end{table}

Вихід машини влаштовано тими самими прийомами, що й усе інше (таблиця~\ref{tab:muscles}): два проводи замість знака, фільтр нижніх частот замість ЦАП, логарифмічна шкала замість лінійної, взаємне гальмування й утома замість тактового генератора. Вхід (розділ~\ref{sec:sensors}) і вихід міряють світ у відносних частках, а між ними працює машина, якій присвячено цю книгу.

\section{Задачі}

\begin{exercise}[\lvlA]
\label{ex:13:1}
\emph{Рухома кома в цифрах.} У м'язі $N=100$ рухових одиниць із силами $f_n=f_1q^{\,n-1}$; найсильніша в $100$ разів сильніша за найслабшу. (а)~Знайдіть $q$, відносний крок~\eqref{eq:sizeprinciple} і динамічний діапазон у децибелах. (б)~Який крок за сили $10f_1$ у «лінійного ЦАП» зі $100$ однакових одиниць тієї самої сумарної сили?
\end{exercise}

\begin{exercise}[\lvlA]
\label{ex:13:2}
\emph{Ціна запасу надійності.} На кожен імпульс синапс на м'язовому волокні вивільняє пуассонівську кількість порцій із середнім $m$, а волокно спрацьовує за $n_{\text{пор}}=20$ порцій і більше; запас надійності $S=m/n_{\text{пор}}$. Скільки збоїв на добу дає одиниця, що працює з частотою $20$ імпульсів за секунду, при $S=2$ і при $S=4$?
\end{exercise}

\begin{exercise}[\lvlB]
\label{ex:13:3}
\emph{М'яз як фільтр.} Посмикування~\eqref{eq:twitch} з $T_c=50\ms$ --- імпульсна характеристика лінійної системи. (а)~Знайдіть її передавальну функцію $H(s)$ і частоту зрізу. (б)~За якої частоти регулярних імпульсів розмах пульсацій сили становить $10\%$ від середньої сили?
\end{exercise}

\begin{exercise}[\lvlB]
\label{ex:13:4}
\emph{Швидше за затримку не виправиш.} Регулятор $\dd x/\dd t=-G\,x(t-d)$. (а)~Покажіть, що найшвидше відхилення загасає без коливань при $Gd=1/e$, зі швидкістю $1/d$. (б)~Знайдіть це $G$ і сталу часу загасання для спинномозкової петлі, $d=30\ms$, і для петлі через око, $d=150\ms$.
\end{exercise}

\begin{exercise}[\lvlB]
\label{ex:13:5}
\emph{Період генератора ритму.} При $\tau\ll\tau_a$ період $T$ моделі~\eqref{eq:halfcentre} задається рівнянням $(\beta-1)(1-E^{2+b})/(\beta-E)=b\,(1-E^{1+b})/(1+b)$, де $E=e^{-T/(2\tau_a)}$. (а)~Знайдіть $T$ при $\beta=b=2$, $\tau_a=0.4$~с. (б)~Покажіть, що $T$ не залежить від входу $I$, а ритм можливий лише при $b>\beta-1$.
\end{exercise}

\begin{exercise}[\lvlB]
\label{ex:13:6}
\emph{Моделювання генератора.} Імпортуйте функцію \texttt{halfcentre} з \texttt{models/\allowbreak muscle\_models.py} (сам файл не запускайте) і виміряйте період, відкинувши перші два цикли, при $b=0.8$, $1.05$, $1.2$, $1.5$, $2$, $4$ і при $I=0.5$, $1$, $2$. Чи може команда «іди швидше» змінювати темп через $I$?
\end{exercise}

\begin{exercise}[\lvlB]
\label{ex:13:7}
\emph{Жорсткість проти рефлексу.} Суглоб охоплено петлею розтягу: $\dd\theta/\dd t=-a\theta(t)-G\theta(t-d)$, $a=K/B$, $d=30\ms$. Зведіть її до задачі~\ref{ex:13:4} і знайдіть найменше $a$, за якого відхилення загасає без коливань зі сталою часу $15\ms$, і потрібне $G$. Як змінювати $G$ зі збільшенням спільного напруження м'язів?
\end{exercise}

\begin{exercise}[\lvlC]
\label{ex:13:8}
\emph{Ручка темпу.} За задачами~\ref{ex:13:5} і~\ref{ex:13:6} вхід $I$ генератора~\eqref{eq:halfcentre} змінює лише розмах, але не темп. Запропонуйте мінімальну зміну моделі з деталей книги, за якої нижче від деякого $I_{\min}$ ритму немає, а вище період спадає зі зростанням $I$, хоча б удвічі при потроєнні $I$. Знайдіть $I_{\min}$ при $\tau\ll\tau_a$ і перевірте моделюванням.
\end{exercise}

\chapter{Біль: сигнал тривоги}
\chaptermark{Біль}
\label{sec:pain}

\section{Тривога, а не вимірювання}

Усі датчики розділу~\ref{sec:sensors} вимірювали світ: скільки світла, якої висоти звук, який тиск. Біль улаштований інакше. Він повідомляє не «що там», а «небезпечно, кидай усе». Для інженера це не вимірювальний канал, а \emph{сигнал тривоги}: лінія з найвищим пріоритетом, яка має пробитися крізь будь-яку поточну роботу машини.

Від вимірювального каналу сигналізація відрізняється трьома властивостями, і всі три біль має. По-перше, його важко приглушити звиканням: датчики болю звикають набагато слабше за датчики дотику, а після легкого ушкодження стають навіть чутливішими~\cite{LaMotte1982hyper}; послаблення відповіді після сильного нагрівання виміряно лише в одного їхнього типу~\cite{Treede1995heat}. Датчик світла на сталому тлі замовкає (рис.~\ref{fig:adaptation}), а сигналізація, яка замовкла через хвилину, марна. По-друге, вона має високий поріг: датчики болю спрацьовують лише тоді, коли дія близька до небезпечної, наприклад на нагрівання --- у різних датчиків пороги від $41^\circ$ до $46$--$48^\circ$, залежно від типу~\cite{Treede1995heat, Weidner1999cfib}. По-третє --- і це головне для цього розділу --- гучність тривоги вирішує не лише датчик, а й сама машина. Та сама рана болить по-різному в різних обставинах. Подивімося, з яких уже знайомих деталей це зібрано.

Датчики болю --- \nt{GLUT}-нейрони, як і решта чутливих нейронів розділу~\ref{sec:sensors}. Частина з них разом із глутаматом виділяє речовину~P з додатка~\ref{sec:othermediators}, яка повільно підсилює передачу.

\section{Два проводи: швидкий і повільний біль}

Ударте палець --- і ви відчуєте біль двічі: спершу короткий і різкий, а через секунду тупий і довгий. Це два різні проводи. Швидкий провід ізольований і проводить імпульси зі швидкістю $v$ від $5$ до $30$ метрів за секунду; повільний тонкий і голий, $0.5$--$2$ метри за секунду. Сигнал від стопи проходить близько метра, і час приходу
\begin{empheq}[box=\keybox]{equation}
t \;=\; \frac{L}{v}
\label{eq:paintime}
\end{empheq}
для швидкого проводу --- десятки мілісекунд, для повільного --- близько секунди. Два проводи несуть два повідомлення. Швидкий каже \emph{де} і \emph{коли}: його імпульси приходять раніше й точніше, як у коді часом розділу~\ref{sec:code}. Повільний каже, \emph{наскільки погано}, і його довгий тупий біль тримає машину в режимі тривоги, доки ушкодження не минуло.

\section{Ворота в спинному мозку}

Перше, куди надходить сигнал болю, --- спинний мозок, і там він проходить через \emph{ворота}~\cite{MelzackWall1965}; конкретні нейрони цих воріт знайдено порівняно недавно~\cite{Duan2014}. На вихідний \nt{GLUT}-нейрон, що передає біль далі в мозок, сходяться провід болю й провід дотику. А поруч сидить гальмівний нейрон-посередник, \nt{GABA} або \nt{GLYC}, якого дотик збуджує (рис.~\ref{fig:gate}). Дотик зачиняє ворота. У вихідній теорії воріт~\cite{MelzackWall1965} біль, навпаки, вимикає цього посередника, і сильний біль ворота відчиняє; це припущення теорії, на знайдених нейронах воріт його прямо не показано.

\begin{figure}[t]
\centering
\begin{tikzpicture}[>=Stealth,
  nrn/.style={circle,draw,very thick,minimum size=1.0cm,inner sep=0pt,font=\scriptsize},
  inh/.style={-{Bar[width=7pt]},thick,red!70!black},
  bc/.style={->,thick,densely dashed,orange!85!black}]
\node[nrn,fill=blue!6] (Z) at (6,0) {\nt{GLUT}};
\node[nrn,fill=red!10] (Q) at (3.6,-1.6) {\nt{GABA}};
\draw[->,very thick,red!70!black] (0,0.6) node[left,font=\small,black,align=right] {біль $\nu$} -- (5.45,0.15);
\draw[->,very thick,blue!60!black] (0,-2.6) node[left,font=\small,black,align=right] {дотик $\mu$} -- (2.9,-2.0);
\draw[->,thick,blue!60!black] (1.8,-2.35) -- (5.5,-0.35) node[pos=0.8,below right=-2pt,font=\scriptsize] {$w_{z\mu}$};
\draw[inh] (1.6,0.5) -- (3.35,-1.1) node[pos=0.5,left=1pt,font=\scriptsize] {$w_{q\nu}$};
\draw[inh] (Q) -- (Z) node[pos=0.5,above left=-1pt,font=\scriptsize] {$\beta$};
\draw[->,very thick] (Z) -- (8.2,0) node[right,font=\small,align=left] {у мозок: $z$};
\draw[bc] (6,2.1) node[above,font=\scriptsize,black,align=center] {згори: \nt{SERO}, \nt{NORA}, опіоїди} -- (Z.north) node[pos=0.45,right,font=\scriptsize,black] {підсилення $g$};
\end{tikzpicture}
\caption{Ворота болю в спинному мозку за моделлю~\eqref{eq:gate}. Вихідний \nt{GLUT}-нейрон отримує біль $\nu$ і слабке збудження від дотику $\mu$. Гальмівний посередник (\nt{GABA} або \nt{GLYC}) збуджується дотиком і, за припущенням теорії воріт, вимикається болем. Згори, широкомовними каналами, надходить підсилення $g$.}
\label{fig:gate}
\end{figure}

Запишімо це для частот, як у розділі~\ref{sec:gaba}. Частота посередника $q$ і вихідна частота $z$ дорівнюють
\begin{empheq}[box=\keybox]{equation}
q \;=\; \bigl[w_{q\mu}\,\mu + q_0 - w_{q\nu}\,\nu\bigr]_+ ,
\qquad
z \;=\; \bigl[\,g\,(\nu + w_{z\mu}\,\mu) - \beta\,q\,\bigr]_+ ,
\label{eq:gate}
\end{empheq}
де $\nu$ --- вхід болю, $\mu$ --- вхід дотику, $w$ --- ваги зв'язків (перший індекс --- кому, другий --- від кого), $\beta$ --- сила гальмування, $q_0$ --- власна фонова активність посередника, $g$ --- підсилення, яке задають згори (про нього нижче). Це заборона~\eqref{eq:veto} з розділу~\ref{sec:gaba}: «біль, але не за дотику», --- з однією поправкою, взятою з теорії воріт як припущення: сильний біль сам вимикає забороняльний нейрон.

Ми порахували модель при $w_{q\mu}=1$, $q_0=0.2$, $w_{q\nu}=0.5$, $w_{z\mu}=0.3$, $\beta=1.5$ (рис.~\ref{fig:gatecurves}). Без дотику біль починає проходити з $\nu\approx0.18$. З дотиком поріг піднімається до $0.86$, і при $\nu=1$ вихід падає вчетверо, з $1$ до $0.25$. А за сильного болю, $\nu=2$, дотик уже нічого не змінює: ворота розчахнуті. Так і в житті: потерти забите місце допомагає, а за серйозної травми --- ні. Сам дотик, без болю, через ворота не проходить: тривога від дотику не здіймається.

\begin{figure}[t]
\centering
\begin{tikzpicture}[>=Stealth,x=5.2cm,y=1.35cm]
\draw[->,gray!70!black] (0,0) -- (2.12,0) node[right,font=\small,black] {біль $\nu$};
\draw[->,gray!70!black] (0,0) -- (0,4.3) node[left,font=\small,black] {$z$};
\foreach \t in {0,0.5,1,1.5,2}{\draw[gray!60!black] (\t,-0.05) -- (\t,0.05); \node[below,font=\scriptsize] at (\t,0) {\t};}
\foreach \f in {1,2,3,4}{\draw[gray!60!black] (-0.01,\f) -- (0.01,\f); \node[left,font=\scriptsize] at (0,\f) {\f};}
\draw[very thick,red!70!black] plot file {figures/pain_gate_notouch.dat};
\draw[very thick,blue!60!black] plot file {figures/pain_gate_touch.dat};
\draw[very thick,orange!85!black,densely dashed] plot file {figures/pain_gate_sens.dat};
\draw[very thick,red!70!black] (0.08,3.9) -- (0.2,3.9) node[right,font=\scriptsize,black] {без дотику};
\draw[very thick,blue!60!black] (0.08,3.45) -- (0.2,3.45) node[right,font=\scriptsize,black] {з дотиком};
\draw[very thick,orange!85!black,densely dashed] (0.08,3.0) -- (0.2,3.0) node[right,font=\scriptsize,black] {після сенсибілізації, $g=2$};
\end{tikzpicture}
\caption{Вихід воріт~\eqref{eq:gate} залежно від сили болю. Дотик піднімає поріг і приглушує слабкий і середній біль, але не сильний. Сенсибілізація (штрихова лінія) вдвічі збільшує підсилення: той самий біль передається вдвічі гучніше, а поріг опускається.}
\label{fig:gatecurves}
\end{figure}

\section{Ручка гучності згори}

Ворота керуються не лише знизу. Зі стовбура мозку до них ідуть низхідні шляхи, які змінюють підсилення $g$ у~\eqref{eq:gate}: \nt{SERO}, \nt{NORA} та опіоїдні пептиди з додатка~\ref{sec:othermediators}~\cite{Fields2004}. Це ті самі широкомовні канали, що й у розділах~\ref{sec:serocomputer}--\ref{sec:acet}, тільки тут їхня ручка --- гучність тривоги. Вони вміють крутити її в обидва боки: ті самі низхідні кола можуть і приглушувати біль, і посилювати його.

Навіщо машині зменшувати тривогу? Тому що тривога --- це переривання, а переривання не завжди доречне. Тварина, що тікає від хижака, не повинна зупинятися через поранену лапу: спершу тікати, потім зализувати рану. Очікування полегшення теж приглушує біль, і частина цього ефекту зникає, якщо заблокувати опіоїдні рецептори~\cite{Levine1978}: очікування, обчислене в мозку, спускається опіоїдним каналом до воріт. Саму цю картину виміряно; як саме мозок вирішує, якою має бути гучність, --- гіпотеза.

\section{Сенсибілізація: тривога, що навчається}

Після ушкодження вихідні нейрони воріт стають чутливішими: той самий вхід викликає більший вихід, а поріг опускається~\cite{Woolf1983}. У моделі~\eqref{eq:gate} це зростання підсилення $g$, і найпростіший його опис --- повільне RC-коло, яке накачує сам біль:
\begin{empheq}[box=\keybox]{equation}
\tau_s\,\frac{\dd g}{\dd t} \;=\; -(g-1) + k_s\,z ,
\label{eq:sensitization}
\end{empheq}
де $\tau_s$ --- час, за який чутливість повертається до норми; він більший за тривалість самого болю, бо сенсибілізація тримається й після того, як ушкоджувальний стимул припинився~\cite{Woolf1983}, а $k_s$ --- сила накачування. Доки тканину ушкоджено, вихід воріт тримає $g$ вище від одиниці, і все, що відбувається поруч із раною, передається гучніше (штрихова лінія на рис.~\ref{fig:gatecurves}: при $g=2$ поріг опускається до $0.11$, а вихід подвоюється). Це розумне налаштування сигналізації: доки ушкодження не загоїлося, краще перестрахуватися.

Для інженера це додатний зворотний зв'язок із повільним витоком: біль підвищує підсилення, підсилення --- біль. Доки петля слабка, $g$ повертається до норми, коли рана загоїлася. Якщо ж петля сильна, сигналізація може застрягнути в увімкненому стані й після того, як ушкодження вже немає, --- як група із сильним зворотним зв'язком на рис.~\ref{fig:bistable}. Модель~\eqref{eq:sensitization} --- спрощення; те, що центральна сенсибілізація існує, виміряно.

\section{Біль як учитель і як переривання}

Біль у машині має дві роботи, крім сигналізації.

\emph{Учитель.} Біль --- від'ємна винагорода: у формулу помилки~\eqref{eq:ctd} він входить як $r<0$. Усе, що говорилося в розділі~\ref{sec:dofa}, працює й тут: провісник болю починає викликати помилку заздалегідь, і мережа вчиться уникати того, що до болю веде. Досліди на людях показують, що активність мозку під час навчання на болю йде саме за помилкою часових різниць~\cite{Seymour2004}, а частина \nt{DOPA}-нейронів відповідає й на неприємні події (розділ~\ref{sec:dofaalt}).

\emph{Переривання.} Біль відриває увагу від поточної задачі --- це його призначення, а не побічний ефект~\cite{EcclestonCrombez1999}. У розділі~\ref{sec:nora} ту саму роль «переривання» обговорювали для уривчастої пачки \nt{NORA}. А найшвидша відповідь навіть не доходить до мозку: відсмикування руки від гарячого замикається просто в спинному мозку тією самою парою м'язів і тією самою забороною антагоніста, що на рис.~\ref{fig:antagonists}.

\begin{keyblock}
\begin{proposition}[Сигнал тривоги]
\label{prop:pain}
Біль --- не вимірювання, а сигнал тривоги з високим пріоритетом. Він звикає набагато слабше за вимірювальні канали, іде двома проводами (швидким --- «де», повільним --- «наскільки погано»), проходить через ворота, які зачиняє дотик (а сильний біль, за припущенням теорії воріт, відчиняє), і його гучність задають згори широкомовні канали. Ці пристрої виміряно, крім відчинення воріт болем; прості моделі~\eqref{eq:gate} і~\eqref{eq:sensitization} --- спрощення.
\end{proposition}
\end{keyblock}

\section{Підсумок}

\begin{table}[t]
\centering
\footnotesize
\setlength{\tabcolsep}{4pt}
\renewcommand{\arraystretch}{1.15}
\begin{tabular}{>{\raggedright}p{3.0cm}>{\raggedright}p{4.9cm}>{\raggedright\arraybackslash}p{4.9cm}}
\toprule
Що робить біль & Як & Що відомо \\
\midrule
здіймає тривогу & високий поріг, звикання слабше, ніж у дотику & поріг виміряно; порівняння звикання --- оцінка \\
повідомляє «де» і «наскільки погано» & два проводи з різною швидкістю~\eqref{eq:paintime} & виміряно \\
проходить через ворота & заборона через \nt{GABA}/\nt{GLYC}~\eqref{eq:gate} & ворота виміряно; відчинення болем --- припущення теорії; модель --- спрощення \\
змінює гучність & низхідні \nt{SERO}, \nt{NORA}, опіоїди & виміряно; правило вибору гучності --- гіпотеза \\
навчається після ушкодження & зростання підсилення~\eqref{eq:sensitization} & сенсибілізацію виміряно; модель --- спрощення \\
навчає уникати & від'ємна винагорода в~\eqref{eq:ctd} & схожість із помилкою часових різниць виміряно \\
перериває роботу & захоплення уваги; рефлекс відсмикування & виміряно \\
\bottomrule
\end{tabular}
\caption{Біль як сигнал тривоги.}
\label{tab:pain}
\end{table}

Біль зібрано з деталей, які ми вже бачили (таблиця~\ref{tab:pain}): \nt{GLUT}-датчики, два проводи, заборона через гальмівного посередника, широкомовна ручка гучності, повільне RC-коло, помилка передбачення й переривання. Нове в ньому одне: це єдиний вхід машини, гучність якого машина задає сама, --- сигналізація, яку власник може і приглушити, і переналаштувати.

\section{Задачі}

\begin{exercise}[\lvlA]
\label{ex:14:1}
\emph{Два проводи.} Сигнал іде від стопи, $L=1$~м. (а)~Залп іде пучком проводів одного типу, швидкості яких заповнюють діапазон~\eqref{eq:paintime}. На скільки він розтягується до приходу в спинний мозок для швидкого й повільного типу? (б)~Хвилі болю проводами $15$ і $1$~м/с розрізненні за різниці від $0.1$~с. З якої відстані вони зливаються?
\end{exercise}

\begin{exercise}[\lvlB]
\label{ex:14:2}
\emph{Коли дотик болить.} Ворота~\eqref{eq:gate} з параметрами тексту, підсилення $g$ зростає під час сенсибілізації. До якого $g$ дотик без болю не проходить через ворота за жодної сили $\mu$? За якого $g$ починає проходити дотик $\mu=1$?
\end{exercise}

\begin{exercise}[\lvlB]
\label{ex:14:3}
\emph{Стійкість сенсибілізації.} Входи сталі, вихід воріт лінійний за $g$: $z=gA-C$, $A=\nu+w_{z\mu}\mu$, $C=\beta q\ge0$. (а)~Знайдіть рівновагу $g^*$ моделі~\eqref{eq:sensitization} і умову її стійкості. (б)~При $k_s=0.5$ знайдіть силу болю, вище від якої модель іде врознос, без дотику і з дотиком $\mu=1$.
\end{exercise}

\begin{exercise}[\lvlA]
\label{ex:14:4}
\emph{Провісник болю.} Сигнал у момент $0$ передбачає біль у момент $T=2$~с, у~\eqref{eq:ctd} $r(t)=-P\,\delta(t-T)$, $\tau_\gamma=2$~с. (а)~Знайдіть площу сплеску $\delta$ у момент сигналу. (б)~Дія в момент $t_e=1$~с скасовує біль. Який сплеск $\delta$ у цей момент і чого він навчить мережу?
\end{exercise}

\begin{exercise}[\lvlB]
\label{ex:14:5}
\emph{Сигналізація, що защіпається.} Доповніть ворота з \texttt{models/\allowbreak pain\_gate.py} динамікою~\eqref{eq:sensitization} і насиченням $z\le4$. Дотик $\mu=1$ є завжди, ушкодження $\nu=2$ триває час $T$, далі $\nu=0$. Моделюванням знайдіть найменше $T$ (в одиницях $\tau_s$), після якого сигналізація лишається ввімкненою, при $k_s=4$, $4.6$, $5$, $6$, $8$.
\end{exercise}

\begin{exercise}[\lvlB]
\label{ex:14:6}
\emph{Проєктування воріт.} При $\beta=1.5$, $w_{z\mu}=0.3$, $g=1$ доберіть у~\eqref{eq:gate} $q_0$, $w_{q\nu}$, $w_{q\mu}$ так, щоб поріг болю був $0.2$ без дотику і $1.0$ з дотиком $\mu=1$, а дотик переставав приглушувати біль при $\nu_\times=2.5$. До якого підсилення $g$ дотик без болю не проходить через ворота?
\end{exercise}

\begin{exercise}[\lvlC]
\label{ex:14:7}
\emph{Ручка гучності.} Робота приносить цінність $V$ за одиницю часу, робота з ушкодженням $\nu$ коштує $c\nu$, тож переривати вигідно при $\nu>V/c$; біль перериває при $z>\theta=0.5$. (а)~Яке підсилення $g^*$ воріт~\eqref{eq:gate} без дотику дає цей поріг при $V/c=0.25$, $0.5$, $2$? (б)~З яких сигналів книги машина могла б оцінювати $V$ і $c$?
\end{exercise}

\chapter{Як машина вчиться рухатися}
\chaptermark{Як машина вчиться рухатися}
\label{sec:motorlearning}
\begin{chaptergoals}
\item чому рухам потрібен третій учитель --- провід помилки до кожного виходу --- і чого він коштує;
\item як \nt{GLUT}-розширювач і провід помилки на кожен вихід дають правило найменших квадратів;
\item як навчене коло стає внутрішньою моделлю --- адаптивним фільтром, що передбачає тіло;
\item чому швидкий і повільний процеси навчання пояснюють післядію і спонтанне відновлення.
\end{chaptergoals}

\section{Три вчителі}

У розділі~\ref{sec:muscles} ми бачили, як машина рухає м'язами, і зупинилися на гіпотезі: швидкі й точні рухи потребують внутрішньої моделі тіла, яка передбачає результат команди~\cite{WolpertGhahramani2000}. Звідки береться ця модель? Її треба вивчити, і вивчити не так, як ми вчилися досі.

У книзі вже було два вчителі. Перший --- \emph{жодного}: правило Гебба (розділ~\ref{sec:dofa}) посилює те, що часто буває разом, і стискає входи. Другий --- \emph{винагорода}: \nt{DOPA} розсилає всім одне число, «краще чи гірше, ніж очікувалося». Рухові потрібен третій учитель --- \emph{помилка}. Рука промахнулася повз чашку на два сантиметри ліворуч: це не «гірше», це вектор, який каже кожному виходу, куди і наскільки він помилився. Інженер назвав би це навчанням з учителем.

Різниця між другим і третім учителем --- це різниця між одним проводом на всіх і окремим проводом на кожен вихід. У твердженні~\ref{prop:broadcastcost} навчання за одним спільним числом сповільнювалося з кожним нейроном, чий шум впливав на результат. Якщо ж кожен вихід $i$ отримує свою помилку $e_i$, ваги можна змінювати за найпростішим правилом:
\begin{empheq}[box=\keybox]{equation}
\frac{\dd w_{ij}}{\dd t} \;=\; -\eta\,e_i(t)\,x_j(t) .
\label{eq:lms}
\end{empheq}
Це правило найменших квадратів, давно відоме в техніці адаптивних фільтрів~\cite{WidrowHoff1960}: вага змінюється пропорційно добуткові свого входу на помилку свого виходу й у середньому йде за градієнтом квадрата помилки. Жодного шуму-пошуку, як у розділі~\ref{sec:dofa}, не потрібно: напрямок поправки надходить просто проводом помилки. І швидкість не падає з кількістю виходів, бо чужі помилки до синапсу не доходять.

\begin{keyblock}
\begin{proposition}[Три вчителі]
\label{prop:teachers}
Правило Гебба без учителя навчає того, що часте; сигнал \nt{DOPA} одним числом на всю мережу навчає того, що вигідне, і тим повільніше, що більше нейронів; провід помилки до кожного виходу навчає за правилом~\eqref{eq:lms} того, що точне, і швидкість не залежить від кількості виходів. Ціна третього способу --- окремий провід на кожен вихід і хтось, хто знає правильну відповідь.
\end{proposition}
\end{keyblock}

\section{Провід помилки}

Хто може знати правильну відповідь для руху? Самі датчики. Якщо рука пішла ліворуч від цілі, про це повідомляють око й датчики розтягу з розділу~\ref{sec:sensors}. Потрібна схема, яка доставить цю помилку до потрібних виходів.

У мозку є ділянка, яка, схоже, влаштована саме під це~\cite{Marr1969,Albus1971}. Її схема незвично правильна, майже кристалічна, і в ній легко впізнати правило~\eqref{eq:lms} (рис.~\ref{fig:errorwire}).

\emph{Розширювач входу.} Сигнали про команду й про стан тіла надходять на величезний шар маленьких \nt{GLUT}-нейронів. Їх близько сімдесяти мільярдів --- більше, ніж у всьому іншому мозку разом узятому~\cite{Azevedo2009}. Кожен із них змішує кілька входів, і сигнал розкладається в дуже довгий вектор. Інженерові цей прийом знайомий: якщо підняти розмірність входу, у новому просторі майже будь-яку потрібну залежність можна отримати простою зваженою сумою~\cite{Cover1965}.

\emph{Вихідні нейрони.} Зважену суму обчислюють близько тридцяти мільйонів великих вихідних нейронів~\cite{Andersen1992cereb}, у кожного --- порядку ста тисяч входів від розширювача. І це \nt{GABA}-нейрони: результат навчання передається далі не збудженням, а гальмуванням, як у забороні розділу~\ref{sec:gaba}.

\emph{Провід помилки.} Кожен вихідний нейрон отримує ще один вхід, особливий: рівно один \nt{GLUT}-провід, що спрацьовує рідко, близько разу на секунду, але щоразу викликає у вихідному нейроні потужний спалах~\cite{Ito2001}. Це і є $e_i$: імовірність спалаху залежить від напрямку помилки руху~\cite{Herzfeld2018}. Збіг спалаху помилки з активністю входу послаблює цей вхід --- точнісінько член $-\eta\,e_i x_j$ у~\eqref{eq:lms}.

\begin{figure}[t]
\centering
\begin{tikzpicture}[>=Stealth,
  gran/.style={circle,draw,thick,fill=blue!6,minimum size=0.32cm,inner sep=0pt},
  outn/.style={circle,draw,very thick,fill=red!10,minimum size=1.1cm,inner sep=0pt,font=\scriptsize},
  inh/.style={-{Bar[width=7pt]},very thick,red!70!black}]
% inputs
\foreach \y/\lab in {1.6/{команда},0.4/{стан тіла}}{
  \draw[->,thick,blue!60!black] (-0.6,\y) node[left,font=\scriptsize,black] {\lab} -- (0.35,\y);}
% expansion layer
\foreach \i in {0,...,11}{\node[gran] (g\i) at ({0.8+0.28*mod(\i,2)},{-0.3+0.23*\i}) {};}
\foreach \i in {0,...,11}{\pgfmathsetmacro{\yy}{mod(\i,3)>0 ? 1.6 : 0.4}\draw[gray!60!black] (0.35,\yy) -- (g\i);}
\node[font=\scriptsize,align=center] at (0.95,-1.05) {розширювач\\$\sim7\cdot10^{10}$ \nt{GLUT}};
% parallel wires to the output neuron
\foreach \i in {0,...,11}{\draw[->,gray!70!black] (g\i.east) -- (4.1,{1.0+0.07*(\i-5.5)});}
\node[outn] (P) at (4.6,1.0) {\nt{GABA}};
\node[font=\scriptsize,align=center,above=2pt] at (P.north) {вихідний нейрон\\$\sim10^5$ входів $x_j$, ваги $w_{ij}$};
\draw[inh] (P) -- (6.8,1.0) node[right,font=\scriptsize,align=left] {поправка\\до руху};
% error wire
\draw[->,very thick,red!70!black,densely dashed] (4.6,-1.4) node[below,font=\scriptsize,black,align=center] {один провід помилки $e_i$\\\nt{GLUT}, $\sim1$ раз/с} -- (P.south);
\end{tikzpicture}
\caption{Схема навчання з учителем, як її, схоже, реалізує мозок. Команда й стан тіла розкладаються величезним шаром \nt{GLUT}-нейронів у довгий вектор $x_j$; вихідний \nt{GABA}-нейрон додає його з вагами $w_{ij}$. Єдиний провід помилки приносить $e_i$; збіг помилки з активним входом послаблює вагу цього входу, як у правилі~\eqref{eq:lms}.}
\label{fig:errorwire}
\end{figure}

Одна трудність знайома нам із розділу~\ref{sec:dofa}: помилка надходить пізніше, ніж вхід, що до неї призвів. Промах видно через десятки--сотні мілісекунд після команди. Отже, синапс має пам'ятати, що був активний, --- потрібен слід, як у правилі~\eqref{eq:threefactor}. І довжина цього сліду, судячи з дослідів, підлаштована під затримку зворотного зв'язку в кожній задачі: там, де помилка надходить через сто мілісекунд, найсильніше послаблюється вхід, що був активним за сто мілісекунд до неї~\cite{Suvrathan2016}.

\begin{keyblock}
\begin{proposition}[Провід помилки]
\label{prop:errorwire}
У схемі, зображеній на рис.~\ref{fig:errorwire}, кожен вихідний нейрон отримує рівно один провід помилки, і збіг помилки з входом послаблює цей вхід. Це правило найменших квадратів~\eqref{eq:lms}, застосоване до входу, розширеного до величезної розмірності. Будову схеми й послаблення за збігу виміряно; що вся схема працює як навчання з учителем, --- провідна гіпотеза.
\end{proposition}
\end{keyblock}

\section{Внутрішня модель як адаптивний фільтр}

Чого вчиться така схема? Найприродніша відповідь --- передбачення. Нехай на вхід надходить команда $u(t)$, пропущена через набір фільтрів із різними сталими часу, як у розширювачі: $x_j(t)=(g_j*u)(t)$. Вихід --- їхня зважена сума, передбачення того, що повідомлять датчики:
\begin{empheq}[box=\keybox]{equation}
\hat y(t) \;=\; \sum_j w_j\,(g_j*u)(t), \qquad e(t) \;=\; \hat y(t) - y(t) ,
\label{eq:adaptivefilter}
\end{empheq}
а ваги вчаться за~\eqref{eq:lms}. Це \emph{адаптивний фільтр}: схема, яка сама добирає свою передавальну функцію так, щоб повторити невідому систему~\cite{Fujita1982}. Тут невідома система --- власне тіло з його масою, пружністю й затримками. Вивчений фільтр і є внутрішньою моделлю розділу~\ref{sec:muscles}: він каже, що повідомлять датчики, не чекаючи на них.

Така модель корисна двічі. По-перше, її передбачення можна подати замість запізнілого зворотного зв'язку, і тоді умова стійкості~\eqref{eq:delaystable} перестає зв'язувати руки: виправляти можна швидко, за моделлю. По-друге, різниця між передбаченим і отриманим --- це сигнал про \emph{несподіване}. Усе, що машина зробила сама, модель передбачила й відняла; лишається те, що зробив світ. Тому, за однією з гіпотез, ми й не можемо полоскотати самі себе~\cite{Blakemore1998}.

\section{Коли світ змінюється: швидке й повільне}

Перевірмо схему дослідом, який легко поставити. Людина тягнеться до цілі, а світ несподівано змінюється: наприклад, на руку діє бічна сила. Перші рухи промахуються, потім промах зменшується. Якщо силу прибрати, рука спершу промахується в \emph{інший} бік --- модель вивчила силу й далі її компенсує. Це пряме свідчення того, що вивчено модель, а не просто «почало виходити».

Досліди показують і тоншу річ: навчаються одночасно два процеси --- швидкий, що швидко вчиться і швидко забуває, і повільний, що вчиться повільно й пам'ятає довго~\cite{Smith2006}. Поправки оновлюються після кожного руху, за подією:
\begin{empheq}[box=\keybox]{equation}
e = p - (x_f + x_s), \qquad x_f \leftarrow A_f\,x_f + B_f\,e, \qquad x_s \leftarrow A_s\,x_s + B_s\,e ,
\label{eq:tworate}
\end{empheq}
де $p$ --- збурення, $x_f$ і $x_s$ --- вивчені поправки двох процесів, $A$ --- наскільки поправка зберігається до наступного руху, $B$ --- наскільки сильно вона вчиться на помилці. У швидкого процесу $B$ велике, а $A$ помітно менше за одиницю; у повільного --- навпаки.

\begin{figure}[t]
\centering
\begin{tikzpicture}[>=Stealth,x=0.034cm,y=1.9cm]
\draw[->,gray!70!black] (0,0) -- (355,0) node[right,font=\small,black] {рух};
\draw[->,gray!70!black] (0,-1.15) -- (0,1.25);
\foreach \v in {-1,1}{\draw[gray!60!black] (-3,\v) -- (3,\v); \node[left,font=\scriptsize] at (0,\v) {$\v$};}
\foreach \n in {100,200,300}{\draw[gray!60!black] (\n,-0.04) -- (\n,0.04); \node[below,font=\scriptsize] at (\n,0) {\n};}
\draw[thin,gray!70!black] plot file {figures/tworate_p.dat};
\draw[thick,orange!85!black] plot file {figures/tworate_fast.dat};
\draw[thick,blue!60!black] plot file {figures/tworate_slow.dat};
\draw[very thick,black] plot file {figures/tworate_net.dat};
\node[font=\scriptsize,gray!50!black] at (120,1.12) {збурення $p$};
\node[font=\scriptsize,black,anchor=west] at (238,0.52) {сума: повернення};
\node[font=\scriptsize,blue!60!black,anchor=west] at (150,0.39) {повільний $x_s$};
\node[font=\scriptsize,orange!85!black,anchor=west] at (60,0.08) {швидкий $x_f$};
\draw[densely dashed,gray!60!black] (226,-1.15) -- (226,1.25);
\node[font=\scriptsize,align=center,anchor=south west] at (228,-1.1) {помилок\\більше немає};
\end{tikzpicture}
\caption{Два процеси навчання за моделлю~\eqref{eq:tworate}, $A_f=0.92$, $B_f=0.1$, $A_s=0.996$, $B_s=0.02$. Двісті рухів зі збуренням $p=1$, далі шість зі зворотним, $p=-1$, доки сума не повернеться до нуля. Після штрихової лінії помилки перестають надходити, і сума сама повертається до старої поправки: швидкий процес забув зворотне збурення, повільний пам'ятає пряме.}
\label{fig:tworate}
\end{figure}

Модель~\eqref{eq:tworate} пояснює дослід, який важко зрозуміти без неї (рис.~\ref{fig:tworate}). У нашому розрахунку за двісті рухів зі збуренням сума поправок сягає $0.84$, причому більшу частину тримає повільний процес, $0.63$. Далі шість рухів зі зворотним збуренням повертають суму до нуля: швидкий процес пішов у мінус, $-0.53$, повільний ще в плюсі, $0.46$. Якщо тепер перестати повідомляти помилки, швидкий процес забуває свою поправку за десятки рухів, повільний --- набагато довше, і сума \emph{сама} піднімається до $0.37$ за сорок рухів: стара навичка повертається без жодного тренування. Таке спонтанне відновлення виміряно в людей; модель з одним процесом його дати не може --- без помилок вона просто згасає до нуля.

Навіщо два процеси? Правдоподібну інженерну відповідь дає оптимальний фільтр розділу~\ref{sec:acet}. Тіло змінюється з різних причин із різною швидкістю: м'язи втомлюються за хвилини, ростуть і слабшають за місяці. Якщо завади бувають швидкими й повільними, найкраща оцінка складається з двох фільтрів із різними сталими часу, і кожен ловить свою~\cite{Kording2007}. Швидкий процес --- це тимчасова поправка, «рука втомилася»; повільний --- «рука стала іншою». Це поки гіпотеза, але вона пояснює, чому навичка, вивчена за один день, частково втрачається до наступного і чому вдруге вона вчиться швидше.

\section{Підсумок}

\begin{table}[t]
\centering
\footnotesize
\setlength{\tabcolsep}{4pt}
\renewcommand{\arraystretch}{1.15}
\begin{tabular}{>{\raggedright}p{3.0cm}>{\raggedright}p{4.6cm}>{\raggedright\arraybackslash}p{5.2cm}}
\toprule
Що робить схема & Як & Що відомо \\
\midrule
вчиться з учителем & провід помилки до кожного виходу, правило~\eqref{eq:lms} & будову схеми й послаблення за збігу виміряно; навчання з учителем --- провідна гіпотеза \\
розширює вхід & сімдесят мільярдів \nt{GLUT}-нейронів & кількість нейронів виміряно; роль розширення --- теорія \\
враховує затримку помилки & слід, підлаштований під затримку & виміряно \\
будує внутрішню модель & адаптивний фільтр~\eqref{eq:adaptivefilter} & добре обґрунтована гіпотеза \\
вчиться двома темпами & швидкий і повільний процеси~\eqref{eq:tworate} & післядію й спонтанне повернення виміряно; два процеси --- модель \\
\bottomrule
\end{tabular}
\caption{Як машина вчиться рухатися.}
\label{tab:motorlearning}
\end{table}

Тепер у машини три вчителі (таблиця~\ref{tab:motorlearning}). Гебб стискає те, що часте; \nt{DOPA} одним числом навчає вибирати вигідне; провід помилки навчає точності, і навчає швидко, бо до кожного виходу йде своя помилка. Мозок, схоже, користується всіма трьома, і кожним там, де він найдешевший: широкомовним сигналом --- для небагатьох рішень, окремими проводами --- для точного підлаштування тіла, де правильну відповідь повідомляють самі датчики.

\section{Задачі}

\begin{exercise}[\lvlA]
\label{ex:15:1}
\emph{Ємність схеми з проводом помилки.} У схемі $3\cdot10^7$ вихідних нейронів по $10^5$ входів, у кожного свій провід помилки ($1$~імп/с). Нейрон з $n$ входами вивчає будь-яку розмітку $\approx2n$ векторів~\cite{Cover1965}. (а)~Скільки асоціацій зберігає схема? (б)~Оцініть за~\eqref{eq:spikeentropy} при $\Delta t=10\ms$ найменший час заповнення ємності нейрона.
\end{exercise}

\begin{exercise}[\lvlB]
\label{ex:15:2}
\emph{Швидкість правила найменших квадратів.} Моди правила~\eqref{eq:lms} згасають зі швидкостями $\eta\lambda_k(C)$. Для п'яти фільтрів~\eqref{eq:adaptivefilter} на білому шумі $C_{jk}=2\sqrt{\tau_j\tau_k}/(\tau_j+\tau_k)$. Знайдіть відношення найшвидшої й найповільнішої швидкостей, якщо $\tau_j$ зростають у $2$ рази ($10$--$160\ms$) і в $4$ рази.
\end{exercise}

\begin{exercise}[\lvlB]
\label{ex:15:3}
\emph{Аналіз двох процесів.} Запишіть модель~\eqref{eq:tworate} з параметрами рис.~\ref{fig:tworate} як $\mathbf x_{n+1}=M\mathbf x_n+\mathbf b\,p$. (а)~Знайдіть за власними числами $M$ сталі часу в рухах. (б)~Знайдіть усталені $x_f$, $x_s$ та їхню суму за сталого $p=1$.
\end{exercise}

\begin{exercise}[\lvlB]
\label{ex:15:4}
\emph{Моделювання: вдруге швидше.} За моделлю~\eqref{eq:tworate} з параметрами з \texttt{models/\allowbreak motor\_learning.py} знайдіть кількість рухів до $x_f+x_s\ge0.8$ при $p=1$: (а)~у ненавченої машини; (б)~після $200$ рухів із $p=1$ і зворотного збурення до нуля суми, як на рис.~\ref{fig:tworate}. Чи може так прискоритися модель з одним процесом?
\end{exercise}

\begin{exercise}[\lvlB]
\label{ex:15:5}
\emph{Регулятор із внутрішньою моделлю.} Об'єкт $\dd x/\dd t=ku$ видно із затримкою $d$; регулятор $u=-G\hat x$ передбачає $\hat x(t)=x(t-d)+\hat k\int_{t-d}^{t}u\,\dd s$. (а)~При $\hat k=k=1$, $d=150\ms$ знайдіть $G$ для сталої часу $20\ms$ і порівняйте з межею~\eqref{eq:delaystable}. (б)~Чи стійка схема при $\hat k=0.4k$?
\end{exercise}

\begin{exercise}[\lvlB]
\label{ex:15:6}
\emph{Адаптивний фільтр навчає м'яз.} Об'єкт --- посмикування~\eqref{eq:twitch} одиничної площі, $T_c=50\ms$; фільтри~\eqref{eq:adaptivefilter} --- RC-кола з $\tau_j=12.5$, $25$, $50$, $100$, $200\ms$; вхід --- нове нормальне число кожні $10\ms$. За скільки секунд правило~\eqref{eq:lms} при $\eta=50$ і $200$ знижує помилку до $0.5\%$? Чому ваги й через $300$~с далекі від найкращих?
\end{exercise}

\begin{exercise}[\lvlC]
\label{ex:15:7}
\emph{Два процеси як оптимальний фільтр.} Збурення --- сума процесів $p^{f,s}_{n+1}=A_{f,s}\,p^{f,s}_n+w^{f,s}_n$ із дисперсіями шумів $q_f$, $q_s$, помилку спостерігають із шумом дисперсії $R$; фільтр Калмана дає~\eqref{eq:tworate}. За яких $q_f/R$, $q_s/R$ з $A_f=0.92$, $A_s=0.996$ виходить $B_f=0.1$, $B_s=0.02$? Як зміняться $B_f$, $B_s$, якщо вчетверо збільшити $q_s$?
\end{exercise}

\chapter{Другий вихід: як мозок керує хімією тіла}
\chaptermark{Хімічний вихід}
\label{sec:chemoutput}
\begin{chaptergoals}
\item як газ \nt{NORA} і гальмо \nt{ACET} підлаштовують ритмоводій і чому гальмо швидше;
\item як кров працює широкомовною шиною, а \nt{ADRE} переводить тіло в режим бігу;
\item чому каскад гормонів зі зворотним зв'язком не може бути водночас точним і стійким: $K<8$;
\item як утомлений рецептор читає частоту порцій і як добовий ритм підстроюється за світлом.
\end{chaptergoals}

\section{Швидкий і повільний виходи}

У розділі~\ref{sec:muscles} ми бачили швидкий вихід машини --- м'язи. Але є в неї й другий, повільний. Мозок тримає в потрібних межах температуру тіла, частоту серця, кількість води та цукру, готує тіло до бігу чи до сну. Для цього він не рухає суглобами, а випускає хімічні сигнали. Шляхів два. Перший --- нерви, що йдуть просто до внутрішніх органів: до серця, судин, кишківника, залоз. Другий --- кров: деякі нейрони й залози викидають речовину не в щілину до сусіда, а просто в кровотік.

Кров --- найширокомовніша шина з усіх, що траплялися в книзі. Широкомовні канали розділів~\ref{sec:serocomputer}--\ref{sec:acet} розсилали сигнал на великі ділянки мозку; кров розсилає його кожній клітині тіла. Кров обходить тіло за час порядку хвилини, тож повідомлення доходить до всіх за десятки секунд і тримається хвилини й години, доки речовина не зруйнується. Адреси в такої розсилки немає зовсім. Як і в твердженні~\ref{prop:receptor}, зміст повідомлення обирає отримувач: відгукуються лише клітини, що мають рецептор до цієї речовини.

\section{Газ і гальмо: два проводи до органів}

Найкращий приклад --- серце. У серці є власний ритмоводій, як у \nt{SERO}-нейрона розділу~\ref{sec:serocomputer}: він сам задає ритм, і якщо відключити всі нерви, серце б'ється приблизно сто разів на хвилину. Мозок ритму не породжує, а лише підлаштовує його, до того ж двома проводами протилежного знака (рис.~\ref{fig:gasbrake}). Одним іде \nt{NORA}: це \emph{газ}, він прискорює ритм. Другим --- \nt{ACET}: це \emph{гальмо}, він сповільнює. Грубо,
\begin{empheq}[box=\keybox]{equation}
f \;\approx\; f_0 \;+\; a_S\,S \;-\; a_P\,P ,
\label{eq:gasbrake}
\end{empheq}
де $f_0\approx100$ ударів на хвилину --- власний ритм ритмоводія, $S$ і $P$ --- активність газу й гальма. У спокої гальмо натиснуте, і серце б'ється зазвичай порядку шістдесяти--сімдесяти разів на хвилину; під навантаженням гальмо відпускають і натискають газ.

\begin{figure}[t]
\centering
\begin{tikzpicture}[>=Stealth,
  blk/.style={draw,very thick,rounded corners=2pt,align=center,font=\small},
  pace/.style={circle,draw,very thick,double,double distance=1.2pt,minimum size=1.8cm,inner sep=0pt,font=\scriptsize,fill=red!8,align=center},
  inh/.style={-{Bar[width=7pt]},thick,red!70!black}]
\node[blk,fill=blue!5,text width=1.6cm] (B) at (0,0) {команда\\мозку};
\node[pace] (H) at (6.2,0) {ритмо-\\водій};
\draw[->,very thick,orange!85!black] (B.north east) to[bend left=20] node[above,font=\scriptsize,black,align=center] {\nt{NORA}: газ, секунди} (H.north west);
\draw[inh,very thick,blue!60!black] (B.south east) to[bend right=20] node[below,font=\scriptsize,black,align=center] {\nt{ACET}: гальмо, частки секунди} (H.south west);
\draw[->,very thick] (H) -- (8.4,0) node[right,font=\small] {частота $f$};
\node[blk,fill=orange!10,text width=1.6cm] (A) at (0,-2.6) {\nt{ADRE}\\у кров};
\draw[->,thick,orange!85!black] (B) -- (A);
\foreach \y/\lab in {-2.0/{серце},-2.6/{судини},-3.2/{печінка, м'язи}}{
  \draw[->,thick,densely dashed,orange!85!black] (A.east) -- (4.3,\y) node[right,font=\scriptsize,black] {\lab};}
\end{tikzpicture}
\caption{Два проводи протилежного знака до ритмоводія серця: газ \nt{NORA} повільний, гальмо \nt{ACET} швидке. Унизу --- той самий газ широкомовною шиною: \nt{ADRE}-клітини викидають адреналін у кров, і сигнал отримують усі органи з відповідним рецептором (штрихові стрілки).}
\label{fig:gasbrake}
\end{figure}

Це двопровідний код розділу~\ref{sec:glutcomputer} і пара м'язів розділу~\ref{sec:muscles}, тільки тепер проводи керують не суглобом, а темпом ритмоводія. І швидкість у проводів різна. Обидва працюють як фільтри нижніх частот, але гальмо пропускає зміни в кілька разів швидше за газ~\cite{Berger1989}: у \nt{ACET} шлях короткий: пов'язаний із рецептором білок сам відкриває калієвий канал, без другого посередника всередині клітини~\cite{Logothetis1987gbg}, а \nt{NORA} діє через ланцюжок реакцій усередині клітини. Інженер упізнає тут прийом із двома приводами різної швидкості: швидкий робить точне миттєве підлаштування, повільний --- великі тривалі зміни.

Тут же знаходить свою роль \nt{ADRE}, про який у таблиці~\ref{tab:types} було написано, що його роль у мозку неясна. Поза мозком вона ясна. Частина клітин газового шляху не посилає провід до органа, а викидає адреналін просто в кров. Це той самий газ, але широкомовною шиною: за десятки секунд він доходить до серця, судин, печінки й м'язів одразу і тримається хвилини. Адресний газ \nt{NORA} підлаштовує один орган; широкомовний газ \nt{ADRE} переводить у режим бігу все тіло.

\section{Каскад зі зворотним зв'язком}

Багато гормонів випускає не одна залоза, а каскад. Невелика група нейронів в основі мозку викидає в кров перший сигнал; він змушує залозу-посередника випустити другий; той --- кінцеву залозу випустити гормон, який і діє на тіло. А кінцевий гормон повертається до перших ступенів і гальмує їх. Виходить підсилювач із трьох каскадів, охоплений від'ємним зворотним зв'язком, --- регулятор рівня, як у \nt{SERO}-системи у формулі~\eqref{eq:feedback}.

Опишімо кожен ступінь RC-колом зі сталою часу $\tau$ і підсиленням $g$, а зворотний зв'язок --- коефіцієнтом $k$:
\begin{equation}
\tau\,\frac{\dd x_1}{\dd t} = -x_1 + g\bigl[u - k\,x_3\bigr]_+ ,\qquad
\tau\,\frac{\dd x_2}{\dd t} = -x_2 + g\,x_1 ,\qquad
\tau\,\frac{\dd x_3}{\dd t} = -x_3 + g\,x_2 ,
\label{eq:cascade}
\end{equation}
де $u$ --- команда мозку, $x_3$ --- рівень кінцевого гормону. Підсилення петлі $K=g^3k$. У рівновазі $x_3=g^3u/(1+K)$, і, як у будь-якого регулятора зі зворотним зв'язком, велика петля робить рівень нечутливим до розкиду підсилень ступенів. Але кожен ступінь зсуває фазу: на частоті $\omega=\sqrt3/\tau$ три однакові ступені разом дають зсув на пів оберту й ослаблення у вісім разів. Звідси класичний результат теорії автоматичного керування:
\begin{empheq}[box=\keybox]{equation}
K \;<\; 8 ,
\label{eq:cascadestable}
\end{empheq}
інакше регулятор починає коливатися з періодом близько $2\pi\tau/\sqrt3\approx3.6\,\tau$.

\begin{keyblock}
\begin{proposition}[Точність проти стійкості]
\label{prop:cascade}
У триступеневого каскаду з пропорційним зворотним зв'язком похибка рівня дорівнює $1/(1+K)$, а стійкий він лише за $K<8$. Тому точніше, ніж приблизно на десять відсотків, такий регулятор рівня не тримає; спроба підняти підсилення перетворює його на генератор коливань.
\end{proposition}
\end{keyblock}

Ми перевірили це на моделі~\eqref{eq:cascade} (рис.~\ref{fig:cascade}). За $K=4$ рівень після ввімкнення трохи коливається й усталюється. За $K=7$ він теж усталюється, але повільно. За $K=12$ коливання не згасають і не ростуть без меж: ступінь не може видати від'ємний сигнал, і каскад виходить на рівні пульсації між $0.83$ і $1.24$ від середнього рівня з періодом $3.7$--$3.9\,\tau$.

\begin{figure}[t]
\centering
\begin{tikzpicture}[>=Stealth,x=0.3cm,y=1.2cm]
\draw[->,gray!70!black] (0,0) -- (41.5,0) node[right,font=\small,black] {$t/\tau$};
\draw[->,gray!70!black] (0,0) -- (0,2.8) node[left,font=\small,black] {$x_3/x_3^*$};
\foreach \t in {0,10,20,30,40}{\draw[gray!60!black] (\t,-0.05) -- (\t,0.05); \node[below,font=\scriptsize] at (\t,0) {\t};}
\foreach \f in {1,2}{\draw[gray!60!black] (-0.3,\f) -- (0.3,\f); \node[left,font=\scriptsize] at (0,\f) {\f};}
\draw[densely dashed,gray!60!black] (0,1) -- (40,1);
\draw[thick,blue!60!black] plot file {figures/cascade_K4.dat};
\draw[thick,red!70!black] plot file {figures/cascade_K12.dat};
\node[font=\scriptsize,blue!60!black,anchor=west] at (16,0.55) {$K=4$: усталюється};
\node[font=\scriptsize,red!70!black,anchor=west] at (16,1.75) {$K=12$: пульсує};
\end{tikzpicture}
\caption{Каскад із трьох ступенів зі зворотним зв'язком за~\eqref{eq:cascade}, рівень кінцевого гормону в частках рівноважного $x_3^*$. Коли підсилення петлі менше за вісім, рівень усталюється; коли більше --- каскад сам породжує рівні пульсації.}
\label{fig:cascade}
\end{figure}

Справжні гормони каскадів справді виходять пульсами: рівень гормону стресу, наприклад, коливається з періодом близько години. За однією з гіпотез, ці пульсації породжує сам зворотний зв'язок між ступенями із запізненням, і окремий генератор ритму для них не потрібен~\cite{Walker2010}. Це та сама причина коливань, що й у петлі \nt{GLUT}--\nt{GABA} у твердженні~\ref{prop:ei} і в петлі розтягу м'яза в умові~\eqref{eq:delaystable}: від'ємний зворотний зв'язок, який запізнюється.

\section{Імпульсний код гормонів}

Пульси бувають не лише побічним ефектом, а й кодом. Нейрони, що керують розмноженням, випускають свій сигнал у кров короткими порціями приблизно раз на годину. Якщо залозі-отримувачу подавати той самий сигнал не порціями, а безперервно, вона не працює на повну силу, як за пульсів, а замовкає; якщо знову подавати порціями раз на годину, робота відновлюється~\cite{Belchetz1978}. Отже, отримувач читає не рівень, а \emph{частоту} порцій.

Для інженера тут немає загадки. Рецептор, який за тривалої дії речовини втомлюється й перестає відповідати, --- фільтр верхніх частот, як синапс, що виснажується~\eqref{eq:depression}, з розділу~\ref{sec:code}. Безперервний сигнал він глушить, а зміни пропускає. Такому отримувачеві можна щось повідомити лише імпульсами, і інформація переходить із рівня в частоту й форму порцій. Різні частоти порцій до того ж по-різному діють на різні клітини однієї й тієї самої залози, тож однією хімічною лінією можна передати більше ніж одну величину --- як одним проводом \nt{DOPA} в розділі~\ref{sec:dofaalt} ішли швидка й повільна складові.

\section{Добовий ритм: повільний генератор режимів}

Найповільніший ритм тіла --- добовий. Його породжує від'ємний зворотний зв'язок із запізненням усередині клітин: гени виробляють білки, які із затримкою в кілька годин пригнічують власне виробництво~\cite{TakahashiJS2017}. Знову запізнілий від'ємний зворотний зв'язок --- вчетверте в цій книзі, --- і знову ритм, цього разу з періодом близько доби. Головний генератор --- невелика група з десятків тисяч нейронів в основі мозку, ритм яких синхронізує решту клітин тіла.

Власний період цього генератора в людини не рівно доба, а в середньому $24.18$ години~\cite{Czeisler1999}. Щодня його підводить світло, що надходить від датчиків ока (розділ~\ref{sec:sensors}). Для електронника це \emph{фазове автопідстроювання частоти}: генератор із власною частотою $\omega_0$ підлаштовується до зовнішнього сигналу з частотою $\omega$, і різниця фаз $\varphi$ підпорядковується рівнянню
\begin{empheq}[box=\keybox]{equation}
\frac{\dd\varphi}{\dd t} \;=\; (\omega-\omega_0) \;-\; K_\varphi\,\sin\varphi .
\label{eq:pll}
\end{empheq}
Підстроювання утримується, якщо $|\omega-\omega_0|<K_\varphi$. Генераторові з періодом $24.18$ години треба щодня зсуватися приблизно на одинадцять хвилин; ранкового світла для такого зсуву зазвичай вистачає. У полярну ніч або в печері без світла підстроювання немає, і ритм переходить на власний період.

Добовий генератор --- не тактовий генератор комп'ютера. Він не відлічує кроків обчислень і не синхронізує імпульсів нейронів. Він перемикає \emph{режими}: сон і неспання, рівні гормонів, температуру тіла, готовність до їжі. Це повільний широкомовний регістр того самого роду, що й регістри розділів~\ref{sec:serocomputer}--\ref{sec:acet}, тільки його значення змінюється за розкладом, підведеним до сонця.

\begin{keyblock}
\begin{proposition}[Хімічний вихід]
\label{prop:chemoutput}
Повільний вихід машини збудовано з тих самих деталей, що й швидкий. Органи отримують два проводи протилежного знака й різної швидкості --- газ \nt{NORA} і гальмо \nt{ACET}, --- а весь організм одразу --- широкомовні сигнали кров'ю, зокрема адреналін \nt{ADRE}. Рівні гормонів тримають каскади з від'ємним зворотним зв'язком, точність яких обмежена стійкістю; частина гормонів кодується частотою порцій; добовий ритм --- повільний генератор із фазовим підстроюванням за світлом. Будову шляхів і досліди з порціями та з періодом виміряно; пояснення пульсацій каскаду самим зворотним зв'язком --- гіпотеза.
\end{proposition}
\end{keyblock}

\section{Підсумок}

\begin{table}[t]
\centering
\footnotesize
\setlength{\tabcolsep}{4pt}
\renewcommand{\arraystretch}{1.15}
\begin{tabular}{>{\raggedright}p{3.1cm}>{\raggedright}p{4.8cm}>{\raggedright\arraybackslash}p{4.9cm}}
\toprule
Що робить хімічний вихід & Як & Що відомо \\
\midrule
підлаштовує органи & газ \nt{NORA} і гальмо \nt{ACET}~\eqref{eq:gasbrake} & виміряно; гальмо швидше за газ \\
переводить усе тіло в режим бігу & адреналін \nt{ADRE} у кров & виміряно \\
тримає рівні гормонів & каскад зі зворотним зв'язком, $K<8$~\eqref{eq:cascadestable} & каскади й зворотний зв'язок виміряно; пульсації від самої петлі --- гіпотеза \\
передає частотою порцій & утомлений рецептор як фільтр верхніх частот & залежність від порцій виміряно \\
перемикає режими за добу & генератор із фазовим підстроюванням за світлом~\eqref{eq:pll} & період і підстроювання світлом виміряно \\
\bottomrule
\end{tabular}
\caption{Як машина керує хімією тіла.}
\label{tab:chemoutput}
\end{table}

У машини два виходи (таблиця~\ref{tab:chemoutput}). Швидкий --- м'язи: мілісекунди, адресно, у кожен суглоб. Повільний --- хімія тіла: секунди, хвилини й доба, двома проводами до органів і кров'ю до всіх одразу. Деталі в обох виходів ті самі --- два проводи протилежного знака, ритмоводії, зворотний зв'язок та його запізнення, --- і ті самі, з яких зібрано машину всередині.

\section{Задачі}

\begin{exercise}[\lvlA]
\label{ex:16:1}
\emph{Кров як фільтр.} Гормон виводиться з крові з періодом напіврозпаду $t_{1/2}=10$~хв і випускається короткими порціями кожні $T=60$~хв. У скільки разів максимум концентрації більший за мінімум і у скільки разів ослаблена основна гармоніка порцій? За якого $t_{1/2}$ максимум перевищує мінімум лише вдвічі?
\end{exercise}

\begin{exercise}[\lvlA]
\label{ex:16:2}
\emph{Добове автопідстроювання.} Власний період генератора~\eqref{eq:pll} дорівнює $24.18$~год, світло зсуває фазу не більше ніж на годину за добу: $K_\varphi=2\pi/24$~рад/доб. Знайдіть усталену різницю фаз $\varphi^*$ у хвилинах і час релаксації. За скільки діб розузгодження після стрибка зовнішнього сигналу на $\pm6$~год стане меншим за годину?
\end{exercise}

\begin{exercise}[\lvlB]
\label{ex:16:3}
\emph{Точність проти стійкості.} Лінеаризований каскад~\eqref{eq:cascade} подовжили до $n$ однакових ступенів. Знайдіть критичне підсилення $K_c(n)$ і найменшу досяжну похибку $1/(1+K_c)$ за $n=3$ і $n=6$. До чого прямує ця похибка за $n\to\infty$?
\end{exercise}

\begin{exercise}[\lvlB]
\label{ex:16:4}
\emph{Газ і гальмо.} У~\eqref{eq:gasbrake} внески газу й гальма --- виходи RC-кіл із $\tau_S=2$~с і $\tau_P=0.4$~с, команди не більші за $80$ і $60$ ударів на хвилину. Ритм $f_0=100$; у спокої гальмо $35$, газ $0$, $f=65$. За який час можна вийти на $f=120$? А якщо гальма не відпускати й працювати лише газом?
\end{exercise}

\begin{exercise}[\lvlB]
\label{ex:16:5}
\emph{Моделювання: каскад із затримкою.} У зворотний зв'язок функції \texttt{cascade} з \texttt{models/\allowbreak chem\_models.py} (скопіюйте її, файл не запускайте) додайте затримку $d$: перший ступінь бачить $x_3(t-d)$. Знайдіть $K_c$ і період на межі за $d/\tau=0$, $0.5$, $1$, $2$ і перевірте моделюванням за $0.9K_c$ і $1.1K_c$.
\end{exercise}

\begin{exercise}[\lvlB]
\label{ex:16:6}
\emph{Отримувач, що читає порції.} Рецептор утомлюється: $y=[c-a]_+$, $\tau_a\,\dd a/\dd t=c-a$, $\tau_a=30$~хв. Гормон надходить порціями по $6$~хв із періодом $T$ за тієї самої середньої дози $\bar c$. Знайдіть середній відгук за $T=15$, $60$, $120$~хв і $T\to\infty$. Який він за сталої концентрації $\bar c$?
\end{exercise}

\begin{exercise}[\lvlC]
\label{ex:16:7}
\emph{Генератор чи зворотний зв'язок?} Запропонуйте дослід, який відрізнить пульсації, породжені запізнілим зворотним зв'язком каскаду~\eqref{eq:cascade}, від окремого ритмоводія. Чи можна розрізнити гіпотези за зсувом періоду, якщо $K$ вдається змінити лише в півтора раза, а період вимірюється з точністю $5\%$?
\end{exercise}

\chapter{Налагодження машини}
\chaptermark{Налагодження машини}
\label{sec:debugging}

\section{Як налагоджують машину без схеми}

Інженер, якому дали робочу плату без схеми, налагоджує її чотирма прийомами. Він ставить \emph{щуп} і дивиться, що йде проводом. Він \emph{подає тестовий сигнал} і дивиться, що змінилося. Він \emph{вимикає} шматок схеми і дивиться, що перестало працювати. І він \emph{читає регістри керування}, щоб дізнатися, в якому режимі працює плата. Уся ця книга --- спроба прочитати схему мозку, і в цьому розділі подивимося, якими з чотирьох прийомів інженери вже вміють користуватися і що кажуть їм попередні розділи про те, чого чекати.

Найпомітніший сьогодні приклад такого налагодження --- інтерфейси мозок--комп'ютер: пристрої, які читають імпульси нейронів і перетворюють їх на команди для зовнішніх машин або, навпаки, подають у мозок сигнали ззовні. Їм і присвячено більшу частину розділу, і насамперед тому, що робить компанія Neuralink.

\section{Щуп: що чує електрод}

Щуп мозку --- тонкий електрод, уведений у тканину. Він чує імпульси нейронів, що лежать у радіусі порядку сотні мікрометрів, --- зазвичай від одного до кількох. Імпульс на електроді --- сплеск у десятки--сотні мікровольтів тривалістю близько мілісекунди, і за його формою можна відрізнити сусідні нейрони один від одного. Найкраще чутно великі \nt{GLUT}-нейрони: їх у корі більшість (таблиця~\ref{tab:types}), і їхні струми сильніші.

Одразу видно головну трудність. Добрий щуп сьогодні --- це сотні або тисячі електродів, а нейронів у мозку $8.6\cdot10^{10}$. Ми підслуховуємо приблизно одну стомільйонну машини. Чому ж із такої мізерної вибірки взагалі можна щось зрозуміти? Відповідь дав розділ~\ref{sec:code}. Сусідні нейрони шумлять узгоджено, і усереднення за групою впирається в межу (твердження~\ref{prop:pooling}): сотня нейронів несе майже стільки ж, скільки тисяча. До того ж задача, яку розв'язує рухова кора, маловимірна: у руки небагато суглобів (розділ~\ref{sec:muscles}), і активність сотень нейронів, що нею керують, укладається в простір лише з одного-двох десятків спільних змінних~\cite{Gallego2017}. Щуп на сотню нейронів чує майже всі ці змінні. Якби мозок зберігав незалежне повідомлення в кожному нейроні, таке підслуховування було б марним.

\section{Потік даних: стиснення просто на щупі}

Щуп видає величезний потік. Щоб побачити форму імпульсу, кожен електрод треба оцифровувати приблизно $20$ тисяч разів на секунду з точністю близько десяти біт. Тисяча електродів дає
\begin{empheq}[box=\keybox]{equation}
\begin{aligned}
R_{\text{сирий}} \;&=\; N\,f_s\,b \;\approx\; 10^3\cdot2\cdot10^4\cdot10 \;\approx\; 2\cdot10^8\ \text{біт/с},\\
R_{\text{імп}} \;&\approx\; N\,r\,\log_2\frac{e}{r\,\Delta t} \;\approx\; 8\cdot10^4\ \text{біт/с}.
\end{aligned}
\label{eq:probedata}
\end{empheq}
Друга оцінка --- пропускна здатність потоку імпульсів~\eqref{eq:spikeentropy} за $r=10$ імпульсів на секунду і точності $\Delta t=1\ms$: усе, що в ньому насправді є, займає у дві з половиною тисячі разів менше. Для вживленого пристрою це вирішальна різниця: сирий потік не можна передати радіо крізь шкіру з тим живленням, яке можна дозволити собі всередині голови, не нагріваючи тканини. Тому вживлений пристрій має виділяти імпульси просто на кристалі поруч з електродами й передавати назовні лише події --- номер каналу й момент імпульсу. У першому описі системи Neuralink до цього ще не дійшли: кристали підсилюють і оцифровують сигнали, увесь сирий потік іде кабелем, а імпульси виділяє програма ззовні; стиснення на кристалі й бездротовий зв'язок названо там планом~\cite{Musk2019}.

Це знайомий прийом. Око стискає сигнал у сто разів, перш ніж надіслати його довгим кабелем (розділ~\ref{sec:sensors}); подієва камера передає не кадри, а зміни. Щуп, щоб умістися в провід, змушений робити те саме, що робить сам мозок: говорити імпульсами і лише про новини.

\section{Що робить Neuralink}

Пристрій Neuralink --- щуп, зібраний під ці обмеження~\cite{Musk2019}. Замість жорсткої ґратки голок --- багато гнучких ниток, тонших за волосину, з електродами вздовж кожної: у першому описі системи --- до $3072$ електродів на $96$ нитках, по $32$ на нитку, а в пристрої, вживленому людям, за повідомленнями компанії, близько тисячі каналів на $64$ нитках. Гнучкі нитки менше ушкоджують тканину й краще переносять те, що мозок усередині черепа весь час трохи зсувається. Нитки вводить робот, оминаючи судини. У першому описі, як сказано вище, сирий потік~\eqref{eq:probedata} ішов назовні кабелем. Пристрій для людей, за повідомленнями компанії, влаштований інакше: корпус повністю закритий шкірою, імпульси виділяються всередині, дані йдуть радіо, живлення надходить бездротово.

Завдання першого пристрою --- читання. Нитки ставлять у ту частину кори, яка планує рухи руки, і декодер перетворює імпульси на рух курсора на екрані. Людина з паралічем рук керує комп'ютером, уявляючи рух. Сама ідея не нова: інтерфейси на жорстких ґратках із сотні електродів давно дозволяли керувати курсором~\cite{Hochberg2006}, писати текст, уявляючи рухи руки під час письма~\cite{Willett2021}, і навіть перетворювати спроби говорити на текст на екрані зі швидкістю близько шістдесяти слів на хвилину~\cite{Willett2023}. Нове в Neuralink --- інженерія: кількість каналів, бездротовість, закритий корпус і роботизоване введення, тобто спроба зробити щуп, який можна носити роками поза лабораторією. Наскільки нам відомо, результати випробувань на людях публікувала переважно сама компанія, а не рецензовані статті; компанія повідомляла, зокрема, що частина ниток після введення змістилася і кількість робочих каналів зменшилася. Для налагодження це звичайна неприємність: щуп, який рухається відносно плати, чує вже інші проводи.

Другий напрям компанії --- запис: пристрій для зору, який подає сигнали в зорову кору людини, що втратила зір. Про нього нижче, у підрозділі про запис.

\section{Читання: декодер як отримувач}

Декодер інтерфейсу --- це отримувач у сенсі твердження~\ref{prop:reader}: він сам вирішує, яку частину повідомлення прочитати. Практично всі інтерфейси читають \emph{частоти груп}: рахують імпульси кожного каналу у вікнах по кілька десятків мілісекунд і додають їх із вагами, дібраними так, щоб виходив задуманий рух. Це частотний код групи з підрозділу~\ref{sec:rate}, і його обрано не випадково: коли на вікно припадає кілька імпульсів, частота одного нейрона не визначена, а сума за сотнею вже працює.

Скільки інформації вдається витягти? Письмо «уявною рукою» йшло зі швидкістю близько дев'яноста знаків на хвилину~\cite{Willett2021}, тобто півтора знака на секунду: навіть за п'яти бітів на знак це менше ніж вісім біт на секунду. Керування курсором, за повідомленнями Neuralink, дає порядку десяти біт на секунду. А верхня межа~\eqref{eq:spikeentropy} для \emph{одного} нейрона --- десятки біт на секунду, для сотні --- тисячі. Інтерфейс витягає з підслуханих нейронів малу частку того, що вони могли б нести. Вузьке місце --- не провід, а те, скільки незалежних змінних несе група (твердження~\ref{prop:pooling}) і наскільки добре декодер розуміє код, який, як ми бачили в розділі~\ref{sec:code}, до кінця не розшифровано.

Є й друга особливість, знайома з розділу~\ref{sec:motorlearning}. Декодер і мозок учаться один в одного. Мозок бачить, куди поїхав курсор, отримує помилку і підлаштовує свої команди --- те саме навчання рухів проводом помилки. А інженери час від часу заново добирають ваги декодера, бо нейрони під нитками змінюються і зв'язки в мережі перенавчаються. Виходять два учні, які підлаштовуються один під одного; щоб така пара не втратила стійкості, швидкості їхнього навчання зручно розвести: нехай один учиться швидко, а другий повільно.

\begin{keyblock}
\begin{proposition}[Чому щуп на тисячу нейронів працює]
\label{prop:probe}
Інтерфейс, що підслуховує мізерну частку нейронів, може керувати рухом тому, що активність групи маловимірна і її шуми узгоджені: кілька сотень нейронів несуть майже всі спільні змінні. З тієї самої причини інтерфейс витягає лише кілька біт на секунду --- стільки, скільки незалежних змінних у задачі, а не стільки, скільки вміщує потік імпульсів. Роботу інтерфейсів виміряно; наскільки кращим стане декодер, коли код буде зрозуміло, --- відкрите питання.
\end{proposition}
\end{keyblock}

\section{Запис: важче, ніж читання}

Подати сигнал у машину важче, ніж прочитати. Електрод, через який пропускають струм, збуджує не один нейрон, а всі нейрони й проводи навколо себе, без розбору типу й знака. Записати ним код, у якому важливо, \emph{який} нейрон і \emph{коли} спрацював (розділ~\ref{sec:code}), не можна: можна лише грубо задати, \emph{де} і \emph{наскільки сильно}.

Для зору цього частково вистачає. Струм через електрод у зоровій корі людина бачить як світну точку в певному місці поля зору; коли точки запалювали одночасно, до шістнадцяти разом, людина, що втратила зір, упізнавала деякі літери й розрізняла межі предметів~\cite{Fernandez2021}. Це код місцем, і його записати можна. Але згадайте розділ~\ref{sec:sensors}: око надсилає в мозок не пікселі, а стиснений опис --- краї, зміни, посвітлішання й потемніння, окремими проводами. Запис «точок світла» в кору --- це запис пікселів у машину, яка чекає на вході зовсім іншого коду. Тому штучний зір поки грубіший, ніж можна було б чекати з кількості електродів. Є й випробувальні сигнали, яким електрод не потрібен: зорові ілюзії подають через око незвичний вхід і виводять машину зі звичайного режиму, а кожна помилка вказує на свій механізм (підрозділ~\ref{sec:illusions}).

\section{Чого щуп не бачить}

Електрод чує адресну шину --- імпульси \nt{GLUT}-нейронів і, гірше, дрібних \nt{GABA}. Але в розділах~\ref{sec:serocomputer}--\ref{sec:acet} ми бачили, що режим роботи всієї машини задають широкомовні регістри: концентрації \nt{SERO}, \nt{DOPA}, \nt{NORA}, \nt{ACET}. Їх електрод не чує: нейронів-джерел мізерно мало, і їхній сигнал --- не імпульс поруч зі щупом, а концентрація в об'ємі. Налагоджувач, який бачить шину даних, але не бачить регістрів керування, завжди дивуватиметься: той самий вхід даватиме різні відповіді, бо десь змінилося $g$, $a$ чи $\tau_\gamma$. Концентрації можна вимірювати хімічними датчиками просто в тканині~\cite{Bunin1998}, але в інтерфейсах такі датчики поки не використовують. Зовсім поза полем зору щупа лишається й другий, повільний вихід машини --- гормони в крові (розділ~\ref{sec:chemoutput}): їх вимірюють у пробах крові, а не електродом.

Не бачить щуп і ваг --- а саме в них зберігається майже вся пам'ять і все вивчене. Їх можна лише вгадувати з того, як змінюються відповіді. І не бачить він болю так, як його відчуває людина: у розділі~\ref{sec:pain} ми бачили, що силу сигналу тривоги задає сама машина --- воротами в спинному мозку й регуляторами згори, тож із датчика не можна прочитати, наскільки боляче.

\section{Підсумок: налагодження і відкриті питання}

\begin{table}[t]
\centering
\footnotesize
\setlength{\tabcolsep}{4pt}
\renewcommand{\arraystretch}{1.15}
\begin{tabular}{>{\raggedright}p{2.2cm}>{\raggedright}p{3.6cm}>{\raggedright}p{3.4cm}>{\raggedright\arraybackslash}p{4.0cm}}
\toprule
Прийом налагодження & Що робить інтерфейс & Що пояснює книга & Що відомо \\
\midrule
щуп & чує сотні--тисячі нейронів & маловимірність і узгоджений шум (розд.~\ref{sec:code}) & виміряно \\
стиснення на щупі & передає події замість сирого сигналу & пропускна здатність потоку імпульсів~\eqref{eq:spikeentropy} & інженерно розв'язано \\
читання & частоти груп $\to$ рух, текст, мовлення & декодер --- отримувач, частотний код групи & працює; кілька біт на секунду \\
спільне навчання & мозок і декодер підлаштовуються один під одного & провід помилки (розд.~\ref{sec:motorlearning}) & спостерігається; правила підлаштування --- предмет досліджень \\
запис & струм запалює точки в полі зору & код місцем записати можна, часовий --- ні (розд.~\ref{sec:sensors}) & точки виміряно; корисного зору --- поки немає \\
читання регістрів & майже не робиться & широкомовні канали (розд.~\ref{sec:serocomputer}--\ref{sec:acet}) & хімічні датчики є в дослідах \\
\bottomrule
\end{tabular}
\caption{Налагодження машини: що вміють інтерфейси мозок--комп'ютер і що про них кажуть розділи книги.}
\label{tab:debugging}
\end{table}

Таблиця~\ref{tab:debugging} підсумовує розділ. Інтерфейси мозок--комп'ютер уже вміють ставити щуп на адресну шину й читати з неї кілька біт на секунду, і те, що це взагалі працює, пояснюється маловимірністю й узгодженим шумом із розділу~\ref{sec:code}. Записувати в машину вони вміють лише грубо, бо її вхідний код стиснений і адресований окремим проводам. А регістри керування й ваги --- те, що робить машину здатною навчатися й налаштовуватися, --- щупові поки майже не видно.

Кожне відкрите питання розділу~\ref{sec:synthesis} --- це й задача для налагоджувача. Який код читати, вирішиться, коли щупи стануть щільнішими, а декодери навчаться користуватися часом імпульсів, а не лише частотами. Як навчається багатошарова мережа, можна перевірити, поставивши щупи одразу в кілька шарів і дивлячись, як змінюються їхні відповіді під час навчання. Що передають широкомовні канали, стане видно, коли до електричних щупів додадуться хімічні. Ця книга --- спроба прочитати схему машини за її поведінкою і за законами, які знає будь-який інженер. Налагоджувачі, що будуються зараз, дадуть змогу перевірити цю схему щупом.

\section{Задачі}

\begin{exercise}[\lvlA]
\label{ex:17:1}
\emph{Бюджет радіоканалу.} Передавання крізь шкіру коштує $1$~нДж на біт, передавачеві можна витрачати не більше $10$~мВт. Яка потужність потрібна для сирого потоку~\eqref{eq:probedata} за $N=1000$ електродів і для потоку імпульсів? Якої енергії на біт вимагав би сирий потік?
\end{exercise}

\begin{exercise}[\lvlA]
\label{ex:17:2}
\emph{Скільки біт у сотні нейронів.} Декодер додає імпульси $N$ нейронів за вікна $50\ms$ за $r=20$~імп/с, попарна кореляція шуму $c=0.1$, сигнал змінює частоту групи на $30\%$ (середньоквадратично). Оцініть швидкість $\tfrac12\log_2(1+\text{SNR})$ за вікно за $N=10$, $100$, $1000$ і $N\to\infty$. А за $c=0$, $N=1000$?
\end{exercise}

\begin{exercise}[\lvlB]
\label{ex:17:3}
\emph{Швидкість інтерфейсу-клавіатури.} Інтерфейс півтора раза на секунду обирає один із $N=32$ знаків: задуманий з імовірністю $P$, помилки рівноймовірні. Скільки біт на секунду він передає за $P=0.99$, $0.94$, $0.8$? Скільки знаків на секунду треба за $P=0.94$ для $10$~біт/с?
\end{exercise}

\begin{exercise}[\lvlB]
\label{ex:17:4}
\emph{Моделювання: декодер групи.} $N$ нейронів із $\lambda_i=[b+m\,\mathbf v\cdot\mathbf p_i]_+$, $b=20$, $m=15$~імп/с, вікна $50\ms$, $\mathbf v$ зі стандартним відхиленням $0.5$ за компонентою; усі бачать спільний шум, $\mathbf v+\boldsymbol\xi$, $\sigma_\xi=0.2$. Змоделюйте незміщений декодер вектором групи. За якого $N$ обидва шуми рівні і скільки біт на компоненту за вікно дає $N\to\infty$?
\end{exercise}

\begin{exercise}[\lvlB]
\label{ex:17:5}
\emph{Два учні.} Мозок видає команду $m=b\,v^*$, декодер --- $v=d\,m$, $\langle v^{*2}\rangle=1$. Обидва після кожної спроби роблять крок градієнта за $\tfrac12(v-v^*)^2$ зі швидкістю $\eta$. Змоделюйте з $b=1.2$, $d=1$ за $\eta=0.8$, $1.1$, $1.5$, $2$. Кожен поодинці стійкий за $\eta<2$; за якого $\eta$ стійка пара?
\end{exercise}

\begin{exercise}[\lvlB]
\label{ex:17:6}
\emph{Проєктування: конвеєр щупа.} $1024$ канали по $20$~тис. відліків на секунду, нейрони $10$~імп/с, радіо $1$~нДж/біт. Який поріг за білим шумом відліків дає не більше $0.1$~Гц хибних імпульсів на канал? Передаємо числа імпульсів за $20\ms$ по $2$~біти: яка потужність і у скільки разів вона менша, ніж для сирих $10$-бітних відліків?
\end{exercise}

\begin{exercise}[\lvlC]
\label{ex:17:7}
\emph{Межа швидкості інтерфейсу.} Оцініть $R\le DW\log_2(1+\text{SNR})$ для $D=10$ змінних зі смугою $W=5$~Гц за $\text{SNR}\approx0.9$ (шум, схожий на сигнал) і $90$ (незалежний шум $1000$ нейронів). Виміряно близько $10$~біт/с. Який дослід відрізнить два пояснення розриву: шум, схожий на сигнал, і незнання коду?
\end{exercise}

\chapter{Типи нейронів за електричною функцією}
\chaptermark{Нейрони як прилади}
\label{sec:eetypes}

У розділі~\ref{sec:neurontypes} ми розсортували нейрони за тим, який медіатор викидає їхній вихід. Електронник розсортував би їх інакше: за тим, що нейрон робить зі своїм сигналом, --- яким приладом він працює. Головний прилад книги ми знаємо з розділу~\ref{sec:glutcomputer}: інтегрувальний нейрон із витоком~\eqref{eq:glutneuron}, RC-коло з компаратором. Але в мозку є й інші. У таблиці~\ref{tab:eetypes} їх зібрано в чотири групи, і майже кожен уже траплявся в книзі.

\begin{center}
\footnotesize
\setlength{\tabcolsep}{3pt}
\renewcommand{\arraystretch}{1.15}
\begin{longtable}{>{\raggedright}p{3.1cm}>{\raggedright}p{3.3cm}>{\raggedright}p{4.7cm}>{\raggedright\arraybackslash}p{2.4cm}}
\caption{Нейрони як прилади: назва, аналог в електроніці, що прилад робить і де він трапляється в книзі.}\label{tab:eetypes}\\
\toprule
Нейрон & Прилад & Що робить & Де в книзі \\
\midrule
\endfirsthead
\toprule
Нейрон & Прилад & Що робить & Де в книзі \\
\midrule
\endhead
\bottomrule
\endfoot
\multicolumn{4}{l}{\emph{Фільтри й інтегратори}} \\
інтегрувальний нейрон із витоком & RC-інтегратор із компаратором & додає входи у вікні $\tau$ і спрацьовує біля порога & розд.~\ref{sec:glutcomputer}, \eqref{eq:glutneuron} \\
детектор збігів & І з часовим вікном & спрацьовує, лише якщо входи надійшли майже разом; дуже мале $\tau$ & розд.~\ref{sec:glutcomputer}, \ref{sec:code}, \eqref{eq:window} \\
фазичний нейрон & фільтр верхніх частот, детектор фронту & відповідає на зміну входу й замовкає & розд.~\ref{sec:sensors} \\
адаптивний нейрон & фільтр верхніх частот з автоматичним регулюванням підсилення & частота падає за сталого входу & розд.~\ref{sec:sensors}, \eqref{eq:nakarushton} \\
резонатор & коливальний контур, смуговий фільтр & найкраще відповідає на вхід своєї частоти & підрозділ~\ref{sec:resonator} \\
інтегратор без витоку & інтегратор на операційному підсилювачі & перетворює швидкість на положення й тримає його & підрозділ~\ref{sec:perfectint} \\
\midrule
\multicolumn{4}{l}{\emph{Перетворювачі}} \\
тонічний нейрон & перетворювач напруга---частота & частота лінійно йде за входом і майже не втомлюється & розд.~\ref{sec:code} \\
градуальний нейрон (без імпульсів) & аналоговий підсилювач, буфер & передає плавну напругу на коротку відстань & розд.~\ref{sec:sensors} \\
випрямний нейрон & однопівперіодний випрямляч & частота не буває від'ємною, $[u]_+$ & розд.~\ref{sec:glutcomputer}, \ref{sec:gaba} \\
пара «увімкнення---вимкнення» & пара випрямлячів, двопровідний код & один відповідає на «більше», другий --- на «менше» & розд.~\ref{sec:glutcomputer}, \eqref{eq:dualrail} \\
нейрон із шунтувальним гальмуванням & аналоговий дільник, регулювання підсилення & ділить вхід, а не віднімає від нього & розд.~\ref{sec:gaba}, \eqref{eq:shunt} \\
\midrule
\multicolumn{4}{l}{\emph{Генератори й час}} \\
ритмоводій & релаксаційний генератор, мультивібратор & сам видає імпульси зі сталою частотою & розд.~\ref{sec:serocomputer}, \ref{sec:dofa} \\
ритмоводій із входом & генератор, керований напругою & власний ритм, який вхід прискорює чи сповільнює & розд.~\ref{sec:chemoutput} \\
пачковий нейрон & стробований генератор пачок & видає пачки частих імпульсів --- «тире» & розд.~\ref{sec:code}, \eqref{eq:burst} \\
нейрон, прив'язаний до фази & фазовий детектор, фазове автопідстроювання & спрацьовує на певній фазі вхідної хвилі & розд.~\ref{sec:sensors}, \ref{sec:chemoutput} \\
аксон із затримкою & лінія затримки & затримує сигнал на заданий час & розд.~\ref{sec:glutcomputer}, рис.~\ref{fig:itd} \\
\midrule
\multicolumn{4}{l}{\emph{Пам'ять і перемикання}} \\
бістабільний нейрон & тригер Шмітта, засувка & короткий вхід переводить його в тривалий стан & підрозділ~\ref{sec:bistable} \\
нейрон із дендритними гілками & мультиплексор, маленька багатошарова мережа & гілка пропускає вхід, лише якщо є дозвільний вхід & розд.~\ref{sec:synthesis} \\
\end{longtable}
\end{center}

Одне застереження важливіше за всю таблицю: це не різні клітини, а різні \emph{режими}. Той самий нейрон може бути інтегратором за повільного входу й детектором збігів, коли гальмування вкоротило його $\tau$ (розділ~\ref{sec:code}); ритмоводієм під час неспання й мовчазним приймачем уві сні. Який прилад вийде, вирішують іонні канали мембрани й широкомовні канали, що їх підлаштовують.

\section{Чотири прилади, яких у книзі ще не було}

\subsection{Резонатор}
\label{sec:resonator}

Деякі нейрони мають повільний струм, який протидіє будь-якій зміні напруги: напруга зростає --- струм тягне її вниз, падає --- тягне вгору, але із запізненням. Для електронника такий струм поводиться як індуктивність, а разом з ємністю мембрани вони утворюють коливальний контур. Нейрон найсильніше відповідає на вхід певної частоти, зазвичай від одиниць до одного-двох десятків герц, і слабше --- на повільніший і швидший~\cite{HutcheonYarom2000}. Це смуговий фільтр в одній клітині: із шумного входу він вибирає ритм своєї частоти, наприклад місцевий ритм розділу~\ref{sec:gaba}.

\subsection{Інтегратор без витоку}
\label{sec:perfectint}

Інтегратор із витоком пам'ятає вхід лише час $\tau$ --- десятки мілісекунд. Але окові, щоб стояти нерухомо, треба пам'ятати своє положення секундами: команда «повернути» надходить як короткочасна швидкість, а тримати око треба в новому положенні. Мережа розв'язує це тим самим зворотним зв'язком, що й пам'ять розділу~\ref{sec:glutcomputer}. Якщо група нейронів зі сталою часу $\tau$ збуджує сама себе з вагою $w$, то $\tau\,\dd r/\dd t=-r+w\,r+I$, і стала часу групи дорівнює
\begin{empheq}[box=\keybox]{equation}
\tau_{\text{еф}} \;=\; \frac{\tau}{1-w} .
\label{eq:perfectint}
\end{empheq}
За $\tau=0.1\,\text{с}$ і $w=0.995$ виходить $20\,\text{с}$: майже ідеальний інтегратор із нейронів, кожен з яких сам по собі забуває за десяту частку секунди~\cite{Seung1996}. Ціна --- точне налаштування: коли $w$ трохи більше за одиницю, інтегратор іде в насичення, коли трохи менше --- швидко забуває. Тому такому інтеграторові потрібне постійне підлаштування навчанням, на кшталт того, що описано в розділі~\ref{sec:motorlearning}.

\subsection{Бістабільний нейрон}
\label{sec:bistable}

У розділах~\ref{sec:glutcomputer} і~\ref{sec:gaba} тригер збирався з групи нейронів. Але буває тригер і в одній клітині. Деякі нейрони мають струм, який, щойно ввімкнувшись, сам підтримує збудження: короткий вхід переводить нейрон у тривалу роботу --- \emph{плато}, --- а гальмування повертає в спокій. Це тригер Шмітта: у нейрона два стійкі стани й гістерезис між ними. Так поводяться, наприклад, \nt{ACET}-нейрони, що керують м'язами, і вмикається ця властивість широкомовним каналом: плато з'являється, коли до нейрона доходить серотонін~\cite{Hounsgaard1988}. Та сама клітина із \nt{SERO} --- засувка, без нього --- звичайний інтегратор.

\subsection{Інтегратори й резонатори}

Теорія поділяє збудливі клітини на два класи~\cite{Izhikevich2007}. \emph{Інтегратори} схожі на генератор, керований напругою, який починає з нуля: трохи вище порога нейрон працює як завгодно повільно, і частота плавно зростає зі входом. Такий нейрон додає все, що надійшло, і добре передає частотою величину входу. \emph{Резонатори} схожі на генератор із мінімальною частотою: повільно працювати вони не вміють, за порогового входу одразу видають імпульси зі скінченною частотою і віддають перевагу входу своєї частоти. Перші --- природні прилади частотного коду розділу~\ref{sec:code}, другі --- часового.

\begin{keyblock}
\begin{proposition}[Один нейрон --- багато приладів]
\label{prop:eetypes}
За електричною функцією нейрони працюють як інтегратори з витоком і без нього, детектори збігів, фільтри верхніх частот і смугові фільтри, перетворювачі напруга---частота, випрямлячі, дільники, генератори, лінії затримки й тригери. Який прилад вийде з даної клітини, задають її іонні канали й широкомовні канали, що їх підлаштовують. Самі режими виміряно; те, наскільки мозок використовує таке переналаштування для обчислень, --- здебільшого гіпотеза.
\end{proposition}
\end{keyblock}

Для інженера це схоже на програмовану логічну матрицю: залізо одне, а прилад задає конфігурація. Тільки конфігурацію тут змінюють не під час прошивання, а на ходу, --- повільними широкомовними каналами, про які йшлося в розділах~\ref{sec:serocomputer}--\ref{sec:acet}.

\section{Задачі}

\begin{exercise}[\lvlA]
\label{ex:18:1}
\emph{Допуск інтегратора без витоку.} Нейрони з $\tau=0.1$~с тримають положення ока за~\eqref{eq:perfectint} протягом $T=1$~с із похибкою не більше $5\%$ в обидва боки. (а)~Знайдіть допустимий діапазон $w$. (б)~$w$ --- сума ваг $K$ синапсів, кожен із незалежною похибкою $10\%$. За якого $K$ розкид суми вкладається в допуск?
\end{exercise}

\begin{exercise}[\lvlB]
\label{ex:18:2}
\emph{Резонатор як контур.} Повільний струм підрозділу~\ref{sec:resonator} опишемо змінною $W$: $C\,\dd V/\dd t=-g_LV-g_wW+I$, $\tau_w\,\dd W/\dd t=V-W$. Покажіть, що це $RC$-коло з паралельною гілкою $R_w=1/g_w$, $L_w=\tau_w/g_w$, і за $C/g_L=10\ms$, $\tau_w=100\ms$, $\alpha=g_w/g_L=4$ знайдіть частоту резонансу й виграш $|Z|$ у ньому над $|Z(0)|$.
\end{exercise}

\begin{exercise}[\lvlB]
\label{ex:18:3}
\emph{Як інтегратор починає працювати.} (а)~Нейрон $\tau\,\dd V/\dd t=-V+\mu$, $\tau=10\ms$, з порогом $\theta$ і скиданням у нуль. За якого $\mu/\theta-1$ частота дорівнює $1$~Гц? (б)~Яку частоту дає модель $\tau\,\dd v/\dd t=v^2+\mu$ (імпульс --- відхід $v$ у $+\infty$, скидання в $-\infty$) за $\mu=0.01$?
\end{exercise}

\begin{exercise}[\lvlB]
\label{ex:18:4}
\emph{Моделювання: інтегратор і резонатор.} Модель задачі~\ref{ex:18:2} з $g_L=1$, $\tau_m=10\ms$, $\tau_w=100\ms$, порогом $1$ і скиданням $V\to0$ ($W$ не змінюється) отримує $\mu=1.5\sin2\pi ft$, $f=1$--$40$~Гц. У якій смузі частот нейрон видає імпульси за $\alpha=0$ і за $\alpha=4$? Порівняйте з умовою $1.5\,|Z(f)|>1$.
\end{exercise}

\begin{exercise}[\lvlB]
\label{ex:18:5}
\emph{Засувка з однієї клітини.} Нейрон зі струмом плато: $\tau\,\dd r/\dd t=-r+F(wr+I)$, $F(x)=\min(1,\max(0,x))$, $\tau=10\ms$, $w=1.5$, фон $I=-0.2$. (а)~Якої найменшої тривалості імпульс $I=0.3$ переводить нейрон з $r=0$ у плато? (б)~Якої найменшої амплітуди $B$ довгий імпульс $I=-0.2-B$ повертає його в спокій?
\end{exercise}

\begin{exercise}[\lvlB]
\label{ex:18:6}
\emph{Самоналаштування інтегратора.} Під час фіксації погляду вхід $I=0$, датчик ковзання дає швидкість відходу $e=\dd r/\dd t$, і $\dd w/\dd t=-\eta\,e\,r$. За $w=0.95$, $\tau=0.1$~с, $\langle r^2\rangle=1$ знайдіть $\eta$, за якого за $60$~с фіксацій $|1-w|$ увійде в допуск задачі~\ref{ex:18:1}. Що зламається, якщо вчитися й під час поворотів ока?
\end{exercise}

\begin{exercise}[\lvlC]
\label{ex:18:7}
\emph{Навіщо переналаштовувати прилад?} Широкомовний канал перемикає $\alpha$ моделі задачі~\ref{ex:18:2} між $0$ і $4$. У скільки разів резонатор виграє в інтегратора у відношенні сигнал/шум для слабкого синуса в білому шумі струму за $11$~Гц і за $1$~Гц? Придумайте задачу, де вигідне саме перемикання режиму.
\end{exercise}

\chapter{Нейрони за кількістю дендритів}
\chaptermark{Кількість дендритів}
\label{sec:dendrites}
\begin{chaptergoals}
\item чому ємність мембрани робить нейрон із небагатьма входами швидким, а з багатьма --- повільним;
\item як датчики, буфери й детектори з двома дендритами тримають окремий сигнал швидким і точним;
\item чому змішувачі з близько чотирма входами дають великому суматору зберегти багато відповідностей;
\item як великі дерева додають і класифікують, а дендрити-виходи роблять із клітини кілька детекторів.
\end{chaptergoals}

\section{Кількість входів як параметр схеми}

У розділі~\ref{sec:neurontypes} ми сортували нейрони за тим, що викидає їхній вихід, а в розділі~\ref{sec:eetypes} --- за тим, яким приладом вони працюють. Є й третя ознака, яку інженер помітить одразу: \emph{скільки в нейрона входів}. Входи збирають дендрити --- гілки, на яких сідають чужі аксони (розділ~\ref{sec:dofa}, підрозділ~\ref{sec:weightstore}). Одні нейрони не мають дендритів зовсім, інші мають один чи два, ще інші --- дерево, що збирає сотню тисяч входів. Для електронника це навантажувальна здатність за входом, fan-in, і вона майже завжди відповідає задачі.

Чому кількість і розмір дендритів такі важливі, видно з простої оцінки. Мембрана --- конденсатор із питомою ємністю $c_m\approx1$~мкФ/см$^2$, і що більша площа $A$ нейрона разом із його дендритами, то більший заряд треба принести, щоб підняти напругу на $\Delta V$. Час, за який вхідний струм $I$ це зробить, якщо не враховувати витоку,
\begin{empheq}[box=\keybox]{equation}
t \;\approx\; \frac{c_m\,A\,\Delta V}{I} .
\label{eq:dendriteload}
\end{empheq}
Маленький нейрон із кількома короткими дендритами має площу порядку $2\cdot10^{-5}$~см$^2$ і ємність близько $20$~пФ. Один великий синапс зі струмом у кілька наноамперів піднімає його на $20\mV$ менш ніж за десяту частку мілісекунди. У нейрона з великим деревом площа разів у п'ятнадцять більша, а ємність близько $300$~пФ; звичайний синапс зі струмом близько $0.1$~нА зарядив би його за $60\ms$ --- набагато довше за сталу часу витоку, тобто не зарядив би ніколи. Такому нейронові потрібні десятки одночасних входів. Числа тут --- порядки величин, але висновок від них не залежить: \emph{мало входів --- швидко й точно, багато входів --- повільно й за сумою}.

\begin{figure}[t]
\centering
\begin{tikzpicture}[>=Stealth,
  soma/.style={circle,draw,very thick,fill=blue!6,minimum size=0.55cm,inner sep=0pt},
  ax/.style={->,very thick,red!70!black},
  den/.style={very thick,blue!60!black}]
% 0 dendrites: sensor on a line
\begin{scope}[xshift=0cm]
\draw[den,decorate,decoration={snake,amplitude=1pt,segment length=4pt}] (-0.9,0) -- (0,0);
\node[font=\scriptsize] at (-0.9,0.3) {датчик};
\draw[ax] (0,0) -- (1.0,0);
\draw[thick] (0.1,0) -- (0.1,0.45);
\node[soma,minimum size=0.4cm] at (0.1,0.65) {};
\node[font=\scriptsize,align=center] at (0.05,-0.75) {$0$\\лінія з датчиком};
\end{scope}
% 1 dendrite
\begin{scope}[xshift=3.3cm]
\draw[den] (-0.9,0) -- (-0.28,0);
\node[soma] at (0,0) {};
\draw[ax] (0.28,0) -- (0.9,0);
\node[font=\scriptsize,align=center] at (0,-0.75) {$1$\\буфер, повторювач};
\end{scope}
% 2 dendrites: one per ear
\begin{scope}[xshift=6.2cm]
\draw[den] (-0.28,0) -- (-0.9,0); \draw[den] (0.28,0) -- (0.9,0);
\node[soma] at (0,0) {};
\draw[ax] (0,-0.28) -- (0,-0.6);
\node[font=\scriptsize] at (-0.9,0.25) {ліве}; \node[font=\scriptsize] at (0.9,0.25) {праве};
\node[font=\scriptsize,align=center] at (0,-1.0) {$2$\\«І» в часі};
\end{scope}
% ~4 short dendrites: mixer
\begin{scope}[xshift=9.0cm]
\foreach \a in {45,135,225,315}{\draw[den] (\a:0.28) -- (\a:0.7); \fill[blue!60!black] (\a:0.72) circle (0.06);}
\node[soma] at (0,0) {};
\draw[ax] (0.28,0) -- (1.0,0);
\node[font=\scriptsize,align=center] at (0,-1.0) {$\sim4$\\змішувач};
\end{scope}
% many: summing tree
\begin{scope}[xshift=12.0cm]
\foreach \a in {60,90,120}{\draw[den] (0,0.28) -- ++(\a:0.55) coordinate (b); \foreach \d in {-20,20}{\draw[den] (b) -- ++(\a+\d:0.35);}}
\foreach \a in {-120,-60}{\draw[den] (0,-0.28) -- ++(\a:0.45);}
\node[soma] at (0,0) {};
\draw[ax] (0.28,0) -- (1.0,0);
\node[font=\scriptsize,align=center] at (0,-1.0) {$10^4$--$10^5$\\суматор};
\end{scope}
\end{tikzpicture}
\caption{П'ять класів за кількістю входів. Синім --- дендрити (входи), червоним --- аксон (вихід). Що менше входів, то швидше й точніше нейрон передає окремий сигнал; що більше, то краще він додає й класифікує. Це схема, а не малюнок форми клітин.}
\label{fig:dendriteclasses}
\end{figure}

\section{Нуль дендритів: датчик на кінці лінії}

Чутливі нейрони дотику, болю й розтягу м'язів (розділи~\ref{sec:sensors}, \ref{sec:pain}, \ref{sec:muscles}) не мають на тілі клітини жодного дендрита. Аксон виходить із тіла однією гілочкою і відразу ділиться надвоє: одна гілка йде до шкіри чи м'яза, і її кінець сам слугує датчиком; друга йде в спинний мозок. Імпульс біжить від датчика просто в спинний мозок, оминаючи тіло клітини. Для інженера це лінія зв'язку, у якої датчик стоїть на дальньому кінці, а тіло клітини висить збоку, як блок живлення: воно живить лінію, але в передаванні не бере участі. Вигода --- швидкість і надійність: на шляху немає жодного місця, де сигнал треба було б підсумовувати й знову вирішувати, чи передавати його.

Датчики світла й звуку --- ще крайніший випадок: це не нейрони-збирачі, а самі перетворювачі, у яких входом слугує не синапс, а білок-рецептор (рис.~\ref{fig:transducer}).

\section{Один дендрит: буфер}

Нейрон з одним дендритом і одним аксоном приймає одну вхідну лінію й передає її далі. Так улаштовані нюхові нейрони: єдиний дендрит виходить до поверхні носа й несе єдиний тип рецептора~\cite{Mombaerts2004}; так улаштовані передавальні клітини сітківки між датчиками світла й вихідними нейронами ока; так улаштовані нейрони, що зв'язують датчики вуха з мозком. Інженер назвав би їх буфером або повторювачем: вони не обчислюють, а доставляють один сигнал, іноді змінюючи його форму --- наприклад, перетворюючи плавну напругу датчика на імпульси, як у розділі~\ref{sec:sensors}.

\section{Два дендрити: «І» в часі}

Детектори збігів, що визначають напрямок на звук (рис.~\ref{fig:itd}), у ссавців часто мають рівно два дендрити, спрямовані в протилежні боки: на один надходять входи від лівого вуха, на другий --- від правого~\cite{Grothe2010}. Навіщо розносити входи? Згадаймо формулу~\eqref{eq:shunt}: коли на одній ділянці мембрани відкрито багато збуджувальних каналів, напруга там наближається до рівноважного потенціалу, і кожен наступний вхід додає дедалі менше. Тому багато входів від одного вуха, зібраних на одному дендриті, насичують його і не можуть самотужки довести нейрон до порога. А по одному входу з кожного дендрита додаються майже лінійно. Два дендрити перетворюють нейрон на суворіше «І»: спрацьовуй, лише якщо сигнали надійшли з обох боків і разом. На моделях показано, що таке рознесення справді поліпшує детекцію збігів~\cite{AgmonSnir1998}.

\section{Кілька коротких дендритів: змішувач і ретранслятор}

\emph{Змішувачі.} У колі навчання рухів із розділу~\ref{sec:motorlearning} близько сімдесяти мільярдів маленьких \nt{GLUT}-нейронів розширюють вхід в дуже довгий вектор. Кожен із них має лише близько чотирьох коротких дендритів, і кожен закінчується «пазуром», який охоплює закінчення одного входу. Кожен змішувач, отже, бачить лише близько чотирьох входів із багатьох і працює як маленький елемент «І» над своєю четвіркою. З $M$ вхідних ліній можна вибрати $\binom{M}{4}$ різних четвірок: за $M=100$ це майже чотири мільйони. Навіщо стільки? Пороговий елемент з $n$ входами може правильно поділити на два класи не більше ніж
\begin{empheq}[box=\keybox]{equation}
P_{\max} \;\approx\; 2n
\label{eq:cover}
\end{empheq}
випадкових образів~\cite{Cover1965}. Вихідний нейрон того самого кола отримує близько $10^5$ входів від змішувачів, і за~\eqref{eq:cover} верхня межа для нього --- порядку $2\cdot10^5$ різних відповідностей. Це межа для ідеального елемента: якщо всі ваги додатні, як у збуджувальних синапсів, і потрібен запас на шум, оцінка для цього самого нейрона --- близько $4\cdot10^4$ відповідностей~\cite{Brunel2004}. Змішувачі з малою кількістю входів потрібні саме для того, щоб великий суматор мав багато різних, майже незалежних входів~\cite{Marr1969,Albus1971}.

\emph{Ретранслятори.} На шляху від вуха до детекторів напрямку є нейрони з небагатьма короткими дендритами, які отримують лише кілька величезних синапсів, а деякі --- по суті один. За~\eqref{eq:dendriteload} це саме те, що треба: мала ємність і великий струм, і нейрон спрацьовує через частки мілісекунди після кожного вхідного імпульсу, майже не втрачаючи точності часу. Один із таких ретрансляторів отримує \nt{GLUT}, а видає \nt{GLYC}: це інвертор із точністю в десятки мікросекунд, що постачає швидке гальмування детекторам напрямку~\cite{BorstSoria2012}.

\section{Тисячі входів: суматор і класифікатор}

На другому кінці шкали --- \nt{GLUT}-нейрони кори з десятком тисяч входів і вихідні \nt{GABA}-нейрони кола навчання рухів із сотнею тисяч. Такий нейрон за~\eqref{eq:dendriteload} повільний і не може відгукнутися на окремий вхід: він додає. Це інтегрувальний нейрон розділу~\ref{sec:glutcomputer} і той суматор, з якого будують класифікатори. Велике дерево додає й нове: його гілки працюють як окремі підсхеми з власними порогами, тож один нейрон поводиться як маленька багатошарова мережа~\cite{Beniaguev2021,Gidon2020}.

\section{Дендрити, які є ще й виходом}

Є нейрони, що не мають аксона зовсім, а дендрити слугують і входом, і виходом: вони отримують сигнали й самі викидають медіатор. Найвивченіший приклад --- \nt{GABA}-нейрони сітківки, що беруть участь у визначенні напрямку руху (розділ~\ref{sec:gaba}, рис.~\ref{fig:direction}). У такого нейрона кожна гілка сама по собі відповідає на рух від тіла клітини до свого кінчика й видає гальмування з цього кінчика~\cite{Euler2002}. Один нейрон --- це кілька незалежних детекторів напрямку, по одному на гілку. Для інженера це не один елемент, а мікросхема з кількома каналами в одному корпусі.

\section{Підсумок}

\begin{table}[t]
\centering
\footnotesize
\setlength{\tabcolsep}{3pt}
\renewcommand{\arraystretch}{1.15}
\begin{tabular}{>{\raggedright}p{1.9cm}>{\raggedright}p{3.4cm}>{\raggedright}p{2.6cm}>{\raggedright}p{1.8cm}>{\raggedright\arraybackslash}p{3.8cm}}
\toprule
Дендритів & Приклад & Прилад & Входів & Що відомо \\
\midrule
$0$ & датчики дотику, болю, розтягу & лінія з датчиком на кінці & --- & встановлено \\
$1$ & нюхові нейрони, передавальні клітини сітківки, нейрони вуха & буфер, повторювач & $1$--небагато & встановлено \\
$2$ & детектори напрямку на звук & «І» в часі & два пучки & будову встановлено; виграш від рознесення входів показано на моделях \\
$\sim4$ & змішувачі кола навчання рухів & маленький елемент «І» & $\sim4$ & встановлено; роль розширення входу --- добре обґрунтована гіпотеза \\
небагато, короткі & ретранслятори на шляху від вуха & ретранслятор, інвертор & $1$--кілька величезних & встановлено \\
тисячі & \nt{GLUT}-нейрони кори, вихідні нейрони навчання рухів & суматор, класифікатор & $10^4$--$10^5$ & встановлено; гілки як підсхеми виміряно на окремих прикладах \\
дендрити-виходи & \nt{GABA}-нейрони напрямку в сітківці & багатоканальна мікросхема & багато & виміряно \\
\bottomrule
\end{tabular}
\caption{Нейрони за кількістю дендритів.}
\label{tab:dendrites}
\end{table}

\begin{keyblock}
\begin{proposition}[Кількість входів іде за задачею]
\label{prop:dendrites}
Де потрібні швидкість і точність окремого сигналу --- на датчиках, ретрансляторах і детекторах збігів, --- нейрони мають мало дендритів, мало входів і великі синапси: мала ємність~\eqref{eq:dendriteload} заряджається швидко. Де треба додавати й класифікувати, нейрони мають великі дерева з тисячами входів. Будову класів виміряно; обчислювальне пояснення добре обґрунтоване, але для багатьох класів лишається гіпотезою.
\end{proposition}
\end{keyblock}

Таблиця~\ref{tab:dendrites} доповнює два попередні сортування: за медіатором (таблиця~\ref{tab:types}) і за приладом (таблиця~\ref{tab:eetypes}). Усі три дивляться на той самий елемент із різних боків: що він видає, що він обчислює і скільки він слухає.

\section{Задачі}

\begin{exercise}[\lvlA]
\label{ex:19:1}
\emph{Скільки входів потрібно великому нейронові.} Нейрон із $C=300$~пФ і $\tau_m=20\ms$ отримує синапси --- джерела струму по $0.1$~нА, поріг на $20\mV$ вищий за спокій. (а)~Скільки синапсів має працювати одночасно, щоб дійти до порога взагалі; за $10\ms$; за $5\ms$? (б)~За який час дійде до порога нейрон із $C=20$~пФ від синапсу в $4$~нА?
\end{exercise}

\begin{exercise}[\lvlA]
\label{ex:19:2}
\emph{Чому чотири.} Змішувач --- «І» над $k$ лініями з $M=100$, в образі активна половина ліній, змішувачів $7\cdot10^{10}$. (а)~Скільки з них відповідає на образ за $k=2$; $4$; $40$? (б)~У скільки разів змішувачі збільшують кількість відповідностей~\eqref{eq:cover} вихідного нейрона з $10^5$ входами проти $M=100$ прямих ліній?
\end{exercise}

\begin{exercise}[\lvlB]
\label{ex:19:3}
\emph{Звідки береться $2n$.} Гіперплощина через початок координат ділить $P$ точок загального положення в $n$-вимірному просторі на два класи $C(P,n)=2\sum_{k=0}^{n-1}\binom{P-1}{k}$ способами. (а)~Доведіть, що за $P=2n$ роздільна рівно половина з $2^P$ розбиттів. (б)~Яка їхня частка роздільна за $n=100$ і $P=180$; $220$?
\end{exercise}

\begin{exercise}[\lvlB]
\label{ex:19:4}
\emph{Два дендрити як суворе «І».} Дендрит --- ділянка з витоком $g_L$, кожен вхід відкриває на ній провідність $g_0=0.5\,g_L$ з потенціалом $E_E$, тіло бачить $\kappa(V_1+V_2)$, поріг $\theta=1.1\,\kappa E_E$. Чому жодна кількість входів на одному дендриті не запускає нейрона і скільки входів потрібно на кожному?
\end{exercise}

\begin{exercise}[\lvlB]
\label{ex:19:5}
\emph{Проєктування ретранслятора.} Тремтіння імпульсу $\sigma_t=\sigma_V/\dot V$, де $\dot V$ --- нахил напруги біля порога. Потрібно $\sigma_t\le20$~мкс за шуму $\sigma_V=1\mV$; витоком знехтуйте. Скільки синхронних синапсів по $0.1$~нА треба для цього нейронові з $C=300$~пФ і яке тремтіння нейрона з $C=20$~пФ та одним синапсом у $4$~нА?
\end{exercise}

\begin{exercise}[\lvlB]
\label{ex:19:6}
\emph{Моделювання: заряд довгого дендрита.} Змоделюйте пасивний кабель $\tau_m\partial_tV=\lambda^2\partial_x^2V-V+i/g$ довжини $L=\ell/\lambda$ із закритими кінцями ($40$ ділянок, Ейлер). У дальній кінець подано короткий імпульс струму: коли й до якої частки напруги від того самого заряду, розмазаного по всьому дендриту, піднімається ближній кінець за $L=0.2$; $1$; $2$?
\end{exercise}

\begin{exercise}[\lvlC]
\label{ex:19:7}
\emph{Дослідження: оптимальна кількість входів.} Ємність $C=C_0+nc_s$, одночасно працюють $m$ входів зі струмом $I_s$, нейрон запам'ятовує $P\approx2n$ відповідностей~\eqref{eq:cover}. Як час відгуку залежить від $P$ за сталого $m$ і за сталої частки $m=\varphi n$? Запропонуйте критерій вартості й знайдіть оптимальне $n$ для ретранслятора й класифікатора.
\end{exercise}

\chapter{Що роблять нейрони уві сні}
\chaptermark{Сон}
\label{sec:sleep}

\section{Сон --- не вимкнення, а інший режим}

Інженер, побачивши вимкнений комп'ютер, скаже: він нічого не робить. Зі сплячим мозком так не виходить. Нейрони уві сні не мовчать: у глибокому сні кора видає імпульси, у швидкому сні --- майже стільки ж, скільки вдень. Змінюється не те, чи \emph{працюють} нейрони, а те, \emph{як} вони працюють. Сон --- це набір режимів машини, і перемикають їх ті самі широкомовні канали, що й удень.

Для кожного режиму можна виписати «стан регістрів» (таблиця~\ref{tab:sleepmodes}). Удень \nt{ACET}, \nt{NORA} і \nt{SERO} працюють, датчики підключено, м'язи слухаються. У повільному, глибокому сні всі три канали затихають, а вхід від органів чуття ослаблено~\cite{Hasselmo1999}: викид ацетилхоліну в корі падає, а \nt{NORA}- і \nt{SERO}-нейрони видають імпульси рідше, ніж удень~\cite{Marrosu1995,AstonJonesBloom1981,McGintyHarper1976}. У швидкому сні картина дивна: \nt{ACET} знову піднімається до денного рівня~\cite{Marrosu1995}, а \nt{NORA} і \nt{SERO} замовкають майже повністю~\cite{AstonJonesBloom1981,McGintyHarper1976,JacobsAzmitia1992}. М'язи тіла у швидкому сні вимкнено: \nt{ACET}-нейрони, що ними керують, сильно гальмуються через \nt{GABA} і \nt{GLYC}~\cite{BrooksPeever2012} --- тією самою забороною з розділу~\ref{sec:gaba}, тільки накладеною на весь вихід одразу.

\begin{table}[t]
\centering
\footnotesize
\setlength{\tabcolsep}{4pt}
\renewcommand{\arraystretch}{1.15}
\begin{tabular}{>{\raggedright}p{3.1cm}>{\raggedright}p{2.9cm}>{\raggedright}p{3.2cm}>{\raggedright\arraybackslash}p{3.4cm}}
\toprule
& неспання & повільний сон & швидкий сон \\
\midrule
\nt{ACET} (запис, розд.~\ref{sec:acet}) & високий & низький & високий \\
\nt{NORA} (підсилення, розд.~\ref{sec:nora}) & середній, спалахи & низький & майже мовчить \\
\nt{SERO} (розд.~\ref{sec:serocomputer}) & середній & низький & майже мовчить \\
вхід від датчиків & підключено & ослаблено & ослаблено \\
вихід на м'язи & підключено & ослаблено & вимкнено гальмуванням \\
активність кори & рівна & повільні гойдалки: робота --- тиша & рівна, як удень \\
\bottomrule
\end{tabular}
\caption{Стан регістрів машини в трьох режимах. Усі рядки виміряно; рівні \nt{ACET}, \nt{NORA} і \nt{SERO} --- прямими записами за стадіями сну~\cite{Marrosu1995,AstonJonesBloom1981,McGintyHarper1976}.}
\label{tab:sleepmodes}
\end{table}

\section{Коли спати: реле з гістерезисом}

Що вирішує, коли машині перемикатися в сон? Добре працює проста схема з двох частин~\cite{Borbely1982}. Перша частина --- \emph{тиск сну} $S$, який накопичується, поки ми не спимо, і скидається, поки спимо. Це конденсатор, що заряджається вдень і розряджається вночі:
\begin{empheq}[box=\keybox]{equation}
\frac{\dd S}{\dd t} \;=\; \frac{1-S}{\tau_r}\ \ \text{(неспання)},\qquad
\frac{\dd S}{\dd t} \;=\; -\frac{S}{\tau_d}\ \ \text{(сон)} .
\label{eq:sleeppressure}
\end{empheq}
Друга частина --- два пороги: машина засинає, коли $S$ доходить до верхнього порога $H(t)$, і прокидається, коли $S$ опускається до нижнього $L(t)$. Обидва пороги плавно гойдаються в такт добовому ритму з розділу~\ref{sec:chemoutput}. Вийшло реле з гістерезисом --- тригер Шмітта з підрозділу~\ref{sec:bistable}, --- а разом із конденсатором воно дає релаксаційний генератор, що підстроюється за фазою добовим ритмом~\eqref{eq:pll}.

\begin{figure}[t]
\centering
\begin{tikzpicture}[>=Stealth,x=0.16cm,y=4.2cm]
\foreach \a/\b in {0/4.3,21.2/29.3,45.4/53.4,69.5/72}{\fill[blue!8] (\a,0) rectangle (\b,1);}
\draw[->,gray!70!black] (0,0) -- (75,0) node[right,font=\small,black] {$t$, год};
\draw[->,gray!70!black] (0,0) -- (0,1.08) node[left,font=\small,black] {$S$};
\foreach \t in {0,24,48,72}{\draw[gray!60!black] (\t,-0.02) -- (\t,0.02); \node[below,font=\scriptsize] at (\t,0) {\t};}
\draw[thick,densely dashed,red!70!black] plot file {figures/sleep_H.dat};
\draw[thick,densely dashed,green!45!black] plot file {figures/sleep_L.dat};
\draw[very thick,blue!60!black] plot file {figures/sleep_S.dat};
\node[font=\scriptsize,red!70!black,anchor=west] at (73,0.72) {$H$: заснути};
\node[font=\scriptsize,green!45!black,anchor=west] at (73,0.24) {$L$: прокинутися};
\node[font=\scriptsize,blue!60!black] at (25.25,0.93) {сон};
\node[font=\scriptsize,blue!60!black] at (49.4,0.93) {сон};
\end{tikzpicture}
\caption{Тиск сну $S$~\eqref{eq:sleeppressure} за $\tau_r=18.2$~год, $\tau_d=4.2$~год. Удень $S$ заряджається до верхнього порога $H$, уночі (зафарбовано) розряджається до нижнього $L$; обидва пороги гойдаються разом із добовим ритмом. Машина сама приходить до режиму близько $16$ годин неспання і $8$ годин сну.}
\label{fig:twoprocess}
\end{figure}

У нашій моделі зі звичайними значеннями $\tau_r=18.2$~год і $\tau_d=4.2$~год машина сама виходить на $16.1$ години неспання і $8.0$ години сну (рис.~\ref{fig:twoprocess}). Модель пояснює й те, що здається дивним: після сорока годин без сну тиск доходить до $0.90$, але відновлювальний сон у моделі триває лише $8.8$ години, а не на добу довше за звичайний. Розряд експоненційний, і високий заряд спадає швидко; «борг» повертається не тривалістю, а глибиною сну.

Що фізично грає роль $S$? Сильний кандидат --- аденозин, «лічильник навантаження» з додатка~\ref{sec:othermediators}: його концентрація зростає за час неспання і падає уві сні~\cite{PorkkaHeiskanen1997,Dunwiddie2001}. Що тиск сну --- це саме аденозин і нічого більше, --- гіпотеза.

\section{Повільний сон: уся мережа як релаксаційний генератор}

У глибокому сні нейрони кори працюють почергово всією групою: рідше ніж раз на секунду вони разом умикаються на кілька сотень мілісекунд (\emph{стан роботи}) і разом замовкають (\emph{стан тиші}): частота цих гойдалок нижча за $1$~Гц, під наркозом найчастіше $0.3$--$0.4$~Гц~\cite{Steriade1993}. Ці повільні гойдалки ми можемо зібрати з деталей, які вже маємо. Візьмімо групу \nt{GLUT}-нейронів із сильним зворотним зв'язком на себе --- пам'ять із розділу~\ref{sec:glutcomputer}, яка має два стійкі стани, --- і додамо повільну втому, як у генераторі ритму розділу~\ref{sec:muscles}:
\begin{empheq}[box=\keybox]{equation}
\tau\,\frac{\dd r}{\dd t} = -r + F\bigl(w\,r + I - a\bigr),
\qquad
\tau_a\,\frac{\dd a}{\dd t} = -a + b\,r .
\label{eq:updown}
\end{empheq}
Група вмикається, працює й утомлюється; втома $a$ з'їдає верхній стан, і група падає в тишу; у тиші втома минає, нижній стан зникає, і група знову вмикається. Виходить той самий релаксаційний генератор, що й тиск сну, тільки приблизно у вісімдесят тисяч разів швидший.

\begin{figure}[t]
\centering
\begin{tikzpicture}[>=Stealth,x=2.9cm,y=1.0cm]
\begin{scope}[yshift=1.9cm]
\draw[->,gray!70!black] (0,0) -- (4.1,0);
\draw[->,gray!70!black] (0,0) -- (0,1.25) node[left,font=\small,black] {$r$};
\draw[thick,blue!60!black] plot file {figures/updown_sleep.dat};
\node[font=\scriptsize,anchor=west] at (4.15,0.5) {сон: $b=5$};
\end{scope}
\draw[->,gray!70!black] (0,0) -- (4.1,0) node[right,font=\small,black] {$t$, с};
\draw[->,gray!70!black] (0,0) -- (0,1.25) node[left,font=\small,black] {$r$};
\draw[thick,red!70!black] plot file {figures/updown_wake.dat};
\node[font=\scriptsize,anchor=west] at (4.15,0.5) {день: $b=1$};
\foreach \t in {0,1,2,3,4}{\draw[gray!60!black] (\t,-0.05) -- (\t,0.05); \node[below,font=\scriptsize] at (\t,0) {\t};}
\end{tikzpicture}
\caption{Та сама група~\eqref{eq:updown} у двох режимах: $w=5$, $I=2$, $\theta=2.5$, $\tau=10\ms$, $\tau_a=0.3$~с. Із сильною втомою ($b=5$, як у повільному сні) група гойдається між роботою й тишею з періодом $1.1$~с (значення моделі; у корі період більший за секунду, під наркозом --- близько $3$~с). Послабмо втому ($b=1$) --- так діє ацетилхолін --- і та сама група працює рівно.}
\label{fig:updown}
\end{figure}

А що переводить мережу з гойдалок у рівну денну роботу? Ацетилхолін пригнічує в нейронах повільні струми, що створюють утому~\cite{McCormick1992}. У нашій моделі за сильної втоми, $b=5$, група гойдається з періодом $1.1$~с --- швидше за виміряні гойдалки, але того самого порядку --- і працює близько $35\%$ часу, якщо усереднювати за цілим числом періодів; за слабкої, $b=1$, та сама група працює рівно (рис.~\ref{fig:updown}). Один регістр --- рівень \nt{ACET} --- перемикає генератор на підсилювач. Поверх повільних гойдалок уві сні йдуть і швидші спалахи ритму $7$--$15$ коливань на секунду; їх породжує петля \nt{GLUT}--\nt{GABA} між корою й проміжним релейним вузлом~\cite{SteriadeMcCormickSejnowski1993}, родичка ритму з розділу~\ref{sec:gaba}.

\section{Повтори: перемотування дня}

У розділі~\ref{sec:acet} ми бачили, що в повільному сні нейрони, які працювали вдень разом, знову спрацьовують разом і в тому самому порядку. Додамо одну інженерну деталь: повтори йдуть \emph{із перемотуванням}. Послідовність, яка вдень займала секунди, програється уві сні за десяті частки секунди~\cite{Nadasdy1999}: у гіпокампі приблизно вдвадцятеро швидше~\cite{LeeWilson2002}, у корі --- у шість-сім разів~\cite{Euston2007}. І програється не будь-коли: короткі спалахи повторів в одній ділянці підлаштовуються під гойдалки роботи й тиші та під швидкі спалахи ритму в корі. Якщо штучно посилити цю узгодженість, пам'ять поліпшується~\cite{Maingret2016} --- це вже причиновий зв'язок, а не збіг.

Навіщо машині перемотувати день? Інженерові знайома трудність, яку називають \emph{катастрофічним забуванням}: мережа, яка вчить нове підряд, одне за одним, псує старе. Рятує \emph{перемішування}: учити нове впереміш зі старим, невеликими порціями, багато разів. Гіпотеза про дві системи навчання~\cite{McClelland1995} полягає в тому, що швидка пам'ять записує день цілком, а уві сні потроху, впереміш і багаторазово програє його повільній корі. Це та сама пара «швидкий буфер --- повільне сховище», що в розділі~\ref{sec:acet}, і та сама пара швидкого й повільного учнів, що в розділі~\ref{sec:motorlearning}. Повтори та їхню узгодженість виміряно; що їхнє призначення --- перемішане навчання повільної пам'яті, --- добре обґрунтована гіпотеза.

\section{Перенормування ваг}

Удень правила Гебба з розділу~\ref{sec:dofa} здебільшого посилюють зв'язки. Якщо лише посилювати, ваги ростуть, а разом із ними ростуть витрати енергії й падає чутливість: нейрон, у якого всі входи сильні, перестає їх розрізняти. За гіпотезою \emph{синаптичного гомеостазу}~\cite{TononiCirelli2014}, сон потрібен, зокрема, для того, щоб повернути ваги до робочого рівня: усі зв'язки слабшають в одну й ту саму кількість разів, тож їхні \emph{відношення} --- вивчене --- зберігаються. Для інженера це нормування, як у правилі~\eqref{eq:oja}, тільки зроблене не на ходу, а в окремий час.

Прямі вимірювання цьому не суперечать: у мишей площа контакту на синапсах кори після сну приблизно на $18\%$ менша, ніж після неспання, зменшення пропорційне розміру, а найбільші й найстійкіші контакти майже не змінюються~\cite{deVivo2017}. Масштабування ваг виміряно; що це головне завдання сну, --- гіпотеза.

\section{Швидкий сон: машина, відключена від входів і виходів}

У швидкому сні кора працює майже як удень, але вхід від органів чуття ослаблено, а вихід на м'язи вимкнено гальмуванням. Для інженера це знайомий режим: \emph{випробування на стенді}. Схема працює на повну потужність, але її входи замкнено на внутрішній генератор, а виходи від'єднано від виконавчих механізмів, щоб нічого не зламати. Сновидіння, за однією з гіпотез, --- саме такий внутрішній прогін моделі світу, зібраної за день; що саме мережа при цьому робить і чи вчиться, невідомо. Виміряно тут лише те, що записано в таблиці~\ref{tab:sleepmodes}: який канал увімкнено, який мовчить, і що вихід заблоковано.

\section{Обслуговування і місцевий сон}

Довго вважалося, що уві сні міжклітинні щілини розширюються і мозок швидше промивається від продуктів обміну речовин~\cite{Xie2013}. Недавні досліди, що вимірювали промивання іншим способом, дістали протилежне: уві сні воно повільніше~\cite{Miao2024}. Питання нині відкрите, і ми його залишимо відкритим.

Надійніший інший факт: сон --- властивість місцевих мереж, а не всієї машини одразу. Якщо тварині довго не давати спати, то окремі ділянки її кори, поки тварина не спить і рухається, ненадовго провалюються в тишу, як у повільному сні, і в ці моменти тварина частіше помиляється~\cite{Vyazovskiy2011}. Кожна мережа втомлюється сама й сама перемикається; загальний режим виходить тоді, коли широкомовні канали переводять усі мережі одразу.

\begin{keyblock}
\begin{proposition}[Сон як набір режимів]
\label{prop:sleep}
Уві сні нейрони не вимикаються: широкомовні канали переводять машину в інші режими. Коли спати, вирішує реле з гістерезисом~\eqref{eq:sleeppressure}, підстроєне добовим ритмом. У повільному сні мережа стає релаксаційним генератором~\eqref{eq:updown} і програє день із перемотуванням; у швидкому --- працює, відключена від входів і виходів. Режими, повтори й масштабування ваг виміряно; їхнє призначення --- перемішане навчання й перенормування --- добре обґрунтовані гіпотези.
\end{proposition}
\end{keyblock}

\section{Підсумок}

\begin{table}[t]
\centering
\footnotesize
\setlength{\tabcolsep}{4pt}
\renewcommand{\arraystretch}{1.15}
\begin{tabular}{>{\raggedright}p{3.2cm}>{\raggedright}p{4.4cm}>{\raggedright\arraybackslash}p{5.2cm}}
\toprule
Що відбувається & Як & Що відомо \\
\midrule
зміна режимів & регістри \nt{ACET}, \nt{NORA}, \nt{SERO} (табл.~\ref{tab:sleepmodes}) & виміряно \\
коли спати & конденсатор $S$ і реле з гістерезисом~\eqref{eq:sleeppressure} & модель добре описує досліди; аденозин як $S$ --- гіпотеза \\
повільні гойдалки & група зі зворотним зв'язком і втомою~\eqref{eq:updown} & гойдалки виміряно; модель --- найпростіша схема \\
перемотування дня & повтори в $6$--$20$ разів швидші, узгоджені з гойдалками & виміряно; посилення узгодженості поліпшує пам'ять \\
навіщо повтори & перемішане навчання повільної пам'яті & добре обґрунтована гіпотеза \\
перенормування ваг & масштабування зв'язків уві сні & зменшення контактів виміряно; головна роль сну --- гіпотеза \\
швидкий сон & робота з відключеними входами й виходом & режим виміряно; призначення невідоме \\
промивання & розширення міжклітинних щілин & суперечливі результати \\
місцевий сон & окремі ділянки провалюються в тишу & виміряно \\
\bottomrule
\end{tabular}
\caption{Що роблять нейрони уві сні.}
\label{tab:sleep}
\end{table}

Таблиця~\ref{tab:sleep} підсумовує розділ. Сон виявився не паузою в роботі машини, а її обслуговуванням на ходу: перемикання режимів тими самими широкомовними каналами, генератор із тих самих деталей, що й ритм кроків, перемотування дня для повільної пам'яті й перенормування ваг. Удень машина пише, уночі переписує й упорядковує записане.

\section{Задачі}

\begin{exercise}[\lvlA]
\label{ex:20:1}
\emph{Реле без добового ритму.} Нехай пороги сталі: $H=0.67$, $L=0.17$ (середні значення порогів моделі рис.~\ref{fig:twoprocess}). Виведіть із~\eqref{eq:sleeppressure} тривалості неспання й сну та період генератора за $\tau_r=18.2$~год, $\tau_d=4.2$~год. Чому модель із гойдальними порогами все ж виходить на $16.1+8.0=24$~год? До якого приладу розділу~\ref{sec:chemoutput} належить цей ефект?
\end{exercise}

\begin{exercise}[\lvlA]
\label{ex:20:2}
\emph{Борг сну.} Сон, розпочатий за тиску $S_0$, триває $t_s=\tau_d\ln(S_0/L)$, $\tau_d=4.2$~год, $L$ --- нижній поріг під час пробудження. (а)~Після $40$~год без сну $S_0=0.90$, а сон тривав $8.8$~год. Яким був $L$? Скільки тривав би сон за звичайного $L\approx0.092$? (б)~За ніч площа синаптичних контактів зменшується на $18\%$. На скільки мають зрости ваги за день?
\end{exercise}

\begin{exercise}[\lvlB]
\label{ex:20:3}
\emph{Проєктування порогів.} Доберіть сталі пороги $H$ і $L$, за яких реле~\eqref{eq:sleeppressure} з $\tau_r=18.2$~год і $\tau_d=4.2$~год дає рівно $16$~год неспання і $8$~год сну. Чим така машина гірша за машину з гойдальними порогами, якщо одного разу вночі її розбудили через $4$~години?
\end{exercise}

\begin{exercise}[\lvlB]
\label{ex:20:4}
\emph{Період повільних гойдалок.} У моделі~\eqref{eq:updown} $F(x)=1/\bigl(1+e^{-(x-\theta)/k}\bigr)$, $w=5$, $I=2$, $\theta=2.5$, $k=0.25$, $b=5$, $\tau\ll\tau_a=0.3$~с. (а)~Знайдіть точки повороту $a_{\text{в}}$ і $a_{\text{н}}$ кривої $r=F(wr+I-a)$, де зникають робота й тиша. (б)~Узявши $r\approx1$ у роботі й $r\approx0$ у тиші, знайдіть період і частку роботи.
\end{exercise}

\begin{exercise}[\lvlB]
\label{ex:20:5}
\emph{Коли гойдалки вимикаються.} У моделі задачі~\ref{ex:20:4} нерухома точка лежить на перетині кривої $a=wr+I-F^{-1}(r)$ з прямою $a=br$. Коливання йдуть, поки перетин на середній гілці, між точками повороту. За яких $b$ це так? Як, змінюючи $b$, \nt{ACET} перемикає генератор на підсилювач?
\end{exercise}

\begin{exercise}[\lvlB]
\label{ex:20:6}
\emph{Моделювання: борг чи фаза.} У \texttt{sleep\_models.py} візьміть перше пробудження після $48$~год і тримайте машину без сну $D=24$, $32$, $40$, $48$, $64$~год (\texttt{stay\_awake}, \texttt{days=8}). Скільки триває відновлювальний сон за кожного $D$? Що визначає його тривалість?
\end{exercise}

\begin{exercise}[\lvlB]
\label{ex:20:7}
\emph{Моделювання: забування й перемішування.} Лінійний нейрон $y=\mathbf w\cdot\mathbf x$, $n=40$, учиться за правилом $\mathbf w\leftarrow\mathbf w-0.5\,(y-y^*)\,\mathbf x$. Задачі A і B --- по $10$ образів, $x_j\sim\mathcal N(0,1/n)$, $y^*=\pm1$. Яка середньоквадратична похибка на A після $200$ епох A, а потім $200$ епох B? А після $200$ епох A і B впереміш?
\end{exercise}

\begin{exercise}[\lvlC]
\label{ex:20:8}
\emph{Дослідження: перемішування чи перенормування?} Нейрон задачі~\ref{ex:20:7} з порогом; дві нічні процедури: (i)~усі ваги множаться на $0.82$; (ii)~старі образи повторюються впереміш із новими. Що кожна робить із відношеннями ваг і з відповідями нейрона? Як розрізнити гіпотези за розмірами синапсів після сну?
\end{exercise}

\chapter{Машина з живих нейронів}
\chaptermark{Машина з живих нейронів}
\label{sec:livingmachine}

\section{Стенд із відкритою схемою}

У розділі~\ref{sec:debugging} ми налагоджували машину, в якої немає схеми: щуп чує тисячу нейронів із вісімдесяти мільярдів, а все інше доводиться вгадувати. Інженер у такому становищі зробив би просту річ --- зібрав би невеликий шматок схеми на макетній платі, де видно кожну деталь. З нейронами це теж можна зробити.

Нейрони кори розділяють і вирощують на чипі з ґраткою електродів --- від кількох десятків до десятків тисяч. За кілька тижнів клітини проростають відростками й з'єднуються в мережу з тисяч або сотень тисяч нейронів, здебільшого \nt{GLUT} і меншої частки \nt{GABA}, яка залежить від препарату: у культурах гіпокампа це близько $6\%$ нейронів~\cite{Benson1994}. Кожен електрод уміє і слухати, як щуп розділу~\ref{sec:debugging}, і подавати струм, як тестовий сигнал. Другий різновид стенда --- \emph{органоїди}: тривимірні грудочки нервової тканини розміром близько міліметра, вирощені зі стовбурових клітин людини~\cite{Lancaster2013}. На таких стендах живі мережі працюють місяцями; є майданчики, де досліди на органоїдах можна ставити віддалено, через мережу, понад сто днів поспіль~\cite{Jordan2024}. Цю галузь називають \emph{біологічними обчисленнями} або \emph{органоїдним інтелектом}~\cite{Smirnova2023}.

Стенд має дві важливі відмінності від мозку. Він крихітний: нейронів у ньому в сто тисяч разів менше. І в ньому немає широкомовних каналів: нейронів \nt{DOPA}, \nt{SERO}, \nt{NORA} і \nt{ACET} у культурі з кори немає. Це машина, яка має шину даних і керування потоком, але не має регістрів керування. Подивімося, на що вона здатна.

\section{Що мережа робить сама}

Якщо культуру залишити в спокої, вона не мовчить і не шумить рівномірно. Час від часу майже вся мережа спалахує разом --- \emph{залпом} --- і знову затихає; інтервали між залпами --- від секунди до хвилин, залежно від культури~\cite{Wagenaar2006}. Для інженера це знайома картина: підсилювач із сильним додатним зворотним зв'язком і без навантаження перетворюється на генератор.

Механізм ми вже знаємо. Сильні взаємні зв'язки \nt{GLUT}-нейронів запускають лавину, як у розділі~\ref{sec:gaba} при порушенні умов твердження~\ref{prop:ei}. Лавина швидко виснажує запас медіатора в синапсах~\eqref{eq:depression}, і мережа замовкає. Запас відновлюється зі сталою часу $\tau_{\text{відн}}$, і коли відновиться частка $x^*$, достатня для нової лавини, настає новий залп. Інтервал між залпами
\begin{empheq}[box=\keybox]{equation}
T_{\text{залп}} \;\approx\; \tau_{\text{відн}}\,\ln\frac{1}{1-x^*} .
\label{eq:burstinterval}
\end{empheq}
При $\tau_{\text{відн}}=0.8\,\text{с}$ із розділу~\ref{sec:code} і $x^*=0.9$ виходить близько двох секунд, при $x^*=0.99$ --- близько чотирьох. Це оцінка моделі з $\tau_{\text{відн}}$ книги, і вона потрапляє в короткий край виміряних інтервалів. Пряме вимірювання на мікрокультурах показує, що запас медіатора після залпу відновлюється за десятки секунд~\cite{CohenSegal2011}, тобто $\tau_{\text{відн}}$ там більша; з такою $\tau_{\text{відн}}$ формула~\eqref{eq:burstinterval} дає десятки секунд і хвилини, як у повільних культурах. Це генератор релаксаційного типу, родич генератора ритму з розділу~\ref{sec:muscles}: там його перезаряджала втома груп, тут --- запас медіатора.

Залпи --- це те, що мережа робить, коли їй нічого робити. У живому мозку схожа картина буває в глибокому сні (розділ~\ref{sec:sleep}) і в мозку, що розвивається, доки до нього не надійшли справжні входи. Ця схожість наводить на думку, але однаковість механізмів --- гіпотеза.

\section{Як навчати мережу без учителя-дофаміну}

Оскільки \nt{DOPA} в культурі немає, сигнал «краще чи гірше» доводиться подавати самим --- електродами. Досліди показують, що мережа на стенді справді змінює відповіді, і найцікавіше в них --- \emph{чим} її навчають.

\emph{Учитель, який замовкає.} Мережу стимулюють через один електрод раз на одну--три секунди, доки на іншому електроді не з'явиться потрібна відповідь через $50\pm10\ms$ після стимулу. Щойно вона з'явилася, стимуляцію вимикають на п'ять хвилин. Після кількох таких циклів потрібна відповідь з'являється на стимул одразу~\cite{Shahaf2001}. Винагорода тут --- тиша: стимул розхитує мережу, а припинення стимулу залишає її в тому стані, який дав правильну відповідь. Це спуск за градієнтом через випадкові відхилення з розділу~\ref{sec:dofa} (твердження~\ref{prop:gradient}), де роль помилки відіграє сам факт трусіння: трусять, доки не стане правильно.

\emph{Мережа, яка грає в теніс.} У досліді на стенді зі щільною ґраткою електродів близько мільйона нейронів підключили до комп'ютерної гри в теніс з однією ракеткою~\cite{Kagan2022}. Положення м'яча подавали струмом на «чутливі» електроди: де м'яч --- кодом місця, як далеко --- частотою. Активність у двох «рухових» ділянках рухала ракетку вгору або вниз. Правило навчання було незвичним. Якщо ракетка відбила м'яч, мережа отримувала короткий \emph{передбачуваний} стимул; якщо схибила --- кілька секунд \emph{непередбачуваного} випадкового струму. За двадцять хвилин гри серії відбитих м'ячів подовжувалися. Значуще зростання всередині сеансу було тільки за повного зворотного зв'язку; коли замість стимулу була тиша, культура наприкінці сеансу теж грала краще, ніж без зворотного зв'язку, але з меншим ефектом, а стійкого навчання від сеансу до сеансу впродовж кількох днів не було. Докладний розбір цього досліду --- у розділі~\ref{sec:dishbrain}.

Автори пояснили результат гіпотезою, що мережа діє так, щоб її вхід був передбачуваним~\cite{Friston2010}. Це та сама ідея, що й економія імпульсів у розділі~\ref{sec:code} та передбачення кадру в розділі~\ref{sec:recognition}: машина витрачає менше, коли вхід очікуваний. Але і сам дослід, і особливо слова авторів про «розумність» культури викликали різку критику: ефект невеликий, а пояснити його можна й простіше~\cite{Balci2023}. Чесна оцінка така: культура на стенді змінює поведінку під дією зворотного зв'язку --- це виміряно, а те, що вона вчиться саме передбачати свій вхід, --- гіпотеза.

\emph{Мережа, яка керує тілом.} Подібно до цього культуру підключали до простого віртуального «тіла» і вчили його рухатися до цілі, добираючи просторово-часовий малюнок тренувальних стимулів~\cite{Bakkum2008}. У жодному з цих дослідів немає дофаміну: учителем є малюнок струму, який обирає інженер. Що змінюється, коли учителем стає програма або сам дофамін, --- у підрозділах~\ref{sec:dishdopamine}--\ref{sec:cartpole}.

\begin{figure}[t]
\centering
\begin{tikzpicture}[>=Stealth,
  blk/.style={draw,very thick,rounded corners=2pt,align=center,font=\footnotesize,minimum height=0.8cm,fill=blue!5}]
% (a) closed loop
\node[font=\small] at (1.5,4.3) {(а) замкнена петля};
\node[blk,text width=2.6cm] (G) at (1.5,3.3) {гра: м'яч і ракетка};
\draw[very thick,fill=orange!10,rounded corners=3pt] (0.5,0.2) rectangle (2.5,1.6);
\foreach \x in {0.7,1.0,...,2.35}{\foreach \y in {0.4,0.7,1.0,1.3}{\fill[gray!60] (\x,\y) circle (0.04);}}
\draw[->,thick,blue!60!black] (G.west) -| (-0.5,0.9) -- (0.5,0.9);
\node[font=\scriptsize,blue!60!black,anchor=east,align=right] at (-0.6,2.1) {м'яч\\$\to$ струм};
\draw[->,thick,red!70!black] (2.5,0.9) -- (3.5,0.9) |- (G.east);
\node[font=\scriptsize,red!70!black,anchor=west] at (3.6,2.1) {імпульси $\to$ ракетка};
\node[font=\scriptsize] at (1.5,-0.2) {нейрони на ґратці електродів};
\node[font=\scriptsize,align=center] at (1.5,-0.85) {відбив: передбачуваний стимул\\промах: випадковий струм};
% (b) reservoir
\begin{scope}[xshift=7.9cm]
\node[font=\small] at (1.6,4.3) {(б) резервуар};
\node[blk,text width=1.1cm] (X) at (-0.4,1.0) {вхід\\$u(t)$};
\draw[very thick,fill=orange!10] (1.6,1.0) circle (0.8);
\foreach \a/\r in {20/0.35,80/0.5,150/0.3,210/0.55,280/0.4,330/0.2,110/0.15}{\fill[red!60!black] ({1.6+\r*cos(\a)},{1.0+\r*sin(\a)}) circle (0.05);}
\node[font=\scriptsize] at (1.6,-0.2) {органоїд: без учителя};
\node[blk,text width=1.8cm,fill=green!8] (W) at (4.0,1.0) {зчитування\\$W_{\text{вих}}$};
\draw[->,thick] (X) -- (0.8,1.0);
\draw[->,thick] (2.4,1.0) -- node[above,font=\scriptsize] {$\mathbf r(t)$} (W);
\draw[->,thick] (W) -- (4.0,2.3) node[above,font=\scriptsize] {$y(t)$};
\node[font=\scriptsize,green!40!black] at (4.0,0.3) {навчається з учителем};
\end{scope}
\end{tikzpicture}
\caption{Два способи змусити живу мережу обчислювати. (а) Замкнена петля: положення м'яча подається струмом, імпульси рухових ділянок рухають ракетку, а передбачуваний або випадковий стимул після удару слугує учителем. (б) Резервуар~\eqref{eq:reservoir}: органоїду не подають навчального сигналу, він розщеплює вхід на безліч нелінійних відгуків із пам'яттю; за помилкою навчається лише одне лінійне зчитування ззовні. Зв'язки самого органоїду при цьому теж змінюються --- без учителя.}
\label{fig:livingmachine}
\end{figure}

\section{Живий резервуар}

Є й інший спосіб використати живу мережу: не навчати її взагалі. У техніці він називається \emph{обчисленням на резервуарі}~\cite{Maass2002,JaegerHaas2004}. Беруть готову нелінійну систему з пам'яттю --- будь-яку, аби її відгук на вхід був багатим і залежав від недавнього минулого. Вхід $u(t)$ розгойдує її, виходить вектор відгуків $\mathbf r(t)$, і навчається тільки лінійне зчитування
\begin{empheq}[box=\keybox]{equation}
y(t) \;=\; W_{\text{вих}}\,\mathbf r(t) ,
\label{eq:reservoir}
\end{empheq}
ваги якого добирають правилом найменших квадратів~\eqref{eq:lms} з розділу~\ref{sec:motorlearning}. Резервуар робить те саме, що робили маленькі \nt{GLUT}-нейрони розділу~\ref{sec:motorlearning}: розщеплює вхід на довгий вектор, у якому потрібні класи стають роздільними площиною (розділ~\ref{sec:dendrites}).

Органоїд на ґратці електродів підійшов на роль такого резервуара. Звуки мовлення, подані на нього малюнками струму, після навчання одного лінійного зчитування розрізнялися за мовцем із точністю близько $78\%$~\cite{Cai2023}. Точність $78\%$ --- цифра, яку повідомляють самі автори і яку повторює преса. Результат корисний, але важливо розуміти, де в ньому «розум»: органоїд дає нелінійність і згасну пам'ять, навчання за помилкою зроблено ззовні, у кремнії (рис.~\ref{fig:livingmachine}). При цьому органоїд не лишається незмінним: за даними авторів, під час тренування змінювалася його власна функціональна зв'язність, і це називають навчанням без учителя в самому органоїді~\cite{Cai2023}. Яку частку точності дає ця зміна, а яку --- зчитування, не розділено.

\section{Що стенд каже про машину}

\emph{Енергія.} За оцінкою~\eqref{eq:energy} імпульс коштує близько $10^{-10}$~Дж. Культура з мільйона нейронів за частоти один імпульс на секунду витрачає на передавання сигналів
\begin{empheq}[box=\keybox]{equation}
P \;\approx\; 10^{6}\times1\,\text{с}^{-1}\times10^{-10}\,\text{Дж} \;=\; 10^{-4}\,\text{Вт},
\label{eq:dishpower}
\end{empheq}
десяту частку мілівата. Електроніка, яка стимулює, записує й обробляє ці імпульси, витрачає на порядки більше. Як і в розділі~\ref{sec:debugging}, вузьке місце --- не обчислення, а зв'язок із ним.

\emph{Відсутні деталі.} Стенд показує, що вміє шина \nt{GLUT} із керуванням потоком \nt{GABA} без регістрів керування: самоорганізовуватися, генерувати залпи, змінювати відповіді під дією зворотного зв'язку і слугувати добрим резервуаром. Чого вона без них не вміє, --- самостійно вирішити, що вважати успіхом. Цей сигнал на стенді подає інженер, а в мозку, як ми бачили в розділах~\ref{sec:dofa}--\ref{sec:acet}, --- широкомовні канали.

\emph{Етика.} Органоїди вирощують із клітин людини, і що складнішими вони стають, то серйозніше постає питання, чи може така тканина щось відчувати. Інженерний погляд підказує часткову відповідь: біль у розділі~\ref{sec:pain} --- це не сигнал одного датчика, а ціла система тривоги з воротами, регулятором гучності та широкомовними каналами, яких на стенді немає. Але це міркування, а не доведення, і сама галузь пропонує вести етичну оцінку разом із дослідами, від самого початку~\cite{Smirnova2023}.

\begin{keyblock}
\begin{proposition}[Жива мережа на стенді]
\label{prop:livingmachine}
Мережа з \nt{GLUT}- і \nt{GABA}-нейронів на ґратці електродів сама породжує залпи~\eqref{eq:burstinterval}, змінює відповіді під дією зворотного зв'язку і може слугувати резервуаром~\eqref{eq:reservoir} для навченого ззовні зчитування. Усе це виміряно. Те, що така мережа вчиться за якимось загальним правилом, наприклад передбачати свій вхід, --- гіпотеза, а гучних слів про її «розумність» більшість фахівців не приймає.
\end{proposition}
\end{keyblock}

\section{Дофамін за спалахом світла}
\label{sec:dishdopamine}

Єдиний відомий нам прямий дослід, де культурі дали дофамін як учителя, поставлено на платформі з органоїдів, до якої дослідники підключаються віддалено~\cite{Jordan2024}. У рідину, що омиває органоїди, додали дофамін у світлочутливій «клітці»: зв'язана молекула не діє на рецептори. Концентрація дофаміну в клітці --- $1$~мМ. Коли треба нагородити мережу, на неї світять ультрафіолетом: клітка розпадається, і дофамін виходить в об'єм.

Петля влаштована так. Той самий стимул подають раз у раз; якщо він викликав більше імпульсів, ніж раніше, вмикається спалах і вивільняється дофамін. За $240$ повторів кількість викликаних імпульсів загалом зростала. Автори самі пишуть, що з одного такого досліду не можна зробити висновок, що зростання спричинив дофамін~\cite{Jordan2024}.

Для інженера це перша спроба зібрати в чашці правило трьох множників~\eqref{eq:threefactor}: активність відправника й отримувача залишає слід, а спільний сигнал вирішує, чи закріпити зміну. Але сигнал тут буває лише додатним: винагорода є або її немає, від'ємної помилки $\delta<0$ установка не видає. До того ж дофамін розпливається й відкачується повільно, як концентрація в~\eqref{eq:lowpass}, тож часті спалахи можуть просто підняти його загальний рівень. Щоб відрізнити від цього навчання, потрібні два контролі: стільки ж спалахів у випадкові моменти, не пов'язані з відповіддю мережі, і ті самі спалахи без дофаміну в клітці. І потрібна перевірка після відпочинку: чи залишилася зміна, коли дофамін пішов.

\section{Дофамін навчає кремній}
\label{sec:biohybrid}

У зворотному досліді живі клітини навчають електроніку~\cite{Keene2020}. Клітини, що виділяють дофамін, вирощують у мікроканалі над органічним нейроморфним елементом --- транзистором, провідність каналу якого відіграє роль ваги синапсу. Дофамін, що вийшов із клітин, окиснюється біля електрода елемента, і це змінює провідність. Пристрій відтворює і механізм прибирання медіатора зі щілини, тому вагу можна не тільки змінити надовго, а й повернути назад.

Схемотехнічно це інтегратор із витоком за хімічним входом: вага накопичує потік дофаміну, а «відкачування» задає, як швидко вона забувається, --- те саме RC-коло, що й~\eqref{eq:lowpass}. Головне тут --- напрям зв'язку. Щуп із розділу~\ref{sec:debugging} не бачить широкомовних регістрів, бо вони хімічні. Цей елемент читає саме хімічний регістр і перетворює його на електричну вагу. Для інтерфейсів мозок---комп'ютер це шлях до датчика, який чує не тільки шину даних, а й регістри керування.

\section{Органоїди з власними дофаміновими нейронами}
\label{sec:midbrainorg}

Зі стовбурових клітин людини вміють вирощувати органоїди, схожі на ту ділянку мозку, де розташовані \nt{DOPA}-нейрони. У них є робочі дофамінові нейрони, які виділяють дофамін~\cite{Jo2016}. Такі органоїди вирощують здебільшого для вивчення хвороб. Дослідів, де їхній дофамін слугував би учителем для іншої мережі, ми не знайшли.

Для машини з живих нейронів це відсутня деталь. Стенд із розділу~\ref{sec:livingmachine} --- шина даних і керування потоком без регістрів керування; тут джерело широкомовного сигналу вже є. Бракує критика: мережі, яка обчислювала б помилку передбачення~\eqref{eq:ctd} і керувала б цими нейронами. З'єднати органоїд кори з таким органоїдом так, щоб дофамін виділявся за помилкою, --- відкрита задача, і поки що це гіпотеза, а не дослід.

\section{Органоїд балансує жердиною без дофаміну}
\label{sec:cartpole}

Найповніший із недавніх дослідів обходиться зовсім без дофаміну~\cite{Robbins2026}. Органоїди кори миші підключили замкненою петлею до задачі «жердина на візку»: активність органоїду рухає візок, положення жердини подається назад стимуляцією. Між спробами органоїд отримує короткі високочастотні навчальні стимули. Які саме --- обирає програма навчання з підкріпленням.

Результати виміряно:
\begin{itemize}
\item у більшості органоїдів навчальні сигнали, обрані програмою, дали кращий результат, ніж випадково обрані сигнали або відсутність сигналу;
\item після $45$ хвилин відпочинку покращення не зберігалося;
\item блокада рецепторів глутамату обох типів, AMPA і NMDA, скасовувала покращення.
\end{itemize}

Останній пункт показує, де живе вивчене: у \nt{GLUT}-синапсах, і через той рецептор, який у підрозділі~\ref{sec:weightstore} працює як детектор збігу входу й виходу. Другий пункт показує, чого бракує: швидка зміна є, закріплення немає, як у швидкого учня з розділу~\ref{sec:motorlearning} без повільного. А роль учителя змінилася: інженера, який обирає малюнок струму, замінила програма, але учитель, як і раніше, стоїть поза чашкою.

Серед раніших дослідів навчання культур на електродних матрицях ми також не знайшли жодного, де використовувався б дофамін: навчальним сигналом скрізь слугувала електрична стимуляція.

\section{Хто навчає культуру}

\begin{table}[t]
\centering
\footnotesize
\setlength{\tabcolsep}{4pt}
\renewcommand{\arraystretch}{1.15}
\begin{tabular}{>{\raggedright}p{2.7cm}>{\raggedright}p{2.5cm}>{\raggedright}p{3.2cm}>{\raggedright\arraybackslash}p{4.4cm}}
\toprule
Дослід & Хто навчає & Навчальний сигнал & Що відомо \\
\midrule
зупинка стимуляції за потрібної відповіді & інженер & кінець стимуляції & відповідь закріплюється --- виміряно \\
DishBrain (розд.~\ref{sec:dishbrain}) & інженер & передбачуваний або випадковий струм & значуще зростання розіграшів за сеанс лише за повного зворотного зв'язку, за тиші ефект менший --- виміряно; від сеансу до сеансу нестійкий; причина дискусійна \\
жердина на візку & програма навчання з підкріпленням & обрані програмою високочастотні стимули & краще за випадкові, але не зберігається після відпочинку --- виміряно \\
дофамін за спалахом & правило «більше імпульсів --- винагорода» & дофамін, вивільнений світлом & зростання імпульсів --- спостереження; роль дофаміну не доведено \\
біогібридний синапс & живі клітини & дофамін від клітин & тривала зміна і повернення ваги елемента --- виміряно \\
органоїди з \nt{DOPA}-нейронами & --- & --- & джерело дофаміну є; як учителя не використано \\
\bottomrule
\end{tabular}
\caption{Хто і чим навчає живі нейрони на чипі.}
\label{tab:dishteachers}
\end{table}

В усіх дослідах, де навчається сама культура, учитель поки що ззовні (таблиця~\ref{tab:dishteachers}): інженер, програма або правило, задане інженером. Дофамін потрапив у чашку лише одного разу, і його роль ще належить довести. Зібрати в чашці справжній широкомовний канал --- із критиком, помилкою і слідом, як у розділі~\ref{sec:dofa}, --- наступний крок, якого ще ніхто не зробив.


\section{Підсумок}

\begin{table}[t]
\centering
\footnotesize
\setlength{\tabcolsep}{4pt}
\renewcommand{\arraystretch}{1.15}
\begin{tabular}{>{\raggedright}p{2.8cm}>{\raggedright}p{4.2cm}>{\raggedright\arraybackslash}p{5.8cm}}
\toprule
Дослід & Що робить жива мережа & Що відомо \\
\midrule
мережа без входу & залпи з інтервалами від секунди до хвилин~\eqref{eq:burstinterval} & виміряно; механізм через виснаження синапсів --- загальноприйнята модель \\
учитель, який замовкає & потрібна відповідь закріплюється, коли стимуляцію вимикають & виміряно \\
гра в теніс & серії відбитих м'ячів подовжуються; значуще зростання за сеанс --- тільки за повного зворотного зв'язку & зміну поведінки виміряно; навчання від сеансу до сеансу нестійке; величина ефекту і пояснення передбаченням входу дискусійні \\
віртуальне тіло & рух до цілі за дібраного малюнка стимулів & виміряно \\
резервуар & нелінійність і пам'ять для зчитування~\eqref{eq:reservoir} & виміряно; навчання за помилкою --- ззовні, у кремнії; зв'язки органоїду теж змінюються \\
енергія & близько $10^{-4}$~Вт на мільйон нейронів~\eqref{eq:dishpower} & оцінка за~\eqref{eq:energy} \\
\bottomrule
\end{tabular}
\caption{Що показали машини з живих нейронів.}
\label{tab:livingmachine}
\end{table}

Машина з живих нейронів (таблиця~\ref{tab:livingmachine}) --- це стенд, на якому видно, що вміють шина даних і керування потоком без регістрів керування. Вони вміють багато, але не вміють самі вирішити, що добре. На стенді цю роль відіграє інженер зі своїм малюнком струму; у мозку --- ті нечисленні широкомовні проводи, яким присвячено середину книги.

\section{Задачі}

\begin{exercise}[\lvlA]
\label{ex:21:1}
\emph{Інтервал між залпами.} Залп виснажує запас до нуля, далі $\tau_{\text{відн}}\,\dd x/\dd t=1-x$, і новий залп починається при $x=x^*$; $\tau_{\text{відн}}=0.8$~с. (а)~Знайдіть інтервал при $x^*=0.9$ і $0.99$. (б)~Поріг $x^*$ випадково змінюється на $\pm0.005$ біля $0.99$. У якому діапазоні лежать інтервали?
\end{exercise}

\begin{exercise}[\lvlA]
\label{ex:21:2}
\emph{Енергія стенда.} (а)~Яку потужність дає оцінка~\eqref{eq:dishpower} для всього мозку ($8.6\cdot10^{10}$ нейронів)? (б)~Стенд записує $1024$ електроди з частотою $20$~кГц по $10$~біт на відлік. У скільки разів передавання цього потоку при $1$~нДж на біт дорожче за потужність самої культури?
\end{exercise}

\begin{exercise}[\lvlB]
\label{ex:21:3}
\emph{Проєктування: придушити залпи.} Рідкісна стимуляція через багато електродів змушує кожен синапс викидати медіатор із частотою $r_b$; запас $\tau_{\text{відн}}\,\dd x/\dd t=1-x-Ur_b\tau_{\text{відн}}x$, $U=0.5$, $\tau_{\text{відн}}=0.8$~с. За якої найменшої $r_b$ запас не доростає до $x^*=0.9$; $0.99$ і наскільки при цьому слабшає синапс?
\end{exercise}

\begin{exercise}[\lvlB]
\label{ex:21:4}
\emph{Пам'ять резервуара не більша за число нейронів.} Резервуар з $N$ виходами $\mathbf r(t)$ отримує білий вхід $u(t)$ з нульовим середнім і одиничною дисперсією; $m_k$ --- найбільший квадрат кореляції зчитування $\mathbf w_k\cdot\mathbf r(t)$ з $u(t-k)$. Доведіть, що ємність пам'яті $\mathrm{MC}=\sum_{k\ge0}m_k\le N$.
\end{exercise}

\begin{exercise}[\lvlB]
\label{ex:21:5}
\emph{Моделювання: резервуар у кремнії.} Резервуар $\mathbf r(t+1)=\tanh(W\mathbf r(t)+\mathbf v\,u(t)+\mathbf c)$, $N=30$, $W$ випадкова зі спектральним радіусом $0.9$, $v_i\sim U(-0.5,0.5)$, $c_i\sim U(-0.2,0.2)$, $u\sim U(-1,1)$, зчитування --- гребенева регресія. Знайдіть ємність пам'яті задачі~\ref{ex:21:4} з $\tanh$ і без нього. Який резервуар пам'ятає більше?
\end{exercise}

\begin{exercise}[\lvlB]
\label{ex:21:6}
\emph{Зчитування для живого резервуара.} Органоїд дає $1024$ канали і $P=240$ навчальних записів двох класів. (а)~Скільки ознак $n$ можна залишити, щоб за формулою задачі~\ref{ex:19:3} випадкова розмітка була роздільною з імовірністю не більше $1\%$? (б)~На скільки сигм вища за вгадування точність $78\%$ на $100$ перевірних записах?
\end{exercise}

\begin{exercise}[\lvlC]
\label{ex:21:7}
\emph{Дослідження: учитель без широкомовного каналу.} Запропонуйте протокол стимуляції, що реалізує на стенді правило трьох множників із розділу~\ref{sec:dofa} через $8$--$1024$ електродів, і контроль із замороженим зчитуванням, який відрізняє навчання ваг від зсуву стану мережі. Який результат спростував би гіпотезу про навчання?
\end{exercise}

\chapter{DishBrain}
\chaptermark{DishBrain}
\label{sec:dishbrain}
\begin{chaptergoals}
\item як DishBrain кодує м'яч місцем і частотою, а ракетку читає як різницю лічильників;
\item чому лічба у вікні $10\ms$ шумна, тож мережа може зсунути лише середнє;
\item як мережа, яку струшують лише після промахів, дрейфує до налаштувань, що рідше помиляються;
\item що виміряли, про що сперечаються і як перебудувати стенд, щоб перевірити механізм.
\end{chaptergoals}

\section{Мережа в чашці грає в теніс}

DishBrain --- так автори назвали стенд, на якому культура живих нейронів, вирощена на чипі з електродами, грає в спрощений теніс з однією ракеткою~\cite{Kagan2022}. Це найобговорюваніший дослід із машиною з живих нейронів (огляд таких машин --- у розділі~\ref{sec:livingmachine}), і для нашої книги він особливо зручний. Тут маленька мережа доступна вся: кожен вхід задає експериментатор, кожен вихід записується. Усі питання попередніх розділів --- як записано сигнал, хто навчає, що бачить щуп --- можна поставити мережі, яка не має ні очей, ні м'язів, ні \nt{DOPA}-нейронів. Розберімо дослід так, як інженер розбирає чужий стенд: схема, вхід, вихід, учитель, результат, суперечки.

\section{Стенд}

Нейрони беруть із кори зародка миші --- близько $800\,000$ клітин на чип --- або вирощують зі стовбурових клітин людини, близько мільйона на чип. Кілька тижнів вони ростуть на чипі, проростають одна в одну й починають самі видавати імпульси та пачки. На майданчику $8$~мм$^2$ під ними лежить $26\,000$ платинових електродів; одночасно можна записувати до $1024$ із них, а стимулювати на практиці --- через вісім.

Навколо культури замкнено петлю (рис.~\ref{fig:dishloop}). Комп'ютер моделює гру: м'яч летить, відбивається від стін, ракетка рухається вгору і вниз. Положення м'яча перетворюється на струм на восьми «чутливих» електродах. Імпульси з двох «рухових» ділянок підраховують вікнами по $10\ms$, і вони рухають ракетку. Після кожного удару чи промаху культура отримує ще один стимул --- зворотний зв'язок. Вікна по $10\ms$ --- це такт комп'ютера стенда, а не культури: сама мережа, як і раніше, працює асинхронно, як і всі мережі цієї книги.

\begin{figure}[t]
\centering
\begin{tikzpicture}[>=Stealth,
  blk/.style={draw,very thick,rounded corners=3pt,align=center,font=\small,minimum height=1.2cm},
  lab/.style={font=\scriptsize,align=center}]
\node[blk,fill=blue!6,text width=3.4cm] (C) at (0,0) {культура на чипі\\\scriptsize $\sim10^6$ нейронів, $26\,000$ електродів};
\node[blk,fill=orange!10,text width=3.0cm] (G) at (7.2,0) {комп'ютер:\\гра в теніс};
\draw[->,very thick,blue!60!black] (G.north west) to[bend right=25] node[lab,above] {м'яч: \emph{місце} --- який із 8 електродів,\\\emph{частота} --- 4--40 імп/с} (C.north east);
\draw[->,very thick,red!70!black] (C.south east) to[bend right=25] node[lab,below] {ракетка: лічба імпульсів\\ділянок «вгору» і «вниз» за $10\ms$} (G.south west);
\draw[->,thick,densely dashed,green!45!black] (G.east) -- ++(0.9,0) |- ++(0,-2.6) -| node[lab,pos=0.25,below] {відбив: передбачуваний стимул; промах: непередбачуваний} (C.south);
\end{tikzpicture}
\caption{Петля DishBrain. Синя стрілка --- вхід: положення м'яча записано кодом місця і частотним кодом (розділ~\ref{sec:code}). Червона --- вихід: різниця лічильників двох ділянок рухає ракетку. Зелена штрихова --- зворотний зв'язок після кожного розіграшу, єдиний «учитель» стенда.}
\label{fig:dishloop}
\end{figure}

\section{Вхід і вихід}

\emph{Вхід} записано одразу двома кодами розділу~\ref{sec:code}. Те, який із восьми електродів стимулюється, каже, на якій висоті м'яч: це код місця, як у датчиків розділу~\ref{sec:sensors}. Те, як часто йдуть імпульси стимуляції, каже, наскільки м'яч далеко: $4$ на секунду, коли він біля дальньої стіни, і лінійно більше, до $40$ на секунду, коли він біля стіни ракетки. Це частотний код; амплітуда імпульсів м'яча --- $75$~мВ. Культура не «знає» заздалегідь жодного з цих кодів: код задав експериментатор, і мережа має навчитися його читати.

\emph{Вихід} читають так само, як в інтерфейсах розділу~\ref{sec:debugging}. Імпульси однієї рухової ділянки штовхають ракетку вгору, іншої --- вниз; у найпростішому записі за кожне вікно
\begin{empheq}[box=\keybox]{equation}
\Delta p \;=\; c\,\bigl(n_\uparrow - n_\downarrow\bigr),
\label{eq:dishreadout}
\end{empheq}
де $n_\uparrow$, $n_\downarrow$ --- кількості імпульсів двох ділянок за $10\ms$, а $c$ --- коефіцієнт стенда. Це двопровідний код розділу~\ref{sec:glutcomputer}: знак руху несе не знак сигналу, а те, який із двох проводів гучніший.

Погляньмо на шум такого виходу. Якщо ділянка видає загалом $R$ імпульсів на секунду, за вікно $T=10\ms$ їх у середньому $RT$, і за~\eqref{eq:rateerror} відносний розкид дорівнює $1/\sqrt{RT}$. При $R=1000$ імпульсів на секунду це близько $30$ відсотків у кожному вікні; при $R=100$ --- понад сто. Ракетка неминуче тремтить, і навчитися мережа може лише в один спосіб --- зсунути \emph{середню} різницю лічильників у потрібний бік. Це ті самі закони, що обмежують будь-який інтерфейс розділу~\ref{sec:debugging}, тільки тепер по обидва боки петлі.

\section{Учитель без дофаміну}

Найцікавіше в стенді --- учитель. У культурі нейронів кори немає \nt{DOPA}-нейронів, і жодного сигналу «краще, ніж очікувалося» стенд не надсилає. Зворотний зв'язок інший. Відбила мережа м'яч --- на всі вісім чутливих електродів одразу подається короткий \emph{передбачуваний} стимул: імпульси по $75$~мВ, $100$ на секунду впродовж $0.1\,\text{с}$. Схибила --- чотири секунди триває стимул, який автори називають \emph{непередбачуваним}: імпульси по $150$~мВ, $5$ на секунду, у випадкові моменти на випадково вибрані електроди; далі чотири секунди паузи без стимуляції, і м'яч стартує знову~\cite{Kagan2022}.

Автори виходили з гіпотези, що мережа діє так, щоб її вхід ставав передбачуваним~\cite{Friston2010}. Для інженера зміст простий: культура має рівно один спосіб позбутися чотирьох секунд випадкового шуму --- відбити м'яч. Ту саму ідею ми зустрічали в економії імпульсів розділу~\ref{sec:code} і в передбаченні кадру в розділі~\ref{sec:recognition}.

Але чи потрібна тут саме передбачуваність? Можливий і грубіший механізм. Нехай непередбачуваний стимул після промаху просто \emph{перемішує} недавні зміни в мережі, а спокійний після удару їх не зачіпає. Опишімо налаштування мережі одним вектором $\boldsymbol\psi$, і нехай у налаштуванні $\boldsymbol\psi$ мережа хибить з імовірністю $P_{\text{пр}}(\boldsymbol\psi)$. Якщо поштовхи симетричні --- з $\boldsymbol\psi$ у $\boldsymbol\psi'$ так само ймовірно, як назад, --- то в усталеному режимі потік виходів із кожного налаштування врівноважений потоком входів, і частка часу, яку мережа проводить у налаштуванні $\boldsymbol\psi$, дорівнює
\begin{empheq}[box=\keybox]{equation}
\rho(\boldsymbol\psi) \;\propto\; \frac{1}{P_{\text{пр}}(\boldsymbol\psi)} .
\label{eq:dishdensity}
\end{empheq}
Налаштування, яке хибить удесятеро рідше, утримується вдесятеро довше. Ніхто не повідомляє мережі, в який бік змінювати ваги: добрі налаштування просто рідше руйнуються.

Ми перевірили це на моделі, звівши культуру до двох чисел: ракетка приходить у точку $a\,y+b$, коли м'яч приходить на висоту $y$, і відбиває, якщо схибила менш ніж на $0.1$. Після промаху обидва числа отримують випадковий поштовх, після удару лишаються незмінними. Частка відбитих м'ячів за $2000$ розіграшів зросла з $0.21$ у перших двохстах до $0.38$ в останніх п'ятистах (сорок «культур», рис.~\ref{fig:dishmodel}). Якщо поштовх давати після \emph{кожного} розіграшу, у зворотному зв'язку немає відомостей, і частка лишається близько $0.12$; якщо поштовхів немає зовсім, вона стоїть на $0.19$. Дослід на стенді з цим узгоджується: значуще зростання всередині сеансу було тільки за повного зворотного зв'язку. Коли замість стимулів була тиша, культура наприкінці сеансу все ж грала краще, ніж без зворотного зв'язку, але з меншим ефектом~\cite{Kagan2022}; зауважмо, що і в цій умові після промаху йде довга пауза, а після удару --- коротка, тож різниця між наслідками не зникає цілком.

\begin{figure}[t]
\centering
\begin{tikzpicture}[x=1cm,y=5cm]
\draw[->,gray!70!black] (0,0) -- (10.6,0);
\draw[->,gray!70!black] (0,0) -- (0,0.5) node[above,font=\small,black] {частка відбитих};
\foreach \v in {0.1,0.2,0.3,0.4}{\draw[gray!60!black] (-0.06,\v) -- (0.06,\v); \node[left,font=\scriptsize] at (0,\v) {\v};}
\foreach \x/\f/\l/\lab in {1.8/0.21/0.38/{шум після\\промаху},5.2/0.13/0.11/{шум після\\кожного розіграшу},8.6/0.19/0.19/{без шуму}}{
  \fill[blue!35] (\x-0.8,0) rectangle (\x-0.05,\f);
  \fill[red!60!black] (\x+0.05,0) rectangle (\x+0.8,\l);
  \node[below,font=\scriptsize,align=center] at (\x,-0.01) {\lab};
  \node[above,font=\scriptsize] at (\x-0.425,\f) {\f};
  \node[above,font=\scriptsize] at (\x+0.425,\l) {\l};}
\fill[blue!35] (7.4,0.47) rectangle (7.7,0.495); \node[right,font=\scriptsize] at (7.75,0.482) {перші 200};
\fill[red!60!black] (7.4,0.41) rectangle (7.7,0.435); \node[right,font=\scriptsize] at (7.75,0.422) {останні 500};
\end{tikzpicture}
\caption{Модель «перемішування при промаху» (\texttt{dishbrain\_model.py}): частка відбитих м'ячів на початку і наприкінці $2000$ розіграшів, середнє за сорока моделями-культурами. Навчання йде, тільки якщо поштовх настає після промаху, але не після удару, --- як і в досліді, де значуще зростання за сеанс було лише за повного зворотного зв'язку.}
\label{fig:dishmodel}
\end{figure}

\begin{keyblock}
\begin{proposition}[Учитель-перемішувач]
\label{prop:dishscramble}
Якщо невдача струшує мережу, а вдача залишає її в спокої, мережа сама дрейфує до налаштувань, які рідше помиляються: час, проведений у налаштуванні, обернено пропорційний імовірності невдачі~\eqref{eq:dishdensity}. Для такого навчання не потрібні ні сигнал винагороди, ні його знак, ні передбачуваність як така --- досить, щоб зворотний зв'язок після вдачі та після невдачі різнився. Зміну поведінки культури виміряно; який механізм працює в чашці --- передбачення входу чи перемішування, --- не встановлено.
\end{proposition}
\end{keyblock}

Порівняйте з розділом~\ref{sec:dofa}. Там шум теж слугував пошуком, але дофамін мав знак: помилка $\delta$ казала, в який бік зсунути ваги, і крок ішов за градієнтом~\eqref{eq:perturbation}. Тут знака немає --- є лише «залишити» або «струснути». Такий пошук грубіший і повільніший, але вимагає від мережі менше: ні \nt{DOPA}-нейронів, ні критика, ні сліду, розтягнутого на секунди.

\section{Що виміряли і про що сперечаються}

Кожна гра тривала $20$ хвилин; порівнювали перші $5$ хвилин з останніми $15$. У культур із повним зворотним зв'язком серії відбитих м'ячів стали довшими, а миттєвих промахів поменшало; у контролях без клітин, без гри та з комп'ютерною «ракеткою» нічого подібного не було. Культури з клітин людини в середньому відбивали довше за мишачі. Покращення статистично надійне --- на дев'яти мишачих і одинадцяти людських культурах, у сотні з гаком ігор кожної групи, --- але невелике: мережа не стала добрим гравцем, вона стала трохи кращою за випадкову. Надійне воно всередині сеансу; стійкого навчання від сеансу до сеансу впродовж кількох днів автори не отримали~\cite{Kagan2022}. У наступній роботі тієї самої групи виявили, що під час гри активність культури наближається до \emph{критичного} режиму --- межі між згасанням і лавиною, того самого життя на межі, що в розділі~\ref{sec:gaba}, --- і що ближче, то краща гра~\cite{Habibollahi2023}.

Головну суперечку викликали не цифри, а слова. Заголовок статті стверджував, що нейрони в чашці «виявляють розумність» (sentience), і група дослідників заперечила: зміна поведінки в замкненій петлі --- ще не розум і не відчуття, а висновки треба формулювати скромніше~\cite{Balci2023}. З інженерного погляду відкриті два питання, і обидва --- про механізм. Перше: чи навчається мережа, тобто чи змінюються її ваги надовго, чи зворотний зв'язок лише зсуває її поточний стан? Друге: чи потрібна саме передбачуваність, чи вистачає будь-якої різниці між відгуком на вдачу і на невдачу, як у твердженні~\ref{prop:dishscramble}? Стенд, на якому доступний кожен електрод, дає змогу відповісти на обидва --- і це, мабуть, його головна цінність.

\section{Специфікація стенда: як його відтворити}
\label{sec:dishspec}

Нижче зібрано все, що потрібно, щоб зібрати такий стенд наново: обладнання, петля, кодування входу, читання виходу, зворотний зв'язок і протокол досліду. Кожен параметр позначено джерелом. «Стаття» --- значення взято з тексту роботи~\cite{Kagan2022}. «Вибір» --- у статті його немає, і для відтворення його доведеться задати самому; ми пропонуємо типове значення і кажемо, як його перевірити.

\subsection*{Обладнання і культура}

\begin{center}
\footnotesize
\setlength{\tabcolsep}{4pt}
\renewcommand{\arraystretch}{1.15}
\begin{tabular}{>{\raggedright}p{4.0cm}>{\raggedright}p{6.9cm}>{\raggedright\arraybackslash}p{1.6cm}}
\toprule
Параметр & Значення & Джерело \\
\midrule
чип & матриця MaxOne: $26\,000$ платинових електродів на майданчику $8$~мм$^2$ & стаття \\
запис & до $1024$ каналів одночасно, $20\,000$ відліків на секунду & стаття \\
стимуляція & до $32$ електродів технічно, у досліді --- $8$ незалежних & стаття \\
клітини & кора зародка миші; нейрони людини зі стовбурових клітин (два способи вирощування); клітини HEK293T (не нейрони) для контролю & стаття \\
висів & близько $10^6$ клітин на чип & стаття \\
вік культури на час досліду & миша: від $\sim10$ діб; людина: $14$--$24$ доби або $\sim80$ діб залежно від способу вирощування & стаття \\
\bottomrule
\end{tabular}
\end{center}

\subsection*{Петля і час}

Петля працює вікнами по $10\ms$ ($200$ відліків): за вікно підраховують імпульси на електродах рухових ділянок, гра оновлюється, і видається чергова стимуляція. Затримка від імпульсу в культурі до відповідного стимулу --- близько $5\ms$. Під час гри сигнал кожного електрода проходить фільтр верхніх частот Бесселя 2-го порядку з частотою зрізу $100$~Гц; модуль сигналу згладжується фільтром нижніх частот Бесселя 1-го порядку з частотою $1$~Гц, і поріг імпульсу пропорційний цьому згладженому модулю. Поріг у шість стандартних відхилень фону стаття вказує для окремих вимірювань активності, а не для гри. Щоб стимуляцію не зараховували як відповідь культури, усі сигнали гасять, коли одночасно надходить понад $15$ великих, вищих за $75$~мВ, імпульсів. Усе це --- стаття.

\subsection*{Кодування м'яча}

\begin{center}
\footnotesize
\setlength{\tabcolsep}{4pt}
\renewcommand{\arraystretch}{1.15}
\begin{tabular}{>{\raggedright}p{3.4cm}>{\raggedright}p{7.5cm}>{\raggedright\arraybackslash}p{1.6cm}}
\toprule
Параметр & Значення & Джерело \\
\midrule
чутлива ділянка & $8$ електродів, розташованих так, щоб сусідні електроди означали сусідні положення м'яча & стаття \\
код місця & номер електрода $k=1,\dots,8$ задає положення м'яча відносно центру ракетки & стаття \\
яка координата & вертикальне положення м'яча; для відтворення --- відносно ракетки, $k=\lceil 8(y-p+\tfrac12)\rceil$ при $y-p\in[-\tfrac12,\tfrac12]$ & стаття; формула --- вибір \\
частотний код & частота стимуляції від $4$ до $40$~імп/с несе відстань до ракетки & стаття \\
напрям шкали & $4$~імп/с, коли м'яч біля протилежної стіни, і лінійно до $40$~імп/с, коли він біля стіни ракетки: $f=4+36\,(1-x/L)$, $x$ --- відстань до стіни ракетки, $L$ --- ширина корту & стаття \\
форма імпульсу & короткий прямокутний двофазний: спершу додатна, потім від'ємна фаза & стаття \\
амплітуда & $75$~мВ; обрано так, щоб не змушувати спрацьовувати загальмовані клітини & стаття \\
тривалість фаз & не вказано; взяти стандартне налаштування системи стимуляції і не змінювати його впродовж досліду & вибір \\
\bottomrule
\end{tabular}
\end{center}

\subsection*{Читання ракетки}

У статті дві рухові ділянки: імпульси в першій рухають ракетку вгору, у другій --- вниз; внесок ділянок зважують, щоб урахувати різну щільність клітин і різну активність~\cite{Kagan2022}. Точної формули не наведено. Для відтворення беремо правило~\eqref{eq:dishreadout}, $\Delta p=c\,(n_\uparrow-n_\downarrow)$ за кожне вікно $10\ms$, з вагою, що вирівнює ділянки за їхньою середньою активністю в спокої (вибір): $n_\uparrow$ і $n_\downarrow$ ділять на середню лічбу своєї ділянки за вікно, виміряну перед грою. Коефіцієнт $c$ обираємо так, щоб різниця в одну середню лічбу зсувала ракетку на $1\%$ висоти корту за вікно (вибір); тоді стала перевага однієї ділянки проводить ракетку через увесь корт за $1$~с.

Розташування ділянок на чипі в тексті статті координатами не задано; є лише схема. Для відтворення досить трьох груп електродів, що не перетинаються: вісім чутливих у центрі та по групі записувальних електродів з обох боків, одна --- «вгору», друга --- «вниз» (вибір).

\subsection*{Гра}

Розмірів корту, довжини ракетки і швидкості м'яча в статті не вказано; сказано лише, що м'яч летить зі сталою швидкістю і відбивається від стін. Для відтворення (вибір): корт $1\times1$, ракетка на правій стінці завдовжки $0.2$ (влучання при $|y-p|<0.1$, як у моделі~\texttt{models/dishbrain\_model.py}), м'яч перетинає корт за $1$~с. Промах без жодного відбитого м'яча в розіграші вважають «ейсом»; розіграш, довший за три удари поспіль, --- «довгим» (стаття).

\subsection*{Зворотний зв'язок}

\begin{center}
\footnotesize
\setlength{\tabcolsep}{4pt}
\renewcommand{\arraystretch}{1.15}
\begin{tabular}{>{\raggedright}p{2.6cm}>{\raggedright}p{8.3cm}>{\raggedright\arraybackslash}p{1.6cm}}
\toprule
Подія & Що подається & Джерело \\
\midrule
удар & передбачуваний стимул: усі $8$ чутливих електродів одночасно, $75$~мВ, $100$~імп/с, $100\ms$; звичайна стимуляція м'яча на цей час переривається & стаття \\
промах & непередбачуваний стимул: $150$~мВ, $5$~імп/с, $4$~с, у випадкові моменти на випадково вибрані з $8$ електроди; потім $4$~с без стимуляції, і м'яч стартує у випадковому напрямку & стаття \\
закон випадковості & не вказано; для відтворення --- моменти імпульсів як пуассонівський потік із середньою частотою $5$~імп/с & вибір \\
\bottomrule
\end{tabular}
\end{center}

\subsection*{Протокол і контроль}

Один сеанс триває $20$~хв; порівнюють перші $5$ і останні $15$~хв (стаття). Три міри результату: середня довжина розіграшу, частка ейсів і частка довгих розіграшів. Умови досліду (стаття):
\begin{itemize}
\item \emph{стимул} --- повна петля, як описано вище;
\item \emph{тиша} --- замість стимулу удару чи промаху стільки ж часу без жодної стимуляції, потім м'яч стартує у випадковому напрямку;
\item \emph{без зворотного зв'язку} --- м'яч після промаху просто відбивається й летить далі, стимуляція положення м'яча триває;
\item \emph{спокій} --- гра йде і ракетка рухається за активністю, але стимуляції немає зовсім;
\item \emph{контролі} --- чип лише із середовищем; клітини HEK293T; «культура в комп'ютері», симуляція замість клітин.
\end{itemize}
Обсяг вибірки в основній серії: $9$ культур миші ($101$ сеанс), $11$ культур людини ($138$), $6$ чипів із середовищем ($80$), $20$ культур у спокої ($42$), $3$ симуляції ($38$), разом $399$ сеансів. У серії про зворотний зв'язок --- $486$ сеансів за $3$ дні по $3$ сеанси на день, умови «стимул», «тиша» і «без зворотного зв'язку» ($n=27$, $17$ і $15$). Зростання середньої довжини розіграшу від перших $5$ хвилин до останніх $15$ отримано лише в живих культур. У серії про зворотний зв'язок значуще зростання в часі --- тільки в умові «стимул»; наприкінці сеансу умова «тиша» все ж перевершувала «без зворотного зв'язку», але з меншим ефектом, а стійкого навчання від сеансу до сеансу за три дні не було~\cite{Kagan2022}.

\subsection*{Цикл стенда крок за кроком}

Кожні $10\ms$ комп'ютер стенда робить таке.
\begin{enumerate}
\item Зчитує кількість імпульсів $n_\uparrow$, $n_\downarrow$ у двох рухових ділянках за минуле вікно і нормує їх на середню лічбу ділянки в спокої.
\item Зсуває ракетку на $\Delta p=c\,(n_\uparrow-n_\downarrow)$ і обмежує її краями корту.
\item Зсуває м'яч на один крок; відбиває його від верхньої, нижньої та лівої стінок.
\item Якщо м'яч дійшов до правої стінки: при $|y-p|<0.1$ --- удар, лічба розіграшу зростає на один, подається стимул удару ($100\ms$); інакше --- промах, розіграш записується, подається стимул промаху ($4$~с), потім $4$~с паузи, і м'яч стартує знову.
\item Інакше обирає електрод $k$ за вертикальним положенням м'яча і частоту $f$ за відстанню до стіни ракетки та подає на електрод $k$ двофазні імпульси $75$~мВ із частотою $f$ до наступного вікна.
\end{enumerate}
Цього досить, щоб зібрати стенд, сумісний з опублікованим, і перевірити головне питання розділу: чи навчає культуру передбачуваність, чи проста різниця між відповіддю на удар і на промах. Для цього до трьох умов статті варто додати четверту, якої в ній немає: після промаху --- \emph{передбачуваний} стимул тієї самої тривалості й енергії, що й непередбачуваний. Якщо культура навчається і так, річ у різниці, а не в передбачуваності.


\section{Підсумок}

\begin{table}[t]
\centering
\footnotesize
\setlength{\tabcolsep}{4pt}
\renewcommand{\arraystretch}{1.15}
\begin{tabular}{>{\raggedright}p{2.1cm}>{\raggedright}p{4.2cm}>{\raggedright}p{1.9cm}>{\raggedright\arraybackslash}p{4.6cm}}
\toprule
Вузол стенда & Як улаштований & Розділ книги & Що відомо \\
\midrule
вхід & місце (який із 8 електродів) і частота (4--40 імп/с) & \ref{sec:code}, \ref{sec:sensors} & задано експериментатором \\
вихід & різниця лічильників двох ділянок за $10\ms$~\eqref{eq:dishreadout} & \ref{sec:glutcomputer}, \ref{sec:debugging} & задано експериментатором; шум за~\eqref{eq:rateerror} \\
учитель & передбачуваний стимул після удару, непередбачуваний після промаху; дофаміну немає & \ref{sec:dofa} & значуще зростання за сеанс --- тільки за повного зворотного зв'язку; за тиші ефект менший; від сеансу до сеансу нестійкий \\
механізм & передбачення входу або перемішування при промаху~\eqref{eq:dishdensity} & \ref{sec:dofa}, \ref{sec:recognition} & не встановлено; обидва пояснюють результат \\
режим мережі & наближення до критичного режиму під час гри & \ref{sec:gaba} & зв'язок з якістю гри виміряно \\
«розумність» & твердження авторів & --- & не показано; оскаржено \\
\bottomrule
\end{tabular}
\caption{DishBrain очима інженера.}
\label{tab:dishbrain}
\end{table}

DishBrain --- машина з живих нейронів у найпростішому вигляді: вхід задано кодом, вихід прочитано лічбою, учитель --- різниця між відгуком на вдачу і на невдачу (таблиця~\ref{tab:dishbrain}). Мережа без очей, м'язів і дофаміну трішки навчилася грати, і вже це саме по собі робить її випробувальним стендом для ідей книги. Яким би не виявився механізм --- передбачення, перемішування чи щось третє, --- відповідь знайдуть так, як радить розділ~\ref{sec:debugging}: щупом, тестовим сигналом і вимиканням частин на машині, яку нарешті видно цілком.

\section{Задачі}

\begin{exercise}[\lvlA]
\label{ex:22:1}
\emph{Наскільки треба зсунути лічбу.} Дві ділянки видають разом $R_\uparrow+R_\downarrow=2000$ імпульсів на секунду, потоки пуассонівські. (а)~Який відносний розкид лічби однієї ділянки за $10\ms$ при $R=1000$ і $R=100$? (б)~За $1$~с польоту м'яча зсуви~\eqref{eq:dishreadout} додаються. За якої найменшої $R_\uparrow-R_\downarrow$ середній зсув удвічі більший за розкид?
\end{exercise}

\begin{exercise}[\lvlA]
\label{ex:22:2}
\emph{Скільки розіграшів треба, щоб побачити навчання.} Розіграші незалежні, частка відбитих біноміальна. Скільки розіграшів потрібно в кожному з двох періодів, щоб зростання частки відбитих перевищило два стандартні відхилення різниці: (а)~з $0.20$ до $0.25$; (б)~з $0.21$ до $0.38$, як у моделі?
\end{exercise}

\begin{exercise}[\lvlB]
\label{ex:22:3}
\emph{Виведення~\eqref{eq:dishdensity}.} Налаштування $\boldsymbol\psi$ пробігає скінченну множину; при промаху (імовірність $P(\boldsymbol\psi)>0$) мережа переходить у $\boldsymbol\psi'$ із симетричною ймовірністю $K(\boldsymbol\psi,\boldsymbol\psi')$, при ударі лишається. (а)~Доведіть, що $\rho\propto1/P$ --- стаціонарний розподіл. (б)~Яка частка промахів за двох налаштувань із $P=0.9$ і $P=0.1$?
\end{exercise}

\begin{exercise}[\lvlB]
\label{ex:22:4}
\emph{Поглинальна область моделі.} Ракетка приходить у $p=\min\bigl(1,\max(0,ay+b)\bigr)$, м'яч --- на висоту $y\sim U(0,1)$, удар --- якщо $|p-y|<h=0.1$. (а)~Доведіть, що при $|b|<h$ і $|a-1+b|<h$ мережа відбиває всі м'ячі, і знайдіть площу цієї області. (б)~Знайдіть чисельно середню частку відбитих при $a\sim U(-1,1)$, $b\sim U(0,1)$.
\end{exercise}

\begin{exercise}[\lvlB]
\label{ex:22:5}
\emph{Моделювання: паркан навколо налаштувань.} У копії \texttt{dishbrain\_model.py} дайте $200$ культурам по $20\,000$ розіграшів. Яка частка відбитих в останніх $500$? А якщо після кожного поштовху обрізати налаштування прямокутником $a\in[-1,2]$, $b\in[-1,1]$? Поясніть різницю за допомогою задачі~\ref{ex:22:3}.
\end{exercise}

\begin{exercise}[\lvlB]
\label{ex:22:6}
\emph{Проєктування вхідного коду.} М'яч відбито, якщо помилка за висотою менша за $h=0.1$; мережа читає вхід за $T=0.25$~с, до частоти $r$ розрізняється близько $\sqrt{rT}$ рівнів, стимуляція не частіше за $40$ імпульсів на секунду. (а)~Чи вистачить коду місця на восьми електродах? (б)~Чи можна передати висоту частотою на одному електроді?
\end{exercise}

\begin{exercise}[\lvlC]
\label{ex:22:7}
\emph{Дослідження: передбачення чи перемішування?} Узагальніть модель на $d$ налаштовуваних чисел ($p=\sum_{i<d}a_iy^i$). Як зростає з $d$ кількість розіграшів до частки відбитих $0.5$ при перемішуванні та за правилом зі знаком помилки~\eqref{eq:perturbation}? Який дослід на стенді розділив би дві гіпотези?
\end{exercise}

\appendix
\renewcommand{\chaptername}{\appendixname}
\renewcommand{\theHchapter}{app\arabic{chapter}} % hyperref anchors: Ukrainian appendix letters are Cyrillic
\chapter{Речовини, які не задають тип нейрона}
\label{sec:othermediators}

У розділі~\ref{sec:neurontypes} ми розсортували нейрони за головним медіатором і отримали десять типів (таблиця~\ref{tab:types}). Тепер час додати ускладнення. Нейрон викидає не тільки головний медіатор, а в мозку, крім головних медіаторів, працюють й інші речовини. Вони змінюють те, як мережа передає і зберігає сигнали.

Крім адресної та широкомовної шин (підрозділ~\ref{sec:buses}) є й третя: \emph{зворотний канал}. Кілька речовин викидає не відправник, а \emph{отримувач}, і вони йдуть через щілину назад, до виходу відправника, і змінюють, скільки медіатора той викине наступного разу. У мережах зв'язку це схоже на підтвердження приймання, яке відправник використовує, щоб налаштувати свою передачу. Саме цим каналом нейрони користуються, коли змінюють ваги своїх зв'язків: ним вихід відправника дізнається про активність отримувача, тобто про другий множник правил навчання розділу~\ref{sec:dofa}.

Повний список відомих медіаторів довгий: понад сотня речовин, більшість із них --- короткі білкові ланцюжки, нейропептиди. Але всі вони вкладаються в кілька класів, і всередині кожного класу речовини поводяться схоже. Речовини, які не бувають головним медіатором, зібрано в таблиці~\ref{tab:othermediators}, з тими самими трьома питаннями, які поставив би інженер: якою шиною йде сигнал, чи швидко він діє і який у нього зазвичай знак. Тип нейрона вони не задають: їх викидають разом із головним медіатором, або їх викидають не нейрони, або вони йдуть у зворотний бік.

\begin{table}[t]
\centering
\footnotesize
\setlength{\tabcolsep}{4pt}
\renewcommand{\arraystretch}{1.12}
\begin{tabular}{>{\raggedright}p{2.25cm}>{\raggedright}p{2.8cm}>{\raggedright}p{1.8cm}>{\raggedright}p{1.5cm}>{\centering}p{0.8cm}>{\raggedright\arraybackslash}p{4.3cm}}
\toprule
Клас & Речовина & Шина & Швид\-кість & Знак & Роль в обчисленні \\
\midrule
Амінокислоти & D-серин & локальна & повільно & І & другий ключ для рецептора глутамату: умова «І» \\
\midrule
Моноаміни & слідові аміни & широко\-мовна & повільно & $\pm$ & тонке налаштування моноамінових каналів \\
\midrule
\multirow{2}{2.25cm}{Пурини}
 & АТФ & адресна & швидко і повільно & $+$ & супутній медіатор; сигнал між нейронами і глією \\
 & аденозин & об'ємна & повільно & $-$ & накопичується під час активності: лічильник навантаження~\cite{Dunwiddie2001} \\
\midrule
\multirow{8}{2.25cm}{Нейропептиди (понад сотню)}
 & енкефаліни, ендорфіни, динорфін & об'ємна & повільно & $-$ & придушення передачі \\
 & речовина P & об'ємна & повільно & $+$ & повільне підсилення \\
 & нейропептид Y & об'ємна & повільно & $-$ & повільне гальмування \\
 & соматостатин & об'ємна & повільно & $-$ & повільне гальмування, супроводжує GABA \\
 & вазоактивний інтестинальний пептид (ВІП) & об'ємна & повільно & $+$ & розгальмовування: гальмування гальмівних нейронів \\
 & холецистокінін & об'ємна & повільно & $\pm$ & супроводжує GABA в частині гальмівних нейронів \\
 & орексин & широко\-мовна & повільно & $+$ & стабілізатор режиму неспання \\
 & окситоцин, вазопресин & широко\-мовна & повільно & $\pm$ & глобальні налаштування поведінки \\
\midrule
\multirow{2}{2.25cm}{Гази}
 & оксид азоту NO & зворотна, об'ємна & повільно & $\pm$ & проходить крізь мембрани; зворотний сигнал під час зміни ваг~\cite{Garthwaite2008} \\
 & CO, H$_2$S & об'ємна & повільно & $\pm$ & роль вивчено мало \\
\midrule
Ліпіди & ендоканабіноїди (анандамід, 2-АГ) & зворотна & повільно & $-$ & отримувач зменшує викид у відправника: зворотний зв'язок за вагою~\cite{WilsonNicoll2001} \\
\bottomrule
\end{tabular}
\caption{Речовини, які не задають тип нейрона. \emph{Шина}, \emph{швидкість} і \emph{знак} --- як у таблиці~\ref{tab:types}; крім того, об'ємна шина --- медіатор розпливається в міжклітинному просторі навколо місця викиду, зворотна --- йде від отримувача назад до відправника, локальна --- діє поруч із місцем викиду. «І» --- дозвільний сигнал: сам по собі струму не викликає, але без нього не спрацьовує інший вхід, як другий вхід логічного елемента «І». Нейропептиди майже завжди викидаються разом із головним медіатором і здебільшого за високої частоти імпульсів~\cite{vandenPol2012}.}
\label{tab:othermediators}
\end{table}

\chapter{Позначення}
\label{sec:notation}

Типи нейронів позначено першими чотирма літерами англійської назви головного медіатора: \nt{GLUT} --- глутамат, \nt{GABA} --- гамма-аміномасляна кислота, \nt{GLYC} --- гліцин, \nt{ACET} --- ацетилхолін, \nt{DOPA} --- дофамін, \nt{SERO} --- серотонін, \nt{HIST} --- гістамін, \nt{NORA} --- норадреналін, \nt{ADRE} --- адреналін, \nt{ASPA} --- аспартат (таблиця~\ref{tab:types}).

Літер менше, ніж величин, і в різних розділах одна літера іноді означає різне. У таблиці нижче кожне значення має свій рядок; у правому стовпці --- формула, де його введено.

\begin{center}
\footnotesize
\setlength{\tabcolsep}{4pt}
\renewcommand{\arraystretch}{1.12}
\begin{longtable}{>{\raggedright}p{2.0cm}>{\raggedright}p{9.6cm}>{\raggedright\arraybackslash}p{2.4cm}}
\toprule
Символ & Значення & Де введено \\
\midrule
\endfirsthead
\toprule
Символ & Значення & Де введено \\
\midrule
\endhead
\bottomrule
\endfoot
$a$ & крутість лінійної передавальної характеристики, $f(u)=au$ & \eqref{eq:feedback} \\
$a$ & рівень \nt{ACET}: $0$ --- читання, $1$ --- запис & \eqref{eq:acetnet} \\
$a_i$, $a$ & утома групи: у генераторі ритму; у гойдалці повільного сну & \eqref{eq:halfcentre}, \eqref{eq:updown} \\
$a_S$, $a_P$ & чутливість ритму серця до газу і до гальма & \eqref{eq:gasbrake} \\
$A(r)$, $A_0$ & відповідь синапсу на імпульс за частоти $r$; у відпочилого синапсу & \eqref{eq:depression} \\
$A_1$ & відповідь отримувача на одну порцію медіатора & \eqref{eq:quantal} \\
$A$ & площа мембрани нейрона разом із дендритами & \eqref{eq:dendriteload} \\
$A_f$, $A_s$ & частка поправки, що зберігається до наступного руху, у швидкого і повільного процесів & \eqref{eq:tworate} \\
$b_i$ & власний приплив ритмоводія & \eqref{eq:seroneuron} \\
$b$ & як швидко робота втомлює групу в генераторі ритму або в гойдалці сну & \eqref{eq:halfcentre}, \eqref{eq:updown} \\
$b$ & кількість біт на один відлік оцифрування & \eqref{eq:probedata} \\
$B_f$, $B_s$ & сила, з якою поправка вчиться на помилці, у швидкого і повільного процесів & \eqref{eq:tworate} \\
$c$ & концентрація медіатора в об'ємі & \eqref{eq:seroneuron} \\
$c$ & коефіцієнт кореляції шумів сусідніх нейронів & \eqref{eq:correlated} \\
$c$ & коефіцієнт, з яким різниця лічильників рухає ракетку & \eqref{eq:dishreadout} \\
$c_f$ & відповідь $C$-нейрона: найбільша з $s_{f,p}$ за всіма положеннями & \eqref{eq:scstage} \\
$c_m$ & питома ємність мембрани & \eqref{eq:dendriteload} \\
$C$ & ціна зусилля: дія за час $\tau_a$ коштує $C/\tau_a$ & \eqref{eq:vigor} \\
$d'$ & розрізненність двох входів: різниця, поділена на шум & \eqref{eq:gainsnr} \\
$d$, $d_j$ & затримка на шляху сигналу & \eqref{eq:glutneuron}, \eqref{eq:vstar} \\
$d$ & затримка петлі зворотного зв'язку з м'язом & \eqref{eq:delaystable} \\
$D$ & коефіцієнт дифузії & \eqref{eq:diffusionlength} \\
$D$ & знаменник у частотах \nt{GLUT}--\nt{GABA}-мережі & \eqref{eq:isn} \\
$D$ & час очікування відкладеної винагороди & \eqref{eq:patience} \\
$e$, $e_{ij}$ & слід у синапсі & \eqref{eq:threefactor} \\
$e$, $e_i$ & помилка виходу або руху, що надходить проводом помилки & \eqref{eq:lms}, \eqref{eq:adaptivefilter}, \eqref{eq:tworate} \\
$E_{\text{імп}}$ & енергія одного імпульсу разом із роботою синапсів & \eqref{eq:energy} \\
$E$ & частота \nt{GLUT}-групи & \eqref{eq:ei} \\
$E_E$ & рівноважний потенціал збуджувального синапсу & \eqref{eq:shunt} \\
$f$, $F$ & передавальна характеристика: частота як функція входу & \eqref{eq:ratenet}, \eqref{eq:gain} \\
$f$, $f_0$ & частота скорочень серця; власний ритм ритмоводія & \eqref{eq:gasbrake} \\
$f_s$ & частота оцифрування одного електрода & \eqref{eq:probedata} \\
$F(t)$, $F_1$ & сила м'яза; сила одного посмикування & \eqref{eq:forcesum} \\
$F_1$, $F_2$ & сили двох м'язів, що тягнуть суглоб у різні боки & \eqref{eq:antagonists} \\
$g$ & чутливість рецептора до концентрації & \eqref{eq:seroneuron} \\
$g_L$, $g_E$, $g_I$ & провідності витоку, збудження і гальмування & \eqref{eq:shunt} \\
$g$ & коефіцієнт підсилення, що його задає \nt{NORA} & \eqref{eq:gain}, \eqref{eq:machine} \\
$g$ & загальне гальмування від \nt{GABA} у мережі пам'яті & \eqref{eq:acetnet} \\
$g$ & підсилення воріт болю, що задається згори & \eqref{eq:gate} \\
$g$ & підсилення одного ступеня гормонального каскаду & \eqref{eq:cascade} \\
$g_j$ & фільтр $j$ зі своєю сталою часу на вході адаптивного фільтра & \eqref{eq:adaptivefilter} \\
$G$ & підсилення петлі зворотного зв'язку & \eqref{eq:feedback} \\
$G$ & швидкість виправлення в петлі із затримкою & \eqref{eq:delaystable} \\
$h$, $h_I$ & зовнішній вхід \nt{GLUT}-групи і \nt{GABA}-групи & \eqref{eq:ei}, \eqref{eq:isn} \\
$h(t)$ & одиночне посмикування м'яза & \eqref{eq:twitch} \\
$H(t)$ & верхній поріг тиску сну: дійшовши до нього, машина засинає & \eqref{eq:sleeppressure} \\
$I$, $I_i$ & зовнішній приплив до нейрона або групи & \eqref{eq:ratenet}, \eqref{eq:wta}, \eqref{eq:updown} \\
$I$ & частота \nt{GABA}-групи & \eqref{eq:ei} \\
$I$ & яскравість або інтенсивність подразника & \eqref{eq:photonnoise}, \eqref{eq:nakarushton} \\
$I$ & вхідний струм, що заряджає мембрану & \eqref{eq:dendriteload} \\
$k$ & кількість імпульсів у пачці & \eqref{eq:burst} \\
$k$ & у скільки разів приріст має перевищувати шум & \eqref{eq:photonnoise} \\
$k$ & коефіцієнт зворотного зв'язку гормонального каскаду & \eqref{eq:cascade} \\
$k_s$ & сила, з якою біль накачує сенсибілізацію & \eqref{eq:sensitization} \\
$K$ & кількість ознак результату & \eqref{eq:vectortd} \\
$K$ & жорсткість суглоба & \eqref{eq:antagonists} \\
$K$ & підсилення петлі гормонального каскаду, $K=g^3k$ & \eqref{eq:cascade}, \eqref{eq:cascadestable} \\
$K_\varphi$ & сила підлаштування добового генератора до зовнішнього сигналу & \eqref{eq:pll} \\
$L$ & довжина лінії затримок & \eqref{eq:itd} \\
$L$ & довжина проводу від датчика болю & \eqref{eq:paintime} \\
$L(t)$ & нижній поріг тиску сну: дійшовши до нього, машина прокидається & \eqref{eq:sleeppressure} \\
$M$ & міра причинного зв'язку провісника і винагороди & \eqref{eq:retro} \\
$M_k$ & очікування ознаки $k$ & \eqref{eq:vectortd} \\
$M_{ik}$ & сила гальмування від \nt{GABA}-нейрона $k$ & \eqref{eq:machine} \\
$M$ & обертальний момент суглоба & \eqref{eq:antagonists} \\
$M$ & кількість вхідних ліній змішувачів & \eqref{eq:cover} \\
$n$, $\bar n$ & кількість імпульсів за вікно; її середнє & \eqref{eq:rateerror} \\
$n$ & кількість входів порогового елемента & \eqref{eq:cover} \\
$n_\uparrow$, $n_\downarrow$ & лічба імпульсів ділянок «вгору» і «вниз» за вікно & \eqref{eq:dishreadout} \\
$N$ & кількість нейронів & \eqref{eq:correlated}, \eqref{eq:energy} \\
$N$ & кількість електродів щупа & \eqref{eq:probedata} \\
$N_s$ & кількість місць викиду між двома нейронами & \eqref{eq:quantal} \\
$p$ & імовірність викиду медіатора на імпульс & \eqref{eq:burst} \\
$p$ & частка активних нейронів у мережі пам'яті & \eqref{eq:acetnet} \\
$p$ & збурення, яке вчаться компенсувати & \eqref{eq:tworate} \\
$P$ & потужність, що витрачається на сигнали & \eqref{eq:energy}, \eqref{eq:dishpower} \\
$P$ & невизначеність збереженої оцінки & \eqref{eq:kalman} \\
$P$ & активність «гальма»: \nt{ACET}-проводів до серця & \eqref{eq:gasbrake} \\
$P_{\max}$ & скільки випадкових образів пороговий елемент може поділити на два класи & \eqref{eq:cover} \\
$P_{\text{пр}}$ & імовірність промаху & \eqref{eq:dishdensity} \\
$q$ & швидкість, з якою змінюється світ & \eqref{eq:kalman} \\
$q_k$ & частота \nt{GABA}-нейрона $k$ & \eqref{eq:machine} \\
$q_0$ & фонова активність гальмівного посередника & \eqref{eq:gate} \\
$q$ & частота гальмівного посередника у воротах болю & \eqref{eq:gate} \\
$q$ & у скільки разів наступна рухова одиниця сильніша за попередню & \eqref{eq:sizeprinciple} \\
$r$, $r_i$ & частота імпульсів нейрона & \eqref{eq:ratenet} \\
$r$, $r_t$ & винагорода (потік винагороди) & \eqref{eq:rpe}, \eqref{eq:ctd} \\
$\mathbf r(t)$ & вектор відгуків резервуара & \eqref{eq:reservoir} \\
$R$, $\bar R$ & отримана винагорода; її очікування або середнє за одиницю часу & \eqref{eq:perturbation}, \eqref{eq:vigor} \\
$R$ & дисперсія шуму на вході фільтра & \eqref{eq:kalman} \\
$R$, $r_0$ & відкладена і негайна винагороди & \eqref{eq:patience} \\
$R_{\max}$ & гранична швидкість передавання, біт/с & \eqref{eq:spikeentropy} \\
$r_{\max}$ & найбільша частота імпульсів & \eqref{eq:gain}, \eqref{eq:nakarushton} \\
$R_{\text{сирий}}$, $R_{\text{імп}}$ & потік даних щупа: сирий і лише імпульси, біт/с & \eqref{eq:probedata} \\
$s_j$ & знак виходів нейрона $j$ & \eqref{eq:dalesign} \\
$s_i$ & знак рецептора отримувача & \eqref{eq:seroneuron} \\
$s$, $\hat s$ & стан світу; його оцінка & \eqref{eq:rpe}, \eqref{eq:kalman} \\
$s_{f,p}$ & відповідь $S$-нейрона на ознаку $f$ у положенні $p$ & \eqref{eq:scstage} \\
$S$ & активність «газу»: \nt{NORA}-проводів до серця & \eqref{eq:gasbrake} \\
$S$ & тиск сну & \eqref{eq:sleeppressure} \\
$t_1$ & момент першого імпульсу & \eqref{eq:latency} \\
$t_k$ & момент $k$-го імпульсу & \eqref{eq:forcesum} \\
$t_{f,i}$ & зразок ознаки $f$ & \eqref{eq:scstage} \\
$T$ & вікно лічби імпульсів; час накопичення фотонів & \eqref{eq:rateerror}, \eqref{eq:photonnoise} \\
$T_c$ & час наростання посмикування м'яза & \eqref{eq:twitch} \\
$T_{\text{залп}}$ & інтервал між залпами культури & \eqref{eq:burstinterval} \\
$u$ & сумарний вхід нейрона & \eqref{eq:gain} \\
$u$, $u(t)$ & команда мозку: на вході адаптивного фільтра; гормональному каскаду & \eqref{eq:adaptivefilter}, \eqref{eq:cascade} \\
$u(t)$ & вхід резервуара & \eqref{eq:reservoir} \\
$U$ & рівень сталого входу & \eqref{eq:latency} \\
$U$ & частка запасу медіатора, що викидається на імпульс & \eqref{eq:depression} \\
$v$ & швидкість сигналу в проводі; швидкість руху плями & \eqref{eq:itd}, \eqref{eq:vstar} \\
$V$ & напруга на мембрані & \eqref{eq:glutneuron}, \eqref{eq:shunt} \\
$V$ & очікування майбутньої винагороди & \eqref{eq:rpe}, \eqref{eq:ctd} \\
$w$, $w_{ij}$, $W_{ij}$ & вага зв'язку & \eqref{eq:glutneuron}, \eqref{eq:ratenet}, \eqref{eq:acetnet}, \eqref{eq:lms}, \eqref{eq:adaptivefilter} \\
$\mathbf w_k$ & ваги входів $C$-нейрона $k$ & \eqref{eq:traceinvariance} \\
$w_{q\mu}$, $w_{q\nu}$, $w_{z\mu}$ & ваги зв'язків у воротах болю & \eqref{eq:gate} \\
$W_{\text{вих}}$ & ваги лінійного зчитування резервуара & \eqref{eq:reservoir} \\
$x$, $y$ & логічні входи і вихід & \eqref{eq:dualrail}, \eqref{eq:veto} \\
$x$, $y$ & активність відправника й отримувача & \eqref{eq:hebb} \\
$x$ & місце детектора на лінії затримок & \eqref{eq:itd} \\
$x_i$ & вхід від органів чуття & \eqref{eq:acetnet} \\
$x_p$ & вхід у положенні $p$ & \eqref{eq:scstage} \\
$\mathbf x$ & виходи $S$-нейронів, вхід $C$-нейрона & \eqref{eq:traceinvariance} \\
$x_j$ & вхід $j$ адаптивного фільтра: команда після фільтра $g_j$ & \eqref{eq:lms}, \eqref{eq:adaptivefilter} \\
$x_f$, $x_s$ & поправки швидкого і повільного процесів & \eqref{eq:tworate} \\
$x_1$, $x_2$, $x_3$ & рівні на трьох ступенях гормонального каскаду; $x_3$ --- кінцевий гормон & \eqref{eq:cascade} \\
$x^*$ & частка відновленого запасу медіатора, за якої починається нова лавина & \eqref{eq:burstinterval} \\
$y$ & зашумлений вхід фільтра & \eqref{eq:kalman} \\
$y$, $\hat y$ & сигнал датчиків; його передбачення адаптивним фільтром & \eqref{eq:adaptivefilter} \\
$y_k$, $\bar y_k$ & активність $C$-нейрона $k$; її слід & \eqref{eq:traceinvariance} \\
$y(t)$ & вихід зчитування резервуара & \eqref{eq:reservoir} \\
$z$ & частота виходу воріт болю & \eqref{eq:gate} \\
\midrule
$\alpha_i^\pm$ & вага додатних і від'ємних помилок & \eqref{eq:asymmetric} \\
$\beta$ & сила взаємного гальмування & \eqref{eq:wta} \\
$\beta$ & сила гальмування виходу воріт болю & \eqref{eq:gate} \\
$\gamma$ & коефіцієнт дисконтування & \eqref{eq:rpe} \\
$\delta(\cdot)$ & дельта-функція: миттєвий стрибок & \eqref{eq:glutneuron} \\
$\delta$, $\delta_t$ & помилка передбачення винагороди & \eqref{eq:rpe}, \eqref{eq:ctd} \\
$\delta t$ & різниця часів приходу звуку у два вуха & \eqref{eq:itd} \\
$\Delta$, $\Delta_{\max}$ & інтервал між імпульсами; вікно збігу & \eqref{eq:window} \\
$\Delta t$ & точність, з якою отримувач розрізняє час & \eqref{eq:spikeentropy} \\
$\Delta F$ & приріст сили від чергової рухової одиниці & \eqref{eq:sizeprinciple} \\
$\Delta I$ & помітний приріст яскравості & \eqref{eq:photonnoise} \\
$\Delta p$ & зсув ракетки за вікно & \eqref{eq:dishreadout} \\
$\Delta u$ & різниця входів двох груп & \eqref{eq:gainsnr} \\
$\Delta V$ & на скільки треба підняти напругу мембрани & \eqref{eq:dendriteload} \\
$\Delta x$ & відстань між сусідніми датчиками & \eqref{eq:vstar} \\
$\varepsilon$ & відносна різниця входів & \eqref{eq:wtagain} \\
$\eta$ & швидкість навчання & \eqref{eq:plasticity}, \eqref{eq:lms}, \eqref{eq:traceinvariance} \\
$\eta$ & шум у мережі після підсилювача & \eqref{eq:softmax} \\
$\mu$ & вхід дотику до воріт болю & \eqref{eq:gate} \\
$\nu$ & вхід болю & \eqref{eq:gate} \\
$\theta$ & поріг спрацювання & \eqref{eq:window}, \eqref{eq:scstage} \\
$\kappa$ & кількість медіатора на один імпульс & \eqref{eq:seroneuron} \\
$\kappa_i$ & відношення ваг додатних і від'ємних помилок & \eqref{eq:expectile} \\
$\kappa_m$ & зв'язок сили м'яза з його жорсткістю & \eqref{eq:antagonists} \\
$\ell$ & відстань, на яку розпливається медіатор & \eqref{eq:diffusionlength} \\
$\ell$ & плече, на якому м'яз тягне суглоб & \eqref{eq:antagonists} \\
$\xi$ & випадкове тремтіння виходу нейрона & \eqref{eq:perturbation} \\
$\rho(\boldsymbol\psi)$ & частка часу, проведеного в налаштуванні $\boldsymbol\psi$ & \eqref{eq:dishdensity} \\
$\sigma$ & розкид, масштаб шуму & \eqref{eq:rateerror}, \eqref{eq:softmax} \\
$\sigma$ & яскравість, за якої відповідь дорівнює половині найбільшої & \eqref{eq:nakarushton} \\
$\sigma_{\text{до}}$, $\sigma_{\text{після}}$ & шум до підсилювача і після нього & \eqref{eq:gainsnr} \\
$\tau$ & стала часу мембрани & \eqref{eq:glutneuron} \\
$\tau$ & стала часу ступеня гормонального каскаду & \eqref{eq:cascade} \\
$\tau_{\text{еф}}$ & стала часу групи, що збуджує сама себе & \eqref{eq:perfectint} \\
$\tau_c$ & час життя медіатора в об'ємі & \eqref{eq:seroneuron} \\
$\tau_E$, $\tau_I$ & сталі часу \nt{GLUT}- і \nt{GABA}-групи & \eqref{eq:ei} \\
$\tau_e$ & час життя сліду & \eqref{eq:threefactor} \\
$\tau_{\text{сл}}$ & час життя сліду активності $C$-нейрона & \eqref{eq:traceinvariance} \\
$\tau_s$ & час повернення чутливості після пошкодження & \eqref{eq:sensitization} \\
$\tau_r$, $\tau_d$ & час накопичення тиску сну під час неспання і його скидання уві сні & \eqref{eq:sleeppressure} \\
$\tau_{\text{відн}}$ & час відновлення запасу медіатора & \eqref{eq:depression} \\
$\tau_\gamma$ & горизонт планування & \eqref{eq:ctd} \\
$\tau_v$ & швидкість, з якою уточнюється очікування & \eqref{eq:ramp} \\
$\tau_a$ & час утоми в генераторі ритму або в гойдалці сну & \eqref{eq:halfcentre}, \eqref{eq:updown} \\
$\tau_a^*$ & найкращий час виконання дії & \eqref{eq:vigor} \\
$\Phi$ & права частина правила навчання & \eqref{eq:plasticity} \\
$\Phi$ & потік фотонів & \eqref{eq:photonnoise} \\
$\varphi_k$ & ознака $k$ отриманого результату & \eqref{eq:vectortd} \\
$\varphi$ & різниця фаз добового генератора і зовнішнього сигналу & \eqref{eq:pll} \\
$\boldsymbol\psi$ & налаштування мережі в моделі DishBrain & \eqref{eq:dishdensity} \\
$\omega$, $\omega_0$ & частота зовнішнього сигналу (світла); власна частота добового генератора & \eqref{eq:pll} \\
\end{longtable}
\end{center}

\chapter{Відповіді}
\label{sec:answers}

Короткі відповіді до задач наприкінці розділів. Повні розв'язки --- в окремому посібнику для викладачів.

\answer{ex:01:1}{(а)~$8\cdot10^{10}$ \nt{GLUT} --- $10$ населень Землі; широкомовних $\approx1.9\cdot10^6$ --- місто завбільшки з Відень. (б)~$\approx8\cdot10^{14}$ закінчень, із них широкомовних $\approx3\cdot10^{11}$, частка $\approx4\cdot10^{-4}$ --- у $\approx16$ разів більша за частку цих нейронів ($2.2\cdot10^{-5}$).}
\answer{ex:01:2}{(а)~$\tau_c\ln(c_0/c_*)\approx4.6\ms$. (б)~Лінія вільна при $T\ge\tau_c\ln101$: до $\approx22$ викидів за секунду замість $\approx220$.}
\answer{ex:01:3}{$V(s_k)=\gamma^{K-1-k}r$; $\delta_{\text{лампа}}=\gamma^Kr=re^{-T/\tau_\gamma}\approx0.60\,r$ за будь-якого $\Delta$, $\tau_\gamma\approx1.95\,\text{с}$.}
\answer{ex:01:4}{$n^*=93$, $47$, $25$, $12$: $n^*\approx5/\alpha$ при малих $\alpha$.}
\answer{ex:01:5}{\nt{GLUT}: $1\to10\to10^4\to10^7$ --- $\approx10^7$ нейронів, $4$ ланки, $4\ms$. \nt{DOPA}: $1\to10\to3.2\cdot10^5$ --- $\approx3\cdot10^5$ нейронів, $3$ ланки, у $30$ разів менше; той самий порядок, що й кількість \nt{DOPA}-нейронів.}
\answer{ex:01:6}{$\lambda=f_EJ_E-f_IJ_I$, звідки $J_I=37.5$; відношення струмів $f_IJ_I/f_EJ_E=0.94$; лавина ($\lambda\ge1$) при послабленні гальмування всього на $6.7\,\%$.}
\answer{ex:01:7}{$n\ge57$ пред'явлень на кожну затримку в кожному стані. Такий самий спад дала б, наприклад, неточність внутрішньої оцінки часу, що зростає з $T$.}

\answer{ex:02:1}{Ряд завдовжки $3.2$~мм із кроком $25\um$: $\approx129$ детекторів. Спрацьовує смуга завширшки $v\,\tau\ln2=1.7$~мм --- $\approx69$ детекторів, більше ніж половина ряду; напрямок --- центр смуги.}
\answer{ex:02:2}{Вікно $-\tau\ln\frac{w_2}{\theta-w_1}\le\delta\le\tau\ln\frac{w_1}{\theta-w_2}$, тобто $[-0.235,\,0.178]\ms$; центр $c=\frac\tau2\ln\frac{w_1(\theta-w_1)}{w_2(\theta-w_2)}=-28$~мкс. Напрямок зсувається на $28$~мкс у бік сильнішого вуха --- майже три кроки роздільності.}
\answer{ex:02:3}{Потрібні $s_i$ з $s_iw_{ij}s_j\ge0$; вони існують тоді й лише тоді, коли цикли парні. Взаємне гальмування зводиться (стійких коливань немає), петля \nt{GLUT}--\nt{GABA} --- ні (може коливатися).}
\answer{ex:02:4}{Шість нейронів $g_k=[x+y+z\ge k]$ і $\bar g_k=[\bar x+\bar y+\bar z\ge4-k]$, $k=1,2,3$; $C=g_2$, $\bar C=\bar g_2$, $S=[g_1+\bar g_2+g_3\ge2]$, $\bar S=[\bar g_1+g_2+\bar g_3\ge2]$: вісім нейронів. З І та АБО --- $12$.}
\answer{ex:02:5}{$T_{\min}(0.6)=0.718\,\tau$; $I_c=0.225$, біля порога $T_{\min}\propto(I-I_c)^{-1/2}$.}
\answer{ex:02:6}{(а)~$500$ ланок, $5\cdot10^4$ нейронів; розкид $4.5\ms$ ($0.45\,\%$), відносна похибка $\propto T^{-1/2}$. (б)~Спільний для всіх ланок випадковий множник затримки з розкидом $5\,\%$.}
\answer{ex:02:7}{Освіжати раз на $\tau_F\ln2=0.69\,\text{с}$: $4.3$ імпульсу за секунду на нейрон проти $20$, у $4.6$ раза ощадливіше ($4\cdot10^{-7}$ проти $2\cdot10^{-6}$~Дж при $10^{-10}$~Дж на імпульс). Тригер утрачає біт; конденсатор зберігає його, якщо глушіння почалося не пізніше ніж через $0.49\,\text{с}$ після освіження (імовірність $0.71$).}

\answer{ex:03:1}{$35\to17.5\mV$ при $g_E=g_L$; $56\to40\mV$ при $g_E=4g_L$, віднімання дало б $38.5\mV$: шунт ділить $g_E$ на $3$. Вікно звужується вдвічі ($\tau_{\text{еф}}$ падає у $2$ рази).}
\answer{ex:03:2}{$\operatorname{tr}=\frac{w_{EE}-1}{\tau_E}-\frac1{\tau_I}$, $\det=\frac{1-w_{EE}+w_{EI}w_{IE}}{\tau_E\tau_I}$; на межі $\omega=\sqrt{\det}$, $\tau_I=20\ms$, період $73\ms$.}
\answer{ex:03:3}{$d_c=\tau_E\arccos(-1/G)/\sqrt{G^2-1}$, період $2\pi\tau_E/\sqrt{G^2-1}\in(2d_c,4d_c)$. $G=2$: $d_c=12.1\ms$, період $36\ms$; $G=5$: $d_c=3.6\ms$, період $12.8\ms$ --- якраз смуга швидких ритмів.}
\answer{ex:03:4}{(а)~При $I_2=\beta I_1=2$ угору й $I_2=I_1/\beta=0.5$ униз, петля завширшки $1.5$. (б)~На $\tau\ln10/(\beta-1)=2.3\tau$ на кожен десятковий порядок $\varepsilon$.}
\answer{ex:03:5}{Три посередники, $d_k=5$, $10$, $15\ms$: заборони мають покрити $15\ms$, по $5\ms$ на кожну.}
\answer{ex:03:6}{$n+M+1=3+4+1=8$ нейронів. Два: \nt{GABA} $=[x+y+z\ge2]$, вихід $=[x+y+z-2\,\nt{GABA}\ge1]$.}
\answer{ex:03:7}{$\binom84=70<100\le\binom94=126$: $9$ посередників (теорема Шпернера). Без прямих зв'язків --- $100$: ранг $\beta(J-I)$ дорівнює $n$.}
\answer{ex:03:8}{$h_c=\frac1{2w_{EE}}+\frac{w_{EI}w_{IE}-w_{EE}}{4w_{EE}^2}=\frac7{18}\approx0.39$; моделювання дає зміну знака між $h=0.38$ і $0.40$.}

\answer{ex:04:1}{(а)~$\approx1.1\cdot10^{11}$ біт/Дж. (б)~$\log_2\sqrt{10}\approx1.7$ біта проти $81$: у $\approx50$ разів.}
\answer{ex:04:2}{$r^*=(P_0/E_{\text{імп}})\ln\bigl(1/(r^*\Delta t)\bigr)$, $P_0/E_{\text{імп}}=1.16\,\text{с}^{-1}$: $r^*\approx6$ імпульсів за секунду, $7.4\cdot10^{10}$ біт/Дж.}
\answer{ex:04:3}{$\mathcal I=\frac{N}{1+(N-1)c}=9.9$ (границя $10$) проти $\mathcal I=\frac{N}{1-c}=1111$: у $\approx110$ разів; кореляція шкідлива, лише коли схожа на сигнал.}
\answer{ex:04:4}{$\theta\approx68.7$ і $19.9$ (в одиницях $w$), $m=19$ і $m=10$: швидкому отримувачеві потрібно майже вдвічі менше збіжних входів, хоча його середній вхід уп'ятеро менший.}
\answer{ex:04:5}{$U\tau_{\text{відн}}\ge0.725\,\text{с}$ і $\tau_{\text{відн}}(1-2U)\le0.05\,\text{с}$, що вимагає $U\ge0.483$; найменше $\tau_{\text{відн}}=0.725\,\text{с}$ при $U=1$ ($1.45\,\text{с}$ при $U=0.5$).}
\answer{ex:04:6}{(а)~$\theta\,e/(e-1)\approx1.58\,\theta$ за будь-яких $\tau$ і $\theta$. (б)~$14$ імпульсів.}
\answer{ex:04:7}{(а)~$1.62$ біта на слово: $0.77$ --- кількість імпульсів, $0.84$ --- їхнє розташування. (б)~$1.13$ і $1.59$ біта: за $30$ словами оцінка втрачає майже третину.}

\answer{ex:05:1}{$\ell\approx17\um$; $\approx100$~діб: розсилку по мозку робить розгалуження проводу, дифузія лише розмиває кожне місце викиду на $\sim20\um$.}
\answer{ex:05:2}{$r=ab/(1+G)$, $S=1/(1+G)$ точно; $+1.7\,\%$ замість $+20\,\%$.}
\answer{ex:05:3}{$T^*=1.21\,\text{с}$ при $G=2$ і $0.19\,\text{с}$ при $G=9$: реально $G\sim2$--$4$, пригнічення в $3$--$5$ разів.}
\answer{ex:05:4}{$\mathrm{CV}=1/\sqrt{2\lambda\tau_c}\approx9\cdot10^{-4}$ за сумарної частоти $\lambda$; одному нейрону потрібно $\tau_c=12.5\,\text{с}$.}
\answer{ex:05:5}{$c\approx0.40$ в обох випадках; $\mathrm{CV}=0.050$ проти $0.158$, у $\sqrt{10}$ разів. Частоти окремих нейронів лишаються пропорційними $a_i$: регулюється лише сума.}
\answer{ex:05:6}{$N_1=\overline{A\vee B}$, $N_2=\overline{A\vee N_1}$, $N_3=\overline{B\vee N_1}$, $N_4=\overline{N_2\vee N_3}$, XOR $=N_5=\overline{N_4}$; кожне перемикання займає $\tau_c\ln2$, шлях із $4$ вентилів --- $2.8\,\text{с}$ на один XOR.}
\answer{ex:05:7}{(а)~$\tau_\gamma=6/\ln3\approx5.5\,\text{с}$. (б)~$\bar D<4\,\text{с}$ проти $D<2.2\,\text{с}$: непередбачуваність затримки робить терплячішим.}
\answer{ex:05:8}{$k_2=1/\tau_d=10\,\text{с}^{-1}$, $k_1=1/\tau_d-1/\tau_\gamma=9.5\,\text{с}^{-1}$; $\tau_\gamma=1/(k_2-mk_1)$: подвоєння при $+2.6\,\%$ (а при $+5.3\,\%$ горизонт нескінченний).}

\answer{ex:06:1}{$e=(1-e^{-T_a/\tau_e})e^{-d/\tau_e}$. (а)~$0.110$ проти $4.5\cdot10^{-5}$, у $\approx2.5\cdot10^3$ разів. (б)~$\tau_e^*=T_a/\ln(1+T_a/d)=0.59\,\text{с}$.}
\answer{ex:06:2}{Швидкість $\nu_x\nu_y(A_+\tau_+-A_-\tau_-)=0.8A_+$ на секунду, $\approx2900A_+$ за годину; $\tau_-=\tau_+/0.6=33\ms$.}
\answer{ex:06:3}{$\mathbf e_1$ при $\mu^2+s_1^2>s_2^2$; тут $1.01>0.25$, і нейрон вивчить $\mathbf e_1$, хоча розкид уздовж $\mathbf e_2$ у $25$ разів більший: правило знаходить головний вектор матриці кореляцій, а не коваріацій.}
\answer{ex:06:4}{$V=Re^{-(t_c+L-t)/\tau_\gamma}$ між лампочкою і соком, тож $\dot V=V/\tau_\gamma$; сплеск площею $Re^{-L/\tau_\gamma}$: $0.61R$ і $0.78R$.}
\answer{ex:06:5}{Відношення сигнал/шум $1/(N+1)$, спроб $\propto N+1$: $5\cdot10^5$ спроб $\approx1$~міс. і $5\cdot10^{11}$ спроб $\approx8\cdot10^4$~років.}
\answer{ex:06:6}{(а)~$k_2=1/\tau_d$, $k_1=1/\tau_d-1/\tau_\gamma$. (б)~Хибна помилка $-(V/\tau_\gamma)\,\varepsilon/(1+\varepsilon)$, $\varepsilon=\tau_d/\tau_\gamma$, звідки $\tau_d\le0.105\,\text{с}$.}
\answer{ex:06:7}{$0.46$, $0.90$, $0.99$, $0.95$ і $0.70$ при $d=3\,\text{с}$. Точки лягають на одну криву від частки сліду $F=(1-e^{-T_{\text{cue}}/\tau_e})e^{-d/\tau_e}$: навчання при $F\gtrsim0.1$, випадковий вибір при $F<10^{-2}$.}
\answer{ex:06:8}{$\eta^*=1/\bigl(2\sigma^2(N+2)\bigr)$, за $N+2$ спроби; при $m$ тремтячих з $N$ --- за $N(m+2)/m\ge N+2$: розрідженість не допомагає.}

\answer{ex:07:1}{$V=\kappa p/\bigl(\kappa p+(1-\kappa)(1-p)\bigr)$: $0.059$, $0.20$, $0.50$; медіана дорівнює $0$, а точка~\eqref{eq:expectile} змінюється з $\kappa$ плавно.}
\answer{ex:07:2}{$p=1/3$, $x=0.75$, $\sigma_r=0.35$, $V(0.2)=0.083$.}
\answer{ex:07:3}{$V=\sqrt\kappa/(\sqrt\kappa+\sqrt{1-\kappa})$, від $0.25$ до $0.75$; розкид $\propto\sqrt\eta$; у крапель $V=0.25+0.5\kappa$, розподіли різняться на $0.05$ у крайніх нейронів, потрібно $\eta\lesssim0.01$.}
\answer{ex:07:4}{(а)~$J^*=2\sqrt{C\bar R}$. (б)~Переплата $\bigl(k^{1/2}+k^{-1/2}\bigr)/2-1$: $6\,\%$ при $k=2$ і $25\,\%$ при $k=4$ --- досить точності «в межах двох разів».}
\answer{ex:07:5}{Сплеск площею $R(e^{-1}-e^{-2})\approx0.23R$ --- як $2.3\,\text{с}$ рампи; якщо дофамін повідомляє $V$, сплеску немає, рівень просто сходинкою зростає в $e$ разів.}
\answer{ex:07:6}{Подія зсуває вагу $\propto A\tau_f/(\tau_e+\tau_f)$, рівень дає помилку $s\tau_f$: $9\,\text{с}\le\tau_f\le17\,\text{с}$.}
\answer{ex:07:7}{Вихідно $\Delta P=0.80$, $M=0.90$. (а)~$\Delta P=0.74$ ($-8\,\%$), $M=0.43$ ($-52\,\%$); (б)~$\Delta P=0.40$ ($-50\,\%$), $M=0.80$ ($-11\,\%$): подвійна дисоціація.}
\answer{ex:07:8}{$N_{\text{еф}}\to2+\rho^2N\approx10^6$ спроб: у $10^4$ разів менше, ніж з одним проводом, але витік в $1\,\%$ усе ще коштує мільйона спроб.}

\answer{ex:08:1}{(а)~$\approx1.7\cdot10^6$. (б)~$g=\frac{0.9}{\sqrt{0.19}}\,\sigma_{\text{після}}/\sigma_{\text{до}}\approx16.5$: відповідь залежить лише від відношення шумів.}
\answer{ex:08:2}{(а)~Власні числа $-(1\pm g\beta)/\tau$; різниця загасає за $\tau/(1-g\beta)$: $40\ms$ при $g\beta=0.5$ (різниця входів підсилена втричі), $200\ms$ при $0.9$ (у $19$ разів). (б)~$0.905$ і $1.105$; між межами перемагає лише сильний вхід.}
\answer{ex:08:3}{$P(k)=e^{gu_k/\sigma}/\sum_le^{gu_l/\sigma}$, тобто~\eqref{eq:softmax}; $g=10\ln9\approx22$.}
\answer{ex:08:4}{$\bar\Delta=1/3$, $g^*\approx3\ln(1/h)$: $21$, $14$, $9$ проти $16$--$23$, $11$ і $8$ у моделі.}
\answer{ex:08:5}{$g^*$ від $16$--$23$ при $h=0.001$ до $6$--$8$ при $h=0.2$; оцінка $3\ln(1/h)$ правильна в межах півтора раза при $h\le0.05$; при $h\gtrsim\eta/10$ правило $Q\leftarrow Q+\eta(R-Q)$ не встигає засвоїти здобуте пошуком, і $g^*$ застрягає близько $8$.}
\answer{ex:08:6}{$T_p=\tau\ln9=44\ms$ при $g_{\text{lo}}=0$ (модель $44\ms$) і $70\ms$ при $g_{\text{lo}}=0.25$: пачка в $2$--$3\tau$ за підсилення близько нуля.}
\answer{ex:08:7}{(а)~Найкраще стале $0.925$ ($g=8$), оракул $0.929$ ($16/8$) --- вигравати майже нічого. (б)~Наприклад, спокійні епохи з рідкісними винагородами ($p\in[0,0.3]$) і мінливі з $h=0.1$: $0.850$ проти $0.872$; підлаштування забирає близько чверті при $\eta=0.2$ і майже все при $\eta=0.05$.}

\answer{ex:09:1}{(а)~$P_\infty=0.2$, пам'ять $20\,\text{с}$. (б)~$P_\infty=1$, пам'ять $1\,\text{с}$. (в)~Пам'ять пропорційна амплітуді шуму: ушестеро.}
\answer{ex:09:2}{$P(t)=P_\infty\coth(t\sqrt{q/R})$. (а)~$\frac12\sqrt{R/q}=5\,\text{с}$ --- удвічі коротше за пам'ять. (б)~$t=\frac12\ln21\,\sqrt{R/q}\approx15\,\text{с}$: стільки має триматися підвищений \nt{ACET}.}
\answer{ex:09:3}{Дисперсія $qT/2+R/(2T)$; $T^*=\sqrt{R/q}$, дисперсія $\sqrt{qR}=P_\infty$, як у фільтра~\eqref{eq:kalman}; зростання в $(\kappa+1/\kappa)/2$ раза: $1.25$ і $5.05$.}
\answer{ex:09:4}{$a>a_w=1-\theta/(mw)$: $0.15$ при $w=0.041$, $0.22$ при $w=0.045$.}
\answer{ex:09:5}{Помилки з'являються при $M\approx40$--$60$, при $M=80$ --- $1.9$ зайвого нейрона, при $M=100$ --- частка образу $0.86$; завада від чужих образів має дисперсію $\approx(M-1)p(1-p)/N$, тож $M_{\max}\approx80$, $M_{\max}\propto N/p$.}
\answer{ex:09:6}{Наприклад, $\eta(a)=\eta\min\bigl(1,\max(0,(a-0.3)/0.6)\bigr)$. З $\eta a$ усі три чужі нейрони ввійшли в образ; з $\eta(a)$ --- жоден, запис при $a\ge0.9$ той самий (частка $1.00$).}
\answer{ex:09:7}{(а)~$40$ повторів; з $\eta(a)$ задачі~\ref{ex:09:6} --- ніколи. (б)~$200$ повторів, $10$ ночей.}

\answer{ex:10:1}{Шина $\approx10^{12}$ біт/с ($11.4$ біта на імпульс), керування $\approx40$ біт/с; $2.5\cdot10^{10}\,\text{с}$ --- близько восьмисот років.}
\answer{ex:10:2}{$a>1/3$ (при $a=0$ --- наростаючі коливання $\approx67$~Гц). Ні: стійкість задає добуток $g(1-a)$.}
\answer{ex:10:3}{Середнє $\sigma^2R'x_j$, дисперсія $(N+1)\,x_j^2\sigma^4R'^2$; відношення $1/(N+1)$, потрібно $n\approx z^2(N+1)\propto N$ спроб.}
\answer{ex:10:4}{$L_v\in[32.9,\,47.9]\ms$: затримка $35.4\ms$, вікно $\pm7.5\ms$; хибних зв'язувань $2\Delta_{\max}r_ar_v\approx0.38$ за секунду.}
\answer{ex:10:5}{$0.098$ проти $0.180$ за найкращого $\alpha=0.2$: на $45\,\%$ менша.}
\answer{ex:10:6}{$n=\ln\bigl(\ln1.5/(0.6g)\bigr)/\ln(1-\alpha)$: $24$ і $4$; при $g=2$ --- $11$ і $2$.}
\answer{ex:10:7}{$S'=S\bigl(1-4\eta\sigma^2+4\eta^2\sigma^4(G+2)\bigr)$, найкраще $\eta\sigma^2=1/(2(G+2))$; $\approx10^6$ спроб --- близько місяця.}

\answer{ex:11:1}{(а)~$T=k^2/\bigl(\Phi(\Delta I/I)^2\bigr)=9\,\text{с}$. (б)~$45$ датчиків, пляма $\approx7\times7$, роздільність гірша в $\approx7$ разів.}
\answer{ex:11:2}{(а)~$2$--$3.3$ біта на імпульс, $22$--$34\,\%$ межі ($9.1$--$9.8$ біта). (б)~$\log_2\binom{400}{20}\approx111$ бітів; у середньому $1$ спільний тип.}
\answer{ex:11:3}{(а)~$\sigma/9\le I\le9\sigma$, $1.9$ порядку. (б)~$6$ і $7$.}
\answer{ex:11:4}{(а)~$f<343/0.44\approx780$~Гц. (б)~$625$ імпульсів, $125$ нейронів.}
\answer{ex:11:5}{$21$ біт на подію, $4.8\cdot10^5$ подій/с, $7$ подій на піксель краю: $136$ пікселів/с. Звичайна камера --- $1.25$ кадру/с.}
\answer{ex:11:6}{У логарифмах $\approx1.0\,\text{с}$ в обидва боки (теорія $0.96\tau$); лінійно $0.2\,\text{с}$ угору і $3.0\,\text{с}$ униз (теорія $0.18\tau$ і $3.0\tau$).}
\answer{ex:11:7}{(а)~$\ln I$ розподілений логістично з медіаною $\ln\sigma$ і масштабом $1$: розкид $\pi/\sqrt3$ у натуральних логарифмах, $0.79$ порядку. (б)~$r=r_{\max}\Phi_I(I)^2$.}

\answer{ex:12:1}{$0.60$ (без кореляції було б $0.18$), межа $\sqrt{c/(rT)}\approx0.58$: сто нейронів розрізняють лише «є/немає».}
\answer{ex:12:2}{(а)~$10^8$ імпульсів, $\approx10\,\text{мДж}$: $0.3\,\%$ від $3$~Дж усього мозку. (б)~$3$~Дж, у $\approx300$ разів більше.}
\answer{ex:12:3}{(а)~$\theta_A\in[2,3)$, $\theta_B\in[1,2)$; чужа фігура дає не більше за $2$ і $1$. (б)~$10\cdot96^2\approx9.2\cdot10^4$.}
\answer{ex:12:4}{(а)~$T_d=0.85\,\tau_a=0.34\,\text{с}$. (б)~$T_d\propto\tau_a$ і не залежить від $I$: $\tau_a\approx3.5\,\text{с}$ (моделювання: $3.1\,\text{с}$).}
\answer{ex:12:5}{$p\approx0.17$. Частка падає повільно: $0.90$, $0.88$, $0.86$, $0.84$, $0.82$, $0.79$ при $N=8$, $12$, $24$, $48$, $96$, $192$.}
\answer{ex:12:6}{Усі десять прогонів безпомилкові при $\tau_{\text{сл}}=20$, $50$ і $200\ms$ (при $30\ms$ один прогін із десяти зірвався); при $5$--$10\ms$ --- близько $0.5$, при $0.5$--$2\,\text{с}$ якість падає до $0.86$. Згори: слід має згаснути за паузу, $e^{-0.3\,\text{с}/\tau_{\text{сл}}}\ll1$.}
\answer{ex:12:7}{Однієї затримки замало: вікно $\pm5\ms$ не покриває $\Delta x/v$ від $5$ до $20\ms$. Дві гілки гальмування з $5<d_1\le10\ms$ і $15\le d_2\le d_1+10\ms$.}
\answer{ex:12:8}{\texttt{two\_stage}: $\Delta=2$, нічия від одного перевертання, зміна відповіді --- від двох. Лічильник одиниць: одне перевертання; $0.57$ при $p=0.1$.}

\answer{ex:13:1}{(а)~$q=100^{1/99}\approx1.048$, крок $\approx4.8\%$; діапазон $\approx2.2\cdot10^3$ ($67$~дБ). (б)~Крок $22f_1$, тобто $220\%$ за $10f_1$.}
\answer{ex:13:2}{$P_{\text{збою}}\approx1.8\cdot10^{-4}$ і $2.8\cdot10^{-16}$: близько $300$ збоїв на добу при $S=2$ і $5\cdot10^{-10}$ на добу при $S=4$.}
\answer{ex:13:3}{(а)~$H(s)=eF_1T_c/(1+sT_c)^2$, фільтр другого порядку з частотою зрізу $1/(2\pi T_c)\approx3.2$~Гц. (б)~$r\approx22$~імп/с (за першою гармонікою $\approx20$).}
\answer{ex:13:4}{(а)~При $Gd=1/e$ корінь $s=-1/d$ подвійний, і за жодного іншого $G$ найправіший корінь не лежить лівіше. (б)~$d=30\ms$: $G\approx12$~с$^{-1}$, $\tau=30\ms$; $d=150\ms$: $G\approx2.5$~с$^{-1}$, $\tau=150\ms$ --- стала часу дорівнює затримці.}
\answer{ex:13:5}{(а)~$E\approx0.427$, $T\approx1.70\,\tau_a=0.68$~с (модель з $\tau=20\ms$ дає $0.84$~с). (б)~Рівняння не містить $I$; ритм існує при $b>\beta-1$.}
\answer{ex:13:6}{При $b=0.8$ ритму немає; при $b=1.05$, $1.2$, $1.5$, $2$, $4$ $T\approx2.73$, $1.74$, $1.19$, $0.84$, $0.45$~с; за будь-якого $I$ $T=0.840$~с: вхід змінює лише розмах, а не темп.}
\answer{ex:13:7}{Підстановка $s=u-a$ дає найбільшу швидкість загасання $a+1/d$ при $G=e^{-1-ad}/d$; $a\ge33.3$~с$^{-1}$, $G=e^{-2}/d\approx4.5$~с$^{-1}$; що сильніше спільне напруження (більше $a$), то менше треба брати $G$.}
\answer{ex:13:8}{Відкрита задача. Наприклад, поріг у накачуванні втоми $b[r_i-r_0]_+$ дає $I_{\min}=\beta b r_0/(b-\beta+1)\approx0.53$ при $b=4$, $r_0=0.2$ і період від $1.8$~с при $I=0.6$ до $0.52$~с при $I=3$.}

\answer{ex:14:1}{(а)~На $0.17$~с і на $1.5$~с. (б)~З відстані, меншої за $11$~см.}
\answer{ex:14:2}{Дотик без болю не проходить за жодного $\mu$, доки $g\le\beta w_{q\mu}/w_{z\mu}=5$; дотик $\mu=1$ проходить при $g>6$: сильна сенсибілізація робить болем звичайний дотик.}
\answer{ex:14:3}{(а)~$g^*=(1-k_sC)/(1-k_sA)$, стійка при $k_sA<1$. (б)~Без дотику при $\nu\ge2$, з дотиком --- уже при $\nu\ge1.7$.}
\answer{ex:14:4}{$V(t)=-Pe^{-(T-t)/\tau_\gamma}$. (а)~$-Pe^{-T/\tau_\gamma}\approx-0.37P$. (б)~$+Pe^{-(T-t_e)/\tau_\gamma}\approx+0.61P$; він зростає з $t_e$ до $+P$ і навчає уникати болю.}
\answer{ex:14:5}{При $k_s=4$ защіпання немає (другий стійкий стан є при $k_s\ge4.58$); $T_{\min}\approx4.35$, $1.40$, $0.65$, $0.32\,\tau_s$ при $k_s=4.6$, $5$, $6$, $8$.}
\answer{ex:14:6}{$w_{q\nu}=4/9\approx0.444$, $q_0=2/9\approx0.222$, $w_{q\mu}=49/45\approx1.089$; дотик без болю безпечний до $g=49/9\approx5.4$.}
\answer{ex:14:7}{(а)~Поріг $\nu_c(g)=\theta/g$ при $\nu_c\ge q_0/w_{q\nu}$, інакше $(\theta+\beta q_0)/(g+\beta w_{q\nu})$; $g^*\approx2.45$, $1.0$, $0.25$. (б)~Відкрите питання.}

\answer{ex:15:1}{(а)~$2\cdot10^5$ на нейрон, $6\cdot10^{12}$ на схему. (б)~$\approx8$~біт/с на провід, не менше за $2.5\cdot10^4$~с $\approx7$~год.}
\answer{ex:15:2}{Крок $2$: сусідня кореляція $0.943$, $\lambda$ від $4.18$ до $7.3\cdot10^{-4}$, відношення $\approx5.7\cdot10^3$; крок $4$: кореляція $0.8$, відношення $\approx94$.}
\answer{ex:15:3}{(а)~$\lambda=0.988$ і $0.808$: $82$ і $4.7$ руху. (б)~$x_s^*=0.690$, $x_f^*=0.172$, сума $0.862$.}
\answer{ex:15:4}{(а)~$116$ рухів. (б)~$19$. Модель з одним процесом за нульової суми економії не дає.}
\answer{ex:15:5}{(а)~$G=50$~с$^{-1}$ проти межі $10.5$~с$^{-1}$ без моделі. (б)~Ні: при $G=50$~с$^{-1}$ критична затримка $d_c\approx103\ms$; при $d=150\ms$ потрібно $\hat k>0.445k$ (за будь-яких $G$ і $d$ --- $\hat k>k/2$).}
\answer{ex:15:6}{Помилка падає нижче від $0.5\%$ за $6$--$30$~с, дно $\approx(6$--$9)\cdot10^{-4}$ за найкращих ваг $7.5\cdot10^{-4}$; ваги при $\eta=50$ ще відрізняються від найкращих на величину до $0.2$ вздовж погано обумовлених напрямків $C$ (задача~\ref{ex:15:2}).}
\answer{ex:15:7}{$\mathbf B=A\mathbf K$; $q_f/R\approx0.035$, $q_s/R\approx8.6\cdot10^{-4}$; при $4q_s$ $B_f\approx0.088$, $B_s\approx0.047$: повільний процес учиться більш ніж удвічі швидше, а швидкий --- повільніше.}

\answer{ex:16:1}{Кров --- RC-коло з $\tau=t_{1/2}/\ln2=14.4$~хв; відношення максимуму до мінімуму $2^{T/t_{1/2}}=64$; основну гармоніку ослаблено до $0.55$; відношення $2$ за $t_{1/2}=T=60$~хв.}
\answer{ex:16:2}{$\varphi^*=\arcsin\bigl((\omega-\omega_0)/K_\varphi\bigr)\approx41$~хв ($\omega-\omega_0\approx0.047$~рад/доб), час релаксації $1/(K_\varphi\cos\varphi^*)\approx3.9$~доби; після зсуву на $+6$ і $-6$~год --- $7.7$ і $7.9$~доби.}
\answer{ex:16:3}{$K_c(n)=\sec^n(\pi/n)$: $8$ і $2.37$, похибка $11\%$ і $30\%$; за $n\to\infty$ $K_c\to1$, похибка $\to50\%$.}
\answer{ex:16:4}{Крайні команди (газ $80$, гальмо $0$) дають $f=120$ через $0.76$~с; «лише газ» --- $2.33$~с, утричі довше.}
\answer{ex:16:5}{$3\arctan(\omega\tau)+\omega d=\pi$, $K_c=(1+\omega^2\tau^2)^{3/2}$: $K_c=8$, $3.54$, $2.50$, $1.76$ і період $3.63$, $5.46$, $6.86$, $9.27\,\tau$; за $0.9K_c$ коливання згасають, за $1.1K_c$ усталюються.}
\answer{ex:16:6}{$0.60\bar c$, $0.87\bar c$, $0.90\bar c$, границя $0.91\bar c$; за сталої концентрації відгук нульовий: такий отримувач читає лише порції.}
\answer{ex:16:7}{Відкрита задача. За зворотного зв'язку ритм є лише при $K>8$, період пропорційний $\tau$ ступенів і майже не залежить від $K$ ($3.64\,\tau$ за $K=8.5$, $4.23\,\tau$ за $K=40$): зміна $K$ у півтора раза зсуває період на $2$--$4\%$, менше за похибку. Вирішує розмикання петлі (сталий кінцевий гормон на вході першого ступеня вбиває такий ритм) або відповідь на коротку добавку гормону в різні фази циклу.}

\answer{ex:17:1}{$0.2$~Вт проти $81$~мкВт; сирому потокові потрібно не більше $50$~пДж/біт.}
\answer{ex:17:2}{За $c=0.1$ $5.6$, $8.7$, $9.2$~біт/с, границя $9.3$~біт/с; за $c=0$ і $N=1000$ $65$~біт/с.}
\answer{ex:17:3}{$B=\log_2N+P\log_2P+(1-P)\log_2\frac{1-P}{N-1}=4.87$, $4.37$, $3.29$~біта на знак, тобто $7.3$, $6.6$, $4.9$~біт/с; для $10$~біт/с треба $2.3$~знака на секунду.}
\answer{ex:17:4}{Дисперсія похибки $2b/(Nm^2T)+\sigma_\xi^2$ на компоненту; внески рівні за $N\approx90$; границя $\tfrac12\log_2(1+0.25/0.04)\approx1.4$~біта на компоненту за вікно.}
\answer{ex:17:5}{Пара збігається за $\eta_bd^2+\eta_db^2<2$, тобто приблизно за $\eta<1$: за $0.8$ збігається, за $1.1$ стійкі коливання періоду $2$, за $1.5$ неперіодичні, за $2$ рознос; швидкості учнів треба розводити.}
\answer{ex:17:6}{Поріг $\approx4.4\sigma$; $102$~кбіт/с, $0.10$~мВт проти $205$~мВт для сирого потоку, у $2000$ разів менше; насичується (більше $3$ імпульсів) лише $5.7\cdot10^{-5}$ відліків.}
\answer{ex:17:7}{Відкрита задача. Близько $46$~біт/с за $\text{SNR}\approx0.9$ і близько $330$~біт/с за незалежного шуму. Якщо заважає шум, схожий на сигнал, витягнута інформація насичується зі зростанням кількості каналів; якщо незнання коду, вона зростає з кращим декодером (наприклад, що враховує часи імпульсів).}

\answer{ex:18:1}{(а)~$0.9949\le w\le1.0049$, тобто $|1-w|\lesssim0.005$. (б)~$K\ge400$ синапсів.}
\answer{ex:18:2}{$(\omega_R\tau_w)^2=\sqrt{\alpha(\alpha+2+2\varrho)}/\varrho-1$, $\varrho=\tau_m/\tau_w$; $f_R\approx11.1$~Гц; $|Z(\omega_R)|/|Z(0)|=4.6$.}
\answer{ex:18:3}{(а)~$f=1/\bigl(\tau\ln\frac{\mu}{\mu-\theta}\bigr)$; для $1$~Гц потрібно $\mu/\theta-1\approx3.7\cdot10^{-44}$. (б)~$T=\pi\tau/\sqrt\mu$, $f\approx3.2$~Гц: $f\propto\sqrt\mu$.}
\answer{ex:18:4}{Інтегратор працює за $f\lesssim17.7$~Гц (фільтр нижніх частот), резонатор --- лише в смузі $5.5$--$22$~Гц (смуговий фільтр).}
\answer{ex:18:5}{(а)~$T_{\min}=\frac{\tau}{w-1}\ln\frac{0.5}{0.3}\approx10.2\ms$. (б)~$B>w-1-0.2=0.3$.}
\answer{ex:18:6}{Стала часу підлаштування $\tau/(\eta\langle r^2\rangle)$, $\eta=\tau\ln10/60\,\text{с}\approx3.8\cdot10^{-3}$; за $I\neq0$ вага зміщується, з $e$ треба відняти копію команди $I/\tau$.}
\answer{ex:18:7}{Відкрита задача. Резонатор виграє лише в $1.35$ раза за $11$~Гц і програє в $17$ разів за $1$~Гц; основний виграш від переналаштування --- у пороговому виборі смуги (задача~\ref{ex:18:4}).}

\answer{ex:19:1}{(а)~Один синапс дає $6.7\mV$ ($R=66.7$~МОм), тож до порога взагалі треба не менше $4$ (трьох вистачає лише асимптотично); $8$ за $10\ms$; $14$ за $5\ms$. (б)~$0.1\ms\ll\tau_m$, поправка від витоку $\sim0.25\%$.}
\answer{ex:19:2}{(а)~$1.75\cdot10^{10}$ (чверть змішувачів відповідає на все підряд), $4.4\cdot10^9$ і $0.06$ (за $k=40$ не відповідає майже ніколи). (б)~У $10^3$ разів: $2\cdot10^5$ проти $200$.}
\answer{ex:19:3}{(а)~$C(2n,n)=2^{2n-1}$ із симетрії біноміальних коефіцієнтів. (б)~$0.93$ і $0.09$; ширина переходу $\sim\sqrt{2n}$ за $P$, відносна $\sim1/\sqrt n$.}
\answer{ex:19:4}{З одного дендрита $V_{\text{тіло}}<\kappa E_E<\theta$ за будь-якої кількості входів; за $g_0=0.5g_L$ потрібно $n\ge3$ на кожному дендриті.}
\answer{ex:19:5}{$I\ge C\sigma_V/\sigma_t=15$~нА, тобто $150$ синхронних синапсів --- нереалістично; маленькому нейронові вистачило б $1$~нА, а із синапсом у $4$~нА тремтіння $5$~мкс.}
\answer{ex:19:6}{Максимум на ближньому кінці: $0.03\,\tau_m$ і $0.97$ ($L=0.2$); $0.32\,\tau_m$ і $0.66$ ($L=1$); $0.79\,\tau_m$ і $0.33$ ($L=2$). Оцінка~\eqref{eq:dendriteload} за повною ємністю правильна лише за $L\ll1$.}
\answer{ex:19:7}{Відкрита задача. За фіксованого $m$ час відгуку зростає лінійно із запам'ятовуваним $P$; за фіксованої частки $m=\varphi n$ він перестає залежати від $n$, і повільність великого нейрона --- наслідок підсумовування багатьох входів, а не розміру як такого.}

\answer{ex:20:1}{$t_{\text{нс}}=\tau_r\ln\frac{1-L}{1-H}=16.8$~год, $t_{\text{с}}=\tau_d\ln\frac HL=5.8$~год, період $22.5$~год; до $24$~год генератор підтягує добове гойдання порогів --- фазове автопідстроювання~\eqref{eq:pll}.}
\answer{ex:20:2}{(а)~$L=0.111$; за $L=0.092$ сон тривав би $9.6$~год. (б)~На $22\%$ ($1/0.82=1.22$), а не на $18\%$.}
\answer{ex:20:3}{$H=\frac{p-1}{p-1/q}=0.623$, $L=H/q=0.093$, де $p=e^{16/\tau_r}$, $q=e^{8/\tau_d}$. Після пробудження через $4$~год фаза зсувається назавжди (засинання на $7.2$~год раніше); із гойдальними порогами вона повертається за дві--три доби.}
\answer{ex:20:4}{(а)~$r_\pm=\frac12\bigl(1\pm\sqrt{1-4k/w}\bigr)=0.947,\ 0.053$; $a_{\text{в}}=3.51$, $a_{\text{н}}=0.49$. (б)~Робота $\approx0.33$~с, тиша $\approx0.59$~с, період $\approx0.93$~с, частка роботи $0.36$.}
\answer{ex:20:5}{$a_{\text{в}}/r_+<b<a_{\text{н}}/r_-$: $3.71<b<9.20$. \nt{ACET} зменшує $b$ нижче за $3.71$: перетин переходить на верхню гілку, і стійким стає стан роботи ($r\approx1$).}
\answer{ex:20:6}{$4.9$, $6.8$, $8.8$, $5.5$, $8.9$~год за $D=24$, $32$, $40$, $48$, $64$: тривалість задає добова фаза засинання, а не борг.}
\answer{ex:20:7}{Після B похибка на A в середньому $0.51$ (від $0.19$ до $1.08$ за $20$ наборами; у нульової відповіді --- $1$). Упереміш обидві похибки менші за $10^{-4}$.}
\answer{ex:20:8}{Відкрита задача. Рівномірне масштабування зберігає відношення ваг, але зберігає відповіді порогового нейрона лише тоді, коли поріг масштабується в ту саму кількість разів; перемішані повтори змінюють відношення. Незмінність великих контактів суперечить суто рівномірному масштабуванню.}

\answer{ex:21:1}{(а)~$T=\tau_{\text{відн}}\ln\frac{1}{1-x^*}$: $1.84$~с і $3.68$~с. (б)~Від $3.36$ до $4.24$~с: чутливість $\dd T/\dd x^*=\tau_{\text{відн}}/(1-x^*)=80$~с, при $x^*\to1$ залпи нерегулярні.}
\answer{ex:21:2}{(а)~$8.6$~Вт, як у~\eqref{eq:energy}. (б)~$205$~Мбіт/с, $0.2$~Вт --- у $2\cdot10^3$ разів більше за культуру ($10^{-4}$~Вт).}
\answer{ex:21:3}{$x_{\text{уст}}=1/(1+Ur_b\tau_{\text{відн}})<x^*$ при $r_b>(1/x^*-1)/(U\tau_{\text{відн}})=0.28$~Гц для $x^*=0.9$ і $0.025$~Гц для $x^*=0.99$; сила синапсу падає до $x_{\text{уст}}$, тобто на $10\%$ і $1\%$.}
\answer{ex:21:4}{$u(t-k)$ ортонормовані, $m_k=\|P_Vu(t-k)\|^2$, де $P_V$ --- проєктор на $N$-вимірну оболонку виходів; за нерівністю Бесселя $\sum_k m_k\le N$.}
\answer{ex:21:5}{$\mathrm{MC}\approx9$--$11$ з $\tanh$ і $15$--$18$ для лінійного резервуара (три затравки), обидва $\le30$: нелінійність платить пам'яттю.}
\answer{ex:21:6}{(а)~$n\le102$. (б)~На $5.6$.}
\answer{ex:21:7}{Відкрита задача. Ключовий контроль: заморозити зчитування і перевірити, чи зберігається покращення після паузи без стимуляції; якщо ні, змінився стан, а не ваги.}

\answer{ex:22:1}{(а)~$32\%$ і $100\%$. (б)~$R_\uparrow-R_\downarrow\ge2\sqrt{2000\,\text{Гц}/1\,\text{с}}\approx89$~Гц, лише $4.5\%$ сумарної частоти.}
\answer{ex:22:2}{(а)~$n\ge4(p_1q_1+p_2q_2)/\Delta^2\approx560$. (б)~$\approx56$.}
\answer{ex:22:3}{(а)~Детальний баланс: $\rho PK$ симетричне. (б)~Частка промахів дорівнює середньому гармонічному $1/\langle1/P\rangle=0.18$ проти $0.5$ без зворотного зв'язку.}
\answer{ex:22:4}{(а)~Паралелограм площею $4h^2=0.04$ (з урахуванням обмеження $p$ --- чисельно $0.045$). (б)~$\approx0.18$.}
\answer{ex:22:5}{$0.49$: частина культур іде в область, де промах майже завжди, і блукає там. З обмеженням --- $1.00$: на обмеженій множині густина $\propto1/P$ збирається в поглинальній області.}
\answer{ex:22:6}{(а)~Потрібно $\ge5$ рівнів; $8$ електродів вистачає. (б)~Ні: потрібно $r_{\max}\ge25/T=100$~Гц.}
\answer{ex:22:7}{Відкрита задача. Час випадкового пошуку зростає з $d$ приблизно як відношення об'ємів $(L/h)^d$, у пошуку зі знаком помилки --- лінійно за $d$. Розділяє гіпотези дослід із перестановкою: після промаху --- передбачуваний, але сильний стимул, після удару --- непередбачуваний, але слабкий; потрібно кілька десятків культур у кожній групі.}


\chapter{Експериментальний додаток}
\label{sec:supplement}

Для кожного розділу тут зібрано його головні твердження і те, на чому вони тримаються: дослід, у якому це виміряно, що саме виміряно і де це опубліковано. Статус у правому стовпці показує, наскільки твердження тверде: \emph{виміряно} --- є прямий дослід; \emph{теорія} --- математичний наслідок, виведений у розділі або в класичній роботі; \emph{модель} --- результат розрахунку за моделлю книги (скрипти в теці \texttt{models/}); \emph{оцінка} --- число порядку величини без прямого вимірювання; \emph{гіпотеза} --- правдоподібне пояснення, яке дослідом ще не доведено. Де досліду немає, так і сказано.

\section{Розділ~\ref{sec:neurontypes}. Типи нейронів}

Числа розділу спираються на прямі підрахунки клітин у мозку людини там, де такі підрахунки є, правило «один нейрон --- один медіатор» --- на клітинні атласи й досліди, які знайшли й винятки, роль \nt{DOPA} --- на записи імпульсів, а ролі решти широкомовних каналів --- на гіпотези.

{\footnotesize
\setlength{\tabcolsep}{3pt}\renewcommand{\arraystretch}{1.15}
\begin{longtable}{>{\raggedright}p{0.5cm}>{\raggedright}p{4.0cm}>{\raggedright}p{5.6cm}>{\raggedright}p{2.3cm}>{\raggedright\arraybackslash}p{1.7cm}}
\toprule
№ & Твердження & Дослід і що виміряно & Джерело & Статус \\
\midrule\endfirsthead
\toprule
№ & Твердження & Дослід і що виміряно & Джерело & Статус \\
\midrule\endhead
1 & У мозку людини близько $8.6\cdot10^{10}$ нейронів, і майже всі \nt{GLUT}-нейрони --- дрібні нейрони мозочка (табл.~\ref{tab:types}) & Ізотропний фракціонатор: тканину розчиняють в однорідну суспензію ядер і рахують ядра з нейронним маркером NeuN. У дорослих чоловіків $86.1\pm8.1$ млрд нейронів; у корі великих півкуль із них лише $19\,\%$, основна маса --- в мозочку & \cite{Azevedo2009} & виміряно \\
2 & Майже кожен нейрон має один головний медіатор (твердження~\ref{prop:dale}), але є винятки & Транскриптомний атлас мозку миші: близько $4\cdot10^6$ клітин, $5322$ кластери; кластерам приписано медіатор за генами синтезу й пакування. Винятки виміряно безпосередньо: \nt{DOPA}-нейрони у зрізах викидають ще й GABA через той самий везикулярний транспортер і гальмують нейрони стріатума; у частини \nt{DOPA}-аксонів глутамат і дофамін виходять із різних закінчень одного проводу (електронна мікроскопія та оптогенетика) & \cite{Strata1999,Yao2023,Tritsch2012,Zhang2015} & виміряно \\
3 & Увесь стовпець матриці ваг одного знака~\eqref{eq:dalesign} & Висновок у розділі з твердження~\ref{prop:dale} і того, що рецептори глутамату майже всі збуджувальні, а GABA --- гальмівні. Прямої перевірки «усі отримувачі одного нейрона отримують вагу одного знака» як загального правила немає; контрприклади --- спільний викид GABA \nt{DOPA}-нейронами (рядок~2) і отримувачі з рецепторами різного знака (рядок~11) & висновок у розділі & теорія \\
4 & Від трьох чвертей до чотирьох п'ятих нейронів кори --- \nt{GLUT}, від п'ятої частини до чверті --- \nt{GABA} & Імунофарбування на GABA у стовпчиках завширшки $50$\,мкм крізь усю товщу кори в $10$ областях кори мавпи: у $8$ областях GABA-нейрони становлять близько $25\,\%$ усіх нейронів, у первинній зоровій корі --- менше $20\,\%$ & \cite{Hendry1987} & виміряно \\
5 & \nt{GLUT}-нейрон має близько $10^4$ виходів & Стереологія на розтині: у новій корі молодих чоловіків $1.64\cdot10^{14}$ синапсів (електронна мікроскопія, дисектор) і близько $2.3\cdot10^{10}$ нейронів. Відношення дає $\sim7\cdot10^3$ синапсів на нейрон; у середньому кількість входів дорівнює кількості виходів & \cite{Tang2001synapses,Pakkenberg1997} & виміряно \\
6 & Вихід \nt{DOPA}-нейрона розгалужується на $10^5$--$10^6$ закінчень; це оцінка за охопленням мішені, а не підрахунок & Вірусна мітка мембрани окремих \nt{DOPA}-нейронів щура й повна реконструкція аксона: один нейрон покриває від $0.45$ до $5.7\,\%$ (у середньому $2.7\,\%$) об'єму стріатума. Виміряно охоплення, а не кількість закінчень: її виведено з охоплення за порядком величини, прямого підрахунку закінчень немає. Для \nt{SERO} і \nt{NORA} кількості закінчень --- грубі оцінки & \cite{Matsuda2009} & виміряно \\
7 & Широкомовних нейронів мало: \nt{DOPA} $\sim5\cdot10^5$ (головна з двох груп, нижня межа), \nt{SERO} $\sim3\cdot10^5$ (сума двох часткових підрахунків), \nt{HIST} $\sim6\cdot10^4$ (табл.~\ref{tab:types}, рис.~\ref{fig:typepie}) & Стереологічні підрахунки в людини: $550\,000$ пігментованих нейронів у чорній субстанції (без вентральної покришки); $235\,000$ нейронів у дорсальному ядрі шва (усі нейрони ядра) і ще $125\,000$ серотонінових нейронів у покришці мосту; близько $64\,000$ гістамінових нейронів гіпоталамуса. Для \nt{NORA}, \nt{GLYC}, \nt{ACET} і \nt{ADRE} прямого підрахунку немає: у таблиці це грубі оцінки за порядком величини & \cite{Pakkenberg1991,Baker1990,Baker1991,Airaksinen1991} & виміряно \\
8 & Глутамат у щілині злітає приблизно до мілімоля й спадає за $\sim1\ms$ & За витісненням швидко від'єднуваного конкурентного блокатора з рецепторів NMDA під час синаптичної передачі (культура нейронів гіпокампа): пік $1.1$\,мМ, стала спаду $1.2\ms$ & \cite{Clements1992} & виміряно \\
9 & У працюючій корі збуджувальний і гальмівний струми майже врівноважені, нейрон сидить біля порога & Теорія: мережа із сильними зв'язками сама приходить у врівноважений режим. Дослід: у префронтальній корі тхора in vivo під час «верхніх» станів потенціал реверсії синаптичного струму тримається біля $-37$\,мВ сотні мілісекунд, хоча провідність змінюється на $\sim21$\,нСм: збудження й гальмування зростають і спадають разом & \cite{vanVreeswijk1996,Haider2006} & виміряно \\
10 & \nt{DOPA}-нейрони повідомляють помилку передбачення винагороди~\eqref{eq:rpe} (рис.~\ref{fig:rpe}) & Імпульси \nt{DOPA}-нейронів мавпи: пачка на несподіваний сік; після навчання --- пачка на сигнал і немає відповіді на передбачений сік; провал частоти, коли обіцяного соку не дано. У кількісному досліді частота пропорційна різниці поточної винагороди та зваженого середнього минулих, але лише для додатних помилок. Саме рівняння~\eqref{eq:rpe} --- алгоритм часових різниць & \cite{Schultz1997,Bayer2005,SuttonBarto2018} & виміряно \\
11 & Знак і швидкість задає рецептор: іонотропні --- частки мілісекунди, метаботропні --- набагато повільніше (твердження~\ref{prop:receptor}) & Швидкі імпульси глутамату на клаптики мембрани нейронів гіпокампа: струм через іонотропні рецептори наростає (від $20$ до $80\,\%$) за $0.2$--$0.6\ms$ і спадає за $2$--$3\ms$. Повільний гальмівний потенціал пірамідних нейронів знімається вибірковим блокатором метаботропного рецептора GABA. Дві родини рецепторів дофаміну з протилежною дією; у стріатумі нейронам із рецептором D1 дофамін потрібен для посилення синапсів, нейронам із D2 --- для послаблення & \cite{Colquhoun1992,DutarNicoll1988,Beaulieu2011,Shen2008} & виміряно \\
12 & Широкомовний канал --- не один провід, а джгут із небагатьох ліній (підрозділ~\ref{sec:buses}) & Вірусно-генетичне маркування серотонінових нейронів дорсального ядра шва миші: нейрони, що посилають вихід у мигдалину й у лобову кору, лежать у різних місцях ядра, розгалужуються в області, які доповнюють одна одну, і протилежно відповідають на неприємний подразник. Огляд: підгрупи нейронів блакитної плями (\nt{NORA}) кодують різні процеси & \cite{Ren2018,Poe2020} & виміряно \\
13 & \nt{NORA} задає коефіцієнт підсилення $g$ & Гіпотеза адаптивного підсилення; модель мережі, у якій катехоламіни підвищують крутість передавальної характеристики. Прямого вимірювання «$g$ зростає з норадреналіном» у всій мережі немає; виміряно, що діаметр зіниці слідує за імпульсами нейронів блакитної плями мавпи & \cite{AstonJones2005,ServanSchreiber1990,Joshi2016} & гіпотеза \\
14 & \nt{ACET} перемикає «запис/читання» і задає швидкість навчання & Гіпотеза. Опора: у зрізах нюхової кори щура агоністи ацетилхоліну ($100$\,мкМ) послаблюють внутрішні, «пригадувальні» синапси більш ніж на $60\,\%$, а вхідні --- менш ніж на $15\,\%$ & \cite{Hasselmo2006,HasselmoBower1992} & гіпотеза \\
15 & \nt{SERO} задає коефіцієнт дисконтування $\gamma$; серотонінові нейрони працюють рівно, кілька імпульсів за секунду, і майже мовчать в одній із фаз сну & Гіпотеза «серотонін $=\gamma$». Найсильніший аргумент: оптогенетична активація серотонінових нейронів дорсального ядра шва миші зменшує кількість передчасних відмов від очікування відкладеної винагороди. Частоту імпульсів і мовчання уві сні зі швидкими рухами очей виміряно (огляд записів) & \cite{Doya2002,Miyazaki2014,JacobsAzmitia1992} & гіпотеза \\
\bottomrule
\end{longtable}}

\section{Розділ~\ref{sec:glutcomputer}. Що обчислюють \nt{GLUT}-нейрони}

Логічні твердження розділу --- теореми про схеми з невід'ємних ваг, виведені в самому розділі; дослід підтверджує модель нейрона й те, що схеми, зібрані за цими теоремами (детектор збігів, лінія затримки, самопідтримувана група, ланцюжок), у мозку справді трапляються.

{\footnotesize
\setlength{\tabcolsep}{3pt}\renewcommand{\arraystretch}{1.15}
\begin{longtable}{>{\raggedright}p{0.5cm}>{\raggedright}p{4.0cm}>{\raggedright}p{5.6cm}>{\raggedright}p{2.3cm}>{\raggedright\arraybackslash}p{1.7cm}}
\toprule
№ & Твердження & Дослід і що виміряно & Джерело & Статус \\
\midrule\endfirsthead
\toprule
№ & Твердження & Дослід і що виміряно & Джерело & Статус \\
\midrule\endhead
1 & Нейрон --- RC-коло з порогом і скиданням~\eqref{eq:glutneuron} & У тіло пірамідних нейронів кори щура (зрізи) подавали шумовий струм, схожий на синаптичний потік у живому мозку. Залежність середньої частоти від середнього й розкиду струму описується інтегрувальним нейроном із витоком, якщо додати повільну адаптацію частоти. Діапазони $\tau$ і затримок наведено в розділі без посилання на дослід & \cite{Rauch2003} & виміряно \\
2 & Модель~\eqref{eq:glutneuron} --- спрощення: гілки входів самі видають місцеві імпульси & Запис безпосередньо з дендритів пірамідних нейронів зорової кори миші in vivo: зоровий подразник викликає місцеві дендритні імпульси, що залежать від рецепторів NMDA; їх придушення розширює налаштування нейрона на орієнтацію & \cite{Smith2013} & виміряно \\
3 & Вікно збігу $\Delta_{\max}=\tau\ln\frac{w}{\theta-w}$~\eqref{eq:window}; при $\theta=1.5w$ воно дорівнює $0.7\tau$ & Наслідок моделі~\eqref{eq:glutneuron}. Пряме вимірювання вузького вікна підсумовування є для нейронів зі швидким гальмуванням (розділ~\ref{sec:gaba}, рядок~3 його таблиці) & висновок у розділі & теорія \\
4 & Мережа з самих лише \nt{GLUT}-нейронів обчислює лише монотонні функції входів (твердження~\ref{prop:monotone}) & Доведення в розділі: ітерація від тиші з неспадною правою частиною. Це твердження про модель; досліду, який перевіряв би його на живій мережі без гальмування, немає & висновок у розділі & теорія \\
5 & Органи чуття дають пару проводів: одні клітини відповідають на посилення подразника, інші --- на послаблення & Сітківка: уже на біполярних клітинах сигнал колбочок розходиться на вмикальний і вимикальний канали. Фармакологічний блок вмикальних біполярних клітин у мавпи: канали йдуть окремо до кори; після блоку погіршується виявлення посилення світла, але не послаблення & \cite{Schiller1992} & виміряно \\
6 & З двопровідним кодом мережа з \nt{GLUT}-нейронів обчислює будь-яку логічну функцію асинхронно, без спільного такту, ціною вдвічі більшої кількості нейронів (твердження~\ref{prop:dualrail}, \eqref{eq:dualrail}, рис.~\ref{fig:dualxor}) & Двопровідний код і визначення готовності за самим сигналом --- стандартна техніка асинхронних схем, нечутливих до затримок. Схему виключного АБО з шести нейронів зібрано в розділі. Досліду, який показував би, що мозок обчислює логічні функції саме двопровідним кодом, немає & \cite{Sparso2001} & теорія \\
7 & Найбільша різниця часу приходу звуку у вуха в людини $\approx0.6\ms$, а мозок розрізняє близько десяти мікросекунд & Слухачі в досліді з двома варіантами відповіді: поріг розрізнення різниці часу (75\,\% правильних) --- $9$\,мкс для шуму в смузі $150$--$1700$\,Гц, $11$\,мкс для тону $1$\,кГц, $28$\,мкс для клацання. Межа $0.6\ms$ --- розрахунок розділу, $0.22\,\text{м}/343\,\text{м/с}$ & \cite{KlumppEady1956} & виміряно \\
8 & У птахів різниця часів перетворюється на місце через лінії затримки й детектори збігів~\eqref{eq:itd} (рис.~\ref{fig:itd}) & Сипуха: аксони від двох вух входять у ядро детекторів збігів із протилежних боків; час приходу імпульсів уздовж аксона змінюється впорядковано з глибиною, і розкид затримок крізь ядро ($160$\,мкс на $700$\,мкм) збігається з діапазоном міжвушних затримок & \cite{Carr1990} & виміряно \\
9 & У ссавців ряду затримок не знайшли; ту саму задачу розв'язують детектори збігів із точно розрахованим гальмуванням від \nt{GLYC}-нейронів & Запис in vivo в медіальній верхній оливі піщанки: відповіді не утворюють карти напрямків; підведення гліцину та його блокатора через мікропіпетку показує, що точно розраховане в часі \emph{гліцинове} гальмування задає робочий діапазон затримок. Гальмування тут від \nt{GLYC}, а не від \nt{GABA} & \cite{Brand2002,Grothe2010} & виміряно \\
10 & Група із сильним зворотним зв'язком --- тригер; самопідтримувана активність триває секунди після подразника (рис.~\ref{fig:bistable}) & Префронтальна кора мавпи в задачі з відкладеним рухом очей до запам'ятованої точки (затримка $1$--$6$\,с): $87$ із $170$ нейронів, пов'язаних із задачею, змінюють частоту під час затримки, у $79\,\%$ із них це залежить від напрямку на запам'ятовану точку. Що цю активність тримає зворотний зв'язок із двома стійкими станами, --- теорія & \cite{Funahashi1989,Wang2001} & виміряно \\
11 & Біт записується, якщо приплив триває довше $\approx0.7\tau$ (рис.~\ref{fig:recurrentgroup}б) & Розрахунок рівняння $\tau\,\dd r/\dd t=-r+f(wr+I)$ методом Ейлера при $w=1$, $I=0.6$; скрипту в \texttt{models/} немає. Досліду немає & розрахунок у розділі & модель \\
12 & Пам'ять можна зберігати в полегшенні синапсів, майже без імпульсів & Теорія: залишковий кальцій у закінченнях тримає запис близько секунди. Опора: у медіальній префронтальній корі тхора (запис із кількох нейронів одночасно) є підмережа пірамідних нейронів, зв'язаних полегшуваними синапсами з тривалою післядією. Огляд спостережень «тихих» проміжків під час утримання в пам'яті. Який спосіб кора використовує коли --- відкрите питання & \cite{Mongillo2008,WangMarkram2006,Stokes2015} & гіпотеза \\
13 & Пам'ять стирається втомою: запас медіатора виснажується & Парні записи пірамідних нейронів кори: відповідь синапсу падає за частих імпульсів, швидкість падіння задається ймовірністю викиду. Що саме це гасить самопідтримувану групу, безпосередньо не показано & \cite{Tsodyks1997} & виміряно \\
14 & Ланцюжок груп --- таймер, що запускається подією; у співочих птахів нейрони спрацьовують суворо по черзі & Запис упізнаних нейронів вищого вокального центру зебрової амадини уві сні й під час співу: кожен нейрон, що посилає вихід у моторне ядро, видає одну коротку пачку в один точний момент пісні, і нейрони спрацьовують послідовно & \cite{Hahnloser2002} & виміряно \\
\bottomrule
\end{longtable}}

\section{Розділ~\ref{sec:gaba}. \nt{GLUT} і \nt{GABA} разом}

Логічні й динамічні твердження розділу виведено в ньому самому з лінійного аналізу й закону Ома; дослід показує, що схеми, на які вони спираються, --- заборона із затримкою, пряме гальмування слідом за збудженням, стабілізація гальмуванням, ритм петлі \nt{GLUT}--\nt{GABA}, --- у мозку справді працюють.

{\footnotesize
\setlength{\tabcolsep}{3pt}\renewcommand{\arraystretch}{1.15}
\begin{longtable}{>{\raggedright}p{0.5cm}>{\raggedright}p{4.0cm}>{\raggedright}p{5.6cm}>{\raggedright}p{2.3cm}>{\raggedright\arraybackslash}p{1.7cm}}
\toprule
№ & Твердження & Дослід і що виміряно & Джерело & Статус \\
\midrule\endfirsthead
\toprule
№ & Твердження & Дослід і що виміряно & Джерело & Статус \\
\midrule\endhead
1 & \nt{GABA}-нейронів у корі від п'ятої частини до чверті & Імунофарбування на GABA у $10$ областях кори мавпи: близько $25\,\%$ нейронів у $8$ областях, менше $20\,\%$ у первинній зоровій корі. Огляд різноманіття гальмівних нейронів кори & \cite{Hendry1987,Markram2004} & виміряно \\
2 & Що робить гальмування, багато в чому вирішує місце посадки: дендрит, тіло, місце народження імпульсу, закінчення чужого проводу & Огляд: класи \nt{GABA}-нейронів кори різняться мішенню на отримувачі. Одночасний запис із тіла й дендрита пірамідного нейрона: пряме гальмування набагато сильніше на тілі, ніж на дендритах. Гальмування на закінченні чужого проводу в спинному мозку зменшує викид медіатора однієї вхідної лінії (огляд вимірювань) & \cite{Markram2004,Pouille2001,Rudomin1999} & виміряно \\
3 & Зміна знака через посередника коштує однієї затримки, близько мілісекунди; гальмування, що прийшло слідом за збудженням, звужує вікно збігу & Пірамідні нейрони гіпокампа щура у зрізах: гальмування через одного посередника приходить слідом за прямим збудженням із короткою затримкою, і вікно, у якому два збуджувальні входи підсумовуються до порога, менше за $2\ms$. На дендритах, де це гальмування слабше, вікно ширше & \cite{Pouille2001} & виміряно \\
4 & Заборона $a\wedge\bar b$~\eqref{eq:veto} з АБО й сталою одиницею функціонально повна (твердження~\ref{prop:gabacomplete}) & Доведення в розділі. Класичний результат: мережа порогових елементів із гальмівними входами реалізує будь-який логічний вираз. Твердження про модель, досліду немає & \cite{McCullochPitts1943} & теорія \\
5 & Заборона із затримкою дає детектор напрямку руху з оптимальною швидкістю $v^*=\Delta x/d$~\eqref{eq:vstar} & Сітківка кроля: вибірковість до напрямку створюється гальмуванням, яке за руху в «нульовому» напрямку гасить збудження. Окремий запис збуджувального й гальмівного входів вибіркових клітин показав, що зірчасті амакринові клітини (\nt{GABA}) гальмують їх несиметрично --- лише відростками, спрямованими в «нульовий» бік. Формула $v^*$ --- висновок розділу & \cite{Barlow1965,Fried2002} & виміряно \\
6 & Шунтувальне гальмування ділить напругу, а не віднімає~\eqref{eq:shunt} (рис.~\ref{fig:shunt}) & Закон Ома для подільника з трьох провідностей за рівноважного потенціалу хлору біля спокою; для нейрона без імпульсів & висновок у розділі & теорія \\
7 & На частоту імпульсів шунт діє відніманням, а ділення виходить від гальмування разом зі врівноваженим фоновим шумом & Модель: у нейрона, що спрацьовує, струм через шунт майже сталий, і ефект на частоту віднімальний. Дослід: у пірамідні нейрони соматосенсорної кори щура (зрізи) через динамічну фіксацію вводили фонові збуджувальні й гальмівні провідності; одночасне врівноважене зростання обох ділить крутість відповіді без помітного зсуву частоти. Огляд: ділення на загальну активність трапляється в багатьох відділах мозку & \cite{HoltKoch1997,Chance2002,Carandini2012} & виміряно \\
8 & При $\beta>1$ взаємне гальмування обирає одного переможця (твердження~\ref{prop:wta}, \eqref{eq:wta}) & Власні значення $-(1\pm\beta)/\tau$ --- висновок у розділі. Частоти $0.73$ і $0.53$ при $\beta=0.5$ --- нерухома точка $r_{1,2}=(I_{1,2}-\beta I_{2,1})/(1-\beta^2)$; скрипту в \texttt{models/} немає. Прямого досліду в розділі не наведено & висновок у розділі & теорія \\
9 & Скидання пам'яті сигналом через \nt{GABA}; дві групи із взаємним гальмуванням --- засувка & Випливає з рис.~\ref{fig:bistable} і твердження~\ref{prop:wta}. Досліду немає & висновок у розділі & теорія \\
10 & Рівновага петлі \nt{GLUT}--\nt{GABA} стійка, якщо гальмування досить сильне й досить швидке (твердження~\ref{prop:ei}) & Модель двох популяцій~\eqref{eq:ei}; умови (i) і (ii) --- визначник і слід її матриці, виведені в розділі & \cite{WilsonCowan1972} & теорія \\
11 & При $w_{EE}=1.5$ і $w_{EI}w_{IE}=2$ швидке гальмування ($\tau_I=2\ms$) тримає $E^*=0.67h$, повільне ($30\ms$) розгойдує коливання (рис.~\ref{fig:ei}) & Чисельний розв'язок~\eqref{eq:ei}, скрипт \texttt{figures/make\_ei.py} & розрахунок & модель \\
12 & Кора працює в режимі стабілізації гальмуванням; підпис --- підштовхнеш \nt{GABA}-нейрони, і їхня частота впаде~\eqref{eq:isn} & Теорія передбачила парадоксальну відповідь. Дослід: внутрішньоклітинні записи в первинній зоровій корі кішки --- при пригніченні відповіді подразником навколо гальмівна провідність не зростає, а падає разом зі збуджувальною. Оптогенетичне збудження PV-нейронів миші в зоровій, соматосенсорній і моторній корі знижує частоту гальмівних нейронів, зокрема без подразника й за легкого наркозу. Коефіцієнт $-1/3$ за чисел рис.~\ref{fig:ei} --- висновок розділу & \cite{Tsodyks1997isn,Ozeki2009,Sanzeni2020} & виміряно \\
13 & Петля «\nt{GLUT} збуджує \nt{GABA}, \nt{GABA} гальмує \nt{GLUT}» породжує швидкий ритм із періодом $12$--$50\ms$ & Оптогенетична стимуляція швидких \nt{GABA}-нейронів у соматосенсорній корі миші на частотах $8$--$200$\,Гц вибірково посилює гамма-ритм ($20$--$80$\,Гц, період $12$--$50\ms$); стимуляція \nt{GLUT}-нейронів посилює лише повільніші ритми & \cite{Cardin2009,Buzsaki2012} & виміряно \\
\bottomrule
\end{longtable}}

\section{Розділ~\ref{sec:code}. Абетка Морзе}

Граничні оцінки розділу --- теорія інформації та рахунок; виміряно, де в каналі виникає шум, як швидко працює мозок і скільки коштує імпульс, а питання про те, яким кодом кора користується насправді, лишається відкритим.

{\footnotesize
\setlength{\tabcolsep}{3pt}\renewcommand{\arraystretch}{1.15}
\begin{longtable}{>{\raggedright}p{0.5cm}>{\raggedright}p{4.0cm}>{\raggedright}p{5.6cm}>{\raggedright}p{2.3cm}>{\raggedright\arraybackslash}p{1.7cm}}
\toprule
№ & Твердження & Дослід і що виміряно & Джерело & Статус \\
\midrule\endfirsthead
\toprule
№ & Твердження & Дослід і що виміряно & Джерело & Статус \\
\midrule\endhead
1 & Частоти \nt{GLUT}-нейронів кори --- від часток до десятків імпульсів за секунду, і більшу частину часу нейрони працюють біля нижньої межі & Запис через піпетку, прикладену до клітини (вибір нейрона не залежить від його активності), у слуховій корі щура, що не спить: у кожен момент частоту вище $20$ імпульсів за секунду дають менше $5\,\%$ нейронів; у частини нейронів фонова частота нижча за $0.01$ за секунду; розподіл частот логнормальний & \cite{Hromadka2008} & виміряно \\
2 & Пропускна здатність імпульсного каналу $R_{\max}\approx r\log_2\frac{e}{r\Delta t}$~\eqref{eq:spikeentropy}, майже вся вона --- у часі імпульсів (твердження~\ref{prop:capacity}) & Ентропія двійкового рядка з клітинок довжини $\Delta t$. Числа розділу: $8$ біт на імпульс при $r=10$ за секунду й $\Delta t=1\ms$, близько $5$ біт при $\Delta t=10\ms$, менше $2$ біт за секунду в лічильника імпульсів & \cite{MacKay1952} & теорія \\
3 & Чутливі нейрони передають від одного до кількох біт на імпульс, у найкращих випадках до половини межі & Відновлення подразника за імпульсами чутливих нейронів, для яких подразник відомий; зведення вимірювань у книзі & \cite{Rieke1997} & виміряно \\
4 & Генератор імпульсів точний; точний час задають швидкі перепади входу & Нейрони кори щура у зрізах: на повторюваний струм зі швидкими коливаннями імпульси повторюються з точністю менше $1\ms$, на сталий струм моменти розповзаються & \cite{Mainen1995} & виміряно \\
5 & Синапс ненадійний: понад половина імпульсів не доходить до отримувача через випадковість викиду & Мінімальна стимуляція входів пірамідних нейронів гіпокампа (зрізи): понад половина імпульсів не викликає відповіді. Виключено коливання порога стимуляції, збої проведення на розгалуженнях аксона й низьку температуру досліду; лишається ймовірнісний викид & \cite{AllenStevens1994} & виміряно \\
6 & Потік імпульсів у живій корі виглядає випадковим: інтервали розкидані як у пуассонівського потоку, дисперсія кількості імпульсів близька до середнього; частина «шуму» --- повідомлення про рухи тварини & Записи в зорових областях V1 і MT мавпи, що не спить: коефіцієнт варіації інтервалів близько $1$, відношення дисперсії кількості імпульсів до середнього --- як у випадкового процесу; модель, що підсумовує незалежні малі входи, дає розкид набагато менший. У зоровій корі миші значну частину розкиду передбачають за відеозаписом рухів самої тварини & \cite{Softky1993,Stringer2019} & виміряно \\
7 & Частотний код повільний: похибка $1/\sqrt{rT}$~\eqref{eq:rateerror}, а мозок упізнає картинку за $\sim150\ms$ & Формула --- пуассонівська статистика, висновок розділу. Дослід: люди вирішували, чи є тварина на незнайомій фотографії, показаній на $20\ms$; потенціали мозку на «так» і «ні» розходяться приблизно через $150\ms$. Розбиття цього часу на десяток ступенів по $10$--$20\ms$ --- оцінка розділу, не вимірювання & \cite{Thorpe1996} & виміряно \\
8 & Усереднення за групою обмежене кореляцією: похибка падає не більше ніж у $1/\sqrt c$ разів~\eqref{eq:correlated} (твердження~\ref{prop:pooling}) & Пари нейронів області MT мавпи, записані одночасно: коефіцієнт кореляції кількостей імпульсів у середньому $0.12$, і теоретичний аналіз показує, що навіть така кореляція обмежує виграш від групи. Теорія: межу ставить лише та частина кореляції, що схожа на сигнал. Модель протилежної позиції: частоту групи з $50$--$100$ нейронів читають за один міжімпульсний інтервал, $10$--$50\ms$, і час окремих імпульсів у корі не використовується & \cite{Zohary1994,MorenoBote2014,ShadlenNewsome1998} & виміряно \\
9 & Число можна писати в час першого імпульсу~\eqref{eq:latency}, у порядок імпульсів групи або у фазу місцевого ритму & Сітківка саламандри: просторову структуру коротко показаної картинки записано у відносних затримках перших імпульсів вихідних нейронів, майже незалежно від контрасту. Клітини місця гіпокампа щура: під час проходу через «своє» місце імпульси зсуваються дедалі раніше в циклі тета-ритму ($7$--$12$\,Гц), зсув від $100^\circ$ до $355^\circ$. Наскільки кора користується часом імпульсів --- відкрите питання (рядок~8) & \cite{Gollisch2008,OKeefe1993} & виміряно \\
10 & Який код прочитано, вирішує стала часу отримувача~\eqref{eq:reader} (твердження~\ref{prop:reader}); в активній корі нейрони ближчі до детекторів збігів & Формула~\eqref{eq:reader} --- висновок розділу. Огляд записів in vivo: в активній корі провідність мембрани в кілька разів вища, ніж у спокої, і стала часу відповідно коротша. Що кора перемикає код саме так, безпосередньо не показано & \cite{Destexhe2003} & гіпотеза \\
11 & Синапс, що виснажується, передає зміну частоти, по суті $\Delta r/r$~\eqref{eq:depression} (рис.~\ref{fig:depression}) & Парні записи пірамідних нейронів кори: швидкість виснаження задається ймовірністю викиду; за повільного виснаження відповідь відображає частоту, за швидкого --- збіги. Модель на основі виміряного виснаження в корі щура: однакові відносні зміни частоти на швидких і повільних входах дають однакову відповідь, тобто синапс передає $\Delta r/r$. Числа $1.7\to2.2$, $6.7$ і $90\ms$ --- розрахунок розділу & \cite{Tsodyks1997,Abbott1997depression} & виміряно \\
12 & Пачка --- «тире»: її ймовірність загубитися $(1-p)^k$~\eqref{eq:burst} мала, полегшення зменшує її ще & Огляд: на ненадійних синапсах одиночні імпульси часто губляться, а пачки передаються надійно завдяки полегшенню; пачка викликає довготривалі зміни синапсів. Що мозок читає пачку як окремий знак, --- гіпотеза & \cite{Lisman1997,AllenStevens1994} & гіпотеза \\
13 & Швидкі \nt{GABA}-нейрони задають ритм і «пунктуацію»; ще один клас гальмує гальмівних і відчиняє ворота (твердження~\ref{prop:punctuation}) & Огляд властивостей класів \nt{GABA}-нейронів кори. Оптогенетика: ритмічне збудження швидких \nt{GABA}-нейронів викликає гамма-ритм, і відповідь на подразник залежить від фази циклу. У зоровій корі миші PV-нейрони гальмують один одного, SST-нейрони --- усі інші класи, VIP-нейрони --- переважно SST. Збудження VIP-нейронів у миші, що не спить, знімає гальмування з \nt{GLUT}-нейронів; їх вмикає сигнал підкріплення. Поділ ролей як загальне правило --- гіпотеза & \cite{Tremblay2016,Cardin2009,Pfeffer2013,Pi2013} & гіпотеза \\
14 & Енергія обмежує середню частоту величиною порядку одного імпульсу за секунду~\eqref{eq:energy} & Енергетичний бюджет сірої речовини гризуна за анатомічними й фізіологічними даними: імпульси й синаптичні струми --- $47\,\%$ і $34\,\%$ енергії сигналізації; одночасно можуть бути активними не більше $\sim15\,\%$ нейронів. Для кори людини: не більше $\sim1\,\%$ нейронів можуть бути сильно активними одночасно. Число $\sim10^9$ АТФ на імпульс і потужність $\sim10$\,Вт на сигналізацію --- оцінка розділу на основі цих розрахунків & \cite{Attwell2001,Lennie2003} & теорія \\
15 & Код розріджений і підлаштований під найбільшу кількість біт на джоуль; кора міняє точність на енергію (твердження~\ref{prop:economy}) & Гіпотеза ощадливого коду. Опора: у мишей за обмеження їжі в зоровій корі падає провідність AMPA-рецепторів і синаптична витрата АТФ ($-29\,\%$), частота імпульсів зберігається, а налаштування на орієнтацію розширюється на $32\,\%$ & \cite{Barlow1961,Padamsey2022} & гіпотеза \\
\bottomrule
\end{longtable}}

\section{Розділ~\ref{sec:serocomputer}. \nt{SERO}-нейрони}

Властивості самих \nt{SERO}-нейронів (ритм, автогальмування, знак за рецептором, підсистеми) виміряно прямо; обчислення розділу --- регулятор, фільтр, логіка --- випливають із його рівнянь; $\tau_c$, коефіцієнт дифузії серотоніну і кількість місць викиду --- оцінки; роль петлі як регулятора рівня в мозку --- гіпотеза (у ядрі шва автогальмування точкове і дає лише короткі паузи), а роль \nt{SERO} як горизонту --- гіпотеза, на користь якої свідчать поведінкові досліди.

{\footnotesize
\setlength{\tabcolsep}{3pt}\renewcommand{\arraystretch}{1.15}
\begin{longtable}{>{\raggedright}p{0.5cm}>{\raggedright}p{4.0cm}>{\raggedright}p{5.6cm}>{\raggedright}p{2.3cm}>{\raggedright\arraybackslash}p{1.7cm}}
\toprule
№ & Твердження & Дослід і що виміряно & Джерело & Статус \\
\midrule\endfirsthead
\toprule
№ & Твердження & Дослід і що виміряно & Джерело & Статус \\
\midrule\endhead
1 & Знак дії серотоніну вибирає рецептор отримувача (твердження~\ref{prop:receptor}) & Рецептори серотоніну поділено на родини за фармакологією і механізмом: 5-HT$_{1A}$ через G-білок відкриває калієві канали і гальмує клітину, 5-HT$_{2A}$ та іонотропний 5-HT$_3$ збуджують. Один медіатор дає протилежні відповіді в різних клітинах. & \cite{BarnesSharp1999} & виміряно \\
2 & \nt{SERO}-нейрон у стані неспання рівно видає кілька імпульсів на секунду & Запис нейронів дорсального ядра шва у вільно рухливих котів: повільний ритмічний розряд $2.8\pm0.2$ імп/с у спокійному неспанні; у повільному сні частота падає на $34$--$68\,\%$, у швидкому --- на $98\,\%$. Під наркозом у щурів --- $1.7\pm0.5$ імп/с, дуже регулярно. & \cite{TrulsonJacobs1979, GagnonParent2014, JacobsAzmitia1992} & виміряно \\
3 & \nt{SERO}-нейрон --- ритмоводій: поштовх на кожен імпульс не потрібен, але потрібен рівний фоновий приток (у зрізі --- норадреналін через $\alpha_1$-рецептори, в організмі --- від \nt{NORA}) & У зрізі мозку клітини розряджаються повільно і рівно, після кожного імпульсу --- слідова гіперполяризація на $200$--$800\ms$. Спонтанно активних клітин у зрізі менше, ніж в організмі: рівному ритму потрібен тонічний вхід норадреналіну через $\alpha_1$-рецептори, блокада яких його гасить. & \cite{VandermaelenAghajanian1983} & виміряно \\
4 & Майже всі рецептори метаботропні; відповідь повільна: наростає і спадає в межах приблизно секунди & Усі родини рецепторів, крім 5-HT$_3$, пов'язані з G-білками. У зрізі викликаний викид серотоніну дає через 5-HT$_{1A}$ гальмівний струм, який наростає і спадає в межах секунди. & \cite{BarnesSharp1999, CourtneyFord2016} & виміряно \\
5 & В областях призначення серотонін виходить в об'єм, а не лише в щілину; у самому ядрі шва гальмування сусідів іде точка-точка & Швидка циклічна вольтамперометрія у зрізах: викид на один імпульс однаковий в області із синапсами і в ядрі шва, де синапсів майже немає; не залежить від блокади захоплення і рецепторів; молекул на закінчення набагато менше, ніж переносників і рецепторів, а позаклітинна концентрація збігається зі спорідненістю головного рецептора. У ядрі шва гальмування через 5-HT$_{1A}$ іде точка-точка: переносник не дає серотоніну накопичитися поза щілиною. & \cite{Bunin1998, CourtneyFord2016} & виміряно \\
6 & Концентрація за~\eqref{eq:seroneuron} спадає через відкачування; $\tau_c$ порядку секунди --- оцінка & Вольтамперометрія в живому мозку: позаклітинний серотонін обмежують насамперед зворотне захоплення і руйнування, а не синтез. Прямого вимірювання $\tau_c$ у цитованих роботах немає; порядок величини --- оцінка розділу. & \cite{Hashemi2012serotonin} & виміряно; $\tau_c$ --- оцінка \\
7 & НЕ і АБО-НЕ на одному ритмоводії; повнота \nt{SERO}-логіки (твердження~\ref{prop:serocomplete}, рис.~\ref{fig:serogates}) & Випливає з~\eqref{eq:seroneuron}: гальмований ритмоводій --- АБО-НЕ, а АБО-НЕ функціонально повний. Досліду, де з \nt{SERO}-нейронів зібрано логіку, немає; умова окремих об'ємів у мозку не виконується. & виведено в розділі & теорія \\
8 & Серотонін гальмує сам \nt{SERO}-нейрон через рецептори на ньому & У щурів під наркозом агоністи 5-HT$_{1A}$ внутрішньовенно та іонофоретично гальмують розряд нейронів дорсального шва з тією самою залежністю доза--відповідь, що й сам серотонін; агоністи 5-HT$_{1B}$ майже не діють. У зрізі ендогенний серотонін сусідніх клітин дає через 5-HT$_{1A}$ струм і коротку паузу в розряді, не змінюючи середньої частоти. & \cite{Sprouse1987, CourtneyFord2016} & виміряно \\
9 & У моделі петля через спільний об'єм --- регулятор рівня: розкид і стала часу зменшуються в $1+G$ разів~\eqref{eq:feedback}; що петля так працює в мозку, --- гіпотеза & Класична формула підсилювача з від'ємним зворотним зв'язком; у розділі виведена заново. Підсилення петлі $G$ у \nt{SERO}-системи не виміряно. Виміряно: у зрізі автогальмування іде точка-точка, дає коротку паузу і не змінює середньої частоти. & \cite{Black1934feedback, CourtneyFord2016} & теорія; у мозку --- гіпотеза \\
10 & Концентрація --- ковзне середнє частоти, фільтр нижніх частот~\eqref{eq:lowpass}, рис.~\ref{fig:lowpass} & Розв'язок лінійного рівняння~\eqref{eq:seroneuron}; рисунок --- розрахунок за цією формулою при $\tau_c=1\,\text{с}$. Окремої моделі в \texttt{models/} немає. & виведено в розділі & теорія \\
11 & Звивистість щілин сповільнює дифузію малої молекули приблизно в $2.6$ раза & Метод точкового джерела: іонофорез тетраметиламонію і запис його концентрації іоноселективним електродом у зрізах і в живому мозку. Звивистість $\lambda\approx1.6$, тобто ефективний $D$ менший за вільний у $\lambda^2\approx2.6$ раза; частка міжклітинного об'єму близько $0.2$. Коефіцієнт дифузії самого серотоніну в тканині в цих роботах не вимірювали; у розділі він --- оцінка. & \cite{Sykova2008} & виміряно \\
12 & Довжина незалежного «проводу» $\ell\sim\sqrt{D\tau}\approx20$~мкм~\eqref{eq:diffusionlength}; кількість проводів обмежена кількістю таких областей (твердження~\ref{prop:volumewires}) & Закон дифузії і підстановка оцінок $D$ і $\tau$ (рядки~6 і~11); твердження~\ref{prop:volumewires} --- наслідок. Прямого досліду, який рахує незалежні канали об'ємної передачі, немає. & виведено в розділі & теорія \\
13 & Вихід одного \nt{SERO}-нейрона розкинутий по величезних ділянках мозку; місць викиду, за оцінкою, $\sim10^5$ & Повна реконструкція окремих нейронів дорсального шва в щура: усі висхідні аксони сильно розгалужуються, загальна довжина аксона однієї клітини до $18.7$~см, гілки однієї клітини сягають стріатума, префронтальної кори і мигдалеподібного тіла. Кількість місць викиду на клітину прямо не підраховано. & \cite{GagnonParent2014} & виміряно; $10^5$ --- оцінка \\
14 & \nt{SERO}-нейрони поділяються на кілька підсистем із різними виходами і відповідями & Вірусно-генетичне мічення в мишей: нейрони, що йдуть до мигдалеподібного тіла і до лобової кори, майже не перетинаються за розгалуженням, отримують різні входи і протилежно відповідають на неприємний стимул; увімкнення і вимкнення кожної групи по-різному змінює поведінку. & \cite{Ren2018} & виміряно \\
15 & Умова терпіння $D<\tau_\gamma\ln(R/r_0)$~\eqref{eq:patience} за експоненційного знецінення; за виміряного, гіперболічного, --- $D<\tau_\gamma(R/r_0-1)$ & Випливає зі знецінення $e^{-D/\tau_\gamma}$ або $1/(1+D/\tau_\gamma)$; виведено в розділі. Мавпи, вибір на затримках у секунди: знецінення ближче до гіперболи, ніж до експоненти; відповідь \nt{DOPA}-нейронів на провісник відкладеної винагороди спадає із затримкою теж за гіперболою. & виведено в розділі; \cite{KobayashiSchultz2008} & теорія \\
16 & \nt{SERO} задає горизонт $\tau_\gamma$ (підрозділ~\ref{sec:horizon}) & Оптогенетичне ввімкнення \nt{SERO}-нейронів дорсального шва в мишей подовжує очікування відкладеної винагороди і підвищує наполегливість у пошуку винагороди, що стала рідкісною. У людей зниження серотоніну (дієта без триптофану) робить знецінення відкладеної винагороди крутішим. Як концентрація перетворюється на $\tau_\gamma$, невідомо. & \cite{Doya2002, Miyazaki2014, Lottem2018, Schweighofer2008} & гіпотеза \\
\bottomrule
\end{longtable}}

\section{Розділ~\ref{sec:dofa}. Навчання з \nt{DOPA}}

Механізми синапсу (порції, детектор збігів, кальцій, вікна пластичності) і поведінку дофаміну як помилки передбачення виміряно прямо; те, що правила ведуть до градієнта і чого коштує широкомовлення, --- теореми; підсумок навчання вибору --- розрахунок за моделлю книги; схема, що обчислює помилку, --- гіпотеза.

{\footnotesize
\setlength{\tabcolsep}{3pt}\renewcommand{\arraystretch}{1.15}
\begin{longtable}{>{\raggedright}p{0.5cm}>{\raggedright}p{4.0cm}>{\raggedright}p{5.6cm}>{\raggedright}p{2.3cm}>{\raggedright\arraybackslash}p{1.7cm}}
\toprule
№ & Твердження & Дослід і що виміряно & Джерело & Статус \\
\midrule\endfirsthead
\toprule
№ & Твердження & Дослід і що виміряно & Джерело & Статус \\
\midrule\endhead
1 & Зворотного поширення помилки в тій формі, що в штучних мережах, у мозку не знайдено & Прямого досліду немає: це твердження про відсутність. Огляди розбирають, чого вимагає алгоритм (зворотні зв'язки, що повторюють прямі, окрема фаза) і які біологічні наближення до нього запропоновано. & \cite{Lillicrap2020, WhittingtonBogacz2019} & гіпотеза \\
2 & Вага розкладається на кількість місць викиду, імовірність викиду і відповідь на порцію: $w=N_s\,p\,A_1$~\eqref{eq:quantal} & Нервово-м'язовий синапс жаби за зниженого викиду: відповідь отримувача коливається сходинками, кратними спонтанній мініатюрній відповіді на одну порцію; кількість порцій на імпульс підлягає закону Пуассона. & \cite{delCastillo1954} & виміряно \\
3 & Рецептор отримувача --- детектор збігів; кальцій запускає зміну ваги & Патч-клемп нейронів миші: іони Mg$^{2+}$ закупорюють канал, відкритий глутаматом, за потенціалу спокою і виходять за деполяризації. Зрізи гіпокампа: хелатор кальцію в отримувачі блокує довготривале посилення, а вивільнення кальцію світлом усередині клітини саме посилює передачу. & \cite{Nowak1984, Malenka1988calcium} & виміряно \\
4 & Рішення виконує будь-який бік: росте $A_1$ в отримувача, змінюється $p$ у відправника & Глутамат, вивільнений світлом над одним шипиком, спричиняє швидкий і стійкий ріст шипика і струму через його рецептори; посилення супроводжується вбудовуванням AMPA-рецепторів. Деполяризація отримувача на певний час пригнічує викид у відправника; ефект переносять ендоканабіноїди, що йдуть назад через щілину (показано на \nt{GABA}-входах). Короткочасне виснаження залежить від імовірності викиду у відправника. & \cite{Matsuzaki2004, Malinow2002, WilsonNicoll2001, Tsodyks1997} & виміряно \\
5 & Синапс посилюється, коли відправник і отримувач активні разом~\eqref{eq:hebb} & Кролик під наркозом: після коротких серій імпульсів вхідним шляхом гіпокампа відповідь посилена на час від $30$~хв до $10$~год. У зрізі гіпокампа імпульси отримувача посилюють лише ті синапси, які в ту саму мить активні. & \cite{Bliss1973LTP, Kelso1986hebb} & виміряно \\
6 & Правило~\eqref{eq:oja} знаходить головний власний вектор кореляцій входів (твердження~\ref{prop:oja}) & Теорема; доведення в розділі. Досліду, де нейрон вивчив головну компоненту своїх входів, немає; механізм обмеження росту ваг у синапсі не встановлено. & \cite{Oja1982} & теорія \\
7 & Гебб за часом: вікно близько $\pm20\ms$ (рис.~\ref{fig:stdp}); з повторюваних послідовностей виростають ланцюжки & Пари нейронів гіпокампа в культурі: імпульс отримувача в межах $20\ms$ після імпульсу відправника дає стійке посилення, у межах $20\ms$ до --- послаблення; обидва залежать від NMDA-рецепторів. Відношення $A_-=0.6A_+$ на рисунку --- ілюстрація. Що так у мережі виростають ланцюжки, прямо не виміряно. & \cite{BiPoo1998} & виміряно \\
8 & Слід~\eqref{eq:threefactor}: дофамін, що надійшов через $1$--$2\,\text{с}$ після спільної активності, посилює зв'язок, пізніше --- ні (твердження~\ref{prop:threefactor}) & Зрізи стріатума, роздільна оптогенетична стимуляція глутаматних і дофамінових входів: шипик росте, лише якщо дофамін надходить через $0.3$--$2\,\text{с}$ після глутаматного входу; вікно задається швидким підйомом і розпадом цАМФ. У корі спарювання лишає окремі сліди для посилення і послаблення, які моноамінові рецептори перетворюють на довготривалі зміни у вікні порядку секунд. & \cite{Yagishita2014, He2015eligibility} & виміряно \\
9 & Вікна завдовжки в секунди бувають і без дофаміну: третій сигнал --- сильне збудження отримувача & Жива миша, поле CA1: одна плато-подія в отримувачі посилює входи, що надійшли за секунди до і після неї, і за один прохід створює поле місця. & \cite{Bittner2017} & виміряно \\
10 & Середній крок правила трьох множників іде за градієнтом винагороди~\eqref{eq:perturbation} (твердження~\ref{prop:gradient}) & Теорема про градієнт для навчання за випадковими відхиленнями; у розділі виведена для малого шуму. & \cite{Williams1992} & теорія \\
11 & Шум --- це пошук: мережа навчається на своїх випадкових відхиленнях & Молоді амадини: вимкнення ядра, що виходить із кола базальних гангліїв, різко і зворотно прибирає мінливість пісні. Дорослі амадини: завада подається на виконання складу з висотою по один бік від середньої, і птахи швидко зсувають висоту в інший бік --- навчаються за власним розкидом. & \cite{Olveczky2005, TumerBrainard2007} & виміряно \\
12 & Вибір з двох навчається при $\tau_e=1\,\text{с}$ і $\delta=R-\bar R$; не навчається при $\tau_e=50\ms$; без віднімання ваги ростуть без меж, а за великої швидкості навчання мережа може закріпити гірший вибір & \texttt{models/dofa\_learning.py}, $\eta=0.05$, $500$ спроб, $6$ зерен. $\tau_e=1\,\text{с}$: частка виборів $A$ $0.65$--$0.79\to0.96$--$1.00$, ваги $0.85$--$1.33$. $\tau_e=50\ms$: $0.38$--$0.55$, ваги не змінюються. $\delta=R$: вага переможця $860$--$1190$. $\delta=R$, $\eta=0.5$, $3$ зерна: одне закріпилося на $B$ ($0.04\to0$). Скрипт друкує всі чотири випадки; перерахунок збігся з розділом. & \texttt{models/\allowbreak dofa\_\allowbreak learning.py} & модель \\
13 & Час навчання за одним спільним сигналом росте пропорційно до кількості шумних нейронів (твердження~\ref{prop:broadcastcost}) & Криві навчання лінійної мережі: збурення нейронів повільніше за точний градієнт у стільки разів, скільки вихідних нейронів, збурення ваг --- ще в стільки разів, скільки входів. & \cite{Werfel2005} & теорія \\
14 & Виходи \nt{DOPA} ідуть насамперед до областей вибору дій & Повна реконструкція аксонів окремих нігростріатних \nt{DOPA}-нейронів щура: гілки однієї клітини покривають $0.45$--$5.7\,\%$ об'єму стріатума (у середньому $2.7\,\%$), поза стріатумом гілок майже немає. & \cite{Matsuda2009} & виміряно \\
15 & Кора навчається передбачати свій вхід, помилка рахується на місці & Огляд: у сенсорній корі є відповіді на розбіжність очікуваного і того входу, що надійшов. Що це загальне правило навчання кори, не доведено. & \cite{Keller2018} & гіпотеза \\
16 & Дофамін поводиться як помилка~\eqref{eq:ctd}: сплеск, перенесення на провісник, провал (рис.~\ref{fig:ctdcases}) & \nt{DOPA}-нейрони мавпи: сплеск на несподіваний сік; після навчання --- сплеск на провісник і немає відповіді на обіцяний сік; провал, коли сік не надійшов. Неперервна форма~\eqref{eq:ctd} --- теорія. & \cite{Schultz1997, Doya2000} & виміряно \\
17 & Схема, що обчислює $\dd V/\dd t$ і віднімає очікування & Миші, запис і оптогенетика у вентральній покришці: \nt{DOPA}-нейрони віднімають очікування від винагороди; сусідні \nt{GABA}-нейрони гальмують їх, коли винагорода очікується, і їхнє збудження чи гальмування змінює помилку. Як обчислюється похідна, не показано. & \cite{Eshel2015arithmetic} & гіпотеза \\
18 & \nt{DOPA}-нейрон --- ритмоводій; провали мілкіші за сплески & У зрізі \nt{DOPA}-нейрони розряджаються спонтанно і дуже регулярно, на повільному пейсмекерному потенціалі. У мавпи частота кількісно йде за додатною помилкою, а від'ємну передає лише слабко. & \cite{GraceOnn1989, Bayer2005} & виміряно \\
19 & Актор і критик (рис.~\ref{fig:actorcritic}) & Алгоритм --- з теорії навчання з підкріпленням. У людей (фМРТ) черевний стріатум поводиться як критик і під час вибору, і без нього, спинний --- лише коли треба вибирати дії. & \cite{SuttonBarto2018, ODoherty2004actor} & гіпотеза \\
20 & Знак $\delta$ у правилі перевернутий у нейронів із другою родиною рецепторів & Зрізи стріатума: у нейронів із D1-рецепторами спарювання дає посилення, якому потрібен D1; у нейронів із D2 --- послаблення, якому потрібен D2. & \cite{Shen2008} & виміряно \\
\bottomrule
\end{longtable}}

\section{Розділ~\ref{sec:dofaalt}. Інші теорії \nt{DOPA}}

Кожна з розширених картин спирається на свої прямі вимірювання дофаміну (розкид точок перевертання, відповіді на неприємне і нове, повільне зростання, відповіді на зміну смаку), але формули, які їх пояснюють, --- теорії, і щодо двох із них досліди поки суперечать одне одному.

{\footnotesize
\setlength{\tabcolsep}{3pt}\renewcommand{\arraystretch}{1.15}
\begin{longtable}{>{\raggedright}p{0.5cm}>{\raggedright}p{4.0cm}>{\raggedright}p{5.6cm}>{\raggedright}p{2.3cm}>{\raggedright\arraybackslash}p{1.7cm}}
\toprule
№ & Твердження & Дослід і що виміряно & Джерело & Статус \\
\midrule\endfirsthead
\toprule
№ & Твердження & Дослід і що виміряно & Джерело & Статус \\
\midrule\endhead
1 & Асиметричне правило~\eqref{eq:asymmetric} сходиться до точки~\eqref{eq:expectile}; при $\kappa=1/2$ це середнє, для краплі з імовірністю $1/2$ --- $V=\kappa$ (рис.~\ref{fig:expectiles}) & Точка~\eqref{eq:expectile} --- експектиль, розв'язок задачі найменших квадратів із різними вагами додатних і від'ємних відхилень; для прикладу розділу виведено в самому розділі. & \cite{NeweyPowell1987}; виведено в розділі & теорія \\
2 & У різних \nt{DOPA}-нейронів різні точки перевертання, і вони впорядковані за асиметрією (твердження~\ref{prop:expectile}) & Миші, оптогенетично мічені \nt{DOPA}-нейрони вентральної покришки, винагороди різного розміру: точка, де сплеск змінюється провалом, у різних нейронів різна; у нейронів із більшою асиметрією відповідей вона вища; за відповідями групи відновлюється форма розподілу винагород. Що це головна функція розкиду --- гіпотеза. & \cite{Dabney2020} & виміряно \\
3 & Частина \nt{DOPA}-нейронів відповідає на неприємне сплеском & Мавпи: одні \nt{DOPA}-нейрони збуджуються на провісник винагороди і гальмуються на провісник струменя повітря в око, інші збуджуються на обидва; групи розташовані в різних частинах середнього мозку. & \cite{Matsumoto2009, BrombergMartin2010} & виміряно \\
4 & Модуль помилки або новизна задає швидкість навчання & Прямого досліду, де сигнал важливості в \nt{DOPA}-нейронів змінює швидкість навчання, у цитованих роботах немає; огляд пропонує роль сигналу «увага, важливо». & \cite{BrombergMartin2010} & гіпотеза \\
5 & Окремий провід для осі небезпеки & Миші: \nt{DOPA}-нейрони, що йдуть до заднього кінця стріатума, відповідають на нові й сильні подразники, а не на винагороду; їхнє руйнування послаблює уникання нового предмета. & \cite{Menegas2018} & виміряно \\
6 & Оптимальний час дії $\tau_a^*=\sqrt{C/\bar R}$~\eqref{eq:vigor} & Мінімум суми $C/\tau_a+\bar R\tau_a$; виведено в розділі. У вихідній моделі швидкість дій задає ціна часу, що дорівнює середній винагороді. & \cite{Niv2007}; виведено в розділі & теорія \\
7 & Повільний рівень дофаміну несе $\bar R$ і задає швидкість дій (твердження~\ref{prop:multiplex}) & Щури, мікродіаліз і вольтамперометрія в прилеглому ядрі: рівень дофаміну від хвилини до хвилини йде за частотою винагород і за швидкістю дій. У людей L-DOPA посилює залежність часу реакції від середньої частоти винагород. Що повільний рівень дорівнює саме $\bar R$, не доведено. & \cite{Hamid2016, Beierholm2013vigor} & гіпотеза \\
8 & Повільний рівень складається з двох джерел: викид змінюється біля виходів без зміни частоти імпульсів, але повільне зростання буває і в самій частоті & Щури, одна задача: викид у прилеглому ядрі йде за змінним очікуванням винагороди, а частота імпульсів упізнаних \nt{DOPA}-нейронів вентральної покришки --- ні. Миші у віртуальному коридорі (рядок~10): плавне зростання під час наближення до винагороди є і в імпульсах упізнаних \nt{DOPA}-нейронів. & \cite{Mohebi2019, Kim2020} & виміряно \\
9 & Дофамін плавно росте, поки тварина наближається до далекої винагороди & Щури в лабіринті, вольтамперометрія в стріатумі: концентрація дофаміну наростає за секунди в міру наближення до мети, крутизна залежить від відстані і розміру винагороди. & \cite{Howe2013ramp, Hamid2016} & виміряно \\
10 & Зростання --- помилка $\delta$ при уточненні оцінки~\eqref{eq:ramp}; стрибок у коридорі дає сплеск (рис.~\ref{fig:ramp}) & Формула і число $\delta=0.75\,V$ на секунду --- виведено в розділі. Дослід: миші у віртуальному коридорі, миттєве перенесення ближче до мети; імпульси тіл, кальцій тіл і виходів та концентрація в стріатумі відповідають як помилка, а не як очікування $V$. Яке з двох пояснень правильне на всіх виходах, обговорюється. & \cite{Kim2020} & виміряно \\
11 & Погляд назад: дофамін повідомляє міру причинного зв'язку~\eqref{eq:retro} & Миші, викид дофаміну в прилеглому ядрі в дослідах, де ця теорія і помилка~\eqref{eq:ctd} передбачають різне: у низці дослідів викид ішов за першою. & \cite{Jeong2022} & гіпотеза \\
12 & Відповідь дофаміну під час навчання поступово зсувається від винагороди до провісника & Миші, сигнал виходів \nt{DOPA}-нейронів у черевному стріатумі в ході навчання: сплеск поступово переходить з моменту винагороди на момент провісника, як вимагає навчання за~\eqref{eq:ctd}. & \cite{Amo2022} & виміряно \\
13 & \nt{DOPA}-нейрони відповідають на зміну смаку винагороди за тієї самої цінності & Щури, запис нейронів вентральної покришки: заміна соку іншим, рівним за цінністю, спричиняє відповідь, хоча помилка~\eqref{eq:ctd} дорівнює нулю. & \cite{Takahashi2017} & виміряно \\
14 & Штучні сплески вчать зв'язку двох нейтральних подразників & Щури, дослід із попереднім поєднанням двох подразників без винагороди: оптогенетичне ввімкнення \nt{DOPA}-нейронів на переході від першого до другого створює зв'язок, вимкнення заважає його вивчити. & \cite{Sharpe2017} & виміряно \\
15 & Вектор помилок за ознаками~\eqref{eq:vectortd} & Теорія помилки передбачення ознак; прямої перевірки векторної форми немає. & \cite{Gardner2018} & гіпотеза \\
16 & Різні \nt{DOPA}-нейрони в одній задачі несуть різні величини & Миші у віртуальному коридорі, двофотонна зйомка \nt{DOPA}-нейронів: одні пов'язані з винагородою, інші зі швидкістю і поворотом, треті з провісниками і точністю вибору. & \cite{Engelhard2019} & виміряно \\
17 & Джгут замість проводу (твердження~\ref{prop:bundle}) & Зведення рядків~2--16: кожне розкладання виміряно окремо; єдиної картини, що всі вони --- частини одного сигналу, дослідом не встановлено. & --- & гіпотеза \\
\bottomrule
\end{longtable}}

\section{Розділ~\ref{sec:nora}. \nt{NORA}}

Будову \nt{NORA}-системи, її два режими, зв'язок із зіницею (не лише через \nt{NORA}), зростання відношення сигналу до фону і перевернуте U в поведінці виміряно прямо; мультиплікативне підсилення $g$ --- модель; те, що підсилення перемикає режим вибору і працює як обернена температура, --- висновки розділу; виграш від підлаштування підсилення --- розрахунок за моделлю книги; «\nt{NORA} --- ручка пошуку» і «\nt{NORA} --- переривання» --- гіпотези.

{\footnotesize
\setlength{\tabcolsep}{3pt}\renewcommand{\arraystretch}{1.15}
\begin{longtable}{>{\raggedright}p{0.5cm}>{\raggedright}p{4.0cm}>{\raggedright}p{5.6cm}>{\raggedright}p{2.3cm}>{\raggedright\arraybackslash}p{1.7cm}}
\toprule
№ & Твердження & Дослід і що виміряно & Джерело & Статус \\
\midrule\endfirsthead
\toprule
№ & Твердження & Дослід і що виміряно & Джерело & Статус \\
\midrule\endhead
1 & \nt{NORA}-нейронів у людини кілька десятків тисяч, майже всі в одній групі & Серійні зрізи стовбура мозку людини, забарвлення на тирозингідроксилазу: $53\,900$ клітин у блакитній плямі і ще $6\,260$ у сусідньому під'ядрі. & \cite{Baker1989LC} & виміряно \\
2 & Виходи розходяться майже по всьому мозку, але частини групи частково обслуговують різні області & Вірусно-генетичне мічення в мишей: \nt{NORA}-нейрони блакитної плями отримують входи з широких областей мозку і надсилають виходи в широкі області головного і спинного мозку. У щурів: група, що йде до мигдалеподібного тіла, посилює вивчення страху, а окрема група, що йде до префронтальної кори, --- його згасання. & \cite{Schwarz2015LC, Uematsu2017LC, Poe2020} & виміряно \\
3 & Фоновий режим: рівні імпульси з частотою від часток до кількох на секунду, що змінюється з рівнем неспання & Щури, запис нейронів блакитної плями в циклі сну і неспання: частота найбільша в неспанні, нижча в повільному сні, майже нуль у швидкому; зміни частоти випереджають зміну стану. & \cite{AstonJonesBloom1981} & виміряно \\
4 & Уривчастий режим: коротка пачка на важливу для задачі подію, приблизно в момент рішення & Мавпи, задача знайти рідкісну ціль: усі нейрони блакитної плями відповідають пачкою лише на ціль, через $\sim90\ms$ після неї і за $\sim200\ms$ до відповіді рукою; за поганої роботи пачки слабші. У задачі з вибором із двох пачка точніше прив'язана до відповіді, ніж до стимулу, і є як за правильних, так і за помилкових відповідей. & \cite{AstonJones1994vigilance, Clayton2004LC} & виміряно \\
5 & Діаметр зіниці --- неточна контрольна точка \nt{NORA}: він іде і за \nt{NORA}, і за \nt{ACET}, і за іншими областями & Мавпи: імпульси блакитної плями, аж до окремих імпульсів однієї клітини, передбачають розширення зіниці; схожий, але пізніший зв'язок є в горбках і поясній корі. Миші: коливання зіниці йдуть за активністю і норадреналінових, і холінергічних виходів у корі. & \cite{Joshi2016, Reimer2016} & виміряно \\
6 & Норадреналін пригнічує фонову активність сильніше за викликану; мультиплікативне підсилення~\eqref{eq:gain} --- модель, що передає цей ефект & Мавпи в стані неспання, слухова кора, іонофорез норадреналіну: фонова активність нейронів пригнічується сильніше, ніж відповідь на голоси родичів, --- відношення сигналу до шуму зростає. Мультиплікативну сигмоїду~\eqref{eq:gain} взято з мережевої моделі; буквального множення крутизни не виміряно. & \cite{Foote1975NE, ServanSchreiber1990} & виміряно; $g$ --- модель \\
7 & Підсилення підвищує $d'$ лише проти шуму після підсилювача~\eqref{eq:gainsnr} (твердження~\ref{prop:gainnoise}) & Виведено в розділі; у мережевій моделі підсилення покращує виявлення сигналу. Досліду, який відокремлював би шум до і після підсилювача і перевіряв би цю різницю, немає. & виведено в розділі; \cite{ServanSchreiber1990} & теорія \\
8 & Межа $g\beta=1$ перемикає мережу вибору з «розмірковує» на «вирішила»~\eqref{eq:wtagain} (твердження~\ref{prop:gainwta}, рис.~\ref{fig:pitchfork}) & Лінійний аналіз двох груп, що гальмують одна одну; виведено в розділі. Прямих вимірювань цього порога в мозку немає. & виведено в розділі & теорія \\
9 & Коефіцієнт підсилення --- обернена температура вибору~\eqref{eq:softmax} & За логістичного шуму після підсилювача ймовірність вибору --- softmax з $T=\sigma/g$; виведено в розділі. & виведено в розділі & теорія \\
10 & \nt{NORA} задає, скільки шукати: висока фонова активність --- мале діюче $g$ (пошук), уривчаста пачка в момент рішення --- велике (рішення) & Люди: перед вибором навмання базова зіниця ширша, ніж перед вибором найкращого відомого; у людей із ширшою зіницею схильність шукати вища. Щури: перехід у режим випадкового вибору керується входом \nt{NORA} у передню поясну кору. Що пачка підвищує діюче $g$, прямо не виміряно. & \cite{JepmaNieuwenhuis2011, Tervo2014stochastic, AstonJones2005} & гіпотеза \\
11 & Винагорода залежить від $g$ як перевернуте U; найкраще $g$ менше у світі, що швидко змінюється (твердження~\ref{prop:explore}, рис.~\ref{fig:invertedU}) & \texttt{models/nora\_explore.py}, $60$ прогонів по $4000$ спроб. Зміна раз на $1000$ спроб: найкраще $g=16$, частка $0.976$; раз на $20$: $g=8$, частка $0.886$; при $g=0.5$ --- $0.77$. Перерахунок збігся з розділом. & \texttt{models/\allowbreak nora\_\allowbreak explore.py} & модель \\
12 & Перевернуте U в поведінці: найкраще за помірної фонової активності з чіткими пачками & Мавпи, та сама задача, що в рядку~4: у періоди доброї роботи фонова частота помірна, а пачки на ціль чіткі; за високої фонової частоти пачок немає і хибних відповідей більше. Зміни фонової і викликаної активності йдуть за коливаннями якості роботи. & \cite{Usher1999LC, AstonJones2005} & виміряно \\
13 & \nt{NORA} повідомляє неочікувану невизначеність, \nt{ACET} --- очікувану & Теорія байєсівського висновку з двома видами невизначеності; прямого досліду, що розділяє ці сигнали за медіаторами, у цитованих роботах немає. & \cite{YuDayan2005} & гіпотеза \\
14 & Після різкої зміни люди навчаються швидше, і зіниця розширюється & Люди, задача передбачення з раптовими змінами: короткі зміни зіниці йдуть за оцінкою ймовірності зміни і разом із базовою зіницею передбачають, наскільки сильно нові дані змінюють оцінку (швидкість навчання); зміна зіниці, не пов'язана зі світлом і задачею, змінює цю швидкість. Що зіниця тут показує саме \nt{NORA}, --- висновок за непрямими даними рядка~5. & \cite{Nassar2012} & виміряно \\
15 & Уривчаста пачка працює як переривання: мережа скидає стан & Прямого досліду, де пачка \nt{NORA} скидає стан мережі, немає; виміряно лише пачки на важливі події (рядок~4). & \cite{BouretSara2005} & гіпотеза \\
16 & Підлаштування $g$ або швидкості навчання за несподіваністю майже не краще за найкращі сталі налаштування & \texttt{models/nora\_explore.py}, світ з епохами по $1000$ спроб (зміна раз на $1000$ і раз на $20$). Найкраще стале $g=8$: $0.925$; підлаштування $g$: до $0.927$; підлаштування швидкості навчання: до $0.926$; стала швидкість $0.4$: $0.928$. Перерахунок збігся з розділом. & \texttt{models/\allowbreak nora\_\allowbreak explore.py} & модель \\
\bottomrule
\end{longtable}}

\section{Розділ~\ref{sec:acet}. Додаємо \nt{ACET}}

Три дії ацетилхоліну виміряно на зрізах і в живому мозку, конфлікт запису і читання показано на моделі книги, а те, що \nt{ACET} слугує лінією дозволу запису і повідомляє очікувану невизначеність, --- гіпотеза з частковою опорою на дослід.

{\footnotesize
\setlength{\tabcolsep}{3pt}\renewcommand{\arraystretch}{1.15}
\begin{longtable}{>{\raggedright}p{0.5cm}>{\raggedright}p{4.0cm}>{\raggedright}p{5.6cm}>{\raggedright}p{2.3cm}>{\raggedright\arraybackslash}p{1.7cm}}
\toprule
№ & Твердження & Дослід і що виміряно & Джерело & Статус \\
\midrule\endfirsthead
\toprule
№ & Твердження & Дослід і що виміряно & Джерело & Статус \\
\midrule\endhead
1 & \nt{ACET}-нейронів мозку небагато, вони зібрані в кількох групах, а виходи розходяться по всій корі: широкомовний канал & Забарвлення антитілами до ферменту синтезу ацетилхоліну і ретроградна мітка в щура: кора, гіпокамп, мигдалеподібне тіло, нюхова цибулина і таламус отримують ацетилхолін переважно від шести компактних груп клітин базального переднього мозку і стовбура & \cite{Mesulam1983} & виміряно \\
2 & Рівень ацетилхоліну в корі високий у неспанні і низький у повільному сні & Мікродіаліз у корі і гіпокампі вільно рухливих котів із записом ЕЕГ: викид ацетилхоліну в спокійному неспанні на $\sim100\%$, в активному --- на $\sim175\%$ вищий, ніж у повільному сні; у швидкому сні в корі --- як у неспанні & \cite{Marrosu1995,Hasselmo1999} & виміряно \\
3 & Зміст сигналу задає рецептор: ацетилхолін має швидкі рецептори-канали і повільні, що запускають реакції в клітині (твердження~\ref{prop:receptor}) & Огляд вимірювань: дія ацетилхоліну на збудливість, передачу і пластичність залежить від місця викиду, підтипу рецептора (нікотинові --- іонні канали, мускаринові --- через G-білки) і клітини-отримувача & \cite{Picciotto2012} & виміряно \\
4 & Перша дія: послабити внутрішні зв'язки, майже не зачіпаючи входів ззовні (множник $1-a$ у~\eqref{eq:acetnet}) & Зрізи нюхової кори щура: за $100$\,мкМ агоністів ацетилхоліну потенціали внутрішніх зв'язків (шар Ib) падають більш ніж на $60\%$, зовнішнього входу (шар Ia) --- менш ніж на $15\%$, у тій самій клітині; пригнічення пресинаптичне. Зрізи гіпокампа (CA1): $89\%$ проти $40\%$ для внутрішнього і вхідного шляхів & \cite{HasselmoBower1992,HasselmoSchnell1994} & виміряно \\
5 & Друга дія: посилити вхід (множник $\frac{1+a}{2}$) & Зрізи неокортексу: нікотинові рецептори посилюють синапси входу з таламуса і не діють на внутрішньокіркові; мускаринові послаблюють обидва типи. Збудливість нейронів ацетилхолін підвищує (огляд) & \cite{Gil1997,Hasselmo2006} & виміряно \\
6 & Третя дія: полегшити зміну ваг (швидкість $\eta a$) & Щур: електрична стимуляція базального ядра (джерела ацетилхоліну для кори) в парі з тоном спричиняє в дорослої тварини сильну перебудову карти слухової кори під звук, що пред'являється & \cite{KilgardMerzenich1998,Hasselmo2006} & виміряно \\
7 & Правило Гебба для рідкісних образів у другому рівнянні~\eqref{eq:acetnet} дає асоціативну пам'ять, що дописує образ за частиною & Розрахунок ємності мережі з розрідженими образами і коваріаційним правилом: за малої частки активних нейронів $p$ мережа зберігає помітно більше образів, ніж при $p=1/2$ & \cite{Tsodyks1988} & теорія \\
8 & Запис за низького $a$ псує пам'ять: мережа записує згадане, а не побачене; при $a=0.1$ псування не слабше, ніж при $a=0$ & \texttt{models/\allowbreak acet\_memory.py}: $200$ нейронів, п'ять образів по $20$, новий образ ділить з одним із них $10$ нейронів, $10$ прогонів. При $a=0$ новий образ не запам'ятано, старий тягне $5.9$ чужого нейрона з $10$; при $a=0.1$ новий образ згадується ($0.97$), але обидва зливаються: новий тягне $7.3$ чужого нейрона, старий --- $7.9$; з $a=0.2$ --- менше за один & виведено в розділі & модель \\
9 & Твердження~\ref{prop:writeread}: немає рівня $a$, за якого однаково добрі запис і читання (рис.~\ref{fig:acetsim}) & Той самий скрипт, швидкість навчання $\eta a$: новий образ запам'ятовується за одне пред'явлення цілком при $a\ge0.9$, на $0.8$ при $a=0.75$, не запам'ятовується при $a\le0.5$. Читання за половиною образу: $1.0$ при $a\le0.2$, $0.96$ при $a=0.3$, $0.04$ при $a=0.4$ & виведено в розділі & модель \\
10 & У мозку один широкомовний провід --- \nt{ACET} --- перемикає запис і читання & Ідею взято з моделей, побудованих за вимірюваннями на зрізах. Найпряміший дослід --- у людини: блокатор мускаринових рецепторів скополамін погіршує заучування нових пар слів, найсильніше тих, що перетинаються зі старими, але не згадування вже вивчених пар. Прямого вимірювання двох режимів мережі немає & \cite{HasselmoAndersonBower1992,Atri2004} & гіпотеза \\
11 & Фільтр~\eqref{eq:kalman} оптимальний; вага входу і швидкість навчання --- одна величина $P/R$, в усталеному режимі $P=\sqrt{qR}$ і пам'ять $\sqrt{R/q}$ (твердження~\ref{prop:kalman}) & Досліду не потребує: оптимальність фільтра Калмана--Б'юсі --- класичний результат теорії оцінювання; усталений режим виведено в розділі & \cite{KalmanBucy1961} & теорія \\
12 & \nt{ACET} повідомляє мережі величину, подібну до $P$, --- очікувану невизначеність & Запропоновано в імовірнісній моделі уваги. Прямого вимірювання того, що рівень ацетилхоліну йде за невизначеністю оцінки, немає & \cite{YuDayan2005} & гіпотеза \\
13 & Частина \nt{ACET}-нейронів відповідає на винагороду і покарання за два десятки мілісекунд, тим сильніше, чим несподіваніше підкріплення & Миша, оптогенетично впізнані холінергічні нейрони базального переднього мозку: відповідь на винагороду і покарання через $18\pm3$\,мс, величина зростає з несподіваністю підкріплення, відповіді схожі в нейронів двох ядер & \cite{Hangya2015} & виміряно \\
14 & \nt{NORA} помічає неочікувану зміну світу, \nt{ACET} тримає мережу в режимі запису & Розподіл праці запропоновано в моделі. Непрямо: у людини діаметр зіниці, пов'язаний із \nt{NORA}, відстежує зміни і швидкість навчання. Прямої перевірки пари \nt{NORA}--\nt{ACET} немає & \cite{YuDayan2005,Nassar2012} & гіпотеза \\
15 & У повільному сні мережа в режимі читання програє записане вдень, і ці повтори переписують його в довготривалу пам'ять & Записи ансамблів нейронів гіпокампа щура: клітини, що працювали разом удень, частіше спрацьовують разом у наступному повільному сні і в тому самому порядку --- виміряно. Для перезапису: пригнічення брижів гіпокампа після навчання погіршує просторову пам'ять, а подовження спонтанних брижів її покращує; що причина --- низький рівень \nt{ACET}, не доведено & \cite{WilsonMcNaughton1994,Skaggs1996,Girardeau2009,FernandezRuiz2019,Hasselmo1999} & гіпотеза \\
\bottomrule
\end{longtable}}

\section{Розділ~\ref{sec:synthesis}. Уся машина}

Розділ-зведення нових дослідів не наводить: кожен його рядок спирається на розділ, де твердження розібрано, і на ті самі джерела; твердо встановлено шину \nt{GLUT} і керування потоком через \nt{GABA}, а ролі широкомовних каналів та їхня спільна робота --- здебільшого гіпотеза.

{\footnotesize
\setlength{\tabcolsep}{3pt}\renewcommand{\arraystretch}{1.15}
\begin{longtable}{>{\raggedright}p{0.5cm}>{\raggedright}p{4.0cm}>{\raggedright}p{5.6cm}>{\raggedright}p{2.3cm}>{\raggedright\arraybackslash}p{1.7cm}}
\toprule
№ & Твердження & Дослід і що виміряно & Джерело & Статус \\
\midrule\endfirsthead
\toprule
№ & Твердження & Дослід і що виміряно & Джерело & Статус \\
\midrule\endhead
1 & На $8.6\cdot10^{10}$ нейронів припадає лише чотири види керувальних сигналів~\eqref{eq:machine} & Підрахунок ядер клітин у гомогенаті мозку (ізотропний фракціонатор): у дорослого чоловіка $86.1\pm8.1$ млрд нейронів & \cite{Azevedo2009} & виміряно \\
2 & Твердження~\ref{prop:architecture}, перші два шари: дані йдуть адресною шиною \nt{GLUT}, знак зв'язку задає відправник, потік спрямовують і нормують \nt{GABA}-нейрони (розділи~\ref{sec:neurontypes}, \ref{sec:gaba}) & Один головний медіатор на нейрон (огляд вимірювань); пряме гальмування звужує вікно, в якому пірамідний нейрон додає входи; ділення на загальну активність виміряно в багатьох ділянках & \cite{Strata1999,Pouille2001,Carandini2012} & виміряно \\
3 & Елементи ненадійні: синапс губить більшість імпульсів (табл.~\ref{tab:computer}, розділ~\ref{sec:code}) & Зрізи гіпокампа, мінімальна стимуляція: понад половина імпульсів не викликає відповіді в CA1; збої порога й проведення аксоном виключено, причина --- імовірнісне вивільнення медіатора & \cite{AllenStevens1994} & виміряно \\
4 & Мозок працює на $\sim20$ ватах, у середньому близько імпульсу за секунду на нейрон (розділ~\ref{sec:code}); такої машини інженери поки не збудували & Оцінка~\eqref{eq:energy} з виміряних витрат: близько $10^9$ молекул АТФ на імпульс разом із роботою синапсів (кора гризунів); детальна оцінка для кори дає ще меншу частоту. Стан техніки --- за оглядом & \cite{Attwell2001,Lennie2003,MehonicKenyon2022} & теорія \\
5 & Навчання більшості синапсів --- добуток місцевого сліду й одного спільного числа $\delta$ (друге рівняння~\eqref{eq:machine}, розділ~\ref{sec:dofa}) & \nt{DOPA}-нейрони мавпи відповідають на помилку передбачення винагороди; у зрізах смугастого тіла дофамін збільшує шипик, лише якщо надходить через $0.3$--$2$\,с після глутаматного входу --- слід виміряно & \cite{Schultz1997,Yagishita2014} & виміряно \\
6 & Окремий провід помилки до кожного виходу є лише в колі навчання рухів (розділ~\ref{sec:motorlearning}) & Мозочок: збіг сигналу ліаноподібного волокна з входом послаблює цей вхід; у мавпи складні імпульси клітин Пуркіньє несуть напрямок помилки руху й викликають поправку & \cite{Ito2001,Herzfeld2018} & виміряно \\
7 & Кожен широкомовний канал --- джгут із кількох ліній (твердження~\ref{prop:bundle}) & Для \nt{DOPA}: в одній задачі різні нейрони несуть винагороду, рух і ознаки стимулу; швидка помилка й повільний рівень дофаміну різні. Для \nt{SERO}, \nt{NORA}, \nt{ACET} такий розклад вивчено гірше & \cite{Engelhard2019,Hamid2016,Mohebi2019} & виміряно \\
8 & \nt{SERO}: регулятор рівня і, за гіпотезою, горизонт $\tau_\gamma$ (розділ~\ref{sec:serocomputer}) & Самогальмування: агоністи власних рецепторів 5-HT$_{1A}$ гальмують серотонінові нейрони шва --- виміряно. Горизонт запропоновано в теорії; найсильніший дослід --- оптогенетична активація цих нейронів у миші подовжує очікування відкладеної винагороди & \cite{Sprouse1987,Doya2002,Miyazaki2014} & гіпотеза \\
9 & \nt{NORA}: підсилення $g$ вибирає між пошуком і рішенням (розділ~\ref{sec:nora}) & Зростання відношення сигнал/шум: у мавпи в стані неспання норадреналін у слуховій корі пригнічує фонову активність сильніше, ніж відповідь на голоси родичів --- виміряно; мультиплікативне підсилення --- модель; роль у виборі між пошуком і рішенням --- теорія & \cite{Foote1975NE,ServanSchreiber1990,AstonJones2005} & гіпотеза \\
10 & \nt{ACET}: перемикає запис і читання (розділ~\ref{sec:acet}) & Три дії (послаблення внутрішніх зв'язків, підсилення входу, полегшення пластичності) виміряно; перемикання режимів опосередковано підтримане дослідом зі скополаміном у людини & \cite{HasselmoBower1992,Gil1997,Atri2004} & гіпотеза \\
11 & Формула~\eqref{eq:machine} описує машину цілком & Це зведення моделей розділів~\ref{sec:glutcomputer}--\ref{sec:acet}, кожна частина спирається на свій розділ; як ціле дослідом не перевірялася & виведення в розділі & гіпотеза \\
12 & Ручки працюють злагоджено без спільного керування: помилка, несподіванка, запис, повернення (рис.~\ref{fig:timeline}) & Досліду немає: схема якісна. Окремі ланки --- теорія розподілу праці \nt{NORA} і \nt{ACET} та зв'язок зіниці зі швидкістю навчання в людини & \cite{YuDayan2005,Nassar2012} & гіпотеза \\
13 & Навчання за одним спільним сигналом тим повільніше, що більше нейронів впливає на результат (твердження~\ref{prop:broadcastcost}) & Розрахунок кривих навчання лінійної мережі: навчання за збуренням виходів повільніше за градієнтний спуск у стільки разів, скільки виходів & \cite{Werfel2005} & теорія \\
14 & Як кора налаштовує багатошарові кола, невідомо; кандидати --- пачки імпульсів, обчислення в гілках нейрона, самонавчання передбаченням & Пачки як носій помилки --- модель; повторити відповідь детальної моделі нейрона кори змогла лише мережа з $5$--$8$ шарів --- модель; кальцієві імпульси в гілках нейронів кори людини, що дають виключне АБО, --- виміряно в зрізах; самонавчання --- огляд свідчень & \cite{Payeur2021,Beniaguev2021,Gidon2020,Keller2018} & гіпотеза \\
\bottomrule
\end{longtable}}

\section{Розділ~\ref{sec:sensors}. Входи машини}

Будову датчиків виміряно прямо --- записами струму поодиноких клітин, коливань завитки й підрахунком клітин сітківки; межа чутливості й формула дробового шуму випливають із фізики, а те, що коди датчиків підлаштовано під найбільшу інформацію, --- гіпотеза з сильною опорою.

{\footnotesize
\setlength{\tabcolsep}{3pt}\renewcommand{\arraystretch}{1.15}
\begin{longtable}{>{\raggedright}p{0.5cm}>{\raggedright}p{4.0cm}>{\raggedright}p{5.6cm}>{\raggedright}p{2.3cm}>{\raggedright\arraybackslash}p{1.7cm}}
\toprule
№ & Твердження & Дослід і що виміряно & Джерело & Статус \\
\midrule\endfirsthead
\toprule
№ & Твердження & Дослід і що виміряно & Джерело & Статус \\
\midrule\endhead
1 & Датчик --- перетворювач: рецептор дає струм, а вихід чутливих клітин ока й вуха --- глутамат на шину \nt{GLUT}; перші клітини імпульсів не видають (рис.~\ref{fig:transducer}) & Палички й колбочки черепахи виділяють збуджувальну амінокислоту: її вловлюють глутаматні канали на відокремленому клаптику мембрани. Внутрішні волоскові клітини завитки щура: струми в закінченнях нервових волокон ідуть через глутаматні AMPA-рецептори, частота вивільнень зростає з плавною деполяризацією клітини до $150$\,Гц & \cite{CopenhagenJahr1989,GlowatzkiFuchs2002} & виміряно \\
2 & Датчик світла в темряві відповідає на один фотон струмом близько пікоампера & Всмоктувальний електрод на зовнішньому сегменті палички ропухи: відповіді на тьмяні спалахи квантовані, подія $\approx1$\,пА ($5\%$ темнового струму) з розкидом $0.2$\,пА, ефективність $0.5$. Людина повідомляє про влучання одного фотона частіше, ніж випадково & \cite{Baylor1979,Tinsley2016} & виміряно \\
3 & Датчик звуку помічає зміщення менше за нанометр & Метод Мессбауера, основна мембрана завитки морської свинки в місці $16$--$19$\,кГц: на порозі відповіді слухового нерва швидкість мембрани близько $0.04$\,мм/с, тобто зміщення $\approx0.35$\,нм (наш перерахунок). Вимірювали мембрану, а не пучок датчика & \cite{Sellick1982,Hudspeth1989} & виміряно \\
4 & Точність у темряві обмежена дробовим шумом фотонів~\eqref{eq:photonnoise}; за слабкого світла накопичення довше & Формула випливає з пуассонівської статистики. У паличці шум за тьмяного світла збігається із сумою незалежних квантових подій; на яскравому тлі відповідь на спалах менша й коротша. Число «десяті частки секунди» тут окремо не перевірено & \cite{Baylor1979} & теорія \\
5 & Діапазон входу --- десять порядків для світла й дванадцять для звуку, а частота імпульсів --- менше за три порядки & Для звуку: відповідь основної мембрани на характеристичній частоті зростає з крутістю до $0.2$\,дБ/дБ у смузі $40$--$80$\,дБ, стиснення видно й вище $100$\,дБ. Самі числа порядків --- стандартні оцінки, у розділі без джерела & \cite{Ruggero1997} & виміряно \\
6 & Відповідь датчика~\eqref{eq:nakarushton}, а $\sigma$ іде за середньою яскравістю, зсуваючи криву логарифмічною віссю (рис.~\ref{fig:adaptation}) & Формулу вперше підібрано до S-потенціалів сітківки риби. Колбочки черепахи: відповіді підкоряються $V=I V_m/(I+\sigma)$ і охоплюють $\sim3.5$ порядку яскравості; після $2$--$3$ хвилин тла крива зсувається по горизонталі, зберігаючи ширину. Зв'язок із діленням на загальну активність --- огляд & \cite{NakaRushton1966,NormannPerlman1979,Carandini2012} & виміряно \\
7 & Датчики передають зміни, а їхні коди підлаштовано під найбільшу інформацію на імпульс (твердження~\ref{prop:sensors}) & Принцип запропоновано теоретично. Найсильніше свідчення: у мухи характеристика «контраст $\to$ відповідь» нейронів першого шару збігається з кумулятивним розподілом контрастів природних сцен --- так досягається найбільша інформація & \cite{Barlow1961,Laughlin1981} & гіпотеза \\
8 & Датчики ввімкнення й вимкнення --- двопровідний код; подієва камера повторює його в кремнії & Огляд дослідів: канали ввімкнення й вимкнення сітківки розділяються вже на першому синапсі й вимикаються окремо. Камера: $128\times128$ пікселів без кадрів, діапазон $120$\,дБ, затримка $15$\,мкс & \cite{Schiller1992,Lichtsteiner2008,Gallego2022} & виміряно \\
9 & В оці $\sim10^8$ датчиків світла й $\sim10^6$ вихідних нейронів: стиснення в $\sim100$ разів & Підрахунок на цілих препаратах сітківки людини: у середньому $92$ млн паличок і $4.6$ млн колбочок; вихідних (гангліозних) клітин від $0.7$ до $1.5$ млн у різних очах & \cite{Curcio1990,CurcioAllen1990} & виміряно \\
10 & У миші вихідних каналів ока понад тридцять & Миша: двофотонна кальцієва зйомка понад $11\,000$ вихідних клітин і кластеризація відповідей --- помітно більше за $30$ функціональних типів. Для людини число не виміряно & \cite{Baden2016} & виміряно \\
11 & Око морської свинки передає $\sim875\,000$ біт/с; перерахунок на око людини дає $\sim10^7$ біт/с & Морська свинка, багатоелектродна матриця, рух у природних сценах: $6$--$13$ біт/с на клітину, $\sim875\,000$ біт/с на $\sim10^5$ вихідних клітин. Для людини ($\sim10^6$ клітин) $\sim10^7$ біт/с --- перерахунок авторів & \cite{Koch2006} & виміряно \\
12 & Вухо --- банк смугових фільтрів із майже логарифмічною картою частот (рис.~\ref{fig:filterbank}) & Місце найбільшого коливання основної мембрани залежить від частоти; функція «частота--місце» майже експоненційна й однією формою описує завитки людини та живих тварин & \cite{Greenwood1990,Robles2001} & виміряно \\
13 & Резонатори активні: частина датчиків підсилює слабкі коливання в тисячі разів за амплітудою ($66$--$76$\,дБ), зі зростанням гучності підсилення падає & Лазерна віброметрія основи завитки шиншили: за слабких тонів на характеристичній частоті підсилення $66$--$76$\,дБ відносно стремінця; смерть знижує чутливість на $60$--$81$\,дБ (у $10^3$--$10^4$ разів за амплітудою) і прибирає стиснення. Джерело сили --- зовнішні волоскові клітини & \cite{Ruggero1997,Robles2001} & виміряно \\
14 & Підсилювач вуха стоїть біля межі самозбудження & Розрахунок: біля точки біфуркації Хопфа неминучі стиснення діапазону, гостре налаштування й комбінаційні тони --- усі відомі нелінійності слуху. Що завитка налаштована саме так, прямо не доведено & \cite{Eguiluz2000} & гіпотеза \\
15 & На низьких частотах імпульси йдуть у фазі зі звуком (у морської свинки прив'язка слабшає від $\sim600$\,Гц і помітна до $\sim3.5$\,кГц), і за ними знаходять напрямок; на високих лишається код місцем & Слуховий нерв морської свинки: прив'язка до фази починає падати близько $600$\,Гц і зникає вище $3.5$\,кГц; у кішки й мавпи межа приблизно на октаву вища. Різниця часу приходу у два вуха --- огляд & \cite{PalmerRussell1986,Grothe2010} & виміряно \\
16 & Інші датчики (табл.~\ref{tab:sensors}): у нюху $\sim400$ типів рецепторів, по одному на нейрон, і комбінаторний код; смак частково через АТФ і серотонін; дотик --- датчики, що звикають швидко й повільно; тепло й холод --- окремі & Миша: один рецептор упізнає багато запахів, один запах --- багато рецепторів, різні запахи --- різні набори. Без рецепторів АТФ нерв смаку не відповідає; смакові бруньки виділяють серотонін. Запис поодиноких волокон шкіри руки людини: чотири типи за швидкістю звикання. Холодовий канал відмінний від теплових & \cite{Mombaerts2004,Malnic1999,Finger2005,Huang2005,JohanssonVallbo1979,McKemy2002} & виміряно \\
\bottomrule
\end{longtable}}

\section{Розділ~\ref{sec:recognition}. Розпізнавання}

Швидкість розпізнавання, будову перших ступенів, лінійне читання відповіді й навчання на неперервності часу виміряно на людині й мавпі; зведення всього ланцюжка до двох операцій і навчання на відео перевірено на моделях книги, а пояснення ілюзій --- від добре обґрунтованих до спірних.

{\footnotesize
\setlength{\tabcolsep}{3pt}\renewcommand{\arraystretch}{1.15}
\begin{longtable}{>{\raggedright}p{0.5cm}>{\raggedright}p{4.0cm}>{\raggedright}p{5.6cm}>{\raggedright}p{2.3cm}>{\raggedright\arraybackslash}p{1.7cm}}
\toprule
№ & Твердження & Дослід і що виміряно & Джерело & Статус \\
\midrule\endfirsthead
\toprule
№ & Твердження & Дослід і що виміряно & Джерело & Статус \\
\midrule\endhead
1 & Порівняння зі зразком за пікселями при зсуві вгадує не краще за монетку; пара ступенів~\eqref{eq:scstage} упізнає фігуру в будь-якому місці & \texttt{models/\allowbreak recognition\_models.py}, «сітківка» з $24$ точок, $4000$ проб: зразок --- $0.49$, $0.51$, $0.52$ за частки перевернутих точок $0$, $0.1$, $0.2$; пара $S$ і $C$ --- $1.00$, $0.86$, $0.70$ & виведення в розділі & модель \\
2 & Тварину на картинці впізнають за $\sim150$\,мс, один прохід десятком ступенів по $10$--$20$\,мс & Людина, задача «є тварина чи ні» на фотографіях, показаних на $20$\,мс: викликані потенціали для двох відповідей розходяться через $\sim150$\,мс. Мавпа: час першої відповіді зростає від ступеня до ступеня (V1 раніше за всіх, далі V2 і V4). Кількість ступенів і «нуль або один імпульс на ступінь» --- висновок із цих чисел & \cite{Thorpe1996,Schmolesky1998} & виміряно \\
3 & Ступінь $S$ --- І за частинами, ступінь $C$ --- найбільше за положеннями (твердження~\ref{prop:conveyor}) & Кішка, первинна зорова кора: прості клітини відповідають на край певної орієнтації в певному місці, складні --- на ту саму орієнтацію в більшій ділянці. Внутрішньоклітинний запис складних клітин: відповідь на дві смужки в багатьох клітин близька до більшої з двох відповідей, у середньому найближча до операції max. Найбільше через взаємне гальмування --- твердження~\ref{prop:wta}, ділення на загальну активність --- огляд & \cite{HubelWiesel1962,Lampl2004,Carandini2012} & виміряно \\
4 & Вище ланцюжком поєднання більші й складніші, відповідь дедалі менше залежить від положення & Мавпа, спрощення стимулів до «критичної ознаки»: клітини V4 і задньої нижньоскроневої кори відповідають на складні поєднання форми, кольору й текстури за малого поля; у передній нижньоскроневій корі поле більше. Чергування $S$ і $C$ у моделі відтворює цю картину & \cite{KobatakeTanaka1994,Riesenhuber1999} & виміряно \\
5 & Згорткові мережі повторюють це чергування й найкраще передбачають відповіді верхніх ступенів; предмет читається лінійно за частотами групи нейронів & Мережа, навчена розпізнавати, без підгонки до мозку: верхній шар передбачає відповіді нейронів нижньоскроневої кори мавпи, середні --- відповіді V4. Лінійний класифікатор за активністю $\sim100$ випадково вибраних клітин у вікнах від $12.5$\,мс упізнає предмет і категорію за різних положення й розміру & \cite{Fukushima1980,Yamins2014,Hung2005} & виміряно \\
6 & Правило Гебба зі слідом~\eqref{eq:traceinvariance} навчає інваріантності на відео без підписів, якщо слід не коротший за $\sim20$\,мс & Ідея повільних ознак і правила зі слідом --- теорія. \texttt{models/\allowbreak recognition\_models.py}, два $C$-нейрони, $16$ входів, $10$ прогонів: пауза між предметами $0.3$\,с. При $\tau_{\text{сл}}=5$\,мс збіг із фігурою $0.52$ у середньому, при $10$\,мс --- $0.56$; при $20$, $50$ і $200$\,мс --- $1.00$ в усіх прогонах (при $30$\,мс один прогін із десяти помиляється) & \cite{Foldiak1991,WiskottSejnowski2002} & модель \\
7 & У мозку інваріантність навчається на неперервності часу & Мавпа: предмет непомітно підміняли під час стрибка ока в певне місце поля; позиційна інваріантність нейронів нижньоскроневої кори змінювалася відповідно до підміни, значуще вже за годину. Що це головне правило навчання зору --- гіпотеза & \cite{LiDiCarlo2008} & виміряно \\
8 & Рух: детектори напрямку, далі та сама пара І/АБО за великими ділянками & Детектори напрямку в сітківці кроля; в ділянці MT мавпи всі $110$ записаних нейронів вибіркові до напрямку, відповідь на рух сильніша, ніж у V1 & \cite{Barlow1965,Albright1984} & виміряно \\
9 & Верхні ступені надсилають униз передбачення, угору йде помилка & Модель пояснює позакласичні ефекти полів первинної зорової кори; окремі спостереження узгоджуються (огляд), як загальну будову конвеєра не доведено & \cite{RaoBallard1999,Keller2018} & гіпотеза \\
10 & Складні картинки потребують зворотних зв'язків і кількох кіл & Мавпа: для «складних» картинок рішення в нижньоскроневій корі складається на $\sim30$\,мс пізніше; ці пізні відповіді погано передбачає пряма мережа й краще --- мережі зі зворотними зв'язками або глибші & \cite{Kar2019} & виміряно \\
11 & Непомітна підробка пікселів обманює мережу; зір людини набагато стійкіший, але не невразливий & На мережах: мале спеціально дібране збурення змінює відповідь; приклад «панда $\to$ гібон» --- із другої праці. У людини за короткого показу підробки, що добре переносяться між мережами, зсувають вибір & \cite{Szegedy2014,Goodfellow2015,Elsayed2018} & виміряно \\
12 & Ілюзії очікування: ненадійний вхід поступається очікуванню --- бліде рухається «повільніше», сукню бачать по-різному (підрозділ~\ref{sec:illusions}) & Ґратки меншого контрасту здаються рухомими повільніше. Опитування $1401$ людини: сукню бачать у трьох стійких варіантах, підказка про освітлення перемикає колір; у $\sim13\,000$ спостерігачів вирішує припущення про освітлення. Баєсова модель з очікуванням повільних рухів пояснює низку ілюзій руху & \cite{Thompson1982,LaferSousa2015,Wallisch2017,Weiss2002,Purves2011} & гіпотеза \\
13 & Регулювання підсилення: водоспад, нахил і контраст; ґратка Германа --- не гальмування сусідів & Детектори напрямку сітківки кроля за тривалого руху знижують частоту, а після зупинки падають нижче за фон. Ілюзію нахилу відтворює модель із діленням на активність сусідів. Зміни ґратки, що не змінюють відповіді вихідних нейронів сітківки, прибирають плями --- пояснення сусідським гальмуванням відкинуто & \cite{BarlowHill1963,Schwartz2009,SchillerCarvey2005} & виміряно \\
14 & Компенсація затримок: рухомий предмет здається попереду спалаху & Ілюзію виміряно. Психофізика суперечить і продовженню руху вперед, і різниці затримок; запропоновано пояснення заднім числом --- за подіями $\sim80$\,мс після спалаху. Суперечку не розв'язано & \cite{Nijhawan1994,EaglemanSejnowski2000} & гіпотеза \\
15 & Нерухомі «змії» обертаються: різниця затримок~\eqref{eq:latency} читається детекторами напрямку & Ілюзія працює й без кольору; пари сусідніх елементів породжують її поодинці; нейрони кори мавпи, вибіркові до напрямку, дають спрямовану відповідь на ці нерухомі пари. У людини обертання запускають мікрострибки очей і кліпання & \cite{Conway2005,OteroMillan2012} & виміряно \\
16 & Двозначні картинки: взаємне гальмування, втома й шум; зміни викликає головно шум & Модель конкурентних груп нейронів: зміни в ній викликає \emph{шум}, без шуму їх немає; вона пояснює зростання частоти змін із силою стимулу. Втома тут відіграє меншу роль & \cite{MorenoBote2007} & гіпотеза \\
17 & Частина ілюзій --- наслідок задачі, якої вчиться машина & Мережа, навчена передбачати наступний кадр відео, «бачить» обертання нерухомих кілець і не бачить його на контрольних картинках; мережі для очищення й підвищення різкості зображень піддаються ілюзіям кольору та яскравості, як людина & \cite{Watanabe2018,GomezVilla2020} & виміряно \\
\bottomrule
\end{longtable}}

\section{Розділ~\ref{sec:muscles}. Вихід на м'язи}

Будову виходу --- рухові одиниці, запас надійності синапсу, увімкнення за розміром, пару м'язів, петлю розтягу й генератори ритму --- виміряно прямо; умова стійкості петлі із затримкою --- теорема; внутрішня модель тіла --- добре обґрунтована гіпотеза; числа рисунків --- розрахунок за моделлю книги.

{\footnotesize
\setlength{\tabcolsep}{3pt}\renewcommand{\arraystretch}{1.15}
\begin{longtable}{>{\raggedright}p{0.5cm}>{\raggedright}p{4.0cm}>{\raggedright}p{5.6cm}>{\raggedright}p{2.3cm}>{\raggedright\arraybackslash}p{1.7cm}}
\toprule
№ & Твердження & Дослід і що виміряно & Джерело & Статус \\
\midrule\endfirsthead
\toprule
№ & Твердження & Дослід і що виміряно & Джерело & Статус \\
\midrule\endhead
1 & Рухова одиниця: один \nt{ACET}-нейрон керує групою волокон, від кількох сотень до пари тисяч; у м'язах тонкого підлаштування одиниці дрібніші & М'язи людини, підрахунок рухових аксонів і м'язових волокон: у середньому близько $340$ волокон на нейрон у міжкістковому м'язі кисті, близько $410$ у плечопроменевому м'язі, близько $1930$ у медіальному литковому. Точного числа для найдрібніших одиниць розділ не наводить. & \cite{Feinstein1955mu} & виміряно \\
2 & Синапс на м'язі вивільняє в кілька разів більше медіатора, ніж потрібно для спрацювання волокна & Огляд вимірювань на нервово-м'язових синапсах: запас надійності (відношення вивільненого медіатора до порогового) у дорослих ссавців зазвичай $3$--$5$; у людини запас тримається більше на складках мембрани одержувача, ніж на величині вивільнення. & \cite{WoodSlater2001} & виміряно \\
3 & Посмикування~\eqref{eq:twitch} наростає за час $T_c$ порядку $50\ms$ & Поодинокі рухові одиниці кисті людини за довільного зусилля, силу усереднено за імпульсами одиниці: час наростання $30$--$100\ms$, у $80\,\%$ одиниць менше за $70\ms$. Функція $(t/T_c)e^{1-t/T_c}$ --- апроксимація з моделі пулу одиниць. & \cite{MilnerBrown1973rec, Fuglevand1993} & виміряно \\
4 & Сума посмикувань~\eqref{eq:forcesum}: за $5$, $15$ і $40$ імп/с сила $0.2$--$1.1F_1$, $1.8$--$2.2F_1$ і близько $5.4F_1$ (рис.~\ref{fig:tetanus}) & Розрахунок за \texttt{models/\allowbreak muscle\_\allowbreak models.py}, $T_c=50\ms$: $0.21$--$1.10$, $1.76$--$2.19$ і $5.28$--$5.49\,F_1$ на останніх $0.3\,\text{с}$. Лінійне додавання --- спрощення: насичення справжнього м'яза в моделі немає. & \texttt{models/\allowbreak muscle\_\allowbreak models.py} & модель \\
5 & Одиниці вмикаються за розміром; найслабші в сотню разів слабші за найсильніші & Вентральні корінці кішки: за наростаючого рефлекторного входу першими розряджаються нейрони з дрібними імпульсами в корінці, тобто малі. Міжкістковий м'яз кисті людини: сила посмикування одиниць від $0.1$ до $10$~г і зростає майже лінійно із силою, за якої одиниця вмикається. & \cite{Henneman1957, MilnerBrown1973rec} & виміряно \\
6 & Крок сили пропорційний силі, $\Delta F/F\approx q-1$~\eqref{eq:sizeprinciple}: рухома кома & Виведено в розділі з геометричної прогресії сил. Узгоджується з дослідом: за великих зусиль на той самий приріст сили вмикається набагато менше нових одиниць. & виведення в розділі; \cite{MilnerBrown1973rec} & теорія \\
7 & Розкид сили зростає пропорційно самій силі & Довільне ізометричне зусилля м'яза великого пальця: стандартне відхилення сили лінійно зростає із середньою силою; за того самого зусилля, викликаного електричною стимуляцією, розкид сталий. Модель пулу: лінійне зростання дає ввімкнення одиниць за порядком сили. & \cite{Jones2002sdn} & виміряно \\
8 & Плавність рухів --- наслідок шуму, що залежить від сигналу & Модель: траєкторія, що мінімізує розкид кінцевої точки за такого шуму, відтворює виміряні траєкторії ока й руки. Прямого досліду, який відрізняв би цю причину від інших, немає. & \cite{HarrisWolpert1998} & гіпотеза \\
9 & Різниця сил пари м'язів задає момент, сума --- жорсткість~\eqref{eq:antagonists} & Теорія керування імпедансом через спільне напруження. Рука людини: жорсткість кожного суглоба, виміряна малими поштовхами, приблизно пропорційна моменту в ньому; різні співвідношення спільного напруження змінюють форму й напрямок еліпса жорсткості кисті. & \cite{Hogan1984, GomiOsu1998stiff} & виміряно \\
10 & Датчик розтягу збуджує свій \nt{ACET}-нейрон і через гальмівного посередника вимикає антагоніста (рис.~\ref{fig:antagonists}) & Спинний мозок кішки: поодинокий посередник, що збуджується волокнами датчиків розтягу, дає моносинаптичні гальмівні потенціали в \nt{ACET}-нейронах м'яза-антагоніста, синаптична затримка $0.28$--$0.42\ms$. Медіатор посередника (гліцин) у цьому досліді не визначали. & \cite{Jankowska1972Ia} & виміряно \\
11 & Затримка петлі: спинномозкової $25$--$30\ms$, через кору в півтора--три рази більша, через око $100$--$200\ms$ & Рука людини, поштовх по суглобу: коротка відповідь м'язів через $20$--$45\ms$, довга $45$--$105\ms$, довільна $120$--$180\ms$. Зміщення видимого положення пальця під час руху: поправка через $160\ms$ у середньому. & \cite{Pruszynski2008rapid, SaundersKnill2003vis} & виміряно \\
12 & $\dd x/\dd t=-Gx(t-d)$ стійке при $Gd<\pi/2$~\eqref{eq:delaystable}; на межі частота $1/(4d)$ & Теорема про корені рівняння із запізненням. Перевірка за \texttt{models/\allowbreak muscle\_\allowbreak models.py}, $d=30\ms$: до $3\,\text{с}$ амплітуда $3\cdot10^{-6}$ при $G=40$, $0.13$ при $G=50$, $38$ при $G=55$ за секунду (межа $52$). & \cite{Hayes1950} & теорія \\
13 & Фізіологічне тремтіння $8$--$12$~Гц; петля розтягу --- одне з його джерел & Огляд вимірювань: дрібне тремтіння в діапазоні близько $10$~Гц є у здорових людей; його походження змішане --- центральні коливання, властивості розрядів одиниць, механічний резонанс і резонанс рефлекторної петлі, і за різних умов переважає різне. & \cite{McAuleyMarsden2000} & гіпотеза \\
14 & Мозок передбачає результат команди за внутрішньою моделлю тіла & Людина тримає вантаж щипком і рухає рукою: за інерційного, в'язкого й пружного навантаження сила стискання змінюється одночасно з навантаженням, а не із запізненням петлі, тобто навантаження передбачено. Огляд зводить такі досліди; прямого спостереження самої моделі в нейронах немає. & \cite{FlanaganWing1997grip, WolpertGhahramani2000} & гіпотеза \\
15 & Ритм ходьби, дихання, жування породжують місцеві мережі за сталого входу & Огляд дослідів: ізольована нервова система (спинний мозок, ганглії) видає ритмічний руховий візерунок без ритмічного входу й без зворотного зв'язку від м'язів. & \cite{MarderBucher2001} & виміряно \\
16 & Взаємне гальмування з утомою~\eqref{eq:halfcentre} дає ритм із періодом $0.84\,\text{с}\approx2\tau_a$ (рис.~\ref{fig:halfcentre}) & Теорема: мережа із взаємним гальмуванням і адаптацією без стійкого стану за сталого входу обов'язково коливається. Розрахунок за \texttt{models/\allowbreak muscle\_\allowbreak models.py} ($I=1$, $\beta=2$, $b=2$, $\tau=20\ms$, $\tau_a=0.4\,\text{с}$): період $0.84\,\text{с}$. Що справжні генератори влаштовані саме так, не показано. & \cite{Matsuoka1985osc} & модель \\
\bottomrule
\end{longtable}}

\section{Розділ~\ref{sec:pain}. Біль}

Датчики болю, два проводи, ворота в спинному мозку, низхідну ручку гучності, сенсибілізацію й захоплення уваги виміряно прямо; формули воріт і сенсибілізації --- спрощені моделі книги; правило, за яким машина вибирає гучність тривоги, --- гіпотеза.

{\footnotesize
\setlength{\tabcolsep}{3pt}\renewcommand{\arraystretch}{1.15}
\begin{longtable}{>{\raggedright}p{0.5cm}>{\raggedright}p{4.0cm}>{\raggedright}p{5.6cm}>{\raggedright}p{2.3cm}>{\raggedright\arraybackslash}p{1.7cm}}
\toprule
№ & Твердження & Дослід і що виміряно & Джерело & Статус \\
\midrule\endfirsthead
\toprule
№ & Твердження & Дослід і що виміряно & Джерело & Статус \\
\midrule\endhead
1 & Високий поріг: пороги нагрівання в різних датчиків болю від $41^\circ$ до $46$--$48^\circ$ & Запис поодиноких волокон. Шкіра мавпи: медіанний поріг повільних (C) датчиків $41^\circ$, швидких датчиків другого типу $46^\circ$, першого типу --- вище за $53^\circ$. Шкіра людини: C-датчики, що відповідають і на тиск, $40.7^\circ$, що не відповідають на тиск --- $48^\circ$. & \cite{Treede1995heat, Weidner1999cfib} & виміряно \\
2 & Датчики болю звикають набагато слабше за датчики дотику й після легкого ушкодження стають чутливішими; послаблення після сильного нагрівання виміряно в одного типу & Легкий опік шкіри ($50^\circ$, $100\,\text{с}$): через $5$--$10$~хв поріг болю в людей і поріг C-датчиків у мавпи знижено на $2$--$6^\circ$; в інших шкірних датчиків такого зсуву немає. У швидких датчиків другого типу відповідь після сильного нагрівання пригнічено. Порівняння звикання з датчиками дотику --- загальна оцінка, окремим дослідом тут не підкріплена. & \cite{LaMotte1982hyper, Treede1995heat} & виміряно \\
3 & Два проводи: швидкий $5$--$30$~м/с, повільний $0.5$--$2$~м/с; час приходу $t=L/v$~\eqref{eq:paintime} & Швидкість проведення: швидкі датчики болю мавпи в середньому $15$ і $25$~м/с (два типи); C-датчики людини $0.8$--$1.0$~м/с. При $L=1$~м це десятки мілісекунд і близько секунди. & \cite{Treede1995heat, Weidner1999cfib} & виміряно \\
4 & Біль приходить двічі: спершу швидкий і точний, потім тупий і довгий & Час реакції на нагрівання передпліччя людини: від $1100$ до $400\ms$ зі зростанням температури; сильне нагрівання волосистої шкіри дає подвійний біль, долоні --- ні. Перший біль можуть нести лише волокна, швидші за $6$~м/с, --- за латентністю йому відповідають швидкі датчики другого типу, яких у долоні немає. Розподіл ролей («де» і «наскільки погано») --- тлумачення. & \cite{CampbellLaMotte1983first, Treede1995heat} & виміряно \\
5 & Ворота: вихідний нейрон спинного мозку отримує біль і дотик, гальмівний посередник дотиком збуджується (рис.~\ref{fig:gate}) & Схему воріт запропоновано як теорію. Миші, генетична мітка й видалення груп нейронів: збуджувальні нейрони із соматостатином потрібні для механічного болю; гальмівні нейрони з динорфіном не дають входу від датчиків дотику дійти до них --- без цих нейронів легкий дотик викликає біль. & \cite{MelzackWall1965, Duan2014} & виміряно \\
6 & Дотик приглушує біль & Людина, лазерні імпульси болю на руці: одночасний дотик знижує оцінку болю й здатність розрізняти енергію імпульсів; ефект слабшає з відстанню між точкою дотику й точкою болю в межах одного шкірного сегмента. & \cite{Mancini2014touch} & виміряно \\
7 & Припущення теорії воріт: сильний біль сам вимикає гальмівного посередника, і ворота розчахуються & У розділі позначено як припущення вихідної теорії. Праця на мишах описує, як ворота не пускають дотик у канал болю; вимкнення посередника болем у ній не показано. & \cite{MelzackWall1965} & гіпотеза \\
8 & Ворота~\eqref{eq:gate}: поріг $0.18$ без дотику, $0.86$ з дотиком, $0.11$ при $g=2$; при $\nu=1$ дотик знижує вихід з $1$ до $0.25$, при $\nu=2$ не діє; дотик сам не проходить (рис.~\ref{fig:gatecurves}) & Розрахунок за \texttt{models/\allowbreak pain\_\allowbreak gate.py} з вагами з тексту відтворює всі ці числа. & \texttt{models/\allowbreak pain\_\allowbreak gate.py} & модель \\
9 & Низхідні шляхи зі стовбура змінюють підсилення воріт в обидва боки & Щур, довгастий мозок, запис під час нагрівання хвоста: одні нейрони розряджаються перед відсмикуванням («увімк»), інші замовкають («вимк»); слабка стимуляція цієї ділянки пригнічує рефлекс. Огляд: ті самі кола і полегшують, і гальмують передачу болю, опіоїди діють через них. & \cite{Fields1983rvm, Fields2004} & виміряно \\
10 & Очікування полегшення приглушує біль через опіоїдний канал & Люди з гострим болем: у тих, кому очікування полегшення зменшило біль, блокада опіоїдних рецепторів повернула біль до рівня решти; в інших блокада болю не змінювала. & \cite{Levine1978} & виміряно \\
11 & Гучність тривоги мозок вибирає за вигодою переривання & Досліду, в якому виміряно правило вибору гучності, немає. & --- & гіпотеза \\
12 & Сенсибілізація: після ушкодження вихід воріт зростає, поріг падає, почасти за рахунок спинного мозку; вона тримається й після припинення ушкоджувального стимулу & Щур: після ушкодження шкіри поріг згинального рефлексу падає, а відповідь зростає; електрофізіологічний аналіз показує, що частина зростання виникає в самому спинному мозку. Ушкоджувальний стимул викликає тривалі порушення чутливості, що переживають сам стимул. Точного часу спаду розділ не наводить. & \cite{Woolf1983} & виміряно \\
13 & Сенсибілізація~\eqref{eq:sensitization} --- додатний зворотний зв'язок; за сильної петлі тривога може защепнутися після загоєння & Випливає з моделі~\eqref{eq:sensitization} (задача наприкінці розділу). Досліду, який показав би, що стійкий біль після загоєння --- защепнута петля саме такого вигляду, немає. & виведення в розділі & гіпотеза \\
14 & Біль --- від'ємна винагорода, і навчання на ньому йде за помилкою часових різниць & Томограф, люди, ланцюжки провісників болю: активність вентрального стріатума й передньої острівцевої кори збігається із сигналами навчання за часовими різницями. Мавпа: частина \nt{DOPA}-нейронів збуджується від неприємного обдування та його провісника, інша частина гальмується. & \cite{Seymour2004, Matsumoto2009} & виміряно \\
15 & Біль перериває поточну роботу & Люди, електричний больовий стимул і звукові проби: відповідь на пробу через $100\ms$ після початку болю помітно сповільнена, через $1.5\,\text{с}$ різниці з контролем уже немає. Огляд зводить такі досліди в модель захоплення уваги. & \cite{Crombez1996disrupt, EcclestonCrombez1999} & виміряно \\
16 & Відсмикування замикається в спинному мозку & Щур: нейрони глибоких шарів заднього рогу повторюють просторовий візерунок рефлексу відсмикування окремих м'язів; їх не вдається збудити антидромно з шийного відділу, тобто це спінальні посередники. & \cite{Schouenborg1995wr} & виміряно \\
\bottomrule
\end{longtable}}

\section{Розділ~\ref{sec:motorlearning}. Навчання рухів}

Правило найменших квадратів і вигода окремого проводу помилки --- теореми; будову схеми з проводом помилки, послаблення за збігу, підлаштування сліду під затримку, післядію й спонтанне повернення навички виміряно; те, що вся схема працює як навчання з учителем і будує внутрішню модель, --- провідна гіпотеза; числа двопроцесної моделі --- розрахунок за моделлю книги.

{\footnotesize
\setlength{\tabcolsep}{3pt}\renewcommand{\arraystretch}{1.15}
\begin{longtable}{>{\raggedright}p{0.5cm}>{\raggedright}p{4.0cm}>{\raggedright}p{5.6cm}>{\raggedright}p{2.3cm}>{\raggedright\arraybackslash}p{1.7cm}}
\toprule
№ & Твердження & Дослід і що виміряно & Джерело & Статус \\
\midrule\endfirsthead
\toprule
№ & Твердження & Дослід і що виміряно & Джерело & Статус \\
\midrule\endhead
1 & Правило~\eqref{eq:lms} у середньому йде за градієнтом квадрата помилки & Результат теорії адаптивних фільтрів: поправка, пропорційна входу й помилці виходу, є стохастичним спуском за градієнтом середнього квадрата помилки. & \cite{WidrowHoff1960} & теорія \\
2 & З проводом помилки на кожен вихід швидкість не залежить від кількості виходів; з одним спільним числом --- падає з кількістю нейронів, що шумлять (твердження~\ref{prop:teachers}) & Криві навчання лінійної мережі: навчання за випадковими збуреннями нейронів повільніше за точний градієнт у стільки разів, скільки нейронів збурюють. & \cite{Werfel2005} & теорія \\
3 & Схема з проводом помилки працює як навчання з учителем (твердження~\ref{prop:errorwire}) & Теорії, виведені з будови схеми. Досліди рядків 7--10 з ними узгоджуються; що вся схема працює саме так, не доведено. & \cite{Marr1969, Albus1971} & гіпотеза \\
4 & Розширювач: близько $7\cdot10^{10}$ дрібних \nt{GLUT}-нейронів, більше, ніж у всьому іншому мозку & Підрахунок ядер у суспензії розчиненої тканини: у мозочку людини $69\cdot10^9$ нейронів із $86\cdot10^9$ в усьому мозку. Стереологія: $101\cdot10^9$ дрібних нейронів у мозочку літніх чоловіків. & \cite{Azevedo2009, Andersen1992cereb} & виміряно \\
5 & Підняття розмірності входу робить потрібну залежність лінійною & Теорема про кількість розбиттів: зважена сума $n$ входів із порогом реалізує довільну розмітку приблизно $2n$ векторів загального положення. Що мозочок користується саме цим, --- частина гіпотези рядка~3. & \cite{Cover1965} & теорія \\
6 & Близько $3\cdot10^7$ вихідних \nt{GABA}-нейронів, у кожного порядку $10^5$ входів & Стереологія мозочка людини: $30.5\cdot10^6$ вихідних нейронів. Електронна мікроскопія, щур: близько $1.75\cdot10^5$ входів від розширювача на один вихідний нейрон. Гальмівний медіатор вихідних нейронів --- в огляді. & \cite{Andersen1992cereb, NapperHarvey1988, Ito2001} & виміряно \\
7 & Рівно один провід помилки на вихідний нейрон, близько разу на секунду, щоразу потужний спалах & Огляд записів: кожен вихідний нейрон отримує один ліаноподібний \nt{GLUT}-вхід, що викликає складний імпульс із частотою порядку $1$~Гц. & \cite{Ito2001} & виміряно \\
8 & Імовірність спалаху залежить від напрямку помилки руху & Мавпа, окорухова частина мозочка, адаптація саккад: напрямок зорової помилки задає ймовірність складного імпульсу, величина --- його час; спалах змінює прості імпульси в наступній спробі й зсуває рух уздовж переважної помилки клітини. & \cite{Herzfeld2018} & виміряно \\
9 & Збіг спалаху помилки з активним входом послаблює вхід ($-\eta e_i x_j$) & Огляд дослідів: спільна стимуляція входу від розширювача й проводу помилки викликає довготривале послаблення цього входу; стимуляція одного входу --- ні. & \cite{Ito2001} & виміряно \\
10 & Довжину сліду підлаштовано під затримку помилки: за затримки $100\ms$ найсильніше послаблюється вхід, що був активним за $100\ms$ & Зрізи й живі миші: у частині, що відповідає за окорухове навчання, послаблення вузько налаштоване на інтервал близько $120\ms$ між входом і помилкою --- стільки йде сигнал помилки в цій задачі; в іншій частині різні клітини налаштовано на різні інтервали. & \cite{Suvrathan2016} & виміряно \\
11 & Схема --- адаптивний фільтр~\eqref{eq:adaptivefilter}, що вивчає внутрішню модель тіла & Модель, виведена з будови схеми. Прямого досліду, де вивчений фільтр виміряно як передавальну функцію тіла, немає. & \cite{Fujita1982} & гіпотеза \\
12 & Передбачення віднімає наслідки власних рухів, тому себе не полоскочеш & Люди: власний дотик здається менш лоскітним, ніж такий самий чужий; у томографі соматосенсорна кора відповідає на власний дотик слабше, а мозочок менш активний, коли рух торкається шкіри. Пояснення через віднімання передбачення --- гіпотеза. & \cite{Blakemore1998} & гіпотеза \\
13 & За зовнішньої сили промах зменшується, а після її зняття рука промахується в інший бік & Люди тягнуться рукояттю робота в силовому полі: траєкторії повертаються до прямих; після зняття поля післядія --- майже дзеркальне відображення перших викривлень; вона переноситься й у нетреновані частини простору. & \cite{ShadmehrMussaIvaldi1994} & виміряно \\
14 & Навчаються два процеси~\eqref{eq:tworate}; після зворотного збурення навичка повертається сама & Люди, силове поле, далі коротке зворотне поле й спроби із затиснутою помилкою: поправка сама повертається до старої; модель із двох процесів це відтворює, модель з одного --- ні. & \cite{Smith2006} & виміряно \\
15 & Числа рис.~\ref{fig:tworate}: $0.84$ (повільний $0.63$) за $200$ рухів; $6$ зворотних; швидкий $-0.53$, повільний $0.46$; повернення до $0.37$ за $40$ рухів & Розрахунок за \texttt{models/\allowbreak motor\_\allowbreak learning.py}, $A_f=0.92$, $B_f=0.1$, $A_s=0.996$, $B_s=0.02$: $0.840$ ($0.634$ і $0.205$), $6$ рухів, $-0.531$ і $0.457$, пік $0.371$ на $40$-му русі. & \texttt{models/\allowbreak motor\_\allowbreak learning.py} & модель \\
16 & Два процеси потрібні, бо тіло змінюється з різними швидкостями & Теорія: оптимальна оцінка за завад із кількома часами життя дає кілька фільтрів із різними сталими часу й пояснює спонтанне повернення та прискорене повторне навчання. & \cite{Kording2007} & гіпотеза \\
\bottomrule
\end{longtable}}

\section{Розділ~\ref{sec:chemoutput}. Хімічний вихід}

Два проводи до серця та їхню різну швидкість, від'ємний зворотний зв'язок у гормональних каскадах, код частотою порцій, добовий генератор і його підстроювання світлом виміряно прямо; межа $K<8$ --- теорема теорії автоматичного керування, виведена в розділі; пульсації, породжувані самим зворотним зв'язком каскаду, --- гіпотеза із сильним дослідом на її користь.

{\footnotesize
\setlength{\tabcolsep}{3pt}\renewcommand{\arraystretch}{1.15}
\begin{longtable}{>{\raggedright}p{0.5cm}>{\raggedright}p{4.0cm}>{\raggedright}p{5.6cm}>{\raggedright}p{2.3cm}>{\raggedright\arraybackslash}p{1.7cm}}
\toprule
№ & Твердження & Дослід і що виміряно & Джерело & Статус \\
\midrule\endfirsthead
\toprule
№ & Твердження & Дослід і що виміряно & Джерело & Статус \\
\midrule\endhead
1 & Ритмоводій серця сам б'ється близько $100$ разів на хвилину; у спокої гальмо натиснуте, і частота зазвичай порядку $60$--$70$ & Люди, повне вимкнення обох проводів до серця: власна частота вища за частоту спокою і спадає з віком, у дорослих порядку $100$ ударів на хвилину. Отже, у спокої переважає гальмо. Частота спокою $60$--$70$ --- типове значення, окремо не цитоване. & \cite{JoseCollison1970ihr} & виміряно \\
2 & Газ \nt{NORA} і гальмо \nt{ACET} --- фільтри нижніх частот, гальмо швидше~\eqref{eq:gasbrake} & Собака, частотно-модульована стимуляція блукаючого або симпатичного нерва серця й передатна функція до частоти передсердь: обидва шляхи --- фільтри нижніх частот, у симпатичного частота зрізу значно нижча й додано чисту затримку близько $1.7\,\text{с}$. Характеристики залежать від середнього тонусу, тож лінійна формула~\eqref{eq:gasbrake} --- наближення. & \cite{Berger1989} & виміряно \\
3 & Гальмо швидке, бо калієвий канал \nt{ACET} відкриває без другого посередника, а \nt{NORA} --- через ланцюжок реакцій & Мембранні клаптики клітин передсердь: калієвий канал, що відкривається ацетилхоліном, відкривають $\beta\gamma$-субодиниці G-білка, пов'язаного з рецептором, без другого посередника всередині клітини. Шлях \nt{NORA} через цАМФ тут не цитовано. & \cite{Logothetis1987gbg} & виміряно \\
4 & Кров обходить тіло за час порядку хвилини; \nt{ADRE} з кровотоку доходить до всіх органів із рецептором & Загальна оцінка порядку величини; у розділі так і сформульовано, джерело не цитовано. & --- & оцінка \\
5 & Гормональний каскад охоплено від'ємним зворотним зв'язком: кінцевий гормон гальмує перші ступені & Огляд дослідів: гормони кори надниркових залоз пригнічують викид АКТГ гіпофізом і рилізинг-гормону гіпоталамусом у трьох часових діапазонах --- швидкому, проміжному й повільному. & \cite{KellerWood1984fb} & виміряно \\
6 & Три однакові ступені стійкі за $K<8$~\eqref{eq:cascadestable}, на межі період $2\pi\tau/\sqrt3\approx3.6\tau$ & Виведено в розділі: на частоті $\sqrt3/\tau$ три ступені дають зсув фази $180^\circ$ і ослаблення у $8$ разів. & виведення в розділі & теорія \\
7 & Похибка рівня $1/(1+K)$, тому точніше приблизно за $10\,\%$ такий регулятор не тримає (твердження~\ref{prop:cascade}) & Наслідок рядка~6 для пропорційного зворотного зв'язку. Справжня вісь має зворотний зв'язок на кількох ступенях і в кількох часових діапазонах (рядок~5), тож межа стосується моделі~\eqref{eq:cascade}, а не виміряна. & виведення в розділі & теорія \\
8 & За $K=4$ і $7$ рівень усталюється, за $K=12$ --- рівні пульсації з періодом $3.7$--$3.9\,\tau$ (рис.~\ref{fig:cascade}) & Розрахунок за \texttt{models/\allowbreak chem\_\allowbreak models.py}: за $K=12$ послідовні періоди $3.9$, $3.7$, $3.8$, $3.7\,\tau$; за $K=4$ розмах наприкінці рахунку $0.003$. & \texttt{models/\allowbreak chem\_\allowbreak models.py} & модель \\
9 & Гормон стресу виходить пульсами приблизно раз на годину; пульси породжує сам зворотний зв'язок, без окремого генератора & Щури, постійне вливання рилізинг-гормону: АКТГ і кортикостерон коливаються з фізіологічною, майже годинною частотою --- ритмічний вхід від гіпоталамуса для пульсів не потрібен. Модель гіпофіз--наднирник із запізнілим зворотним зв'язком дає такі коливання. & \cite{Walker2012glc, Walker2010} & гіпотеза \\
10 & Отримувач читає частоту порцій, а не рівень & Мавпи без власного викиду рилізинг-гормону гонадотропінів: постійне вливання не відновлює викиду гормонів гіпофіза, порції раз на годину відновлюють; перехід на постійне вливання знову глушить відповідь. Утома рецептора як фільтр верхніх частот --- тлумачення. & \cite{Belchetz1978} & виміряно \\
11 & Різні частоти порцій по-різному діють на різні клітини однієї залози & Ті самі мавпи: почастішання порцій до $2$--$5$ на годину знижує обидва гормони гіпофіза; порідшання до однієї за $3$~години знижує один (ЛГ) і підвищує другий (ФСГ), тож їхнє відношення задається частотою. & \cite{Wildt1981gnrh} & виміряно \\
12 & Добовий ритм породжує внутрішньоклітинний від'ємний зворотний зв'язок із затримкою в кілька годин & Огляд: білки годинникових генів із затримкою пригнічують транскрипцію власних генів; петля з транскрипції й трансляції дає період близько доби. & \cite{TakahashiJS2017} & виміряно \\
13 & Головний генератор --- група з десятків тисяч нейронів, що синхронізує решту тіла & Огляд: супрахіазматичне ядро --- головний добовий генератор; окремі його нейрони в культурі --- самостійні генератори, мережа їх синхронізує; у миші близько $10^4$ нейронів із кожного боку. & \cite{Welsh2010scn} & виміряно \\
14 & Власний період у людини $24.18$~години & Люди в тьмяному світлі за розкладом, відірваним від доби: період ритмів мелатоніну, температури й кортизолу в середньому $24.18$~години в молодих і літніх, із вузьким розкидом. & \cite{Czeisler1999} & виміряно \\
15 & Світло підводить генератор як фазове автопідстроювання~\eqref{eq:pll}; потрібен зсув близько $11$~хв на добу & Люди, одиночне яскраве світло $6.7$~години в різний час доби: крива фазової відповіді першого типу з розмахом $5$~годин --- затримка до мінімуму температури тіла, випередження після. Рівняння~\eqref{eq:pll} --- стандартна модель; $11$~хв $=0.18$~години --- арифметика. & \cite{Khalsa2003prc} & виміряно \\
16 & Без світла підстроювання немає, і ритм переходить на власний період & Цілком незрячі люди без сприйняття світла, що живуть у звичайному суспільстві: приблизно в половини ритми мелатоніну й кортизолу йдуть зі своїм періодом, що не збігається з добою. & \cite{Sack1992blind} & виміряно \\
\bottomrule
\end{longtable}}

\section{Розділ~\ref{sec:debugging}. Налагодження машини}

Що чує електрод, маловимірність активності, роботу інтерфейсів на жорстких ґратках, спільне навчання мозку й декодера та зорові точки від стимуляції виміряно в рецензованих дослідах; оцінки потоків даних --- арифметика розділу; про вживлений людям пристрій Neuralink, наскільки вдалося перевірити, публікувала дані переважно сама компанія.

{\footnotesize
\setlength{\tabcolsep}{3pt}\renewcommand{\arraystretch}{1.15}
\begin{longtable}{>{\raggedright}p{0.5cm}>{\raggedright}p{4.0cm}>{\raggedright}p{5.6cm}>{\raggedright}p{2.3cm}>{\raggedright\arraybackslash}p{1.7cm}}
\toprule
№ & Твердження & Дослід і що виміряно & Джерело & Статус \\
\midrule\endfirsthead
\toprule
№ & Твердження & Дослід і що виміряно & Джерело & Статус \\
\midrule\endhead
1 & Електрод чує імпульси нейронів у радіусі порядку сотні мікрометрів, зазвичай від одного до кількох; їх розрізняють за формою & Щур, поле CA1, одночасний внутрішньоклітинний і позаклітинний запис: форма позаклітинного імпульсу повторює ширину й амплітуду внутрішньоклітинного. Оцінка: тетрод міг би розрізняти $60$--$100$ нейронів, а насправді виділяють близько шести. & \cite{Henze2000extra} & виміряно \\
2 & У мозку $8.6\cdot10^{10}$ нейронів; щуп чує приблизно одну стомільйонну & Підрахунок ядер у суспензії розчиненої тканини: $86\cdot10^9$ нейронів у мозку людини. Частка --- арифметика розділу. & \cite{Azevedo2009} & виміряно \\
3 & Активність сотень нейронів рухової кори вкладається в один-два десятки спільних змінних & Огляд записів у руховій корі мавп: активність популяції описується небагатьма прихованими змінними, спільними для багатьох нейронів. & \cite{Gallego2017} & виміряно \\
4 & Шум сусідів узгоджений, і сотня нейронів несе майже стільки ж, скільки тисяча (твердження~\ref{prop:pooling}) & Зорова кора мавпи: коефіцієнт кореляції шуму сусідніх нейронів близько $0.1$, і виграш від усереднення насичується на сотні нейронів. Теорія кореляцій, що обмежують інформацію. & \cite{Zohary1994, MorenoBote2014} & виміряно \\
5 & Сирий потік $2\cdot10^8$~біт/с проти $8\cdot10^4$~біт/с в імпульсах~\eqref{eq:probedata} & Арифметика розділу за ємністю потоку імпульсів~\eqref{eq:spikeentropy}. Параметри оцифрування реальні: кристал Neuralink --- $19.3$~кГц, $10$~біт на канал. & виведення в розділі; \cite{Rieke1997, Musk2019} & теорія \\
6 & Перша система Neuralink: кристали підсилюють і оцифровують, сирий потік іде кабелем, імпульси виділяє програма ззовні; виділення на кристалі, закритий корпус, радіо й бездротове живлення --- у пристрою для людей, за повідомленнями компанії & Стаття компанії: повний сирий потік кабелем USB-C, імпульси виділяє програма в реальному часі поза кристалом; стиснення на борту, бездротові живлення й передавання названо планом для клінічних пристроїв. Для вживленого людям пристрою --- лише повідомлення компанії. Що виділяти імпульси на кристалі необхідно, --- висновок розділу з~\eqref{eq:probedata}. & \cite{Musk2019}; звіти Neuralink, рецензованої статті немає & виміряно (стаття компанії); для людей --- звіт компанії \\
7 & До $3072$ електродів на $96$ нитках по $32$; нитки вводить робот, оминаючи судини & Стаття компанії: $3072$ електроди на $96$ нитках по $32$; робот уводить $6$ ниток ($192$ електроди) за хвилину; до $70\,\%$ електродів у щурів із довготривало вживленою ґраткою чують імпульси. & \cite{Musk2019} & виміряно (стаття компанії) \\
8 & У людей близько тисячі каналів на $64$ нитках; частина ниток змістилася; курсор порядку $10$~біт/с & Лише повідомлення компанії. Пошук у PubMed не знайшов рецензованої статті з результатами на людях; це узгоджується із застереженням у тексті розділу. & звіти Neuralink; рецензованої статті немає & виміряно (звіт компанії) \\
9 & Інтерфейс на ґратці із сотні електродів керує курсором & Людина з паралічем, $96$ електродів у первинній руховій корі: намір рухати рукою змінює імпульси через три роки після травми; декодер дає «нейронний курсор» (пошта, телевізор), керування протезом кисті й роботом-рукою. & \cite{Hochberg2006} & виміряно \\
10 & Письмо «уявною рукою»: близько $90$ знаків на хвилину, менше ніж $8$~біт/с & Людина з паралічем кисті, спроби писати від руки, рекурентний декодер: $90$ знаків на хвилину з точністю $94.1\,\%$ у реальному часі, понад $99\,\%$ після автокорекції. Оцінка в бітах --- арифметика розділу. & \cite{Willett2021} & виміряно \\
11 & Спроби говорити перетворюються на текст зі швидкістю близько $60$ слів на хвилину & Людина, що втратила розбірливе мовлення, ґратки в корі: $62$ слова на хвилину; $9.1\,\%$ помилок у словах за словника з $50$ слів, $23.8\,\%$ за словника зі $125\,000$. & \cite{Willett2023} & виміряно \\
12 & Інтерфейс витягає малу частку ємності: в одного нейрона десятки біт на секунду, у сотні --- тисячі (твердження~\ref{prop:probe}) & Верхня межа~\eqref{eq:spikeentropy} --- теорія; порівняння з рядками~8, 10 --- висновок розділу. Що вузьке місце --- маловимірність і незнання коду, а не провід, прямим дослідом не перевірено. & \cite{Rieke1997} & теорія \\
13 & Мозок і декодер учаться один в одного & Мавпи керують курсором; декодер підлаштовується в замкненій петлі, а нейрони водночас перенавчаються: точність зростає, навичка зберігається й менше страждає від рухів власної руки. Порада розвести швидкості навчання --- інженерне правило, окремого досліду в розділі немає. & \cite{Orsborn2014coad} & виміряно \\
14 & Струм через електрод збуджує навколишні нейрони й проводи без розбору типу & Кора миші й кішки, двофотонний запис кальцію: навіть слабкий струм збуджує рідкісні нейрони, розкидані на відстань до міліметрів, через пряме збудження аксонів в об'ємі діаметром у десятки мікрометрів. Тобто збудження невибіркове, але не суцільне. & \cite{Histed2009stim} & виміряно \\
15 & Стимуляція зорової кори запалює точки; до $16$ точок одночасно дають змогу впізнати деякі літери й межі предметів & Незряча людина, $96$ електродів у зоровій корі на $6$ місяців: поріг одного електрода $66.8\pm36.5$~мкА; одночасна стимуляція груп (до $16$ електродів --- межа стимулятора) знижує поріг і дає розрізненні візерунки, деякі літери й межі предметів упізнаються. & \cite{Fernandez2021} & виміряно \\
16 & Другий напрям Neuralink --- пристрій, що подає сигнали в зорову кору & Лише повідомлення компанії про розроблення; рецензованих даних про його роботу не знайдено. & звіти Neuralink & гіпотеза \\
17 & Концентрації широкомовних медіаторів можна вимірювати хімічним датчиком у тканині & Зрізи мозку щура, швидка циклічна вольтамперометрія на вуглецевому волокні: виміряно викид і зворотне захоплення серотоніну; речовина виходить за межі синапсу. & \cite{Bunin1998} & виміряно \\
\bottomrule
\end{longtable}}

\section{Розділ~\ref{sec:eetypes}. Нейрони як прилади}

Режими роботи окремих клітин виміряно прямо, струмом через електрод і записом напруги; математика інтегратора й поділу на класи --- класична теорія, а те, що мозок користується переналаштуванням режимів для обчислень, лишається гіпотезою.

{\footnotesize
\setlength{\tabcolsep}{3pt}\renewcommand{\arraystretch}{1.15}
\begin{longtable}{>{\raggedright}p{0.5cm}>{\raggedright}p{4.0cm}>{\raggedright}p{5.6cm}>{\raggedright}p{2.3cm}>{\raggedright\arraybackslash}p{1.7cm}}
\toprule
№ & Твердження & Дослід і що виміряно & Джерело & Статус \\
\midrule\endfirsthead
\toprule
№ & Твердження & Дослід і що виміряно & Джерело & Статус \\
\midrule\endhead
1 & Резонатор: повільний струм, що протидіє зміні напруги, робить нейрон смуговим фільтром із максимумом на одиницях--десятках герц (підрозділ~\ref{sec:resonator}) & У зрізах у \nt{GLUT}-нейрони вводили синусоїдальний струм із плавно зростаючою частотою і вимірювали імпеданс $|Z(f)|$: максимум на $2$--$5$~Гц за $33^\circ$C (близько $7$~Гц за $38^\circ$C); резонанс створюють повільні калієвий і катіонний струми, а підсилює постійний натрієвий струм. Огляд дає той самий прийом для багатьох класів клітин & \cite{HuVervaekeStorm2002,HutcheonYarom2000} & виміряно \\
2 & Повільний струм поводиться як індуктивність, разом з ємністю мембрани --- коливальний контур & Підпорогову відповідь аксона на струм виміряно й описано лінеаризованими рівняннями провідностей; еквівалентна схема містить «феноменологічну» індуктивність, величина якої залежить від напруги & \cite{Mauro1970} & теорія \\
3 & Стала часу групи зі зворотним зв'язком $\tau_{\text{еф}}=\tau/(1-w)$~\eqref{eq:perfectint}; за $\tau=0.1$~с і $w=0.995$ --- $20$~с & Виводиться в розділі з лінійного рівняння групи; як механізм утримання положення ока запропоновано в теоретичній роботі & виведення в розділі; \cite{Seung1996} & теорія \\
4 & Інтегратор положення ока тримає вхід мережею, а не властивістю окремої клітини & Внутрішньоклітинний запис у риби у вузлі, що тримає положення ока: після швидкого повороту частота нейрона сходинкою змінюється й тримається; короткий струм в одну клітину спричиняє лише короткочасний зсув, а сходинки напруги зберігаються й нижче порога --- їх підтримують синаптичні входи & \cite{Aksay2001} & виміряно \\
5 & Інтеграторові без витоку потрібне постійне підлаштування навчанням; коли його розладнано, він або забуває, або йде в насичення & Рибу навчали зоровим зворотним зв'язком, пропорційним $\pm$ положенню ока: утримання ставало нестійким або з витоком, стала часу нейронів падала більш ніж на порядок, до $<1$~с; звичайний зоровий зворотний зв'язок поступово повертав стійкість & \cite{Major2004} & виміряно \\
6 & Бістабільний нейрон: короткий вхід умикає тривале плато, гальмування вимикає; властивість умикається серотоніном (підрозділ~\ref{sec:bistable}) & Внутрішньоклітинний запис \nt{ACET}-нейронів, що керують м'язами, у кішки: коротке синаптичне збудження переводить нейрон у тривалу роботу, коротке гальмування скидає; у спінальної кішки стан з'являється після введення попередника серотоніну & \cite{Hounsgaard1988} & виміряно \\
7 & Два класи: інтегратори (частота зростає від нуля) і резонатори (на порозі одразу скінченна частота) & Класи виведено з теорії біфуркацій. Дослід: у зрізах кори \nt{GLUT}-нейрони починають працювати з як завгодно малою частотою (тип~1), швидкі \nt{GABA}-нейрони --- стрибком зі скінченної частоти (тип~2) & \cite{Izhikevich2007,Tateno2004} & виміряно \\
8 & Один нейрон --- інтегратор або детектор збігів: гальмування вкорочує вікно додавання & У зрізах гіпокампа вікно, у якому збуджувальні входи додаються до порога, коротше за $2$~мс; його вкорочує пряме гальмування, що надходить одразу за збудженням; у дендритах, де гальмування слабше, вікно ширше & \cite{Pouille2001} & виміряно \\
9 & Лінія затримки з аксонів (таблиця~\ref{tab:eetypes}) & У сипухи виміряно затримки в аксонах, що йдуть до детекторів збігів: різна довжина аксона дає різну затримку, і детектори налаштовані на різницю часів надходження звуку & \cite{Carr1990} & виміряно \\
10 & Ритмоводій під час неспання й мовчазна клітина уві сні & \nt{SERO}-нейрони працюють майже регулярно вдень, сповільнюються в повільному сні й майже мовчать у швидкому (докладніше --- розділ~\ref{sec:sleep}) & \cite{JacobsAzmitia1992,McGintyHarper1976} & виміряно \\
11 & Прилад задають іонні канали й широкомовні канали, що їх підлаштовують (твердження~\ref{prop:eetypes}) & Показано на малих мережах: речовини-регулятори змінюють власні властивості нейронів і силу синапсів, тож та сама схема видає різні ритми & \cite{Marder2012} & виміряно \\
12 & Мозок використовує переналаштування режимів для обчислень (твердження~\ref{prop:eetypes}) & Прямого досліду, де переналаштування режиму клітини було б потрібне для розв'язання задачі, немає; є лише досліди, де регулятори змінюють вихід схеми & \cite{Marder2012} & гіпотеза \\
\bottomrule
\end{longtable}}

\section{Розділ~\ref{sec:dendrites}. Кількість дендритів}

Будову класів і кількість входів виміряно анатомічно й електричними записами; оцінка~\eqref{eq:dendriteload} --- проста фізика конденсатора, а обчислювальне пояснення кількості входів спирається на теорію й моделі.

{\footnotesize
\setlength{\tabcolsep}{3pt}\renewcommand{\arraystretch}{1.15}
\begin{longtable}{>{\raggedright}p{0.5cm}>{\raggedright}p{4.0cm}>{\raggedright}p{5.6cm}>{\raggedright}p{2.3cm}>{\raggedright\arraybackslash}p{1.7cm}}
\toprule
№ & Твердження & Дослід і що виміряно & Джерело & Статус \\
\midrule\endfirsthead
\toprule
№ & Твердження & Дослід і що виміряно & Джерело & Статус \\
\midrule\endhead
1 & Питома ємність мембрани $c_m\approx1$~мкФ/см$^2$ & Із тіла нейрона знімали мембрану на піпетку, подавали сходинку напруги й за спадом ємнісного струму знаходили повну ємність, потім ділили на площу: $0.9$~мкФ/см$^2$ у трьох класів нейронів & \cite{Gentet2000} & виміряно \\
2 & Час заряджання $t\approx c_mA\Delta V/I$~\eqref{eq:dendriteload}: великому нейронові потрібні десятки одночасних входів, маленькому вистачає одного великого синапсу & Оцінка за законом заряджання конденсатора; числа ($20$ і $300$~пФ, $0.1$ і кілька нА) --- порядки величин & виведення в розділі & теорія \\
3 & Один величезний синапс дає струм у кілька наноамперів і відразу доводить маленький нейрон до порога & Одночасний запис закінчення й отримувача у зрізі: струм одного синапсу $2$--$13$~нА, кожен імпульс входу викликає надпорогову відповідь, затримка близько $0.4$~мс за $36^\circ$C; отримувач видає \nt{GLYC} & \cite{Borst1995,BorstSoria2012} & виміряно \\
4 & Нуль дендритів: у чутливих нейронів дотику й болю імпульс іде від датчика в спинний мозок повз тіло клітини & Будову (аксон, що ділиться надвоє біля тіла) встановлено анатомічно. Те, що проведення не потребує збудливості тіла, показано на перевіреній моделі такого нейрона: без збудливого тіла імпульс проходить розгалуження так само надійно & \cite{AmirDevor2003} & теорія \\
5 & Один дендрит: нюховий нейрон несе один тип рецептора й передає одну вхідну лінію & Огляд дослідів із генами рецепторів: кожен нюховий нейрон експресує один ген рецептора з великої родини & \cite{Mombaerts2004} & виміряно \\
6 & Два дендрити: у детекторів напрямку на звук входи від лівого й правого вуха сідають на різні дендрити & Анатомія й записи в ссавців, зібрані в огляді: входи від двох вух рознесено по двох протилежних дендритах & \cite{Grothe2010} & виміряно \\
7 & Рознесення входів по двох дендритах робить «І» суворішим & Показано на моделі: насичення~\eqref{eq:shunt} на одному дендриті не дає входам з одного боку дійти до порога, і детекція збігів поліпшується. Прямого досліду, де змінювали б розкладку входів, немає & \cite{AgmonSnir1998} & теорія \\
8 & Близько $7\cdot10^{10}$ маленьких \nt{GLUT}-нейронів кола навчання рухів, по $\sim4$ входи в кожного & Підрахунок клітин ізотропним фракціонуванням: у цій частині мозку людини близько $69$ млрд нейронів. Запис однієї такої клітини в живого щура: на дотик надходять пачки дискретних струмів від небагатьох входів, і клітина спрацьовує, коли вони додаються & \cite{Azevedo2009,Chadderton2004} & виміряно \\
9 & Пороговий елемент з $n$ входами розділяє не більше $P_{\max}\approx2n$ випадкових образів~\eqref{eq:cover}; для вихідного нейрона кола навчання рухів це межа $\sim2\cdot10^5$, а реалістична оцінка --- $\sim4\cdot10^4$ відповідностей & Межа --- класичний результат комбінаторної геометрії. Реалістична оцінка --- теорія ємності лише з додатними вагами й запасом на шум, застосована до виміряної кількості входів цього нейрона & \cite{Cover1965,Brunel2004} & теорія \\
10 & Змішувачі з малою кількістю входів потрібні, щоб розширити вхід; «чотири» близько до оптимуму & Теорія: вимірність подання, яку дає шар змішувачів, максимальна приблизно за тієї кількості входів, що спостерігається анатомічно; вихідна гіпотеза розширення --- у класичних працях & \cite{Marr1969,Albus1971,LitwinKumar2017} & гіпотеза \\
11 & Великі дерева: $10^4$--$10^5$ входів & Підрахунок синапсів на електронних знімках: $\approx1.7\cdot10^5$ синапсів на одному вихідному \nt{GABA}-нейроні кола навчання рухів у щура; близько $30\,000$ збуджувальних і $1\,700$ гальмівних входів на \nt{GLUT}-нейроні гіпокампа & \cite{NapperHarvey1988,Megias2001} & виміряно \\
12 & Гілки великого дерева --- окремі підсхеми з власними порогами; нейрон --- маленька багатошарова мережа & Точкова стимуляція двох місць на тонких дендритах: входи на одній гілці додаються за сигмоїдою, на різних гілках --- лінійно. У дендритах нейронів кори людини знайдено градуальні кальцієві імпульси, що дають змогу одній клітині розв'язувати лінійно нероздільну задачу. Модель: щоб відтворити один нейрон, потрібна глибока мережа & \cite{Polsky2004,Gidon2020,Beniaguev2021} & виміряно \\
13 & Дендрити-виходи: кожна гілка \nt{GABA}-нейрона сітківки --- окремий детектор напрямку & Двофотонний запис кальцію в окремих гілках: сигнал у гілці вибірковий до руху від тіла клітини до кінчика, а напруга тіла --- ні & \cite{Euler2002} & виміряно \\
14 & Кількість входів іде за задачею (твердження~\ref{prop:dendrites}) & Будову класів виміряно (рядки 3--13); що кількість входів обрано заради швидкості чи класифікації, показано лише теорією й моделями для окремих класів & \cite{LitwinKumar2017,AgmonSnir1998} & гіпотеза \\
\bottomrule
\end{longtable}}

\section{Розділ~\ref{sec:sleep}. Сон}

Режими сну, повтори й зменшення синапсів виміряно записами нейронів, мікродіалізом і електронною мікроскопією; розклад сну й повільні гойдалки описано моделями книги, а призначення сну --- здебільшого гіпотези.

{\footnotesize
\setlength{\tabcolsep}{3pt}\renewcommand{\arraystretch}{1.15}
\begin{longtable}{>{\raggedright}p{0.5cm}>{\raggedright}p{4.0cm}>{\raggedright}p{5.6cm}>{\raggedright}p{2.3cm}>{\raggedright\arraybackslash}p{1.7cm}}
\toprule
№ & Твердження & Дослід і що виміряно & Джерело & Статус \\
\midrule\endfirsthead
\toprule
№ & Твердження & Дослід і що виміряно & Джерело & Статус \\
\midrule\endhead
1 & Уві сні нейрони кори не мовчать; змінюється режим роботи & Тривалі записи нейронів кори в щурів: у повільному сні частота лишається ненульовою, періоди роботи чергуються з періодами тиші; після тривалого неспання частота вища в усіх станах, після сну --- нижча & \cite{Vyazovskiy2009} & виміряно \\
2 & \nt{ACET} високий удень і у швидкому сні, низький у повільному (таблиця~\ref{tab:sleepmodes}) & Мікродіаліз у корі вільно рухливих кішок: під час неспання викид ацетилхоліну на $100$--$175\%$ вищий, ніж у повільному сні, а у швидкому сні --- приблизно на рівні спокійного неспання & \cite{Marrosu1995} & виміряно \\
3 & \nt{NORA} і \nt{SERO} затихають у повільному сні й майже мовчать у швидкому & Запис окремих \nt{NORA}-нейронів у щурів і \nt{SERO}-нейронів у кішок за стадіями сну: частота максимальна під час неспання, нижча в повільному сні й майже нульова у швидкому & \cite{AstonJonesBloom1981,McGintyHarper1976} & виміряно \\
4 & У швидкому сні м'язи вимкнено гальмуванням через \nt{GABA} і \nt{GLYC} & У щурів вимірювали активність \nt{ACET}-нейронів жувального м'яза й подавали блокатори просто до них: параліч знімається, лише якщо заблоковано і \nt{GABA}-рецептори обох типів, і \nt{GLYC}-рецептори & \cite{BrooksPeever2012} & виміряно \\
5 & Тиск сну $S$ --- конденсатор із двома порогами~\eqref{eq:sleeppressure}, $\tau_r=18.2$~год, $\tau_d=4.2$~год & Модель двох процесів; сталі часу отримано з того, як потужність повільних хвиль ЕЕГ зростає з неспанням і падає уві сні, а форму порогів --- з часу спонтанного пробудження & \cite{Borbely1982,Daan1984} & теорія \\
6 & Машина сама виходить на $16.1$~год неспання і $8.0$~год сну (рис.~\ref{fig:twoprocess}) & Розрахунок за~\eqref{eq:sleeppressure} з гойдальними порогами, скрипт \texttt{models/sleep\_models.py}: $16.1$~год неспання і $7.9$--$8.0$~год сну & --- & модель \\
7 & Після $40$~год без сну $S=0.90$, але сон триває лише $8.8$~год: борг повертається глибиною, а не тривалістю & Модель: \texttt{models/sleep\_models.py} дає $S=0.90$ і $8.8$~год. Дослід: після $40.5$~год без сну у восьми людей у першу відновлювальну ніч потужність повільних хвиль ЕЕГ значуще вища за звичайну & \cite{Borbely1981} & модель \\
8 & Фізично $S$ --- аденозин & Мікродіаліз у кішок: позаклітинний аденозин в одній із ділянок, що керують неспанням, зростає за час неспання, зростає далі під час позбавлення сну й падає під час сну. Що $S$ --- лише аденозин, не показано & \cite{PorkkaHeiskanen1997,Dunwiddie2001} & гіпотеза \\
9 & У повільному сні кора гойдається: робота кілька сотень мс --- тиша, рідше ніж раз на секунду ($<1$~Гц, під наркозом частіше $0.3$--$0.4$~Гц) & Внутрішньоклітинний запис у кішок під наркозом: деполяризація $0.8$--$1.5$~с і тривала гіперполяризація, ритм $<1$~Гц, найчастіше $0.3$--$0.4$~Гц; те саме чергування роботи й тиші видно в природному сні (рядок~1) & \cite{Steriade1993,Vyazovskiy2009} & виміряно \\
10 & Гойдалки --- група зі зворотним зв'язком і втомою~\eqref{eq:updown}; за $b=5$ період $1.1$~с (значення моделі) і робота близько $35\%$ часу, за $b=1$ робота рівна (рис.~\ref{fig:updown}) & Розрахунок, скрипт \texttt{models/sleep\_models.py}: період $1.11$~с, частка часу роботи $0.35$ (середнє за цілим числом періодів); за $b=1$ частка роботи $1.0$. Період моделі коротший за виміряний (рядок~9) & --- & модель \\
11 & \nt{ACET} вимикає гойдалки й переводить кору в рівну роботу & Стимуляція ядра \nt{ACET}-нейронів у кішки блокувала повільний ритм у $79\%$ нейронів кори й переводила їх у рівну роботу, переважно через зникнення тривалих гіперполяризацій. Ацетилхолін пригнічує повільні калієві струми, що створюють утому & \cite{SteriadeAmzica1993,McCormick1992} & виміряно \\
12 & Спалахи ритму $7$--$15$~Гц породжує петля \nt{GLUT}--\nt{GABA} між корою й релейним вузлом & У зрізах релейного вузла зору в тхора спалахи створюються відповідними пачками релейних клітин на гальмування від сусіднього \nt{GABA}-ядра, яке вони самі збуджують & \cite{vonKrosigk1993,SteriadeMcCormickSejnowski1993} & виміряно \\
13 & Повтори дня йдуть із перемотуванням: у гіпокампі приблизно в $20$ разів швидше, у корі --- в $6$--$7$ разів & У повільному сні повторювані послідовності нейронів гіпокампа йдуть стисненими в часі. Після бігу доріжкою послідовності нейронів місця повторювалися в тому самому порядку короткими спалахами по $\sim100$~мс, стисненими приблизно в $20$ разів; у корі стиснення в $6$--$7$ разів & \cite{Nadasdy1999,LeeWilson2002,Euston2007} & виміряно \\
14 & Посилення узгодженості повторів із корою поліпшує пам'ять (причиновий зв'язок) & У щурів після навчання стимуляцією, прив'язаною до повторів, посилювали їхній зв'язок із повільними хвилями й спалахами ритму в корі: наступного дня пам'ять була високою, у контрольних --- на рівні випадку & \cite{Maingret2016} & виміряно \\
15 & Призначення повторів --- перемішане навчання повільної пам'яті & Теорія двох систем навчання й розрахунки на штучних мережах: послідовне навчання псує старе, перемішане --- ні. Прямого досліду на мозку, що повтори потрібні саме для цього, немає & \cite{McClelland1995} & гіпотеза \\
16 & Після сну синапси зменшуються пропорційно розміру, великі майже не змінюються & Тривимірна електронна мікроскопія $6920$ синапсів кори миші: площа контакту після сну на $\sim18\%$ менша, ніж після неспання, зменшення пропорційне розміру, великі стійкі синапси не змінюються & \cite{deVivo2017} & виміряно \\
17 & Головне завдання сну --- перенормування ваг & Гіпотеза синаптичного гомеостазу; рядки 1 і 16 з нею узгоджуються, але не доводять, що це головна функція & \cite{TononiCirelli2014} & гіпотеза \\
18 & Швидкий сон --- випробування моделі світу на стенді & Режим (таблиця~\ref{tab:sleepmodes}) виміряно; що мережа проганяє модель світу, --- теоретичне припущення, досліду немає & \cite{Hobson2009} & гіпотеза \\
19 & Промивання мозку уві сні: питання відкрите & Один дослід: уві сні міжклітинний простір на $60\%$ ширший і очищення швидше. Інший, з іншим способом вимірювання: уві сні й під наркозом очищення помітно повільніше. Результати суперечать один одному & \cite{Xie2013,Miao2024} & гіпотеза \\
20 & Місцевий сон: ділянки кори тварини, що не спить, ненадовго провалюються в тишу & У щурів після тривалого неспання нейрони окремої ділянки кори коротко замовкають, як у повільному сні, з повільною хвилею в місцевій ЕЕГ; у такі моменти щур частіше помиляється в задачі & \cite{Vyazovskiy2011} & виміряно \\
\bottomrule
\end{longtable}}

\section{Розділ~\ref{sec:livingmachine}. Машина з живих нейронів}

Поведінку культур на ґратках електродів виміряно в прямих дослідах із записом і стимуляцією; механізм залпів спирається на модель виснаження синапсів і досліди на мікрокультурах, а твердження про дофамін у чашці --- на поодинокі досліди, частину яких самі автори вважають лише демонстрацією.

{\footnotesize
\setlength{\tabcolsep}{3pt}\renewcommand{\arraystretch}{1.15}
\begin{longtable}{>{\raggedright}p{0.5cm}>{\raggedright}p{4.0cm}>{\raggedright}p{5.6cm}>{\raggedright}p{2.3cm}>{\raggedright\arraybackslash}p{1.7cm}}
\toprule
№ & Твердження & Дослід і що виміряно & Джерело & Статус \\
\midrule\endfirsthead
\toprule
№ & Твердження & Дослід і що виміряно & Джерело & Статус \\
\midrule\endhead
1 & Культура на чипі: здебільшого \nt{GLUT}, менша частка \nt{GABA}, що залежить від препарату; у культурах гіпокампа близько $6\%$ & Мережа з тисяч нейронів на ґратці електродів --- стандартний дослід (рядок~3). Частку \nt{GABA}-нейронів у культурах гіпокампа рахували за фарбуванням на GABA: близько $6\%$ нейронів. Для культур кори перевіреного числа ми не наводимо & \cite{Benson1994} & виміряно \\
2 & Органоїди: тривимірна нервова тканина розміром близько міліметра зі стовбурових клітин людини & Зі стовбурових клітин виростили тривимірні органоїди з кількома ділянками мозку, зокрема корою з шарами попередників і зрілими нейронами & \cite{Lancaster2013} & виміряно \\
3 & Без входу культура спалахує залпами з інтервалами від секунди до хвилин, залежно від культури & $58$ культур кори щільністю від $3\,000$ до $50\,000$ нейронів записували п'ять тижнів ($963$ півгодинні записи): після двох тижнів майже в усіх культурах переважають залпи всієї мережі; інтервали між залпами --- від $1$ до $300$~с і сильно залежать від препарату & \cite{Wagenaar2006} & виміряно \\
4 & Залп закінчується виснаженням медіатора, інтервал $T_{\text{залп}}\approx\tau_{\text{відн}}\ln\frac{1}{1-x^*}$~\eqref{eq:burstinterval}; $2$--$4$~с --- оцінка моделі з $\tau_{\text{відн}}=0.8$~с, виміряне відновлення повільніше & Формулу виведено в розділі. Модель: мережа з синапсами, що виснажуються, сама породжує залпи всієї мережі. Дослід на мікрокультурах із $4$--$30$ нейронів: тривалість залпу залежить від часу після попереднього, виснаження пухирців медіатора скасовує залпи; висновок авторів --- залп обмежений запасом медіатора. Відновлення запасу там --- десятки секунд, що з тією самою формулою дає інтервали в десятки секунд і хвилини & висновок у розділі; \cite{Tsodyks2000,CohenSegal2011,Tsodyks1997} & теорія \\
5 & «Учитель, який замовкає»: стимуляцію вимикають, коли з'явилася потрібна відповідь, і відповідь закріплюється & Культуру стимулювали з частотою $0.3$--$1$~Гц, доки через $50\pm10$~мс після стимулу не з'явиться обрана відповідь, потім стимуляцію вимикали на $5$~хв; після кількох циклів потрібна відповідь викликалася одразу; будували криві навчання & \cite{Shahaf2001} & виміряно \\
6 & Культура грає в теніс; значуще зростання серій за сеанс --- тільки за повного зворотного зв'язку & Докладно --- у розділі~\ref{sec:dishbrain} і додатку до нього: коли замість стимулу тиша, культура наприкінці сеансу теж краща, ніж без зворотного зв'язку, але з меншим ефектом; навчання від сеансу до сеансу нестійке & \cite{Kagan2022} & виміряно \\
7 & Те, що культура вчиться передбачати свій вхід, --- гіпотеза; гучні слова про «розумність» оскаржено & Принцип вільної енергії --- теоретична пропозиція; прямого досліду, що відрізняв би його від простіших пояснень, немає. Тридцять дослідників у статті-відповіді вказали, що зміна поведінки в петлі не показує ні розуму, ні відчуттів & \cite{Friston2010,Balci2023} & гіпотеза \\
8 & Культура веде віртуальне тіло до цілі за дібраного малюнка стимулів & Програма сама добирала навчальні малюнки стимуляції за поточним результатом; за десятки хвилин мережа досягала заданих станів, пластичність трималася $80$~хв після $2$~год навчання. Ті самі стимули, подані без зв'язку з поведінкою, не діяли & \cite{Bakkum2008} & виміряно \\
9 & Резервуар: навчається тільки лінійне зчитування $y=W_{\text{вих}}\mathbf r$~\eqref{eq:reservoir} & Теорія обчислення на резервуарі: будь-яка нелінійна система зі згасною пам'яттю плюс навчений лінійний вихід & \cite{Maass2002,JaegerHaas2004} & теорія \\
10 & Органоїд як резервуар розрізняє мовців із точністю близько $78\%$ (за даними авторів); за помилкою навчається зчитування, а зв'язки органоїду змінюються без учителя & Звуки мовлення кількох мовців подавали органоїду малюнками струму на ґратці електродів; навчене зчитування розрізняло мовця з точністю близько $78\%$ --- так повідомляють автори. Вони ж повідомляють, що під час тренування змінювалася функціональна зв'язність органоїду, і називають це навчанням без учителя в самому органоїді & \cite{Cai2023} & виміряно \\
11 & Мільйон нейронів витрачає на імпульси близько $10^{-4}$~Вт~\eqref{eq:dishpower} & Оцінка: ціна імпульсу $10^{-10}$~Дж з енергетичного бюджету кори~\eqref{eq:energy}, помножена на кількість імпульсів; прямого вимірювання потужності культури немає & висновок у розділі; \cite{Attwell2001} & теорія \\
12 & Дофамін за спалахом світла: $1$~мМ дофаміну в клітці, $240$ повторів, кількість викликаних імпульсів зростала & На віддаленій платформі з органоїдами дофамін у світлочутливій клітці ($1$~мМ) вивільняли ультрафіолетом ($365$~нм, $0.8$~с), коли стимул викликав більше імпульсів; за $240$ кроків (близько $5$~год) кількість імпульсів загалом зростала. Автори пишуть, що за одним дослідом зв'язок із дофаміном встановити не можна; контролів немає & \cite{Jordan2024} & гіпотеза \\
13 & Дофамін від живих клітин змінює вагу органічного транзистора надовго й оборотно & Клітини, що виділяють дофамін, вирощено над органічним нейроморфним елементом; окиснення дофаміну біля електрода змінює провідність каналу; показано довготривалу зміну ваги та її повернення & \cite{Keene2020} & виміряно \\
14 & Органоїди з робочими \nt{DOPA}-нейронами існують; як учителя їхній дофамін не використовували & В органоїдах зі стовбурових клітин людини знайдено електрично активні зрілі \nt{DOPA}-нейрони й вироблення дофаміну. Дослідів, де цей дофамін навчає іншу мережу, ми в літературі не знайшли & \cite{Jo2016} & виміряно \\
15 & Органоїд балансує жердиною: навчальні стимули від програми кращі за випадкові, без закріплення, через \nt{GLUT}-синапси & Органоїди кори миші в замкненій петлі із задачею «жердина на візку»: у більшості органоїдів сигнали, обрані навчанням із підкріпленням, дали кращий результат, ніж випадкові або жодні; після $45$~хв відпочинку покращення не зберігалося; блокада AMPA- і NMDA-рецепторів його скасовувала & \cite{Robbins2026} & виміряно \\
16 & Без регістрів керування мережа на стенді не може сама вирішити, що вважати успіхом (твердження~\ref{prop:livingmachine}) & В усіх дослідах рядків 5--15 навчальний сигнал задає експериментатор або програма; досліду, де культура сама виробила б критерій успіху, немає --- це висновок із відсутності, а не дослід & --- & гіпотеза \\
\bottomrule
\end{longtable}}

\section{Розділ~\ref{sec:dishbrain}. DishBrain}

Будову стенда й результати гри взято з однієї експериментальної роботи та її продовження; формули шуму і дрейфу до добрих налаштувань виведено в розділі й перевірено моделлю книги, а механізм навчання в чашці не встановлено.

{\footnotesize
\setlength{\tabcolsep}{3pt}\renewcommand{\arraystretch}{1.15}
\begin{longtable}{>{\raggedright}p{0.5cm}>{\raggedright}p{4.0cm}>{\raggedright}p{5.6cm}>{\raggedright}p{2.3cm}>{\raggedright\arraybackslash}p{1.7cm}}
\toprule
№ & Твердження & Дослід і що виміряно & Джерело & Статус \\
\midrule\endfirsthead
\toprule
№ & Твердження & Дослід і що виміряно & Джерело & Статус \\
\midrule\endhead
1 & Стенд: близько $800\,000$ клітин кори миші або близько $10^6$ клітин людини на чипі; $26\,000$ електродів на $8$~мм$^2$, запис до $1024$, стимуляція через $8$ & Методи роботи: чип MaxOne, $26\,000$ платинових електродів на $8$~мм$^2$, запис $1024$ каналів по $20\,000$ відліків на секунду; стимулювати технічно можна до $32$ електродів, незалежно вдалося $8$; мишачі культури використовували приблизно з $10$~діб & \cite{Kagan2022} & виміряно \\
2 & Петля з тактом $10$~мс: імпульси рухових ділянок рухають ракетку (рис.~\ref{fig:dishloop}) & Імпульси підраховують вікнами по $10$~мс; затримка від імпульсу до відповідного стимулу близько $5$~мс, її задає буфер апаратури & \cite{Kagan2022} & виміряно \\
3 & Вхід: код місця (який із $8$ електродів) і частота $4$--$40$~імп/с & В основному досліді частота стимуляції зростає лінійно від $4$~Гц, коли м'яч біля дальньої стінки, до $40$~Гц біля ракетки; амплітуда імпульсів м'яча --- $75$~мВ; імпульси двофазні прямокутні & \cite{Kagan2022} & виміряно \\
4 & Вихід $\Delta p=c\,(n_\uparrow-n_\downarrow)$~\eqref{eq:dishreadout}, шум лічби за вікно $1/\sqrt{RT}$ & Правило різниці двох ділянок описано в роботі (з вагами ділянок за активністю), точної формули не дано. Оцінка шуму --- наслідок пуассонівської лічби, виведена в розділі & \cite{Kagan2022}; висновок у розділі & теорія \\
5 & Учитель: після удару --- передбачуваний стимул $75$~мВ, $100$~імп/с, $0.1$~с; після промаху --- непередбачуваний, $150$~мВ, $5$~імп/с, $4$~с у випадкові місця й моменти, потім $4$~с паузи & Методи: після удару $75$~мВ, $100$~Гц, $100$~мс на всі $8$ електродів одразу; після промаху $150$~мВ, $5$~Гц у випадкові електроди у випадкові моменти впродовж $4$~с, потім $4$~с без стимуляції. Дофаміну в культурі немає & \cite{Kagan2022} & виміряно \\
6 & Мережа діє так, щоб її вхід став передбачуваним & Принцип вільної енергії --- теоретична пропозиція, з якої автори виводили протокол; дослід із нею узгоджується, але не відрізняє її від простіших механізмів (рядки~7--8) & \cite{Friston2010,Kagan2022} & гіпотеза \\
7 & Якщо невдача струшує мережу, а вдача ні, час у налаштуванні $\rho\propto1/P_{\text{пр}}$~\eqref{eq:dishdensity} (твердження~\ref{prop:dishscramble}) & Випливає з балансу потоків у марковському ланцюгу із симетричними поштовхами, виведено в розділі. Прямої перевірки на культурі немає & висновок у розділі & теорія \\
8 & Модель «перемішування при промаху» навчається: частка ударів $0.21\to0.38$; за поштовху після кожного розіграшу $\approx0.12$, без поштовхів $0.19$ (рис.~\ref{fig:dishmodel}) & Розрахунок, скрипт \texttt{models/dishbrain\_model.py}, $40$ моделей-культур по $2000$ розіграшів: $0.21\to0.38$; $0.13\to0.11$; $0.19\to0.19$ & --- & модель \\
9 & Серії відбитих м'ячів подовжуються тільки в живих культур у повній петлі; контролі не навчаються & $399$ сеансів по $20$~хв: $9$ культур миші ($101$ сеанс), $11$ людини ($138$), $6$ чипів із середовищем ($80$), $20$ культур у спокої ($42$), $3$ симуляції ($38$). Середня довжина розіграшу за останні $15$~хв більша, ніж за перші $5$, тільки в живих культур; у клітин, що не є нейронами (HEK293T), і в чипів із середовищем ефекту немає & \cite{Kagan2022} & виміряно \\
10 & Значуще зростання за сеанс --- тільки за повного зворотного зв'язку; коли замість стимулу тиша, наприкінці сеансу гра краща, ніж без зворотного зв'язку, але ефект менший & $486$ сеансів за $3$ дні: «стимул» ($n=27$), «тиша» ($17$), «без зворотного зв'язку» ($15$). Зростання в часі значуще тільки в умові «стимул»; наприкінці сеансу умова «тиша» перевершувала «без зворотного зв'язку» з меншим ефектом & \cite{Kagan2022} & виміряно \\
11 & Ефект невеликий; чи навчається мережа надовго --- відкрите питання & Усередині сеансу покращення надійне, але, за словами самих авторів, стійкого навчання між сеансами впродовж кількох днів не спостерігали & \cite{Kagan2022} & виміряно \\
12 & Під час гри мережа наближається до критичного режиму, і що ближче, то краща гра & Аналіз лавин активності тих самих культур: структурований вхід наближає мережу до критичного режиму, якість гри корелює з близькістю до нього; але самої лише критичності без відомостей про наслідки дій для навчання замало & \cite{Habibollahi2023} & виміряно \\
13 & «Розумність» культури не показано & Стаття-відповідь тридцяти дослідників: зміна поведінки в замкненій петлі не доводить ні розуму, ні відчуттів; досліду, який би це показав, немає & \cite{Balci2023} & гіпотеза \\
14 & Запропонований четвертий контроль: передбачуваний стимул після промаху тієї самої тривалості й енергії (підрозділ~\ref{sec:dishspec}) & Такого досліду не поставлено; це пропозиція книги & --- & гіпотеза \\
\bottomrule
\end{longtable}}


\begin{thebibliography}{99}
\bibitem{Azevedo2009} F. A. C. Azevedo et al., \emph{Equal numbers of neuronal and nonneuronal cells make the human brain an isometrically scaled-up primate brain}, J. Comp. Neurol. \textbf{513}, 532 (2009).
\bibitem{Schultz1997} W. Schultz, P. Dayan, and P. R. Montague, \emph{A neural substrate of prediction and reward}, Science \textbf{275}, 1593 (1997).
\bibitem{SuttonBarto2018} R. S. Sutton and A. G. Barto, \emph{Reinforcement Learning: An Introduction}, 2nd ed. (MIT Press, Cambridge, MA, 2018).
\bibitem{Doya2002} K. Doya, \emph{Metalearning and neuromodulation}, Neural Networks \textbf{15}, 495 (2002).
\bibitem{Strata1999} P. Strata and R. Harvey, \emph{Dale's principle}, Brain Res. Bull. \textbf{50}, 349 (1999).
\bibitem{vanVreeswijk1996} C. van Vreeswijk and H. Sompolinsky, \emph{Chaos in neuronal networks with balanced excitatory and inhibitory activity}, Science \textbf{274}, 1724 (1996).
\bibitem{AstonJones2005} G. Aston-Jones and J. D. Cohen, \emph{An integrative theory of locus coeruleus-norepinephrine function: adaptive gain and optimal performance}, Annu. Rev. Neurosci. \textbf{28}, 403 (2005).
\bibitem{Beaulieu2011} J.-M. Beaulieu and R. R. Gainetdinov, \emph{The physiology, signaling, and pharmacology of dopamine receptors}, Pharmacol. Rev. \textbf{63}, 182 (2011).
\bibitem{Sparso2001} J. Spars{\o} and S. Furber (eds.), \emph{Principles of Asynchronous Circuit Design: A Systems Perspective} (Kluwer, Boston, 2001).
\bibitem{Carr1990} C. E. Carr and M. Konishi, \emph{A circuit for detection of interaural time differences in the brain stem of the barn owl}, J. Neurosci. \textbf{10}, 3227 (1990).
\bibitem{Wang2001} X.-J. Wang, \emph{Synaptic reverberation underlying mnemonic persistent activity}, Trends Neurosci. \textbf{24}, 455 (2001).
\bibitem{Hahnloser2002} R. H. R. Hahnloser, A. A. Kozhevnikov, and M. S. Fee, \emph{An ultra-sparse code underlies the generation of neural sequences in a songbird}, Nature \textbf{419}, 65 (2002).
\bibitem{Barlow1965} H. B. Barlow and W. R. Levick, \emph{The mechanism of directionally selective units in rabbit's retina}, J. Physiol. \textbf{178}, 477 (1965).
\bibitem{Carandini2012} M. Carandini and D. J. Heeger, \emph{Normalization as a canonical neural computation}, Nat. Rev. Neurosci. \textbf{13}, 51 (2012).
\bibitem{Pouille2001} F. Pouille and M. Scanziani, \emph{Enforcement of temporal fidelity in pyramidal cells by somatic feed-forward inhibition}, Science \textbf{293}, 1159 (2001).
\bibitem{Tsodyks1997isn} M. V. Tsodyks, W. E. Skaggs, T. J. Sejnowski, and B. L. McNaughton, \emph{Paradoxical effects of external modulation of inhibitory interneurons}, J. Neurosci. \textbf{17}, 4382 (1997).
\bibitem{Buzsaki2012} G. Buzs\'aki and X.-J. Wang, \emph{Mechanisms of gamma oscillations}, Annu. Rev. Neurosci. \textbf{35}, 203 (2012).
\bibitem{Rieke1997} F. Rieke, D. Warland, R. de Ruyter van Steveninck, and W. Bialek, \emph{Spikes: Exploring the Neural Code} (MIT Press, Cambridge, MA, 1997).
\bibitem{Mainen1995} Z. F. Mainen and T. J. Sejnowski, \emph{Reliability of spike timing in neocortical neurons}, Science \textbf{268}, 1503 (1995).
\bibitem{Softky1993} W. R. Softky and C. Koch, \emph{The highly irregular firing of cortical cells is inconsistent with temporal integration of random EPSPs}, J. Neurosci. \textbf{13}, 334 (1993).
\bibitem{Thorpe1996} S. Thorpe, D. Fize, and C. Marlot, \emph{Speed of processing in the human visual system}, Nature \textbf{381}, 520 (1996).
\bibitem{Zohary1994} E. Zohary, M. N. Shadlen, and W. T. Newsome, \emph{Correlated neuronal discharge rate and its implications for psychophysical performance}, Nature \textbf{370}, 140 (1994).
\bibitem{MorenoBote2014} R. Moreno-Bote et al., \emph{Information-limiting correlations}, Nat. Neurosci. \textbf{17}, 1410 (2014).
\bibitem{Gollisch2008} T. Gollisch and M. Meister, \emph{Rapid neural coding in the retina with relative spike latencies}, Science \textbf{319}, 1108 (2008).
\bibitem{OKeefe1993} J. O'Keefe and M. L. Recce, \emph{Phase relationship between hippocampal place units and the EEG theta rhythm}, Hippocampus \textbf{3}, 317 (1993).
\bibitem{Destexhe2003} A. Destexhe, M. Rudolph, and D. Par\'e, \emph{The high-conductance state of neocortical neurons in vivo}, Nat. Rev. Neurosci. \textbf{4}, 739 (2003).
\bibitem{Tsodyks1997} M. V. Tsodyks and H. Markram, \emph{The neural code between neocortical pyramidal neurons depends on neurotransmitter release probability}, Proc. Natl. Acad. Sci. USA \textbf{94}, 719 (1997).
\bibitem{Lisman1997} J. E. Lisman, \emph{Bursts as a unit of neural information: making unreliable synapses reliable}, Trends Neurosci. \textbf{20}, 38 (1997).
\bibitem{Attwell2001} D. Attwell and S. B. Laughlin, \emph{An energy budget for signaling in the grey matter of the brain}, J. Cereb. Blood Flow Metab. \textbf{21}, 1133 (2001).
\bibitem{Lennie2003} P. Lennie, \emph{The cost of cortical computation}, Curr. Biol. \textbf{13}, 493 (2003).
\bibitem{Yagishita2014} S. Yagishita et al., \emph{A critical time window for dopamine actions on the structural plasticity of dendritic spines}, Science \textbf{345}, 1616 (2014).
\bibitem{Werfel2005} J. Werfel, X. Xie, and H. S. Seung, \emph{Learning curves for stochastic gradient descent in linear feedforward networks}, Neural Comput. \textbf{17}, 2699 (2005).
\bibitem{Doya2000} K. Doya, \emph{Reinforcement learning in continuous time and space}, Neural Comput. \textbf{12}, 219 (2000).
\bibitem{Oja1982} E. Oja, \emph{Simplified neuron model as a principal component analyzer}, J. Math. Biol. \textbf{15}, 267 (1982).
\bibitem{BiPoo1998} G.-Q. Bi and M.-M. Poo, \emph{Synaptic modifications in cultured hippocampal neurons: dependence on spike timing, synaptic strength, and postsynaptic cell type}, J. Neurosci. \textbf{18}, 10464 (1998).
\bibitem{Williams1992} R. J. Williams, \emph{Simple statistical gradient-following algorithms for connectionist reinforcement learning}, Mach. Learn. \textbf{8}, 229 (1992).
\bibitem{Bayer2005} H. M. Bayer and P. W. Glimcher, \emph{Midbrain dopamine neurons encode a quantitative reward prediction error signal}, Neuron \textbf{47}, 129 (2005).
\bibitem{Shen2008} W. Shen, M. Flajolet, P. Greengard, and D. J. Surmeier, \emph{Dichotomous dopaminergic control of striatal synaptic plasticity}, Science \textbf{321}, 848 (2008).
\bibitem{Dabney2020} W. Dabney, Z. Kurth-Nelson, N. Uchida, C. K. Starkweather, D. Hassabis, R. Munos, and M. Botvinick, \emph{A distributional code for value in dopamine-based reinforcement learning}, Nature \textbf{577}, 671 (2020).
\bibitem{Matsumoto2009} M. Matsumoto and O. Hikosaka, \emph{Two types of dopamine neuron distinctly convey positive and negative motivational signals}, Nature \textbf{459}, 837 (2009).
\bibitem{BrombergMartin2010} E. S. Bromberg-Martin, M. Matsumoto, and O. Hikosaka, \emph{Dopamine in motivational control: rewarding, aversive, and alerting}, Neuron \textbf{68}, 815 (2010).
\bibitem{Menegas2018} W. Menegas, K. Akiti, R. Amo, N. Uchida, and M. Watabe-Uchida, \emph{Dopamine neurons projecting to the posterior striatum reinforce avoidance of threatening stimuli}, Nat. Neurosci. \textbf{21}, 1421 (2018).
\bibitem{Niv2007} Y. Niv, N. D. Daw, D. Joel, and P. Dayan, \emph{Tonic dopamine: opportunity costs and the control of response vigor}, Psychopharmacology \textbf{191}, 507 (2007).
\bibitem{Mohebi2019} A. Mohebi et al., \emph{Dissociable dopamine dynamics for learning and motivation}, Nature \textbf{570}, 65 (2019).
\bibitem{Hamid2016} A. A. Hamid et al., \emph{Mesolimbic dopamine signals the value of work}, Nat. Neurosci. \textbf{19}, 117 (2016).
\bibitem{Kim2020} H. R. Kim, A. N. Malik et al., \emph{A unified framework for dopamine signals across timescales}, Cell \textbf{183}, 1600 (2020).
\bibitem{Jeong2022} H. Jeong et al., \emph{Mesolimbic dopamine release conveys causal associations}, Science \textbf{378}, eabq6740 (2022).
\bibitem{Amo2022} R. Amo et al., \emph{A gradual temporal shift of dopamine responses mirrors the progression of temporal difference error in machine learning}, Nat. Neurosci. \textbf{25}, 1082 (2022).
\bibitem{Takahashi2017} Y. K. Takahashi et al., \emph{Dopamine neurons respond to errors in the prediction of sensory features of expected rewards}, Neuron \textbf{95}, 1395 (2017).
\bibitem{Sharpe2017} M. J. Sharpe et al., \emph{Dopamine transients are sufficient and necessary for acquisition of model-based associations}, Nat. Neurosci. \textbf{20}, 735 (2017).
\bibitem{Gardner2018} M. P. H. Gardner, G. Schoenbaum, and S. J. Gershman, \emph{Rethinking dopamine as generalized prediction error}, Proc. R. Soc. B \textbf{285}, 20181645 (2018).
\bibitem{Engelhard2019} B. Engelhard et al., \emph{Specialized coding of sensory, motor and cognitive variables in VTA dopamine neurons}, Nature \textbf{570}, 509 (2019).
\bibitem{Clements1992} J. D. Clements, R. A. J. Lester, G. Tong, C. E. Jahr, and G. L. Westbrook, \emph{The time course of glutamate in the synaptic cleft}, Science \textbf{258}, 1498 (1992).
\bibitem{Hasselmo2006} M. E. Hasselmo, \emph{The role of acetylcholine in learning and memory}, Curr. Opin. Neurobiol. \textbf{16}, 710 (2006).
\bibitem{JacobsAzmitia1992} B. L. Jacobs and E. C. Azmitia, \emph{Structure and function of the brain serotonin system}, Physiol. Rev. \textbf{72}, 165 (1992).
\bibitem{Schiller1992} P. H. Schiller, \emph{The ON and OFF channels of the visual system}, Trends Neurosci. \textbf{15}, 86 (1992).
\bibitem{Grothe2010} B. Grothe, M. Pecka, and D. McAlpine, \emph{Mechanisms of sound localization in mammals}, Physiol. Rev. \textbf{90}, 983 (2010).
\bibitem{Markram2004} H. Markram et al., \emph{Interneurons of the neocortical inhibitory system}, Nat. Rev. Neurosci. \textbf{5}, 793 (2004).
\bibitem{Joshi2016} S. Joshi, Y. Li, R. M. Kalwani, and J. I. Gold, \emph{Relationships between pupil diameter and neuronal activity in the locus coeruleus, colliculi, and cingulate cortex}, Neuron \textbf{89}, 221 (2016).
\bibitem{ServanSchreiber1990} D. Servan-Schreiber, H. Printz, and J. D. Cohen, \emph{A network model of catecholamine effects: gain, signal-to-noise ratio, and behavior}, Science \textbf{249}, 892 (1990).
\bibitem{YuDayan2005} A. J. Yu and P. Dayan, \emph{Uncertainty, neuromodulation, and attention}, Neuron \textbf{46}, 681 (2005).
\bibitem{Nassar2012} M. R. Nassar et al., \emph{Rational regulation of learning dynamics by pupil-linked arousal systems}, Nat. Neurosci. \textbf{15}, 1040 (2012).
\bibitem{BouretSara2005} S. Bouret and S. J. Sara, \emph{Network reset: a simplified overarching theory of locus coeruleus noradrenaline function}, Trends Neurosci. \textbf{28}, 574 (2005).
\bibitem{Hebb1949} D. O. Hebb, \emph{The Organization of Behavior: A Neuropsychological Theory} (Wiley, New York, 1949).
\bibitem{MacKay1952} D. M. MacKay and W. S. McCulloch, \emph{The limiting information capacity of a neuronal link}, Bull. Math. Biophys. \textbf{14}, 127 (1952).
\bibitem{WilsonCowan1972} H. R. Wilson and J. D. Cowan, \emph{Excitatory and inhibitory interactions in localized populations of model neurons}, Biophys. J. \textbf{12}, 1 (1972).
\bibitem{AllenStevens1994} C. Allen and C. F. Stevens, \emph{An evaluation of causes for unreliability of synaptic transmission}, Proc. Natl. Acad. Sci. USA \textbf{91}, 10380 (1994).
\bibitem{Barlow1961} H. B. Barlow, \emph{Possible principles underlying the transformations of sensory messages}, in \emph{Sensory Communication}, ed. W. A. Rosenblith (MIT Press, Cambridge, MA, 1961), pp.~217--234.
\bibitem{Cardin2009} J. A. Cardin et al., \emph{Driving fast-spiking cells induces gamma rhythm and controls sensory responses}, Nature \textbf{459}, 663 (2009).
\bibitem{Lillicrap2020} T. P. Lillicrap, A. Santoro, L. Marris, C. J. Akerman, and G. Hinton, \emph{Backpropagation and the brain}, Nat. Rev. Neurosci. \textbf{21}, 335 (2020).
\bibitem{Fremaux2016} N. Fr\'emaux and W. Gerstner, \emph{Neuromodulated spike-timing-dependent plasticity, and theory of three-factor learning rules}, Front. Neural Circuits \textbf{9}, 85 (2016).
\bibitem{Haider2006} B. Haider, A. Duque, A. R. Hasenstaub, and D. A. McCormick, \emph{Neocortical network activity in vivo is generated through a dynamic balance of excitation and inhibition}, J. Neurosci. \textbf{26}, 4535 (2006).
\bibitem{HoltKoch1997} G. R. Holt and C. Koch, \emph{Shunting inhibition does not have a divisive effect on firing rates}, Neural Comput. \textbf{9}, 1001 (1997).
\bibitem{Chance2002} F. S. Chance, L. F. Abbott, and A. D. Reyes, \emph{Gain modulation from background synaptic input}, Neuron \textbf{35}, 773 (2002).
\bibitem{Ozeki2009} H. Ozeki, I. M. Finn, E. S. Schaffer, K. D. Miller, and D. Ferster, \emph{Inhibitory stabilization of the cortical network underlies visual surround suppression}, Neuron \textbf{62}, 578 (2009).
\bibitem{HasselmoBower1992} M. E. Hasselmo and J. M. Bower, \emph{Cholinergic suppression specific to intrinsic not afferent fiber synapses in rat piriform (olfactory) cortex}, J. Neurophysiol. \textbf{67}, 1222 (1992).
\bibitem{HasselmoAndersonBower1992} M. E. Hasselmo, B. P. Anderson, and J. M. Bower, \emph{Cholinergic modulation of cortical associative memory function}, J. Neurophysiol. \textbf{67}, 1230 (1992).
\bibitem{Hasselmo1999} M. E. Hasselmo, \emph{Neuromodulation: acetylcholine and memory consolidation}, Trends Cogn. Sci. \textbf{3}, 351 (1999).
\bibitem{Tsodyks1988} M. V. Tsodyks and M. V. Feigel'man, \emph{The enhanced storage capacity in neural networks with low activity level}, Europhys. Lett. \textbf{6}, 101 (1988).
\bibitem{KalmanBucy1961} R. E. Kalman and R. S. Bucy, \emph{New results in linear filtering and prediction theory}, J. Basic Eng. \textbf{83}, 95 (1961).
\bibitem{WilsonMcNaughton1994} M. A. Wilson and B. L. McNaughton, \emph{Reactivation of hippocampal ensemble memories during sleep}, Science \textbf{265}, 676 (1994).
\bibitem{BarnesSharp1999} N. M. Barnes and T. Sharp, \emph{A review of central 5-HT receptors and their function}, Neuropharmacology \textbf{38}, 1083 (1999).
\bibitem{Bunin1998} M. A. Bunin and R. M. Wightman, \emph{Quantitative evaluation of 5-hydroxytryptamine (serotonin) neuronal release and uptake: an investigation of extrasynaptic transmission}, J. Neurosci. \textbf{18}, 4854 (1998).
\bibitem{Sprouse1987} J. S. Sprouse and G. K. Aghajanian, \emph{Electrophysiological responses of serotoninergic dorsal raphe neurons to 5-HT$_{1A}$ and 5-HT$_{1B}$ agonists}, Synapse \textbf{1}, 3 (1987).
\bibitem{Sykova2008} E. Syková and C. Nicholson, \emph{Diffusion in brain extracellular space}, Physiol. Rev. \textbf{88}, 1277 (2008).
\bibitem{WilsonNicoll2001} R. I. Wilson and R. A. Nicoll, \emph{Endogenous cannabinoids mediate retrograde signalling at hippocampal synapses}, Nature \textbf{410}, 588 (2001).
\bibitem{Garthwaite2008} J. Garthwaite, \emph{Concepts of neural nitric oxide-mediated transmission}, Eur. J. Neurosci. \textbf{27}, 2783 (2008).
\bibitem{vandenPol2012} A. N. van den Pol, \emph{Neuropeptide transmission in brain circuits}, Neuron \textbf{76}, 98 (2012).
\bibitem{Dunwiddie2001} T. V. Dunwiddie and S. A. Masino, \emph{The role and regulation of adenosine in the central nervous system}, Annu. Rev. Neurosci. \textbf{24}, 31 (2001).
\bibitem{Skaggs1996} W. E. Skaggs and B. L. McNaughton, \emph{Replay of neuronal firing sequences in rat hippocampus during sleep following spatial experience}, Science \textbf{271}, 1870 (1996).
\bibitem{Baylor1979} D. A. Baylor, T. D. Lamb, and K.-W. Yau, \emph{Responses of retinal rods to single photons}, J. Physiol. \textbf{288}, 613 (1979).
\bibitem{Hudspeth1989} A. J. Hudspeth, \emph{How the ear's works work}, Nature \textbf{341}, 397 (1989).
\bibitem{NakaRushton1966} K. I. Naka and W. A. H. Rushton, \emph{S-potentials from colour units in the retina of fish (Cyprinidae)}, J. Physiol. \textbf{185}, 536 (1966).
\bibitem{Lichtsteiner2008} P. Lichtsteiner, C. Posch, and T. Delbruck, \emph{A $128\times128$ 120 dB 15 $\mu$s latency asynchronous temporal contrast vision sensor}, IEEE J. Solid-State Circuits \textbf{43}, 566 (2008).
\bibitem{Koch2006} K. Koch et al., \emph{How much the eye tells the brain}, Curr. Biol. \textbf{16}, 1428 (2006).
\bibitem{Robles2001} L. Robles and M. A. Ruggero, \emph{Mechanics of the mammalian cochlea}, Physiol. Rev. \textbf{81}, 1305 (2001).
\bibitem{Mombaerts2004} P. Mombaerts, \emph{Genes and ligands for odorant, vomeronasal and taste receptors}, Nat. Rev. Neurosci. \textbf{5}, 263 (2004).
\bibitem{WoodSlater2001} S. J. Wood and C. R. Slater, \emph{Safety factor at the neuromuscular junction}, Prog. Neurobiol. \textbf{64}, 393 (2001).
\bibitem{Henneman1957} E. Henneman, \emph{Relation between size of neurons and their susceptibility to discharge}, Science \textbf{126}, 1345 (1957).
\bibitem{HarrisWolpert1998} C. M. Harris and D. M. Wolpert, \emph{Signal-dependent noise determines motor planning}, Nature \textbf{394}, 780 (1998).
\bibitem{Hogan1984} N. Hogan, \emph{Adaptive control of mechanical impedance by coactivation of antagonist muscles}, IEEE Trans. Autom. Control \textbf{29}, 681 (1984).
\bibitem{McAuleyMarsden2000} J. H. McAuley and C. D. Marsden, \emph{Physiological and pathological tremors and rhythmic central motor control}, Brain \textbf{123}, 1545 (2000).
\bibitem{WolpertGhahramani2000} D. M. Wolpert and Z. Ghahramani, \emph{Computational principles of movement neuroscience}, Nat. Neurosci. \textbf{3}, 1212 (2000).
\bibitem{MarderBucher2001} E. Marder and D. Bucher, \emph{Central pattern generators and the control of rhythmic movements}, Curr. Biol. \textbf{11}, R986 (2001).
\bibitem{Miyazaki2014} K. W. Miyazaki et al., \emph{Optogenetic activation of dorsal raphe serotonin neurons enhances patience for future rewards}, Curr. Biol. \textbf{24}, 2033 (2014).
\bibitem{Fuglevand1993} A. J. Fuglevand, D. A. Winter, and A. E. Patla, \emph{Models of recruitment and rate coding organization in motor-unit pools}, J. Neurophysiol. \textbf{70}, 2470 (1993).
\bibitem{Curcio1990} C. A. Curcio, K. R. Sloan, R. E. Kalina, and A. E. Hendrickson, \emph{Human photoreceptor topography}, J. Comp. Neurol. \textbf{292}, 497 (1990).
\bibitem{Hayes1950} N. D. Hayes, \emph{Roots of the transcendental equation associated with a certain difference-differential equation}, J. London Math. Soc. \textbf{25}, 226 (1950).
\bibitem{MelzackWall1965} R. Melzack and P. D. Wall, \emph{Pain mechanisms: a new theory}, Science \textbf{150}, 971 (1965).
\bibitem{Fields2004} H. Fields, \emph{State-dependent opioid control of pain}, Nat. Rev. Neurosci. \textbf{5}, 565 (2004).
\bibitem{Levine1978} J. D. Levine, N. C. Gordon, and H. L. Fields, \emph{The mechanism of placebo analgesia}, Lancet \textbf{312}, 654 (1978).
\bibitem{Woolf1983} C. J. Woolf, \emph{Evidence for a central component of post-injury pain hypersensitivity}, Nature \textbf{306}, 686 (1983).
\bibitem{Seymour2004} B. Seymour et al., \emph{Temporal difference models describe higher-order learning in humans}, Nature \textbf{429}, 664 (2004).
\bibitem{EcclestonCrombez1999} C. Eccleston and G. Crombez, \emph{Pain demands attention: a cognitive-affective model of the interruptive function of pain}, Psychol. Bull. \textbf{125}, 356 (1999).
\bibitem{WidrowHoff1960} B. Widrow and M. E. Hoff, \emph{Adaptive switching circuits}, IRE WESCON Convention Record, part 4, 96 (1960).
\bibitem{Marr1969} D. Marr, \emph{A theory of cerebellar cortex}, J. Physiol. \textbf{202}, 437 (1969).
\bibitem{Albus1971} J. S. Albus, \emph{A theory of cerebellar function}, Math. Biosci. \textbf{10}, 25 (1971).
\bibitem{Cover1965} T. M. Cover, \emph{Geometrical and statistical properties of systems of linear inequalities with applications in pattern recognition}, IEEE Trans. Electron. Comput. \textbf{EC-14}, 326 (1965).
\bibitem{Ito2001} M. Ito, \emph{Cerebellar long-term depression: characterization, signal transduction, and functional roles}, Physiol. Rev. \textbf{81}, 1143 (2001).
\bibitem{Suvrathan2016} A. Suvrathan, H. L. Payne, and J. L. Raymond, \emph{Timing rules for synaptic plasticity matched to behavioral function}, Neuron \textbf{92}, 959 (2016).
\bibitem{Fujita1982} M. Fujita, \emph{Adaptive filter model of the cerebellum}, Biol. Cybern. \textbf{45}, 195 (1982).
\bibitem{Blakemore1998} S.-J. Blakemore, D. M. Wolpert, and C. D. Frith, \emph{Central cancellation of self-produced tickle sensation}, Nat. Neurosci. \textbf{1}, 635 (1998).
\bibitem{Smith2006} M. A. Smith, A. Ghazizadeh, and R. Shadmehr, \emph{Interacting adaptive processes with different timescales underlie short-term motor learning}, PLoS Biol. \textbf{4}, e179 (2006).
\bibitem{Kording2007} K. P. K\"ording, J. B. Tenenbaum, and R. Shadmehr, \emph{The dynamics of memory as a consequence of optimal adaptation to a changing body}, Nat. Neurosci. \textbf{10}, 779 (2007).
\bibitem{Yao2023} Z. Yao et al., \emph{A high-resolution transcriptomic and spatial atlas of cell types in the whole mouse brain}, Nature \textbf{624}, 317 (2023).
\bibitem{Sanzeni2020} A. Sanzeni, B. Akitake, H. C. Goldbach, C. E. Leedy, N. Brunel, and M. H. Histed, \emph{Inhibition stabilization is a widespread property of cortical networks}, eLife \textbf{9}, e54875 (2020).
\bibitem{Padamsey2022} Z. Padamsey, D. Katsanevaki, N. Dupuy, and N. L. Rochefort, \emph{Neocortex saves energy by reducing coding precision during food scarcity}, Neuron \textbf{110}, 280 (2022).
\bibitem{Stringer2019} C. Stringer, M. Pachitariu, N. Steinmetz, C. B. Reddy, M. Carandini, and K. D. Harris, \emph{Spontaneous behaviors drive multidimensional, brainwide activity}, Science \textbf{364}, eaav7893 (2019).
\bibitem{Bittner2017} K. C. Bittner, A. D. Milstein, C. Grienberger, S. Romani, and J. C. Magee, \emph{Behavioral time scale synaptic plasticity underlies CA1 place fields}, Science \textbf{357}, 1033 (2017).
\bibitem{WhittingtonBogacz2019} J. C. R. Whittington and R. Bogacz, \emph{Theories of error back-propagation in the brain}, Trends Cogn. Sci. \textbf{23}, 235 (2019).
\bibitem{Reimer2016} J. Reimer et al., \emph{Pupil fluctuations track rapid changes in adrenergic and cholinergic activity in cortex}, Nat. Commun. \textbf{7}, 13289 (2016).
\bibitem{Beniaguev2021} D. Beniaguev, I. Segev, and M. London, \emph{Single cortical neurons as deep artificial neural networks}, Neuron \textbf{109}, 2727 (2021).
\bibitem{Gidon2020} A. Gidon et al., \emph{Dendritic action potentials and computation in human layer 2/3 cortical neurons}, Science \textbf{367}, 83 (2020).
\bibitem{Gallego2022} G. Gallego et al., \emph{Event-based vision: a survey}, IEEE Trans. Pattern Anal. Mach. Intell. \textbf{44}, 154 (2022).
\bibitem{Berger1989} R. D. Berger, J. P. Saul, and R. J. Cohen, \emph{Transfer function analysis of autonomic regulation. I. Canine atrial rate response}, Am. J. Physiol. \textbf{256}, H142 (1989).
\bibitem{Walker2010} J. J. Walker, J. R. Terry, and S. L. Lightman, \emph{Origin of ultradian pulsatility in the hypothalamic--pituitary--adrenal axis}, Proc. R. Soc. B \textbf{277}, 1627 (2010).
\bibitem{Belchetz1978} P. E. Belchetz, T. M. Plant, Y. Nakai, E. J. Keogh, and E. Knobil, \emph{Hypophysial responses to continuous and intermittent delivery of hypothalamic gonadotropin-releasing hormone}, Science \textbf{202}, 631 (1978).
\bibitem{Czeisler1999} C. A. Czeisler et al., \emph{Stability, precision, and near-24-hour period of the human circadian pacemaker}, Science \textbf{284}, 2177 (1999).
\bibitem{TakahashiJS2017} J. S. Takahashi, \emph{Transcriptional architecture of the mammalian circadian clock}, Nat. Rev. Genet. \textbf{18}, 164 (2017).
\bibitem{Musk2019} E. Musk and Neuralink, \emph{An integrated brain-machine interface platform with thousands of channels}, J. Med. Internet Res. \textbf{21}, e16194 (2019).
\bibitem{Hochberg2006} L. R. Hochberg et al., \emph{Neuronal ensemble control of prosthetic devices by a human with tetraplegia}, Nature \textbf{442}, 164 (2006).
\bibitem{Willett2021} F. R. Willett et al., \emph{High-performance brain-to-text communication via handwriting}, Nature \textbf{593}, 249 (2021).
\bibitem{Willett2023} F. R. Willett et al., \emph{A high-performance speech neuroprosthesis}, Nature \textbf{620}, 1031 (2023).
\bibitem{Gallego2017} J. A. Gallego, M. G. Perich, L. E. Miller, and S. A. Solla, \emph{Neural manifolds for the control of movement}, Neuron \textbf{94}, 978 (2017).
\bibitem{Fernandez2021} E. Fern\'andez et al., \emph{Visual percepts evoked with an intracortical 96-channel microelectrode array inserted in human occipital cortex}, J. Clin. Invest. \textbf{131}, e151331 (2021).
\bibitem{Tritsch2012} N. X. Tritsch, J. B. Ding, and B. L. Sabatini, \emph{Dopaminergic neurons inhibit striatal output through non-canonical release of GABA}, Nature \textbf{490}, 262 (2012).
\bibitem{Zhang2015} S. Zhang et al., \emph{Dopaminergic and glutamatergic microdomains in a subset of rodent mesoaccumbens axons}, Nat. Neurosci. \textbf{18}, 386 (2015).
\bibitem{Mongillo2008} G. Mongillo, O. Barak, and M. Tsodyks, \emph{Synaptic theory of working memory}, Science \textbf{319}, 1543 (2008).
\bibitem{Stokes2015} M. G. Stokes, \emph{`Activity-silent' working memory in prefrontal cortex: a dynamic coding framework}, Trends Cogn. Sci. \textbf{19}, 394 (2015).
\bibitem{Smith2013} S. L. Smith, I. T. Smith, T. Branco, and M. H\"ausser, \emph{Dendritic spikes enhance stimulus selectivity in cortical neurons in vivo}, Nature \textbf{503}, 115 (2013).
\bibitem{Baden2016} T. Baden et al., \emph{The functional diversity of retinal ganglion cells in the mouse}, Nature \textbf{529}, 345 (2016).
\bibitem{Lottem2018} E. Lottem et al., \emph{Activation of serotonin neurons promotes active persistence in a probabilistic foraging task}, Nat. Commun. \textbf{9}, 1000 (2018).
\bibitem{Payeur2021} A. Payeur, J. Guerguiev, F. Zenke, B. A. Richards, and R. Naud, \emph{Burst-dependent synaptic plasticity can coordinate learning in hierarchical circuits}, Nat. Neurosci. \textbf{24}, 1010 (2021).
\bibitem{MehonicKenyon2022} A. Mehonic and A. J. Kenyon, \emph{Brain-inspired computing needs a master plan}, Nature \textbf{604}, 255 (2022).
\bibitem{Poe2020} G. R. Poe et al., \emph{Locus coeruleus: a new look at the blue spot}, Nat. Rev. Neurosci. \textbf{21}, 644 (2020).
\bibitem{Hangya2015} B. Hangya, S. P. Ranade, M. Lorenc, and A. Kepecs, \emph{Central cholinergic neurons are rapidly recruited by reinforcement feedback}, Cell \textbf{162}, 1155 (2015).
\bibitem{FernandezRuiz2019} A. Fern\'andez-Ruiz et al., \emph{Long-duration hippocampal sharp wave ripples improve memory}, Science \textbf{364}, 1082 (2019).
\bibitem{Herzfeld2018} D. J. Herzfeld, Y. Kojima, R. Soetedjo, and R. Shadmehr, \emph{Encoding of error and learning to correct that error by the Purkinje cells of the cerebellum}, Nat. Neurosci. \textbf{21}, 736 (2018).
\bibitem{Duan2014} B. Duan et al., \emph{Identification of spinal circuits transmitting and gating mechanical pain}, Cell \textbf{159}, 1417 (2014).
\bibitem{Tremblay2016} R. Tremblay, S. Lee, and B. Rudy, \emph{GABAergic interneurons in the neocortex: from cellular properties to circuits}, Neuron \textbf{91}, 260 (2016).
\bibitem{Ren2018} J. Ren et al., \emph{Anatomically defined and functionally distinct dorsal raphe serotonin sub-systems}, Cell \textbf{175}, 472 (2018).
\bibitem{Pfeffer2013} C. K. Pfeffer, M. Xue, M. He, Z. J. Huang, and M. Scanziani, \emph{Inhibition of inhibition in visual cortex: the logic of connections between molecularly distinct interneurons}, Nat. Neurosci. \textbf{16}, 1068 (2013).
\bibitem{Pi2013} H.-J. Pi et al., \emph{Cortical interneurons that specialize in disinhibitory control}, Nature \textbf{503}, 521 (2013).
\bibitem{Keller2018} G. B. Keller and T. D. Mrsic-Flogel, \emph{Predictive processing: a canonical cortical computation}, Neuron \textbf{100}, 424 (2018).
\bibitem{HubelWiesel1962} D. H. Hubel and T. N. Wiesel, \emph{Receptive fields, binocular interaction and functional architecture in the cat's visual cortex}, J. Physiol. \textbf{160}, 106 (1962).
\bibitem{Riesenhuber1999} M. Riesenhuber and T. Poggio, \emph{Hierarchical models of object recognition in cortex}, Nat. Neurosci. \textbf{2}, 1019 (1999).
\bibitem{Fukushima1980} K. Fukushima, \emph{Neocognitron: a self-organizing neural network model for a mechanism of pattern recognition unaffected by shift in position}, Biol. Cybern. \textbf{36}, 193 (1980).
\bibitem{Yamins2014} D. L. K. Yamins et al., \emph{Performance-optimized hierarchical models predict neural responses in higher visual cortex}, Proc. Natl. Acad. Sci. USA \textbf{111}, 8619 (2014).
\bibitem{Hung2005} C. P. Hung, G. Kreiman, T. Poggio, and J. J. DiCarlo, \emph{Fast readout of object identity from macaque inferior temporal cortex}, Science \textbf{310}, 863 (2005).
\bibitem{WiskottSejnowski2002} L. Wiskott and T. J. Sejnowski, \emph{Slow feature analysis: unsupervised learning of invariances}, Neural Comput. \textbf{14}, 715 (2002).
\bibitem{Foldiak1991} P. F\"oldi\'ak, \emph{Learning invariance from transformation sequences}, Neural Comput. \textbf{3}, 194 (1991).
\bibitem{LiDiCarlo2008} N. Li and J. J. DiCarlo, \emph{Unsupervised natural experience rapidly alters invariant object representation in visual cortex}, Science \textbf{321}, 1502 (2008).
\bibitem{RaoBallard1999} R. P. N. Rao and D. H. Ballard, \emph{Predictive coding in the visual cortex: a functional interpretation of some extra-classical receptive-field effects}, Nat. Neurosci. \textbf{2}, 79 (1999).
\bibitem{Kar2019} K. Kar, J. Kubilius, K. Schmidt, E. B. Issa, and J. J. DiCarlo, \emph{Evidence that recurrent circuits are critical to the ventral stream's execution of core object recognition behavior}, Nat. Neurosci. \textbf{22}, 974 (2019).
\bibitem{Szegedy2014} C. Szegedy et al., \emph{Intriguing properties of neural networks}, in \emph{Proc. ICLR} (2014), arXiv:1312.6199.
\bibitem{Weiss2002} Y. Weiss, E. P. Simoncelli, and E. H. Adelson, \emph{Motion illusions as optimal percepts}, Nat. Neurosci. \textbf{5}, 598 (2002).
\bibitem{Purves2011} D. Purves, W. T. Wojtach, and R. B. Lotto, \emph{Understanding vision in wholly empirical terms}, Proc. Natl. Acad. Sci. USA \textbf{108}, Suppl.~3, 15588 (2011).
\bibitem{LaferSousa2015} R. Lafer-Sousa, K. L. Hermann, and B. R. Conway, \emph{Striking individual differences in color perception uncovered by `the dress' photograph}, Curr. Biol. \textbf{25}, R545 (2015).
\bibitem{Wallisch2017} P. Wallisch, \emph{Illumination assumptions account for individual differences in the perceptual interpretation of a profoundly ambiguous stimulus in the color domain: ``The dress''}, J. Vis. \textbf{17}(4), 5 (2017).
\bibitem{Schwartz2009} O. Schwartz, T. J. Sejnowski, and P. Dayan, \emph{Perceptual organization in the tilt illusion}, J. Vis. \textbf{9}(4), 19 (2009).
\bibitem{SchillerCarvey2005} P. H. Schiller and C. E. Carvey, \emph{The Hermann grid illusion revisited}, Perception \textbf{34}, 1375 (2005).
\bibitem{Nijhawan1994} R. Nijhawan, \emph{Motion extrapolation in catching}, Nature \textbf{370}, 256 (1994).
\bibitem{Conway2005} B. R. Conway, A. Kitaoka, A. Yazdanbakhsh, C. C. Pack, and M. S. Livingstone, \emph{Neural basis for a powerful static motion illusion}, J. Neurosci. \textbf{25}, 5651 (2005).
\bibitem{OteroMillan2012} J. Otero-Millan, S. L. Macknik, and S. Martinez-Conde, \emph{Microsaccades and blinks trigger illusory rotation in the ``rotating snakes'' illusion}, J. Neurosci. \textbf{32}, 6043 (2012).
\bibitem{MorenoBote2007} R. Moreno-Bote, J. Rinzel, and N. Rubin, \emph{Noise-induced alternations in an attractor network model of perceptual bistability}, J. Neurophysiol. \textbf{98}, 1125 (2007).
\bibitem{Watanabe2018} E. Watanabe, A. Kitaoka, K. Sakamoto, M. Yasugi, and K. Tanaka, \emph{Illusory motion reproduced by deep neural networks trained for prediction}, Front. Psychol. \textbf{9}, 345 (2018).
\bibitem{GomezVilla2020} A. G\'omez-Villa, A. Mart\'in, J. Vazquez-Corral, M. Bertalm\'io, and J. Malo, \emph{Color illusions also deceive CNNs for low-level vision tasks: analysis and implications}, Vision Res. \textbf{176}, 156 (2020).
\bibitem{delCastillo1954} J. del Castillo and B. Katz, \emph{Quantal components of the end-plate potential}, J. Physiol. \textbf{124}, 560 (1954).
\bibitem{Nowak1984} L. Nowak, P. Bregestovski, P. Ascher, A. Herbet, and A. Prochiantz, \emph{Magnesium gates glutamate-activated channels in mouse central neurones}, Nature \textbf{307}, 462 (1984).
\bibitem{Malinow2002} R. Malinow and R. C. Malenka, \emph{AMPA receptor trafficking and synaptic plasticity}, Annu. Rev. Neurosci. \textbf{25}, 103 (2002).
\bibitem{Matsuzaki2004} M. Matsuzaki, N. Honkura, G. C. R. Ellis-Davies, and H. Kasai, \emph{Structural basis of long-term potentiation in single dendritic spines}, Nature \textbf{429}, 761 (2004).
\bibitem{Rudomin1999} P. Rudomin and R. F. Schmidt, \emph{Presynaptic inhibition in the vertebrate spinal cord revisited}, Exp. Brain Res. \textbf{129}, 1 (1999).
\bibitem{HutcheonYarom2000} B. Hutcheon and Y. Yarom, \emph{Resonance, oscillation and the intrinsic frequency preferences of neurons}, Trends Neurosci. \textbf{23}, 216 (2000).
\bibitem{Seung1996} H. S. Seung, \emph{How the brain keeps the eyes still}, Proc. Natl. Acad. Sci. USA \textbf{93}, 13339 (1996).
\bibitem{Hounsgaard1988} J. Hounsgaard, H. Hultborn, B. Jespersen, and O. Kiehn, \emph{Bistability of alpha-motoneurones in the decerebrate cat and in the acute spinal cat after intravenous 5-hydroxytryptophan}, J. Physiol. \textbf{405}, 345 (1988).
\bibitem{Izhikevich2007} E. M. Izhikevich, \emph{Dynamical Systems in Neuroscience: The Geometry of Excitability and Bursting} (MIT Press, Cambridge, MA, 2007).
\bibitem{AgmonSnir1998} H. Agmon-Snir, C. E. Carr, and J. Rinzel, \emph{The role of dendrites in auditory coincidence detection}, Nature \textbf{393}, 268 (1998).
\bibitem{BorstSoria2012} J. G. G. Borst and J. Soria van Hoeve, \emph{The calyx of Held synapse: from model synapse to auditory relay}, Annu. Rev. Physiol. \textbf{74}, 199 (2012).
\bibitem{Euler2002} T. Euler, P. B. Detwiler, and W. Denk, \emph{Directionally selective calcium signals in dendrites of starburst amacrine cells}, Nature \textbf{418}, 845 (2002).
\bibitem{AstonJonesBloom1981} G. Aston-Jones and F. E. Bloom, \emph{Activity of norepinephrine-containing locus coeruleus neurons in behaving rats anticipates fluctuations in the sleep-waking cycle}, J. Neurosci. \textbf{1}, 876 (1981).
\bibitem{BrooksPeever2012} P. L. Brooks and J. H. Peever, \emph{Identification of the transmitter and receptor mechanisms responsible for REM sleep paralysis}, J. Neurosci. \textbf{32}, 9785 (2012).
\bibitem{Borbely1982} A. A. Borb\'ely, \emph{A two process model of sleep regulation}, Hum. Neurobiol. \textbf{1}, 195 (1982).
\bibitem{PorkkaHeiskanen1997} T. Porkka-Heiskanen et al., \emph{Adenosine: a mediator of the sleep-inducing effects of prolonged wakefulness}, Science \textbf{276}, 1265 (1997).
\bibitem{Steriade1993} M. Steriade, A. N\'u\~nez, and F. Amzica, \emph{A novel slow ($<1$ Hz) oscillation of neocortical neurons in vivo: depolarizing and hyperpolarizing components}, J. Neurosci. \textbf{13}, 3252 (1993).
\bibitem{McCormick1992} D. A. McCormick, \emph{Neurotransmitter actions in the thalamus and cerebral cortex and their role in neuromodulation of thalamocortical activity}, Prog. Neurobiol. \textbf{39}, 337 (1992).
\bibitem{SteriadeMcCormickSejnowski1993} M. Steriade, D. A. McCormick, and T. J. Sejnowski, \emph{Thalamocortical oscillations in the sleeping and aroused brain}, Science \textbf{262}, 679 (1993).
\bibitem{Nadasdy1999} Z. N\'adasdy et al., \emph{Replay and time compression of recurring spike sequences in the hippocampus}, J. Neurosci. \textbf{19}, 9497 (1999).
\bibitem{Maingret2016} N. Maingret et al., \emph{Hippocampo-cortical coupling mediates memory consolidation during sleep}, Nat. Neurosci. \textbf{19}, 959 (2016).
\bibitem{McClelland1995} J. L. McClelland, B. L. McNaughton, and R. C. O'Reilly, \emph{Why there are complementary learning systems in the hippocampus and neocortex}, Psychol. Rev. \textbf{102}, 419 (1995).
\bibitem{TononiCirelli2014} G. Tononi and C. Cirelli, \emph{Sleep and the price of plasticity: from synaptic and cellular homeostasis to memory consolidation and integration}, Neuron \textbf{81}, 12 (2014).
\bibitem{deVivo2017} L. de Vivo et al., \emph{Ultrastructural evidence for synaptic scaling across the wake/sleep cycle}, Science \textbf{355}, 507 (2017).
\bibitem{Xie2013} L. Xie et al., \emph{Sleep drives metabolite clearance from the adult brain}, Science \textbf{342}, 373 (2013).
\bibitem{Miao2024} A. Miao et al., \emph{Brain clearance is reduced during sleep and anesthesia}, Nat. Neurosci. \textbf{27}, 1046 (2024).
\bibitem{Vyazovskiy2011} V. V. Vyazovskiy et al., \emph{Local sleep in awake rats}, Nature \textbf{472}, 443 (2011).
\bibitem{Lancaster2013} M. A. Lancaster et al., \emph{Cerebral organoids model human brain development and microcephaly}, Nature \textbf{501}, 373 (2013).
\bibitem{Jordan2024} F. D. Jordan et al., \emph{Open and remotely accessible Neuroplatform for research in wetware computing}, Front. Artif. Intell. \textbf{7}, 1376042 (2024).
\bibitem{Smirnova2023} L. Smirnova et al., \emph{Organoid intelligence (OI): the new frontier in biocomputing and intelligence-in-a-dish}, Front. Sci. \textbf{1}, 1017235 (2023).
\bibitem{Wagenaar2006} D. A. Wagenaar, J. Pine, and S. M. Potter, \emph{An extremely rich repertoire of bursting patterns during the development of cortical cultures}, BMC Neurosci. \textbf{7}, 11 (2006).
\bibitem{Shahaf2001} G. Shahaf and S. Marom, \emph{Learning in networks of cortical neurons}, J. Neurosci. \textbf{21}, 8782 (2001).
\bibitem{Kagan2022} B. J. Kagan et al., \emph{In vitro neurons learn and exhibit sentience when embodied in a simulated game-world}, Neuron \textbf{110}, 3952 (2022).
\bibitem{Friston2010} K. Friston, \emph{The free-energy principle: a unified brain theory?}, Nat. Rev. Neurosci. \textbf{11}, 127 (2010).
\bibitem{Balci2023} F. Balci et al., \emph{A response to claims of emergent intelligence and sentience in a dish}, Neuron \textbf{111}, 604 (2023).
\bibitem{Bakkum2008} D. J. Bakkum, Z. C. Chao, and S. M. Potter, \emph{Spatio-temporal electrical stimuli shape behavior of an embodied cortical network in a goal-directed learning task}, J. Neural Eng. \textbf{5}, 310 (2008).
\bibitem{Maass2002} W. Maass, T. Natschl\"ager, and H. Markram, \emph{Real-time computing without stable states: a new framework for neural computation based on perturbations}, Neural Comput. \textbf{14}, 2531 (2002).
\bibitem{JaegerHaas2004} H. Jaeger and H. Haas, \emph{Harnessing nonlinearity: predicting chaotic systems and saving energy in wireless communication}, Science \textbf{304}, 78 (2004).
\bibitem{Cai2023} H. Cai et al., \emph{Brain organoid reservoir computing for artificial intelligence}, Nat. Electron. \textbf{6}, 1032 (2023).
\bibitem{Habibollahi2023} F. Habibollahi, B. J. Kagan, A. N. Burkitt, and C. French, \emph{Critical dynamics arise during structured information presentation within embodied in vitro neuronal networks}, Nat. Commun. \textbf{14}, 5287 (2023).
\bibitem{Robbins2026} A. Robbins et al., \emph{Goal-directed learning in cortical organoids}, Cell Rep. \textbf{45}, 116984 (2026).
\bibitem{Keene2020} S. T. Keene et al., \emph{A biohybrid synapse with neurotransmitter-mediated plasticity}, Nat. Mater. \textbf{19}, 969 (2020).
\bibitem{Jo2016} J. Jo et al., \emph{Midbrain-like organoids from human pluripotent stem cells contain functional dopaminergic and neuromelanin-producing neurons}, Cell Stem Cell \textbf{19}, 248 (2016).
\bibitem{Hendry1987} S. H. Hendry, H. D. Schwark, E. G. Jones, and J. Yan, \emph{Numbers and proportions of GABA-immunoreactive neurons in different areas of monkey cerebral cortex}, J. Neurosci. \textbf{7}, 1503 (1987).
\bibitem{Tang2001synapses} Y. Tang, J. R. Nyengaard, D. M. G. De Groot, and H. J. G. Gundersen, \emph{Total regional and global number of synapses in the human brain neocortex}, Synapse \textbf{41}, 258 (2001).
\bibitem{Pakkenberg1997} B. Pakkenberg and H. J. G. Gundersen, \emph{Neocortical neuron number in humans: effect of sex and age}, J. Comp. Neurol. \textbf{384}, 312 (1997).
\bibitem{Matsuda2009} W. Matsuda, T. Furuta, K. C. Nakamura, H. Hioki, F. Fujiyama, R. Arai, and T. Kaneko, \emph{Single nigrostriatal dopaminergic neurons form widely spread and highly dense axonal arborizations in the neostriatum}, J. Neurosci. \textbf{29}, 444 (2009).
\bibitem{Pakkenberg1991} B. Pakkenberg, A. M{\o}ller, H. J. G. Gundersen, A. Mouritzen Dam, and H. Pakkenberg, \emph{The absolute number of nerve cells in substantia nigra in normal subjects and in patients with Parkinson's disease estimated with an unbiased stereological method}, J. Neurol. Neurosurg. Psychiatry \textbf{54}, 30 (1991).
\bibitem{Baker1990} K. G. Baker, G. M. Halliday, and I. T\"ork, \emph{Cytoarchitecture of the human dorsal raphe nucleus}, J. Comp. Neurol. \textbf{301}, 147 (1990).
\bibitem{Baker1991} K. G. Baker, G. M. Halliday, P. Halasz, J.-P. Hornung, L. B. Geffen, R. G. H. Cotton, and I. T\"ork, \emph{Cytoarchitecture of serotonin-synthesizing neurons in the pontine tegmentum of the human brain}, Synapse \textbf{7}, 301 (1991).
\bibitem{Airaksinen1991} M. S. Airaksinen, A. Paetau, L. Palj\"arvi, K. Reinikainen, P. Riekkinen, R. Suomalainen, and P. Panula, \emph{Histamine neurons in human hypothalamus: anatomy in normal and Alzheimer diseased brains}, Neuroscience \textbf{44}, 465 (1991).
\bibitem{Colquhoun1992} D. Colquhoun, P. Jonas, and B. Sakmann, \emph{Action of brief pulses of glutamate on AMPA/kainate receptors in patches from different neurones of rat hippocampal slices}, J. Physiol. \textbf{458}, 261 (1992).
\bibitem{DutarNicoll1988} P. Dutar and R. A. Nicoll, \emph{A physiological role for GABA$_B$ receptors in the central nervous system}, Nature \textbf{332}, 156 (1988).

\bibitem{Rauch2003} A. Rauch, G. La Camera, H.-R. L\"uscher, W. Senn, and S. Fusi, \emph{Neocortical pyramidal cells respond as integrate-and-fire neurons to in vivo-like input currents}, J. Neurophysiol. \textbf{90}, 1598 (2003).
\bibitem{KlumppEady1956} R. G. Klumpp and H. R. Eady, \emph{Some measurements of interaural time difference thresholds}, J. Acoust. Soc. Am. \textbf{28}, 859 (1956).
\bibitem{Brand2002} A. Brand, O. Behrend, T. Marquardt, D. McAlpine, and B. Grothe, \emph{Precise inhibition is essential for microsecond interaural time difference coding}, Nature \textbf{417}, 543 (2002).
\bibitem{Funahashi1989} S. Funahashi, C. J. Bruce, and P. S. Goldman-Rakic, \emph{Mnemonic coding of visual space in the monkey's dorsolateral prefrontal cortex}, J. Neurophysiol. \textbf{61}, 331 (1989).
\bibitem{WangMarkram2006} Y. Wang, H. Markram, P. H. Goodman, T. K. Berger, J. Ma, and P. S. Goldman-Rakic, \emph{Heterogeneity in the pyramidal network of the medial prefrontal cortex}, Nat. Neurosci. \textbf{9}, 534 (2006).

\bibitem{McCullochPitts1943} W. S. McCulloch and W. Pitts, \emph{A logical calculus of the ideas immanent in nervous activity}, Bull. Math. Biophys. \textbf{5}, 115 (1943).
\bibitem{Fried2002} S. I. Fried, T. A. M\"unch, and F. S. Werblin, \emph{Mechanisms and circuitry underlying directional selectivity in the retina}, Nature \textbf{420}, 411 (2002).

\bibitem{Hromadka2008} T. Hrom\'adka, M. R. DeWeese, and A. M. Zador, \emph{Sparse representation of sounds in the unanesthetized auditory cortex}, PLoS Biol. \textbf{6}, e16 (2008).
\bibitem{Abbott1997depression} L. F. Abbott, J. A. Varela, K. Sen, and S. B. Nelson, \emph{Synaptic depression and cortical gain control}, Science \textbf{275}, 221 (1997).
\bibitem{ShadlenNewsome1998} M. N. Shadlen and W. T. Newsome, \emph{The variable discharge of cortical neurons: implications for connectivity, computation, and information coding}, J. Neurosci. \textbf{18}, 3870 (1998).

\bibitem{TrulsonJacobs1979} M. E. Trulson and B. L. Jacobs, \emph{Raphe unit activity in freely moving cats: correlation with level of behavioral arousal}, Brain Res. \textbf{163}, 135 (1979).
\bibitem{GagnonParent2014} D. Gagnon and M. Parent, \emph{Distribution of VGLUT3 in highly collateralized axons from the rat dorsal raphe nucleus as revealed by single-neuron reconstructions}, PLoS One \textbf{9}, e87709 (2014).
\bibitem{VandermaelenAghajanian1983} C. P. Vandermaelen and G. K. Aghajanian, \emph{Electrophysiological and pharmacological characterization of serotonergic dorsal raphe neurons recorded extracellularly and intracellularly in rat brain slices}, Brain Res. \textbf{289}, 109 (1983).
\bibitem{CourtneyFord2016} N. A. Courtney and C. P. Ford, \emph{Mechanisms of 5-HT$_{1A}$ receptor-mediated transmission in dorsal raphe serotonin neurons}, J. Physiol. \textbf{594}, 953 (2016).
\bibitem{Hashemi2012serotonin} P. Hashemi et al., \emph{Brain dopamine and serotonin differ in regulation and its consequences}, Proc. Natl. Acad. Sci. USA \textbf{109}, 11510 (2012).
\bibitem{Black1934feedback} H. S. Black, \emph{Stabilized feedback amplifiers}, Bell Syst. Tech. J. \textbf{13}, 1 (1934).
\bibitem{KobayashiSchultz2008} S. Kobayashi and W. Schultz, \emph{Influence of reward delays on responses of dopamine neurons}, J. Neurosci. \textbf{28}, 7837 (2008).
\bibitem{Schweighofer2008} N. Schweighofer et al., \emph{Low-serotonin levels increase delayed reward discounting in humans}, J. Neurosci. \textbf{28}, 4528 (2008).

\bibitem{Malenka1988calcium} R. C. Malenka, J. A. Kauer, R. S. Zucker, and R. A. Nicoll, \emph{Postsynaptic calcium is sufficient for potentiation of hippocampal synaptic transmission}, Science \textbf{242}, 81 (1988).
\bibitem{Bliss1973LTP} T. V. P. Bliss and T. L{\o}mo, \emph{Long-lasting potentiation of synaptic transmission in the dentate area of the anaesthetized rabbit following stimulation of the perforant path}, J. Physiol. \textbf{232}, 331 (1973).
\bibitem{Kelso1986hebb} S. R. Kelso, A. H. Ganong, and T. H. Brown, \emph{Hebbian synapses in hippocampus}, Proc. Natl. Acad. Sci. USA \textbf{83}, 5326 (1986).
\bibitem{He2015eligibility} K. He et al., \emph{Distinct eligibility traces for LTP and LTD in cortical synapses}, Neuron \textbf{88}, 528 (2015).
\bibitem{Olveczky2005} B. P. \"Olveczky, A. S. Andalman, and M. S. Fee, \emph{Vocal experimentation in the juvenile songbird requires a basal ganglia circuit}, PLoS Biol. \textbf{3}, e153 (2005).
\bibitem{TumerBrainard2007} E. C. Tumer and M. S. Brainard, \emph{Performance variability enables adaptive plasticity of `crystallized' adult birdsong}, Nature \textbf{450}, 1240 (2007).
\bibitem{Eshel2015arithmetic} N. Eshel et al., \emph{Arithmetic and local circuitry underlying dopamine prediction errors}, Nature \textbf{525}, 243 (2015).
\bibitem{GraceOnn1989} A. A. Grace and S.-P. Onn, \emph{Morphology and electrophysiological properties of immunocytochemically identified rat dopamine neurons recorded in vitro}, J. Neurosci. \textbf{9}, 3463 (1989).
\bibitem{ODoherty2004actor} J. O'Doherty et al., \emph{Dissociable roles of ventral and dorsal striatum in instrumental conditioning}, Science \textbf{304}, 452 (2004).

\bibitem{NeweyPowell1987} W. K. Newey and J. L. Powell, \emph{Asymmetric least squares estimation and testing}, Econometrica \textbf{55}, 819 (1987).
\bibitem{Beierholm2013vigor} U. Beierholm et al., \emph{Dopamine modulates reward-related vigor}, Neuropsychopharmacology \textbf{38}, 1495 (2013).
\bibitem{Howe2013ramp} M. W. Howe, P. L. Tierney, S. G. Sandberg, P. E. M. Phillips, and A. M. Graybiel, \emph{Prolonged dopamine signalling in striatum signals proximity and value of distant rewards}, Nature \textbf{500}, 575 (2013).

\bibitem{Baker1989LC} K. G. Baker, I. T\"ork, J.-P. Hornung, and P. Halasz, \emph{The human locus coeruleus complex: an immunohistochemical and three dimensional reconstruction study}, Exp. Brain Res. \textbf{77}, 257 (1989).
\bibitem{Schwarz2015LC} L. A. Schwarz et al., \emph{Viral-genetic tracing of the input--output organization of a central noradrenaline circuit}, Nature \textbf{524}, 88 (2015).
\bibitem{Uematsu2017LC} A. Uematsu et al., \emph{Modular organization of the brainstem noradrenaline system coordinates opposing learning states}, Nat. Neurosci. \textbf{20}, 1602 (2017).
\bibitem{AstonJones1994vigilance} G. Aston-Jones, J. Rajkowski, P. Kubiak, and T. Alexinsky, \emph{Locus coeruleus neurons in monkey are selectively activated by attended cues in a vigilance task}, J. Neurosci. \textbf{14}, 4467 (1994).
\bibitem{Clayton2004LC} E. C. Clayton, J. Rajkowski, J. D. Cohen, and G. Aston-Jones, \emph{Phasic activation of monkey locus ceruleus neurons by simple decisions in a forced-choice task}, J. Neurosci. \textbf{24}, 9914 (2004).
\bibitem{Foote1975NE} S. L. Foote, R. Freedman, and A. P. Oliver, \emph{Effects of putative neurotransmitters on neuronal activity in monkey auditory cortex}, Brain Res. \textbf{86}, 229 (1975).
\bibitem{JepmaNieuwenhuis2011} M. Jepma and S. Nieuwenhuis, \emph{Pupil diameter predicts changes in the exploration--exploitation trade-off: evidence for the adaptive gain theory}, J. Cogn. Neurosci. \textbf{23}, 1587 (2011).
\bibitem{Tervo2014stochastic} D. G. R. Tervo et al., \emph{Behavioral variability through stochastic choice and its gating by anterior cingulate cortex}, Cell \textbf{159}, 21 (2014).
\bibitem{Usher1999LC} M. Usher, J. D. Cohen, D. Servan-Schreiber, J. Rajkowski, and G. Aston-Jones, \emph{The role of locus coeruleus in the regulation of cognitive performance}, Science \textbf{283}, 549 (1999).

\bibitem{Mesulam1983} M.-M. Mesulam, E. J. Mufson, B. H. Wainer, and A. I. Levey, \emph{Central cholinergic pathways in the rat: an overview based on an alternative nomenclature (Ch1--Ch6)}, Neuroscience \textbf{10}, 1185 (1983).
\bibitem{Marrosu1995} F. Marrosu et al., \emph{Microdialysis measurement of cortical and hippocampal acetylcholine release during sleep-wake cycle in freely moving cats}, Brain Res. \textbf{671}, 329 (1995).
\bibitem{Picciotto2012} M. R. Picciotto, M. J. Higley, and Y. S. Mineur, \emph{Acetylcholine as a neuromodulator: cholinergic signaling shapes nervous system function and behavior}, Neuron \textbf{76}, 116 (2012).
\bibitem{HasselmoSchnell1994} M. E. Hasselmo and E. Schnell, \emph{Laminar selectivity of the cholinergic suppression of synaptic transmission in rat hippocampal region CA1: computational modeling and brain slice physiology}, J. Neurosci. \textbf{14}, 3898 (1994).
\bibitem{Gil1997} Z. Gil, B. W. Connors, and Y. Amitai, \emph{Differential regulation of neocortical synapses by neuromodulators and activity}, Neuron \textbf{19}, 679 (1997).
\bibitem{KilgardMerzenich1998} M. P. Kilgard and M. M. Merzenich, \emph{Cortical map reorganization enabled by nucleus basalis activity}, Science \textbf{279}, 1714 (1998).
\bibitem{Atri2004} A. Atri et al., \emph{Blockade of central cholinergic receptors impairs new learning and increases proactive interference in a word paired-associate memory task}, Behav. Neurosci. \textbf{118}, 223 (2004).
\bibitem{Girardeau2009} G. Girardeau, K. Benchenane, S. I. Wiener, G. Buzs\'aki, and M. B. Zugaro, \emph{Selective suppression of hippocampal ripples impairs spatial memory}, Nat. Neurosci. \textbf{12}, 1222 (2009).

\input{supplement/bib_10}
\bibitem{CopenhagenJahr1989} D. R. Copenhagen and C. E. Jahr, \emph{Release of endogenous excitatory amino acids from turtle photoreceptors}, Nature \textbf{341}, 536 (1989).
\bibitem{GlowatzkiFuchs2002} E. Glowatzki and P. A. Fuchs, \emph{Transmitter release at the hair cell ribbon synapse}, Nat. Neurosci. \textbf{5}, 147 (2002).
\bibitem{Tinsley2016} J. N. Tinsley et al., \emph{Direct detection of a single photon by humans}, Nat. Commun. \textbf{7}, 12172 (2016).
\bibitem{Sellick1982} P. M. Sellick, R. Patuzzi, and B. M. Johnstone, \emph{Measurement of basilar membrane motion in the guinea pig using the M\"ossbauer technique}, J. Acoust. Soc. Am. \textbf{72}, 131 (1982).
\bibitem{Ruggero1997} M. A. Ruggero, N. C. Rich, A. Recio, S. S. Narayan, and L. Robles, \emph{Basilar-membrane responses to tones at the base of the chinchilla cochlea}, J. Acoust. Soc. Am. \textbf{101}, 2151 (1997).
\bibitem{NormannPerlman1979} R. A. Normann and I. Perlman, \emph{The effects of background illumination on the photoresponses of red and green cones}, J. Physiol. \textbf{286}, 491 (1979).
\bibitem{Laughlin1981} S. Laughlin, \emph{A simple coding procedure enhances a neuron's information capacity}, Z. Naturforsch. C \textbf{36}, 910 (1981).
\bibitem{CurcioAllen1990} C. A. Curcio and K. A. Allen, \emph{Topography of ganglion cells in human retina}, J. Comp. Neurol. \textbf{300}, 5 (1990).
\bibitem{Greenwood1990} D. D. Greenwood, \emph{A cochlear frequency-position function for several species --- 29 years later}, J. Acoust. Soc. Am. \textbf{87}, 2592 (1990).
\bibitem{Eguiluz2000} V. M. Egu\'iluz, M. Ospeck, Y. Choe, A. J. Hudspeth, and M. O. Magnasco, \emph{Essential nonlinearities in hearing}, Phys. Rev. Lett. \textbf{84}, 5232 (2000).
\bibitem{PalmerRussell1986} A. R. Palmer and I. J. Russell, \emph{Phase-locking in the cochlear nerve of the guinea-pig and its relation to the receptor potential of inner hair-cells}, Hear. Res. \textbf{24}, 1 (1986).
\bibitem{Malnic1999} B. Malnic, J. Hirono, T. Sato, and L. B. Buck, \emph{Combinatorial receptor codes for odors}, Cell \textbf{96}, 713 (1999).
\bibitem{Finger2005} T. E. Finger et al., \emph{ATP signaling is crucial for communication from taste buds to gustatory nerves}, Science \textbf{310}, 1495 (2005).
\bibitem{Huang2005} Y.-J. Huang et al., \emph{Mouse taste buds use serotonin as a neurotransmitter}, J. Neurosci. \textbf{25}, 843 (2005).
\bibitem{JohanssonVallbo1979} R. S. Johansson and A. B. Vallbo, \emph{Tactile sensibility in the human hand: relative and absolute densities of four types of mechanoreceptive units in glabrous skin}, J. Physiol. \textbf{286}, 283 (1979).
\bibitem{McKemy2002} D. D. McKemy, W. M. Neuhausser, and D. Julius, \emph{Identification of a cold receptor reveals a general role for TRP channels in thermosensation}, Nature \textbf{416}, 52 (2002).

\bibitem{Schmolesky1998} M. T. Schmolesky et al., \emph{Signal timing across the macaque visual system}, J. Neurophysiol. \textbf{79}, 3272 (1998).
\bibitem{Lampl2004} I. Lampl, D. Ferster, T. Poggio, and M. Riesenhuber, \emph{Intracellular measurements of spatial integration and the MAX operation in complex cells of the cat primary visual cortex}, J. Neurophysiol. \textbf{92}, 2704 (2004).
\bibitem{KobatakeTanaka1994} E. Kobatake and K. Tanaka, \emph{Neuronal selectivities to complex object features in the ventral visual pathway of the macaque cerebral cortex}, J. Neurophysiol. \textbf{71}, 856 (1994).
\bibitem{Albright1984} T. D. Albright, \emph{Direction and orientation selectivity of neurons in visual area MT of the macaque}, J. Neurophysiol. \textbf{52}, 1106 (1984).
\bibitem{Goodfellow2015} I. J. Goodfellow, J. Shlens, and C. Szegedy, \emph{Explaining and harnessing adversarial examples}, in \emph{Proc. ICLR} (2015), arXiv:1412.6572.
\bibitem{Elsayed2018} G. F. Elsayed et al., \emph{Adversarial examples that fool both computer vision and time-limited humans}, in \emph{Advances in Neural Information Processing Systems 31} (2018), arXiv:1802.08195.
\bibitem{Thompson1982} P. Thompson, \emph{Perceived rate of movement depends on contrast}, Vision Res. \textbf{22}, 377 (1982).
\bibitem{BarlowHill1963} H. B. Barlow and R. M. Hill, \emph{Evidence for a physiological explanation of the waterfall phenomenon and figural after-effects}, Nature \textbf{200}, 1345 (1963).
\bibitem{EaglemanSejnowski2000} D. M. Eagleman and T. J. Sejnowski, \emph{Motion integration and postdiction in visual awareness}, Science \textbf{287}, 2036 (2000).

\bibitem{Feinstein1955mu} B. Feinstein, B. Lindeg{\aa}rd, E. Nyman, and G. Wohlfart, \emph{Morphologic studies of motor units in normal human muscles}, Acta Anat. \textbf{23}, 127 (1955).
\bibitem{MilnerBrown1973rec} H. S. Milner-Brown, R. B. Stein, and R. Yemm, \emph{The orderly recruitment of human motor units during voluntary isometric contractions}, J. Physiol. \textbf{230}, 359 (1973).
\bibitem{Jones2002sdn} K. E. Jones, A. F. de C. Hamilton, and D. M. Wolpert, \emph{Sources of signal-dependent noise during isometric force production}, J. Neurophysiol. \textbf{88}, 1533 (2002).
\bibitem{GomiOsu1998stiff} H. Gomi and R. Osu, \emph{Task-dependent viscoelasticity of human multijoint arm and its spatial characteristics for interaction with environments}, J. Neurosci. \textbf{18}, 8965 (1998).
\bibitem{Jankowska1972Ia} E. Jankowska and W. J. Roberts, \emph{Synaptic actions of single interneurones mediating reciprocal Ia inhibition of motoneurones}, J. Physiol. \textbf{222}, 623 (1972).
\bibitem{Pruszynski2008rapid} J. A. Pruszynski, I. Kurtzer, and S. H. Scott, \emph{Rapid motor responses are appropriately tuned to the metrics of a visuospatial task}, J. Neurophysiol. \textbf{100}, 224 (2008).
\bibitem{SaundersKnill2003vis} J. A. Saunders and D. C. Knill, \emph{Humans use continuous visual feedback from the hand to control fast reaching movements}, Exp. Brain Res. \textbf{152}, 341 (2003).
\bibitem{FlanaganWing1997grip} J. R. Flanagan and A. M. Wing, \emph{The role of internal models in motion planning and control: evidence from grip force adjustments during movements of hand-held loads}, J. Neurosci. \textbf{17}, 1519 (1997).
\bibitem{Matsuoka1985osc} K. Matsuoka, \emph{Sustained oscillations generated by mutually inhibiting neurons with adaptation}, Biol. Cybern. \textbf{52}, 367 (1985).

\bibitem{Treede1995heat} R.-D. Treede, R. A. Meyer, S. N. Raja, and J. N. Campbell, \emph{Evidence for two different heat transduction mechanisms in nociceptive primary afferents innervating monkey skin}, J. Physiol. \textbf{483}, 747 (1995).
\bibitem{Weidner1999cfib} C. Weidner et al., \emph{Functional attributes discriminating mechano-insensitive and mechano-responsive C nociceptors in human skin}, J. Neurosci. \textbf{19}, 10184 (1999).
\bibitem{CampbellLaMotte1983first} J. N. Campbell and R. H. LaMotte, \emph{Latency to detection of first pain}, Brain Res. \textbf{266}, 203 (1983).
\bibitem{LaMotte1982hyper} R. H. LaMotte, J. G. Thalhammer, H. E. Torebj{\"o}rk, and C. J. Robinson, \emph{Peripheral neural mechanisms of cutaneous hyperalgesia following mild injury by heat}, J. Neurosci. \textbf{2}, 765 (1982).
\bibitem{Mancini2014touch} F. Mancini, T. Nash, G. D. Iannetti, and P. Haggard, \emph{Pain relief by touch: a quantitative approach}, Pain \textbf{155}, 635 (2014).
\bibitem{Fields1983rvm} H. L. Fields, J. Bry, I. Hentall, and G. Zorman, \emph{The activity of neurons in the rostral medulla of the rat during withdrawal from noxious heat}, J. Neurosci. \textbf{3}, 2545 (1983).
\bibitem{Crombez1996disrupt} G. Crombez, C. Eccleston, F. Baeyens, and P. Eelen, \emph{The disruptive nature of pain: an experimental investigation}, Behav. Res. Ther. \textbf{34}, 911 (1996).
\bibitem{Schouenborg1995wr} J. Schouenborg, H.-R. Weng, J. Kalliom{\"a}ki, and H. Holmberg, \emph{A survey of spinal dorsal horn neurones encoding the spatial organization of withdrawal reflexes in the rat}, Exp. Brain Res. \textbf{106}, 19 (1995).

\bibitem{Andersen1992cereb} B. B. Andersen, L. Korbo, and B. Pakkenberg, \emph{A quantitative study of the human cerebellum with unbiased stereological techniques}, J. Comp. Neurol. \textbf{326}, 549 (1992).
\bibitem{ShadmehrMussaIvaldi1994} R. Shadmehr and F. A. Mussa-Ivaldi, \emph{Adaptive representation of dynamics during learning of a motor task}, J. Neurosci. \textbf{14}, 3208 (1994).

\bibitem{JoseCollison1970ihr} A. D. Jose and D. Collison, \emph{The normal range and determinants of the intrinsic heart rate in man}, Cardiovasc. Res. \textbf{4}, 160 (1970).
\bibitem{Logothetis1987gbg} D. E. Logothetis, Y. Kurachi, J. Galper, E. J. Neer, and D. E. Clapham, \emph{The $\beta\gamma$ subunits of GTP-binding proteins activate the muscarinic K$^+$ channel in heart}, Nature \textbf{325}, 321 (1987).
\bibitem{KellerWood1984fb} M. E. Keller-Wood and M. F. Dallman, \emph{Corticosteroid inhibition of ACTH secretion}, Endocr. Rev. \textbf{5}, 1 (1984).
\bibitem{Walker2012glc} J. J. Walker et al., \emph{The origin of glucocorticoid hormone oscillations}, PLoS Biol. \textbf{10}, e1001341 (2012).
\bibitem{Wildt1981gnrh} L. Wildt et al., \emph{Frequency and amplitude of gonadotropin-releasing hormone stimulation and gonadotropin secretion in the rhesus monkey}, Endocrinology \textbf{109}, 376 (1981).
\bibitem{Welsh2010scn} D. K. Welsh, J. S. Takahashi, and S. A. Kay, \emph{Suprachiasmatic nucleus: cell autonomy and network properties}, Annu. Rev. Physiol. \textbf{72}, 551 (2010).
\bibitem{Khalsa2003prc} S. B. S. Khalsa, M. E. Jewett, C. Cajochen, and C. A. Czeisler, \emph{A phase response curve to single bright light pulses in human subjects}, J. Physiol. \textbf{549}, 945 (2003).
\bibitem{Sack1992blind} R. L. Sack, A. J. Lewy, M. L. Blood, L. D. Keith, and H. Nakagawa, \emph{Circadian rhythm abnormalities in totally blind people: incidence and clinical significance}, J. Clin. Endocrinol. Metab. \textbf{75}, 127 (1992).

\bibitem{Henze2000extra} D. A. Henze et al., \emph{Intracellular features predicted by extracellular recordings in the hippocampus in vivo}, J. Neurophysiol. \textbf{84}, 390 (2000).
\bibitem{Histed2009stim} M. H. Histed, V. Bonin, and R. C. Reid, \emph{Direct activation of sparse, distributed populations of cortical neurons by electrical microstimulation}, Neuron \textbf{63}, 508 (2009).
\bibitem{Orsborn2014coad} A. L. Orsborn et al., \emph{Closed-loop decoder adaptation shapes neural plasticity for skillful neuroprosthetic control}, Neuron \textbf{82}, 1380 (2014).

\bibitem{HuVervaekeStorm2002} H. Hu, K. Vervaeke, and J. F. Storm, \emph{Two forms of electrical resonance at theta frequencies, generated by M-current, h-current and persistent Na$^+$ current in rat hippocampal pyramidal cells}, J. Physiol. \textbf{545}, 783 (2002).
\bibitem{Mauro1970} A. Mauro, F. Conti, F. Dodge, and R. Schor, \emph{Subthreshold behavior and phenomenological impedance of the squid giant axon}, J. Gen. Physiol. \textbf{55}, 497 (1970).
\bibitem{Aksay2001} E. Aksay et al., \emph{In vivo intracellular recording and perturbation of persistent activity in a neural integrator}, Nat. Neurosci. \textbf{4}, 184 (2001).
\bibitem{Major2004} G. Major et al., \emph{Plasticity and tuning of the time course of analog persistent firing in a neural integrator}, Proc. Natl. Acad. Sci. USA \textbf{101}, 7745 (2004).
\bibitem{Tateno2004} T. Tateno, A. Harsch, and H. P. C. Robinson, \emph{Threshold firing frequency-current relationships of neurons in rat somatosensory cortex: type 1 and type 2 dynamics}, J. Neurophysiol. \textbf{92}, 2283 (2004).
\bibitem{Marder2012} E. Marder, \emph{Neuromodulation of neuronal circuits: back to the future}, Neuron \textbf{76}, 1 (2012).
\bibitem{McGintyHarper1976} D. J. McGinty and R. M. Harper, \emph{Dorsal raphe neurons: depression of firing during sleep in cats}, Brain Res. \textbf{101}, 569 (1976).

\bibitem{Gentet2000} L. J. Gentet, G. J. Stuart, and J. D. Clements, \emph{Direct measurement of specific membrane capacitance in neurons}, Biophys. J. \textbf{79}, 314 (2000).
\bibitem{Borst1995} J. G. G. Borst, F. Helmchen, and B. Sakmann, \emph{Pre- and postsynaptic whole-cell recordings in the medial nucleus of the trapezoid body of the rat}, J. Physiol. \textbf{489}, 825 (1995).
\bibitem{AmirDevor2003} R. Amir and M. Devor, \emph{Electrical excitability of the soma of sensory neurons is required for spike invasion of the soma, but not for through-conduction}, Biophys. J. \textbf{84}, 2181 (2003).
\bibitem{Chadderton2004} P. Chadderton, T. W. Margrie, and M. H\"ausser, \emph{Integration of quanta in cerebellar granule cells during sensory processing}, Nature \textbf{428}, 856 (2004).
\bibitem{LitwinKumar2017} A. Litwin-Kumar et al., \emph{Optimal degrees of synaptic connectivity}, Neuron \textbf{93}, 1153 (2017).
\bibitem{NapperHarvey1988} R. M. A. Napper and R. J. Harvey, \emph{Number of parallel fiber synapses on an individual Purkinje cell in the cerebellum of the rat}, J. Comp. Neurol. \textbf{274}, 168 (1988).
\bibitem{Megias2001} M. Meg\'{\i}as, Z. Emri, T. F. Freund, and A. I. Guly\'as, \emph{Total number and distribution of inhibitory and excitatory synapses on hippocampal CA1 pyramidal cells}, Neuroscience \textbf{102}, 527 (2001).
\bibitem{Polsky2004} A. Polsky, B. W. Mel, and J. Schiller, \emph{Computational subunits in thin dendrites of pyramidal cells}, Nat. Neurosci. \textbf{7}, 621 (2004).
\bibitem{Brunel2004} N. Brunel et al., \emph{Optimal information storage and the distribution of synaptic weights: perceptron versus Purkinje cell}, Neuron \textbf{43}, 745 (2004).

\bibitem{Vyazovskiy2009} V. V. Vyazovskiy et al., \emph{Cortical firing and sleep homeostasis}, Neuron \textbf{63}, 865 (2009).
\bibitem{Daan1984} S. Daan, D. G. Beersma, and A. A. Borb\'ely, \emph{Timing of human sleep: recovery process gated by a circadian pacemaker}, Am. J. Physiol. \textbf{246}, R161 (1984).
\bibitem{Borbely1981} A. A. Borb\'ely et al., \emph{Sleep deprivation: effect on sleep stages and EEG power density in man}, Electroencephalogr. Clin. Neurophysiol. \textbf{51}, 483 (1981).
\bibitem{SteriadeAmzica1993} M. Steriade, F. Amzica, and A. Nu\~nez, \emph{Cholinergic and noradrenergic modulation of the slow ($\approx0.3$ Hz) oscillation in neocortical cells}, J. Neurophysiol. \textbf{70}, 1385 (1993).
\bibitem{vonKrosigk1993} M. von Krosigk, T. Bal, and D. A. McCormick, \emph{Cellular mechanisms of a synchronized oscillation in the thalamus}, Science \textbf{261}, 361 (1993).
\bibitem{LeeWilson2002} A. K. Lee and M. A. Wilson, \emph{Memory of sequential experience in the hippocampus during slow wave sleep}, Neuron \textbf{36}, 1183 (2002).
\bibitem{Euston2007} D. R. Euston, M. Tatsuno, and B. L. McNaughton, \emph{Fast-forward playback of recent memory sequences in prefrontal cortex during sleep}, Science \textbf{318}, 1147 (2007).
\bibitem{Hobson2009} J. A. Hobson, \emph{REM sleep and dreaming: towards a theory of protoconsciousness}, Nat. Rev. Neurosci. \textbf{10}, 803 (2009).

\bibitem{CohenSegal2011} D. Cohen and M. Segal, \emph{Network bursts in hippocampal microcultures are terminated by exhaustion of vesicle pools}, J. Neurophysiol. \textbf{106}, 2314 (2011).
\bibitem{Tsodyks2000} M. Tsodyks, A. Uziel, and H. Markram, \emph{Synchrony generation in recurrent networks with frequency-dependent synapses}, J. Neurosci. \textbf{20}, RC50 (2000).
\bibitem{Benson1994} D. L. Benson, F. H. Watkins, O. Steward, and G. Banker, \emph{Characterization of GABAergic neurons in hippocampal cell cultures}, J. Neurocytol. \textbf{23}, 279 (1994).

\input{supplement/bib_22}
\end{thebibliography}


\end{document}
